
\documentclass[11pt]{article}
\usepackage[a4paper,margin=25mm]{geometry}
\usepackage[T1]{fontenc}
\usepackage[utf8]{inputenc}
\usepackage[ngerman]{babel}
\usepackage{lmodern}
\usepackage{amsmath,amssymb,amsthm,mathtools,array,booktabs,graphicx,xcolor,enumitem,pdflscape}
\usepackage{microtype}
\usepackage[colorlinks=true,urlcolor=black,linkcolor=black]{hyperref}
\setlength{\parindent}{1.2em}
\setlength{\parskip}{0.15em}
\emergencystretch=3em
\allowdisplaybreaks
\hbadness=10000
\newcommand{\grp}{\mathfrak{H}}
\newcommand{\Hgrp}{\mathfrak{H}}
\newcommand{\Ggrp}{\mathfrak{G}}
\newcommand{\Egrp}{\mathfrak{E}}
\newcommand{\zchar}{\mathfrak{z}}
\newcommand{\Hgr}{\mathfrak{H}}
\newcommand{\Ggr}{\mathfrak{G}}
\newcommand{\Dgr}{\mathfrak{D}}
\newcommand{\Fgr}{\mathfrak{F}}
\newcommand{\FrobP}{\mathfrak{P}}
\newcommand{\Char}{\chi}
\newcommand{\eps}{\varepsilon}
\newcommand{\ten}{\mathrm{ten}}
\newcommand{\te}{\mathrm{te}}
\newcommand{\Class}[1]{#1^{\mathrm{te}}}
\newcommand{\class}[1]{\textsuperscript{#1}}
\newcommand{\Cl}[1]{#1^{\mathrm{ten}}}
\DeclareMathOperator{\rang}{Rang}
\newcommand{\resizeMath}[1]{\resizebox{\textwidth}{!}{$#1$}}
\newcommand{\half}{\tfrac{1}{2}}
\newcommand{\frobtitle}[1]{\begin{center}{\Large #1}\\[0.75em]Von G. \textsc{Frobenius}.\end{center}\vspace{0.5em}}
\newcommand{\sign}{\operatorname{sign}}
\newcommand{\coeff}[1]{\left[#1\right]}
\newcommand{\CharTwo}[4]{\left(\begin{matrix}#1&\cdots&#2\\#3&\cdots&#4\end{matrix}\right)}
\newcommand{\ch}[2]{\left(\begin{smallmatrix}#1\\#2\end{smallmatrix}\right)}
\newcommand{\vtheta}{\vartheta}
\newcommand{\fnum}[1]{\textup{(\ensuremath{#1}.)}\quad}
\newcommand{\tr}{\operatorname{tr}}
\newcommand{\rank}{\operatorname{Rang}}
\newcommand{\grS}{\mathfrak{S}}
\newcommand{\modeta}{\pmod{\eta}}
\newcommand{\frobtag}[2]{\[\begin{array}{@{}r@{\quad}l@{}}\textup{(#1)} & \displaystyle #2\end{array}\]}
\newcommand{\frobtagc}[2]{\[\begin{array}{@{}r@{\quad}c@{}}\textup{(#1)} & \displaystyle #2\end{array}\]}
\begin{document}
\begin{center}{\huge Frobenius: ausgewählte Arbeiten zur Gruppentheorie und hyperkomplexen Größen}\end{center}
\bigskip

% ===== BEGIN WORK 053: 053_ber_Gruppencharaktere =====
\begin{center}
{\Large \textbf{Über Gruppencharaktere.}}\\[0.8em]
{\large Von G. \textsc{Frobenius}.}\\[0.6em]
\rule{6em}{0.4pt}
\end{center}

Bei dem Beweise des Satzes, dass jede lineare Function einer Variabeln unendlich viele Primzahlen darstellt, wenn ihre Coefficienten theilerfremde ganze Zahlen sind, benutzte \textsc{Dirichlet} zum ersten Male gewisse Systeme von Einheitswurzeln, die auch in der nahe verwandten Frage nach der Anzahl der Idealclassen in einem Kreiskörper auftreten (vergl. die Bemerkung von \textsc{Dedekind} in \textsc{Dirichlet}'s Vorlesungen über Zahlentheorie, 4. Aufl. S. 625), sowie bei der Verallgemeinerung jenes Satzes auf quadratische Formen und in den Untersuchungen über deren Eintheilung in Geschlechter. Die charakteristische Eigenschaft dieser Ausdrücke besteht nach \textsc{Dedekind} darin, dass sie von einer variabeln positiven ganzen Zahl $n$ abhängige Grössen $\chi(n)$ sind, die nur eine endliche Anzahl von Werthen haben und der Bedingung
\[
\chi(m)\chi(n)=\chi(mn)
\]
genügen. Wie er in rein abstracter Form ausführt, lassen sich den Elementen $A,B,C,\ldots$ jeder endlichen Gruppe $\grp$ vertauschbarer Elemente (\textsc{Abel}'schen Gruppe) solche Einheitswurzeln $\chi(A),\chi(B),\chi(C),\ldots$ zuordnen, welche die Gleichungen
\[
\chi(A)\chi(B)=\chi(AB)
\]
befriedigen, und die er nach dem Vorgange von \textsc{Gauss} die \emph{Charaktere der Gruppe} nannte.

Unter einem \emph{Charakter} einer quadratischen Form versteht \textsc{Gauss}, \emph{Disqu. arithm.} Art. 230 eine Relation der durch die Form darstellbaren Zahlen zu den in ihrer Determinante aufgehenden ungeraden Primzahlen $p$ (oder 4 oder 8). Er drückt jene Beziehung durch die Zeichen $Rp$ und $Np$ aus. Diese Symbole ersetzt \textsc{Dirichlet}, \emph{Recherches sur diverses applications de l'analyse infinitésimale à la théorie des nombres}, \S~3 (\textsc{Crelle}'s Journal Bd. 19) durch das \textsc{Legendre}'sche (und \textsc{Jacobi}'sche) Zeichen $\left(\frac{m}{p}\right)$, welches (nächst der Resolvente von \textsc{Lagrange}) wohl das älteste Beispiel der Anwendung von Charakteren commutativer Gruppen


darbietet. Der Vorzug dieser Umwandlung besteht darin, dass die Geschlechtscharaktere von \textsc{Gauss} nur Beziehungen, die von \textsc{Dirichlet} aber Zahlen sind, mit denen man rechnen kann. So wird durch die Multiplication dieser charakteristischen Zahlen die der Composition der Geschlechter entsprechende Composition der Charaktere (Art. 246 bis 248) ersetzt\footnote{Die in diesem Absatze enthaltenen Bemerkungen entnehme ich einem Briefe \textsc{Dedekind}'s vom 8. Juli 1896.}.

Die Anzahl $h$ der Charaktere einer \textsc{Abel}'schen Gruppe $\grp$ ist der Ordnung der Gruppe gleich. Man kann die $h$ Charaktere erhalten, indem man die Elemente von $\grp$ durch eine Basis unabhängiger Elemente darstellt, und den Basiselementen beliebige Einheitswurzeln zuordnet, deren Grad ihrer Ordnung gleich ist. Sie lassen sich (in mehrfacher Art) den Elementen der Gruppe zuordnen und können demnach mit $\chi_B(A)$ bezeichnet werden. Da das Product zweier Charaktere wieder ein Charakter ist, so bilden sie eine Gruppe, und diese ist mit $\grp$ isomorph. Ihre Beziehungen zu den Untergruppen von $\grp$ sind am ausführlichsten von \textsc{Weber} erörtert (\emph{Theorie der \textsc{Abel}'schen Zahlkörper}, I. \S~3, IV. \S~2 und 3, Acta Math. Bd. 8 und 9).

Im April dieses Jahres theilte mir \textsc{Dedekind} eine Aufgabe mit, auf die er im Jahre 1880 gekommen war, und die, weil sie sowohl der Gruppentheorie wie der Determinantentheorie angehöre, mich seiner Meinung nach wohl interessiren dürfte, während ihn selbst ein näheres Eingehen darauf zu weit von seinen arithmetischen Untersuchungen abziehen würde. Ihre Lösung, die ich nächstens mittheilen zu können hoffe, brachte mich auf eine Verallgemeinerung des Begriffs der Charaktere auf beliebige endliche Gruppen. Diesen Begriff will ich hier entwickeln in der Meinung, dass durch seine Einführung die Gruppentheorie eine wesentliche Förderung und Bereicherung erfahren dürfte. Ein besonderes Interesse gewinnt die Theorie der Charaktere noch durch ihre merkwürdigen Beziehungen zu der Theorie der aus mehreren Haupteinheiten gebildeten complexen Grössen.

\section*{\S~1.}

Zwei Elemente $A$ und $B$ einer endlichen Gruppe $\grp$ heissen \emph{conjugirt} (in Bezug auf $\grp$), wenn es in $\grp$ ein Element $T$ giebt, das der Bedingung $T^{-1}AT=B$ genügt. Sind zwei Elemente einem dritten conjugirt, so sind sie es auch unter einander. Daher kann man die $h$ Elemente von $\grp$ in \emph{Classen conjugirter Elemente} eintheilen, eine Eintheilung, von der ich mehrfach, besonders in meinem neuen Be-


weise des \textsc{Sylow}'schen Satzes, \textsc{Crelle}'s Journal Bd. 100 vortheilhaft Gebrauch gemacht habe. Das Hauptelement $E$ bildet für sich eine Classe, die \emph{Hauptclasse}. Sie werde mit $(0)$, die übrigen mit $(1),(2),\ldots,(k-1)$ bezeichnet, wenn $k$ die Anzahl der Classen ist. Ist $A$ irgend ein Element der $\alpha^{\mathrm{ten}}$ Classe, so bilden die mit $A$ vertauschbaren Elemente von $\grp$ eine in $\grp$ enthaltene Gruppe. Ist $\frac{h}{h_\alpha}$ ihre Ordnung, so ist $h_\alpha$ die Anzahl der verschiedenen Elemente der $\alpha^{\mathrm{ten}}$ Classe, also $h_0=1$. Da jedes Element von $\grp$ einer und nur einer dieser $k$ Classen angehört, so ist
\begin{equation*}\tag{1.}
\sum_\alpha h_\alpha=h.
\end{equation*}

Die Classenanzahl $k$ kann man auch so erhalten: Durchläuft jedes der beiden veränderlichen Elemente $R$ und $S$ unabhängig von dem anderen die $h$ Elemente von $\grp$, so zähle man ab, wie oft $SR=RS$ ist. Setzt man für $R$ ein bestimmtes Element der $\alpha^{\mathrm{ten}}$ Classe, so giebt es $\frac{h}{h_\alpha}$ mit $R$ vertauschbare Elemente $S$. Setzt man also für $R$ der Reihe nach jedes der $h_\alpha$ Elemente der $\alpha^{\mathrm{ten}}$ Classe, so erhält man $\frac{h}{h_\alpha}h_\alpha$ Lösungen für jene Gleichung. Da diese Zahl für jede der $k$ Classen dieselbe ist, so ist $hk$ die Anzahl der Lösungen der Gleichung $SR=RS$.

Durchläuft $A$ die $h_\alpha$ Elemente der $\alpha^{\mathrm{ten}}$ Classe, so durchläuft auch $A^{-1}$ die sämmtlichen Elemente einer Classe. Sie möge die \emph{inverse Classe} von $(\alpha)$ heissen und mit $(\alpha')$ bezeichnet werden. Inverse Classen enthalten gleich viele Elemente,
\begin{equation*}\tag{2.}
h_{\alpha'}=h_\alpha.
\end{equation*}
Ist, wie z. B. bei der Hauptclasse, $(\alpha')=(\alpha)$, so wird die $\alpha^{\mathrm{te}}$ Classe eine \emph{zweiseitige} genannt.

Seien $(\alpha),(\beta),(\gamma)$ irgend drei verschiedene oder gleiche Classen. Durchläuft $A$ die $h_\alpha$ verschiedenen Elemente der $\alpha^{\mathrm{ten}}$ Classe, $B$ die $h_\beta$ Elemente der $\beta^{\mathrm{ten}}$ und $C$ die $h_\gamma$ Elemente der $\gamma^{\mathrm{ten}}$, so soll die Zahl $h_{\alpha\beta\gamma}$ (die auch Null sein kann) angeben, wie viele der $h_\alpha h_\beta h_\gamma$ Elemente $ABC$ gleich dem Hauptelemente sind, also der Gleichung
\begin{equation*}\tag{3.}
ABC=E
\end{equation*}
genügen. Da dann $AB=C^{-1}$ ist, so giebt $h_{\alpha\beta\gamma}$ auch an, wie viele der $h_\alpha h_\beta$ Elemente $AB$ der Classe $(\gamma')$ angehören. Die Gleichung (3.) ist identisch mit $BCA=E$ und $CAB=E$. Daher sind auch $h_{\alpha\beta\gamma}$ der $h_\beta h_\gamma$ Producte $BC$ in $(\alpha')$ und $h_{\alpha\beta\gamma}$ der $h_\gamma h_\alpha$ Producte $CA$ in $(\beta')$ enthalten. Mithin ist $h_{\alpha\beta\gamma}$ nicht grösser als die kleinste der drei Zahlen $h_\beta h_\gamma$, $h_\gamma h_\alpha$ und $h_\alpha h_\beta$.


Nun sind aber die beiden Elemente $AB$ und
\begin{equation*}\tag{4.}
BA=A^{-1}(AB)A=B(AB)B^{-1}
\end{equation*}
conjugirt. Sie gehören daher der Gruppe $(\gamma')$ entweder beide an, oder beide nicht. Mithin ist $h_{\beta\alpha\gamma}=h_{\alpha\beta\gamma}$. In Verbindung mit den obigen Bemerkungen folgt daraus, dass die Zahl $h_{\alpha\beta\gamma}$ bei jeder Vertauschung der drei Indices ungeändert bleibt. Da aus der Gleichung (3.) auch $C^{-1}B^{-1}A^{-1}=E$ folgt, so ist
\begin{equation*}\tag{5.}
h_{\alpha'\beta'\gamma'}=h_{\alpha\beta\gamma}.
\end{equation*}

Setzt man in der Gleichung (3.) für $A$ ein bestimmtes Element der $\alpha^{\mathrm{ten}}$ Classe, während $B$ und $C$ ebenso veränderlich bleiben wie oben, so möge sie $m$ Lösungen haben. Ist $A'=T^{-1}AT$ irgend ein anderes bestimmtes Element der $\alpha^{\mathrm{ten}}$ Classe, und setzt man $B'=T^{-1}BT$ und $C'=T^{-1}CT$, so ist auch $A'B'C'=E$. Nun durchläuft $B'$ gleichzeitig mit $B$ die $h_\beta$ Elemente der $\beta^{\mathrm{ten}}$ Classe, und $C'$ gleichzeitig mit $C$ die $h_\gamma$ Elemente der $\gamma^{\mathrm{ten}}$ Classe. Daher hat auch diese Gleichung $m$ Lösungen. Setzt man also in (3.) für $A$ der Reihe nach die $h_\alpha$ Elemente der $\alpha^{\mathrm{ten}}$ Classe, so hat sie im Ganzen $h_\alpha m=h_{\alpha\beta\gamma}$ Lösungen. Ebenso hat jene Gleichung $\frac{h_{\alpha\beta\gamma}}{h_\beta}$ Lösungen, wenn $B$ ein bestimmtes Element der $\beta^{\mathrm{ten}}$ Classe ist, während $A$ und $C$ veränderliche Elemente der $\alpha^{\mathrm{ten}}$ und der $\gamma^{\mathrm{ten}}$ Classe sind. Mithin ist $h_{\alpha\beta\gamma}$ durch jede der drei Zahlen $h_\alpha,h_\beta,h_\gamma$ theilbar, also auch durch ihr kleinstes gemeinschaftliches Vielfaches.

Setzt man für $A$ der Reihe nach die $h_\alpha$ Elemente der $\alpha^{\mathrm{ten}}$ Classe und für $B$ die $h_\beta$ Elemente der $\beta^{\mathrm{ten}}$ Classe, so erhält man $h_\alpha h_\beta$ Elemente $AB$, die nicht verschieden zu sein brauchen. Jedes derselben gehört entweder der Classe $(0)$ oder der Classe $(1)\ldots$ oder der Classe $(k-1)$ an. Daher ist
\begin{equation*}\tag{6.}
\sum_\gamma h_{\alpha\beta\gamma}=h_\alpha h_\beta.
\end{equation*}

Die für drei Classen angestellte Betrachtung lässt sich in derselben Art für beliebig viele Classen durchführen. Seien etwa $(\alpha),(\beta),(\gamma),(\delta)$ vier verschiedene oder gleiche Classen, und sei $h_{\alpha\beta\gamma\delta}$ die Anzahl der Lösungen der Gleichung
\begin{equation*}\tag{7.}
ABCD=E,
\end{equation*}
falls $A$ (bez. $B,C,D$) die $h_\alpha$ (bez. $h_\beta,h_\gamma,h_\delta$) Elemente der Classe $(\alpha)$ (bez. $(\beta),(\gamma),(\delta)$) durchläuft. Dann ist auch $BCDA=CDAB=DABC=E$, so dass $h_{\alpha\beta\gamma\delta}$ bei cyklischer Vertauschung der Indices ungeändert bleibt. Ferner ist $BCD=A^{-1}$, $CDA=B^{-1}$, u. s. w. Mithin giebt $h_{\alpha\beta\gamma\delta}$ an, wie viele der $h_\beta h_\gamma h_\delta$ Elemente $BCD$ der Classe $(\alpha')$ angehören, oder wie viele der $h_\gamma h_\delta h_\alpha$ Elemente $CDA$ der Classe $(\beta')$ angehören. Folglich ist $h_{\alpha\beta\gamma\delta}$ nicht grösser als die kleinste der vier Zahlen $h_\beta h_\gamma h_\delta$, $h_\gamma h_\delta h_\alpha$,


$h_\delta h_\alpha h_\beta$, $h_\alpha h_\beta h_\gamma$. Ist $B$ ein festes Element der $\beta^{\mathrm{ten}}$ Classe, während $A,C,D$ die $\alpha^{\mathrm{te}},\gamma^{\mathrm{te}},\delta^{\mathrm{te}}$ Classe durchlaufen, so hat die Gleichung (7.) nur $\frac{h_{\alpha\beta\gamma\delta}}{h_\beta}$ Lösungen. Demnach ist $h_{\alpha\beta\gamma\delta}$ durch jede der vier Zahlen $h_\alpha,h_\beta,h_\gamma,h_\delta$ theilbar. Endlich ist
\begin{equation*}\tag{8.}
\sum_\delta h_{\alpha\beta\gamma\delta}=h_\alpha h_\beta h_\gamma.
\end{equation*}

Ist $A,B,C,D$ eine Lösung der Gleichung (7.), und setzt man
\[
B^{-1}AB=A',\qquad BA'B^{-1}=A,
\]
so ist $BA'CD=E$. Da $A'$ ebenso wie $A$ der $\alpha^{\mathrm{ten}}$ Classe angehört, so kann man auf diese Weise die Lösungen dieser Gleichung und der Gleichung (7.) einander paarweise eindeutig zuordnen. Mithin ist $h_{\beta\alpha\gamma\delta}=h_{\alpha\beta\gamma\delta}$. Auf diese Weise erkennt man, dass $h_{\alpha\beta\gamma\delta}$ bei jeder Vertauschung der vier Indices ungeändert bleibt. Daher ist auch
\begin{equation*}\tag{9.}
h_{\alpha'\beta'\gamma'\delta'}=h_{\alpha\beta\gamma\delta}.
\end{equation*}

Dieselbe Bedeutung und die analogen Eigenschaften hat für $n$ Classen $(\alpha),(\beta),(\gamma),\ldots,(\nu)$ die Zahl $h_{\alpha\beta\gamma\ldots\nu}$. Für den Fall $n=1$ soll dies Zeichen nicht benutzt werden, wohl aber noch für $n=2$. Es ist also $h_{\alpha\beta}$ die Anzahl der Lösungen der Gleichung $AB=E$ oder $B=A^{-1}$. Sind demnach $(\alpha)$ und $(\beta)$ nicht inverse Classen, so ist $h_{\alpha\beta}=0$. Ist aber $(\beta)=(\alpha')$, so ist
\begin{equation*}\tag{10.}
h_{\alpha\alpha'}=h_\alpha=h_{\alpha'}.
\end{equation*}
Ist $\nu=0$, so ist $h_{\alpha\beta\gamma\ldots\mu0}$ die Anzahl der Lösungen der Gleichung $ABC\ldots ME=E$. Daher ist
\begin{equation*}\tag{11.}
h_{\alpha\beta\gamma\ldots\mu0}=h_{\alpha\beta\gamma\ldots\mu}.
\end{equation*}
Speciell ist
\begin{equation*}\tag{12.}
h_{\alpha\beta0}=h_{\alpha\beta},
\end{equation*}
also Null, ausser wenn $(\beta)=(\alpha')$ ist, dann aber gleich $h_\alpha$.

\section*{\S~2.}

Die wichtigste Eigenschaft der Zahlen $h_{\alpha\beta\gamma}$ erhält man, indem man die Gleichung $ABCD=E$ in die beiden Gleichungen
\[
AB=L^{-1},\qquad CD=L
\]
zerlegt. Setzt man für $L$ ein bestimmtes Element der $\lambda^{\mathrm{ten}}$ Classe, so hat die erste Gleichung $\frac{h_{\alpha\beta\lambda}}{h_\lambda}$, die zweite $\frac{h_{\lambda'\gamma\delta}}{h_\lambda}$ Lösungen. So ergeben sich $\frac{h_{\alpha\beta\lambda}h_{\lambda'\gamma\delta}}{h_\lambda^2}=m_\lambda$ Lösungen der Gleichung (7.) \S~1. Setzt man für $L$ der Reihe nach jedes der $h_\lambda$ Elemente der $\lambda^{\mathrm{ten}}$ Classe, so erhält man




$h_\lambda m_\lambda$ Lösungen, und setzt man endlich für $(\lambda)$ jede der $k$ Classen, so erkennt man, dass sie im Ganzen
\begin{equation*}\tag{1.}
h_{\alpha\beta\gamma\delta}=\sum_\lambda \frac{1}{h_\lambda}h_{\alpha\beta\lambda}h_{\lambda'\gamma\delta}
\end{equation*}
Lösungen hat. In derselben Weise ergiebt sich
\begin{equation*}\tag{2.}
h_{\alpha\beta\ldots\xi\eta\ldots\sigma\tau}=\sum_\lambda \frac{1}{h_\lambda}h_{\alpha\beta\ldots\xi\lambda}h_{\lambda'\eta\ldots\sigma\tau}.
\end{equation*}
Berücksichtigt man noch die Gleichungen (10.) und (12.), \S~1, so kann man daher die Zahlen $h_{\alpha\beta\gamma\ldots}$ alle aus den Zahlen $h_{\alpha\beta\gamma}$ zusammensetzen.

Die Summe auf der rechten Seite der Gleichung (1.) bleibt folglich ungeändert, nicht nur wenn man $\alpha$ mit $\beta$ vertauscht oder $\gamma$ mit $\delta$, sondern auch wenn man $\beta$ mit $\delta$ vertauscht oder irgend eine Permutation unter den vier Zahlen $\alpha,\beta,\gamma,\delta$ ausführt. Dieselben Schlüsse benutzt \textsc{Gauss}, \emph{Theoria residuorum biquadraticorum}, Comm. prima, \S~17 (Ges. Werke Bd. II, S. 81). Setzt man also für einen Augenblick
\[
\frac{1}{h_\alpha}h_{\alpha\beta\gamma}=a_{\alpha\beta\gamma},
\]
so ist
\begin{equation*}\tag{3.}
a_{\alpha\beta\gamma}=a_{\alpha\gamma\beta},\qquad
\sum_\lambda a_{\alpha\lambda\gamma}a_{\lambda\beta\delta}=\sum_\lambda a_{\alpha\lambda\delta}a_{\lambda\beta\gamma}.
\end{equation*}
Daher kann man auf diese Grössen die Sätze anwenden, die \textsc{Weierstrass} und \textsc{Dedekind} in ihren Arbeiten \emph{Zur Theorie der aus $n$ Haupteinheiten gebildeten complexen Grössen}, Göttinger Nachrichten 1884 und 1885 entwickelt haben, und die ich, da sie die Grundlage dieser Untersuchung bilden, in meiner Arbeit \emph{Über vertauschbare Matrizen} (S. 601 dieses Bandes) von Neuem hergeleitet und verallgemeinert habe: Ist die aus den $k^2$ Grössen
\[
p_{\alpha\beta}=\sum_{\varkappa,\lambda}a_{\varkappa\lambda\alpha}a_{\lambda\varkappa\beta}=\sum_{\varkappa,\lambda}a_{\varkappa\varkappa\lambda}a_{\lambda\alpha\beta}
\]
gebildete Determinante $k^{\text{ten}}$ Grades von Null verschieden, so haben die Gleichungen
\[
r_\beta r_\gamma=\sum_\alpha a_{\alpha\beta\gamma}r_\alpha
\]
genau $k$ verschiedene Lösungen $r_\alpha=r_\alpha^{(\varkappa)}$, und die aus diesen Lösungen gebildete Determinante $k^{\text{ten}}$ Grades ist von Null verschieden. Sind $x_0,x_1,\ldots,x_{k-1}$ Variable, und setzt man $a_{\alpha\beta}=\sum_\gamma a_{\alpha\beta\gamma}x_\gamma$, so ist die Determinante $k^{\text{ten}}$ Grades
\[
\left|a_{\alpha\beta}-r\epsilon_{\alpha\beta}\right|=\prod_\varkappa\left(r_0^{(\varkappa)}x_0+\cdots+r_{k-1}^{(\varkappa)}x_{k-1}-r\right).
\]
Ist $(s_\alpha^{(\varkappa)})$ das complementäre System zu $(r_\alpha^{(\varkappa)})$, so sind die Verhältnisse der $k$ Grössen
\[
s_0=s_0^{(\varkappa)},\qquad s_1=s_1^{(\varkappa)},\ldots\qquad s_{k-1}=s_{k-1}^{(\varkappa)}
\]


vollständig bestimmt durch die linearen Gleichungen
\[
s_\alpha r=\sum_\beta a_{\alpha\beta}s_\beta,
\]
falls man darin $r=r_0^{(\varkappa)}x_0+\cdots+r_{k-1}^{(\varkappa)}x_{k-1}$ setzt (S. 613 und 614).

Im vorliegenden Falle ist
\[
p_{\alpha\beta}=\sum_{\varkappa,\lambda}\frac{h_{\varkappa\lambda\alpha}h_{\lambda\varkappa\beta}}{h_\varkappa h_\lambda}.
\]
Da $(\varkappa')$ zugleich mit $(\varkappa)$ alle $k$ Classen durchläuft, so kann man in dieser Summe auch $(\varkappa)$ durch $(\varkappa')$ ersetzen, und erhält dann nach (5.), \S~1
\begin{equation*}\tag{4.}
p_{\alpha\beta'}=\sum_{\varkappa,\lambda}\frac{h_{\varkappa\lambda\alpha}h_{\varkappa\lambda\beta}}{h_\varkappa h_\lambda},
\end{equation*}
also
\begin{equation*}\tag{5.}
p_{\alpha\beta'}=p_{\beta\alpha'}=p_{\alpha'\beta}=p_{\beta'\alpha}.
\end{equation*}
Die Determinante $k^{\text{ten}}$ Grades $|p_{\alpha\beta'}|$ wird aus $|p_{\alpha\beta}|$ durch eine Vertauschung unter den Spalten erhalten, unterscheidet sich also von ihr nur etwa im Vorzeichen. Betrachtet man nun die $k$ Grössen
\[
\frac{h_{\varkappa\lambda0}}{\sqrt{h_\varkappa h_\lambda}},\qquad
\frac{h_{\varkappa\lambda1}}{\sqrt{h_\varkappa h_\lambda}},\ldots\qquad
\frac{h_{\varkappa\lambda,k-1}}{\sqrt{h_\varkappa h_\lambda}}.
\]
Setzt man für $\varkappa$ und für $\lambda$ die Werthe $0,1,\ldots,k-1$, so erhält man ein System von $k^3$ Grössen, die in $k^2$ Zeilen und $k$ Spalten geordnet sind. Aus diesem System kann man, indem man irgend $k$ Zeilen auswählt, eine Determinante $k^{\text{ten}}$ Grades bilden, also auf $\binom{k^2}{k}$ Arten. Die Summe der Quadrate aller dieser Determinanten ist auf Grund des allgemeinen Multiplicationstheorems der Determinantentheorie nach Gleichung (4.) gleich $|p_{\alpha\beta'}|$. Setzt man $\varkappa=0$ und $\lambda'=0,1,\ldots,k-1$, so erhält man ein bestimmtes System von $k$ Zeilen. In der aus ihren Elementen gebildeten Determinante sind die Elemente der Diagonale $\frac{h_{0\lambda'\lambda}}{\sqrt{h_\lambda}}=\sqrt{h_\lambda}$ von Null verschieden, die übrigen Elemente Null. Folglich ist die Determinante $k^{\text{ten}}$ Grades $|p_{\alpha\beta}|$ von Null verschieden.

Sei $f$ ein vorläufig unbestimmt gelassener Proportionalitätsfactor und
\[
r_\alpha=\frac{h_\alpha\chi_\alpha}{f}.
\]
Dann haben die Gleichungen
\begin{equation*}\tag{6.}
h_\beta h_\gamma\chi_\beta\chi_\gamma=f\sum_\alpha h_{\alpha'\beta\gamma}\chi_\alpha
\end{equation*}
$k$ verschiedene Systeme von Lösungen
\begin{equation*}\tag{7.}
\chi_\alpha=\chi_\alpha^{(\varkappa)},\qquad f=f^{(\varkappa)}\qquad(\varkappa=0,1,\ldots,k-1),
\end{equation*}


und die Determinante $k^{\text{ten}}$ Grades
\begin{equation*}\tag{8.}
\left|\chi_\alpha^{(\varkappa)}\right|
\end{equation*}
ist von Null verschieden. Ich werde später über die Wahl des Factors $f$, d. h. der $k$ Factoren $f^{(\varkappa)}$, eine bestimmte Verfügung treffen und dann die Grössen $\chi_\alpha^{(\varkappa)}$ die $k$ \emph{Charaktere der Gruppe $\grp$} nennen. Ein Charakter $\chi$ ist ein System von $k$ Zahlen $\chi_0,\chi_1,\ldots,\chi_{k-1}$, die den $k$ Classen $(0),(1),\ldots,(k-1)$ entsprechen. Der $\varkappa^{\text{te}}$ Charakter $\chi^{(\varkappa)}$ ist das System der $k$ Zahlen
\[
\chi_0=\chi_0^{(\varkappa)},\qquad \chi_1=\chi_1^{(\varkappa)},\ldots\qquad \chi_{k-1}=\chi_{k-1}^{(\varkappa)}\qquad(\varkappa=0,1,\ldots,k-1).
\]
Die Bedeutung der Gleichungen (6.) ergiebt sich aus folgender Betrachtung: Sind $x_0,x_1,\ldots,x_{k-1}$ und $y_0,y_1,\ldots,y_{k-1}$ $2k$ unabhängige Variable, und sind
\begin{equation*}\tag{9.}
h_\gamma z_\gamma=\sum_{\alpha,\beta}h_{\alpha\beta\gamma}x_\alpha y_\beta
\end{equation*}
$k$ bilineare Functionen derselben, so ist
\[
f\sum_\gamma h_\gamma\chi_\gamma z_\gamma
=f\sum_{\alpha,\beta,\gamma}h_{\alpha\beta\gamma}\chi_\gamma x_\alpha y_\beta
=\sum_{\alpha,\beta}h_\alpha\chi_\alpha x_\alpha\,h_\beta\chi_\beta y_\beta,
\]
also
\begin{equation*}\tag{10.}
f\left(\sum_\gamma h_\gamma\chi_\gamma z_\gamma\right)=
\left(\sum_\alpha h_\alpha\chi_\alpha x_\alpha\right)
\left(\sum_\beta h_\beta\chi_\beta y_\beta\right).
\end{equation*}
Ich definire nun eine dem Charakter $\chi^{(\varkappa)}$ entsprechende lineare Function $\xi^{(\varkappa)}$ von $x_0,x_1,\ldots,x_{k-1}$ durch die Gleichung
\begin{equation*}\tag{11.}
\sum_\alpha h_\alpha\chi_\alpha^{(\varkappa)}x_\alpha=f^{(\varkappa)}\xi^{(\varkappa)}\qquad(\varkappa=0,1,\ldots,k-1).
\end{equation*}
Dann ist die Determinante $k^{\text{ten}}$ Grades
\begin{equation*}\tag{12.}
\left|\sum_\gamma h_{\alpha\beta'\gamma}x_\gamma-h_{\alpha\beta'}r\right|=
\prod_\varkappa h_\varkappa\left(\xi^{(\varkappa)}-r\right).
\end{equation*}
Setzt man $x_1=\cdots=x_{k-1}=0$ und $x_0=1$, so ist daher
\[
(1-r)^k=\prod_\varkappa\left(\frac{\chi_0^{(\varkappa)}}{f^{(\varkappa)}}-r\right)
\]
und mithin
\begin{equation*}\tag{13.}
\chi_0^{(\varkappa)}=f^{(\varkappa)},\qquad \chi_0=f,
\end{equation*}
da ich bei Gleichungen, die für jeden Werth des oberen Index $\varkappa$ gelten, diesen auch weglassen werde. (Vergl. die Erörterungen von \textsc{Dedekind} über mehrwerthige Systeme von Grössen, a. a. O. S. 144.) Demnach ist die Gleichung (12.) um nichts allgemeiner als die Gleichung
\begin{equation*}\tag{14.}
\left|\sum_\gamma h_{\alpha\beta'\gamma}x_\gamma\right|=\prod_\varkappa h_\varkappa\xi^{(\varkappa)},
\end{equation*}
weil sie aus dieser hervorgeht, indem man $x_0$ durch $x_0-r$ ersetzt. Schreibt man jene in der Form
\[
\left|\sum_\gamma \frac{h_{\alpha\beta'\gamma}}{h_\alpha}x_\gamma-\frac{h_{\alpha\beta'}}{h_\alpha}r\right|=\prod_\varkappa(\xi^{(\varkappa)}-r),
\]


so ist $\frac{h_{\alpha\beta'}}{h_\alpha}$ in den Elementen der Diagonale 1, in den übrigen 0. Ferner sind die Grössen $\frac{h_{\alpha\beta'\gamma}}{h_\alpha}$ ganze Zahlen. Setzt man also alle Variabeln $x_\gamma=0$ ausser einer, so zeigt sie, dass die Grössen
\begin{equation*}\tag{15.}
\frac{h_\alpha\chi_\alpha}{f}
\end{equation*}
ganze algebraische Zahlen sind.

\section*{\S~3.}

Die Elemente jeder Zeile des zu $(r_\alpha^{(\varkappa)})$ complementären Systems $(s_\alpha^{(\varkappa)})$ sind bis auf einen gemeinsamen Factor durch die Gleichungen
\[
\sum_\beta \frac{h_{\alpha'\beta\gamma}}{h_\alpha}s_\beta=\frac{h_\gamma\chi_\gamma}{f}s_\alpha
\]
bestimmt. Infolge der Symmetrieeigenschaften der Grössen $h_{\alpha\beta\gamma}$ werden dieselben nach (6.), \S~2 durch die Werthe $s_\beta=\chi_{\beta'}$ befriedigt. Daher ist $s_\beta=\frac{e}{h}\chi_{\beta'}$, wo $e$ ein neuer Proportionalitätsfactor ist. Demnach sind die beiden Systeme
\begin{equation*}\tag{1.}
\left(\frac{h_\alpha\chi_\alpha^{(\varkappa)}}{f^{(\varkappa)}}\right),\qquad
\left(\frac{e^{(\varkappa)}\chi_{\alpha'}^{(\varkappa)}}{h}\right)
\end{equation*}
complementäre. Es bestehen also die Gleichungen
\begin{equation*}\tag{2.}
\sum_\alpha h_\alpha\chi_\alpha^{(\varkappa)}\chi_{\alpha'}^{(\lambda)}=0,
\end{equation*}
falls $\varkappa$ und $\lambda$ verschieden sind, aber
\begin{equation*}\tag{3.}
\sum_\alpha h_\alpha\chi_\alpha^{(\varkappa)}\chi_{\alpha'}^{(\varkappa)}=\frac{h f^{(\varkappa)}}{e^{(\varkappa)}}.
\end{equation*}
Ferner ist
\begin{equation*}\tag{4.}
\sum_\varkappa\frac{e^{(\varkappa)}}{f^{(\varkappa)}}\chi_\alpha^{(\varkappa)}\chi_\beta^{(\varkappa)}=\frac{h h_{\alpha\beta}}{h_\alpha h_\beta},
\end{equation*}
also nach (13.), \S~2 für $\beta=0$
\begin{equation*}\tag{5.}
\sum_\varkappa e^{(\varkappa)}\chi_\alpha^{(\varkappa)}=0,
\end{equation*}
falls $\alpha$ von 0 verschieden ist, dagegen für $\alpha=0$
\begin{equation*}\tag{6.}
\sum_\varkappa e^{(\varkappa)}f^{(\varkappa)}=h.
\end{equation*}
Die Gleichung (3.) schreibe ich auch in der einfacheren Form
\begin{equation*}\tag{7.}
\sum_\alpha h_\alpha\chi_\alpha\chi_{\alpha'}=\frac{hf}{e}
\end{equation*}
und die Gleichung (2.) in der Form
\begin{equation*}\tag{8.}
\sum_\alpha h_\alpha\chi_\alpha\psi_{\alpha'}=0,
\end{equation*}
wo $\psi_\alpha$ einen von $\chi_\alpha$ verschiedenen Charakter bedeutet.


Nach Gleichung (5.), \S~1 und (6.), \S~2 ist
\[
h_\alpha h_\gamma\chi_\alpha\chi_{\gamma'}=f\sum_\beta h_{\alpha\beta\gamma'}\chi_\beta=f\sum_\beta h_{\alpha'\beta\gamma}\chi_\beta.
\]
Multiplicirt man mit $\chi_\gamma$ und summirt nach $\gamma$, so ergiebt sich nach (7.) die Formel
\begin{equation*}\tag{9.}
\frac{h}{e}h_\alpha\chi_\alpha=\sum_{\beta,\gamma}h_{\alpha'\beta\gamma}\chi_\beta\chi_\gamma,
\end{equation*}
die für $\alpha=0$ mit (7.) übereinstimmt. In derselben Weise findet man die Formel
\begin{equation*}\tag{10.}
\sum_{\beta,\gamma}h_{\alpha\beta\gamma}\chi_\beta\psi_\gamma=0,
\end{equation*}
die für $\alpha=0$ in (8.) übergeht.

Eliminirt man aus den Gleichungen (6.), \S~2 und (9.) die Producte $\chi_\beta\chi_\gamma$, so erhält man die linearen Gleichungen
\begin{equation*}\tag{11.}
\frac{h}{ef}h_\alpha\chi_\alpha=\sum_\beta p_{\alpha\beta'}\chi_\beta.
\end{equation*}
Daher sind die $k$ Grössen $g^{(\varkappa)}=e^{(\varkappa)}f^{(\varkappa)}$ die Wurzeln der Gleichung $k^{\text{ten}}$ Grades
\begin{equation*}\tag{12.}
\left|g p_{\alpha\beta'}-h h_{\alpha\beta'}\right|=0.
\end{equation*}
Nach Gleichung (4.), \S~2 ist
\begin{equation*}\tag{13.}
\sum p_{\alpha\beta'}x_\alpha x_\beta=\sum_{\varkappa,\lambda}\frac{1}{h_\varkappa h_\lambda}\left(\sum_\alpha h_{\varkappa\lambda\alpha}x_\alpha\right)^2.
\end{equation*}
Diese quadratische Form, deren Determinante von Null verschieden ist, ist also eine positive bestimmte Form. Daher sind die $k$ Wurzeln $g$ der Gleichung (12.) alle reelle positive Grössen, und die Elementartheiler jener Determinante sind alle linear. Für eine $m$fache Wurzel $g=ef$ verschwinden also auch ihre Unterdeterminanten $(k-m+1)^{\text{ten}}$ und höheren Grades. Mithin ist der Rang des Systems linearer Gleichungen (11.) gleich $k-m$, und sie besitzen $m$ unabhängige Lösungen. Die einer $m$fachen Wurzel $g=ef$ entsprechenden $m$ Charaktere $\chi_\alpha$ können dann zunächst als lineare Verbindungen von solchen $m$ unabhängigen Lösungen dargestellt werden, und wenn man diese in die Gleichungen (6.), \S~2 einsetzt, findet man die Coefficienten dieser Verbindungen durch Auflösung einer Gleichung $m^{\text{ten}}$ Grades. Die ursprünglich zur Bestimmung der $k$ Charaktere gefundene Gleichung $k^{\text{ten}}$ Grades (12.), \S~2 zerfällt also nach Auflösung der Gleichung (12.) in so viele Factoren, als diese Gleichung verschiedene Wurzeln hat, und zwar entspricht einer $m$fachen Wurzel der Gleichung (12.) ein Factor $m^{\text{ten}}$ Grades.

Was nun aber die Gleichung (12.) anbetrifft, so lässt sich weiter zeigen, dass ihre Wurzeln nicht nur reell und positiv sind, sondern


dass sie ganze Zahlen (\S~6) und sogar die Quadrate von ganzen Zahlen sind, die in $h$ aufgehen. Es ist mir aber bis jetzt nicht gelungen, diese Sätze durch so einfache Betrachtungen zu beweisen, wie die obigen Ergebnisse.

Die Factoren $f^{(\varkappa)}$ sind bisher beliebig, aber von Null verschieden angenommen. Mittels der Formeln (3.) ergeben sich aus ihnen die Factoren $e^{(\varkappa)}$. Da $g^{(\varkappa)}=e^{(\varkappa)}f^{(\varkappa)}$ eine reelle positive Grösse ist, so soll auch $f^{(\varkappa)}$ reell und positiv gewählt werden. Dann gilt dasselbe von $e^{(\varkappa)}$. Zu den einfachsten Formeln gelangt man, wenn man $f=\sqrt{g}$ setzt. Dann ist $e=f=\chi_0$. Wie schon erwähnt, ist dann $e$ eine positive in $h$ aufgehende ganze Zahl, und es lässt sich zeigen, dass $\chi_\alpha$ eine Summe von $f$ Einheitswurzeln des Grades $n$ ist, falls $n$ die Ordnung der Elemente der $\alpha^{\text{ten}}$ Classe bedeutet.

Nach Gleichung (1.) und (6.), \S~1 und (4.), \S~2 ist
\begin{equation*}\tag{14.}
\sum_\beta p_{\alpha\beta'}=h h_\alpha.
\end{equation*}
Nun genügt man nach (6.), \S~1 den Gleichungen (6.), \S~2, indem man alle Grössen $\chi_\alpha=f$ setzt. Der entsprechende Werth von $ef$ ist 1. Daher kann man $e=f=1$, also $\chi_\alpha=1$ setzen. Dieser Charakter möge der Hauptcharakter genannt werden. Wählt man in der Gleichung (8.) diesen für $\psi$, so erkennt man, dass jeder andere Charakter der Gleichung
\begin{equation*}\tag{15.}
\sum_\alpha h_\alpha\chi_\alpha=0
\end{equation*}
genügt.

Da die Determinante (8.), \S~2 von Null verschieden ist, so hat die Gleichung (12.), \S~2 keine mehrfachen Wurzeln. Auch können, wenn $\chi$ und $\psi$ zwei verschiedene Charaktere sind, die Grössen $\chi_0,\chi_1,\ldots,\chi_{k-1}$ nicht den Grössen $\psi_0,\psi_1,\ldots,\psi_{k-1}$ proportional sein. Ebenso ist jede Classe $(\alpha)$ durch die $k$ entsprechenden Werthe $\chi_\alpha^{(\varkappa)}$ vollständig bestimmt, es kann nicht, wenn $\alpha$ und $\beta$ verschieden sind, $\chi_\alpha^{(\varkappa)}=\chi_\beta^{(\varkappa)}$ für alle Werthe von $\varkappa$ sein. Da die Coefficienten der Gleichung (12.), \S~2 reell sind, so entspricht, falls $x_0,x_1,\ldots,x_{k-1}$ reelle Variable sind, jeder complexen Wurzel eine conjugirte complexe Wurzel. Sind also $\chi_\alpha$ und $\chi'_{\alpha}$ conjugirte complexe Grössen, so entspricht jedem Charakter $\chi_\alpha$ ein conjugirter complexer Charakter $\chi'_{\alpha}$. Es muss aber $\chi'_{\alpha}=\chi_{\alpha'}$ sein. Denn zunächst ist nach (6.), \S~2
\[
h_\beta h_\gamma\chi_\beta\chi_\gamma=f\sum_\alpha h_{\alpha\beta\gamma}\chi_\alpha=f\sum_\alpha h_{\alpha'\beta\gamma}\chi_{\alpha'},
\]
und folglich ist $\psi_\alpha=\chi_{\alpha'}$ ein Charakter. Wäre dieser von $\chi'_\alpha$ verschieden, so wäre nach (8.) $\sum h_\alpha\chi'_{\alpha}\psi_{\alpha'}=0$, demnach $\sum h_\alpha\chi_\alpha\chi'_{\alpha}=0$, während jedes Glied als Product von zwei conjugirten complexen Grössen


positiv ist. In jedem Charakter entsprechen also (falls $f$ reell angenommen wird) inversen Classen conjugirte complexe Werthe $\chi_\alpha$ und $\chi_{\alpha'}$, und mithin einer zweiseitigen Classe ein reeller Werth. Aus jedem complexen Charakter $\chi_0,\chi_1,\ldots,\chi_{k-1}$ ergiebt sich ein conjugirter complexer Charakter, und dieser ist $\chi_{0'},\chi_{1'},\ldots,\chi_{(k-1)'}$. Ein reeller Charakter wird auch ein \emph{zweiseitiger}, zwei conjugirte complexe Charaktere werden auch \emph{inverse} genannt (vergl. \textsc{Weber}, \emph{Beweis des Satzes, dass jede eigentlich primitive quadratische Form unendlich viele Primzahlen darzustellen fähig ist}. Math. Ann. Bd. 20, S. 308).

\section*{\S~4.}

Wenn man die Gleichung (6.), \S~2
\[
h_\beta h_\gamma\chi_\beta^{(\varkappa)}\chi_\gamma^{(\varkappa)}=f^{(\varkappa)}\sum_\lambda h_{\lambda\beta\gamma}\chi_\lambda^{(\varkappa)}
\]
mit $\frac{e^{(\varkappa)}\chi_\alpha^{(\varkappa)}}{f^{(\varkappa)2}}$ multiplicirt und dann nach $\varkappa$ summirt, so erhält man
\begin{equation*}\tag{1.}
\frac{h h_{\alpha\beta\gamma}}{h_\alpha h_\beta h_\gamma}=\sum_\varkappa \frac{e^{(\varkappa)}}{f^{(\varkappa)2}}\chi_\alpha^{(\varkappa)}\chi_\beta^{(\varkappa)}\chi_\gamma^{(\varkappa)},
\end{equation*}
und ebenso gilt allgemein für $n$ Indices $\alpha,\beta,\gamma\ldots\nu$ die Formel
\begin{equation*}\tag{2.}
\frac{h h_{\alpha\beta\gamma\ldots\nu}}{h_\alpha h_\beta h_\gamma\cdots h_\nu}
=\sum_\varkappa \frac{e^{(\varkappa)}}{f^{(\varkappa)n-1}}\chi_\alpha^{(\varkappa)}\chi_\beta^{(\varkappa)}\chi_\gamma^{(\varkappa)}\cdots\chi_\nu^{(\varkappa)},
\end{equation*}
die für $n=2$ in (4.), \S~3 übergeht. Daher ist nach (11.), \S~2
\begin{equation*}\tag{3.}
h\sum_{\alpha,\beta,\gamma,\ldots,\nu}h_{\alpha\beta\gamma\ldots\nu}x_\alpha x_\beta x_\gamma\cdots x_\nu=\sum_\varkappa e^{(\varkappa)}f^{(\varkappa)}\xi^{(\varkappa)n}.
\end{equation*}
Nach (4.), \S~2 und (5.), \S~1 ist
\[
p_{\alpha\beta}=\sum_{\varkappa,\lambda}\frac{h_{\lambda\varkappa\alpha}h_{\varkappa'\lambda'\beta}}{h_\varkappa h_\lambda},
\]
also wenn man die Summation nach $\varkappa$ mittelst der Formel (1.), \S~2 ausführt,
\begin{equation*}\tag{4.}
p_{\alpha\beta}=\sum_\lambda \frac{1}{h_\lambda}h_{\alpha\beta\lambda'},
\end{equation*}
und folglich nach (3.), \S~3 und (2.)
\begin{equation*}\tag{5.}
\frac{p_{\alpha\beta}}{h_\alpha h_\beta}=\sum_\varkappa \frac{\chi_\alpha^{(\varkappa)}\chi_\beta^{(\varkappa)}}{f^{(\varkappa)2}}.
\end{equation*}
Mithin ist
\begin{equation*}\tag{6.}
\sum_{\alpha,\beta}p_{\alpha\beta}x_\alpha x_\beta=\sum_\varkappa(\xi^{(\varkappa)})^2
\end{equation*}
und
\begin{equation*}\tag{7.}
\sum_{\alpha,\beta}p_{\alpha\beta'}x_\alpha x_\beta=\sum_\varkappa \xi^{(\varkappa)}\xi^{(\varkappa)'}.
\end{equation*}



Da $\xi^{(\varkappa)}$ und $\xi^{(\varkappa)'}$ conjugirte complexe Grössen sind, so zeigt diese Formel wieder, dass diese Form eine positive ist.

Die Coefficienten der Gleichung (12.), \S~3, die zur Bestimmung der Grössen $g=ef$ dient, kann man auch berechnen, indem man die Potenzsummen ihrer Wurzeln bestimmt. Betrachtet man $\frac{h}{g}$ als die Unbekannte, so ist die Summe ihrer Wurzeln
\[
\begin{aligned}
\sum_\varkappa \frac{h}{e^{(\varkappa)}f^{(\varkappa)}}
&=\sum_\alpha \frac{p_{\alpha\alpha'}}{h_\alpha}
=\sum_{\alpha,\beta}\frac{h_{\alpha\alpha'\beta\beta'}}{h_\alpha h_\beta} \\
&=\sum_{\alpha,\beta,\gamma}\frac{h_{\alpha\beta\gamma}h_{\alpha'\beta'\gamma'}}{h_\alpha h_\beta h_\gamma}
=\sum_{\alpha,\beta,\gamma}\frac{h_{\alpha\beta\gamma}^{2}}{h_\alpha h_\beta h_\gamma}.
\end{aligned}
\]
Setzt man das System $\frac{p_{\alpha\beta'}}{h_\alpha}$ mit sich selbst zusammen, so erhält man
\[
\sum_\varkappa \frac{p_{\alpha\varkappa'}p_{\varkappa\beta'}}{h_\alpha h_\varkappa}
=\sum_{\varkappa,\lambda,\mu}\frac{h_{\alpha\varkappa\lambda\lambda'}h_{\varkappa\beta'\mu\mu'}}{h_\alpha h_\varkappa h_\lambda h_\mu}
=\sum_{\lambda,\mu}\frac{h_{\alpha\beta\lambda\lambda'\mu\mu'}}{h_\alpha h_\lambda h_\mu}
=\frac{1}{h_\alpha}\sum_{\varkappa,\lambda,\mu}\frac{h_{\alpha\varkappa\lambda\mu}h_{\beta\varkappa\lambda\mu}}{h_\varkappa h_\lambda h_\mu}.
\]
Daher ist die Summe der Quadrate der Wurzeln
\[
\sum_\varkappa\left(\frac{h}{e^{(\varkappa)}f^{(\varkappa)}}\right)^2
=\sum_{\alpha,\beta,\gamma}\frac{h_{\alpha\alpha'\beta\beta'\gamma\gamma'}}{h_\alpha h_\beta h_\gamma}
=\sum_{\alpha,\beta,\gamma,\delta}\frac{h_{\alpha\beta\gamma\delta}^{2}}{h_\alpha h_\beta h_\gamma h_\delta}.
\]
Indem man das System $n-2$ Mal mit sich selbst componirt, findet man in derselben Weise für die Summe der $(n-2)^{\text{ten}}$ Potenzen ihrer Wurzeln
\begin{equation*}\tag{8.}
\sum_\varkappa\left(\frac{h}{e^{(\varkappa)}f^{(\varkappa)}}\right)^{n-2}
=\sum_{\alpha,\beta,\gamma\ldots\mu}\frac{h_{\alpha\alpha'\beta\beta'\ldots\mu\mu'}}{h_\alpha h_\beta\cdots h_\mu}
=\sum_{\alpha,\beta,\gamma\ldots\mu,\nu}\frac{h_{\alpha\beta\cdots\mu\nu}^{2}}{h_\alpha h_\beta\cdots h_\mu h_\nu},
\end{equation*}
wo $n$ die Anzahl der Indices $\alpha,\beta,\ldots\mu,\nu$ ist. (Vergl. \textsc{Borchardt}, Crelle's Journal Bd.~30, S.~38.) Man kann auch die $k^2$ Grössen $p_{\alpha\beta}$ auf die $k$ Grössen
\begin{equation*}\tag{9.}
p_{\alpha0}=p_\alpha=p_{\alpha'}=\sum_\lambda\frac{h_{\alpha\lambda\lambda'}}{h_\lambda}=h_\alpha\sum_\varkappa\frac{\chi_\alpha^{(\varkappa)}}{f^{(\varkappa)}}
\end{equation*}
zurückführen (von \textsc{Dedekind} a. a. O. S.~147 (25.) mit $\sigma_r$ bezeichnet). Denn nach (1.), \S~2 und (4.) ist
\[
p_{\alpha\beta}=\sum_\lambda\frac{1}{h_\lambda}\left(\sum_\gamma\frac{1}{h_\gamma}h_{\alpha\beta\gamma}h_{\gamma\lambda\lambda'}\right),
\]
und mithin
\begin{equation*}\tag{10.}
p_{\alpha\beta}=\sum_\gamma\frac{1}{h_\gamma}h_{\alpha\beta\gamma}p_\gamma.
\end{equation*}
Die Gleichung (12.), \S~3 geht also in die Gleichung (12.), \S~2 über, wenn man $x_\gamma=\frac{p_\gamma}{h_\gamma}$ und $\xi=\frac{h}{g}$ setzt. In der That ist nach (11.), \S~3 und (13.), \S~2
\begin{equation*}\tag{11.}
\sum_\alpha p_\alpha\chi_\alpha=\frac{h}{e}.
\end{equation*}
Wählt man also $f$ so, dass $\chi_0,\chi_1,\ldots\chi_{k-1}$ ganze algebraische Zahlen


werden, so ist $e$ ein Divisor von $h$. Die Formel (3.) ergiebt für diese Werthe der Variabeln
\begin{equation*}\tag{12.}
\sum_\varkappa\left(\frac{h}{e^{(\varkappa)}f^{(\varkappa)}}\right)^{n-1}
=\sum_{\alpha,\beta,\gamma\ldots\nu}h_{\alpha\beta\gamma\ldots\nu}\frac{p_\alpha p_\beta p_\gamma\cdots p_\nu}{h_\alpha h_\beta h_\gamma\cdots h_\nu}.
\end{equation*}
Die Zahl $p_\alpha$ hat folgende Bedeutung: Ist $A$ ein bestimmtes Element der $\alpha^{\text{ten}}$ Classe, und durchläuft jedes der beiden veränderlichen Elemente $R$ und $S$ unabhängig von dem anderen die $h$ Elemente von $\grp$, so ist $\frac{hp_\alpha}{h_\alpha}$ die Anzahl der Lösungen der Gleichung
\begin{equation*}\tag{13.}
SR=RSA.
\end{equation*}
Denn ist $A$ ein bestimmtes Element der $\alpha^{\text{ten}}$ Classe, durchläuft $R^{-1}$ zunächst die $h_\lambda$ Elemente der Classe $(\lambda)$, und $R'$ die $h_\lambda$ Elemente der inversen Classe $(\lambda')$, so hat die Gleichung $R^{-1}R'A=E$ genau $\frac{h_{\alpha\lambda\lambda'}}{h_\alpha}$ Lösungen.

Setzt man aber für $R'$ jedes Element der Classe $(\lambda')$ $\frac{h}{h_\lambda}$ Mal, so hat sie $\frac{h h_{\alpha\lambda\lambda'}}{h_\alpha h_\lambda}$ Lösungen. Dies erreicht man, indem man für jedes $R$ setzt $R'=S^{-1}RS$ und dann $S$ alle $h$ Elemente von $\grp$ durchlaufen lässt. Bewegt sich endlich noch $\lambda$ von $0$ bis $k-1$, so ist
\[
\sum_\lambda\frac{h h_{\alpha\lambda\lambda'}}{h_\alpha h_\lambda}=\frac{hp_\alpha}{h_\alpha}
\]
die Anzahl der Lösungen der Gleichung $R^{-1}(S^{-1}RS)A=E$, d. h. der Gleichung (13.). Speciell ist $p_0=k$, also $hk$ die Anzahl der Lösungen der Gleichung $SR=RS$.

Ebenso ist allgemeiner $hp_{\alpha\beta}$ die Anzahl der Lösungen der Gleichung
\begin{equation*}\tag{14.}
SR=RSAB,
\end{equation*}
falls $A$ die $h_\alpha$ Elemente der $\alpha^{\text{ten}}$ Classe, $B$ die $h_\beta$ der $\beta^{\text{ten}}$ Classe, $R$ und $S$ alle $h$ Elemente von $\grp$ durchlaufen. Setzt man aber für $A$ ein bestimmtes Element der $\alpha^{\text{ten}}$ Classe, so hat jene Gleichung $\frac{hp_{\alpha\beta}}{h_\alpha}$ Lösungen. Mit Hülfe der Formeln (5.) und (6.), \S~3 kann man diesen und ähnliche Sätze auch durch Rechnung beweisen.

\section*{\S~5.}

Für die weitere Entwicklung dieser Theorie ist es vortheilhaft, das Zeichen $\chi_\alpha$, falls $A$ ein Element der $\alpha^{\text{ten}}$ Classe ist, durch $\chi_A$ zu ersetzen, oder auch, indem man $\chi$ als Functionszeichen benutzt, durch $\chi(A)$. Dann ist, da diese Grösse für alle Elemente der $\alpha^{\text{ten}}$ Classe denselben Werth hat,
\begin{equation*}\tag{1.}
\chi(B^{-1}AB)=\chi(A),
\end{equation*}


oder wenn man $A$ durch $BA$ ersetzt,
\begin{equation*}\tag{2.}
\chi(AB)=\chi(BA).
\end{equation*}
Speciell ist $\chi(E)=f$. Ebenso bezeichne ich die Variabeln $x_\alpha$ und $y_\alpha$ mit $x_A$ und $y_A$ und die Constanten $h_\alpha$ und $p_\alpha$ mit $h_A$ und $p_A$. Für alle diese Grössen gilt die zu (2.) analoge Gleichung. Ferner ist $h_{A^{-1}}=h_A$, während $\chi(A^{-1})$ und $\chi(A)$ conjugirte complexe Grössen sind.

Die Gleichung (6.), \S~2, aus der sich die Verhältnisse der $\chi_\alpha$ ergeben, kann dann in der Form
\begin{equation*}\tag{3.}
h_B\chi(A)\chi(B)=f\sum_S'\chi(AS)
\end{equation*}
geschrieben werden, wo $S$ die $h_B$ mit $B$ conjugirten Elemente durchläuft. Denn sind $(\alpha)$ und $(\beta)$ die Classen von $A$ und $B$, so sind für $S$ die $h_\beta$ Elemente der $\beta^{\text{ten}}$ Classe zu setzen. Da $A$ ein bestimmtes Element der $\alpha^{\text{ten}}$ Classe ist, so kommt es $\frac{h_{\alpha\beta\gamma'}}{h_\alpha}$ Mal vor, dass $AS$ der $\gamma^{\text{ten}}$ Classe angehört, also $\chi(AS)=\chi_\gamma$ wird. Mithin ist $\sum_S'\chi(AS)=\frac{1}{h_\alpha}\sum_\gamma h_{\alpha\beta\gamma'}\chi_\gamma$.

Da $\chi(E)=f$ und $\chi(S)=\chi(B)$ ist, kann man die Gleichung (3.) auch auf die Form
\begin{equation*}\tag{4.}
\sum_S'\left(\chi(E)\chi(AS)-\chi(A)\chi(S)\right)=0
\end{equation*}
bringen, wo $S$ die Elemente einer Classe durchläuft. Setzt man $S=R^{-1}BR$ und für $R$ alle $h$ Elemente von $\grp$, so wird $S$ jedem Elemente der $\beta^{\text{ten}}$ Classe gleich und zwar jedem $\frac{h}{h_\beta}$ Mal. Daher ist
\begin{equation*}\tag{5.}
h\chi(A)\chi(B)=f\sum_R\chi(AR^{-1}BR).
\end{equation*}
Durch wiederholte Anwendung dieser Relation erhält man
\begin{equation*}\tag{6.}
\left(\frac{h}{f}\right)^{n-1}\chi(A_1)\chi(A_2)\cdots\chi(A_n)=\sum'\chi(A_1R_1A_2R_2\cdots A_nR_n),
\end{equation*}
wo $R_1,R_2,\ldots R_n$ alle Systeme von $n$ Elementen der Gruppe $\grp$ durchlaufen, welche die Gleichung
\[
R_1R_2\cdots R_n=E
\]
befriedigen. In derselben Weise erhält man die Formel
\begin{equation*}\tag{7.}
p_{\alpha\beta}=\sum_{R,S}'\frac{p_{RS}}{h_{RS}},
\end{equation*}
wo $R$ die $h_\alpha$ Elemente der $\alpha^{\text{ten}}$ Classe durchläuft und $S$ die $h_\beta$ der $\beta^{\text{ten}}$.

Durchlaufen $R$ und $S$ alle Elemente, die der Bedingung $RS=A$ genügen, so ist nach (9.) und (10.), \S~3
\begin{equation*}\tag{8.}
\frac{h}{e}\chi(A)=\sum\chi(R)\chi(S),\qquad 0=\sum\chi(R)\psi(S)\qquad(RS=A).
\end{equation*}
Denn ist $A$ ein festes Element der $\alpha^{\text{ten}}$ Classe, so kommt es, während $R$


die $\beta^{\text{te}}$ und $S$ die $\gamma^{\text{te}}$ Classe durchläuft, $\frac{h_{\alpha'\beta\gamma}}{h_\alpha}$ Mal vor, dass $RS=A$ wird. Dafür kann man auch schreiben
\begin{equation*}\tag{9.}
\frac{h}{e}\chi(AB^{-1})=\sum_R\chi(AR^{-1})\chi(RB^{-1}),\qquad 0=\sum_R\chi(AR^{-1})\psi(RB^{-1}),
\end{equation*}
wo $R$ alle Elemente von $\grp$ durchläuft. Speciell ist
\begin{equation*}\tag{10.}
\frac{hf}{e}=\sum_R\chi(R^{-1})\chi(R),\qquad 0=\sum_R\chi(R^{-1})\psi(R).
\end{equation*}
Die $h$ Variabeln $x_R$ (bez. $y_R$) reduciren sich in Folge der Bedingungen
\begin{equation*}\tag{11.}
x_{AB}=x_{BA},\qquad y_{AB}=y_{BA}
\end{equation*}
auf nur je $k$ unabhängige Variabeln, sie haben für alle Elemente $R$, die derselben Classe angehören, denselben Werth. Ich bilde aus ihnen die beiden Systeme von je $h^2$ Elementen
\begin{equation*}\tag{12.}
(x_{PQ^{-1}}),\qquad (y_{PQ^{-1}}).
\end{equation*}
Die $h$ Zeilen des ersten Systems erhält man, indem man für $P$ die $h$ Elemente von $\grp$ in irgend einer Reihenfolge setzt, die $h$ Spalten, indem man für $Q$ dieselben Elemente in derselben Reihenfolge setzt. Das System hat gewisse durch $\grp$ bestimmte Symmetrieeigenschaften, auf Grund deren unter seinen $h^2$ Elementen nur $k$ verschiedene sind. Das aus den beiden Systemen (12.) zusammengesetzte System hat dieselben Symmetrieeigenschaften und ist ausserdem unabhängig davon, in welcher Reihenfolge die Zusammensetzung erfolgt. Es ist
\begin{equation*}\tag{13.}
z_{PQ^{-1}}=\sum_R x_{PR^{-1}}y_{RQ^{-1}}=\sum_S y_{PS^{-1}}x_{SQ^{-1}}.
\end{equation*}
Denn setzt man $PR^{-1}=SQ^{-1}$, so wird $S=PR^{-1}Q$, $S$ durchläuft also gleichzeitig mit $R$ alle $h$ Elemente von $\grp$, und es ist $S^{-1}P=Q^{-1}R$, also $y_{PS^{-1}}=y_{S^{-1}P}=y_{Q^{-1}R}=y_{RQ^{-1}}$. Ersetzt man ferner $R$ durch $RP$, so erkennt man, dass die erste Summe nur von dem Producte $PQ^{-1}$ abhängt. Man hätte sie zunächst nur mit $z_{P,Q}$ bezeichnen dürfen, darf sie aber nun einer von nur einem Elemente $R=PQ^{-1}$ abhängigen Grösse $z_R=z_{PQ^{-1}}$ gleichsetzen. Endlich ist
\[
z_{AB}=\sum x_{AR^{-1}}y_{RB}=\sum y_{AS^{-1}}x_{SB}=\sum x_{BS}y_{S^{-1}A}=\sum x_{BR^{-1}}y_{RA}=z_{BA}.
\]
Denn $R=S^{-1}$ durchläuft gleichzeitig mit $S$ alle $h$ Elemente von $\grp$.

Je zwei Systeme (12.), welche die hier vorausgesetzten Symmetrieeigenschaften besitzen, sind mit einander vertauschbar. Setzt man $x_\varkappa=1$ und die anderen $k-1$ der Variabeln $x_0,x_1,\ldots x_{k-1}$ Null, so erhält man $k$ specielle Systeme dieser Art für $\varkappa=0,1,\ldots k-1$. Von diesen $k$ Systemen sind also je zwei vertauschbar. Aus ihnen geht wieder das allgemeinste System hervor, indem man sie mit den $k$ Variabeln


$x_0,x_1,\ldots x_{k-1}$ multiplicirt und addirt. Folglich ist die Determinante $h^{\text{ten}}$ Grades
\begin{equation*}\tag{14.}
\left|x_{PQ^{-1}}\right|=\Theta(x_0,x_1,\ldots x_{k-1})=\Theta((x))
\end{equation*}
ein Product von $h$ linearen Functionen der $k$ Variabeln (\emph{Über vertauschbare Matrizen}, S.~601 dieses Bandes). In dieser Determinante sind die $h$ Elemente der Diagonale und nur diese gleich $x_0$, also hat $x_0^h$ den Coefficienten 1. Daher kann man festsetzen, dass auch in jedem linearen Factor von $\Theta$
\begin{equation*}\tag{15.}
\xi=\frac{1}{f}\sum h_\alpha\chi_\alpha x_\alpha
\end{equation*}
der Coefficient von $x_0$ gleich 1 ist. Hier soll $h_\alpha$ dieselbe Bedeutung haben wie bisher, $f$ von Null verschieden sein, und $\chi_0,\chi_1,\ldots\chi_{k-1}$ sollen die noch unbekannten Coefficienten bezeichnen. Demnach ist $\chi_0=f$. Nun ist aber nach (13.)
\[
\left|z_{PQ^{-1}}\right|=\left|x_{PQ^{-1}}\right|\left|y_{PQ^{-1}}\right|.
\]
Folglich muss auch jeder lineare Factor von $\Theta((z))$ in das Product aus einer linearen Function von $x_0,x_1,\ldots x_{k-1}$ und einer von $y_0,y_1,\ldots y_{k-1}$ zerfallen,
\[
\frac{1}{f}\sum h_\gamma\chi_\gamma z_\gamma=\left(\sum a_\alpha x_\alpha\right)\left(\sum b_\beta y_\beta\right),
\]
wo $a_0=b_0=1$ ist. Setzt man $y_0=y_E=1$ und $y_1=\cdots=y_{k-1}=0$, so wird $z_R=x_R$, also $\frac{1}{f}\sum h_\gamma\chi_\gamma x_\gamma=\sum a_\alpha x_\alpha$. Mithin ist
\begin{equation*}\tag{16.}
f\sum h_\gamma\chi_\gamma z_\gamma=\left(\sum h_\alpha\chi_\alpha x_\alpha\right)\left(\sum h_\beta\chi_\beta y_\beta\right).
\end{equation*}
Nun ist aber die Gleichung (13.) oder
\begin{equation*}\tag{17.}
z_C=\sum x_Ry_S\qquad(RS=C)
\end{equation*}
identisch mit
\begin{equation*}\tag{18.}
h_\gamma z_\gamma=\sum_{\alpha,\beta}h_{\alpha\beta\gamma'}x_\alpha y_\beta.
\end{equation*}
Denn ist $C$ ein Element der $\gamma^{\text{ten}}$ Classe, so kommt es, während $R$ die Elemente der $\alpha^{\text{ten}}$ Classe durchläuft und $S$ die der $\beta^{\text{ten}}$, $\frac{h_{\alpha\beta\gamma'}}{h_\gamma}$ Mal vor, dass $RS=C$ wird. Durch Vergleichung der Coefficienten von $x_\alpha y_\beta$ ergiebt sich daher aus (16.)
\begin{equation*}\tag{19.}
h_\alpha h_\beta\chi_\alpha\chi_\beta=f\sum_\gamma h_{\alpha\beta\gamma'}\chi_\gamma,
\end{equation*}
und folglich ist das Werthsystem $\chi_0,\chi_1,\ldots\chi_{k-1}$ einem der $k$ verschiedenen Charaktere gleich. Die Gleichung (15.) lässt sich auf die Form
\begin{equation*}\tag{20.}
f\xi=\sum\chi(R)x_R
\end{equation*}
bringen.


Ist umgekehrt $\chi$ irgend einer der $k$ Charaktere, so ist $\xi$ ein linearer Factor der Determinante $\Theta$. Setzt man nämlich in den Gleichungen (17.) und (18.) $y_\beta=\chi_{\beta'}$, so wird
\[
z_C=\sum_S x_{CS^{-1}}y_S=\sum_R x_{CR}\chi(R),\qquad h_\gamma z_\gamma=\sum_{\alpha,\beta}h_{\alpha\beta\gamma'}x_\alpha\chi_{\beta'}=\frac{1}{f}\sum_\alpha h_\alpha h_\gamma\chi_\alpha\chi_{\gamma'}x_\alpha,
\]
also
\begin{equation*}\tag{21.}
\sum_R \chi(R)x_{CR}=\chi(C^{-1})\xi,
\qquad
\sum_R \chi(CR)x_R=\chi(C)\xi.
\end{equation*}
In der Determinante (14.) multiplicire man die Elemente der ersten Zeile $(P=E)$ mit der von Null verschiedenen Zahl $\chi(E)=f$, die Elemente der Zeile, die durch den Index $P$ charakterisirt ist, mit $\chi(P)$, und addire dann alle Zeilen zur ersten. Dann wird irgend ein Element der ersten Zeile
\[
\sum_P\chi(P)x_{PQ^{-1}}=\chi(Q)\xi,
\]
und folglich ist $\Theta$ durch $\xi$ theilbar. Daher ist
\begin{equation*}\tag{22.}
\Theta=\left|x_{PQ^{-1}}\right|=\prod_\varkappa\left(\xi^{(\varkappa)}\right)^{g^{(\varkappa)}}
=\prod_\varkappa\left(\frac{1}{f^{(\varkappa)}}\sum h_\alpha\chi_\alpha^{(\varkappa)}x_\alpha\right)^{g^{(\varkappa)}},
\end{equation*}
wo $g^{(\varkappa)}$ eine von Null verschiedene ganze Zahl ist.

Mit Hülfe der Sätze über die Matrizen (linearen Systeme), die ich in meiner Arbeit \emph{Über lineare Substitutionen und bilineare Formen}, Crelle's Journal Bd.~84 entwickelt habe, ergiebt sich für diesen Satz ein zweiter Beweis, aus dem zugleich die Bedeutung der Zahlen $g^{(\varkappa)}$ erhellt. Das bisher mit $(x_{PQ^{-1}})$ bezeichnete System von $h$ Zeilen und Spalten will ich noch kürzer mit $(x)$ bezeichnen. Ist $\varepsilon_R=0$, wenn $R$ von $E$ verschieden ist, aber $\varepsilon_E=1$, so ist $(\varepsilon_{PQ^{-1}})=(\varepsilon)$ das Einheitssystem, das aus $(x)$ hervorgeht, indem man $x_0=1$, $x_1=x_2=\cdots=x_{k-1}=0$ setzt. Dann wird der Inhalt der Gleichungen (5.) und (6.), \S~3 ausgedrückt durch die Formel
\begin{equation*}\tag{23.}
\sum_\varkappa\left(\frac{1}{h}e^{(\varkappa)}\chi^{(\varkappa)}\right)=(\varepsilon)
\end{equation*}
und der der Gleichungen (9.) durch die Formeln
\begin{equation*}\tag{24.}
\left(\frac{1}{h}e^{(\varkappa)}\chi^{(\varkappa)}\right)^2=\left(\frac{1}{h}e^{(\varkappa)}\chi^{(\varkappa)}\right),
\qquad
\left(\frac{1}{h}e^{(\varkappa)}\chi^{(\varkappa)}\right)\left(\frac{1}{h}e^{(\lambda)}\chi^{(\lambda)}\right)=0.
\end{equation*}
Lässt man in der ersten den Index $\varkappa$ weg, so zeigt sie: Die Gleichung niedrigsten Grades, der das System $\left(\frac{e\chi}{h}\right)$ genügt, ist $\psi\left(\left(\frac{e\chi}{h}\right)\right)=0$, wenn $\psi(r)=r(r-1)$ ist. Daher kann auch die charakteristische Determinante $\varphi(r)$ dieses Systems nur für $r=0$ und $r=1$ verschwinden, und weil $\psi(r)$ keinen mehrfachen Linearfactor hat, so sind ihre Elementartheiler alle vom ersten Grade. Ist also
\begin{equation*}\tag{25.}
\varphi(r)=\left|r(\varepsilon)-\left(\frac{e\chi}{h}\right)\right|=r^{h-g}(r-1)^g,
\end{equation*}


so verschwinden für $r=0$ alle Unterdeterminanten von $\varphi(0)$ von dem Grade $h,h-1,\ldots g+1$, die vom Grade $g$ aber nicht alle. Folglich ist $g$ der Rang des Systems $(\chi)$.

Durch Auflösung der linearen Gleichungen (11.), \S~2 erhält man
\begin{equation*}\tag{26.}
hx_\alpha=\sum_\varkappa \chi_\alpha^{(\varkappa)}e^{(\varkappa)}\xi^{(\varkappa)}.
\end{equation*}
Mithin ist
\begin{equation*}\tag{27.}
(x)=\sum_\varkappa \xi^{(\varkappa)}\left(\frac{1}{h}e^{(\varkappa)}\chi^{(\varkappa)}\right)
\end{equation*}
und folglich
\[
\prod_\varkappa\left(r(\varepsilon)-\xi^{(\varkappa)}\left(\frac{1}{h}e^{(\varkappa)}\chi^{(\varkappa)}\right)\right)=r^h(\varepsilon)-r^{h-1}(x),
\]
weil in der Entwicklung dieses symbolischen Productes alle anderen Glieder nach (24.) verschwinden. Mithin sind auch die Determinanten dieser beiden Systeme gleich
\[
\prod_\varkappa\left(r^{h-g^{(\varkappa)}}(r-\xi^{(\varkappa)})^{g^{(\varkappa)}}\right)=r^{h(h-1)}\left|r(\varepsilon)-(x)\right|.
\]
Daher müssen zunächst die Potenzen von $r$ auf beiden Seiten dieser Gleichung gleiche Exponenten haben. Hebt man sie auf, und setzt man dann $r=0$, so ergiebt sich die Gleichung (22.). Folglich ist darin $g^{(\varkappa)}$ der Rang des Systems
\begin{equation*}\tag{28.}
\left(\chi^{(\varkappa)}(PQ^{-1})\right).
\end{equation*}

\section*{\S~6.}

In der Determinante $h^{\text{ten}}$ Grades $\left|x_{P,Q}\right|$ seien zunächst die $h^2$ Elemente $x_{P,Q}$ unabhängige Variable. Ist $T$ ein bestimmtes Element von $\grp$, und setzt man
\[
T^{-1}PT=P',\qquad T^{-1}QT=Q',
\]
so durchläuft $P'$ zugleich mit $P$ die $h$ Elemente von $\grp$, und $Q'$ dieselben Elemente in derselben Reihenfolge. Daher ist
\begin{equation*}\tag{1.}
\Theta=\left|x_{P,Q}\right|=\left|x_{P',Q'}\right|,
\end{equation*}
und auch die Unterdeterminante $\Theta_{A,B}$, die dem Elemente $x_{A,B}$ in der ersten Determinante complementär ist, ist gleich der Unterdeterminante, welche demselben Elemente $x_{A,B}$, das dort aber an einer anderen Stelle steht, in der zweiten Determinante complementär ist.

Ich beschränke nun die Veränderlichkeit der $h^2$ Elemente der Determinante $\Theta$ zunächst dadurch, dass ich $x_{P,Q}=x_{PQ^{-1}}$ setze. Dann hängt sie nur noch von $h$ unabhängigen Variabeln ab, je $h$ der $h^2$ Elemente sind einander gleich, in jeder Zeile stehen die $h$ verschiedenen Variabeln sämmtlich, ebenso in jeder Spalte. Die ver-


schiedenen Zeilen unterscheiden sich von einander nur durch die Reihenfolge der Variabeln, die durch die Constitution der Gruppe $\grp$ bedingt ist. Ich behaupte nun, dass auch die Unterdeterminante $\Theta_{A,B}$ nicht von $A$ und $B$ einzeln, sondern nur von $AB^{-1}=C$ abhängt. Oder: Die Variable $x_C$ findet sich in $\Theta$ an $h$ Stellen. An jeder ist ihr dieselbe Unterdeterminante complementär. Zunächst ist für eine Variable der Diagonale $\Theta_{A,A}=\left|x_{R,S}\right|$, wo $R$ die Elemente von $\grp$ mit Ausschluss von $A$ durchläuft, und $S$ dieselben $h-1$ Elemente in derselben Reihenfolge. Setzt man $R=UA$, $S=VA$, so ist
\[
x_{R,S}=x_{UA,VA}=x_{UAA^{-1}V^{-1}}=x_{UV^{-1}}=x_{U,V},
\]
also $\Theta_{A,A}=\left|x_{U,V}\right|$, wo $U$ und $V$ die Elemente von $\grp$ mit Ausschluss von $E$ durchlaufen. Mithin ist $\left|x_{U,V}\right|=\Theta_{E,E}$.

Sind $A$ und $B$ zwei bestimmte Elemente von $\grp$, und ist $C=AB^{-1}$, so durchläuft $CP$ die $h$ Elemente von $\grp$ gleichzeitig mit $P$, nur in einer anderen Reihenfolge. Daher ist
\[
\pm\left|x_{P,Q}\right|=\left|x_{CP,Q}\right|=\left|y_{P,Q}\right|,
\]
falls man $x_{CPQ^{-1}}=y_{PQ^{-1}}=y_{P,Q}$ setzt. Nun ist $\Theta_{B,B}(y)=\Theta_{E,E}(y)$ und mithin
\begin{equation*}\tag{2.}
\Theta_{A,B}=\Theta_{C,E}\qquad(C=AB^{-1}).
\end{equation*}
Aus diesem Satze ergiebt sich die Relation
\begin{equation*}\tag{3.}
\Theta_{A,B}=\frac{1}{h}\frac{\partial\Theta}{\partial x_{AB^{-1}}}.
\end{equation*}
Die Theorie dieser und noch allgemeinerer Determinanten, deren Grundlage die vorliegende Untersuchung bildet, behandle ich in der Eingangs angekündigten Arbeit. Jetzt aber beschränke ich die Veränderlichkeit der $h$ Grössen $x_R$ weiter durch die Bedingungen (11.), \S~5, ich setze also $x_R=x_S$, wenn $R$ und $S$ conjugirt sind. Gehört $C$ der $\gamma^{\text{ten}}$ Classe an, so findet sich jetzt die Variable $x_C=x_\gamma$ in jeder Zeile $h_\gamma$ Mal, in der ganzen Determinante $h h_\gamma$ Mal. Aber auch hier hat die Unterdeterminante $\Theta_{A,B}$ für jeden dieser $h h_\gamma$ Plätze denselben Werth.

Unter der gemachten Voraussetzung haben die beiden Determinanten (1.) nicht nur denselben Werth, sondern es sind auch die an entsprechenden Stellen stehenden Variabeln gleich, $x_{P,Q}=x_{P',Q'}$. Daher ist auch die Unterdeterminante, die der Variabeln $x_{P,Q}$ in der ersten complementär ist, gleich der Unterdeterminante, die der Variabeln $x_{P',Q'}$ in der zweiten complementär ist; diese aber ist, wie schon oben bemerkt, gleich der Unterdeterminante, die der Variabeln $x_{P',Q'}$ in der ersten complementär ist. In Verbindung mit (2.) folgt daraus die Richtigkeit der Behauptung
\begin{equation*}\tag{4.}
\Theta_{P,Q}=\Theta_{R,S},\quad \text{wenn } RS^{-1}=T^{-1}PQ^{-1}T.
\end{equation*}



Denn es ist $RS^{-1}=P'Q'^{-1}$, also
\[
\Theta_{P,Q}=\Theta_{P',Q'}=\Theta_{R,S}.
\]
Mithin ist jetzt
\begin{equation*}\tag{5.}
\Theta_{A,B}=\frac{1}{h h_{AB^{-1}}}\frac{\partial\Theta}{\partial x_{AB^{-1}}}.
\end{equation*}
Aus den identischen Gleichungen zwischen den Elementen einer Determinante und den ihnen complementären Unterdeterminanten ergiebt sich daher
\[
\sum_R \frac{1}{h_R}\frac{\partial\log\Theta}{\partial x_R}x_{RA^{-1}}=h\varepsilon_A,
\]
also wenn man für $\Theta$ seinen Ausdruck (22.), \S~5 einsetzt,
\[
\sum_R\sum_\varkappa \frac{g^{(\varkappa)}}{f^{(\varkappa)}\xi^{(\varkappa)}}\chi_R^{(\varkappa)}x_{RA^{-1}}=h\varepsilon_A,
\]
und weil nach (21.), \S~5
\[
\sum_R \chi_R^{(\varkappa)}x_{RA^{-1}}=\chi^{(\varkappa)}(A)\xi^{(\varkappa)}
\]
ist,
\[
\sum_\varkappa \frac{g^{(\varkappa)}}{f^{(\varkappa)}}\chi^{(\varkappa)}(A)=h\varepsilon_A,
\]
und weil nach (5.) und (6.), \S~3 auch
\begin{equation*}\tag{6.}
\sum_\varkappa e^{(\varkappa)}\chi^{(\varkappa)}(A)=h\varepsilon_A
\end{equation*}
ist,
\[
\sum_\varkappa\left(\frac{g^{(\varkappa)}}{f^{(\varkappa)}}-e^{(\varkappa)}\right)\chi_\alpha^{(\varkappa)}=0.
\]
Da die Determinante (8.), \S~2 von Null verschieden ist, so folgt daraus
\begin{equation*}\tag{7.}
g^{(\varkappa)}=e^{(\varkappa)}f^{(\varkappa)},
\end{equation*}
und damit ist bewiesen, dass die Wurzeln $g=ef$ der Gleichung (12.), \S~3 ganze Zahlen sind.

Die charakteristische Function von $\Theta$ erhält man, indem man $x_0$ durch $x_0-r$, also auch $\xi^{(\varkappa)}$ durch $\xi^{(\varkappa)}-r$ ersetzt. Daher ist
\[
\Theta(x_0-r,x_1,\ldots x_{k-1})=\prod_\varkappa (\xi^{(\varkappa)}-r)^{g^{(\varkappa)}}
\]
und folglich nach (3.), \S~4
\begin{equation*}\tag{8.}
\begin{aligned}
\frac{1}{h}\frac{\partial\log\Theta(x_0-r,x_1,\ldots x_{k-1})}{\partial r}
&=\frac{1}{h}\sum_\varkappa \frac{e^{(\varkappa)}f^{(\varkappa)}}{r-\xi^{(\varkappa)}} \\
&=r^{-1}+x_0r^{-2}+\left(\sum_{\alpha,\beta}h_{\alpha\beta}x_\alpha x_\beta\right)r^{-3}
+\left(\sum_{\alpha,\beta,\gamma}h_{\alpha\beta\gamma}x_\alpha x_\beta x_\gamma\right)r^{-4}\\
&\quad+\left(\sum_{\alpha,\beta,\gamma,\delta}h_{\alpha\beta\gamma\delta}x_\alpha x_\beta x_\gamma x_\delta\right)r^{-5}+\cdots .
\end{aligned}
\end{equation*}
Demnach ist dieser Ausdruck die erzeugende Function der sämmtlichen in \S~1 definirten Zahlen $h_{\alpha\beta\gamma\ldots}$. Nachdem aus dieser Entwicklung die Zahlen $h_{\alpha\beta\gamma}$ erhalten sind, liefert die Determinante (14.), \S~2 das Product der $k$ verschiedenen Linearfactoren von $\Theta$. Die Beziehung zwischen diesen beiden Determinanten $h^{\text{ten}}$ und $k^{\text{ten}}$ Grades ist wohl eins der merkwürdigsten Ergebnisse der entwickelten Theorie.

Aus der Gleichung (5.) folgt noch, dass, wenn $\xi^g$ ein Factor von $\Theta$ ist, die Unterdeterminanten $(h-1)^{\text{ten}}$ Grades von $\Theta$ alle durch $\xi^{g-1}$ theilbar sind. Daher sind die Unterdeterminanten $(h-2)^{\text{ten}}$ Grades alle durch $\xi^{g-2}$, \ldots{} die $(h-g+1)^{\text{ten}}$ Grades alle durch $\xi$ theilbar, die $(h-g)^{\text{ten}}$ Grades aber nicht mehr alle durch $\xi$ theilbar. Sind also $x_0,x_1,\ldots x_{k-1}$ unabhängige Variable, so hat $\Theta$ nur lineare Elementartheiler.

Sind je zwei Elemente von $\grp$ vertauschbar, so ist $k=h$, die Charaktere, die durch die Bedingungen
\begin{equation*}\tag{9.}
\chi(AB)=\chi(A)\chi(B)
\end{equation*}
bestimmt sind, sind alle vom Grade $f=1$, die Determinante $h^{\text{ten}}$ Grades $\Theta$, die mit der Determinante $k^{\text{ten}}$ Grades (14.), \S~2 identisch wird, enthält jeden Linearfactor in der Potenz $e=1$, und ihre Zerlegung ist mit Hülfe der Relationen (9.) leicht direct auszuführen. Für cyklische Gruppen, bei denen $\xi$ in die Resolvente von \textsc{Lagrange} übergeht, ist sie schon 1853 von \textsc{Spottiswoode}, \emph{Elementary Theorems relating to Determinants}, \textsc{Crelle}'s Journal Bd.~51, S.~375 angegeben. Andere Fälle haben \textsc{Nöther}, \emph{Notiz über eine Classe symmetrischer Determinanten}, Math. Ann. Bd.~16, \textsc{Gegenbauer}, \emph{Über eine specielle symmetrische Determinante}, Wiener Ber. 1880, \textsc{Puchta}, \emph{Ein neuer Satz aus der Theorie der Determinanten}, Wiener Denkschriften Bd.~43 behandelt. Die allgemeine Formel für die Determinante einer commutativen Gruppe hat \textsc{Dedekind} im Jahre 1880 bei Gelegenheit der Eingangs erwähnten Untersuchung über Gruppendeterminanten durch Multiplication der Determinanten $\left|x_{PQ^{-1}}\right|$ und $\left|\chi^{(\varkappa)}(R)\right|$ erhalten. Wie ich aus dem 25. Bande der \emph{Fortschritte der Mathematik}, S.~220 ersehe, hat \textsc{Burnside} in der (mir nicht zugänglichen) Zeitschrift \emph{Messenger of Math.} (2) XXIII, p.~112 in einer Arbeit \emph{On a property of certain determinants} ebenfalls diese Formel hergeleitet.

\section*{\S~7.}

Die bisher entwickelten Sätze und Formeln behalten alle ihre Gültigkeit, wenn man den Begriff der Classe, auf dem sie fussen, weiter fasst. Ist die Gruppe $\grp$ eine invariante Untergruppe einer anderen Gruppe $\grp'$, so mögen zwei Elemente $R$ und $R'$ von $\grp$ conjugirt heissen (in Bezug auf $\grp'$), wenn es in $\grp'$ ein Element $T$ giebt, das der Bedingung
\begin{equation*}\tag{1.}
R'=T^{-1}RT
\end{equation*}

genügt. Sei $T$ ein bestimmtes Element von $\grp'$. Dann entspricht auf Grund der Gleichung (1.) jedem Elemente $R$ von $\grp$ ein Element $R'$ von $\grp$. Durchläuft $R$ alle Elemente von $\grp$, so durchläuft auch $R'$ dieselben, nur in einer anderen Reihenfolge. So erhält man einen Isomorphismus von $\grp$ in sich, und zwar jeden möglichen, indem man die Gruppe $\grp'$ und darin das Element $T$ passend wählt (\emph{Über endliche Gruppen} \S~5, Sitzungsber. 1895, S.~22).

Seien $A$, $B$ und $P$ Elemente von $\grp$, und sei $B=P^{-1}AP$ conjugirt mit $A$ in Bezug auf $\grp$. Dann sind auch $A'$ und $B'=P'^{-1}A'P'$ conjugirt in Bezug auf $\grp$, und umgekehrt. Durch jenen Isomorphismus gehen also die Elemente einer Classe $(\alpha)$ in die Elemente derselben oder einer anderen Classe $(\beta)$ über. Für zwei solche Classen, die ich conjugirte nennen will, ist daher $h_\alpha=h_\beta$. Vereinigt man also jetzt alle Elemente von $\grp$, die einander in Bezug auf $\grp'$ conjugirt sind, zu einer Classe, so ist die Anzahl der neuen Classen $k'\le k$, und jede neue Classe $(\rho)'$ entsteht durch Vereinigung einer gewissen Anzahl $r$ von conjugirten alten Classen $(\alpha)$, $(\beta)$, $(\gamma)$, \ldots. Sie enthält, da $h_\alpha=h_\beta=h_\gamma=\cdots$ ist
\begin{equation*}\tag{2.}
h'_\rho=h_\alpha+h_\beta+h_\gamma+\cdots=rh_\alpha
\end{equation*}
Elemente von $\grp$.

Die $k$ Grössen $\chi_0,\chi_1,\ldots \chi_{k-1}$, die einen Charakter $\chi$ bilden, sind bis auf einen gemeinsamen Factor $f$ dadurch bestimmt, dass
\begin{equation*}\tag{3.}
f\xi=\sum h_\alpha\chi_\alpha x_\alpha
\end{equation*}
ein Linearfactor der Determinante $\Theta$ ist. Ebenso ist jeder neue Charakter $\chi'$ dadurch bestimmt, dass
\[
f'\xi'=\sum h'_\rho\chi'_\rho x'_\rho
\]
ein Linearfactor einer analogen Determinante $\Theta'$ ist. Diese geht aus $\Theta$ hervor, indem man für je $r$ conjugirte Classen $(\alpha)$, $(\beta)$, $(\gamma)$, \ldots, die sich zu einer Classe $(\rho)'$ vereinigen, $x_\alpha=x_\beta=x_\gamma=\cdots=x'_\rho$ setzt. Mithin ist
\[
\frac{1}{f'}h'_\rho\chi'_\rho=\frac{1}{f}(h_\alpha\chi_\alpha+h_\beta\chi_\beta+h_\gamma\chi_\gamma+\cdots)
\]
oder nach (2.)
\begin{equation*}\tag{4.}
\frac{f}{f'}\chi'_\rho=\frac{1}{r}(\chi_\alpha+\chi_\beta+\chi_\gamma+\cdots).
\end{equation*}
Ist $R$ ein Element einer jener $r$ Classen, so kann man diese Gleichung auch durch
\begin{equation*}\tag{5.}
\frac{f}{f'}\chi'(R)=\frac{1}{h'}\sum_U'\chi(U^{-1}RU)
\end{equation*}
ersetzen, wo $U$ die sämmtlichen $h'$ Elemente von $\grp'$ durchläuft. Denn

in dieser Summe von $h'$ Gliedern kommt jedes der $r$ Glieder der Summe (4.) vor und jedes gleich oft.

Indem man aber gewisse Complexe der Variabeln von $\Theta$ einander gleich setzt, können auch zwei verschiedene Linearfactoren $\xi$ und $\eta$ von $\Theta$ einander gleich werden. Es können also zwei verschiedene Charaktere $\chi$ und $\psi$ denselben Charakter $\chi'$ erzeugen. Dies tritt in folgendem Falle ein:

Ist $T$ ein bestimmtes Element von $\grp'$, so entspricht nach (1.) jedem Elemente $R$ von $\grp$ ein Element $R'$ von $\grp$. Durchläuft $S$ die Elemente einer Classe $(\alpha)$, so durchläuft auch $S'$ die Elemente einer Classe $(\beta)$. Daher ist nach Gleichung (5.), \S~5
\[
\sum_S\bigl(\chi(E)\chi(A'S')-\chi(A')\chi(S')\bigr)=0.
\]
Setzt man also
\[
\chi(R')=\chi(T^{-1}RT)=\psi(R),
\]
so ist, weil $A'S'=(AS)'$ ist,
\[
\sum_S\bigl(\psi(E)\psi(AS)-\psi(A)\psi(S)\bigr)=0.
\]
Aus dieser Relation, die mit der Gleichung (6.), \S~2 identisch ist, folgt aber, dass $\psi(R)$ ein Charakter ist. Zwei solche Charaktere sollen conjugirt (in Bezug auf $\grp'$) genannt werden. Führt die Substitution $T$ nicht jede der alten Classen in sich selbst über, so kann man den Charakter $\chi$ so wählen, dass $\psi$ davon verschieden ist. Denn wenn etwa $A$ und $A'$ verschiedene Classen repraesentiren, so kann nicht für jeden Werth von $\varkappa$ $\chi^{(\varkappa)}(A)=\chi^{(\varkappa)}(A')$ sein. Wählt man $\varkappa$ so, dass diese Gleichung nicht besteht, und setzt dann $\chi^{(\varkappa)}=\chi$, so ist $\psi(A)=\chi(A')$ von $\chi(A)$ verschieden, also ist $\psi$ ein anderer Charakter als $\chi$.

Zwei conjugirte Charaktere unterscheiden sich nur durch die Anordnung der $k$ Grössen $\chi_0,\chi_1,\ldots \chi_{k-1}$. Nun ist aber nach Formel (10.), \S~5
\begin{equation*}\tag{6.}
\sum_R \chi(R^{-1})\chi(R)=\frac{hf}{e},
\end{equation*}
also ist auch
\[
\sum_R \psi(R^{-1})\psi(R)=\sum_{R'}\chi(R'^{-1})\chi(R')=\frac{hf}{e},
\]
weil $R'$ zugleich mit $R$ alle Elemente von $\grp$ durchläuft. Für zwei conjugirte Charaktere hat daher, weil $f=\chi(E)=\psi(E)$ ist, auch die Zahl $g=ef$, der Exponent des Linearfactors $\xi$ in $\Theta$, denselben Werth.

Zwei conjugirte Charaktere $\chi$ und $\psi$ erzeugen denselben neuen Charakter $\chi'$. Denn es ist
\begin{equation*}\tag{7.}
\frac{h'f}{f'}\chi'(R)=\sum_U'\chi(U^{-1}RU)=\sum_U'\psi(U^{-1}RU).
\end{equation*}

Die letzte Summe ist nämlich gleich $\sum_U'\chi((UT)^{-1}R(UT))$, unterscheidet sich also von der ersten nur durch die Reihenfolge der Summanden, da $UT$ zugleich mit $U$ alle Elemente von $\grp'$ durchläuft.

Umgekehrt müssen zwei Charaktere $\chi$ und $\psi$, die denselben Charakter $\chi'$ erzeugen, stets conjugirt sein. Es muss also, wenn die Gleichung (7.) besteht, ein solches Element $T$ in $\grp'$ geben, dass $\psi(R)=\chi(T^{-1}RT)$ ist für jedes Element $R$ von $\grp$. Denn aus jener Gleichung folgt
\[
\sum_U'\sum_R \chi(U^{-1}RU)\psi(R^{-1})=\sum_U'\sum_R \psi(U^{-1}RU)\psi(R^{-1}).
\]
Sei $U$ ein bestimmtes Element von $\grp'$. Ist dann $\psi(U^{-1}RU)=\psi(R)$ für jedes Element $R$ von $\grp$, so ist nach (6.) $\sum_R\psi(U^{-1}RU)\psi(R^{-1})$ eine positive von Null verschiedene Grösse. Dieser Bedingung genügt jedes Element $U$ von $\grp'$, das in $\grp$ enthalten ist. Ist sie aber nicht erfüllt, so ist $\psi(U^{-1}RU)=\vartheta(R)$ ein von $\psi(R)$ verschiedener Charakter, und mithin verschwindet diese Summe nach Gleichung (10.), \S~5. Die rechte Seite der letzten Gleichung hat also einen von Null verschiedenen Werth, und mithin auch die Summe auf der linken Seite. Nun ist aber, falls wieder $U$ ein bestimmtes Element von $\grp'$ und $R$ ein veränderliches Element von $\grp$ ist, $\chi(U^{-1}RU)=\vartheta(R)$ ein Charakter. Ist dieser von $\psi(R)$ verschieden, so ist $\sum_R\chi(U^{-1}RU)\psi(R^{-1})=0$. Mithin kann $\vartheta$ nicht für jedes $U$ von $\psi$ verschieden sein, es muss also ein Element $U=T$ geben, für das $\chi(T^{-1}RT)=\psi(R)$ ist, und folglich sind $\chi$ und $\psi$ conjugirte Charaktere.

In der Summe (5.) finden sich alle mit $\chi$ conjugirten Charaktere und jeder gleich oft. Denn diejenigen Elemente $U$ von $\grp'$, die der Gleichung $\chi(U^{-1}RU)=\chi(R)$ für jedes Element $R$ von $\grp$ genügen, bilden eine in $\grp'$ enthaltene Gruppe, die $\grp$ enthält. Daher ist die Anzahl $s$ der Charaktere $\chi^{(\varkappa)},\chi^{(\lambda)},\chi^{(\mu)},\ldots$ die mit $\chi$ conjugirt sind, ein Divisor von $\frac{h'}{h}$, und man kann jene Gleichung auch auf die Form
\begin{equation*}\tag{8.}
\frac{f}{f'}\chi'=\frac{1}{s}(\chi^{(\varkappa)}+\chi^{(\lambda)}+\chi^{(\mu)}+\cdots)
\end{equation*}
bringen. Die Formeln (4.) und (8.) stellen die beiden verschiedenen Arten vor Augen, wie man sich die Entstehung der neuen Charaktere, die ich die relativen Charaktere von $\grp$ in Bezug auf $\grp'$ nennen will, aus den alten denken kann.

Wenn beim Übergange von $\Theta$ zu $\Theta'$ mehrere verschiedene lineare Factoren $\xi,\eta,\ldots$ von $\Theta$ einander gleich werden, so entsprechen sie conjugirten Charakteren $\chi,\psi,\ldots$ und umgekehrt. Daher haben die Potenzen von $\xi,\eta,\ldots$, die in $\Theta$ aufgehen, alle denselben Exponenten $g$.

Der Exponent $g'$ der Potenz, in welcher der entsprechende Linearfactor $\xi'$ in $\Theta'$ enthalten ist, ist folglich
\begin{equation*}\tag{9.}
g'=sg.
\end{equation*}
Daher kann man etwa $f'=sf$, $e'=e$ setzen. Während nun, wie oben erwähnt, $g$ ein Quadrat ist, braucht $g'$ nicht ein Quadrat zu sein. Daraus erklärt es sich, dass der Beweis für jenen Satz mit den bisher benutzten Mitteln allein nicht geführt werden kann.

Die relativen Charaktere von $\grp$ haben für die Gruppe $\grp'$ selbst eine einfache Bedeutung. Die neuen Classen, in die ich die Elemente von $\grp$ getheilt habe, sind auch Classen von $\grp'$, weil $\grp$ eine invariante Untergruppe von $\grp'$ ist. Da ich jetzt von den alten Classen und der Function $\Theta$ keinen Gebrauch mehr mache, will ich die Bezeichnung durchgehend ändern, und das, was ich bisher $k'$ und $\Theta'$ genannt habe, mit $k$ und $\Theta$ bezeichnen, die neuen Classen selbst mit $(0),(1)\ldots(k-1)$, die ihnen entsprechenden Variabeln mit $x_0,x_1,\ldots x_{k-1}$. Dagegen sei jetzt $k'$ die Anzahl der Classen von $\grp'$. Dann enthalten die Classen $(k),(k+1),\ldots(k'-1)$ kein Element von $\grp$. Die der Gruppe $\grp'$ entsprechende Determinante (14.), \S~5 des Grades $h'$ sei
\[
\Theta'=\left|x_{P,Q}\right|=\left|x_{PQ^{-1}}\right|,
\]
wo $P$ und $Q$ die $h'=nh$ Elemente von $\grp'$ in derselben Reihenfolge durchlaufen. Zerfällt $\grp'$ (mod. $\grp$) in die $n$ Complexe
\[
\grp'=\grp+\grp T+\grp U+\cdots,
\]
so ordne ich die Zeilen und Spalten jener Determinante so, dass $P$ zuerst die $h$ Elemente von $\grp$ durchläuft, dann die $h$ Elemente von $\grp T$, dann die von $\grp U$, u. s. w. Setzt man nun in $\Theta'$ $x_k=\cdots=x_{k'-1}=0$, so bleiben nur die $k$ Variabeln $x_R$ übrig, deren Index $R$ ein Element von $\grp$ ist.

In den ersten $h$ Zeilen bleiben nur die Elemente der ersten $h$ Spalten, und diese bilden die Determinante $h^{\text{ten}}$ Grades
\[
\Theta=\left|x_{R,S}\right|=\left|x_{RS^{-1}}\right|,
\]
wo $R$ und $S$ die Elemente von $\grp$ durchlaufen. In den folgenden $h$ Zeilen bleiben nur die Elemente $x_{RT,ST}=x_{RT(ST)^{-1}}=x_{RS^{-1}}$, die in den Spalten $h+1$ bis $2h$ stehen, und diese bilden wieder die Determinante $\Theta$, u. s. w. Daher wird, wenn man $x_k=\cdots=x_{k'-1}=0$ setzt,
\begin{equation*}\tag{10.}
\Theta'=\Theta^n.
\end{equation*}
Nun sind aber die Coefficienten der linearen Factoren von $\Theta$ die relativen Charaktere $\chi_0,\chi_1,\ldots \chi_{k-1}$ von $\grp$, und die der linearen Factoren von $\Theta'$ die Charaktere von $\grp'$. In jedem Charakter $\chi_0,\ldots \chi_{k-1},\chi_k,\ldots \chi_{k'-1}$ von $\grp'$ bilden also die ersten $k$ Werthe $\chi_0,\ldots \chi_{k-1}$ (von einem gemeinsamen

Factor $f$ abgesehen) einen relativen Charakter von $\grp$, und umgekehrt lässt sich jeder relative Charakter von $\grp$, $\chi_0,\ldots \chi_{k-1}$, auf eine oder mehrere Arten durch Hinzufügung passender Werthe $\chi_k,\ldots \chi_{k'-1}$ zu einem Charakter von $\grp'$ ergänzen.

\section*{\S~8.}

Ich will nun die Theorie der Gruppencharaktere an einigen Beispielen erläutern. Die geraden Permutationen von 4 Symbolen bilden eine Gruppe $\grp$ der Ordnung $h=12$. Ihre Elemente zerfallen in 4 Classen, die Elemente der Ordnung 2 bilden eine zweiseitige Classe (1), die der Ordnung 3 zwei inverse Classen (2) und $(3)=(2')$. Sei $\rho$ eine primitive cubische Wurzel der Einheit.

\begin{center}
Tetraeder. $h=12$.
\[
\begin{array}{c|rrrr|r}
 & \chi^{(0)} & \chi^{(1)} & \chi^{(2)} & \chi^{(3)} & h_\alpha\\ \hline
\chi_0 & 1 & 3 & 1 & 1 & 1\\
\chi_1 & 1 & -1 & 1 & 1 & 3\\
\chi_2 & 1 & 0 & \rho & \rho^2 & 4\\
\chi_3 & 1 & 0 & \rho^2 & \rho & 4
\end{array}
\]
\end{center}

Die Werthe von $\chi_0$ sind zugleich die von $f=e$.

Alle Permutationen von 4 Symbolen bilden eine Gruppe $\grp'$ der Ordnung $h'=24$, von der $\grp$ eine invariante Untergruppe ist. Die Classen (2) und (3) sind conjugirt und vereinigen sich zu einer Classe $(2)'$. Ebenso sind die Charaktere $\chi^{(2)}$ und $\chi^{(3)}$ conjugirt. Die relativen Charaktere sind
\[
\begin{array}{c|rrr|r}
 & \chi^{(0)} & \chi^{(1)} & \chi^{(2)} & h_\alpha\\ \hline
\chi_0 & 1 & 3 & 2 & 1\\
\chi_1 & 1 & -1 & 2 & 3\\
\chi_2 & 1 & 0 & -1 & 8
\end{array}
\]
Um diese Tabelle zu erhalten, addire man nach Formel (8), \S~7 in der vorigen die beiden letzten Spalten, und behalte von den beiden letzten Zeilen, die dann einander gleich werden, nur die eine bei. Oder man nehme nach Formel (4), \S~7 aus den Elementen der beiden letzten Zeilen das arithmetische Mittel, und behalte von den beiden letzten Spalten nur die eine bei. Um ganze Zahlen zu erhalten, multiplicire man noch die Elemente der letzten Spalte mit $f'=2$. Für die beiden ersten Charaktere ist $e=f=\chi_0$; für den letzten aber $e=1$, $f=2$, also ist $g=2$ kein Quadrat.

In der Gruppe $\grp'$ bilden die geraden Permutationen der Ordnung 2 eine Classe (1), die ungeraden eine Classe (2), die Permuta-

tionen der Ordnungen 3 und 4 je eine Classe (3) und (4). Jede der 5 Classen ist eine zweiseitige. Die Classen (2) und (4) enthalten die ungeraden Permutationen, die Classen (0), (1) und (3) enthalten die geraden und bilden zusammen die Gruppe $\grp$.

\begin{center}
Octaeder. $h=24$.
\[
\begin{array}{c|rrrrr|r}
 & \chi^{(0)} & \chi^{(1)} & \chi^{(2)} & \chi^{(3)} & \chi^{(4)} & h_\alpha\\ \hline
\chi_0 & 1 & 3 & 3 & 2 & 1 & 1\\
\chi_1 & 1 & -1 & -1 & 2 & 1 & 3\\
\chi_2 & 1 & 1 & -1 & 0 & -1 & 6\\
\chi_3 & 1 & 0 & 0 & -1 & 1 & 8\\
\chi_4 & 1 & -1 & 1 & 0 & -1 & 6
\end{array}
\]
\end{center}

Für jeden Charakter ist $e=f=\chi_0$.

2. Die Permutationen von 5 Symbolen bilden eine Gruppe $\grp'$ der Ordnung $h'=120$, die geraden allein eine Gruppe $\grp$ der Ordnung $h=60$. Die Elemente von $\grp'$ zerfallen in 7 Classen. Die Elemente der Ordnung 3, 4, 5, 6 bilden je eine Classe (3), (4), (5), (6). Von den Permutationen der Ordnung 2 bilden die geraden die Classe (1), die ungeraden die Classe (2). Die Classen (2), (4) und (6) enthalten die ungeraden Permutationen, die übrigen enthalten die geraden und bilden zusammen die Untergruppe $\grp$. Jede der 7 Classen ist eine zweiseitige. Betrachtet man $\grp$ für sich, so zerfällt die Classe (5) von $\grp'$ in 2 Classen. Daher sollen die 5 Classen von $\grp$ mit (0), (1), (3), (4), (5) bezeichnet werden, wo (4) und (5) die Elemente der Ordnung 5 enthalten. Sowohl in $\grp$, wie in $\grp'$ ist jede Classe eine zweiseitige.

\begin{center}
Icosaeder. $h=60$.
\[
\begin{array}{c|rrrrr|r}
 & \chi^{(0)} & \chi^{(1)} & \chi^{(3)} & \chi^{(4)} & \chi^{(5)} & h_\alpha\\ \hline
\chi_0 & 1 & 5 & 4 & 3 & 3 & 1\\
\chi_1 & 1 & 1 & 0 & -1 & -1 & 15\\
\chi_3 & 1 & -1 & 1 & 0 & 0 & 20\\
\chi_4 & 1 & 0 & -1 & \frac{1}{2}(1+\sqrt{5}) & \frac{1}{2}(1-\sqrt{5}) & 12\\
\chi_5 & 1 & 0 & -1 & \frac{1}{2}(1-\sqrt{5}) & \frac{1}{2}(1+\sqrt{5}) & 12
\end{array}
\]
\end{center}

Hier ist stets $e=f=\chi_0$. Die relativen Charaktere in Bezug auf $\grp'$ aber sind
\[
\begin{array}{c|rrrr|r}
 & \chi^{(0)} & \chi^{(1)} & \chi^{(3)} & \chi^{(5)} & h_\alpha\\ \hline
\chi_0 & 1 & 5 & 4 & 6 & 1\\
\chi_1 & 1 & 1 & 0 & -2 & 15\\
\chi_3 & 1 & -1 & 1 & 0 & 20\\
\chi_5 & 1 & 0 & -1 & 1 & 24
\end{array}
\]
Hier ist $e^{(5)}=3$, $f^{(5)}=6$, also $g^{(5)}=18$ kein Quadrat. Endlich ergiebt sich für die Gruppe $\grp'$:

\begin{center}
$h=120$
\[
\begin{array}{c|rrrrrrr|r}
 & \chi^{(0)} & \chi^{(1)} & \chi^{(2)} & \chi^{(3)} & \chi^{(4)} & \chi^{(5)} & \chi^{(6)} & h_\alpha\\ \hline
\chi_0 & 1 & 5 & 5 & 4 & 4 & 6 & 1 & 1\\
\chi_1 & 1 & 1 & 1 & 0 & 0 & -2 & 1 & 15\\
\chi_2 & 1 & 1 & -1 & 2 & -2 & 0 & -1 & 10\\
\chi_3 & 1 & -1 & -1 & 1 & 1 & 0 & 1 & 20\\
\chi_4 & 1 & -1 & 1 & 0 & 0 & 0 & -1 & 30\\
\chi_5 & 1 & 0 & 0 & -1 & -1 & 1 & 1 & 24\\
\chi_6 & 1 & 1 & -1 & -1 & 1 & 0 & -1 & 20
\end{array}
\]
\end{center}

Auch für die symmetrische Gruppe $n^{\text{ten}}$ Grades $\grp$ sind alle Charaktere ganze rationale Zahlen. Der Grund davon liegt darin, dass $R$ und $R^a$ stets in $\grp$ conjugirt sind, wenn $a$ zur Ordnung von $R$ theilerfremd ist.

\section*{\S~9.}

Um noch ein allgemeineres Beispiel durchzuführen, wähle ich die Gruppe $\grp$ der Ordnung
\begin{equation*}\tag{1.}
h=\frac{1}{2}p(p^2-1),
\end{equation*}
die von den linearen Substitutionen
\begin{equation*}\tag{2.}
y=\frac{\gamma+\delta x}{\alpha+\beta x}\qquad(\operatorname{mod}.p)
\end{equation*}
der Determinante
\begin{equation*}\tag{3.}
\alpha\delta-\beta\gamma\equiv 1\qquad(\operatorname{mod}.p)
\end{equation*}
gebildet wird, falls $p$ eine ungerade Primzahl bedeutet. Ihre Eigenschaften, die hier in Betracht kommen, sind am ausführlichsten von \textsc{Gierster}, \emph{Die Untergruppen der \textsc{Galois}'schen Gruppe der Modulargleichungen für den Fall eines primzahligen Transformationsgrades}, Math. Ann. Bd.~18, S.~319 abgeleitet. Dieser interessanten und wichtigen Arbeit entnehme ich die folgenden Resultate:

Die $h$ Elemente der Gruppe $\grp$ zerfallen in
\begin{equation*}\tag{4.}
k=\frac{1}{2}(p-1)+3
\end{equation*}
Classen conjugirter Elemente. Die von dem Elemente $E$ gebildete Hauptclasse möge nicht, wie in der allgemeinen Theorie, mit (0), sondern mit $(\lambda)$ bezeichnet werden. Die Elemente der Ordnung $p$ bilden zwei Classen, die ich auch nicht mit Ziffern, sondern mit den Buchstaben $(\mu)$ und $(\nu)$ bezeichnen werde. Seien $P$ und $Q$ zwei nicht conjugirte Elemente der Ordnung $p$. Ist $a$ ein quadratischer Rest von $p$ und $b$ ein Nichtrest, so ist $P^a$ mit $P$ conjugirt, $P^b$ nicht. Ist daher $p\equiv3\pmod{4}$, so sind $(\mu)$ und $(\nu)$ inverse Classen, und man kann $Q=P^{-1}$ wählen. Ist aber $p\equiv1\pmod{4}$, so sind $P$ und $P^{-1}$ conjugirt, ebenso $Q$ und $Q^{-1}$, also ist jede dieser beiden Classen eine zweiseitige.

Die Ordnung jedes anderen Elementes der Gruppe ist ein Divisor von $\frac12(p-1)$ oder von $\frac12(p+1)$. Es giebt ein Element $R$ der Ordnung $\frac12(p-1)$ und ein Element $S$ der Ordnung $\frac12(p+1)$. Unter den Potenzen von $R$ sind je zwei mit entgegengesetzten Exponenten $R^a$ und $R^{-a}$ conjugirt, ebenso unter den Potenzen von $S$. Dagegen repraesentiren, je nachdem $p\equiv1$ oder $3\pmod{4}$ ist,
\begin{equation*}\tag{5.}
\begin{array}{c}
R,\ R^2,\ldots,\ R^{\frac{p-1}{4}}\quad\left(\text{bez. }R^{\frac{p-3}{4}}\right),\\[0.35em]
S,\ S^2,\ldots,\ S^{\frac{p-1}{4}}\quad\left(\text{bez. }S^{\frac{p+1}{4}}\right),
\end{array}
\end{equation*}
$\frac12(p-1)$ verschiedene Classen. Fügt man dazu noch die drei Elemente $E$, $P$ und $Q$, so hat man für jede der $k$ Classen einen und nur einen Repraesentanten. Jede Classe ist eine zweiseitige. Nur wenn $p\equiv3\pmod{4}$ ist, sind $(\mu)$ und $(\nu)$ inverse Classen. Nachdem $R$ und $S$ in irgend einer bestimmten Weise gewählt sind, möge die Zahl $\pm a$ $\left(\operatorname{mod}.\frac12(p-1)\right)$ der Index der durch $R^a$ oder $R^{-a}$ repraesentirten Classe genannt werden, und ebenso die Zahl $\pm b$ $\left(\operatorname{mod}.\frac12(p+1)\right)$ der Index der durch $S^b$ oder $S^{-b}$ repraesentirten Classe.

Für je zwei conjugirte Substitutionen (2.) hat die Zahl
\begin{equation*}\tag{6.}
\varkappa\equiv \frac12(\alpha+\delta)\qquad(\operatorname{mod}.p)
\end{equation*}
abgesehen vom Vorzeichen denselben Werth, also auch die beiden Wurzeln jeder der beiden Congruenzen
\begin{equation*}\tag{7.}
x^2\pm2\varkappa x+1\equiv0\qquad(\operatorname{mod}.p).
\end{equation*}
Daher nenne ich $\pm\varkappa$ die Invariante der Classe, der die Substitution (2.) angehört. Haben umgekehrt zwei Substitutionen dieselbe Invariante $\pm\varkappa$, so sind sie conjugirt, ausser wenn $\varkappa=\pm1$ ist. Die Substitutionen der Invariante $\varkappa=\pm1$ zerfallen in drei Classen nach dem quadratischen Charakter $(\operatorname{mod}.p)$ der Zahlen, die durch die Form
\begin{equation*}\tag{8.}
\frac12(\alpha+\delta)\bigl(\beta x^2+(\alpha-\delta)xy-\gamma y^2\bigr)
\end{equation*}
dargestellt werden. Sind diese alle durch $p$ theilbar, so ist $\beta=\gamma=0$, $\alpha=\delta=\pm1$, und wir erhalten die Hauptclasse $(\lambda)$. Sind sie nicht alle durch $p$ theilbar, so sind die Zahlen, die durch die Form darstellbar und zu $p$ theilerfremd sind, entweder alle Reste oder alle Nichtreste von $p$. Im ersten Falle ist $\varkappa\beta$ Rest, und falls $\beta=0$ ist, $-\varkappa\gamma$ Rest. Diese Classe möge $(\mu)$ heissen. Im zweiten Falle ist $\varkappa\beta$ Nichtrest, und falls $\beta=0$ ist, $-\varkappa\gamma$ Nichtrest. Diese Classe möge $(\nu)$ heissen. Ist aber $\varkappa$ von $\pm1$ verschieden, so giebt es nur eine Classe der Invariante $\pm\varkappa$, sie möge mit $(\varkappa)$ oder auch mit $(+\varkappa)$ bezeichnet werden $\left(\varkappa=0,2,3,\ldots,\frac{p-1}{2}\right)$. Setzt man
\begin{equation*}\tag{9.}
\varepsilon_\varkappa=\left(\frac{\varkappa^2-1}{p}\right),\qquad
\varepsilon_0=\varepsilon=\left(\frac{-1}{p}\right),\qquad
\varepsilon_1=0,
\end{equation*}
so ergeben sich für die Anzahl der Elemente jeder dieser $k$ Classen die Formeln
\begin{equation*}\tag{10.}
h_\lambda=1,\qquad
h_\mu=h_\nu=\frac12(p^2-1),\qquad
h_0=\frac12p(p+\varepsilon),\qquad
h_\varkappa=p(p+\varepsilon_\varkappa).
\end{equation*}

Die Elemente der Classen $(\mu)$ und $(\nu)$ haben alle die Ordnung $p$. Die Classe $(0)$ besteht aus allen Elementen von $\grp$, deren Ordnung gleich $2$ ist. Ihr etwas abweichendes Verhalten gegen die anderen Classen $(\varkappa)$ erklärt sich daraus: Ist $A$ ein Element der Classe $(\varkappa)$, so sind $A$ und $A^{-1}$ zwar conjugirt aber verschieden. Ist aber $\varkappa=0$, so sind diese beiden Elemente gleich. Hat $R^a$ nicht die Ordnung $2$ (oder $1$), so sind die $\frac12(p-1)$ Potenzen von $R$ die einzigen Elemente von $\grp$, die mit $R^a$ vertauschbar sind. Mit der Gruppe dieser Potenzen ist aber noch ein Element $T$ der Ordnung $2$ vertauschbar, das der Bedingung $T^{-1}RT=R^{-1}$ genügt. Das Analoge gilt von $S^b$. Hat aber $R^a$ (bez. $S^b$) die Ordnung $2$, so ist ausser den Potenzen von $R$ auch noch $T$ mit $R^a$ vertauschbar.

Ist $\varkappa^2-1$ quadratischer Rest von $p$, also $\varepsilon_\varkappa=+1$, so ist die Ordnung eines Elements der Classe $(\varkappa)$ der Exponent, zu dem $(\varkappa+\sqrt{\varkappa^2-1})^2$ $(\operatorname{mod}.p)$ gehört, also ein Divisor von $\frac12(p-1)$. Ist aber $\varkappa^2-1$ Nichtrest, also $\varepsilon_\varkappa=-1$, so ist diese Ordnung der Exponent, zu dem $x^2\ (\operatorname{mod}.p,\ x^2-2\varkappa x+1)$ gehört, also ein Divisor von $\frac12(p+1)$, demnach in beiden Fällen ein Divisor von $\frac12(p-\varepsilon_\varkappa)$. Umgekehrt ist daher für die Classe, der $R^a$ angehört, $\varepsilon_\varkappa=+1$, und für die Classe, der $S^b$ angehört, $\varepsilon_\varkappa=-1$. Nach (5.) zerfallen die Elemente, deren Ordnung aufgeht in
\begin{equation*}\tag{11.}
\frac12(p-1),\quad \text{in}\quad \frac14(p-\varepsilon)-\frac12(1-\varepsilon)
\end{equation*}
Classen, und die, deren Ordnung aufgeht in
\begin{equation*}\tag{12.}
\frac12(p+1),\quad \text{in}\quad \frac14(p-\varepsilon)
\end{equation*}
Classen. Dabei ist das Hauptelement ausgeschlossen. Oder anders zusammengefasst: Es zerfallen die Elemente, deren Ordnung aufgeht in
\begin{equation*}\tag{13.}
\frac12(p-\varepsilon),\quad \text{in}\quad \frac14(p-\varepsilon)
\end{equation*}
Classen, und die, deren Ordnung aufgeht in
\begin{equation*}\tag{14.}
\frac12(p+\varepsilon),\quad \text{in}\quad \frac14(p-\varepsilon)-\frac12(1-\varepsilon)
\end{equation*}
Classen. Ist $\pm\varkappa$ die Invariante einer Classe, so will ich ihren Index jetzt so definiren: Sind $r$ und $s$ primitive Wurzeln der Congruenzen
\begin{equation*}\tag{15.}
r^{p-1}\equiv1,\qquad s^{p+1}\equiv1\qquad(\operatorname{mod}.p),
\end{equation*}
so ist $r$ reell, und es ist zwar $s$ imaginär, aber $s^b+s^{-b}$ reell. Ist nun $\varepsilon_\alpha=+1$, so hat jede der beiden Congruenzen (7.) zwei reelle Wurzeln, $r^{\pm a}$, $r^{\frac12(p-1)\pm a}$. Dann nenne ich die nach dem Modul $\frac12(p-1)$ genommene Zahl $\pm a$ den Index der Classe $(\pm\alpha)$. Ist aber $\varepsilon_\beta=-1$, so hat jede zwei imaginäre Wurzeln $s^{\pm b}$, $s^{\frac12(p+1)\pm b}$. Dann nenne ich die nach dem Modul $\frac12(p+1)$ genommene Zahl $\pm b$ den Index der Classe $(\pm\beta)$. Diese Definition geht in die obige über, wenn man für $R$ die Substitution $y=\frac{rx}{r^{-1}}$, und für $S$ eine Substitution wählt, die in imaginärer Form $y=\frac{sx}{s^{-1}}$ lautet. (\textsc{Gierster}, a. a. O. \S~3.) Denn sind $r$ und $r^{-1}$ die Wurzeln der charakteristischen Gleichung (7.) für die Substitution $R$, so sind $r^a$ und $r^{-a}$ die für die Substitution $R^a$.

\section*{\S~10.}

Drei (verschiedene oder gleiche) Classen mögen \emph{concordant} heissen, wenn zwischen ihren Invarianten $\alpha$, $\beta$, $\gamma$ die Beziehung
\begin{equation*}\tag{1.}
\alpha^2+\beta^2+\gamma^2\pm2\alpha\beta\gamma\equiv1\qquad(\operatorname{mod}.p)
\end{equation*}
besteht, sonst \emph{discordant}. Ist z. B. $\gamma=\pm1$, so reducirt sich diese Relation auf $\beta=\pm\alpha$. Ist $\gamma=0$, so lautet sie $\alpha^2+\beta^2=1$. Ich schliesse nun den Fall aus, wo eine der Invarianten gleich $\pm1$ ist. Schreibt man die Gleichung (1.) in der Form
\[
(\alpha^2-1)(\beta^2-1)\equiv(\alpha\beta\pm\gamma)^2,
\]
so folgt daraus
\begin{equation*}\tag{2.}
\varepsilon_\alpha=\varepsilon_\beta=\varepsilon_\gamma.
\end{equation*}
Sind $\alpha$ und $\beta$ gegeben, so sind damit concordant
\begin{equation*}\tag{3.}
\pm\gamma\equiv\alpha\beta+\sqrt{(\alpha^2-1)(\beta^2-1)},\qquad
\pm\delta\equiv\alpha\beta-\sqrt{(\alpha^2-1)(\beta^2-1)},
\end{equation*}
und es ist $\varepsilon_\alpha=\varepsilon_\beta=\varepsilon_\gamma=\varepsilon_\delta$. Ist z. B. $\beta=\pm\alpha$, so ist $\pm\gamma=2\alpha^2-1$ und $\pm\delta=1$, und mithin ist stets $\varepsilon_\alpha=\varepsilon_{(2\alpha^2-1)}$. Die Ordnungen der Elemente von drei concordanten Classen gehen alle in $\frac12(p-1)$ auf oder alle in $\frac12(p+1)$. Sind also $a$, $b$, $c$ ihre Indices, so beziehen sich diese alle auf dieselbe primitive Wurzel $r$ (oder $s$). Ist
\[
2\alpha\equiv r^a+r^{-a},\qquad
2\beta\equiv r^b+r^{-b},\qquad
2\gamma\equiv r^c+r^{-c},
\]
so geht die Congruenz (1.) (für das obere Vorzeichen) über in
\[
(r^{a+b+c}+1)(r^{a-b-c}+1)(r^{-a+b-c}+1)(r^{-a-b+c}+1)\equiv0,
\]
ist also identisch mit der Bedingung
\begin{equation*}\tag{4.}
a\pm b\pm c\equiv0\qquad\left(\operatorname{mod}.\frac12(p-1),\ \text{bez. }\frac12(p+1)\right).
\end{equation*}
Die Indices der beiden Classen (3.) sind folglich
\begin{equation*}\tag{5.}
c\equiv a+b,\qquad d\equiv a-b.
\end{equation*}
Ich setze ferner zur Abkürzung
\begin{equation*}\tag{6.}
2\eta_\varkappa=
\left(\frac{-2}{p}\right)
\left(\left(\frac{1+\varkappa}{p}\right)+\left(\frac{1-\varkappa}{p}\right)\right),
\qquad
\eta_0=\left(\frac{-2}{p}\right),\qquad
2\eta_1=\varepsilon.
\end{equation*}
Ist $\varepsilon_\beta=-\varepsilon$, so ist $\eta_\beta=0$. Ist aber $\varepsilon_\alpha=+\varepsilon$, so ist
\begin{equation*}\tag{7.}
\eta_\alpha=
\left(\frac{-2}{p}\right)\left(\frac{1+\alpha}{p}\right)
=
\left(\frac{-2}{p}\right)\left(\frac{1-\alpha}{p}\right)
=(-1)^a\varepsilon.
\end{equation*}
Denn ist $\varepsilon=+1$ und $2\alpha\equiv r^a+r^{-a}$, so ist $2(\alpha+1)\equiv(r^{-a}+1)^2r^a$, also
\[
\left(\frac{2}{p}\right)\left(\frac{1+\alpha}{p}\right)=\left(\frac{r^a}{p}\right)=(-1)^a.
\]
Ist aber $\varepsilon=-1$ und $2\alpha\equiv s^a+s^{-a}$, so ist $2\alpha+2\equiv s^a+s^{-a}+2$. Ist also $a=2b$ gerade, so ist $2\alpha+2\equiv(s^b+s^{-b})^2$, also quadratischer Rest, da $s^b+s^{-b}$ reell ist. Ist umgekehrt $2\alpha+2=4\beta^2$ quadratischer Rest, so ist $2\alpha-2=4(\beta^2-1)$ Nichtrest, da $\varepsilon_\alpha=-1$, also $\alpha^2-1=(\alpha-1)(\alpha+1)$ Nichtrest ist. Sind dann $s^b$ und $s^{-b}$ die Wurzeln der Congruenz $x^2-2\beta x+1\equiv0$, so ist $2\beta\equiv s^b+s^{-b}$, also $s^a+s^{-a}+2\equiv(s^b+s^{-b})^2$ oder $s^a+s^{-a}\equiv s^{2b}+s^{-2b}$. Mithin ist $a\equiv\pm2b$ $(\operatorname{mod}.p+1)$, also $a$ gerade.

Nun seien $\alpha$, $\beta$, $\gamma$ Zahlen von $2$ bis $\frac12(p-1)$. Dann ist
\[
h_{000}=\frac12h,\qquad
h_{00\alpha}=h,\qquad
h_{0\alpha\beta}=2h,\qquad
h_{\alpha\beta\gamma}=4h,
\]
falls die drei Classen discordant sind. Sind sie aber concordant, so ist
\[
h_{0\alpha\beta}=2h+\varepsilon p^2(p+\varepsilon),\qquad
h_{\alpha\beta\gamma}=4h+\varepsilon_\alpha p^2(p+\varepsilon_\alpha).
\]
Ferner ist
\[
\begin{array}{lll}
h_{\lambda\lambda\lambda}=1, &
h_{\lambda\mu\mu}=h_{\lambda\nu\nu}=h_\mu\frac12(1+\varepsilon), &
h_{\lambda\mu\nu}=h_\mu\frac12(1-\varepsilon),\\[0.35em]
h_{00\lambda}=\frac12p(p+\varepsilon), &
h_{\alpha\alpha\lambda}=p(p+\varepsilon_\alpha), &
h_{\alpha\beta\lambda}=0,\\[0.35em]
h_{\mu\mu\mu}=h_{\nu\nu\nu}=h_\mu\left(\frac34(p-\varepsilon)+\frac12(p-1)\right),&
\multicolumn{2}{l}{
h_{\mu\mu\nu}=h_{\mu\nu\nu}=h_\mu\left(\frac14(p-\varepsilon)-\frac12(1-\varepsilon)\right),}\\[0.35em]
h_{00\mu}=h_{00\nu}=h\frac12(1+\varepsilon),&
h_{0\alpha\mu}=h_{0\alpha\nu}=h,&
h_{\alpha\beta\mu}=h_{\alpha\beta\nu}=2h,\\[0.35em]
h_{\alpha\alpha\mu}=h_{\alpha\alpha\nu}=2h+\varepsilon_\alpha h,&&\\[-0.2em]
h_{0\mu\mu}=h_{0\nu\nu}=h\frac12(1+\varepsilon\eta),&
h_{0\mu\nu}=h\frac12(1-\varepsilon\eta),&\\[0.35em]
h_{\alpha\mu\mu}=h_{\alpha\nu\nu}=h(1+\varepsilon\eta_\alpha),&
h_{\alpha\mu\nu}=h(1-\varepsilon\eta_\alpha).&
\end{array}
\]
Den Weg, auf dem ich diese Zahlen berechnet habe, will ich nur kurz andeuten: Um $\frac{h_{0\alpha\beta}}{h_0}$ zu erhalten, nehme man eine bestimmte Substitution der Classe $(0)$ und setze sie mit allen Substitutionen der Classe $(\alpha)$ zusammen,
\[
\begin{pmatrix}0&1\\-1&0\end{pmatrix}
\begin{pmatrix}\xi&\eta\\ \zeta&2\alpha-\xi\end{pmatrix}
=
\begin{pmatrix}\zeta&2\alpha-\xi\\ -\xi&-\eta\end{pmatrix}.
\]
Damit die zusammengesetzte Substitution der Classe $(\beta)'=(\beta)$ angehöre, muss sein
\[
\zeta-\eta=\pm2\beta,\qquad
\xi(2\alpha-\xi)-\eta\zeta\equiv1,
\]
also
\[
(\xi-\alpha)^2+(\eta\pm\beta)^2\equiv\alpha^2+\beta^2-1.
\]
Ist $\beta$ von Null verschieden, so ist die Anzahl der Lösungen dieser beiden Congruenzen doppelt so gross, wie die der Congruenz $\xi^2+\eta^2\equiv\omega$, wo $\omega\equiv\alpha^2+\beta^2-1$ ist. Ist $\omega$ von Null verschieden, so ist diese Anzahl $p-\varepsilon$. Ist aber $\omega=0$, sind also $(0)$, $(\alpha)$, $(\beta)$ concordante Classen, so ist sie $p-\varepsilon+\varepsilon p$. Da nun $h_0=\frac12p(p+\varepsilon)$ ist, so ist $h_{0\alpha\beta}=2h$, ausser wenn $\alpha^2+\beta^2=1$ ist, dann aber $h_{0\alpha\beta}=2h+\varepsilon p^2(p+\varepsilon)$.

Die Gruppe $\grp$ kann zu einer Gruppe der Ordnung $2h$ erweitert werden, worin $P$ und $Q$ conjugirte Elemente sind. Sie hat also einen Isomorphismus in sich, durch den die Classen $(\mu)$ und $(\nu)$ in einander übergeführt werden. Daher ergeben sich für die Classe $(\mu)$ dieselben Beziehungen, wie für $(\nu)$. Z. B. ist $h_{\mu\nu\nu}=h_{\nu\mu\mu}$. Nach den Formeln (10.), \S~9 ist $h_{\lambda\mu\nu}=h_\mu\frac12(1-\varepsilon)$. Um nun $\frac{1}{h_\mu}(h_{\mu\nu\mu}+h_{\mu\nu\nu}+h_{\mu\nu\lambda})$ zu finden, nehme man eine bestimmte Substitution der Classe $(\mu)$, setze sie mit allen Substitutionen von $(\nu)$ zusammen
\[
\begin{pmatrix}0&1\\-1&2\end{pmatrix}
\begin{pmatrix}\xi&\eta\\ \zeta&2-\xi\end{pmatrix}
\]
und setze die Invariante der zusammengesetzten Substitution gleich $\pm1$, so ergeben sich die Gleichungen
\[
\zeta-\eta+4-2\xi=\pm2,\qquad
2\xi-\xi^2-\eta\zeta\equiv1.
\]
Nimmt man das untere Zeichen, so erhält man durch Elimination von $\xi$
\[
(\eta+\zeta+4)^2\equiv16\eta.
\]
Diese Congruenz hat keine zulässige Lösung. Denn $\eta$ soll Nichtrest sein, und falls $\eta=0$ ist, soll $-2$ Nichtrest sein. Nimmt man das obere Zeichen, so ergiebt sich $(\eta+\zeta)^2\equiv0$, also $\eta=-\zeta$. Daher kann für $\eta$ jeder Nichtrest gesetzt werden, und man erhält $\frac12(p-1)$ Lösungen. Mithin ist $2h_{\mu\mu\nu}+h_\mu\frac12(1-\varepsilon)=h_\mu\frac12(p-1)$, also
\[
h_{\mu\mu\nu}=h_\mu\left(\frac14(p-\varepsilon)-\frac12(1-\varepsilon)\right).
\]

Zur Bestimmung der Charaktere der Gruppe $\grp$ hat man demnach die folgenden Gleichungen zu lösen. Ich setze zur Abkürzung
\[
x=\chi_0+2\sum\chi_\varkappa+\chi_\mu+\chi_\nu,\qquad
y=\eta_0\chi_0+2\sum\eta_\varkappa\chi_\varkappa+\eta_1(\chi_\mu+\chi_\nu),
\]
wo $\varkappa=2,3,\ldots,\frac12(p-1)$ ist. Seien $\alpha$ und $\beta$ Zahlen der Reihe $0,2,3,\ldots,\frac12(p-1)$.

\[
\begin{array}{ll}
1)&\text{Ist }\alpha\text{ von }\beta\text{ verschieden und }\varepsilon_\beta=-\varepsilon_\alpha,\\[0.25em]
& p\chi_\alpha\chi_\beta=fx.\\[0.65em]
2)&\text{Ist aber }\varepsilon_\beta=+\varepsilon_\alpha\\[0.25em]
& \displaystyle \frac{p(p+\varepsilon_\alpha)}{f}\chi_\alpha\chi_\beta
=x(p-\varepsilon_\alpha)+\varepsilon_\alpha p(\chi_\gamma+\chi_\delta),
\end{array}
\]
wo $\gamma$ und $\delta$ durch (3.) bestimmt sind.
\[
\begin{array}{ll}
3)&\displaystyle \frac{(p+\varepsilon_\alpha)}{f}\chi_\alpha\chi_\mu
=x+\varepsilon_\alpha\chi_\alpha+\eta_\alpha(\chi_\mu-\chi_\nu),\\[0.75em]
&\displaystyle \frac{(p+\varepsilon_\alpha)}{f}\chi_\alpha\chi_\nu
=x+\varepsilon_\alpha\chi_\alpha+\eta_\alpha(\chi_\nu-\chi_\mu).\\[0.75em]
4)&\displaystyle \frac{p(p+\varepsilon)}{f}\chi_0^2
=x(p-\varepsilon)+\varepsilon(p-\varepsilon)(\chi_\mu+\chi_\nu)+2f,\\[0.75em]
&\displaystyle \frac{p(p+\varepsilon_\alpha)}{f}\chi_\alpha^2
=x(p-\varepsilon_\alpha)+\frac12\varepsilon_\alpha(p-\varepsilon_\alpha)(\chi_\mu+\chi_\nu)+f+\varepsilon_\alpha p\chi_{(2\alpha^2-1)},\\[0.75em]
5)&\displaystyle \frac{p^2-1}{pf}\chi_\mu^2
=x+\varepsilon y+\frac{f}{p}(1+\varepsilon)-\frac1p\left(1+\frac{\varepsilon}{2}\right)(\chi_\mu+\chi_\nu)+\frac{\varepsilon(p-\varepsilon)}{p}(\chi_\mu-\chi_\nu),\\[0.75em]
&\displaystyle \frac{p^2-1}{pf}\chi_\nu^2
=x+\varepsilon y+\frac{f}{p}(1+\varepsilon)-\frac1p\left(1+\frac{\varepsilon}{2}\right)(\chi_\mu+\chi_\nu)+\frac{\varepsilon(p-\varepsilon)}{p}(\chi_\nu-\chi_\mu),\\[0.75em]
&\displaystyle \frac{p^2-1}{pf}\chi_\mu\chi_\nu
=x-\varepsilon y+\frac{f}{p}(1-\varepsilon)-\frac1p\left(1-\frac{\varepsilon}{2}\right)(\chi_\mu+\chi_\nu).
\end{array}
\]

I. Ich untersuche zuerst, ob es Lösungen giebt, bei denen $\chi_\mu$ und $\chi_\nu$ verschieden sind. Aus 3) folgt $(p+\varepsilon_\alpha)\chi_\alpha=2f\eta_\alpha$. Ist $\varepsilon_\alpha=-\varepsilon$, so ist $\eta_\alpha=0$, also auch $\chi_\alpha=0$. Setzt man also den vorläufig willkürlichen Proportionalitätsfactor $f=\frac12(p+\varepsilon)$, so ist
\[
\chi_\alpha=\eta_\alpha\qquad\left(\alpha=0,2,3,\ldots,\frac12(p-1)\right).
\]
Nach 1) ist $x=0$. Nach Gleichung (2.) ist $\varepsilon_\alpha=\varepsilon_{(2\alpha^2-1)}$. Ist also in 4) $\varepsilon_\alpha=-\varepsilon$, so ergiebt sich $\chi_\mu+\chi_\nu=\varepsilon$. Nach 5) ist dann $y=\frac12(p-\varepsilon-1)$ und $4\chi_\mu\chi_\nu=1-\varepsilon p$, mithin
\[
\chi_\mu=\frac12(\varepsilon\pm\sqrt{\varepsilon p}),\qquad
\chi_\nu=\frac12(\varepsilon\mp\sqrt{\varepsilon p}).
\]
Dass die beiden erhaltenen Lösungen wirklich allen Gleichungen genügen, z. B. den Gleichungen 2)
\begin{equation*}\tag{8.}
2\eta_\alpha\eta_\beta=\varepsilon_\alpha(\eta_\gamma+\eta_\delta)
\end{equation*}
bestätigt man am einfachsten mittels der Formel (7.), \S~9 und (5.).

II. Für jede andere Lösung ist $\chi_\mu=\chi_\nu$. Dadurch vereinfachen sich die Gleichungen beträchtlich, z. B. ist 5) zu ersetzen durch
\[
5^*)\qquad
\frac{p^2-1}{f}\chi_\mu^2=px+f-2\chi_\mu,\qquad
py=\chi_\mu-f.
\]

Ich untersuche nun, ob es Lösungen giebt, bei denen $x$ von Null verschieden ist. Dann sind nach 1) $\chi_0$, $\chi_2$, $\chi_3,\ldots,\chi_{\frac12(p-1)}$ von Null verschieden, und zwar sind alle $\chi_\alpha$ einander gleich, für welche $\varepsilon_\alpha$ denselben Werth hat. Sei $\varepsilon_\alpha=+\varepsilon$ und $\varepsilon_\delta=-\varepsilon$. Dann ist
\[
y=\eta_0\chi_0+2\chi_0\sum\eta_\alpha+\varepsilon\chi_\mu,
\]
also weil $\eta_0+2\sum\eta_\alpha=-\varepsilon$ ist,
\[
y=\varepsilon(\chi_\mu-\chi_0)=\frac{\chi_\mu-f}{p}.
\]
Ferner ist in der Summe $x=\chi_0+2\chi_\mu+2\sum(\chi_\alpha+\chi_\delta)$ die Anzahl der Charaktere $\chi_\alpha$ gleich $\frac14(p-\varepsilon)-1$, die der Charaktere $\chi_\delta$ aber gleich $\frac14(p-\varepsilon)-\frac12(1-\varepsilon)$. Daher ist
\[
x=2\chi_\mu+\left(\frac12(p-\varepsilon)-1\right)\chi_0+
\left(\frac12(p+\varepsilon)-1\right)\chi_\delta.
\]
Nach 1) und 4) ist
\[
fx=p\chi_0\chi_\delta,\qquad
\frac{p(p-\varepsilon)}{f}\chi_\delta^2=x(p+\varepsilon)-2\varepsilon p\chi_\delta,
\]
also
\[
(p-\varepsilon)\chi_\delta=(p+\varepsilon)\chi_0-2\varepsilon f,\qquad
px=(p^2+1)\chi_\mu-f.
\]
Setzt man diese Werthe in $5^*)$ ein, so erhält man
\[
\chi_\mu(\chi_\mu-f)=0,
\]
also entweder $\chi_\mu=f$ oder $\chi_\mu=0$. Im ersten Falle sei $f=1$. Dann ist $\chi_\mu=\chi_\nu=1$, $\chi_0=1$, $\chi_\delta=1$, $\chi_\alpha=1$ $(x=p,\ y=0)$. Im zweiten Falle sei $f=p$. Dann ist $x=-1$, $y=-1$, und $\chi_0=\varepsilon$, $\chi_\delta=-\varepsilon$, $\chi_\alpha=\varepsilon$ oder allgemein $\chi_\varkappa=\varepsilon_\varkappa$, auch $\chi_\mu=\chi_\nu=\varepsilon_1=0$.

III. Für alle anderen Lösungen ist $x=0$. Nach 1) sind daher entweder alle $\chi_\alpha=0$, für die $\varepsilon_\alpha=-1$ ist oder alle, für die $\varepsilon_\alpha=+1$ ist. Sei zuerst $\chi_\alpha=0$, falls $\varepsilon_\alpha=-1$ ist, und $f=p+1$. Nach 3) ist, falls $\varepsilon_\alpha=+1$ ist, $\chi_\alpha\chi_\mu=\chi_\alpha$, und da nicht alle $\chi_\alpha$ verschwinden können, $\chi_\mu=+1$. Ist $\varepsilon_\alpha=\varepsilon_\beta=+1$, also auch $\varepsilon_\gamma=\varepsilon_\delta=+1$, so ist nach 2) und 4)
\[
\chi_\alpha\chi_\beta=\chi_\gamma+\chi_\delta,\qquad
\chi_\alpha^2=\chi_{(2\alpha^2-1)}+2.
\]
Setzt man also, falls $a$ und $b$ die Indices der Classen $(\alpha)$ und $(\beta)$ sind, $\chi_\alpha=\xi_a$, $\chi_\beta=\xi_b$, und $\xi_0=2$, so ist nach (5.)
\[
\xi_a\xi_b=\xi_{a+b}+\xi_{a-b},
\]
auch wenn $b=a$ ist. Ist $\rho$ eine neue Unbekannte, und $\xi_1=\rho+\rho^{-1}$, so ergiebt sich aus $\xi_1\xi_1=\xi_2+\xi_0$, dass $\xi_2=\rho^2+\rho^{-2}$, dann aus $\xi_1\xi_2=\xi_3+\xi_1$, dass $\xi_3=\rho^3+\rho^{-3}$ ist, allgemein, dass $\xi_a=\rho^a+\rho^{-a}$ ist. Aus $\xi_{\frac12(p-1)}=\xi_0$ folgt dann, dass
\begin{equation*}\tag{9.}
\rho^{\frac12(p-1)}=1
\end{equation*}
ist. Man gelangt so zu den Lösungen
\[
\begin{array}{ll}
\chi_\lambda=p+1,\quad \chi_\mu=\chi_\nu=1,\quad
\chi_\alpha=\rho^a+\rho^{-a},&
\text{falls }\varepsilon_\alpha=+1\text{ und }2\alpha\equiv r^a+r^{-a},\\
\chi_\beta=0,&
\text{falls }\varepsilon_\beta=-1\text{ ist.}
\end{array}
\]
Die erhaltenen Werthe genügen auch allen Gleichungen, nur wenn $\rho=1$ ist, nicht der Gleichung $x=0$, und wenn $\frac12(p-1)$ gerade und $\rho=-1$ ist, nicht der Gleichung $y=-1$. $(5^*)$. Ihre Anzahl ist daher
\begin{equation*}\tag{10.}
n=\frac14(p-\varepsilon)-1.
\end{equation*}

IV. Sei endlich $\chi_\alpha=0$, falls $\varepsilon_\alpha=+1$ ist, und $f=p-1$. Dann ist $\chi_\beta\chi_\mu=-\chi_\beta$ und mithin $\chi_\mu=-1$. Ist $\varepsilon_\alpha=\varepsilon_\beta=-1$, so ist
\[
\chi_\alpha\chi_\beta=-\chi_\gamma-\chi_\delta,\qquad
\chi_\alpha^2=-\chi_{(2\alpha^2-1)}+2.
\]
Setzt man dann, falls $a$ und $b$ die Indices der Classen $(\alpha)$ und $(\beta)$ sind, $\chi_\alpha=-\xi_a$, $\chi_\beta=-\xi_b$ und $\xi_0=2$, so ist
\[
\xi_a\xi_b=\xi_{a+b}+\xi_{a-b}.
\]
Ist also $\xi_1=\sigma+\sigma^{-1}$, so ist $\xi_b=\sigma^b+\sigma^{-b}$, und weil $\xi_{\frac12(p+1)}=\xi_0$ ist, so ist
\begin{equation*}\tag{11.}
\sigma^{\frac12(p+1)}=1.
\end{equation*}
So gelangt man zu den Lösungen
\[
\begin{array}{ll}
\chi_\lambda=p-1,\quad \chi_\mu=\chi_\nu=-1,\quad
\chi_\beta=-\sigma^b-\sigma^{-b},&
\text{falls }\varepsilon_\beta=-1\text{ und }2\beta\equiv s^b+s^{-b},\\
\chi_\alpha=0,&
\text{falls }\varepsilon_\alpha=+1\text{ ist.}
\end{array}
\]
Für $\sigma$ ist der Werth $+1$, und wenn $\frac12(p+1)$ gerade ist, auch der Werth $-1$ unzulässig. Die Anzahl dieser Lösungen ist mithin
\begin{equation*}\tag{12.}
n=\frac14(p-\varepsilon)-\frac12(1-\varepsilon).
\end{equation*}
Damit sind die $k=\frac12(p-1)+3$ Charaktere sämmtlich ermittelt. Zu jedem ist noch eine Zahl $e$ mit Hülfe der Formel
\[
\frac{hf}{e}=\sum h_\alpha\chi_\alpha\chi_\alpha'
\]
zu bestimmen. Man findet, dass in allen $k$ Fällen $e=f$ ist. Die Proportionalitätsfactoren sind also der in \S~3 angegebenen Regel entsprechend gewählt.

Ich will die $k$ Charaktere noch einmal zusammenstellen, indem ich mich der in \S~5 eingeführten Bezeichnung bediene:
\[
\begin{array}{c|ccccc}
n & 1 & 1 & 2 & \frac14(p-\varepsilon)-1 & \frac14(p-\varepsilon)-\frac12(1-\varepsilon)\\ \hline
\chi(E) & 1 & p & \frac12(p+\varepsilon) & p+1 & p-1\\
\chi(P) & 1 & 0 & \frac12(\varepsilon\pm\sqrt{\varepsilon p}) & 1 & -1\\
\chi(Q) & 1 & 0 & \frac12(\varepsilon\mp\sqrt{\varepsilon p}) & 1 & -1\\
\chi(R^a) & 1 & 1 & \frac12(1+\varepsilon)(-1)^a & \rho^a+\rho^{-a} & 0\\
\chi(S^b) & 1 & -1 & -\frac12(1-\varepsilon)(-1)^b & 0 & -\sigma^b-\sigma^{-b}
\end{array}
\]
Für $p=3$ und $p=5$ ergeben sich daraus die in \S~8 angegebenen Charaktere der Gruppe des Tetraeders und des Icosaeders.

\begin{center}
\rule{11em}{0.4pt}\\[-0.2em]
Ausgegeben am 27. August.\\[-0.2em]
\rule{11em}{0.4pt}

{\small Berlin, gedruckt in der Reichsdruckerei.}
\end{center}
% ===== END WORK 053 =====

\clearpage

% ===== BEGIN WORK 054: 054_ber_die_Primfactoren_der_Gruppendetermina =====
\begin{center}
{\Large Über die Primfactoren der Gruppendeterminante.}\\[0.65em]
Von G. Frobenius.\\[0.6em]
\rule{0.22\linewidth}{0.4pt}
\end{center}

Die Theorie der Charaktere einer Gruppe, deren Grundlagen ich in meiner letzten Arbeit entwickelt habe, erfordert zu ihrer weiteren Ausgestaltung die Untersuchung einer Determinante, deren Grad der Ordnung der Gruppe gleich ist. Nach dem Vorgange von Dedekind, der zuerst ihre Bedeutung für die Theorie der Gruppen erkannt und meine Aufmerksamkeit auf sie gelenkt hat, nenne ich sie die der Gruppe entsprechende \emph{Gruppendeterminante}. Die $h$ Elemente $A,B,C,\ldots$ der Gruppe $\Hgrp$ benutze ich als Indices für $h$ unabhängige Variable $x_A,x_B,x_C,\ldots$. Indem ich diese Bezeichnung wähle, treffe ich die Festsetzung, dass, wenn $L=MN$ ist, auch $x_L=x_{MN}$ sein soll. Aus diesen $h$ Grössen, die durch einen Index von einander unterschieden sind, bilde ich $h^2$ Grössen, die mit zwei Indices versehen sind, indem ich $x_{P,Q}=x_{PQ^{-1}}$ setze. Sind $G_1,G_2,\ldots,G_h$ die $h$ Elemente von $\Hgrp$ in irgend einer bestimmten Reihenfolge, so betrachte ich die Matrix $(x_{P,Q})=(x_{PQ^{-1}})$, deren $h$ Zeilen man erhält, indem man für $P$ der Reihe nach die $h$ Elemente $G_1,G_2,\ldots,G_h$ setzt, und deren $h$ Spalten man erhält, indem man für $Q$ dieselben $h$ Elemente in derselben Reihenfolge setzt. Diese Matrix besitzt gewisse, durch die Constitution der Gruppe $\Hgrp$ bedingte Symmetrieeigenschaften. In jeder Zeile finden sich die $h$ Variabelen sämmtlich und ebenso in jeder Spalte. Die verschiedenen Zeilen (Spalten) unterscheiden sich von einander nur durch die Anordnung der Variabelen. Die Gruppendeterminante, die der Gruppe $\Hgrp$ entspricht, ist die Determinante dieser Matrix
\[
\Theta=|x_{P,Q}|=|x_{PQ^{-1}}|.
\]

Addirt man zu den Elementen der ersten Zeile die aller anderen Zeilen, so werden jene Elemente alle gleich $\sum x_R=\xi$. Daher ist die ganze Function $h^{\mathrm{ten}}$ Grades $\Theta$ der $h$ Variabeln $x_A,x_B,x_C,\ldots$ durch die lineare Function $\xi$ theilbar. Mithin zerfällt $\Theta$, von dem trivialen Falle $h=1$ abgesehen, stets in Factoren niedrigeren Grades. Die Anzahl $k$ der verschiedenen irreducibelen Factoren oder Primfactoren von $\Theta$ ist gleich der Anzahl der Classen conjugirter Elemente, worin die Elemente von $\Hgrp$ zerfallen. Ist $f$ der Grad eines solchen Primfactors $\Phi$, so ist $\Theta$ durch die $f^{\mathrm{te}}$ und durch keine höhere Potenz von $\Phi$ theilbar. Der Grad $f$ ist ein Divisor der Ordnung $h$. Durch eine lineare Substitution lässt sich $\Phi$ in eine Function von $f^2$, aber nicht von weniger Variabelen transformiren, und wenn man jeden Primfactor von
\[
\Theta=\Pi \Phi^f
\]
in dieser Weise umformt, so sind die $\sum f^2=h$ neuen Variabelen alle unter einander unabhängig. Setzt man immer diejenigen Variabelen $x_R$ einander gleich, deren Indices Elemente derselben Classe sind, so wird
\[
\Phi=\xi^f
\]
die $f^{\mathrm{te}}$ Potenz einer linearen Function $\xi$ von $k$ unabhängigen Variabelen, und die $k$ linearen Functionen, die so aus den $k$ Primfactoren von $\Theta$ entspringen, sind linear unabhängig. Aus den Coefficienten der linearen Function $\xi$ ergiebt sich ein Charakter $\chi$ der Gruppe $\Hgrp$, und aus seinen $k$ Werthen $\chi(R)$ lassen sich die Coefficienten der Primfunction $\Phi$ sämmtlich berechnen. Die Theorie der allgemeinen Gruppendeterminante, worin die $h$ Grössen $x_R$ unbeschränkt veränderlich sind, wird so auf die Theorie der speciellen Gruppendeterminante zurückgeführt, worin die Veränderlichkeit dieser Grössen durch die Bedingungen $x_{BA}=x_{AB}$ beschränkt ist. Die Berechnung dieser Determinante $h^{\mathrm{ten}}$ Grades aber
\[
\Theta=\Pi \xi^{f^2}
\]
lässt sich auf die einer Determinante $k^{\mathrm{ten}}$ Grades
\[
\left|\sum_\gamma \frac{1}{h_\alpha}h_{\alpha\beta\gamma}x_\gamma\right|=\Pi\xi
\]
reduciren, worin der lineare Factor $\xi$, der in $\Theta$ zur Potenz $f^2$ erhoben vorkommt, nur einfach enthalten ist. Die Definition der positiven ganzen Zahlen $h_{\alpha\beta\gamma}$ und die Entwicklung ihrer Beziehungen zu den Charakteren bildet den Hauptinhalt meiner Arbeit \emph{Über Gruppencharaktere}, Sitzungsberichte 1896, die ich im Folgenden mit \emph{Ch.} citiren werde.

Ganz analoge Eigenschaften, wie ein solcher Primfactor $\Phi$ einer Gruppendeterminante, hat die ganze Function $2^p$ten Grades von $2^{2p}$ Variabelen, die ich in meiner Arbeit \emph{Über Thetafunctionen mehrerer Variabelen} (Crelle's Journal Bd. 96) untersucht habe, auf die ich dort aber nicht durch die Betrachtung der Gruppe der zwischen den Thetafunctionen bestehenden Relationen, sondern durch das für diese Functionen geltende Additionstheorem geführt worden bin. Sonst ist die Gruppendeterminante bisher nur für den Fall commutativer Gruppen untersucht worden, wo ihre Primfactoren sämmtlich linear sind. Für einige besonders einfache, nicht commutative Gruppen hat Dedekind im Jahre 1886 die Determinante $\Theta$ durch Rechnung in Primfactoren zerfällt, und seine interessanten Ergebnisse, die er mir vor kurzem mitgetheilt hat, haben mich veranlasst, die Zerlegung der Gruppendeterminante in Primfactoren allgemein für eine beliebig gegebene Gruppe zu untersuchen.

\section*{§ 1.}

Die Determinante der Matrix $h^{\mathrm{ten}}$ Grades
\frobtagc{1.}{(x_{P,Q})=(x_{PQ^{-1}})=(x)}
bezeichne ich mit
\frobtagc{2.}{|x_{P,Q}|=|x_{PQ^{-1}}|=\Theta(x_E,x_A,x_B,x_C,\ldots)=\Theta(x_R)=\Theta(x)=\Theta.}
Unter $E$ verstehe ich immer das Hauptelement. In dem Zeichen $\Theta(x_R)$ bedeutet $R$ ein veränderliches Element, für das die $h$ Elemente $E,A,B,C,\ldots$ der Gruppe $\Hgrp$ zu setzen sind. Bei Anwendung der Bezeichnung $(x)$ oder $\Theta(x)$ ist $x$ ein leeres Zeichen, das erst dadurch eine Bedeutung erhält, dass daran die Indices $E,A,B,C,\ldots$ angehängt werden.

Nun sei $y_E,y_A,y_B,y_C,\ldots$ ein zweites System von $h$ unabhängigen Variabelen. Aus ihnen bilde ich die Matrix
\frobtagc{3.}{(y_{P,Q})=(y_{PQ^{-1}})=(y).}
Ihre Zeilen (Spalten) erhält man, indem man für $P$ ($Q$) die $h$ Elemente $G_1,G_2,\ldots,G_h$ von $\Hgrp$ in derselben Reihenfolge setzt, in der sie bei der Bildung der Matrix (1.) benutzt sind.

Aus jenen beiden Systemen von je $h$ Variabeln $x_R$ und $y_R$ bilde ich ein drittes System $z_R$, indem ich
\frobtagc{4.}{z_A=\sum x_Ry_S\qquad (RS=A)}
setze. In dieser Summe sind für $R$ die $h$ Elemente von $\Hgrp$ zu setzen, und jedes Element $R$ ist mit dem Elemente $S$ $(=R^{-1}A)$ zu verbinden, das der Bedingung $RS=A$ (nicht $SR=A$) genügt, so dass auch $S$ die $h$ Elemente von $\Hgrp$ durchläuft, jedes Mal verbunden mit $R=AS^{-1}$. Dann ergiebt sich durch Zusammensetzung der beiden Matrizen (1.) und (3.), welche die hier vorausgesetzten, durch die Constitution der Gruppe $\Hgrp$ bedingten Symmetrieeigenschaften besitzen, die Matrix
\frobtagc{5.}{(z_{P,Q})=(z_{PQ^{-1}})=(z)=(x)(y)}
mit denselben Symmetrieeigenschaften. Denn es ist
\[
z_{P,Q}=\sum_R x_{P,R}y_{R,Q}=\sum_R x_{PR^{-1}}y_{RQ^{-1}}.
\]
Setzt man in dieser Summe $R=SQ$, so durchläuft $S$ gleichzeitig mit $R$ die $h$ Elemente von $\Hgrp$, nur in einer anderen Reihenfolge, und es wird nach (4.)
\[
z_{P,Q}=\sum_S x_{PQ^{-1}S^{-1}}y_S=z_{PQ^{-1}}.
\]

Derselbe Satz gilt, wenn man beliebig viele derartige Matrizen zusammensetzt. Sind z. B. $z_E,z_A,z_B,z_C,\ldots$ $h$ beliebige Grössen, und setzt man
\[
v_A=\sum x_Ry_Sz_T,\qquad (RST=A)
\]
so ist die Matrix $(v_{PQ^{-1}})=(v)=(x)(y)(z)$ aus den drei Matrizen $(x)$, $(y)$ und $(z)$ in dieser Reihenfolge zusammengesetzt. Setzt man die Matrix $(x)$ $n$ Mal mit sich selbst zusammen, so möge sich die Matrix
\[
(x_{P,Q})^n=(x)^n=(x_{PQ^{-1}}^{(n)})=(x^{(n)})
\]
ergeben. Dann ist
\frobtagc{6.}{x_R^{(n)}=\sum x_{R_1}x_{R_2}\cdots x_{R_n}\qquad (R_1R_2\cdots R_n=R)}

Demnach besitzt auch jede Function der Matrix $(x)$ die hier vorausgesetzten Symmetrieeigenschaften, z. B. die zu $(x)$ adjungirte Matrix (\emph{Über lineare Substitutionen und bilineare Formen}, Crelle's Journal Bd. 84 S. 7.), ebenso die Hauptmatrix (Einheitsmatrix)
\[
(x)^0=(\varepsilon_{P,Q})=(\varepsilon_{PQ^{-1}})=(\varepsilon),
\]
wo $\varepsilon_R=0$ ist, ausser wenn $R=E$ ist, und $\varepsilon_E=1$ ist.

Nun seien $\Phi,\Phi',\Phi'',\ldots$ die verschiedenen in
\frobtagc{7.}{\Theta=\Phi^e\Phi'^{e'}\Phi''^{e''}\cdots=\Pi\Phi^e}
aufgehenden unzerlegbaren Functionen (Primfunctionen), und seien $f,f',f'',\ldots$ die Grade dieser ganzen homogenen Functionen der $h$ Variabeln $x_E,x_A,x_B,x_C,\ldots$. In der Gruppendeterminante $\Theta$ sind die Elemente der Diagonale und nur diese gleich $x_E$. Daher reducirt sich $\Theta$ auf $x_E^h$, wenn man alle Variabelen ausser $x_E$ gleich Null setzt, und folglich reducirt sich dann auch $\Phi$ auf die $f^{\mathrm{te}}$ Potenz von $x_E$. Daher kann man den noch unbestimmten constanten Factor von $\Phi$ so wählen, dass in dieser Function $x_E^f$ den Coefficienten 1 hat.

Nach dem Multiplicationstheorem der Determinanten folgt aus den Gleichungen (2.) und (5.) die Relation
\frobtagc{8.}{\Theta(z)=\Theta(x)\Theta(y).}
Daraus ergiebt sich für jeden Primfactor von $\Theta$
\[
\Phi=\Phi(x_E,x_A,x_B,x_C,\ldots)=\Phi(x_R)=\Phi(x)
\]
die analoge Relation
\frobtagc{9.}{\Phi(z)=\Phi(x)\Phi(y),\quad \text{wenn }(z)=(x)(y)}
ist, durch welche die Function $\Phi$ unabhängig von ihrer Beziehung zur Gruppendeterminante $\Theta$ charakterisirt werden kann. Denn zerlegt man die rechte Seite der Gleichung (8.) in Primfactoren, so folgt daraus, dass $\Phi(z)$ in das Product einer Function $\Lambda(x)$ der $h$ Variabeln $x_R$ allein und einer Function $M(y)$ der $h$ Variabeln $y_R$ allein zerfällt. Setzt man dann in der Gleichung $\Phi(z)=\Lambda(x)M(y)$, $y_R=\varepsilon_R$, so wird $z_R=x_R$, also $\Phi(x)=\Lambda(x)M(\varepsilon)$. Ebenso ist $M(y)=\Lambda(\varepsilon)\Phi(y)$ und $\Lambda(\varepsilon)M(\varepsilon)=\Phi(\varepsilon)=1$.

Umgekehrt muss jede unzerlegbare ganze homogene Function $\Phi$ von $x_E,x_A,x_B,\ldots$, die der Bedingung (9.) genügt, ein Factor der Gruppendeterminante $\Theta(x)$ sein. Denn setzt man in dieser Gleichung für $(y)$ die zu $(x)$ adjungirte Matrix, so wird $z_R=\varepsilon_R\Theta(x)$, also
\[
\Phi(z)=\Theta(x)^f=\Phi(x)\Phi(y),
\]
wo $y_A$ eine ganze Function der $h$ Variabeln $x_R$ ist. Daher muss die Function $\Phi(x)$, weil sie unzerlegbar ist, ein Factor von $\Theta(x)$ sein.

Mit Hülfe der Relation (9.) lassen sich alle Eigenschaften der Determinanten, die aus dem Multiplicationstheorem fliessen, auf die Primfactoren der Gruppendeterminante übertragen, namentlich die Eigenschaften, welche ich in meiner im Folgenden mit \emph{V.} citirten Arbeit \emph{Über vertauschbare Matrizen} (S. 601 dieses Bandes) entwickelt habe.

\section*{§ 2.}

Jeder Primfactor $\Phi(x)$ der Gruppendeterminante genügt der Bedingung
\frobtagc{1.}{\Phi(z)=\Phi(x)\Phi(y),}
falls
\frobtagc{2.}{z_C=\sum x_Ay_B\qquad (AB=C)}
gesetzt wird. Mit Hülfe dieser Beziehung lassen sich die linearen Factoren
\frobtagc{3.}{\Phi(x)=\sum\chi(A)x_A}
vollständig bestimmen. Denn aus der Gleichung
\[
\left(\sum\chi(A)x_A\right)\left(\sum\chi(B)y_B\right)=\left(\sum\chi(C)z_C\right)=\sum\chi(AB)x_Ay_B
\]
ergiebt sich für die Coefficienten $\chi(A)$ dieser Functionen die Relation
\frobtagc{4.}{\chi(AB)=\chi(A)\chi(B).}
Mithin ist $\chi(E)=1$, $\chi(A)\chi(A^{-1})=1$ und allgemeiner
\[
\chi(ABCD\cdots)=\chi(A)\chi(B)\chi(C)\chi(D)\cdots,
\]
und folglich
\[
\begin{aligned}
\chi(B^{-1}A^{-1}BA)
&=\chi(B^{-1})\chi(A^{-1})\chi(B)\chi(A)\\
&=1.
\end{aligned}
\]
Das Element $F$, das sich mittels der Gleichung
\frobtagc{5.}{BA=ABF}
aus $A$ und $B$ ergiebt, nenne ich nach Dedekind den \emph{Commutator} von $A$ und $B$. Demnach ist $\chi(F)=1$ für jeden Commutator $F$ von irgend zwei Elementen der Gruppe $\Hgrp$. Ist $T$ ein beliebiges Element von $\Hgrp$, und ist
\[
T^{-1}AT=A',\qquad T^{-1}BT=B',\qquad T^{-1}FT=F',
\]
so ist auch $B'A'=A'B'F'$. Ist also $F$ ein Commutator, so ist auch jedes mit $F$ conjugirte Element $F'$ ein solcher. Theilt man die Elemente von $\Hgrp$ in Classen conjugirter Elemente, so werden die Commutatoren von den sämmtlichen Elementen einiger dieser Classen gebildet. Die von ihnen erzeugte Gruppe $\Ggrp$ ist daher eine invariante Untergruppe von $\Hgrp$. (Sie kann auch gleich $\Hgrp$ sein oder auch gleich der Hauptgruppe $\Egrp$, das letztere stets und nur dann, wenn $\Hgrp$ eine commutative Gruppe ist.) Ist $G$ ein Element von $\Ggrp$, so giebt es solche Commutatoren $F,F',F'',\ldots$ (die nicht verschieden zu sein brauchen), dass $G=FF'F''\cdots$ ist. Daher ist $\chi(G)=\chi(F)\chi(F')\chi(F'')\cdots=1$.

Nun sei
\[
\Hgrp=\Ggrp A+\Ggrp B+\Ggrp C+\cdots,
\]
seien also $A,B,C,\ldots$ die (mod. $\Ggrp$) verschiedenen Elemente von $\Hgrp$. Ihre Anzahl ist $\frac{h}{g}$, wenn $g$ die Ordnung von $\Ggrp$ ist. Die $\frac{h}{g}$ Complexe $\Ggrp A=A\Ggrp$, $\Ggrp B=B\Ggrp$, $\ldots$ bilden eine Gruppe, die mit $\frac{\Hgrp}{\Ggrp}$ bezeichnet wird. Ist $\Ggrp$ die Commutatorgruppe, so ist $\frac{\Hgrp}{\Ggrp}$ eine commutative (Abel'sche) Gruppe, und damit $\frac{\Hgrp}{\Ggrp}$ eine commutative Gruppe sei, ist nothwendig und hinreichend, dass $\Ggrp$ durch die Commutatorgruppe theilbar ist. Denn sind $\Ggrp A$ und $\Ggrp B$ zwei Elemente von $\frac{\Hgrp}{\Ggrp}$, so giebt es in $\Ggrp$ ein solches Element $F$, dass $BA=ABF$ ist, also auch $\Ggrp BA=\Ggrp ABF$. Nun ist $\Ggrp ABF=(\Ggrp A)(\Ggrp B)(\Ggrp F)$ und $\Ggrp F=\Ggrp$, also
\[
(\Ggrp B)(\Ggrp A)=(\Ggrp A)(\Ggrp B).
\]
Diese Eigenschaften der Commutatorgruppe hat Dedekind im Jahre 1880 gefunden. Veröffentlicht aber sind sie zuerst von Miller, \emph{The regular substitution groups whose order is less than 48}. Quarterly Journal of Math. 1896, vol. 28, p. 266.

Ist $G$ irgend ein Element von $\Ggrp$, so ist $\chi(GA)=\chi(G)\chi(A)=\chi(A)$. Daher hat $\chi(R)$ für alle Elemente $R$ des Complexes $\Ggrp A$ denselben Werth. Mithin kann man auch die Zahl $\chi(A)$ dem Complexe $\Ggrp A$ zuordnen. Da diese Complexe eine commutative Gruppe bilden, und die ihnen zugeordneten Zahlen $\chi(A)$ die Eigenschaft (4.) haben, so bilden die Zahlen $\chi(A),\chi(B),\chi(C),\ldots$ einen Charakter der commutativen Gruppe $\frac{\Hgrp}{\Ggrp}$. Für eine solche giebt es bekanntlich immer $\frac{h}{g}$ verschiedene Charaktere, deren Werthe sämmtlich Einheitswurzeln sind. Ist $\psi$ einer derselben, und setzt man für jedes in dem Complex $\Ggrp A$ enthaltene Element $R$ $\chi(R)=\psi(\Ggrp A)$, so ist für jedes Element $G$ der Gruppe $\Ggrp$ $\chi(G)=1$, und es gilt für je zwei Elemente von $\Hgrp$ die Gleichung (4.). Ferner ist dann die Function (3.), deren Coefficienten diese Werthe $\chi(A)$ sind, ein linearer Factor der Gruppendeterminante $\Theta$. Denn setzt man $y_R=\chi(R)x_R$, so ist
\[
y_{PQ^{-1}}=\chi(PQ^{-1})x_{PQ^{-1}}=\chi(P)\chi(Q^{-1})x_{PQ^{-1}},
\]
und mithin ist $|y_{PQ^{-1}}|=|x_{PQ^{-1}}|$. Diese Determinante enthält aber den Factor $\sum y_R=\sum\chi(R)x_R=\Phi$, und zwar nur in der ersten Potenz. Denn addirt man die Elemente aller Zeilen zu denen der ersten Zeile, so werden dieselben alle gleich $\sum y_R=\Phi$, und wenn man dann den Factor $\Phi$ aufhebt, alle gleich 1. Zieht man nun die Elemente der ersten Spalte von denen der folgenden ab, so erkennt man, dass $\Theta:\Phi$ nur von den Differenzen $y_R-y_A$ abhängt. Mithin kann dieser Quotient nicht noch einmal durch die Summe $\sum y_R$ theilbar sein.

Folglich ist die Anzahl der linearen Factoren der Gruppendeterminante gleich dem Quotienten aus der Ordnung der Gruppe und der Ordnung ihrer Commutatorgruppe, und jeder lineare Factor ist nur in der ersten Potenz in $\Theta$ enthalten.

Diesen Satz hat Dedekind durch Induction gefunden. Einen Linearfactor, nämlich $\sum x_R$, giebt es immer. Der entsprechende Charakter, $\chi(R)=1$ für jedes Element $R$, heisst der \emph{Hauptcharakter}. Ist $\Ggrp=\Hgrp$, so giebt es keinen anderen Linearfactor. Dies muss stets eintreten, wenn $\Hgrp$ eine einfache Gruppe ist, deren Ordnung eine zusammengesetzte Zahl ist.

\section*{§ 3.}

Man wähle jetzt eine beliebige ganze Zahl $f\le h$ und versuche eine ganze Function $f^{\mathrm{ten}}$ Grades $\Phi$ der $h$ Variabeln $x_E,x_A,x_B,\ldots$ zu bilden, die der Bedingung (9.) § 1 genügt. In dieser muss der Coefficient von $x_E^f$ gleich 1 sein. Denn setzt man $y_R=\varepsilon_R$, so wird $z_R=x_R$, also $\Phi(x)=\Phi(x)\Phi(\varepsilon)$, und mithin ist $\Phi(\varepsilon)=1$. Ich bezeichne nun, wenn $R$ von $E$ verschieden ist, den Coefficienten von $x_E^{f-1}x_R$ in $\Phi(x)$ mit $\chi(R)$, setze aber
\frobtagc{1.}{\chi(E)=f.}
Für jedes Element $R$ ist dann $\chi(R)$ der Coefficient von $x_E^{f-1}$ in $\frac{\partial\Phi}{\partial x_R}$. Ist $u$ eine Variable, so setze ich zur Abkürzung
\[
\Phi(x+u\varepsilon)=\Phi(x_E+u,x_A,x_B,x_C,\ldots).
\]
Das Zeichen $u$, das eine veränderliche Grösse bezeichnet, ist hier scharf zu unterscheiden von den leeren Zeichen $x$ und $\varepsilon$, die erst durch Anhängen der Indices $E,A,B,C,\ldots$ eine Bedeutung erhalten.

Ist nun
\frobtagc{2.}{\Phi(x+u\varepsilon)=u^f+\Phi_1u^{f-1}+\Phi_2u^{f-2}+\cdots+\Phi_f,}
so ist $\Phi_n$ eine ganze homogene Function $n^{\mathrm{ten}}$ Grades der $h$ Variabeln $x_E,x_A,x_B,x_C,\ldots$, und zwar ist $\Phi_f=\Phi$ und
\frobtagc{3.}{\Phi_n=\frac{1}{(f-n)!}\frac{\partial^{f-n}\Phi}{\partial x_E^{f-n}},\qquad \frac{\partial\Phi_{n+1}}{\partial x_E}=(f-n)\Phi_n.}
Endlich ist
\frobtagc{4.}{\Phi_1=\sum\chi(R)x_R.}
Die $h$ Constanten $\chi(R)$ betrachte ich als die Werthe einer Function $\chi$. Ist $\Phi$ unzerlegbar, so nenne ich $\chi$ den Charakter $f^{\mathrm{ten}}$ Grades der Gruppe $\Hgrp$, welcher der Primfunction $f^{\mathrm{ten}}$ Grades $\Phi$ entspricht.

Ist speciell $\sum\psi(R)x_R$ ein linearer Factor von $\Theta$, so heisst $\psi(R)$ ein Charakter ersten Grades von $\Hgrp$. Dafür ist, wie oben gezeigt, die Relation
\[
\psi(A)\psi(B)=\psi(AB)
\]
die nothwendige und hinreichende Bedingung. Setzt man $\psi(R)x_R=y_R$, so ist $\Theta(y)=\Theta(x)$. Ist also $\Phi(x)$ ein Primfactor $f^{\mathrm{ten}}$ Grades von $\Theta(x)$, so ist auch $\Phi(y)$ ein solcher. In diesem ist der Coefficient von $x_E^{f-1}x_R$ gleich $\chi(R)\psi(R)$. Es gilt also der Satz:

\emph{Ist $\chi(R)$ ein Charakter $f^{\mathrm{ten}}$ Grades und $\psi(R)$ ein Charakter ersten Grades von $\Hgrp$, so ist auch $\chi(R)\psi(R)$ ein Charakter $f^{\mathrm{ten}}$ Grades von $\Hgrp$.}

Dieser neue Charakter ist gleich $\chi(R)$, wenn $\psi(R)$ der Hauptcharakter ist. Er braucht aber auch in anderen Fällen nicht von $\chi(R)$ verschieden zu sein, nämlich, wenn $\chi(R)=0$ ist für jedes solche Element $R$, wofür $\psi(R)$ von 1 verschieden ist.

Seien $u_1,u_2,\ldots,u_f$ die $f$ Werthe von $-u$, wofür
\frobtagc{5.}{\Phi(x+u\varepsilon)=(u+u_1)(u+u_2)\cdots(u+u_f)}
verschwindet. Seien $a,v_1,v_2,\ldots,v_n$ Constanten und
\[
g(u)=a(u+v_1)(u+v_2)\cdots(u+v_n)
\]
eine ganze Function von $u$. Dann ist auch, falls man die Variable $u$ durch die Matrix $(x)$ ersetzt,
\[
(y)=g((x))=a((x)+v_1(\varepsilon))((x)+v_2(\varepsilon))\cdots((x)+v_n(\varepsilon)).
\]
Mithin findet man durch wiederholte Anwendung der Relation (9.) § 1
\[
\Phi(y)=a^f\Phi(x+v_1\varepsilon)\Phi(x+v_2\varepsilon)\cdots\Phi(x+v_n\varepsilon).
\]
Setzt man hier
\[
\Phi(x+v\varepsilon)=(u_1+v)(u_2+v)\cdots(u_f+v),
\]
so ergiebt sich der Satz: Ist die Matrix $(y)=g((x))$ eine ganze Function der Matrix $(x)$, so ist
\[
\Phi(y)=g(u_1)g(u_2)\cdots g(u_f).
\]
Ersetzt man hier $g(u)$ durch $g(u)+v$, wo $v$ ein Parameter ist, behält aber die Abkürzung $(y)$ für $g((x))$ bei, so ergiebt sich die Gleichung
\[
\Phi(y+v\varepsilon)=(v+g(u_1))(v+g(u_2))\cdots(v+g(u_f)).
\]
Ist z. B. $g(u)=u^n$, so ist
\[
\Phi(x^{(n)}+v\varepsilon)=(v+u_1^n)(v+u_2^n)\cdots(v+u_f^n).
\]
Durch Vergleichung der Coefficienten von $v^{f-1}$ erhält man daraus nach (4.)
\frobtagc{6.}{S_n=\sum_R \chi(R)x_R^{(n)},}
wo $S_n$ die Summe der $n^{\mathrm{ten}}$ Potenzen der $f$ Grössen $u_1,u_2,\ldots,u_f$ ist. Nach Formel (6.) § 1 kann man dafür auch schreiben
\frobtagc{7.}{S_n=\sum_{R_1,R_2,\ldots,R_n}\chi(R_1R_2\cdots R_n)x_{R_1}x_{R_2}\cdots x_{R_n},}
wo jeder der $n$ Summationsbuchstaben $R_1,R_2,\ldots,R_n$ unabhängig von den anderen die $h$ Elemente von $\Hgrp$ durchläuft.

Aus den Potenzsummen $S_n$ kann man aber die Coefficienten $\Phi_n$ der Function (2.) berechnen nach der Formel
\frobtag{8.}{(-1)^n\Phi_n=\sum\frac{(-1)^{a+b+c+\cdots}S_1^aS_2^bS_3^c\cdots}{1^a2^b3^c\cdots a!b!c!\cdots},\qquad (a+2b+3c+\cdots=n)}
wo $a,b,c,\ldots$ alle ganzen Zahlen $(\ge0)$ durchlaufen, die der Bedingung $a+2b+3c+\cdots=n$ genügen. Diese Formel gilt auch, wenn $n>f$ ist, falls dann $\Phi_n=0$ gesetzt wird. Mit ihrer Hülfe werden wir die Functionen $\Phi_2,\Phi_3,\ldots$ und besonders $\Phi_f=\Phi$ darstellen. Die Coefficienten von $\Phi$ sind ganze Functionen der $h$ Constanten $\chi(R)$, deren Coefficienten rationale Zahlen sind. Wählt man $n>f$, so ergeben sich aus jener Formel Relationen, denen die $h$ Grössen $\chi(R)$ genügen müssen. Ehe ich aber zu diesen Rechnungen übergehe, muss ich eine wichtige Eigenschaft der Function $\chi(R)$ vorausschicken.

Setzt man in der Gleichung (9.) § 1 die Variabelen $y_R=0$ ausser einer, $y_A$, so erhält man
\frobtagc{9.}{\Phi(x_{RA^{-1}})=\zchar(A)\Phi(x_R),}
wo $\zchar(A)$ den Coefficienten von $x_A^f$ in $\Phi(x_R)$ bezeichnet. Setzt man aber die Variabelen $x_R=0$ ausser einer, $x_A$ (und ersetzt dann den Buchstaben $y$ durch $x$), so erhält man
\frobtagc{10.}{\Phi(x_{A^{-1}R})=\zchar(A)\Phi(x_R).}
Ersetzt man hier, falls $B$ ein festes Element ist, für jeden Index $R$ die Variable $x_R$ durch $x_{B^{-1}R}$, so findet man
\[
\Phi(x_{B^{-1}A^{-1}R})=\zchar(A)\Phi(x_{B^{-1}R})=\zchar(A)\zchar(B)\Phi(x_R).
\]
Ersetzt man dagegen in der Gleichung (10.) $A$ durch $AB$, so wird
\[
\Phi(x_{B^{-1}A^{-1}R})=\zchar(AB)\Phi(x_R).
\]
Mithin ist $\zchar(AB)=\zchar(A)\zchar(B)$, und demnach ist $\zchar(R)$ ein Charakter ersten Grades der Gruppe $\Hgrp$, eine Einheitswurzel. Ist $A$ ein festes Element, und setzt man für jedes $R$ $y_R=x_{RA^{-1}}$, so ist
\[
y_{P,Q}=y_{PQ^{-1}}=x_{PQ^{-1}A^{-1}}=x_{P,AQ}.
\]
Wenn man aber in der Determinante $|x_{P,Q}|$ $Q$ durch $AQ$ ersetzt, so wird dadurch nur die Reihenfolge der $h$ Spalten geändert. Ist $m$ die Ordnung des Elementes $A$, so besteht jene Permutation der $h$ Spalten aus lauter Cyklen der Ordnung $m$. Die Anzahl dieser Cyklen ist mithin $\frac{h}{m}$, und folglich ist die Permutation eine gerade oder ungerade, je nachdem $h-\frac{h}{m}$ gerade oder ungerade ist. Daher ist
\[
\Theta(x_{RA^{-1}})=(-1)^{h-\frac{h}{m}}\Theta(x_R),
\]
also
\frobtagc{11.}{\Pi(\zchar A)^e=(-1)^{h-\frac{h}{m}}.}
Ersetzt man in der Gleichung (10.) $x_R$ durch $x_{RA}$, so erhält man
\[
\Phi(x_{A^{-1}RA})=\zchar(A)\Phi(x_{RA})=\zchar(A)\zchar(A^{-1})\Phi(x_R),
\]
also
\frobtagc{12.}{\Phi(x_{A^{-1}RA})=\Phi(x_R).}
Die Function $\Phi(x_R)$ bleibt also ungeändert, wenn man für jeden Index $R$ die Variable $x_R$ durch $x_{A^{-1}RA}$ ersetzt, wo $A$ ein festes Element von $\Hgrp$ ist. Dabei bleibt die Variable $x_E$ ungeändert. Durch Vergleichung der Coefficienten von $x_E^{f-1}x_R$ ergiebt sich aber aus (12.) $\chi(ARA^{-1})=\chi(R)$ oder, wenn man $R=BA$ setzt,
\frobtagc{13.}{\chi(AB)=\chi(BA).}
Theilt man also die $h$ Elemente von $\Hgrp$ in Classen conjugirter Elemente, so hat $\chi(R)$ für alle Elemente $R$ einer Classe denselben Werth.

Nunmehr wende ich mich zur Berechnung der Functionen $\Phi_n$ mit Hülfe der Formel (8.). Zunächst ist nach (6.)
\[
\begin{aligned}
\Phi_1&=S_1=\sum\chi(A)x_A,\\
2\Phi_2&=S_1^2-S_2=\sum(\chi(A)\chi(B)-\chi(AB))x_Ax_B,\\
6\Phi_3&=S_1^3-3S_1S_2+2S_3\\
&=\sum\bigl(\chi(A)\chi(B)\chi(C)-\chi(A)\chi(BC)-\chi(B)\chi(AC)\\
&\qquad-\chi(C)\chi(AB)+\chi(ABC)+\chi(ACB)\bigr)x_Ax_Bx_C.
\end{aligned}
\]
Ich setze daher
\frobtag{14.}{\begin{aligned}
\chi(A,B)&=\chi(A)\chi(B)-\chi(AB),\\
\chi(A,B,C)&=\chi(A)\chi(B)\chi(C)-\chi(A)\chi(BC)-\chi(B)\chi(AC)\\
&\qquad-\chi(C)\chi(AB)+\chi(ABC)+\chi(ACB).
\end{aligned}}
Dieser Ausdruck ist symmetrisch in $A,B,C$, weil $\chi(ABC)$ bei cyklischer Vertauschung der Elemente $A,B,C$ nach (13.) ungeändert bleibt. Das allgemeine Bildungsgesetz der Coefficienten der Function
\frobtagc{15.}{n!\Phi_n(x)=\sum_{R_1,R_2,\ldots,R_n}\chi(R_1,R_2,\ldots,R_n)x_{R_1}x_{R_2}\cdots x_{R_n}}
ist etwas complicirt. Seien $A,B,C,D,F,G,H,\ldots,L,M$ irgend $n$ verschiedene oder gleiche unter den Elementen von $\Hgrp$. Man bilde die $n!$ Permutationen von $n$ Symbolen und zerlege jede derselben in cyklische Factoren. Setzt man dann für die $n$ Symbole die $n$ Elemente $A,B,C,\ldots,L,M$, so sei etwa
\frobtagc{16.}{(ABCD)(FGH)\cdots(LM)}
eine dieser $n!$ Permutationen. Man ordne ihr das Product
\frobtagc{17.}{\pm\chi(ABCD)\chi(FGH)\cdots\chi(LM)}
zu, wo das Vorzeichen $+$ oder $-$ zu wählen ist, je nachdem die Permutation (16.) gerade oder ungerade ist. In der Permutation (16.) bedeutet das Zeichen $(FGH)$, dass die drei Symbole $F,G,H$ cyklisch vertauscht werden sollen, in dem Ausdruck (17.) aber bedeutet $FGH$ das Product der drei Elemente $F,G,H$.

Eine gegebene Permutation kann nur in einer Weise als Product von cyklischen Factoren dargestellt werden. Doch kann man die einzelnen Cyklen beliebig anordnen und innerhalb eines jeden Cyklus, ohne dass er seine Bedeutung ändert, die Symbole cyklisch vertauschen. Andere Umstellungen aber sind nicht zulässig. So ist die Permutation (16.) gleich
\[
(GHF)(ML)\cdots(CDAB).
\]
In dem entsprechenden Producte (17.) sind aber dieselben Änderungen gestattet, denn $\chi(ABCD)$ bleibt ungeändert, wenn man die Elemente $A,B,C,D$ cyklisch vertauscht. Die Summe der den $n!$ Permutationen entsprechenden $n!$ Producte sei
\frobtagc{18.}{\chi(A,B,C,\cdots L,M)=\sum^{(n)}(\pm\Pi\chi).}
Diese Function bleibt ungeändert, wenn man die $n$ Elemente
\[
A,B,C,\cdots L,M
\]
beliebig unter einander vertauscht. Besteht die Permutation (16.) aus $l$ Cyklen, so ist sie gerade oder ungerade, je nachdem $n-l$ gerade oder ungerade ist. Daher kann man auch schreiben
\frobtagc{19.}{(-1)^n\chi(A,B,C,\cdots L,M)=\sum^{(n)}(\Pi-\chi).}
Z. B. ist, gleichgültig ob die durch $A,B,C,D$ bezeichneten Elemente verschieden sind oder nicht,
\frobtag{20.}{\begin{aligned}
\chi(A,B,C,D)={}&\chi(A)\chi(B)\chi(C)\chi(D)-\chi(B)\chi(C)\chi(AD)\nobreak\\
&{}-\chi(A)\chi(C)\chi(BD)-\chi(A)\chi(B)\chi(CD)\nobreak\\
&{}-\chi(A)\chi(D)\chi(BC)-\chi(B)\chi(D)\chi(AC)\nobreak\\
&{}-\chi(C)\chi(D)\chi(AB)+\chi(BC)\chi(AD)+\chi(AC)\chi(BD)\nobreak\\
&{}+\chi(AB)\chi(CD)+\chi(A)\chi(BCD)+\chi(B)\chi(ACD)\nobreak\\
&{}+\chi(C)\chi(ABD)+\chi(D)\chi(ABC)+\chi(A)\chi(BDC)\nobreak\\
&{}+\chi(B)\chi(ADC)+\chi(C)\chi(ADB)+\chi(D)\chi(ACB)\nobreak\\
&{}-\chi(ABCD)-\chi(ACBD)-\chi(BACD)\nobreak\\
&{}-\chi(BCAD)-\chi(CABD)-\chi(CBAD).
\end{aligned}}
Ich bilde nun die Summe
\[
(-1)^n\sum\chi(A,B,C,\cdots L,M)x_Ax_Bx_C\cdots x_Lx_M,
\]
worin jeder der $n$ Summationsbuchstaben $A,B,C,\cdots L,M$ unabhängig von den anderen die $h$ Elemente von $\Hgrp$ durchläuft. Der Ausdruck $(-1)^n\chi(A,B,C,\cdots L,M)$ ist eine Summe von $n!$ Producten $\Pi(-\chi)$. Die Permutation der $n$ Symbole, aus welcher eins dieser Producte gebildet ist, möge aus $a$ Cyklen von 1 Symbole, $b$ von 2 Symbolen, $c$ von 3 Symbolen u. s. w. bestehen, so dass $a+2b+3c+\cdots=n$ ist. Multiplicirt man dann dies Product $\Pi(-\chi)$ mit $x_Ax_Bx_C\cdots x_Lx_M$ und summirt, so erhält man nach (7.)
\[
(-1)^{a+b+c+\cdots}S_1^aS_2^bS_3^c\cdots .
\]
Dies Glied ergiebt sich so oft, als es Permutationen giebt, die sich in der angegebenen Art als Product von cyklischen Factoren darstellen lassen, also (\textsc{Cauchy}, \emph{Comptes rendus} tom. 21, p. 604)
\[
\frac{n!}{1^a2^b3^c\cdots a!b!c!\cdots}
\]
Mal. Mithin ist die betrachtete Summe gleich
\[
n!\sum\frac{(-1)^{a+b+c+\cdots}S_1^aS_2^bS_3^c\cdots}{1^a2^b3^c\cdots a!b!c!\cdots}=n!(-1)^n\Phi_n.
\]
Damit ist die Formel (15.) bewiesen. Ist also $n>f$, so ist
\frobtagc{21.}{\chi(R_1,R_2,\cdots R_n)=0\qquad(n>f).}
In jedem der $(n+1)!$ Producte der Summe $\chi(R,R_1,R_2,\ldots R_n)$ stelle man den Factor, der das mit $R$ bezeichnete Element enthält, an die erste Stelle, und in diesem Factor selbst stelle man mit Hülfe einer cyklischen Vertauschung $R$ an die erste Stelle. Dann nehme man zuerst die Producte, die den Factor $\chi(R)$ enthalten, dann die, worin auf $R$ das Element $R_1$ folgt, dann die, worin auf $R$ das Element $R_2$ folgt u. s. w. Auf diese Weise erhält man die Recursionsformel
\frobtag{22.}{\begin{aligned}
\chi(R,R_1,R_2,R_3,\cdots R_n)={}&\chi(R)\chi(R_1,R_2,R_3,\cdots R_n)\nobreak\\
&{}-\chi(RR_1,R_2,R_3,\cdots R_n)-\chi(R_1,RR_2,R_3,\cdots R_n)\nobreak\\
&{}-\chi(R_1,R_2,RR_3,\cdots R_n)-\cdots-\chi(R_1,R_2,R_3,\cdots RR_n).
\end{aligned}}
Daraus geht hervor, dass, wenn für einen Werth von $n$ die Grössen $\chi(R_1,R_2,\cdots R_n)$ sämmtlich verschwinden, dasselbe auch für jeden grösseren Werth von $n$ eintreten muss. Speciell ist
\frobtagc{23.}{\chi(E,R_1,R_2,\cdots R_n)=(f-n)\chi(R_1,R_2,\cdots R_n).}

\section*{§ 4.}

Differentiirt man die Gleichung $\Phi(z)=\Phi(x)\Phi(y)$ nach $y_A$, so erhält man
\[
\sum_R\frac{\partial\Phi(z)}{\partial z_R}x_{RA^{-1}}=\Phi(x)\frac{\partial\Phi(y)}{\partial y_A},
\]
und wenn man $y_R=\varepsilon_R$ setzt,
\frobtagc{1.}{\sum_R\frac{\partial\Phi(x)}{\partial x_R}x_{RA^{-1}}=\chi(A)\Phi(x).}
Differentiirt man aber jene Gleichung nach $x_A$, so findet man auf demselben Wege
\frobtagc{2.}{\sum_R\frac{\partial\Phi(x)}{\partial x_R}x_{A^{-1}R}=\chi(A)\Phi(x).}
In diesen Gleichungen ersetze ich $x_E$ durch $x_E-u$. Dann ergiebt sich daraus, da
\[
\frac{\partial\log\Phi(x-u\varepsilon)}{\partial u}=\frac{1}{u-u_1}+\cdots+\frac{1}{u-u_f}=\sum_{n=0}^{\infty}\frac{S_n}{u^{n+1}}
\]
ist, die Recursionsformel
\frobtag{3.}{\frac{1}{n+1}\frac{\partial S_{n+1}}{\partial x_A}=\frac{1}{n}\sum_R\frac{\partial S_n}{\partial x_R}x_{RA^{-1}}=\frac{1}{n}\sum_R\frac{\partial S_n}{\partial x_R}x_{A^{-1}R},\qquad(n>0)}
aus der sich ein neuer Beweis für die Formel (7.) § 3 ableiten lässt. Differentiirt man die Gleichung (1.) nach $x_S$, multiplicirt sie dann mit $x_{SB^{-1}}$ und summirt nach $S$, so erhält man die Gleichung
\[
\sum_{R,S}\frac{\partial^2\Phi}{\partial x_R\partial x_S}x_{RA^{-1}}x_{SB^{-1}}=(\chi(A)\chi(B)-\chi(BA))\Phi,
\]
aus der sich direct die Formel (13.) § 3 ergiebt. Ebenso ist allgemein, wenn $A,B,\cdots M$ irgend $n$ Elemente von $\Hgrp$ sind,
\[
\sum_{R,S,\ldots,V}\frac{\partial^n\Phi}{\partial x_R\partial x_S\cdots\partial x_V}x_{RA^{-1}}x_{SB^{-1}}\cdots x_{VM^{-1}}=\chi(A,B,\cdots M)\Phi.
\]
Hier mache ich von jenen Relationen eine andere Anwendung. Setzt man
\[
y_{R^{-1}}=\frac{\partial\Phi(x-u\varepsilon)}{\partial x_R},
\]
so lauten sie
\[
\sum_R y_{PR^{-1}}(x_{RQ^{-1}}-u\varepsilon_{RQ^{-1}})=\sum_R(x_{PR^{-1}}-u\varepsilon_{PR^{-1}})y_{RQ^{-1}}=\chi(QP^{-1})\Phi(x-u\varepsilon).
\]
Ich setze noch $\chi(R^{-1})=\chi'(R)$ und bezeichne die Matrix $(\chi'(PQ^{-1}))$ kurz mit $(\chi')$. Dann drückt diese Formel die folgende Beziehung zwischen Matrizen aus
\frobtagc{4.}{(y)((x)-u(\varepsilon))=((x)-u(\varepsilon))(y)=(\chi')\Phi(x-u\varepsilon).}
Mithin ist $(y)(x)=(x)(y)$. Ist also, nach Potenzen von $u$ entwickelt, $(y)=(p)+(q)u+(r)u^2+\cdots$, so ist $(x)$ mit jeder der Matrizen $(p),(q),(r),\cdots$ vertauschbar. Entwickelt man nun in der Relation (4.) auch $\Phi(x-u\varepsilon)$ nach Potenzen von $u$, so müssen die Coefficienten der einzelnen Potenzen von $u$ auf beiden Seiten übereinstimmen. Die so erhaltenen Gleichungen füge man wieder zusammen, nachdem man sie, statt mit den Potenzen der Variabelen $u$, mit den entsprechenden Potenzen der Matrix $(x)$ multiplicirt hat. Dies Verfahren führt zu demselben Resultate, wie wenn man direct in der Gleichung (4.) die Variable $u$ durch die Matrix $(x)$ ersetzt. (Ausführlicher ist diese Schlussweise entwickelt V., S. 605.) Dann ergiebt sich $(\chi')\Phi(x-(x)\varepsilon)=0$ oder deutlicher $(\chi')\Phi(x_E-(x),x_A,x_B,x_C,\cdots)=0$ und noch ausführlicher nach (2.) § 3
\frobtag{5.}{(\chi')\bigl((x)^f-(x)^{f-1}\Phi_1+(x)^{f-2}\Phi_2-\cdots+(-1)^f(x)^0\Phi_f\bigr)=0.}
Multiplicirt man noch mit $(x)^{n-f}$, so erhält man
\frobtag{6.}{\sum_R\chi(AR)\bigl(x_R^{(n)}-x_R^{(n-1)}\Phi_1+x_R^{(n-2)}\Phi_2-\cdots+(-1)^f x_R^{(n-f)}\Phi_f\bigr)=0.}
Setzt man $n=f$, und bestimmt man den Coefficienten von $x_S^f$, so findet man
\frobtag{7.}{\chi(AS^f)-\mathfrak{z}_1(S)\chi(AS^{f-1})+\mathfrak{z}_2(S)\chi(AS^{f-2})-\cdots+(-1)^f\mathfrak{z}_f(S)\chi(A)=0,}
wo
\frobtagc{8.}{\mathfrak{z}_n(S)=\frac{1}{n!}\chi(S,S,\cdots S)}
der Coefficient von $x_S^n$ in $\Phi_n$ ist, also von $A$ unabhängig ist. Speciell ist $\mathfrak{z}_1(S)=\chi(S)$ und $\mathfrak{z}_f(S)=\mathfrak{z}(S)$. Setzt man in der Function $\Phi(x_E+u,x_A,x_B,\cdots,x_S,\cdots)$ alle Variabelen gleich Null ausser $u$ und $x_S$, so wird
\frobtag{9.}{\Phi(u,0,0,\cdots x_S,0,\cdots)=u^f+\mathfrak{z}_1(S)u^{f-1}x_S+\mathfrak{z}_2(S)u^{f-2}x_S^2+\cdots+\mathfrak{z}_f(S)x_S^f.}
Daher ist $\mathfrak{z}_n(E)=\binom{f}{n}$.

\section*{§ 5.}

Die bisherigen Ergebnisse habe ich allein aus der Relation (9.) § 1 abgeleitet, ohne dabei die Unzerlegbarkeit von $\Phi$ und den Exponenten $e$ der Potenz von $\Phi$ zu benutzen, durch welche die Gruppendeterminante $\Theta$ theilbar ist. Jede der $h$ Variabelen $x_R$ kommt in jeder Zeile und in jeder Spalte von $\Theta$ einmal vor, im Ganzen also an $h$ Stellen. An jeder dieser $h$ Stellen ist ihr aber dieselbe Unterdeterminante complementär, wie ich \emph{Ch.} § 6 gezeigt habe. Ist also $\Theta_{P,Q}$ die Unterdeterminante, die dem Elemente $x_{P,Q}$ in der Determinante $\Theta=|x_{P,Q}|$ complementär ist, so ist
\frobtagc{1.}{\Theta_{P,Q}=\frac{1}{h}\frac{\partial\Theta}{\partial x_{PQ^{-1}}}=\Theta_{PQ^{-1}},}
falls man $h\Theta_R=\frac{\partial\Theta}{\partial x_R}$ setzt. Die Unterdeterminanten von $\Theta$ bilden demnach eine Matrix, welche dieselben Symmetrieeigenschaften hat, wie die Matrix $(x)$. Nach den bekannten Relationen zwischen den Elementen einer Determinante und den ihnen complementären Unterdeterminanten ist
\[
\sum_R x_{P,R}\Theta_{Q,R}=\sum_R x_{R,P}\Theta_{R,Q}=\varepsilon_{P,Q}\Theta
\]
oder
\frobtagc{2.}{\sum_R x_{AR}\frac{\partial\Theta}{\partial x_R}=\sum_R x_{RA}\frac{\partial\Theta}{\partial x_R}=\varepsilon_A h\Theta.}
Nach (7.) § 1 ist $\Theta=\Phi^e\Psi$, wo $\Psi$ zu $\Phi$ theilerfremd ist. Mithin ist
\[
\sum_R x_{AR}\left(e\frac{\partial\log\Phi}{\partial x_R}+\frac{\partial\log\Psi}{\partial x_R}\right)=\varepsilon_A h
\]
oder
\[
\sum_R x_{AR}\left(e\Psi\frac{\partial\Phi}{\partial x_R}+\Phi\frac{\partial\Psi}{\partial x_R}\right)=\varepsilon_A h\Phi\Psi.
\]
Da $\Psi$ zu $\Phi$ theilerfremd ist, so ist folglich $\sum_Rx_{AR}\frac{\partial\Phi}{\partial x_R}$ durch $\Phi$ theilbar, und weil beide Functionen von demselben Grade sind, so können sie sich nur um einen constanten Factor unterscheiden. Dieser ergiebt sich durch die Vergleichung der Coefficienten von $x_E^f$, und mithin besteht die Formel
\frobtagc{3.}{\sum_R x_{AR}\frac{\partial\Phi}{\partial x_R}=\sum_R x_{RA}\frac{\partial\Phi}{\partial x_R}=\chi(A^{-1})\Phi,}
die auf einem anderen Wege schon in § 4 abgeleitet ist. Multiplicirt man die Gleichung
\[
\sum_S x_{RS}\frac{\partial\Phi}{\partial x_S}=\chi(R^{-1})\Phi
\]
mit $\Theta_{RA}$ und summirt nach $R$, so erhält man
\[
h\Theta\frac{\partial\Phi}{\partial x_A}=\Phi\sum_R\chi(R^{-1})\frac{\partial\Theta}{\partial x_{RA}}=\Phi\sum_R\chi(AR^{-1})\frac{\partial\Theta}{\partial x_R}.
\]
Setzt man hier $\Theta=\Phi^e\Psi$, so erhält man
\[
h\Psi\frac{\partial\Phi}{\partial x_A}=\sum_R\chi(AR^{-1})\left(e\Psi\frac{\partial\Phi}{\partial x_R}+\Phi\frac{\partial\Psi}{\partial x_R}\right),
\]
und folglich ist
\frobtagc{4.}{\frac{h}{e}\frac{\partial\Phi}{\partial x_A}=\sum_R\chi(AR^{-1})\frac{\partial\Phi}{\partial x_R},}
weil die Differenz dieser beiden Functionen durch $\Phi$ theilbar und nur vom Grade $f-1$ ist. Setzt man aber $\Theta=\Phi^{e'}\Psi'$, wo $\Phi'$ ein von $\Phi$ verschiedener Primfactor von $\Theta$ ist, so erhält man
\[
h\Phi'\Psi'\frac{\partial\Phi}{\partial x_A}=\Phi\sum_R\chi(AR^{-1})\left(e'\Psi'\frac{\partial\Phi'}{\partial x_R}+\Phi'\frac{\partial\Psi'}{\partial x_R}\right),
\]
und folglich ist
\frobtagc{5.}{\sum_R\chi(AR^{-1})\frac{\partial\Phi'}{\partial x_R}=0,}
weil diese Function durch $\Phi'$ theilbar ist. Vergleicht man in diesen Relationen die Coefficienten von $x_E^{f-1}$ (bez. $x_E^{f'-1}$), so ergiebt sich
\[
\sum_R\chi(AR^{-1})\chi(R)=\frac{h}{e}\chi(A),\qquad \sum_R\chi(AR^{-1})\psi(R)=0,
\]
wo $\psi(R)$ der der Primfunction $\Phi'$ entsprechende Charakter ist. Man kann diese Gleichungen auch schreiben
\frobtag{6.}{\sum\chi(R)\chi(S)=\frac{h}{e}\chi(A),\qquad \sum\chi(R)\psi(S)=0\qquad(RS=A)}
oder
\frobtag{7.}{\sum_R\chi(PR^{-1})\chi(RQ^{-1})=\frac{h}{e}\chi(PQ^{-1}),\qquad \sum_R\chi(PR^{-1})\psi(RQ^{-1})=0.}
Für $A=E$ ist
\frobtag{8.}{\sum_R\chi(R)\chi(R^{-1})=\frac{hf}{e},\qquad \sum_R\chi(R)\psi(R^{-1})=0.}
Daraus fliesst die Folgerung, dass die Werthe $\chi(E),\chi(A),\chi(B),\cdots$ des Charakters $\chi$ den entsprechenden Werthen $\psi(E),\psi(A),\psi(B),\cdots$ des Charakters $\psi$ nicht proportional sein können. Bezeichnet man die Matrix $(\chi(PQ^{-1}))$ kurz mit $(\chi)$, so kann man die Relationen (7.) auch auf die Form
\frobtagc{9.}{\left(\frac{e}{h}\chi\right)^2=\left(\frac{e}{h}\chi\right),\qquad (\chi)(\psi)=0}
bringen.

\section*{§ 6.}

Gleichzeitig mit $P$ durchläuft auch $P^{-1}$ die $h$ Elemente von $\Hgrp$, nur in einer anderen Reihenfolge. Daher ist
\[
|x_{P,Q}|=|x_{P^{-1},Q^{-1}}|=|x_{Q^{-1},P^{-1}}|,
\]
also
\frobtagc{1.}{|x_{PQ^{-1}}|=|x_{Q^{-1}P}|.}
In jeder dieser beiden Determinanten erhält man die Zeilen, indem man für $P$, die Spalten, indem man für $Q$ die Elemente $G_1,G_2,\cdots G_h$ von $\Hgrp$ setzt. Sind also $x_R$ und $y_R$ zwei Systeme von je $h$ Variabelen, so ist
\[
|x_{PQ^{-1}}|=\Pi(\Phi(x)^e),\qquad |y_{Q^{-1}P}|=\Pi(\Phi(y)^e).
\]
Die beiden Matrizen $(x_{PQ^{-1}})$ und $(y_{Q^{-1}P})$ sind aber mit einander vertauschbar; es ist
\[
\sum_R x_{PR^{-1}}y_{Q^{-1}R}=\sum_S y_{S^{-1}P}x_{SQ^{-1}}.
\]
Denn setzt man $SQ^{-1}=PR^{-1}$, also $S=PR^{-1}Q$, so durchläuft $S$ gleichzeitig mit $R$ die $h$ Elemente von $\Hgrp$, und es ist auch $S^{-1}P=Q^{-1}R$. Seien $a_1,a_2,a_3,\cdots$ die $h$ Wurzeln der (charakteristischen Gleichung der) Matrix $(x_{PQ^{-1}})$, also die Wurzeln der Gleichung $|x_{PQ^{-1}}-u\varepsilon_{PQ^{-1}}|=0$, und $b_1,b_2,b_3,\cdots$ die $h$ Wurzeln der Gleichung $|y_{Q^{-1}P}-u\varepsilon_{PQ^{-1}}|=0$, also auch der Gleichung $|y_{PQ^{-1}}-u\varepsilon_{PQ^{-1}}|=0$. Dann lassen sich (V., S. 602, III) diese beiden Reihen von je $h$ Wurzeln einander so zuordnen, dass $a_1+b_1,a_2+b_2,a_3+b_3,\cdots$ die Wurzeln der Matrix $(x_{PQ^{-1}}+y_{Q^{-1}P})$ werden.

Auf diesen allgemeinen Satz komme ich später (§ 10) zurück. Hier mache ich jetzt die Voraussetzung, dass für je zwei Elemente von $\Hgrp$
\frobtagc{2.}{y_{BA}=y_{AB}}
ist. Theilt man also die $h$ Elemente von $\Hgrp$ in Classen conjugirter Elemente, so hat $y_R$ für alle Elemente einer Classe denselben Werth, etwa für die Elemente $R$ der $\rho^{\mathrm{ten}}$ Classe den Werth $y_R=y_\rho$. Ist $k$ die Anzahl der Classen, so seien die $k$ Variabelen $y_0,y_1,\cdots,y_{k-1}$, die den $k$ Classen $(0),(1),\ldots,(k-1)$ entsprechen, von einander unabhängig. Da nun $y_{Q^{-1}P}=y_{PQ^{-1}}$ ist, so ist die Matrix $(y_{PQ^{-1}})$ mit der Matrix $(x_{PQ^{-1}})$ vertauschbar. Jede Matrix $(x)$, welche die in § 1 definirten Symmetrieeigenschaften besitzt, ist mit jeder anderen Matrix $(y)$ vertauschbar, deren Elemente ausserdem noch den Bedingungen (2.) genügen. Sind daher $u,v,w$ Variabele, so ist die Determinante
\[
|ux_{PQ^{-1}}+vy_{PQ^{-1}}+w\varepsilon_{PQ^{-1}}|=\Pi\bigl(\Phi(ux+vy+w\varepsilon)^e\bigr)
\]
ein Product von linearen Functionen von $u,v,w$, und mithin ist auch
\frobtagc{3.}{\Phi(ux+vy+w\varepsilon)=(u_1u+v_1v+w)(u_2u+v_2v+w)\cdots(u_fu+v_fv+w),}
wo $u_1,v_1,\cdots,u_f,v_f$ von $u,v,w$ unabhängig sind. Den Coefficienten von $w$ kann man in jeder dieser $f$ linearen Functionen gleich 1 voraussetzen, weil die linke Seite für $u=v=0$ gleich $w^f$ wird. Setzt man $v=0$ und $u=1$, so erhält man
\frobtagc{4.}{\Phi(x+w\varepsilon)=(u_1+w)(u_2+w)\cdots(u_f+w),}
setzt man $u=0$ und $v=1$,
\frobtagc{5.}{\Phi(y+w\varepsilon)=(v_1+w)(v_2+w)\cdots(v_f+w).}
Daher hängen $u_1,u_2,\cdots,u_f$ nur von den $h$ Variabelen $x_R$ ab und haben dieselbe Bedeutung wie in § 3, während $v_1,v_2,\cdots,v_f$ nur von den $k$ Variabelen $y_\rho$ abhängen.

Da $\Phi(x)$ unzerlegbar ist, so ist auch $\Phi(x+w\varepsilon)$ als Function von $w$ irreducibel, d. h. dieser Ausdruck kann nicht als Product zweier ganzen Functionen von $w$ dargestellt werden, deren Coefficienten rationale Functionen der $h$ unabhängigen Variabelen $x_R$ sind. Betrachtet man die $k$ Grössen $y_\rho$ als constant, so ist auch $v_1$ eine Constante, und mithin ist auch
\[
\Phi(x_E+v_1+w,x_A,x_B,\cdots)=(u_1+v_1+w)(u_2+v_1+w)\cdots(u_f+v_1+w)
\]
als Function von $w$ irreducibel. Diese Function hat aber mit der Function
\[
\Phi(x+y+w\varepsilon)=(u_1+v_1+w)(u_2+v_2+w)\cdots(u_f+v_f+w)
\]
den linearen Factor $u_1+v_1+w$ gemeinsam. Folglich müssen beide Functionen identisch sein. Setzt man $w=0$ und $v_1=\eta$, so ist also
\frobtagc{6.}{\Phi(x_E+y_E,x_A+y_A,x_B+y_B,\cdots)=\Phi(x_E+\eta,x_A,x_B,\cdots),}
wo $\eta$ nur von den Variabelen $y_R$ abhängt. Setzt man die Variabelen $x_R$ alle gleich Null ausser $x_E=u$, so erhält man
\frobtagc{7.}{\Phi(y_E+u,y_A,y_B,\cdots)=(u+\eta)^f,\qquad(y_{BA}=y_{AB})}
und daraus durch Vergleichung der Coefficienten von $u^{f-1}$
\frobtagc{8.}{f\eta=\sum_R\chi(R)y_R.}
So ergiebt sich der Satz:

\emph{Ist für je zwei Elemente $A$ und $B$ der Gruppe $\Hgrp$ $x_{BA}=x_{AB}$, so wird die Primfunction $f^{\mathrm{ten}}$ Grades $\Phi(x)$ gleich der $f^{\mathrm{ten}}$ Potenz einer linearen Function $\xi=\frac{1}{f}\sum_R\chi(R)x_R$.}

Für alle Elemente $R$ der $\rho^{\mathrm{ten}}$ Classe, deren Anzahl $h_R=h_\rho$ sei, hat $\chi(R)$ denselben Werth $\chi_\rho$. Ist $x_{BA}=x_{AB}$, so hat auch $x_R$ für diese $h_\rho$ Elemente denselben Werth $x_\rho$. Daher ist $\sum\chi(R)x_R=\sum h_\rho\chi_\rho x_\rho$. Ist $\Phi'$ ein von $\Phi$ verschiedener Primfactor von $\Theta$, und $\psi(R)$ der ihm entsprechende Charakter, so können nach § 5 die $k$ Grössen $\psi_\rho$ den $k$ Grössen $\chi_\rho$ nicht proportional sein, und mithin können die beiden linearen Functionen $\sum h_\rho\psi_\rho x_\rho$ und $\sum h_\rho\chi_\rho x_\rho$ sich nicht etwa nur um einen constanten Factor unterscheiden. Wendet man nun den obigen Satz auf jeden Primfactor von $\Theta$ an, so erhält man
\frobtagc{9.}{|x_{PQ^{-1}}|=\Pi\left(\frac{1}{f}\sum_R\chi(R)x_R\right)^f=\Pi(\xi^f).\qquad(x_{BA}=x_{AB})}
Jedem Primfactor $\Phi$ der Determinante (7.) § 1, worin die $h$ Variabelen $x_R$ von einander unabhängig sind, entspricht ein Linearfactor $\xi$ der Determinante (9.), worin $x_{BA}=x_{AB}$ gesetzt ist. Sind die $k$ Variabelen $x_\rho$ unabhängig, so entsprechen zwei verschiedenen Primfactoren jener Determinante zwei wesentlich verschiedene Linearfactoren der Determinante (9.). Ist $f$ der Grad der Primfunction $\Phi$, und $e$ der Exponent der in $\Theta$ aufgehenden Potenz von $\Phi$, so ist $g=ef$ der Exponent der Potenz des entsprechenden Linearfactors $\xi$, wodurch die Determinante (9.) theilbar ist.

Vergleicht man das erhaltene Resultat mit der Formel (22.) \emph{Ch.} § 5, so erkennt man, dass die Grössen $\chi_\rho$, welche dort auf einem ganz anderen Wege als die Charaktere der Gruppe $\Hgrp$ eingeführt sind, mit den hier definirten Grössen $\chi(R)$ völlig übereinstimmen, und ebenso die Zahlen $g=ef$. Nur waren dort die beiden Factors $e$ und $f$ von $g$ willkürlich gelassen, während sie hier einzeln eine bestimmte Bedeutung haben. Man könnte daher jetzt von allen dort entwickelten Eigenschaften der Charaktere Gebrauch machen. Indessen ziehe ich es vor, die Ergebnisse jener Arbeit von dem hier gewählten Ausgangspunkt aus noch einmal abzuleiten.

\section*{§ 7.}

Wenn man die Gleichung (6.) § 6 nach $y_\rho$ differentiirt und dann die $k$ Variabelen $y_\rho$ alle gleich Null setzt, so findet man
\[
\sum_{(\rho)}'\frac{\partial\Phi(x)}{\partial x_R}=\frac{h_\rho\chi_\rho}{f}\frac{\partial\Phi(x)}{\partial x_E},
\]
wo der Summationsbuchstabe $R$ die $h_\rho$ Elemente der $\rho^{\mathrm{ten}}$ Classe durchläuft. Ist $A$ ein festes Element, und ersetzt man für jedes $R$ die Variable $x_R$ durch $x_{AR}$, so ändert sich $\Phi$ nach (10.) § 3 nur um einen constanten, von Null verschiedenen Factor. Mithin ist auch
\frobtagc{1.}{\sum_{(\rho)}'\frac{\partial\Phi(x)}{\partial x_{AR}}=\frac{h_\rho\chi_\rho}{f}\frac{\partial\Phi(x)}{\partial x_A},}
wo wieder $R$ die $h_\rho$ Elemente der Classe $(\rho)$ durchläuft. Vergleicht man die Coefficienten von $x_E^{f-1}$, so erhält man
\frobtagc{2.}{\sum'\chi(AS)=\frac{h_B}{f}\chi(A)\chi(B),}
wo $S$ die $h_B$ mit $B$ conjugirten Elemente durchläuft. Direct erhält man diese Relation, indem man für die Function (7.) § 6 den Ausdruck $S_2(y)=\sum\chi(PQ)y_Py_Q=f\eta^2$ berechnet. Setzt man $S=R^{-1}BR$ und für $R$ alle $h$ Elemente von $\Hgrp$, so wird $S$ jedem Elemente der Classe von $B$ gleich und zwar jedem $\frac{h}{h_B}$ Mal. Daher ist (\emph{Ch.} § 5, (5.))
\frobtagc{3.}{h\chi(A)\chi(B)=f\sum_R\chi(AR^{-1}BR).}
Zu demselben Resultate gelangt man direct von der Formel (9.) § 6 aus auf dem \emph{Ch.} S. 1001 angegebenen Wege, also mittelst derselben Schlüsse, die in § 1 zu der Formel (9.) geführt haben.

Die Anzahl der verschiedenen Primfactoren der Gruppendeterminante $\Theta$ sei $l$. Diese $l$ Functionen $\Phi$ und die ihnen entsprechenden Charaktere $\chi$ mögen durch obere Indices $\lambda=0,1,\cdots,l-1$ von einander unterschieden werden. Nun ist in der Entwicklung der Determinante $\Theta(x+u\varepsilon)$ nach Potenzen von $u$ der Coefficient von $u^{h-1}$ gleich $hx_E$. Ersetzt man daher in dem Producte
\frobtagc{4.}{\Theta=\Pi_\lambda\bigl(\Phi^{(\lambda)e^{(\lambda)}}\bigr)}
$x_E$ durch $x_E+u$ und vergleicht dann die Coefficienten von $x_E^{h-1}$, so erhält man
\frobtagc{5.}{\sum_\lambda e^{(\lambda)}\chi^{(\lambda)}(R)=h\varepsilon_R.}
Nun ist nach (8.) § 5 und (3.)
\[
\sum_R\chi(S^{-1}R^{-1}SR)=\frac{h}{f}\chi(S)\chi(S^{-1}),\qquad \sum_S\chi(S)\chi(S^{-1})=\frac{hf}{e},
\]
und mithin für jeden Werth von $\lambda$
\[
\sum_{R,S}\frac{e^{(\lambda)}}{h}\chi^{(\lambda)}(S^{-1}R^{-1}SR)=h,
\]
und folglich, da dieser Werth von $\lambda$ unabhängig ist,
\[
\sum_\lambda\sum_{R,S}\frac{e^{(\lambda)}}{h}\chi^{(\lambda)}(S^{-1}R^{-1}SR)=hl.
\]
Jetzt kehre ich die Reihenfolge der beiden Summationen um. Dann ist nach (5.)
\[
\sum_\lambda\frac{e^{(\lambda)}}{h}\chi^{(\lambda)}(S^{-1}R^{-1}SR)=0,
\]
ausser wenn $S^{-1}R^{-1}SR=E$, also $SR=RS$ ist; dann ist die Summe gleich 1. Daher ist $hl$ gleich der Anzahl der Lösungen der Gleichung $SR=RS$. Diese aber ist, wie ich \emph{Ch.} S. 987 durch die einfachsten Betrachtungen gezeigt habe, gleich $hk$, und folglich ist $l=k$.

\emph{Die Anzahl der verschiedenen Primfactoren der Gruppendeterminante ist gleich der Anzahl der Classen conjugirter Elemente, worin die Elemente der Gruppe zerfallen.}

Durchläuft $R$ die Classe $(\alpha)$, so durchläuft $R^{-1}$ die inverse Classe, die ich mit $(\alpha')$ bezeichne. Daher sind die Gleichungen (8.) § 5 identisch mit
\frobtag{6.}{\sum_\alpha h_\alpha\chi_\alpha^{(\kappa)}\chi_{\alpha'}^{(\kappa)}=\frac{h f^{(\kappa)}}{e^{(\kappa)}},\qquad \sum_\alpha h_\alpha\chi_\alpha^{(\kappa)}\chi_{\alpha'}^{(\lambda)}=0.}
Demnach sind die beiden Matrizen des $k^{\mathrm{ten}}$ Grades
\frobtagc{7.}{\left(\frac{h_\alpha}{h}\chi_{\alpha'}^{(\kappa)}\right),\qquad \left(\frac{e^{(\kappa)}}{f^{(\kappa)}}\chi_\alpha^{(\kappa)}\right),}
complementär, und folglich bestehen auch die Gleichungen
\[
\sum_\kappa\frac{e^{(\kappa)}}{f^{(\kappa)}}\chi_\alpha^{(\kappa)}\chi_{\alpha'}^{(\kappa)}=\frac{h}{h_\alpha},\qquad \sum_\kappa\frac{e^{(\kappa)}}{f^{(\kappa)}}\chi_\alpha^{(\kappa)}\chi_\beta^{(\kappa)}=0.
\]
Setzt man also $h_{\alpha\beta}=h_\alpha$ oder 0, je nachdem $(\beta)=(\alpha')$ ist oder nicht, so ist
\frobtagc{8.}{\sum_\kappa\frac{e^{(\kappa)}}{f^{(\kappa)}}\chi_\alpha^{(\kappa)}\chi_\beta^{(\kappa)}=\frac{h\,h_{\alpha\beta}}{h_\alpha h_\beta}.}
Diese Gleichung erhält man unmittelbar aus der Relation (2.). Ist nämlich $\chi(A)=\chi_\alpha$ und $\chi(B)=\chi_\beta$, so lautet diese
\[
\sum_S'\chi^{(\kappa)}(AS)=\frac{h_\beta}{f^{(\kappa)}}\chi_\alpha^{(\kappa)}\chi_\beta^{(\kappa)},
\]
wo $S$ die $h_\beta$ Elemente der $\beta^{\mathrm{ten}}$ Classe durchläuft. Daher ist
\[
\sum_\kappa\sum_S'\frac{e^{(\kappa)}}{h}\chi^{(\kappa)}(AS)=\frac{h_\beta}{h}\sum_\kappa\frac{e^{(\kappa)}}{f^{(\kappa)}}\chi_\alpha^{(\kappa)}\chi_\beta^{(\kappa)}.
\]
Nun ist aber $\sum_\kappa\frac{e^{(\kappa)}}{h}\chi^{(\kappa)}(AS)=0$, ausser wenn $AS=E$, also $S=A^{-1}$ ist. Dann ist die Summe gleich 1. Das letztere tritt aber stets und nur dann ein, wenn $(\beta)=(\alpha')$ ist, und so ergiebt sich die Gleichung (8.). Sowohl aus dieser wie aus (6.) folgt, dass die Determinante $k^{\mathrm{ten}}$ Grades
\frobtagc{9.}{|\chi_\alpha^{(\kappa)}|}
von Null verschieden ist.

Sei allgemeiner $A$ ein bestimmtes Element der $\alpha^{\mathrm{ten}}$ Classe und $R$ ein veränderliches Element, das die $h_\alpha$ mit $A$ conjugirten Elemente durchläuft. Dieselbe Bedeutung mögen für die $\beta^{\mathrm{te}}$ Classe die Zeichen $B$ und $S$, für die $\gamma^{\mathrm{te}}$ die Zeichen $C$ und $T$ haben. Dann ist nach (2.)
\[
\sum_R'\chi(RSC)=\frac{h_\alpha}{f}\chi(A)\chi(SC),\qquad \sum_S'\chi(SC)=\frac{h_\beta}{f}\chi(B)\chi(C),
\]
also
\[
\sum_{R,S}'\chi(RSC)=\frac{h_\alpha h_\beta}{f^2}\chi_\alpha\chi_\beta\chi_\gamma
\]
und mithin
\[
\sum_\kappa\sum_{R,S}'\frac{e^{(\kappa)}}{h}\chi^{(\kappa)}(RSC)=\frac{h_\alpha h_\beta}{h}\sum_\kappa\frac{e^{(\kappa)}}{f^{(\kappa)2}}\chi_\alpha^{(\kappa)}\chi_\beta^{(\kappa)}\chi_\gamma^{(\kappa)}.
\]
Die linke Seite ist gleich der Anzahl der Lösungen der Gleichung $RSC=E$. Die rechte Seite zeigt, dass diese Anzahl nicht von dem Elemente $C$, sondern nur von der Classe $(\gamma)$ abhängt, der $C$ angehört. Bezeichnet man sie mit $\frac{h_{\alpha\beta\gamma}}{h_\gamma}$, so ist demnach
\frobtagc{10.}{\frac{h\,h_{\alpha\beta\gamma}}{h_\alpha h_\beta h_\gamma}=\sum_\kappa\frac{e^{(\kappa)}}{f^{(\kappa)2}}\chi_\alpha^{(\kappa)}\chi_\beta^{(\kappa)}\chi_\gamma^{(\kappa)}.}
Setzt man für $C$ der Reihe nach die $h_\gamma$ Elemente $T$ der Classe $(\gamma)$, so ist folglich $h_{\alpha\beta\gamma}$ die Anzahl der Lösungen der Gleichung $RST=E$, falls $R$ die $h_\alpha$ Elemente der Classe $(\alpha)$ durchläuft, $S$ die $h_\beta$ der Classe $(\beta)$ und $T$ die der Classe $(\gamma)$. Wie die rechte Seite dieser Gleichung zeigt, bleibt die Zahl $h_{\alpha\beta\gamma}$ bei jeder Vertauschung ihrer drei Indices unter einander ungeändert. Daher ist auch $\frac{h_{\alpha\beta\gamma}}{h_\beta}$ die Anzahl der Lösungen der Gleichung $RBT=E$, worin an Stelle von $S$ ein festes Element $B$ der $\beta^{\mathrm{ten}}$ Classe getreten ist.

Aus (6.) und (10.) ergeben sich die Relationen
\[
h_\beta h_\gamma\chi_\beta^{(\kappa)}\chi_\gamma^{(\kappa)}=f^{(\kappa)}\sum_\alpha h_{\alpha\beta\gamma}\chi_\alpha^{(\kappa)}.
\]
Diese zeigen, dass die Grössen $\chi_\alpha$ den Gleichungen
\frobtagc{11.}{h_\beta h_\gamma\chi_\beta\chi_\gamma=f\sum_\alpha h_{\alpha'\beta\gamma}\chi_\alpha}
genügen. Dieselben haben daher $k$ Systeme von Lösungen
\[
\chi_\alpha=\chi_\alpha^{(\kappa)}\qquad(\kappa=0,1,\cdots k-1)
\]
aber nicht mehr. Denn sind $x_0,x_1,\cdots,x_{k-1}$ Variabele und setzt man
\frobtagc{12.}{\sum_\alpha h_\alpha\chi_\alpha^{(\kappa)}x_\alpha=f^{(\kappa)}\xi^{(\kappa)},}
so folgt aus (11.)
\[
h_\alpha\chi_\alpha^{(\kappa)}\xi^{(\kappa)}=\sum_\beta\left(\sum_\gamma h_{\alpha\beta\gamma}x_\gamma\right)\chi_\beta^{(\kappa)}
\]
oder
\[
\sum_\beta\left(\left(\sum_\gamma h_{\alpha\beta\gamma}x_\gamma\right)-h_{\alpha\beta}\xi^{(\kappa)}\right)\chi_\beta^{(\kappa)}=0.
\]
Folglich verschwindet die Determinante
\frobtagc{13.}{\left|\left(\sum_\gamma h_{\alpha\beta\gamma}x_\gamma\right)-h_{\alpha\beta}r\right|,\qquad(\alpha,\beta=0,1,\cdots k-1)}
die eine ganze Function $k^{\mathrm{ten}}$ Grades von $r$ ist, für die $k$ Werthe $r=\xi^{(\kappa)}$, die unter einander verschieden sind (vergl. \textsc{Dedekind}, \emph{Zur Theorie der aus $n$ Haupteinheiten gebildeten complexen Grössen}. Göttinger Nachrichten 1885, S. 146).

Die Gleichungen (11.) bestimmen die Grössen $\frac{\chi_\alpha}{f}$. Alsdann liefert die Gleichung (6.)
\[
\frac{h}{ef}=\sum_\alpha h_\alpha\frac{\chi_\alpha}{f}\frac{\chi_{\alpha'}}{f}
\]
zu jedem der $k$ Werthsysteme $\frac{\chi_\alpha}{f}$ den entsprechenden Werth von $ef=g$.

Nach § 3 ist demnach zur vollständigen Berechnung der $k$ Primfunctionen $\Phi$ weiter nichts mehr erforderlich, als die positiven ganzen Zahlen $e$ und $f$, deren Product $g$ bereits bekannt ist, einzeln zu bestimmen. Diese Aufgabe, von allen die Gruppendeterminante betreffenden Fragen die schwierigste, wird in § 9 durch den Satz gelöst, dass stets $e=f=\sqrt{g}$ ist. Nimmt man dazu aus § 12 das Resultat, dass $\chi_\alpha$ und $\chi_{\alpha'}$ conjugirte complexe Grössen sind, so folgt aus (7.), dass die Matrix
\frobtagc{14.}{\left(\sqrt{\frac{h_\alpha}{h}}\,\chi_\alpha^{(\kappa)}\right)}
der conjugirten complexen Matrix complementär ist.

\section*{§ 8.}

In der Gleichung (7.) § 6 gebe ich den Variabelen $y_R$ die Werthe $\chi(R^{-1})=\chi'(R)$, die der Bedingung $y_{BA}=y_{AB}$ nach (13.) § 3 genügen. Dann ist nach (8.) § 5 $\eta=\frac{h}{e}$, und mithin ist
\frobtagc{1.}{\Phi\left(\frac{e}{h}\chi'\right)=1.}

In Folge der Eigenschaft (13.) § 3 ist die Matrix $(\chi')$ mit jeder Matrix $(x)$ vertauschbar. Ich setze
\frobtag{2.}{\frac{h}{e}\xi_A=\sum\chi(R^{-1})x_S=\sum\chi(S^{-1})x_R,\qquad (RS=A)}
und bilde aus den Grössen $\xi_{PQ^{-1}}$ die Matrix $(\xi)$. Dann ist
\[
((x)+u(\varepsilon))\left(\frac{e}{h}\chi'\right)=(\xi)+u\left(\frac{e}{h}\chi'\right)
\]
und mithin nach (9.) § 1
\[
\Phi(x+u\varepsilon)=\Phi\left(\xi+u\frac{e}{h}\chi'\right).
\]
Nach Gleichung (6.) § 6 ist aber
\[
\Phi\left(\xi+u\frac{e}{h}\chi'\right)=\Phi(\xi+u\varepsilon),
\]
weil der den Grössen $y_R=\frac{e}{h}\chi'(R)$ entsprechende Werth von $\eta$ nach (8.) § 5 gleich 1 ist. Folglich ist
\frobtagc{3.}{\Phi(x+u\varepsilon)=\Phi(\xi+u\varepsilon).}
Ist dagegen $\Phi'$ ein von $\Phi$ verschiedener Primfactor der Gruppendeterminante, $f'$ sein Grad und $\Psi$ der entsprechende Charakter, so ist nach (7.) § 6
\[
\Phi'(y_R+u\varepsilon_R)=(u+\eta')^{f'},\qquad (y_{BA}=y_{AB})
\]
wo $f'\eta'=\sum\Psi(R)y_R$ ist. Setzt man also $y_R=\frac{e}{h}\chi'(R)$, so ist nach (8.) § 5 $\eta'=0$. Demnach ist nach (6.) § 6 für $h$ unabhängige Variable $z_R$
\[
\Phi'\left(uz+\frac{e}{h}\chi'\right)=\Phi'(uz)=u^{f'}\Phi'(z).
\]
Sei $(z)=(x)^{-1}$, also $(z)(x)=(\varepsilon)$ und $\Phi(z)\Phi(x)=1$. Dann ist
\[
\Phi'(x)\,\Phi'\left(uz+\frac{e}{h}\chi'\right)=u^{f'},
\]
und daher, weil $(x)(u(z)+(\frac{e}{h}\chi'))=u(\varepsilon)+(\xi)$ ist,
\frobtagc{4.}{\Phi'(\xi+u\varepsilon)=u^{f'}.}
Aus diesen beiden Gleichungen, die nur specielle Fälle einer allgemeineren Formel sind, ergiebt sich nach (7.) § 1
\frobtagc{5.}{\left|\xi_{PQ^{-1}}+u\varepsilon_{PQ^{-1}}\right|=\Phi(x+u\varepsilon)^e\,u^{h-ef}}
und für $x_R=\varepsilon_R$
\frobtagc{6.}{\left|\frac{e}{h}\chi(PQ^{-1})+u\varepsilon_{PQ^{-1}}\right|=(u+1)^f\,u^{h-ef}.}

Daher verschwindet die charakteristische Determinante der Matrix $(\xi)$ für die $f+1$ Werthe $u=0,u_1,u_2,\cdots u_f$. Weil nämlich die Gruppendeterminante $\Theta(x)$ stets den linearen Factor $\sum x_R$, und zwar nur in der ersten Potenz, enthält, so ist nothwendig $h-ef>0$, falls man den trivialen Fall $h=1$ ausschliesst.

Nun sei $G((\xi))=0$ die Gleichung niedrigsten Grades, der die Matrix $(\xi)$ genügt (V. § 1, VI). Dann muss die ganze Function $G(u)$, die ein Theiler der charakteristischen Determinante $\left|\xi_{PQ^{-1}}-u\varepsilon_{PQ^{-1}}\right|$ ist, für jede Wurzel der Matrix $(\xi)$, also für jeden der $f+1$ Werthe $u=0,u_1,u_2,\cdots u_f$, verschwinden und mithin durch $u\Phi(x-u\varepsilon)$ theilbar sein. Da ferner die Matrix $(\chi')$ mit jeder Matrix $(x)$ vertauschbar ist, so ist nach (9.) § 5
\frobtag{7.}{(\xi)^n=\left(\left(\frac{e}{h}\chi'\right)(x)\right)^n=\left(\frac{e}{h}\chi'\right)^n(x)^n=\left(\frac{e}{h}\chi'\right)(x)^n.}
Multiplicirt man daher die Relation (5.) § 4 mit $(x)$, so erhält man
\frobtag{8.}{(\xi)^{f+1}-(\xi)^f\Phi_1+(\xi)^{f-1}\Phi_2-\cdots+(-1)^f(\xi)\Phi_f=0,}
wo nach (3.)
\frobtagc{9.}{\Phi_n=\Phi_n(x)=\Phi_n(\xi)}
ist, oder kürzer
\[
(\xi)\Phi(x-(\xi)\varepsilon)=0.
\]
Mithin ist $u\Phi(x-u\varepsilon)$ durch $G(u)$ theilbar, also gleich $G(u)$, und folglich ist (8.) die Gleichung niedrigsten Grades, der die Matrix $(\xi)$ genügt.

Ohne die Relation (5.) § 4 zu benutzen, kann man diesen Satz auch so einsehen: Die Function $G(u)$ wird erhalten, indem man die Determinante $h^{\mathrm{ten}}$ Grades $(\xi_{PQ^{-1}}-u\varepsilon_{PQ^{-1}})$ durch den grössten gemeinsamen Divisor ihrer Unterdeterminanten $(h-1)^{\mathrm{ten}}$ Grades dividirt. Sind die $h$ Variabelen $x_R$ unabhängig, so folgt aus der Gleichung (1.) § 5, dass die Unterdeterminanten $(h-1)^{\mathrm{ten}}$ Grades von $\Theta(x)$ alle durch $\Phi(x)^{e-1}$ theilbar sind. Daher sind die Unterdeterminanten von $\Theta(\xi-u\varepsilon)$ alle durch $\Phi(\xi-u\varepsilon)^{e-1}=\Phi(x-u\varepsilon)^{e-1}$, also auch durch $(u-u_1)^{e-1}$, theilbar und nicht durch eine höhere Potenz von $u-u_1$, weil sonst $\Theta(\xi-u\varepsilon)$ durch eine höhere als die $e^{\mathrm{te}}$ Potenz von $u-u_1$ theilbar sein müsste. Mittels derselben Sätze ergiebt sich aus den Gleichungen (9.), § 5 und (6.), dass der Rang der Matrix $(\chi)$ gleich $ef$ ist, wie ich Ch. § 5 ausführlicher gezeigt habe. Daher ist auch der Rang der Matrix
\frobtagc{10.}{(\xi)=\left(\frac{e}{h}\chi'\right)(x)=(x)\left(\frac{e}{h}\chi'\right)}
gleich $ef$. Folglich verschwinden $h-ef$ Elementartheiler der charakteristischen Determinante von $(\xi)$ für $u=0$, und weil das Product derselben nach (6.) gleich $u^{h-ef}$ ist, so muss jeder von ihnen linear sein. Mithin enthält der grösste gemeinsame Divisor der Unterdeterminanten $(h-1)^{\mathrm{ten}}$ Grades jener Determinante den Factor $u$ genau in der $(h-ef-1)^{\mathrm{ten}}$ Potenz, und demnach muss
\[
(-1)^fG(u)=u\Phi(x-u\varepsilon)
\]
sein.

\section*{§ 9.}

Nach diesen Vorbereitungen wende ich mich nun zum Beweise des Fundamentalsatzes der Theorie der Gruppendeterminanten:

\emph{Der Exponent der Potenz, worin die Gruppendeterminante einen Primfactor enthält, ist dem Grade des Factors gleich.}

Den Fall $f=1$ habe ich bereits in § 2 erledigt. Wegen der Schwierigkeit des allgemeinen Beweises schicke ich noch die besonderen Fälle $f=2$ und 3 voraus.

Ist $f=2$, so ist
\frobtagc{1.}{2\Phi(x)=S_1^2-S_2,}
und der entsprechende Charakter $\chi$ genügt den Relationen $\chi(A,B,C)=0$ oder
\frobtagc{2.}{\begin{aligned}
&\chi(A)\chi(B)\chi(C)-\chi(A)\chi(BC)-\chi(B)\chi(AC)\nobreak\\
&\qquad{}-\chi(C)\chi(AB)+\chi(ABC)+\chi(ACB)=0.
\end{aligned}}
In dieser Gleichung ersetze ich $B$ durch $BC^{-1}$ und summire dann nach $C$ über die $h$ Elemente von $\Hgrp$ (oder man summire in (2.) über alle Elemente $B,C$, die der Bedingung $BC=B'$ genügen, wo $B'$ ein festes Element ist). Mit Hülfe der Formeln (7.) § 5 und (3.) § 7 findet man
\[
\frac{h}{e}\chi(A)\chi(B)-h\chi(A)\chi(B)-\frac{h}{e}\chi(AB)-\frac{h}{e}\chi(AB)+h\chi(AB)+\frac{h}{f}\chi(A)\chi(B)=0,
\]
also weil $f=2$ ist,
\[
\left(\frac{h}{e}-\frac{h}{f}\right)(\chi(A)\chi(B)-2\chi(AB))=0.
\]
Daher ist $e=f$, weil nicht für je zwei Elemente
\[
\chi(A)\chi(B)-2\chi(AB)=0
\]
sein kann. Denn sonst erhielte man, indem man diese Gleichung mit $x_Ax_B$ multiplicirt und nach $A$ und $B$ summirt, $S_1^2-2S_2=0$. Nach (1.) wäre also $4\Phi=S_1^2$, während $\Phi$ unzerlegbar ist.

Ist $f=3$, so ist
\frobtagc{3.}{6\Phi=S_1^3-3S_1S_2+2S_3,}
und der entsprechende Charakter genügt den Relationen $\chi(A,B,C,D)=0$ ((20.) § 3). Ersetzt man darin $C$ durch $CD^{-1}$ und summirt dann nach $D$, so erhält man
\[
\begin{aligned}
&\frac{h}{e}\bigl(\chi(A)\chi(B)\chi(C)-\chi(B)\chi(AC)-\chi(A)\chi(BC)-\chi(A)\chi(BC)\nobreak\\
&\qquad{}-\chi(B)\chi(AC)-\chi(C)\chi(AB)+\chi(ABC)+\chi(ACB)\nobreak\\
&\qquad{}+\chi(ABC)+\chi(ABC)+\chi(ACB)+\chi(ACB)\bigr)\nobreak\\
&{}+h\bigl(-\chi(A)\chi(B)\chi(C)+\chi(C)\chi(AB)+\chi(A)\chi(BC)+\chi(B)\chi(AC)\nobreak\\
&\qquad{}-\chi(ABC)-\chi(ACB)\bigr)\nobreak\\
&{}+\frac{h}{f}\bigl(\chi(A)\chi(B)\chi(C)+\chi(A)\chi(B)\chi(C)-\chi(B)\chi(AC)-\chi(A)\chi(BC)\nobreak\\
&\qquad{}-\chi(C)\chi(AB)-\chi(C)\chi(AB)\bigr)=0,
\end{aligned}
\]
also wenn man den Factor $h$ durch $3\frac{h}{f}$ ersetzt,
\frobtag{4.}{\begin{aligned}
&\left(\frac{h}{e}-\frac{h}{f}\right)\bigl(\chi(A)\chi(B)\chi(C)-2\chi(A)\chi(BC)-2\chi(B)\chi(AC)\nobreak\\
&\qquad{}-\chi(C)\chi(AB)+3\chi(ABC)+3\chi(ACB)\bigr)=0.
\end{aligned}}
Wäre nun der zweite Factor immer Null, so erhielte man, indem man mit $x_Ax_Bx_C$ multiplicirt und summirt, $S_1^3-5S_1S_2+6S_3=0$, und indem man mittelst dieser Gleichung $S_3$ aus (3.) eliminirt, $9\Phi=S_1(S_1^2-2S_2)$, während $\Phi$ unzerlegbar ist. Daher ist $e=f=3$.

Im allgemeinen Falle genügt der Charakter $\chi$ den Relationen
\[
\chi(A,B,C,\cdots Q,R,S)=0,
\]
wo $A,B,\cdots R,S$ irgend $f+1$ Elemente sind, oder kurz
\[
\sum^{(f+1)}(\Pi-\chi)=0.
\]
In jedem der $(f+1)!$ Glieder dieser Summe ersetze ich $R$ durch $RS^{-1}$ und summire dann nach $S$. Jedes Glied entspricht einer gewissen Permutation von $f+1$ Symbolen, die in cyklische Factoren zerlegt ist. In Bezug auf diese Permutation unterscheide ich drei Fälle:

1. $R$ und $S$ kommen in zwei verschiedenen Cyklen der Permutation vor, z. B.
\[
(-\chi(ABCD\cdots FR))(-\chi(S))(-\chi(G\cdots K))\cdots .
\]
Ersetzt man $R$ durch $RS^{-1}$ und summirt dann nach $S$, so erhält man nach (7.) § 5
\[
-\frac{h}{e}(-\chi(ABCD\cdots FR))(-\chi(G\cdots K))\cdots .
\]
Dasselbe Resultat ergiebt sich in derselben Weise aus dem Gliede
\[
(-\chi(BCD\cdots FR))(-\chi(SA))(-\chi(G\cdots K))\cdots
\]
und aus
\[
(-\chi(CD\cdots FR))(-\chi(SAB))(-\chi(G\cdots K))\cdots
\]
u. s. w. und schliesslich aus
\[
(-\chi(R))(-\chi(SABCD\cdots F))(-\chi(G\cdots K))\cdots,
\]
aber aus keinem anderen Gliede. Ist also $r$ die Anzahl der Elemente $ABCD\cdots FR$, so erhält man auf diesem Wege
\frobtagc{5.}{-\frac{h}{e}\sum^{(f)}r\,(\Pi-\chi).}
Die $f!$ Glieder dieser Summe sind in analoger Weise wie die der Summe (19.) § 3 aus den $f!$ Permutationen der $f$ Elemente $A,B,C,\cdots Q,R$ gebildet. Nur erhält in dieser Summe, worin das Element $R$ bevorzugt ist, jedes Glied $(\Pi-\chi)$ noch einen Zahlenfactor $r$. Dieser ist gleich der Anzahl der Elemente des Cyklus, worin $R$ vorkommt (vergl. (4.)).

2. $R$ und $S$ kommen beide in demselben Cyklus der Permutation vor, und zwar folgt $S$ in dem Cyklus unmittelbar auf $R$ (es kann also auch $S$ das erste und $R$ das letzte Element des Cyklus sein), z. B.
\[
(-\chi(AB\cdots FRS))(-\chi(G\cdots K))\cdots .
\]
Ersetzt man $R$ durch $RS^{-1}$ und summirt dann nach $S$, so erhält man
\[
h(-\chi(AB\cdots FR))(-\chi(G\cdots K))\cdots,
\]
und zwar jedes Glied nur einmal, also im Ganzen
\frobtagc{6.}{h\sum^{(f)}(\Pi-\chi).}
3. $R$ und $S$ kommen beide in demselben Cyklus vor, ohne dass $S$ unmittelbar auf $R$ folgt, z. B.
\[
(-\chi(A\cdots RBCD\cdots FS))(-\chi(G\cdots K))((-\chi(L\cdots N)))\cdots .
\]
Ersetzt man $R$ durch $RS^{-1}$ und summirt dann nach $S$, so erhält man nach (3.) § 7
\[
-\frac{h}{f}(-\chi(A\cdots R))(-\chi(BCD\cdots F))(-\chi(G\cdots K))(-\chi(L\cdots N))\cdots .
\]
Dasselbe Resultat ergiebt sich in derselben Weise aus dem Gliede
\[
(-\chi(A\cdots RCD\cdots FBS))(-\chi(G\cdots K))(-\chi(L\cdots N))\cdots,
\]
das durch cyklische Vertauschung der zwischen $R$ und $S$ stehenden Elemente $BCD\cdots F$ aus dem obigen hervorgeht. Die Anzahl der cyklischen Vertauschungen, die man so ausführen kann, ist gleich der Anzahl der Elemente $BCD\cdots F$. Ferner ergiebt sich dasselbe Resultat aus dem Gliede
\[
(-\chi(A\cdots RG\cdots KS))(-\chi(BCD\cdots F))(-\chi(L\cdots N))\cdots .
\]
Auch hier kann man noch die zwischen $R$ und $S$ stehenden Elemente $G\cdots K$ cyklisch vertauschen, was auf so viele Arten möglich ist, wie die Anzahl der Elemente $G\cdots K$ beträgt.

Dasselbe Resultat ergiebt sich aus dem Gliede
\[
(-\chi(A\cdots RL\cdots NS))(-\chi(BCD\cdots F))(-\chi(G\cdots K))\cdots
\]
u. s. w., im Ganzen also auf so viele Arten, wie die Anzahl der Elemente $BCD\cdots\allowbreak FG\cdots\allowbreak KL\cdots\allowbreak N\cdots$ beträgt und nicht auf mehr Arten, also auf $f-r$ Arten, wenn, wie oben, die Anzahl der Elemente $A\cdots R$ mit $r$ bezeichnet wird. Demnach erhält man
\frobtagc{7.}{-\frac{h}{f}\sum^{(f)}(f-r)(\Pi-\chi).}
Vereinigt man die drei Ausdrücke (5.), (6.) und (7.), so ergiebt sich die Gleichung
\frobtagc{8.}{\left(\frac{h}{f}-\frac{h}{e}\right)\left(\sum^{(f)}r\,(\Pi-\chi)\right)=0.}
Mithin ist $e=f$, wenn man zeigen kann, dass nicht für je $f$ Elemente $AB\cdots QR$
\[
\sum^{(f)}r\,(\Pi-\chi)=0
\]
ist. Multiplicirt man diese Gleichung mit $x_Ax_B\cdots x_Qx_R$ und summirt nach jedem der $f$ Elemente $A,B,\cdots Q,R$, so erhält man eine Relation zwischen $S_1,S_2,\cdots S_f$. Diese ist nicht identisch (ohne Rücksicht auf die Bedeutung von $S_1,S_2,\cdots S_f$) erfüllt, da sie in Bezug auf $S_f$ linear ist und der Coefficient von $S_f$ eine nicht verschwindende ganze Zahl ist. Ich werde aber zeigen, dass $S_1,S_2,\cdots S_f$ $f$ unabhängige Functionen der $h$ Variabelen $x_R$ sind, dass also zwischen ihnen keine Relation besteht, deren Coefficienten von den Variabelen $x_R$ unabhängig sind. Daraus folgt dann, dass auch $\Phi_1,\Phi_2,\cdots\Phi_f$ unabhängig sind, und ebenso $u_1,u_2,\cdots u_f$.

Bestände zwischen den $f$ Functionen $S_1,S_2,\cdots S_f$ der $h$ unabhängigen Variabelen $x_R$ eine Gleichung, so würde sich, indem man sie nach $x_R$ differentiirt, eine Relation der Form
\[
\Psi_1\frac{\partial S_1}{\partial x_R}+\Psi_2\frac{\partial S_2}{\partial x_R}+\cdots+\Psi_f\frac{\partial S_f}{\partial x_R}=0
\]
ergeben, wo $\Psi_1,\Psi_2,\cdots\Psi_f$ ganze Functionen der $h$ Variabelen sind, die von $R$ unabhängig sind. Nun ist aber
\[
S_n=\sum_{R_1,R_2,R_3,\cdots R_n}\chi(R_1R_2R_3\cdots R_n)x_{R_1}x_{R_2}x_{R_3}\cdots x_{R_n}
\]
und folglich
\[
\begin{aligned}
\frac{\partial S_n}{\partial x_R}={}&\sum_{R_2,R_3,\cdots R_n}\chi(RR_2R_3\cdots R_n)x_{R_2}x_{R_3}\cdots x_{R_n}\nobreak\\
&{}+\sum_{R_1,R_3,\cdots R_n}\chi(R_1RR_3\cdots R_n)x_{R_1}x_{R_3}\cdots x_{R_n}+\cdots\nobreak\\
&{}+\sum_{R_1,R_2,\cdots R_{n-1}}\chi(R_1R_2\cdots R_{n-1}R)x_{R_1}x_{R_2}\cdots x_{R_{n-1}},
\end{aligned}
\]
also da $\chi(ABC\cdots F)$ bei cyklischer Vertauschung der Elemente ungeändert bleibt,
\frobtagc{9.}{\frac{\partial S_n}{\partial x_R}=n\sum_{R_1,R_2,\cdots R_{n-1}}\chi(RR_1R_2\cdots R_{n-1})x_{R_1}x_{R_2}\cdots x_{R_{n-1}}}
oder nach (6.) § 1 und (7.) § 8
\frobtagc{10.}{\frac{\partial S_n}{\partial x_R}=n\sum_S\chi(RS)x_S^{(n-1)}=\frac{nh}{e}\xi_{R^{-1}}^{(n-1)},}
wo $\xi_R^{(n)}$ aus den Grössen $\xi_R$ in derselben Weise gebildet ist wie $x_R^{(n)}$ aus den Grössen $x_R$. Speciell ist
\frobtagc{11.}{\frac{\partial S_n}{\partial x_E}=nS_{n-1}.}
Setzt man diese Ausdrücke in die obige Relation ein, so erhält man, falls man noch $R$ durch $R^{-1}$ ersetzt,
\[
\Psi_1\frac{e}{h}\chi(R^{-1})+2\Psi_2\xi_R+3\Psi_3\xi_R^{(2)}+\cdots+f\Psi_f\xi_R^{(f-1)}=0.
\]
Setzt man $R=PQ^{-1}$, so wird dies eine Gleichung zwischen Matrizen, die, mit $(x)$ multiplicirt, lautet
\[
(\xi)\Psi_1+2(\xi)^2\Psi_2+3(\xi)^3\Psi_3+\cdots+f(\xi)^f\Psi_f=0.
\]
Ich habe aber in § 8 gezeigt, dass die Gleichung niedrigsten Grades, der die Matrix $(\xi)$ genügt, vom Grade $f+1$ ist. Der zweite Factor des Ausdrucks (8.) kann also nicht für jedes System von $f$ Elementen verschwinden, und mithin muss
\frobtagc{12.}{e=f}
sein.

\section*{§ 10.}

Sind $x_R$ und $y_R$ zwei Systeme von je $h$ Variabelen, so ist nach § 6 die Matrix $(x_{PQ^{-1}})$ mit der Matrix $(y_{Q^{-1}P})$ vertauschbar, und folglich ist die Determinante
\frobtagc{1.}{\left|ux_{PQ^{-1}}+vy_{Q^{-1}P}+w\varepsilon_{PQ^{-1}}\right|}
ein Product von $h$ linearen Functionen der drei Variabelen $u,v,w$ von der Form $ua_\alpha+vb_\alpha+w$. Hier sind $a_1,a_2,a_3\cdots$ die $h$ Wurzeln der Matrix $(x_{PQ^{-1}})$ und $b_1,b_2,b_3,\cdots$ die der Matrix $(y_{Q^{-1}P})$ oder, was nach § 6 dasselbe ist, der Matrix $(y_{PQ^{-1}})$. Es fragt sich nun, in welcher Weise die Wurzeln dieser beiden Matrizen einander zugeordnet werden müssen, damit $ua_\alpha+vb_\alpha+w$ ein Linearfactor der Determinante (1.) sei. Ich setze
\[
\sum_R x_R\chi(RS^{-1})=\frac{h}{e}\xi_S,
\qquad
\sum_R y_R\chi(RS^{-1})=\frac{h}{e}\eta_S,
\]
ferner
\[
(x_{PQ^{-1}})=(x),\qquad (\xi_{PQ^{-1}})=(\xi),\qquad (y_{Q^{-1}P})=(\overline{y}),\qquad (\eta_{Q^{-1}P})=(\overline{\eta}),
\]
wobei immer $P$ die Zeilen und $Q$ die Spalten der Matrix bezeichnet.

\clearpage
Dann ist
\[
\left(\frac{e}{h}\chi'\right)(x)=(x)\left(\frac{e}{h}\chi'\right)=(\xi),
\qquad
\left(\frac{e}{h}\chi'\right)(\overline{y})=(\overline{y})\left(\frac{e}{h}\chi'\right)=(\overline{\eta}).
\]
Die $k$ verschiedenen Charaktere $\chi$ unterscheide ich durch obere Indices $\chi^{(\varkappa)}$ $(\varkappa=0,1,\cdots k-1)$. Die dem Charakter $\chi^{(\varkappa)}$ entsprechenden Matrizen $(\xi)$ und $(\overline{\eta})$ bezeichne ich mit $(\xi^{(\varkappa)})$ und $(\overline{\eta}^{(\varkappa)})$. Sind dann $\varkappa$ und $\lambda$ verschieden, so folgt aus (9.) § 5
\frobtagc{2.}{(\xi^{(\varkappa)})(\xi^{(\lambda)})=0,\quad (\overline{\eta}^{(\varkappa)})(\overline{\eta}^{(\lambda)})=0,\quad (\xi^{(\varkappa)})(\overline{\eta}^{(\lambda)})=0,\quad (\overline{\eta}^{(\varkappa)})(\xi^{(\lambda)})=0.}
Nach Gleichung (5.) § 7 ist
\[
\sum_\varkappa\left(\frac{e^{(\varkappa)}}{h}\chi'\right)=(\varepsilon)
\]
und mithin
\[
\sum_\varkappa(\xi^{(\varkappa)})=(x),\qquad \sum_\varkappa(\overline{\eta}^{(\varkappa)})=(\overline{y}).
\]
Entwickelt man also das Product der $k$ Matrizen
\[
\Pi_\varkappa\left(u(\xi^{(\varkappa)})+v(\overline{\eta}^{(\varkappa)})+w(\varepsilon)\right)
\]
nach Potenzen von $w$, so erhält man
\[
w^k(\varepsilon)+w^{k-1}\bigl(u(x)+v(\overline{y})\bigr),
\]
während die übrigen Glieder nach (2.) verschwinden. Zwischen den Determinanten dieser Matrizen ergiebt sich daher die Beziehung
\frobtag{3.}{\Pi_\varkappa\left|u\xi_{PQ^{-1}}^{(\varkappa)}+v\eta_{Q^{-1}P}^{(\varkappa)}+w\varepsilon_{PQ^{-1}}\right|=w^{h(k-1)}\left|ux_{PQ^{-1}}+vy_{Q^{-1}P}+w\varepsilon_{PQ^{-1}}\right|}
(vergl. die analoge Entwicklung V. S. 610). Irgend einer der $k$ Factoren der linken Seite sei
\frobtagc{4.}{\left|u\xi_{PQ^{-1}}+v\eta_{Q^{-1}P}+w\varepsilon_{PQ^{-1}}\right|.}
Wie die rechte Seite zeigt, ist diese Determinante gleich einer Potenz von $w$, multiplicirt mit einer Anzahl der linearen Factoren $ua_\alpha+vb_\alpha+w$ der Determinante (1.). Andererseits kann man die Determinante (4.) als einen speciellen Fall der Determinante (1.) betrachten: Die Wurzeln der Matrix $(\xi)$ sind nach (5.) § 8 die $f$ Wurzeln $u_1,u_2,\cdots u_f$ der Gleichung $\Phi(x-u\varepsilon)=0$, jede $e$ Mal gezählt, und ausserdem $(h-ef)$ Mal gezählt die Zahl $0$. Ebenso sind die Wurzeln der Matrix $(\overline{\eta})$, die Zahl $0$ und die Wurzeln $v_1,v_2,\cdots v_f$ der Gleichung $\Phi(y-v\varepsilon)=0$. Daher ist die Determinante (4.) ein Product von linearen Factoren $au+bv+w$, wo $a$ eine der $f+1$ Grössen $0,u_1,u_2,\cdots u_f$ und $b$ eine der $f+1$ Grössen $0,v_1,v_2,\cdots v_f$ ist. Eine Combination, wie $a=u_1,b=0$, kann aber, wie die rechte Seite der Gleichung (3.) zeigt, nicht vorkommen. Abgesehen von einer Potenz von $w$ enthält daher die Determinante (4.) nur noch lineare Factoren der Form $uu_\alpha+vv_\beta+w$, worin $u_\alpha$ eine der $f$ Grössen $u_1,u_2,\cdots u_f$ und $v_\beta$ eine der $f$ Grössen $v_1,v_2,\cdots v_f$ ist.

Nun ist
\[
u(\xi)+v(\overline{\eta})=\left(\frac{e}{h}\chi'\right)\bigl(u(x)+v(\overline{y})\bigr),
\]
und daher hat diese Matrix den Rang $ef$. Mithin enthält die Determinante den Factor $w$ mindestens in der Potenz $h-ef$, aber auch in keiner höheren, weil dies nach (5.) § 8 nicht einmal für $v=0$ der Fall ist. Die übrigen $ef$ Linearfactoren sind demnach alle von der Form $uu_\alpha+vv_\beta+w$. Sei $uu_1+vv_1+w$ einer derselben. Betrachtet man die $h$ Grössen $y_R$, also die $f$ Grössen $v_\beta$ als constant, so hat die Determinante (4.), als Function von $w$ betrachtet, mit der irreducibelen Function $\Phi(ux_E+vv_1+w,ux_A,ux_B,\cdots)$ den Linearfactor $uu_1+vv_1+w$ gemeinsam. Folglich hat sie alle Factoren $uu_\alpha+vv_1+w$ $(\alpha=1,2,\cdots f)$ mit ihr gemeinsam. Ebenso erkennt man, dass die Determinante die $f^2$ linearen Functionen
\[
uu_\alpha+vv_\beta+w\qquad (\alpha,\beta=1,2,\cdots f)
\]
sämmtlich enthält, und jeden gleich oft. Kommt jeder Factor $m$ Mal vor, so ist $ef=mf^2$, also $e=fm$. Auf diese Weise kann man daher, ohne das Resultat des § 9 zu benutzen, nachweisen, dass $e$ durch $f$ theilbar ist. Nach diesem ist aber $e=f$, und mithin ist $m=1$, also
\frobtag{5.}{\left|u\xi_{PQ^{-1}}+v\eta_{Q^{-1}P}+w\varepsilon_{PQ^{-1}}\right|=w^{h-ef}\prod_{\alpha,\beta}'(uu_\alpha+vv_\beta+w).}
Durch diese Betrachtung ist nun die Art bestimmt, wie man die Wurzeln der beiden Matrizen $(x)$ und $(y)$ einander zuordnen muss, um die linearen Factoren der Determinante (1.) zu erhalten. Sind $\Phi$ und $\Phi'$ zwei verschiedene Primfactoren von $\Theta$, so ist jede Wurzel der Gleichung $\Phi(x-w\varepsilon)=0$ mit jeder Wurzel der Gleichung $\Phi(y-w\varepsilon)=0$ zu combiniren, aber mit keiner Wurzel der Gleichung $\Phi'(y-w\varepsilon)=0$. Die Allgemeinheit der erhaltenen Formel wird nicht vermindert, wenn man $u=v=1$ und $w=0$ setzt. Ist
\frobtagc{6.}{\Psi(x,y)=\prod_{\alpha,\beta}'(u_\alpha+v_\beta)}
die Resultante der beiden Functionen $\Phi(x-\varepsilon w)$ und $\Phi(y+\varepsilon w)$ der Variabelen $w$, so ist
\frobtagc{7.}{\left|x_{PQ^{-1}}+y_{Q^{-1}P}\right|=\Pi\Psi(x,y).}
Könnte man direct beweisen, dass diese Determinante, als Function der $2h$ unabhängigen Variabelen $x_R,y_R$ betrachtet, keinen mehrfachen Factor besitzt, so wäre damit für die Gleichung $e=f$ ein neuer Beweis geliefert. Setzt man die $h$ Grössen $y_R=0$, so wird $\Psi(x,y)=\Phi(x)^f$. Auf diesem Wege erlangt man eine tiefere Einsicht in den Grund der merkwürdigen Erscheinung, dass die Gruppendeterminante jeden Primfactor in einer Potenz enthält, deren Exponent dem Grade des Factors gleich ist.

Von besonderem Interesse ist der specielle Fall, wo $y_R=-x_R$ ist. Dann ist die Determinante
\frobtagc{8.}{\left|x_{PQ^{-1}}-x_{Q^{-1}P}+w\varepsilon_{PQ^{-1}}\right|=w^s\Pi\Psi,}
wo
\frobtagc{9.}{\Psi=\prod_{\alpha>\beta}\left(w-(u_\beta-u_\alpha)^2\right)}
und
\frobtagc{10.}{s=\sum f}
ist. Der Gleichung $\Psi=0$ genügen also die Quadrate der Differenzen der Wurzeln der Gleichung $\Phi(x-w\varepsilon)=0$. Ich will nun zeigen, dass die für $w=0$ verschwindenden Elementartheiler jener Determinante alle linear sind, oder, was dasselbe ist, dass der Rang $r$ der Matrix
\frobtagc{11.}{(x_{PQ^{-1}}-x_{Q^{-1}P})}
gleich
\frobtagc{12.}{r=h-s}
ist. Nach Formel (8.) ist $r\ge h-s$. Da die beiden Matrizen $(x)$ und $(x)'$ mit einander vertauschbar sind, so ist
\[
x_A^{(n+1)}=\sum x_Rx_S^{(n)}=\sum x_R^{(n)}x_S\qquad (RS=A)
\]
oder
\[
x_A^{(n+1)}=\sum_R x_{AR^{-1}}x_R^{(n)}=\sum_R x_R^{(n)}x_{R^{-1}A},
\]
also
\[
\sum_R\left(x_{AR^{-1}}-x_{R^{-1}A}\right)x_R^{(n)}=0.
\]
Setzt man für $A$ der Reihe nach alle $h$ Elemente von $\Hgrp$, so ist
\frobtagc{13.}{\sum_R\left(x_{AR^{-1}}-x_{R^{-1}A}\right)y_R=0}
ein System von $h$ linearen Gleichungen zwischen den $h$ Unbekannten $y_R$. Der Rang der von ihren Coefficienten gebildeten Matrix ist $r$. Mithin bildet das vollständige System ihrer Lösungen eine Matrix vom Range $h-r$, und der Rang irgend eines Systems ihrer Lösungen, z. B. des Systems
\[
y_R=x_R^{(n)}\qquad (n=0,1,2,\cdots)
\]
ist $\le h-r$. Enthält die charakteristische Determinante $\left|x_{PQ^{-1}}-u\varepsilon_{PQ^{-1}}\right|$ der Matrix $(x)$ irgend einen Linearfactor $u-u_i$ in der $e^{\mathrm{ten}}$ Potenz, so enthält ihn nach § 8 der grösste gemeinsame Divisor ihrer Unterdeterminanten $(h-1)^{\mathrm{ten}}$ Grades in der $(e-1)^{\mathrm{ten}}$ Potenz. Folglich ist die Gleichung niedrigsten Grades, der die Matrix $(x)$ genügt, vom Grade $\sum f=s$, nämlich $G((x))=0$, wo
\[
(-1)^sG(u)=\Pi\Phi(x-u\varepsilon)
\]
ist. Daher sind die Matrizen $(x)^0,(x)^1,\cdots(x)^{s-1}$ linear unabhängig, also sind auch die $s$ Lösungen
\frobtagc{14.}{y_R=x_R^{(n)}\qquad (n=0,1,2,\cdots s-1)}
unabhängig, und mithin ist $s\le h-r$. Folglich ist $r=h-s$, und die Grössen (14.) bilden ein vollständiges System unabhängiger Lösungen der $h$ linearen Gleichungen (13.), unter denen $r$ unabhängig sind.

\section*{§ 11.}

Nach den Gleichungen (9.) § 8 ist
\frobtagc{1.}{\Phi(x)=\Phi(\xi),\quad \Phi_n(x)=\Phi_n(\xi),\quad S_n(x)=S_n(\xi).}
also
\frobtag{2.}{n!\Phi_n(x)=\sum_{R_1,R_2,\cdots R_n}\chi(R_1,R_2,\cdots R_n)\xi_{R_1}\xi_{R_2}\cdots\xi_{R_n}}
und
\frobtag{3.}{S_n(x)=\sum_{R_1,R_2,\cdots R_n}\chi(R_1R_2\cdots R_n)\xi_{R_1}\xi_{R_2}\cdots\xi_{R_n}.}
Eine andere Darstellung ergiebt sich aus den Formeln (10.) § 9, nämlich
\frobtag{4.}{\frac{e}{h}S_n(x)=\xi_E^{(n)}=\sum_{R_1,R_2,\cdots R_n}\xi_{R_1}\xi_{R_2}\cdots\xi_{R_n}\qquad (R_1R_2\cdots R_n=E)}
Demnach lassen sich die Functionen $S_n$ und $\Phi_n$ und speciell $\Phi$ selbst durch die $h$ Variabelen
\frobtagc{5.}{\xi_S=\frac{e}{h}\sum_R\chi(RS^{-1})x_R}
ausdrücken, unter denen nur $ef$ unabhängig sind, weil nach § 8 der Rang der Matrix, die von den Coefficienten dieser $h$ linearen Functionen der $h$ Variabelen $x_R$ gebildet wird, gleich $ef$ ist. Führt man diese Umformung für jeden Primfactor von $\Theta$ aus, so wird die Gruppendeterminante durch
\frobtagc{6.}{\sum ef=h}
neue Variabele ausgedrückt.

Man transformire jede der $k$ Primfunctionen $\Phi,\Phi',\cdots$ einzeln durch eine lineare Substitution in eine Function von möglichst wenig neuen Variabelen. Ist ihre Anzahl für $\Phi,\Phi',\cdots$ gleich $g,g',\cdots$, so ist $g\le ef$, $g'\le e'f'$, $\cdots$. Es könnte dann sein, dass sich die Functionen $\Phi,\Phi',\cdots$ insgesamt durch noch weniger als $g+g'+\cdots$ neue Variabele darstellen liessen, lineare Verbindungen der $h$ unabhängigen Variabelen $x_R$. Wäre dies der Fall, oder wäre $g<ef$, oder $g'<e'f'$, $\cdots$, so liesse sich $\Theta$ durch eine lineare Substitution in eine Function von weniger als $h$ Variabelen transformiren. Nun ist aber
\[
\sum_R\frac{\partial l\Theta}{\partial x_R}x_{RA^{-1}}=\varepsilon_A h.
\]
Differentiirt man diese Gleichung nach $x_{B^{-1}}$, so erhält man
\[
\sum_R\frac{\partial^2 l\Theta}{\partial x_R\partial x_{B^{-1}}}x_{RA^{-1}}=-\frac{\partial l\Theta}{\partial x_{B^{-1}A}}=-\frac{h}{\Theta}\Theta_{B^{-1},A^{-1}}
\]
und mithin
\[
\left|\frac{\partial^2 l\Theta}{\partial x_P\partial x_{Q^{-1}}}\right|\left|x_{PQ^{-1}}\right|=\left(-\frac{h}{\Theta}\right)^h\left|\Theta_{P,Q}\right|,
\]
also
\frobtagc{7.}{\left|\frac{\partial^2 l\Theta}{\partial x_P\partial x_{Q^{-1}}}\right|=\frac{(-h)^h}{\Theta^2}}
Könnte man aber $\Theta$ durch eine lineare Substitution in eine Function von weniger als $h$ Variabelen transformiren, so müsste diese Determinante verschwinden. Folglich lässt sich $\Phi$ durch $ef$, aber nicht durch weniger als $ef$ Variabele ausdrücken, die lineare Verbindungen der $h$ Variabeln $x_R$ sind, und die $ef$ Variabeln von $\Phi$, die $e'f'$ von $\Phi',\cdots$ sind alle von einander unabhängig.

In einer besonders einfachen Weise lässt sich $\Phi^e$ durch die Variabeln $\xi_R$ darstellen. Dazu benutze ich den folgenden Determinantensatz (vergl. meine Arbeit \emph{Über das Pfaff'sche Problem}, Crelle's Journal Bd. 82; § 4, I):

\emph{Ist $r$ der Rang der Matrix}
\[
a_{\alpha\beta},\qquad (\alpha,\beta=1,2,\cdots n)
\]
\emph{so verhalten sich die Determinanten $r^{\mathrm{ten}}$ Grades, die sich aus den Elementen von $r$ Spalten dieser Matrix bilden lassen, wie die entsprechenden Determinanten $r^{\mathrm{ten}}$ Grades, die sich aus den Elementen von irgend $r$ anderen Spalten dieser Matrix bilden lassen.}

Dabei heissen zwei Determinanten entsprechende, wenn zu ihrer Bildung dieselben Zeilen und zwar in derselben Reihenfolge benutzt sind. Derselbe Satz gilt, wenn man die Zeilen und die Spalten vertauscht. Er lässt sich so verallgemeinern:

\emph{Ist die Matrix $(c_{\alpha\beta})$ aus den beiden Matrizen $(a_{\alpha\beta})$ und $(b_{\alpha\beta})$ zusammengesetzt, und ist $r$ der Rang der Matrix $(a_{\alpha\beta})$, so verhalten sich die Determinanten $r^{\mathrm{ten}}$ Grades, die sich aus den Elementen von $r$ Spalten der Matrix $(a_{\alpha\beta})$ bilden lassen, wie die entsprechenden Determinanten $r^{\mathrm{ten}}$ Grades, die sich aus den Elementen von irgend $r$ Spalten der Matrix $(c_{\alpha\beta})$ bilden lassen.}

Der Rang der Matrix
\[
c_{\alpha\gamma}=a_{\alpha1}b_{1\gamma}+a_{\alpha2}b_{2\gamma}+\cdots+a_{\alpha n}b_{n\gamma}
\]
ist dann höchstens gleich $r$. Eine Bedeutung hat aber dieser Satz nur, wenn der Rang von $(c_{\alpha\gamma})$ gleich $r$ ist. Dies muss der Fall sein, wenn die Determinante $n^{\mathrm{ten}}$ Grades $(b_{\alpha\beta})$ von Null verschieden ist.

Mit dem Zeichen
\[
\left|\,a\begin{array}{cccc}
\alpha_1&\alpha_2&\cdots&\alpha_r\\
\beta_1&\beta_2&\cdots&\beta_r
\end{array}\,\right|
\]
bezeichne ich die Determinante $r^{\mathrm{ten}}$ Grades, gebildet aus den Elementen der Zeilen $\alpha_1,\alpha_2,\cdots\alpha_r$ und der Spalten $\beta_1,\beta_2,\cdots,\beta_r$ der Matrix $(a_{\alpha\beta})$, in der angegebenen Reihenfolge. Dann ist also das Verhältniss
\[
\left|\,a\begin{array}{cccc}
\rho_1&\rho_2&\cdots&\rho_r\\
\alpha_1&\alpha_2&\cdots&\alpha_r
\end{array}\,\right|:
\left|\,c\begin{array}{cccc}
\rho_1&\rho_2&\cdots&\rho_r\\
\beta_1&\beta_2&\cdots&\beta_r
\end{array}\,\right|
\]
von der Wahl der Indices $\rho_1,\rho_2,\cdots\rho_r$ unabhängig.

Ist wieder $(c_{\alpha\beta})=(a_{\alpha\beta})(b_{\alpha\beta})$, und ist $r$ der Rang der Matrix $(b_{\alpha\beta})$, so verhalten sich die Determinanten $r^{\mathrm{ten}}$ Grades, die sich aus den Elementen von $r$ Zeilen der Matrix $b_{\alpha\beta}$ bilden lassen, wie die entsprechenden Determinanten $r^{\mathrm{ten}}$ Grades, die sich aus den Elementen von irgend $r$ Zeilen der Matrix $(c_{\alpha\beta})$ bilden lassen.

Diesen Satz wende ich auf die Matrix
\frobtagc{8.}{(\xi)=\left(\frac{e}{h}\chi'\right)(x)=(x)\left(\frac{e}{h}\chi'\right)}
an. Der Rang der Matrix $(\chi')$ ist $g=ef$, ebenso der der Matrix $(\xi)$. Daher gilt der obige Satz sowohl für die Zeilen, wie für die Spalten, und es ist
\[
\left|\,\xi\begin{array}{cccc}
A_1&A_2&\cdots&A_g\\
B_1&B_2&\cdots&B_g
\end{array}\,\right|
=\left(\frac{e}{h}\right)^g
\left|\,\chi'\begin{array}{cccc}
A_1&A_2&\cdots&A_g\\
B_1&B_2&\cdots&B_g
\end{array}\,\right|\Psi,
\]
wo $\Psi$ von der Wahl der Elemente $A_1,A_2,\cdots A_g,B_1,B_2,\cdots B_g$ unabhängig ist. Vergleicht man in den Relationen
\[
\left|\xi_{PQ^{-1}}+u\varepsilon_{PQ^{-1}}\right|=\Phi(x+u\varepsilon)^e u^{h-g},\qquad
\left|\frac{e}{h}\chi'(PQ^{-1})+u\varepsilon_{PQ^{-1}}\right|=(u+1)^g u^{h-g}
\]
die Coefficienten von $u^{h-g}$, so ergiebt sich: Die Summe aller Hauptunterdeterminanten $r^{\mathrm{ten}}$ Grades ist für die Matrix $(\xi)$ gleich $\Phi(x)^e$ und für die Matrix $\left(\frac{e}{h}\chi'\right)$ gleich 1. Folglich ist $\Psi=\Phi^e$, also
\frobtag{9.}{\left|\,\xi\begin{array}{cccc}
A_1&A_2&\cdots&A_g\\
B_1&B_2&\cdots&B_g
\end{array}\,\right|
=\left(\frac{e}{h}\right)^g
\left|\,\chi\begin{array}{cccc}
B_1&B_2&\cdots&B_g\\
A_1&A_2&\cdots&A_g
\end{array}\,\right|\Phi(x)^e.}
Wählt man die $2g$ Elemente $A_1,B_1,\cdots A_g,B_g$ so, dass die hierin auftretende Determinante $g^{\mathrm{ten}}$ Grades der Matrix $(\chi)$ von Null verschieden ist, so ist die entsprechende Determinante der Matrix $(\xi)$ bis auf einen constanten Factor gleich $\Phi^e$. Unter den von Null verschiedenen Determinanten $r^{\mathrm{ten}}$ Grades dieser beiden Matrizen giebt es auch Hauptunterdeterminanten (worin $B_1=A_1,\cdots B_g=A_g$ ist), weil die Summe aller Hauptunterdeterminanten $g^{\mathrm{ten}}$ Grades nicht verschwindet.

In derselben Weise ergiebt sich die allgemeinere Formel
\frobtag{10.}{\left|\xi_{PQ^{-1}}+\eta_{Q^{-1}P}\right|=\left(\frac{e}{h}\right)^g\left|\chi(QP^{-1})\right|\Psi(x,y),\qquad \left(\begin{array}{l}P=A_1,A_2,\cdots A_g\\ Q=B_1,B_2,\cdots B_g\end{array}\right)}
wo $\Psi(x,y)$ dieselbe Bedeutung hat, wie in der Gleichung (6.) § 10.

\section*{§ 12.}

Die Ermittlung der $k$ Primfactoren, worin die Gruppendeterminante zerfällt, ist auf die Bestimmung der Constanten $\chi_\alpha^{(\varkappa)}$ zurückgeführt, die von der Auflösung einer Gleichung $k^{\mathrm{ten}}$ Grades abhängt. Ich will nun die algebraische und arithmetische Natur dieser Grössen näher untersuchen. Zunächst bestimme ich den algebraischen Körper, dem diese Zahlen angehören.

Sei $\Ggrp$ eine Untergruppe von $\Hgrp$, $g$ ihre Ordnung, sei $h=gm$ und $H$ die zu $\Ggrp$ gehörige Gruppendeterminante. Seien $E,A,B,\cdots$ die $g$ Elemente von $\Ggrp$, und $L,M,\cdots$ die nicht in $\Ggrp$ enthaltenen Elemente von $\Hgrp$. Setzt man dann in $\Theta$ die Variabelen $x_L,x_M,\cdots$ alle gleich Null, so wird, wie ich \emph{Ch.} § 7, (10.) gezeigt habe,
\frobtagc{1.}{\Theta=H^m.}
Daher sind die Coefficienten derjenigen Glieder von $\Theta$, die nur von $x_E,x_A,x_B,\cdots$ abhängen, den Coefficienten der entsprechenden Glieder von $H^m$ gleich.

Ist $\Ggrp$ eine commutative Gruppe, so ist $H$ ein Product von $g$ linearen Factoren
\[
H=\left(\sum\psi_1(R)x_R\right)\left(\sum\psi_2(R)x_R\right)\cdots,
\]
und die Charaktere $\psi_1(R),\psi_2(R),\cdots$ sind alle vom ersten Grade, also Einheitswurzeln. Speciell ist $\psi_1(E)=\psi_2(E)=\cdots=1$. Ist demnach $\Phi$ ein Primfactor $f^{\mathrm{ten}}$ Grades von $\Theta$, so wird diese Function, wenn man darin $x_L=x_M=\cdots=0$ setzt, gleich dem Producte von $f$ dieser linearen Factoren, etwa
\frobtagc{2.}{\Phi=\left(\sum\psi_1(R)x_R\right)\left(\sum\psi_2(R)x_R\right)\cdots\left(\sum\psi_f(R)x_R\right).}
Ersetzt man $x_E$ durch $x_E-u$, so erkennt man, dass die $f$ Wurzeln der Gleichung $\Phi(x-u\varepsilon)=0$, falls man $x_L=x_M=\cdots=0$ setzt, ganze lineare Functionen der Variabelen $x_E,x_A,x_B,\cdots$ werden, etwa
\frobtagc{3.}{u_\lambda=\sum_R\psi_\lambda(R)x_R,\qquad (\lambda=1,2,\cdots f)}
deren Coefficienten Einheitswurzeln sind. Ist also $A$ eins der Elemente von $\Ggrp$, so ist der Coefficient von $x_E^{f-1}x_A$ in $\Phi(x)$ gleich
\frobtagc{4.}{\chi(A)=\psi_1(A)+\psi_2(A)+\cdots+\psi_f(A),}
und diese Gleichung gilt auch für $A=E$. Sind ferner $A$ und $B$ zwei Elemente von $\Ggrp$, so ist, weil $\psi_\lambda$ ein Charakter ersten Grades ist, $\psi_\lambda(AB)=\psi_\lambda(A)\psi_\lambda(B)$ und mithin
\frobtagc{5.}{\chi(AB)=\psi_1(A)\psi_1(B)+\psi_2(A)\psi_2(B)+\cdots+\psi_f(A)\psi_f(B).}
Daher hat die Matrix $(\chi(PQ^{-1}))$, die für die Gruppe $\Hgrp$ den Rang $ef$ hat, höchstens noch den Rang $f$, falls $P$ und $Q$ nur die Elemente von $\Ggrp$ durchlaufen.

Ist $A$ irgend ein Element von $\Hgrp$, und $m$ seine Ordnung, so bilden die Potenzen von $A$ eine commutative Gruppe der Ordnung $m$. Wählt man diese für $\Ggrp$, so werden die $m$ Charaktere $\psi_1(A)=\rho_1$, $\psi_2(A)=\rho_2,\cdots$ alle $m^{\mathrm{ten}}$ Wurzeln der Einheit. Mithin ist
\frobtagc{6.}{\chi(A)=\rho_1+\rho_2+\cdots+\rho_f,\qquad \chi(A^n)=\rho_1^n+\rho_2^n+\cdots+\rho_f^n}
für jeden Werth von $n$. Für $n=m-1$ folgt daraus, dass $\chi(A)$ und $\chi(A^{-1})$ conjugirt complexe Grössen sind (\emph{Ch.} § 3).

\emph{Ist $A$ ein Element der $m^{\mathrm{ten}}$ Ordnung und $\chi$ ein Charakter $f^{\mathrm{ten}}$ Grades, so lässt sich $\chi(A)$ als eine Summe von $f$ $m^{\mathrm{ten}}$ Wurzeln der Einheit darstellen.}

Dieselben sind einzeln dadurch bestimmt, dass $\chi(A^n)$ gleich der Summe ihrer $n^{\mathrm{ten}}$ Potenzen ist. Setzt man in $\Phi$ alle Variabelen gleich Null ausser $x_E,x_A,x_{A^2},\cdots x_{A^{m-1}}$, so wird nach (2.)
\[
\Phi=(x_E+\rho_1x_A+\rho_1^2x_{A^2}+\cdots+\rho_1^{m-1}x_{A^{m-1}})\cdots(x_E+\rho_fx_A+\rho_f^2x_{A^2}+\cdots+\rho_f^{m-1}x_{A^{m-1}}).
\]
Setzt man daher in $\Phi$ alle Variabelen gleich Null, ausser $x_E$ und $x_A$, so wird
\frobtagc{7.}{\Phi(x_E,x_A,0,0,\cdots)=(x_E+\rho_1x_A)(x_E+\rho_2x_A)\cdots(x_E+\rho_fx_A)}
und mithin nach (9.) § 4
\frobtagc{8.}{u^f+\zchar_1(A)u^{f-1}+\zchar_2(A)u^{f-2}+\cdots+\zchar_f(A)=(u+\rho_1)(u+\rho_2)\cdots(u+\rho_f).}
Speciell ist, da $\zchar_f(A)=\zchar(A)$ ist,
\frobtagc{9.}{\zchar(A)=\rho_1\rho_2\cdots\rho_f.}
Wendet man die Formel (1.) auf die Gruppe $\Ggrp$ an, die von den Potenzen von $A$ gebildet wird, so erhält man
\[
\Theta=\left(\Pi_\rho(x_E+\rho x_A+\rho^2x_{A^2}+\cdots+\rho^{m-1}x_{A^{m-1}})\right)^{\frac{h}{m}},
\]
wo $\rho$ alle $m^{\mathrm{ten}}$ Wurzeln der Einheit durchläuft. Setzt man also in $\Theta$ alle Variabelen gleich Null, ausser $x_E$ und $x_A$, so wird
\frobtagc{10.}{\Theta(x_E,x_A,0,0,\cdots)=(x_E^m+(-x_A)^m)^{\frac{h}{m}}.}
Folglich ist, wie ich in § 3 auf anderem Wege gezeigt habe, der Coefficient von $x_A^h$ in $\Theta$ gleich
\frobtagc{11.}{\Pi\zchar(A)^e=(-1)^{h-\frac{h}{m}}.}
Sei $B$ ein zweites Element von $\Hgrp$, und seien $\sigma_1,\sigma_2,\cdots\sigma_f$ die $f$ dem Charakter $\chi(B)$ entsprechenden Einheitswurzeln, also
\frobtagc{12.}{\chi(B)=\sigma_1+\sigma_2+\cdots+\sigma_f,\qquad \chi(B^n)=\sigma_1^n+\sigma_2^n+\cdots+\sigma_f^n.}
Sind nun $A$ und $B$ mit einander vertauschbar, so erzeugen sie eine Gruppe $\Ggrp$. Daher kann man die Formel (5.) anwenden und erkennt: Die Einheitswurzeln der beiden Systeme lassen sich einander so zuordnen, dass
\frobtagc{13.}{\chi(AB)=\rho_1\sigma_1+\cdots+\rho_f\sigma_f,\\ \chi(A^rB^s)=\rho_1^r\sigma_1^s+\cdots+\rho_f^r\sigma_f^s}
wird. Setzt man in der Primfunction $\Phi$, die dem Charakter $\chi$ entspricht, alle Variabelen gleich Null, ausser $x_E,x_A$ und $x_B$, so erhält man nach (2.)
\frobtagc{14.}{\Phi(x_E,x_A,x_B,0,0,\cdots)=(x_E+\rho_1x_A+\sigma_1x_B)\cdots(x_E+\rho_fx_A+\sigma_fx_B),}
wodurch zugleich die Zuordnung der Einheitswurzeln bestimmt ist.

Um für die entwickelten Sätze ein Beispiel zu geben, betrachte ich ein invariantes Element $B$ der Gruppe $\Hgrp$, d. h. ein solches, das mit jedem Elemente $R$ von $\Hgrp$ vertauschbar ist. In der Formel (3.) § 7 ist dann $R^{-1}BR=B$, und mithin, falls $A$ irgend ein anderes Element von $\Hgrp$ ist,
\[
\chi(A)\chi(B)=f\chi(AB).
\]
Alle invarianten Elemente von $\Hgrp$ bilden eine commutative Gruppe $\Ggrp$. Setzt man für jedes Element $G$ derselben $\chi(G)=f\psi(G)$, so ist demnach für je zwei Elemente $A$ und $B$ von $\Ggrp$ $\psi(A)\psi(B)=\psi(AB)$. Mithin ist $\psi(G)$ ein Charakter von $\Ggrp$, also eine Einheitswurzel $\rho$. Ferner ist $\psi(G^n)=\psi(G)^n$, also $\chi(G)=f\rho$ und $\chi(G^n)=f\rho^n$. Für ein invariantes Element $G$ von $\Hgrp$ sind folglich die $f$ in der Formel (6.) auftretenden Einheitswurzeln alle einander gleich.

Zu demselben Resultate führt die Bemerkung, dass ein invariantes Element $A$ für sich allein eine Classe conjugirter Elemente bildet. Setzt man daher in $\Phi$ alle Variabelen gleich Null ausser $x_E$ und $x_A$, so wird $\Phi$ nach Formel (7.) § 6 die $f^{\mathrm{te}}$ Potenz einer linearen Function von $x_E$ und $x_A$, und folglich ist nach Formel (7.) $\rho_1=\rho_2=\cdots=\rho_f$.

Nun sei wieder $A$ ein beliebiges Element von $\Hgrp$, und sei $m$ seine Ordnung. Ist $\rho$ eine primitive $m^{\mathrm{te}}$ Wurzel der Einheit, so sind die Grössen $\rho_1,\rho_2,\cdots\rho_f$ in der Formel (6.) alle Potenzen von $\rho$, und daher ist $\chi(A)$ eine Zahl des Körpers $K(\rho)$, der von allen rationalen Functionen von $\rho$ gebildet wird. Unter den mit $A$ conjugirten Elementen der Gruppe $\Hgrp$ können sich auch Potenzen von $A$ befinden $A^r,A^s,A^t,\cdots$. Ihre Exponenten sind zu $m$ theilerfremd und bilden eine Gruppe, d. h. einer unter ihnen ist $\equiv rs\pmod m$, wenn $r$ und $s$ irgend zwei von ihnen sind. Da $A$ und $A^r$ conjugirt sind, so ist $\chi(A)=\chi(A^r)$, also $\rho_1+\cdots+\rho_f=\rho_1^r+\cdots+\rho_f^r$. Drückt man $\chi(A)$ durch $\rho$ aus, so bleibt demnach diese Zahl ungeändert, falls man $\rho$ durch $\rho^r$ ersetzt. Diejenigen Zahlen des Körpers $K(\rho)$, die ungeändert bleiben, falls man $\rho$ durch $\rho^r$ oder $\rho^s$, oder $\rho^t,\cdots$ ersetzt, bilden einen Körper $\Lambda(\rho)$, einen Divisor von $K(\rho)$. Die Zahl $\chi(A)$ gehört folglich diesem Körper $\Lambda(\rho)$ an.

Ist z. B. $\Hgrp$ die symmetrische Gruppe des Grades $n$, also $h=n!$, so ist $A$ mit jeder Potenz $A^r$ conjugirt (ähnlich), deren Exponent $r$ zu $m$ theilerfremd ist. Daher sind die Charaktere der symmetrischen Gruppe sämmtlich ganze rationale Zahlen (vergl. die Beispiele $n=4$ und 5, \emph{Ch.} § 8).

In dem Körper $\Lambda(\rho)$ ist $\chi(A)$ als Summe von Einheitswurzeln eine ganze algebraische Zahl. Eine solche ist aber auch jeder Coefficient von $\Phi$, also auch von $\Phi_n$. Denn wenn in einem Producte $\Theta=\Phi\Psi$ von zwei ganzen Functionen von beliebig vielen Variabelen, deren Coefficienten algebraische Zahlen sind, alle Coefficienten ganze algebraische Zahlen sind, so ist auch das Product aus jedem Coefficienten von $\Phi$ und jedem von $\Psi$ eine ganze algebraische Zahl (vergl. Dedekind, \emph{Über einen arithmetischen Satz von Gauss}. Prager Math. Ges. 1892). Sind $A,B,C,\cdots$ verschiedene Elemente von $\Hgrp$ und ist $r+s+t+\cdots=n$, so hat $x_A^rx_B^sx_C^t\cdots$ in $\Phi_n$ den Coefficienten
\frobtagc{15.}{\frac{1}{r!s!t!\cdots}\chi(A,\cdots A,B,\cdots B,C,\cdots C,\cdots).}
Folglich ist dieser Ausdruck eine ganze algebraische Zahl. Wendet man denselben Satz auf die Factoren des Productes (9.) § 6 an, so erkennt man, dass auch
\frobtagc{16.}{\frac{h_\alpha\chi_\alpha}{f}}
eine ganze algebraische Zahl ist. Folglich ist auch nach (6.) § 7
\[
\sum\frac{h_\alpha\chi_\alpha}{f}\chi_{\alpha'}=\frac{h}{e}
\]
eine ganze Zahl. Daher ist die Zahl $e=f$ ein Divisor der Ordnung $h$.

\begin{center}
\vspace{1.0em}
\rule{0.32\linewidth}{0.4pt}\\[0.35em]
Ausgegeben am 10. December.\\[1.8em]
{\scriptsize Berlin, gedruckt in der Reichsdruckerei.}
\end{center}
% ===== END WORK 054 =====

\clearpage

% ===== BEGIN WORK 056: 056_ber_die_Darstellung_der_endlichen_Gruppen =====
\begin{center}
{\Large Über die Darstellung der endlichen Gruppen\\durch lineare Substitutionen.}\par
\vspace{1.2em}
Von G. \textsc{Frobenius}.
\end{center}
\vspace{1em}

In meiner Arbeit \emph{Über die Primfactoren der Gruppendeterminante} (Sitzungsberichte 1896; im Folgenden mit \emph{Pr.} citirt) habe ich jeder endlichen Gruppe $\Hgrp$ der Ordnung $h$ eine Matrix des Grades $h$ zugeordnet, deren Elemente von $h$ Variabelen abhängen. Ihre Wichtigkeit beruht auf dem Zusammenhange, worin sie selbst und die Primfactoren ihrer Determinante mit den linearen Substitutionen stehen, durch die sich die Gruppe $\Hgrp$ und die ihr isomorphen Gruppen darstellen lassen. Aus jeder solchen Darstellung kann man eine zur Gruppe $\Hgrp$ gehörige Matrix ableiten, deren Determinante in einer Potenz der Gruppendeterminante enthalten ist (§ 2). Ist die Determinante unzerlegbar, also einem Primfactor der Gruppendeterminante gleich, so nenne ich die Darstellung eine primitive. Umgekehrt entspricht jedem Primfactor $f^{\ten}$ Grades der Gruppendeterminante eine, und, abgesehen von der Wahl der Variabelen, nur eine primitive Darstellung der Gruppe durch Substitutionen von $f$ Variabelen (§ 4).

Die Gruppenmatrix kann in eine ähnliche Matrix transformirt werden, die in Theilmatrizen zerfällt. Benutzt man dabei allein die $k$ Gruppencharaktere, so kann man sie in $k$ Matrizen zerlegen, deren jede die $f^{\mathrm{te}}$ Potenz einer Primfunction $\Phi$ des $f^{\ten}$ Grades zur Determinante hat (§ 3). Benutzt man aber höhere Irrationalitäten, so kann man sie in $\sum f$ Matrizen zerlegen, deren jede eine Primfunction $\Phi$ selbst zur Determinante hat (§ 5). Mit Hülfe einiger merkwürdigen Sätze über Determinanten $n^{\ten}$ Grades, deren Elemente $n^2$ unabhängige Variable sind (§ 7), zeige ich dann: man kann die Transformation so einrichten, dass je $f$ Theilmatrizen, deren Determinanten demselben Primfactor $f^{\ten}$ Grades $\Phi$ gleich sind, einander identisch gleich sind. Dann sind die Elemente aller Theilmatrizen zusammen $\sum f^2=h$ von einander unabhängige Variable. Aus einer solchen Theilmatrix, deren Determinante $\Phi$ ist, ergeben sich $h$ lineare Substitutionen, die eine mit $\Hgrp$ isomorphe Gruppe bilden. Der Isomorphismus kann auch ein meroedrischer sein. Dies hängt von einer besonderen Beziehung ab, worin der Charakter $\chi$ der Gruppe $\Hgrp$, welcher der Primfunction $\Phi$ entspricht, zu einer invarianten Untergruppe von $\Hgrp$ stehen kann (§ 1). Die primitiven Darstellungen einer Gruppe durch lineare Substitutionen werfen ein neues Licht auf die Bedeutung der Relationen, aus denen die Charaktere der Gruppe und damit die Coefficienten der Primfactoren der Gruppendeterminante berechnet werden (§ 6). Die wichtigsten Ergebnisse dieser Arbeit habe ich \textsc{Dedekind}, dem ich die Anregung zu diesen Untersuchungen verdanke, im April dieses Jahres mitgetheilt.

\section*{§ 1.}

Unter den Charakteren einer Gruppe $\Hgrp$ nimmt der Hauptcharakter, dessen Werthe alle gleich $1$ sind, eine bevorzugte Stellung ein. Ist $\Hgrp$ zusammengesetzt, und ist $\Ggrp$ eine invariante Untergruppe von $\Hgrp$, so giebt es ferner gewisse Charaktere von $\Hgrp$, die in Bezug auf $\Ggrp$ ein besonderes Verhalten zeigen. Ein solcher Charakter hat nämlich für je zwei Elemente von $\Hgrp$, die mod. $\Ggrp$ aequivalent sind, denselben Werth. Ich sage daher, er gehöre zur Gruppe $\frac{\Hgrp}{\Ggrp}$ auf Grund des Satzes:

\begin{quote}\itshape
Ist $\Ggrp$ eine invariante Untergruppe der Gruppe $\Hgrp$, so ist jeder Charakter von $\frac{\Hgrp}{\Ggrp}$ auch ein Charakter von $\Hgrp$.
\end{quote}

Dies ist so zu verstehen: Bilden $A,B,C,\cdots$ ein vollständiges Restsystem von $\Hgrp$ (mod. $\Ggrp$), so ist
\[
\Hgrp=A\Ggrp+B\Ggrp+C\Ggrp+\cdots.
\]
Dann können die Complexe $A\Ggrp=\Ggrp A$, $B\Ggrp=\Ggrp B$, $\ldots$ als die Elemente der Gruppe $\frac{\Hgrp}{\Ggrp}$ aufgefasst werden. Sei $\psi$ ein Charakter dieser Gruppe, und $\psi(A\Ggrp)$ sein Werth für das Element $A\Ggrp$. Jedes Element $R$ der Gruppe $\Hgrp$ gehört einem und nur einem dieser Complexe an. Für alle Elemente $R$ des Complexes $A\Ggrp$ sei $\chi(R)=\psi(A\Ggrp)$. Die so definirten $h$ Grössen $\chi(R)$ bilden einen Charakter von $\Hgrp$. Sei umgekehrt $\chi(R)$ ein Charakter von $\Hgrp$, der den gleichen Werth hat für je zwei Elemente von $\Hgrp$, die mod. $\Ggrp$ aequivalent sind. Setzt man dann für jedes Element $R$ von $\Hgrp$, das aequivalent $A$ (mod. $\Ggrp$) ist, also dem Complexe $A\Ggrp$ angehört, $\psi(A\Ggrp)=\chi(R)$, so ist $\psi$ ein Charakter von $\frac{\Hgrp}{\Ggrp}$. Die beiden entsprechenden Charaktere haben denselben Grad
\[
f=\chi(E)=\psi(\Ggrp).
\]
Gehört der Charakter $\chi$ zu der Gruppe $\frac{\Hgrp}{\Ggrp}$, so gehört er auch zu $\frac{\Hgrp}{\Ggrp'}$, wo $\Ggrp'$ eine in $\Ggrp$ enthaltene invariante Untergruppe von $\Hgrp$ ist.

Damit nämlich $h$ Grössen $\chi(R)$ einen Charakter von $\Hgrp$ bilden, sind die folgenden Bedingungen nothwendig und hinreichend:
\begin{equation}\tag{1.}
\chi(E)=f
\end{equation}
ist eine positive (ganze) Zahl. Für je zwei Elemente $A$ und $B$ von $\Hgrp$ ist
\begin{equation}\tag{2.}
\chi(AB)=\chi(BA)
\end{equation}
und
\begin{equation}\tag{3.}
h\chi(A)\chi(B)=f\sum_R\chi(AR^{-1}BR).
\end{equation}
Endlich ist
\begin{equation}\tag{4.}
h=\sum_R\chi(R)\chi(R^{-1}).
\end{equation}
Nach der zweiten Bedingung hat nämlich $\chi(R)$ für alle Elemente einer Classe conjugirter Elemente denselben Werth, etwa für die Elemente der $\rho^{\ten}$ Classe den Werth $\chi_\rho$ $(\rho=0,1,\cdots k-1)$. Daher sind die Gleichungen (3.) und (4.) identisch mit (vergl. \emph{Pr.} § 7)
\begin{equation}\tag{3.}
h_\alpha h_\beta\chi_\alpha\chi_\beta=f\sum_\gamma h_\gamma c_{\alpha\beta\gamma}\chi_\gamma,
\end{equation}
\begin{equation}\tag{4.}
h=\sum_\rho h_\rho\chi_\rho\chi_{\rho'}.
\end{equation}
Jede Lösung der Gleichungen (3.) ergiebt die Verhältnisse der Werthe eines Charakters $\chi$. Dann liefert die Gleichung (4.) das Quadrat des noch unbestimmt gebliebenen Proportionalitätsfactors und endlich die Gleichung (1.) sein Vorzeichen.

Nun ist $\chi(E)=\psi(\Ggrp)=f$. Ferner ist $A\Ggrp B\Ggrp=AB\Ggrp$. Ist also $R$ ein Element des Complexes $A\Ggrp$, und $S$ ein Element von $B\Ggrp$, so ist $RS$ ein Element von $AB\Ggrp$. Daher ist $\chi(RS)=\psi(AB\Ggrp)=\psi(A\Ggrp B\Ggrp)$ und $\chi(SR)=\psi(B\Ggrp A\Ggrp)$ und mithin, weil $\psi$ die Eigenschaft (2.) besitzt, $\chi(RS)=\chi(SR)$. Ferner ist
\[
\frac{h}{g}\psi(A\Ggrp)\psi(B\Ggrp)=f\sum_S\psi(AS^{-1}BS\Ggrp),
\]
wo $S$ ein vollständiges Restsystem von $\Hgrp$ (mod. $\Ggrp$) durchläuft. Nun ist $AS^{-1}BS\Ggrp=A\Ggrp S^{-1}BS\Ggrp$ und bleibt daher ungeändert, wenn man $S$ durch $SG$ ersetzt, wo $G$ irgend ein Element von $\Ggrp$ ist. Daher ist auch
\[
h\psi(A\Ggrp)\psi(B\Ggrp)=f\sum_R\psi(AR^{-1}BR\Ggrp),
\]
wo $R$ alle Elemente von $\Hgrp$ durchläuft, also auch
\[
h\chi(A)\chi(B)=f\sum_R\chi(AR^{-1}BR).
\]
Ebenso folgt aus der Relation
\[
\frac{h}{g}=\sum_S\psi(S\Ggrp)\psi(S^{-1}\Ggrp)
\]
die Gleichung
\[
h=\sum_R\chi(R)\chi(R^{-1}).
\]
Daher ist $\chi(R)$ ein Charakter von $\Hgrp$.

Umgekehrt sei $\chi(R)$ ein Charakter von $\Hgrp$, und sei stets $\chi(R)=\chi(S)$, wenn $R\sim S$ (mod. $\Ggrp$) ist. Dann erkennt man auf demselben Wege, dass $\psi(R\Ggrp)=\chi(R)$ ein Charakter von $\frac{\Hgrp}{\Ggrp}$ ist.

Der Hauptcharakter gehört zu der Gruppe $\frac{\Hgrp}{\Hgrp}$, die Charaktere ersten Grades zu der commutativen Gruppe $\frac{\Hgrp}{\Ggrp}$, falls $\Ggrp$ die Commutatorgruppe von $\Hgrp$ ist.

Durchläuft $G$ die Elemente der Gruppe $\Ggrp$, so sei
\begin{equation}\tag{5.}
\sum_G x_{RG}=y_{R\Ggrp}.
\end{equation}
Dann ist unter den obigen Voraussetzungen (vergl. \emph{Pr.} § 8)
\[
\frac{h}{f}\xi_A=\sum_R\chi(R)x_{RA}=\sum_S\psi(S\Ggrp)y_{SA\Ggrp}=\frac{h:g}{f}\eta_{A\Ggrp}.
\]
Der Rang der Matrix $(\xi_{PQ^{-1}})$ ist gleich $f^2$. Ist $\Phi(x)$ der Primfactor der Gruppendeterminante $\Theta(x)$ von $\Hgrp$, der dem Charakter $\chi$ entspricht, so ist jede Unterdeterminante des Grades $f^2$ von jener Matrix bis auf einen constanten Factor gleich $\Phi(x)^f$. Ist $\Psi(y)$ der Primfactor der Gruppendeterminante $H(y)$ von $\frac{\Hgrp}{\Ggrp}$, der dem Charakter $\psi$ entspricht, so ist jede Unterdeterminante des Grades $f^2$ von der Matrix $\eta_{PQ^{-1}\Ggrp}$ gleich $\Psi(y)^f$. Nun ist $g\xi_A=\eta_{A\Ggrp}$, und mithin ist unter der Voraussetzung (5.) $\Phi(x)=\Psi(y)$. Da $\Theta(x)$ und $H(y)$ die Factoren $\Phi(x)$ und $\Psi(y)$ genau in der $f^{\ten}$ Potenz enthalten, so ist $\Theta(x)$ durch $H(y)$ theilbar, und der Quotient ist zu $H(y)$ theilerfremd (vergl. \emph{Pr.} § 2).

\section*{§ 2.}

Eine endliche Anzahl linearer Substitutionen
\begin{align*}
(A)\qquad x_\alpha&=a_{\alpha1}y_1+a_{\alpha2}y_2+\cdots+a_{\alpha n}y_n,\qquad (\alpha=1,2,\cdots n)\\
(B)\qquad x_\alpha&=b_{\alpha1}y_1+b_{\alpha2}y_2+\cdots+b_{\alpha n}y_n,
\end{align*}
\[
\cdot\quad\cdot\quad\cdot\quad\cdot\quad\cdot\quad\cdot\quad\cdot\quad\cdot\quad\cdot\quad\cdot,
\]
deren Determinanten von Null verschieden sind, bilden eine Gruppe $\Hgrp'$, wenn die aus irgend zwei von ihnen, $(A)$ und $(B)$, zusammengesetzte Substitution $(C)=(A)(B)$ ebenfalls unter ihnen enthalten ist. Die Coefficienten von $(C)$ sind
\[
c_{\alpha\beta}=a_{\alpha1}b_{1\beta}+a_{\alpha2}b_{2\beta}+\cdots+a_{\alpha n}b_{n\beta}.
\]
Man kann auch nach dem Vorgange von \textsc{Gauss} von der Bezeichnung der Variabelen ganz absehen und unter $(A),(B),(C),\cdots$ die Matrizen $n^{\ten}$ Grades verstehen, die von den Coefficientensystemen $a_{\alpha\beta},b_{\alpha\beta},c_{\alpha\beta},\cdots$ gebildet werden.

Sei $\Hgrp$ eine abstracte Gruppe, $A,B,C,\cdots$ ihre Elemente. Ordnet man dem Elemente $A$ die Matrix $(A)$, dem Elemente $B$ die Matrix $(B)$, u. s. w. zu, so sei die Gruppe $\Hgrp'$ der Gruppe $\Hgrp$ isomorph, d. h. es sei $(A)(B)=(AB)$. Dann sage ich, dass die Substitutionen oder die Matrizen $(A),(B),(C),\cdots$ die Gruppe $\Hgrp$ darstellen. Der Isomorphismus kann auch ein meroedrischer sein. Dann ist $\Hgrp'$ einer Gruppe $\frac{\Hgrp}{\Ggrp}$ holoedrisch isomorph, wo $\Ggrp$ eine gewisse invariante Untergruppe von $\Hgrp$ ist, gebildet von den Elementen von $\Hgrp$, denen das Hauptelement $(E)$ von $\Hgrp'$ entspricht. Diese Substitution $(E)$ entspricht dem Hauptelement $E$ von $\Hgrp$. Da $E^2=E$ ist, so ist auch $(E)^2=(E)$. Mithin ist $(E)$ die identische Substitution, da keine andere Substitution von nicht verschwindender Determinante dieser Bedingung genügt.

Ist $(P)$ irgend eine Matrix $n^{\ten}$ Grades von nicht verschwindender Determinante, so stellen auch die Matrizen $(P)^{-1}(A)(P),(P)^{-1}(B)(P),\cdots$ die Gruppe $\Hgrp$ dar. Die entsprechenden Substitutionen gehen aus den ursprünglichen hervor, indem jedes System von $n$ Variabelen $x_1,\cdots x_n;y_1,\cdots y_n;\cdots$ der Substitution $(P)$ unterworfen wird. Zwei solche Darstellungen von $\Hgrp$ bezeichne ich als aequivalent. Alle Darstellungen, die einer gegebenen aequivalent sind, bilden eine Classe aequivalenter Darstellungen von $\Hgrp$.

Den $h$ Elementen $A,B,C,\cdots$ von $\Hgrp$ ordne ich $h$ unabhängige Variable $x_A,x_B,x_C,\cdots$ zu, und ich setze dann aus den Matrizen $(A),(B),(C),\cdots$ die Matrix
\[
(A)x_A+(B)x_B+(C)x_C+\cdots=\sum_R(R)x_R
\]
zusammen, deren Elemente lineare Functionen der unabhängigen Variabelen sind. Jede auf diese Art einer Darstellung von $\Hgrp$ entsprechende Matrix nenne ich eine zur Gruppe $\Hgrp$ gehörige Matrix.

Sei $y_A,y_B,y_C,\cdots$ ein zweites System von $h$ unabhängigen Variabelen, und sei (\emph{Pr.} § 1)
\begin{equation}\tag{1.}
z_T=\sum x_Ry_S,\qquad (RS=T)
\end{equation}
Dann ist das Product der beiden Matrizen
\begin{equation}\tag{2.}
\left(\sum(R)x_R\right)\left(\sum(S)y_S\right)=\sum(RS)x_Ry_S=\sum(T)z_T.
\end{equation}
Umgekehrt charakterisirt diese Eigenschaft die Matrix $\sum(R)x_R$ als eine zur Gruppe $\Hgrp$ gehörige Matrix. Eine solche ist daher z. B. die Gruppenmatrix $(x_{PQ^{-1}})$. Ist nun die Determinante $\left|\sum(R)x_R\right|=F(x)$, so ist $F(x)F(y)=F(z)$. Nach \emph{Pr.} § 1 folgt daraus, dass $F(x)$ in einer Potenz der Gruppendeterminante
\[
\Theta(x)=\left|x_{PQ^{-1}}\right|=\Phi^f{\Phi'}^{f'}{\Phi''}^{f''}\cdots
\]
aufgeht, also ein Product von Potenzen der Primfactoren $\Phi,\Phi',\Phi'',\cdots$ von $\Theta$ ist
\begin{equation}\tag{3.}
\left|\sum(R)x_R\right|=\Phi^s{\Phi'}^{s'}{\Phi''}^{s''}\cdots,
\end{equation}
deren Exponenten $s,s',s'',\cdots$ auch zum Theil Null sein können.

Ist $\Phi$ ein Primfactor $f^{\ten}$ Grades von $\Theta$, so ist $\Theta$ genau durch die $f^{\mathrm{te}}$ Potenz von $\Phi$ theilbar. Nach \emph{Pr.} § 5 ist die Unterdeterminante $(h-1)^{\ten}$ Grades von $\Theta$, die dem Elemente $x_{PQ^{-1}}$ complementär ist, gleich
\[
\frac1h\frac{\partial\Theta}{\partial x_{PQ^{-1}}}.
\]
Mithin enthält der grösste gemeinsame Divisor aller Unterdeterminanten $(h-1)^{\ten}$ Grades von $\Theta$ den Primfactor $\Phi$ genau in der $(f-1)^{\ten}$ Potenz. Nach den bekannten Eigenschaften der Elementartheiler einer Determinante enthält daher der grösste gemeinsame Divisor aller Unterdeterminanten $(h-m)^{\ten}$ Grades von $\Theta$ den Factor $\Phi$ genau in der $(f-m)^{\ten}$ Potenz $(m<f)$, und die Unterdeterminanten $(h-f)^{\ten}$ Grades von $\Theta$ sind nicht alle durch $\Phi$ theilbar. Die charakteristische Determinante $\Theta(x-u\eps)$ der Gruppenmatrix hat demnach, als Function von $u$, lauter lineare Elementartheiler, und falls man für $u$ eine Wurzel der Gleichung $\Phi(x-u\eps)=0$ setzt, hat jene Determinante den Rang $h-f$.

Setzt man in der Formel (1.) $y_R=\dfrac{1}{h\Phi^{f-1}}\dfrac{\partial\Theta}{\partial x_R}$, so wird $z_R=\dfrac{\Theta}{\Phi^{f-1}}\eps_R$, also weil $\sum(R)\eps_R=(E)$ ist.
\begin{equation}\tag{4.}
\left(\sum(R)x_R\right)\left(\sum(S)y_S\right)=\frac{\Theta}{\Phi^{f-1}}(E).
\end{equation}
Nun sind die Elementartheiler der Determinante der letzteren Matrix, die gleich Potenzen von $\Phi$ sind, alle gleich der ersten Potenz von $\Phi$. Auf Grund von (4.) sind aber diese Elementartheiler durch die entsprechenden von $\left|\sum(R)x_R\right|$ theilbar (\emph{Über die Elementartheiler der Determinanten}, Satz IX, Sitzungsberichte 1894). Mithin sind auch die Elementartheiler der Determinante $\left|\sum x_R(R)\right|$, die Potenzen von $\Phi$ sind, alle gleich der ersten Potenz von $\Phi$, und ihre Anzahl ist gleich $s$. Ersetzt man $x_R$ durch $x_R-u\eps_R$, so erkennt man, dass die charakteristische Determinante jener Matrix $\left|\left(\sum x_R(R)\right)-u(E)\right|$ lauter lineare Elementartheiler hat.

\section*{§ 3.}

Zerfällt eine zur Gruppe $\Hgrp$ gehörige Matrix in Theilmatrizen, so ist nach Formel (2.) § 2 jede einzelne von ihnen auch eine zur Gruppe gehörige Matrix. Das Product der Determinanten der Theilmatrizen ist gleich der Determinante der ganzen Matrix. Wie ich zeigen werde, kann man eine der Gruppenmatrix aequivalente Matrix finden, die in $\sum f$ Theilmatrizen zerfällt. Ihre Determinanten sind den $\sum f$ Primfactoren von $\Theta$ gleich. Weiter lässt sich die Zerlegung der Gruppenmatrix unmöglich treiben. Ehe ich aber zur Darstellung dieser Transformation übergehe, will ich zunächst eine andere einfachere durchführen, bei der die Determinante jeder Theilmatrix gleich der $f^{\ten}$ Potenz eines Primfactors $f^{\ten}$ Grades $\Phi$ ist. Mithin haben die einzelnen Theilmatrizen lauter verschiedene Determinanten, von denen je zwei theilerfremd sind. Um diese Zerlegung durchzuführen, reicht die Kenntniss der Charaktere aus.

Sei $\Phi$ ein Primfactor $f^{\ten}$ Grades von $\Theta$, und sei $\chi$ der entsprechende Charakter. Dann hat die Matrix $h^{\ten}$ Grades $(\chi(PQ^{-1}))$ den Rang $r=f^2$, und es giebt darin eine von Null verschiedene Hauptunterdeterminante $r^{\ten}$ Grades (\emph{Pr.} § 11). Eine solche möge erhalten werden, indem man für $P$ und $Q$ die $r$ Elemente $A_1,A_2,\ldots,A_r$ setzt. Sei $\psi$ ein von $\chi$ verschiedener Charakter, $s$ der Rang der Matrix $(\psi(PQ^{-1}))$, und es möge darin eine von Null verschiedene Hauptunterdeterminante $s^{\ten}$ Grades erhalten werden, indem man $P$ und $Q$ gleich $B_1,B_2,\ldots,B_s$ setzt. Diese $s$ Elemente können auch ganz oder zum Theil mit $A_1,A_2,\ldots,A_r$ übereinstimmen. Trifft man eine solche Bestimmung für jeden der $k$ Charaktere von $\Hgrp$, so ist die Summe der $k$ Zahlen
\[
r+s+\cdots=\sum r=\sum f^2=h.
\]
Nun sei $M$ die Matrix $h^{\ten}$ Grades, deren $h$ Zeilen aus der Zeile
\[
\chi(PA_1^{-1}),\ldots,\chi(PA_r^{-1}),\psi(PB_1^{-1}),\ldots,\psi(PB_s^{-1}),\ldots
\]
erhalten werden, indem man für $P$ die $h$ Elemente von $\Hgrp$ setzt. Ferner sei $L'$ die Matrix, deren Zeilen in analoger Weise aus
\[
\chi(A_1Q^{-1}),\ldots,\chi(A_rQ^{-1}),\psi(B_1Q^{-1}),\ldots,\psi(B_sQ^{-1}),\ldots
\]
hervorgehen. Vertauscht man in $L'$ Zeilen und Spalten, so erhält man die conjugirte Matrix $L$. Ich bilde nun die componirte Matrix $LM$. Sind $\alpha$ und $\beta$ zwei der Indices $1,2,\ldots,r$, so ist darin das $\beta^{\te}$ Element der $\alpha^{\ten}$ Zeile
\[
\sum_R\chi(A_\alpha R^{-1})\chi(RA_\beta^{-1})=\frac{h}{f}\chi(A_\alpha A_\beta^{-1}).
\]
Die aus diesen $r^2$ Elementen gebildete Determinante $r^{\ten}$ Grades ist nach Voraussetzung von Null verschieden. Ist ferner $\alpha$ einer der Indices $1,2,\ldots,r$ und $\beta$ einer der Indices $1,2,\ldots,s$, so ist das $(r+\beta)^{\te}$ Element der $\alpha^{\ten}$ Zeile
\[
\sum_R\chi(A_\alpha R^{-1})\psi(RB_\beta^{-1})=0.
\]
Folglich zerfällt die Matrix $LM$ in $k$ Theilmatrizen, und ihre Determinante ist
\[
\prod\left|\frac{h}{f}\chi(A_\alpha A_\beta^{-1})\right|,
\]
also von Null verschieden. Mithin sind die Determinanten $|L|$ und $|M|$ beide von Null verschieden.

Nun sei $X$ die Gruppenmatrix $(x_{PQ^{-1}})$ und $Y=LXM$. Sind $\alpha$ und $\beta$ zwei der Zahlen $1,2,\ldots,r$, so ist darin das $\beta^{\te}$ Element der $\alpha^{\ten}$ Zeile
\[
\sum_{R,S}\chi(A_\alpha R^{-1})x_{RS^{-1}}\chi(SA_\beta^{-1}).
\]
Weil aber $\chi(PQ)=\chi(QP)$ ist, so sind die beiden Matrizen $(x_{PQ^{-1}})$ und $(\chi(PQ^{-1}))$ mit einander vertauschbar (\emph{Pr.} § 6), und jene Summe ist gleich
\[
\sum_{R,S}\chi(A_\alpha R^{-1})\chi(RS^{-1})x_{SA_\beta^{-1}}=\frac{h}{f}\sum_S\chi(A_\alpha S^{-1})x_{SA_\beta^{-1}}.
\]
Ist $\alpha$ eine der Zahlen $1,2,\ldots,r$, und $\beta$ eine der Zahlen $1,2,\ldots,s$, so ist das $(r+\beta)^{\te}$ Element der $\alpha^{\ten}$ Zeile
\[
\sum_{R,S}\chi(A_\alpha R^{-1})x_{RS^{-1}}\psi(SB_\beta^{-1})
=\sum_{R,S}\chi(A_\alpha R^{-1})\psi(RS^{-1})x_{SB_\beta^{-1}}=0.
\]
Folglich zerfällt $Y$ in $k$ Theilmatrizen, deren erste von den $r^2$ Elementen
\[
\frac{h}{f}\sum_R\chi(A_\alpha R^{-1})x_{RA_\beta^{-1}}
=\sum_R\frac{h}{f}\chi(A_\alpha A_\beta^{-1}R^{-1})x_R
\]
gebildet wird. In dieser Summe ist $x_E$ mit der Matrix
\[
N_1=\left(\frac{h}{f}\chi(A_\alpha A_\beta^{-1})\right)\qquad(\alpha,\beta=1,2,\ldots,r)
\]
multiplicirt, deren Determinante von Null verschieden ist. Ebenso sei
\[
N_2=\left(\frac{h}{f'}\psi(B_\alpha B_\beta^{-1})\right)\qquad(\alpha,\beta=1,2,\ldots,s)
\]
u. s. w. Aus diesen Matrizen der Grade $r,s,\ldots$ bilde man die Matrix
\[
\begin{pmatrix}
N_1&0&0&\cdots\\
0&N_2&0&\cdots\\
0&0&N_3&\cdots\\
\cdot&\cdot&\cdot&\cdots
\end{pmatrix}=N
\]
des Grades $h$; dann zerfällt auch $YN^{-1}=Z$ in $k$ Theilmatrizen, deren erste ist
\begin{equation}\tag{1.}
\left(\sum_R\frac{h}{f}\chi(A_\alpha A_\beta^{-1}R^{-1})x_R\right)
\left(\frac{h}{f}\chi(A_\alpha A_\beta^{-1})\right)^{-1}\qquad(\alpha,\beta=1,2,\ldots,f^2).
\end{equation}
Darin ist $x_E$ mit der Hauptmatrix multiplicirt. Vergleicht man also in der Gleichung $LXMN^{-1}=Z$ die mit $x_E$ multiplicirten Matrizen, so erhält man $LMN^{-1}=E$, und mithin ist
\begin{equation}\tag{2.}
LX L^{-1}=Z
\end{equation}
eine mit $X$ aequivalente Matrix. Daher gehört $Z$ und jede Theilmatrix von $Z$, wie (1.), zur Gruppe $\Hgrp$. Nach \emph{Pr.} § 11 ist die Determinante von (1.) gleich $\Psi^f$, wo $\Psi$ der zu $\Phi$ conjugirte Primfactor von $\Theta$ ist.

Ist $\Ggrp$ eine invariante Untergruppe von $\Hgrp$, gehört der Charakter $\chi$ zur Gruppe $\frac{\Hgrp}{\Ggrp}$, und ist die Matrix (1.) gleich $\sum(R)x_R$, so ist $(R)=(S)$, wenn $R\sim S$ (mod. $\Ggrp$) ist. Daher stellen die Matrizen $(A),(B),(C),\cdots$ die Gruppe $\frac{\Hgrp}{\Ggrp}$ dar.

Die obige Umformung lässt sich in derselben Weise ohne Anwendung neuer Hülfsmittel für eine beliebige zur Gruppe gehörige Matrix durchführen. Auch eine solche kann, wenn ihre Determinante (3.) § 2 durch verschiedene Primfactoren der Gruppendeterminante theilbar ist, allein mit Benutzung der Charaktere in eine aequivalente zerfallende Matrix transformirt werden, deren einzelne Theile die Determinanten $\Phi^s,{\Phi'}^{s'},\ldots$ haben.

\section*{§ 4.}

Ich wende mich nun zu der im Anfang des vorigen Paragraphen besprochenen vollständigen Zerlegung der Gruppenmatrix. Ihre charakteristische Determinante $\Theta(x-u\eps)$ hat nur einfache Elementartheiler und hat daher für eine Wurzel $u$ der Gleichung $\Phi(x-u\eps)=0$ den Rang $h-f$. Ich setze nun für die Variabelen $x_R$ solche Constanten $k_R$, dass die Gleichung $\Phi(k-u\eps)=0$ eine einfache Wurzel $\rho$ hat, für die keine der Functionen $\Phi'(k-u\eps),\Phi''(k-u\eps),\ldots$ verschwindet. Dann hat auch die Matrix $(k_{PQ^{-1}}-\rho\eps_{PQ^{-1}})$ den Rang $h-f$. Folglich haben die $h$ linearen Gleichungen
\begin{equation}\tag{1.}
\sum_Q k_{PQ^{-1}}a_Q=\rho a_P
\end{equation}
$f$ unabhängige Lösungen $a'_Q,a''_Q,\ldots,a^{(f)}_Q$. Aus ihnen lässt sich jede andere Lösung zusammensetzen, indem man sie mit gewissen Factoren multiplicirt und dann addirt. Ersetzt man $P$ und $Q$ durch $PR^{-1}$ und $QR^{-1}$, so wird
\[
\sum_Q k_{PQ^{-1}}a_{QR^{-1}}=\rho a_{PR^{-1}},
\]
und folglich ist, wenn $x_A,x_B,x_C,\ldots$ unabhängige Variable sind,
\[
\sum_{Q,R} k_{PQ^{-1}}a_{QR^{-1}}x_R=\rho\sum_R a_{PR^{-1}}x_R.
\]
Mithin ist auch $a_Q=\sum_R a^{(\kappa)}_{QR^{-1}}x_R$ eine Lösung der Gleichungen (1.). Demnach giebt es solche Factoren $x_{\kappa\lambda}$, dass
\begin{equation}\tag{2.}
\sum_R a^{(\kappa)}_{QR^{-1}}x_R=\sum_\lambda x_{\kappa\lambda}a^{(\lambda)}_Q,\qquad(\kappa,\lambda=1,2,\ldots,f).
\end{equation}
Die Factoren $x_{\kappa\lambda}$ sind durch diese Bedingungen vollständig bestimmt und sind folglich lineare Functionen der Variabelen $x_R$. Die so erhaltenen $f$ neuen Lösungen der Gleichungen (1.) sind linear unabhängig oder nicht, je nachdem die Determinante $f^{\ten}$ Grades $|x_{\kappa\lambda}|$ von Null verschieden ist oder nicht. Da $x_{\kappa\lambda}=\eps_{\kappa\lambda}$ ist für $x_R=\eps_R$, so kann jene Determinante nicht identisch verschwinden.

Geht $x_{\kappa\lambda}$ in $y_{\kappa\lambda}$ oder in $z_{\kappa\lambda}$ über, wenn man die Variabelen $x_R$ durch andere Variable $y_R$ oder $z_R$ ersetzt, so ist auch
\[
\sum_R a^{(\kappa)}_{QR^{-1}}y_R=\sum_\lambda y_{\kappa\lambda}a^{(\lambda)}_Q,
\qquad
\sum_R a^{(\kappa)}_{QR^{-1}}z_R=\sum_\lambda z_{\kappa\lambda}a^{(\lambda)}_Q,
\]
und mithin
\[
\sum_{Q,R} a^{(\kappa)}_{PQ^{-1}}x_{QR^{-1}}y_R
=\sum_\lambda x_{\kappa\lambda}\sum_R a^{(\lambda)}_{PR^{-1}}y_R
=\sum_{\lambda,\mu}x_{\kappa\lambda}y_{\lambda\mu}a^{(\mu)}_Q,
\]
also unter der Voraussetzung (1.), § 2
\[
\sum_\mu z_{\kappa\mu}a^{(\mu)}_Q
=\sum_Q a^{(\kappa)}_{PQ^{-1}}z_Q
=\sum_{\lambda,\mu}x_{\kappa\lambda}y_{\lambda\mu}a^{(\mu)}_Q
\]
und mithin
\[
z_{\kappa\mu}=\sum_\lambda x_{\kappa\lambda}y_{\lambda\mu}.
\]
Daher ist $(x_{\kappa\lambda})$ eine zu $\Hgrp$ gehörige Matrix $f^{\ten}$ Grades, und mithin ist ihre Determinante
\[
|x_{\kappa\lambda}|=\Phi^s{\Phi'}^{s'}\cdots.
\]
Ist $\chi$ der Charakter, welcher der Primfunction $\Phi$ entspricht, und setzt man $\frac{h}{f}\chi(R^{-1})=c_R$, so ist nach \emph{Pr.} § 8
\[
\Phi(c)=1,
\quad \Phi'(c)=0,\ldots,
\quad \Phi(c-u\eps)=(1-u)^f,
\quad \Phi'(c-u\eps)=(-u)^{f'},\ldots,
\]
wo $\Phi'$ irgend eine von $\Phi$ verschiedene Primfunction des Grades $f'$ ist. Sind $\rho,\rho',\rho'',\ldots$ die $f$ Wurzeln der Gleichung $\Phi(k-u\eps)=0$, so sind nach \emph{Pr.} § 6 $\rho-\lambda,\rho'-\lambda,\rho''-\lambda,\ldots$ die der Gleichung $\Phi(k-\lambda c-u\eps)=0$. Ist also $\lambda$ eine von $0,\rho,\rho',\rho'',\ldots$ verschiedene Constante, so ist $\Phi(k-\lambda c-\rho\eps)$ von Null verschieden. Ferner ist $\Phi'(k-\lambda c-\rho\eps)=\Phi'(k-\rho\eps)$ von Null verschieden, und daher kann die Determinante $\Theta(k-\lambda c-\rho\eps)$ nicht verschwinden.

Für $x_R=c_R$ verschwindet die Determinante $|x_{\kappa\lambda}|$ nicht. Denn sonst wären die $f$ Lösungen
\[
\sum_Q a^{(\kappa)}_{PQ^{-1}}c_Q\qquad(\kappa=1,2,\ldots,f)
\]
nicht linear unabhängig. Man könnte also aus ihnen eine Lösung $\sum_Q a_{PQ^{-1}}c_Q=0$ zusammensetzen, worin die Grössen $a_P$ den Gleichungen (1.) genügen und nicht alle Null sind. Da $c_{PQ}=c_{QP}$ ist, so sind die Matrizen $a_{PQ^{-1}}$ und $c_{PQ^{-1}}$ vertauschbar, und mithin ist auch $\sum_Q c_{PQ^{-1}}a_Q=0$. Daher wäre auch
\[
\sum_Q(k_{PQ^{-1}}-\lambda c_{PQ^{-1}}-\rho\eps_{PQ^{-1}})a_Q=0.
\]
Da aber die Determinante dieser $h$ Gleichungen von Null verschieden ist, so können ihnen nur die Werthe $a_Q=0$ genügen.

Das Product $\Phi(c)^s{\Phi'}(c)^{s'}{\Phi''}(c)^{s''}\cdots$ kann aber nur dann von Null verschieden sein, wenn $s'=s''=\cdots=0$ ist. Da $|x_{\kappa\lambda}|$ von Grade $f$ ist, so ist folglich
\begin{equation}\tag{3.}
|x_{\kappa\lambda}|=\Phi(x),\qquad |x_{\kappa\lambda}-u\eps_{\kappa\lambda}|=\Phi(x-u\eps).
\end{equation}
Nun habe ich \emph{Pr.} § 11 gezeigt, dass sich $\Phi(x)$ durch $f^2$, aber nicht durch weniger lineare Verbindungen der Variabelen $x_R$ darstellen lässt. Mithin sind die $f^2$ linearen Functionen $x_{\kappa\lambda}$ der $h$ Variabelen $x_R$ von einander unabhängig.

Diesen merkwürdigen Satz, dass es eine zur Gruppe $\Hgrp$ gehörige Matrix giebt, deren $f^2$ Elemente unabhängige Variable sind, hat auch \textsc{Molien} gefunden in seiner ausgezeichneten Arbeit \emph{Über Systeme höherer complexer Zahlen} (Math. Ann. Bd. 41, S. 124), auf die mich \textsc{Study} vor kurzem aufmerksam gemacht hat. In einer weiteren Arbeit \emph{Eine Bemerkung zur Theorie der homogenen Substitutionsgruppen}, Sitzungsberichte der Naturforscher-Gesellschaft zu Dorpat 1897, Jahrg. 18, S. 259 hat \textsc{Molien} die dort gefundenen allgemeinen Resultate speciell auf die Gruppendeterminante angewendet.

Ist $X=(x_{\kappa\lambda})=\sum(R)x_R$, so ergiebt sich durch Vergleichung der Coefficienten von $u^{f-1}$ in (3.)
\begin{equation}\tag{4.}
\sum_R\chi(R)x_R=\sum_\lambda x_{\lambda\lambda}.
\end{equation}
Ist also $(R)=(r_{\kappa\lambda})$, so ist
\begin{equation}\tag{5.}
\chi(R)=\sum_\lambda r_{\lambda\lambda}.
\end{equation}
Ist nun die von den Matrizen $(A),(B),(C)\cdots$ gebildete Gruppe der Gruppe $\frac{\Hgrp}{\Ggrp}$ isomorph, so ist $(R)=(S)$, wenn $R\sim S$ (mod. $\Ggrp$) ist. Daher ist auch $\chi(R)=\chi(S)$. Der Charakter $\chi$ gehört also ebenfalls zur Gruppe $\frac{\Hgrp}{\Ggrp}$.

Sei $L=(l_{\kappa\lambda})$ irgend eine constante Matrix $f^{\ten}$ Grades von nicht verschwindender Determinante. Wählt man dann für die $f$ Lösungen $a'_Q,a''_Q,\ldots$ der Gleichungen (1.) irgend $f$ andere unabhängige Lösungen, so erhält man statt $X$ immer eine mit $X$ ähnliche Matrix $L^{-1}XL$ und bei passender Wahl der $f$ Lösungen jede solche Matrix.

Man kann aber auch für $\rho$ eine andere Wurzel der Gleichung $\Phi(k-u\eps)=0$ wählen, und ferner ist für $k_R$ jedes System von $h$ Zahlen zulässig, das gewisse Ungleichheiten befriedigt. Dann treten an die Stelle der Grössen $a'_Q,a''_Q,\ldots$ andere, und statt der Matrix $X$ erhält man eine Matrix $U=(u_{\kappa\lambda})$, deren Elemente $u_{\kappa\lambda}$ lineare Functionen der Variabelen $x_R$ sind. Bei jeder Wahl der willkürlichen Grössen ist aber $U$ eine zu $\Hgrp$ gehörige Matrix $f^{\ten}$ Grades, deren charakteristische Determinante $\Phi(x-u\eps)$ ist. Da $\Phi(x)$ die $h$ Grössen $x_R$ nur in den $f^2$ unabhängigen linearen Verbindungen $x_{\kappa\lambda}$ enthält, so müssen die $f^2$ Grössen $u_{\kappa\lambda}$ lineare Verbindungen der Variabelen $x_{\kappa\lambda}$ sein. Nach einem Satze, den ich in § 7 entwickeln werde, kann man daher eine constante Matrix $L$ so bestimmen, dass entweder $U=L^{-1}XL$ oder $U=L^{-1}X'L$ ist, wo $X'$ die zu $X$ conjugirte Matrix ist, und zwar ist, wenn $f>1$ ist, nur der eine dieser beiden Fälle möglich.

Für die hier betrachteten Matrizen kann nun aber, wenn $f>1$ ist, nicht $U=L^{-1}X'L$ sein. Denn ersetzt man die Variabelen $x_R$ durch $y_R$ oder $z_R$, so möge $X$ in $Y$ oder $Z$, und $U$ in $V$ oder $W$ übergehen. Dann ist $XY=Z$ und $UV=W$. Ist nun $U=L^{-1}X'L$, $V=L^{-1}Y'L$, $W=L^{-1}Z'L$, so ist auch $X'Y'=Z'$, und folglich, weil $X'Y'=(YX)'$ ist, $Z=YX=XZ$. Ist also $X=\sum(R)x_R$, so sind je zwei der Matrizen $(A),(B),(C)\cdots$ mit einander vertauschbar. Die von ihnen gebildete Gruppe und die ihr holoedrisch isomorphe Gruppe $\frac{\Hgrp}{\Ggrp}$ ist folglich eine commutative, und der zu ihr gehörige Charakter $\chi$ hat den Grad $f=1$ (\emph{Pr.} § 2). Demnach ist $U=L^{-1}XL$, und man kann bei jeder Verfügung über die willkürlichen Grössen die $f$ Lösungen $a'_Q,a''_Q,\ldots$ so wählen, dass $U=X$ wird.

Gehört der Charakter $\chi$ zur Gruppe $\frac{\Hgrp}{\Ggrp}$, so giebt es eine Primfunction $f^{\ten}$ Grades $\Psi(y)$ der Determinante dieser Gruppe, die durch die Substitution (5.) § 1 in $\Phi(x)$ übergeht. Diese kann als die Determinante einer Matrix $f^{\ten}$ Grades $(y_{\kappa\lambda})$ dargestellt werden, deren Elemente lineare Functionen der Variabelen $y_{R\Ggrp}$ sind, und die zur Gruppe $\frac{\Hgrp}{\Ggrp}$ gehört. Macht man darin die Substitution (5.) § 1, so geht sie in eine Matrix $X=(x_{\kappa\lambda})$ über, die zur Gruppe $\Hgrp$ gehört, und deren Determinante gleich $\Phi(x)$ ist. Da jede andere Matrix, welche dieselben Eigenschaften besitzt, gleich $L^{-1}XL$ ist, so ist damit die Umkehrung des oben erhaltenen Satzes bewiesen, nämlich dass, wenn $\chi$ zur Gruppe $\frac{\Hgrp}{\Ggrp}$ gehört, auch immer $(R)=(S)$ ist, falls $R\sim S$ (mod. $\Ggrp$) ist.

\section*{§ 5.}

Jetzt mögen die Grössen $k_R$ so gewählt werden, dass die Gleichung
\[
\Phi(k-u\eps)\Phi'(k-u\eps)\Phi''(k-u\eps)\cdots=0
\]
keine mehrfachen Wurzeln hat. Dann hat die Determinante der Schaar bilinearer Formen
\begin{equation}\tag{1.}
\sum_{P,Q}(k_{PQ^{-1}}-u\eps_{PQ^{-1}})u_Pv_Q
\end{equation}
lauter lineare Elementartheiler. Die Form $\sum k_{PQ^{-1}}u_Pv_Q$ kann daher nach einem Satze von \textsc{Weierstrass} durch eine lineare Substitution von nicht verschwindender Determinante
\begin{equation}\tag{2.}
u'_\nu=\sum_R a^{(\nu)}_R u_R,\qquad v_R=\sum_\nu a^{(\nu)}_R v'_\nu\qquad(\nu=1,2,\ldots,h),
\end{equation}
die $\sum u_Rv_R$ in $\sum u'_\nu v'_\nu$ überführt, in
\begin{equation}\tag{3.}
\rho u'_1v'_1+\rho u'_2v'_2+\cdots+\rho u'_fv'_f+
ho'u'_{f+1}v'_{f+1}+\cdots+
ho'u'_{f+f'}v'_{f+f'}+\cdots
\end{equation}
transformirt werden, falls $\rho$ eine $f$-fache, $\rho'$ eine $f'$-fache, $\ldots$ Wurzel der charakteristischen Gleichung $\Theta(k-u\eps)=0$ ist. Dann ist
\[
\sum_{P,Q}k_{PQ^{-1}}u_P\sum_\nu a^{(\nu)}_Qv'_\nu
=\rho\left(\sum_P a^{(1)}_P u_P\right)v'_1+
ho\left(\sum_P a^{(2)}_P u_P\right)v'_2+\cdots,
\]
und mithin
\[
\sum_Q k_{PQ^{-1}}a^{(\kappa)}_Q=\rho a^{(\kappa)}_P\qquad(\kappa=1,2,\ldots,f).
\]
Da die Determinante $h^{\ten}$ Grades $|a^{(\nu)}_R|$ von Null verschieden ist, so verschwinden in dem System der $fh$ Grössen $a^{(\kappa)}_Q$ nicht alle Determinanten $f^{\ten}$ Grades, und mithin bilden sie $f$ unabhängige Lösungen der Gleichungen (1.) § 4. Ersetzt man sie durch irgend $f$ andere unabhängige Lösungen, so behält die Substitution (2.) die doppelte Eigenschaft, dass ihre Determinante von Null verschieden ist, und dass sie die Formenschaar (1.) in die Normalform (3.) transformirt. Daher kann man die $h^2$ Grössen $a^{(\nu)}_R$ durch die linearen Gleichungen (1.) § 4 definiren und durch die analogen Gleichungen, die den Wurzeln $\rho',\rho'',\ldots$ entsprechen, und der wesentliche Inhalt des oben benutzten Satzes von \textsc{Weierstrass} besteht darin, dass dann die Determinante $h^{\ten}$ Grades $|a^{(\nu)}_R|$ von Null verschieden ist. Nun sei, wie oben
\[
\sum_Q a^{(\kappa)}_{PQ^{-1}}x_Q=\sum_\lambda x_{\kappa\lambda}a^{(\lambda)}_P
\]
oder
\[
\sum_P a^{(\kappa)}_{QP^{-1}}x_{PQ^{-1}}=\sum_\lambda x_{\kappa\lambda}a^{(\lambda)}_{Q^{-1}}\qquad(\kappa,\lambda=1,2,\ldots,f).
\]
In derselben Weise sei, entsprechend der $f'$fachen Wurzel $\rho'$,
\[
\sum_P a^{(\kappa)}_{QP^{-1}}x_{PQ^{-1}}=\sum_\lambda x_{\kappa\lambda}a^{(\lambda)}_{Q^{-1}}\qquad(\kappa,\lambda=f+1,f+2,\ldots,f+f').
\]
In dem speciellen Falle, wo $\rho'$ eine Wurzel derselben Gleichung $\Phi(x-u\eps)=0$ wie $\rho$ (also $f'=f$) ist, kann man durch geeignete Wahl der Lösung der zu (1.) § 4 analogen Gleichungen erreichen, dass $x_{f+\kappa,f+\lambda}=x_{\kappa\lambda}$ ist.

Dann geht die bilineare Form
\begin{equation}\tag{4.}
\sum_{P,Q}x_{PQ^{-1}}u_Pv_Q
\end{equation}
durch die Substitution
\begin{equation}\tag{5.}
u_R=\sum_\nu a^{(\nu)}_{R^{-1}}u'_\nu,
\qquad
v'_\nu=\sum_R a^{(\nu)}_{R^{-1}}v_R
\end{equation}
in
\begin{equation}\tag{6.}
\sum_{\kappa,\lambda}^{1\ldots f}x_{\kappa\lambda}u'_\kappa v'_\lambda
+\sum_{\kappa,\lambda}^{f+1\ldots f+f'}x_{\kappa\lambda}u'_\kappa v'_\lambda+\cdots
\end{equation}
über, die in $f$ Formen von $f^2$ Variabelen, $f'$ Formen von ${f'}^2$ Variabelen, u. s. w. zerfällt. Durch dieselbe Substitution geht $\sum u_Rv_R$ in $\sum u'_\nu v'_\nu$ über. Ist also $X$ die Gruppenmatrix $(x_{PQ^{-1}})$, so giebt es eine constante Matrix $h^{\ten}$ Grades $L$ von nicht verschwindender Determinante der Art, dass die aequivalente Matrix $L^{-1}XL$ zerfällt in $f$ einander gleiche Theilmatrizen des Grades $f$, deren charakteristische Determinanten gleich $\Phi(x-u\eps)$ sind, in $f'$ einander gleiche Theilmatrizen des Grades $f'$, deren charakteristische Determinanten gleich $\Phi'(x-u\eps)$ sind, u. s. w., und die $f^2+{f'}^2+\cdots=h$ Elemente dieser Theilmatrizen sind $h$ von einander unabhängige Variable, weil sich $\Theta(x)$ nicht durch weniger als $h$ lineare Verbindungen der Variabelen $x_R$ ausdrücken lässt.

Eine Darstellung einer Gruppe durch lineare Substitutionen, für welche die entsprechende Determinante (3.) § 2 unzerlegbar ist, nenne ich eine \emph{primitive Darstellung}. Dann ist die Anzahl der Classen primitiver Darstellungen für die Gruppe $\Hgrp$ und für die mit ihr isomorphen Gruppen gleich der Anzahl $k$ der Classen conjugirter Elemente, worin die Elemente von $\Hgrp$ zerfallen. Ist $f$ die Anzahl der Variabelen, die eine der Substitutionen transformirt, so werden die $k$ Zahlen $g=f^2$ durch Auflösung einer Gleichung $k^{\ten}$ Grades gefunden, die ich in meiner Arbeit \emph{Über Gruppencharaktere} (Sitzungsberichte 1896, § 4, (12.)) entwickelt habe.

Zur Erläuterung dieser Transformation der Gruppenmatrix wähle ich das Beispiel, das \textsc{Dedekind} im Jahre 1886 gefunden und mir im April 1896 mitgetheilt hat. Seien
\[
\begin{array}{cccccccccccc}
1.&abc&2.&bca&3.&cab&4.&acb&5.&cba&6.&bac
\end{array}
\]
die 6 Permutationen von 3 Symbolen. Die Substitutionen, die $abc$ in diese 6 Permutationen überführen, mögen statt mit $A,B,C$, mit den Ziffern $1,2,\ldots,6$ bezeichnet werden. Sei $\rho$ eine dritte Wurzel der Einheit und
\[
\begin{array}{lll}
u=x_1+x_2+x_3, &\qquad& v=x_4+x_5+x_6,\\
u_1=x_1+\rho x_2+\rho^2x_3, && v_1=x_4+\rho x_5+\rho^2x_6,\\
u_2=x_1+\rho^2x_2+\rho x_3, && v_2=x_4+\rho^2x_5+\rho x_6.
\end{array}
\]
Ferner seien $X$, $L$ und $U$ die drei Matrizen
\[
\begin{pmatrix}
x_1&x_3&x_2&x_4&x_5&x_6\\
x_2&x_1&x_3&x_5&x_6&x_4\\
x_3&x_2&x_1&x_6&x_4&x_5\\
x_4&x_5&x_6&x_1&x_3&x_2\\
x_5&x_6&x_4&x_2&x_1&x_3\\
x_6&x_4&x_5&x_3&x_2&x_1
\end{pmatrix},
\qquad
\begin{pmatrix}
1&-1&1&0&0&1\\
1&-1&\rho^2&0&0&\rho\\
1&-1&\rho&0&0&\rho^2\\
1&1&0&1&1&0\\
1&1&0&\rho&\rho^2&0\\
1&1&0&\rho^2&\rho&0
\end{pmatrix},
\]
\[
\begin{pmatrix}
u+v&0&0&0&0&0\\
0&u-v&0&0&0&0\\
0&0&u_1&v_1&0&0\\
0&0&v_2&u_2&0&0\\
0&0&0&0&u_1&v_1\\
0&0&0&0&v_2&u_2
\end{pmatrix}.
\]
Dann ist $X$ die Gruppenmatrix, und es ist
\[
XL=LU,
\qquad
L^{-1}XL=U,
\]
und indem man $L'L$ bildet, erkennt man, dass die Determinante von $L$ nicht verschwindet.

\section*{§ 6.}

Die Formel (5.) § 4 führt zu einer tieferen Einsicht in die Bedeutung der Gleichungen, die ich zur Berechnung der Charaktere einer Gruppe $\Hgrp$ entwickelt habe. Ist die Determinante der zur Gruppe gehörigen Matrix
\[
X=(x_{\kappa\lambda})=\sum_R(R)x_R\qquad(\kappa,\lambda=1,2,\ldots,f)
\]
ein Primfactor $f^{\ten}$ Grades $\Phi$ der Gruppendeterminante, so sind die $f^2$ Grössen $x_{\kappa\lambda}$ unabhängige Variable. Daher kann man den $h$ Variabelen $x_R$ solche Werthe geben, dass $X$ einer beliebigen Matrix $f^{\ten}$ Grades gleich wird.

Die beiden ähnlichen Matrizen $(A)$ und $(R)^{-1}(A)(R)=(B)$ haben dieselbe charakteristische Function $|(A)-u(E)|=|(B)-u(E)|$. Vergleicht man auf beiden Seiten dieser Gleichung die Coefficienten von $u^{f-1}$, so erhält man nach (5.) § 4 $\chi(A)=\chi(B)$. Demnach ist $\chi(R^{-1}AR)=\chi(A)$ oder, was dasselbe ist, $\chi(QP)=\chi(PQ)$.

Theilt man die Elemente von $\Hgrp$ in die $k$ Classen conjugirter Elemente, so möge die $\alpha^{\te}$ Classe aus $h_\alpha$ Elementen bestehen $(\alpha=0,1,\ldots,k-1)$. Durchläuft $R$ alle $h$ Elemente von $\Hgrp$, und ist $A$ ein bestimmtes Element der $\alpha^{\ten}$ Classe, so stellt $R^{-1}AR$ jedes Element der Classe und jedes $\frac{h}{h_\alpha}$ Mal dar. Daher ist
\[
\sum_R(R^{-1}AR)=\frac{h}{h_\alpha}\sum_{(\alpha)}'(R),
\]
wo die Summe rechts über die $h_\alpha$ Elemente der $\alpha^{\ten}$ Classe zu erstrecken ist. Nun ist aber
\[
(S)\sum_R(R^{-1}AR)=\sum_R(SR^{-1}AR)=\sum_R(R^{-1}ARS)=\left(\sum_R(R^{-1}AR)\right)(S).
\]
Die dritte Summe erhält man aus der zweiten, indem man $R$ durch $RS$ ersetzt. Mithin ist die Matrix $\sum_{(\alpha)}'(R)$ mit $(S)$ vertauschbar, also auch mit $\sum(S)x_S$, demnach mit jeder Matrix $f^{\ten}$ Grades. Diese Eigenschaft hat aber nur die Hauptmatrix, mit einem scalaren Factor multiplicirt. Folglich ist
\begin{equation}\tag{1.}
f\sum_{(\alpha)}'(R)=(E)h_\alpha\chi_\alpha,
\qquad
f\sum_R(R^{-1}AR)=(E)h\chi_\alpha.
\end{equation}
Um den noch unbekannten Factor $\chi_\alpha$ zu bestimmen, bilde man auf beiden Seiten die Summe der Diagonalelemente. Diese ist für jede der $h_\alpha$ mit $(A)$ ähnlichen Matrizen $(R)$ gleich $\chi(R)=\chi(A)$, und für $(E)$ gleich $f$. Mithin ist
\[
\chi_\alpha=\chi(A).
\]
Componirt man die Matrix $(A)$ mit der Matrix
\[
f\sum_{(\beta)}'(S)=(E)h_\beta\chi_\beta,
\]
so erhält man
\[
f\sum_{(\beta)}'(AS)=(A)h_\beta\chi_\beta,
\]
und, wenn man wieder auf beiden Seiten die Summe der Diagonalelemente bildet,
\begin{equation}\tag{2.}
f\sum_{(\beta)}'\chi(AS)=h_\beta\chi(A)\chi(B),
\qquad
f\sum_R\chi(AR^{-1}BR)=h\chi(A)\chi(B).
\end{equation}
Dies sind die Gleichungen, welche die Verhältnisse der $k$ Werthe jedes der $k$ Charaktere bestimmen.

Man kann aber diesen Formeln noch eine andere Deutung geben: Seien $e_A,e_B,e_C,\cdots$ die $h$ unabhängigen Einheiten eines Systems hypercomplexer Zahlen, wofür die gewöhnlichen Regeln der Addition und das distributive und associative Gesetz der Multiplication, aber nicht das commutative vorausgesetzt werden. Diesen Bedingungen genügt das Multiplicationsgesetz
\begin{equation}\tag{3.}
e_Ae_B=e_{AB}.
\end{equation}
Demnach vertritt $e_E$ die Zahl 1. Setzt man nun
\begin{equation}\tag{4.}
\sum_{(\alpha)}'e_R=\frac{h_\alpha}{f}e_\alpha,
\end{equation}
so sind die $k$ complexen Zahlen $e_0,e_1,\ldots,e_{k-1}$ mit jeder Zahl des Systems, also auch unter einander vertauschbar. Das Product von zweien dieser commutativen Zahlen ist
\[
\frac{h_\alpha h_\beta}{f^2}e_\alpha e_\beta
=\left(\sum_{(\alpha)}'e_R\right)\left(\sum_{(\beta)}'e_S\right)
=\sum_\tau' h_{\alpha,\beta,\tau} e_T,
\]
wo $h_{\alpha,\beta,\tau}$ angiebt, wie viele der $h_\alpha h_\beta$ Producte $RS$ dem bestimmten Elemente $T$ gleich sind. Nach \emph{Pr.} § 7 ist diese Anzahl, falls $T$ der $\gamma^{\ten}$ Classe angehört, gleich $\frac{1}{h_\gamma}h_{\alpha\beta\gamma}$, hat also für conjugirte Elemente $T$ denselben Werth. Mithin ist
\begin{equation}\tag{5.}
\frac{h_\alpha h_\beta}{f^2}e_\alpha e_\beta
=\sum_\gamma^{k-1}\frac{1}{h_\gamma}h_{\alpha\beta\gamma}e_\gamma.
\end{equation}
Demnach bilden die $k$ unabhängigen Zahlen $e_0,e_1,\ldots,e_{k-1}$ für sich die Basis eines Systems complexer Zahlen, für die aber auch das commutative Gesetz der Multiplication gilt. Wie ich \emph{Über Gruppencharaktere} S. 991 gezeigt habe, gelten dafür überdies die Einschränkungen, unter denen \textsc{Weierstrass} und \textsc{Dedekind} solche Zahlensysteme untersucht haben. Demnach giebt es wirkliche Zahlen $\chi_\alpha$, die, für $e_\alpha$ gesetzt, den Gleichungen (5.) genügen:
\begin{equation}\tag{6.}
h_\alpha h_\beta\chi_\alpha\chi_\beta=f\sum_\gamma h_{\alpha\beta\gamma}\chi_\gamma.
\end{equation}
Man kann daher nicht auf Widersprüche kommen, wenn man die bisherige Unabhängigkeit der $h$ complexen Einheiten $e_R$ durch die linearen Gleichungen $e_\alpha=\chi_\alpha e_0$ oder
\begin{equation}\tag{7.}
f\sum_{(\alpha)}'e_R=h_\alpha\chi_\alpha e_E
\end{equation}
einschränkt. Durch Multiplication mit $e_B$ fliessen daraus nach (3.) die weiteren linearen Gleichungen
\begin{equation}\tag{8.}
f\sum_{(\alpha)}'e_{BR}=h_\alpha\chi_\alpha e_B.
\end{equation}
In Folge dieser Relationen reducirt sich die Anzahl der linear unabhängigen unter den Einheiten $e_R$ auf $f^2$, und diese kann man so wählen, dass die Formeln für ihre Multiplication mit denen für die Composition aller Matrizen $f^{\ten}$ Grades übereinstimmen.

Die hier angedeutete Rechnung hat \textsc{Dedekind} für das in § 5 mitgetheilte Beispiel einer Gruppe der Ordnung $h=6$ durchgeführt und für einige andere Gruppen der kleinsten Ordnungszahlen, so namentlich für die Gruppe $\Omega$ der Ordnung 8, die er die Quaternionengruppe genannt hat, weil das aus ihr abgeleitete System hypercomplexer Zahlen mit dem System der \textsc{Hamilton}'schen Quaternionen übereinstimmt. In seiner Arbeit \emph{Über Gruppen, deren sämmtliche Theiler Normaltheiler sind}, Math. Ann. Bd. 48, spielt er S. 551 auf diese Beziehungen an mit den Worten: »Es findet aber, wie ich schon im Februar 1886 erkannt habe, eine noch tiefer liegende Beziehung zwischen der Gruppe $\Omega$ und \textsc{Hamilton}'s Quaternionen statt«. Die Rechnung mit Quaternionen ist ja auch, wenn man gewöhnliche complexe Zahlen als scalare Coefficienten zulässt, der Rechnung mit Matrizen zweiten Grades völlig aequivalent, weil die beiden quaternären quadratischen Formen $x^2+y^2+z^2+t^2$ und $xt-yz$ in einander transformirt werden können. Nur wenn man sich auf reelle Coefficienten beschränkt, können die Quaternionen eine selbständige Bedeutung beanspruchen, deren Wesen darin besteht, dass es ausser ihnen und den gewöhnlichen complexen Zahlen kein Zahlensystem giebt, worin ein Product nicht verschwinden kann, ohne dass einer der Factoren Null ist, wie ich in meiner Arbeit \emph{Über lineare Substitutionen und bilineare Formen}, Crelle's Journ. Bd. 84, § 14 zuerst dargelegt habe.

\section*{§ 7.}

Die Elemente $x_{\alpha\beta}$ $(\alpha,\beta=1,2,\ldots,n)$ einer Matrix $n^{\ten}$ Grades $X$ seien $n^2$ von einander unabhängige Variable. Vertauscht man in $X$ die Zeilen mit den Spalten, so erhält man die zu $X$ conjugirte Matrix $X'$. Die Elemente der beiden Matrizen $A$ und $B$ seien constante Grössen. Ihre Determinanten $|A|$ und $|B|$ seien von Null verschiedene Grössen, deren Product gleich $k$ ist. Dann haben die beiden Matrizen $AXB$ und $AX'B$ die Determinante $k|X|$, und in jeder von ihnen sind die Elemente lineare Functionen der $n^2$ Variabelen $x_{\alpha\beta}$. Umgekehrt gilt der Satz:

\begin{quote}\itshape
I. Sind die Elemente der Matrix $X$ unabhängige Variable und die der Matrix $Y$ lineare Functionen dieser Variabelen, und unterscheidet sich die Determinante der Matrix $Y$ von der der Matrix $X$ nur um einen constanten von Null verschiedenen Factor, so ist entweder $Y=AXB$ oder $Y=AX'B$, wo $A$ und $B$ constante Matrizen sind; und zwar tritt, wenn der Grad von $X$ grösser als 1 ist, nur einer dieser beiden Fälle ein, und die Matrizen $A$ und $B$ sind bis auf einen scalaren Factor vollständig bestimmt.
\end{quote}

Der zweite Theil dieses Satzes ist leicht zu beweisen. Denn sei $Y=AXB=CXD$. Setzt man darin $x_{\alpha\beta}=0$ oder 1, je nachdem $\alpha$ und $\beta$ verschieden oder gleich sind, so folgt daraus $AB=CD$, und wenn man $BD^{-1}=A^{-1}C=F$ setzt, $XF=FX$. Demnach ist $F$ mit jeder Matrix vertauschbar, und mithin ist $F=hE$, wo $h$ ein scalarer Factor und $E$ die Hauptmatrix ist. Folglich ist $C=hA$ und $D=\frac{1}{h}B$. Ist ferner $n>1$, so kann auch nicht $AXB=CX'D$ sein. Denn daraus folgt in derselben Weise $XF=FX'$, also z. B. für $n=2$
\[
\begin{pmatrix}x&y\\z&t\end{pmatrix}
\begin{pmatrix}a&b\\c&d\end{pmatrix}
=\begin{pmatrix}a&b\\c&d\end{pmatrix}
\begin{pmatrix}x&z\\y&t\end{pmatrix}.
\]
Diesen vier linearen Gleichungen kann man, wenn die Unbekannten $a,b,c,d$ von den Variabelen $x,y,z,t$ unabhängig sein sollen, nur durch verschwindende Werthe der Constanten genügen.

Zwischen den Unterdeterminanten $m^{\ten}$ Grades $(0<m<n)$ der Matrix $X$ besteht keine lineare Relation mit constanten Coefficienten. Denn seien $u,v,w,\ldots$ diese Unterdeterminanten, $a,b,c,\ldots$ Constante, und sei $au+bv+cw+\cdots=0$. Setzt man darin alle Variabelen $x_{\alpha\beta}=0$ ausser den $m^2$ in $u$ vorkommenden, so verschwinden $v,w,\ldots$, und mithin muss $a=0$ sein. Die Unterdeterminanten $(n-1)^{\ten}$ Grades sind die ersten Ableitungen von $|X|$ nach den Variabelen $x_{\alpha\beta}$. Da zwischen ihnen keine lineare Relation mit constanten Coefficienten besteht, so kann man $|X|$ nicht als Function von weniger als $n^2$ Variabelen darstellen, die lineare Functionen der $n^2$ Variabelen $x_{\alpha\beta}$ sind. Ist also $|Y|=k|X|$, so sind die $n^2$ Variabelen $y_{\alpha\beta}$ von einander unabhängig. Daher besteht auch zwischen den Unterdeterminanten $m^{\ten}$ Grades der Matrix $Y$ keine lineare Relation mit constanten Coefficienten.

Der Coefficient von $x_{\kappa\kappa}$ in $y_{\alpha\beta}$ sei $c_{\alpha\beta}^{(\kappa)}$ oder, da zunächst $\kappa$ ein fest bleibender Index ist, kurz $c_{\alpha\beta}$. Ist $r$ eine neue Variable, und ersetzt man $x_{\kappa\kappa}$ durch $x_{\kappa\kappa}+r$, so geht $y_{\alpha\beta}$ in $y_{\alpha\beta}+rc_{\alpha\beta}$ über. Daher ist $|y_{\alpha\beta}+rc_{\alpha\beta}|$ gleich der Determinante, die aus $k|x_{\alpha\beta}|$ hervorgeht, indem man $x_{\kappa\kappa}$ durch $x_{\kappa\kappa}+r$ ersetzt. Diese aber ist eine Function ersten Grades von $r$, und folglich verschwindet auch in jener Determinante der Coefficient von $r^2$. Derselbe ist gleich der Summe der Producte jeder Determinante zweiten Grades der Matrix $c_{\alpha\beta}$ und der complementären Determinante $(n-2)^{\ten}$ Grades der Matrix $y_{\alpha\beta}$. Zwischen den letzteren aber besteht, wenn $n>2$ ist, keine lineare Relation mit constanten Coefficienten. Folglich müssen die Determinanten zweiten Grades der Matrix $c_{\alpha\beta}$ sämmtlich verschwinden. Für $n=2$ ist dies ohne Weiteres ersichtlich. Daher kann man $2n$ Grössen $p_{\alpha},q_{\alpha}$ so bestimmen, dass $c_{\alpha\beta}=p_\alpha q_\beta$ ist.

Sei $c_{\alpha\beta}^{(\kappa)}=p_{\alpha\kappa}q_{\kappa\beta}$. Setzt man die Variabelen $x_{\alpha\beta}$, deren Indices $\alpha$ und $\beta$ verschieden sind, gleich Null, so mögen $X$ und $Y$ in $X_0$ und $Y_0$ übergehen. Die Elemente der Matrix $Y_0$ sind dann
\[
\sum_\kappa c_{\alpha\beta}^{(\kappa)}x_{\kappa\kappa}=\sum_\kappa p_{\alpha\kappa}x_{\kappa\kappa}q_{\kappa\beta}.
\]
Ist also $P$ die von den Grössen $p_{\alpha\beta}$, und $Q$ die von den Grössen $q_{\alpha\beta}$ gebildete Matrix, so ist
\[
Y_0=PX_0Q,
\]
also weil $|Y_0|=k|X_0|=kx_{11}x_{22}\cdots x_{nn}$ ist, $|P|\,|Q|=k$. Demnach sind $|P|$ und $|Q|$ von Null verschieden.

Die Elemente $z_{\alpha\beta}$ der Matrix $Z=P^{-1}YQ^{-1}$ werden also gleich denen von $X_0$, falls man die Variabelen $x_{\alpha\beta}$, deren Indices verschieden sind, gleich Null setzt, also gleich $x_{\alpha\alpha}$ oder 0, je nachdem $\beta=\alpha$ ist oder nicht. Oder die Grössen $z_{\alpha\beta}$ $(\alpha\ne\beta)$ und $z_{\alpha\alpha}-x_{\alpha\alpha}=v_\alpha$ hängen allein von den Variabelen $x_{\alpha\beta}$ mit verschiedenen Indices ab. Entwickelt man die identische Gleichung $|X|=|Z|$ nach Potenzen von $x_{11},x_{22},\ldots,x_{nn}$, so ergiebt sich durch Vergleichung der Coefficienten des Productes $x_{22}x_{33}\cdots x_{nn}$, dass $v_1=0$ ist. Ebenso ist $v_\alpha=0$. Vergleicht man dann die mit $x_{33}x_{44}\cdots x_{nn}$ multiplicirten Glieder, so findet man $x_{12}x_{21}=z_{12}z_{21}$ und allgemein
\begin{equation}\tag{1.}
x_{\alpha\beta}x_{\beta\alpha}=z_{\alpha\beta}z_{\beta\alpha}.
\end{equation}
Vergleicht man endlich die Coefficienten von $x_{44}x_{55}\cdots x_{nn}$, so erhält man
\[
\left|\begin{array}{ccc}
0&x_{12}&x_{13}\\
x_{21}&0&x_{23}\\
x_{31}&x_{32}&0
\end{array}\right|=
\left|\begin{array}{ccc}
0&z_{12}&z_{13}\\
z_{21}&0&z_{23}\\
z_{31}&z_{32}&0
\end{array}\right|
\]
und ebenso allgemein
\[
x_{\beta\gamma}x_{\gamma\alpha}x_{\alpha\beta}+x_{\gamma\beta}x_{\alpha\gamma}x_{\beta\alpha}
=z_{\beta\gamma}z_{\gamma\alpha}z_{\alpha\beta}+z_{\gamma\beta}z_{\alpha\gamma}z_{\beta\alpha}.
\]
Nach (1.) ist auch das Product der beiden Summanden links gleich dem der Summanden rechts. Daher ist entweder
\[
x_{\beta\gamma}x_{\gamma\alpha}x_{\alpha\beta}=z_{\beta\gamma}z_{\gamma\alpha}z_{\alpha\beta},
\qquad
x_{\gamma\beta}x_{\alpha\gamma}x_{\beta\alpha}=z_{\gamma\beta}z_{\alpha\gamma}z_{\beta\alpha},
\]
oder
\[
x_{\beta\gamma}x_{\gamma\alpha}x_{\alpha\beta}=z_{\gamma\beta}z_{\alpha\gamma}z_{\beta\alpha},
\qquad
x_{\gamma\beta}x_{\alpha\gamma}x_{\beta\alpha}=z_{\beta\gamma}z_{\gamma\alpha}z_{\alpha\beta}.
\]
Da die Grössen $z_{\alpha\beta}$ lineare Functionen der $n^2$ unabhängigen Variabelen $x_{\alpha\beta}$ sind, so ist nach (1.) bis auf einen constanten Factor, von dem wir zunächst absehen wollen, entweder $z_{\alpha\beta}=x_{\alpha\beta}$, $z_{\beta\alpha}=x_{\beta\alpha}$ oder $z_{\alpha\beta}=x_{\beta\alpha}$, $z_{\beta\alpha}=x_{\alpha\beta}$.

Sei $z_{12}=x_{12}$, $z_{21}=x_{21}$. Dann kann nicht $x_{12}x_{2\alpha}x_{\alpha1}=z_{21}z_{\alpha2}z_{1\alpha}$ $(\alpha>2)$ sein, weil die rechte Seite nicht durch $x_{12}$ theilbar ist. Daher ist $x_{12}x_{2\alpha}x_{\alpha1}=z_{12}z_{2\alpha}z_{\alpha1}$, also $z_{\alpha1}=x_{\alpha1}$, $z_{2\alpha}=x_{2\alpha}$, und mithin $z_{1\alpha}=x_{1\alpha}$, $z_{\alpha2}=x_{\alpha2}$. Ferner muss dann $x_{1\alpha}x_{\alpha\beta}x_{\beta1}=z_{1\alpha}z_{\alpha\beta}z_{\beta1}$ und folglich $z_{\alpha\beta}=x_{\alpha\beta}$ sein. Ist dagegen $z_{12}=x_{21}$, so erkennt man in derselben Weise, dass allgemein $z_{\alpha\beta}=x_{\beta\alpha}$ ist.

Diese Gleichungen sind aber nur bis auf constante Factoren genau. Ist $z_{\alpha1}=\frac{k_\alpha}{k_1}x_{\alpha1}$, so ist nach (1.) $z_{1\alpha}=\frac{k_1}{k_\alpha}x_{1\alpha}$ und nach der Gleichung $z_{1\alpha}z_{\alpha\beta}z_{\beta1}=x_{1\alpha}x_{\alpha\beta}x_{\beta1}$ allgemein $z_{\alpha\beta}=\frac{k_\alpha}{k_\beta}x_{\alpha\beta}$. Setzt man also
\[
R=\begin{pmatrix}
k_1&0&0&\cdots\\
0&k_2&0&\cdots\\
0&0&k_3&\cdots\\
\cdot&\cdot&\cdot&\cdots
\end{pmatrix},
\]
so ist $Z=RXR^{-1}$ und $Y=PRXP^{-1}Q=AXB$. Ebenso ist in dem anderen Falle $z_{\alpha\beta}=\frac{k_\alpha}{k_\beta}x_{\beta\alpha}$, $Y=PRX'R^{-1}Q=AX'B$.

\begin{quote}\itshape
II. Sind die Voraussetzungen des Satzes I. erfüllt, und sind auch die charakteristischen Functionen der Matrizen $X$ und $Y$ einander gleich, so ist entweder $Y=AXA^{-1}$ oder $Y=AX'A^{-1}$.
\end{quote}

Sei $e_{\alpha\beta}=1$ oder 0, je nachdem $\alpha=\beta$ ist oder nicht. Dann sind $|X-rE|=|Y-rE|$ die charakteristischen Functionen von $X$ und $Y$. Setzt man $x_{\alpha\beta}=e_{\alpha\beta}$, so möge $y_{\alpha\beta}=c_{\alpha\beta}$ werden. Ersetzt man dann in der Gleichung $|Y|=|X|$ jedes $x_{\alpha\beta}$ durch $x_{\alpha\beta}-re_{\alpha\beta}$, so erhält man
\[
|Y-rC|=|X-rE|=|Y-rE|.
\]
Vergleicht man auf beiden Seiten die Coefficienten der ersten Potenzen von $r$, so findet man
\[
\sum c_{\alpha\beta}Y_{\alpha\beta}=\sum e_{\alpha\beta}Y_{\alpha\beta},
\]
wo $Y_{\alpha\beta}$ die Unterdeterminante $(n-1)^{\ten}$ Grades von $Y$ ist, die dem Elemente $y_{\alpha\beta}$ complementär ist. Da zwischen diesen Unterdeterminanten keine lineare Relation besteht, so ist $c_{\alpha\beta}=e_{\alpha\beta}$. Nun kann man nach Satz I die Matrizen $A$ und $B$ so bestimmen, dass $Y=AXB$ oder $Y=AX'B$ wird. Setzt man hier $X=E$, so wird, wie eben gezeigt, auch $Y=E$, und mithin ist $AB=E$, also $B=A^{-1}$.

Die entwickelten Sätze bleiben gültig, wenn die Veränderlichkeit der $n^2$ Variabelen $x_{\alpha\beta}$ durch die Relationen $x_{\alpha\beta}=x_{\beta\alpha}$ beschränkt wird:

\begin{quote}\itshape
III. Sind in einer symmetrischen Matrix $X$ die Elemente $x_{\alpha\beta}$ $(\beta\geqq\alpha)$ unabhängige Variable, und sind die Elemente der symmetrischen Matrix $Y$ lineare Functionen dieser Variabelen, und unterscheidet sich die Determinante $|Y|$ von $|X|$ nur um einen constanten von Null verschiedenen Factor, so ist $Y=AXA'$, wo $A$ eine bis auf das Vorzeichen völlig bestimmte constante Matrix ist. Sind ausserdem die charakteristischen Functionen von $X$ und $Y$ einander gleich, so ist $A$ eine orthogonale Matrix.
\end{quote}

Ist zunächst $AXA'=BXB'$ , so ergiebt sich, indem man $X=E$ setzt, $AA'=BB'$. Setzt man $A^{-1}B=A'B'^{-1}=F$, so ist $XF=FX$ und $FF'=E$. Setzt man die Variabelen $x_{\alpha\beta}$, deren Indices verschieden sind, gleich Null, so wird $x_{\alpha\alpha}f_{\alpha\beta}=f_{\beta\alpha}x_{\beta\beta}$, also $f_{\alpha\beta}=0$, wenn $\alpha$ von $\beta$ verschieden ist. Aus $FF'=E$ folgt dann $f_{\alpha\alpha}=\pm1$. Die allgemeine Gleichung $XF=FX$ ergiebt daher $x_{\alpha\beta}f_{\beta\beta}=f_{\alpha\alpha}x_{\alpha\beta}$. Mithin ist $F=\pm E$, und $B=\pm A$, und folglich ist $A$ bis auf das Vorzeichen völlig bestimmt.

Zwischen den ersten Ableitungen von $|X|$ nach den Variabelen $x_{\alpha\beta}$ $(\beta\geqq\alpha)$ besteht keine Relation, weil sich umgekehrt die Elemente $x_{\alpha\beta}$ als Functionen dieser Ableitungen darstellen lassen. Zwischen den Unterdeterminanten $m^{\ten}$ Grades $u,v,w,\ldots$ bestehen zwar lineare Relationen (\textsc{Kronecker}, \emph{Über die Subdeterminante symmetrischer Systeme}, Sitzungsberichte 1882). Ist $au+bv+cw+\cdots=0$ eine solche, und ist $u$ eine Hauptunterdeterminante, so muss $a=0$ sein. Denn setzt man alle Variabelen $x_{\alpha\beta}=0$ ausser denen, die in $u$ vorkommen, so wird $v=w=\cdots=0$.

Ist nun $c_{\alpha\beta}^{(\kappa)}$ der Coefficient von $x_{\kappa\kappa}$ in $y_{\alpha\beta}$, so ergiebt sich hieraus in derselben Weise wie oben, dass in der symmetrischen Matrix $c_{\alpha\beta}^{(\kappa)}$ alle Hauptunterdeterminanten zweiten und dritten Grades verschwinden. Wie ich in meiner Arbeit \emph{Über das Trägheitsgesetz der quadratischen Formen}, § 2, Satz 2 (Sitzungsberichte 1894) gezeigt habe, verschwinden folglich alle Unterdeterminanten zweiten Grades, und mithin ist $c_{\alpha\beta}^{(\kappa)}=p_{\alpha\kappa}p_{\beta\kappa}$. Nun folgen dieselben Schlüsse wie oben. Nur kann man, weil $z_{\alpha\beta}^2=x_{\alpha\beta}^2$ ist, die Grössen $k_\alpha=\pm1$ setzen. Dann ist $R^{-1}=R=R'$ und $Y=PRXRP'=AXA'$.

Sind die charakteristischen Functionen von $X$ und $Y$ einander gleich, so zeigt man, wie oben, dass für $X=E$ auch $Y=E$ wird, und folglich ist $AA'=E$, also ist $A$ eine orthogonale Matrix.

\begin{center}
\rule{40mm}{0.4pt}

Ausgegeben am 25. November.

\rule{40mm}{0.4pt}
\end{center}

\begin{flushright}
Berlin, gedruckt in der Reichsdruckerei.
\end{flushright}
% ===== END WORK 056 =====

\clearpage

% ===== BEGIN WORK 057: 057_ber_Relationen_zwischen_den_Charakteren_e =====
\begin{center}
{\Large Über Relationen zwischen den Charakteren einer Gruppe und denen ihrer Untergruppen}\\[0.75em]
Von G. \textsc{Frobenius}.
\end{center}

In meiner Arbeit \emph{Über Gruppencharaktere} (Sitzungsberichte 1896) habe ich zur Berechnung der Charaktere einer endlichen Gruppe von bekannter Constitution eine allgemeine Methode entwickelt und ihre praktische Verwendbarkeit an einer Reihe von einfachen Beispielen dargethan. Da aber ihre Anwendung auf complicirtere Gruppen mit erheblichen Schwierigkeiten verknüpft ist, so habe ich nach anderen Wegen gesucht, um die Charaktere einer Gruppe und damit ihre primitiven Darstellungen durch lineare Substitutionen zu erhalten, und ich habe zwei ganz verschiedene Methoden gefunden, die in speciellen Fällen leichter zu diesem Ziele führen können, als jene allgemeine Methode.

Die erste, die ich hier darlegen will, stützt sich auf die Betrachtung der in der gegebenen Gruppe $\Hgr$ enthaltenen Gruppen $\Ggr$ und auf die Beziehungen, die zwischen den Charakteren von $\Ggr$ und $\Hgr$ bestehen. Diese Relationen ergeben sich auf zwei verschiedenen Wegen. Der eine (§ 1) führt von den Primfactoren der Determinante der Gruppe $\Hgr$ zu denen der Determinante der Gruppe $\Ggr$, der andere (§ 3) umgekehrt von den letzteren zu den ersteren. Zu besonders einfachen Ergebnissen gelangt man durch diese Betrachtungen in dem Falle, wo $\Ggr$ eine invariante Untergruppe von $\Hgr$ ist (§§ 2, 4). Die erhaltenen Formeln stehen in naher Beziehung zu der Zerlegung der Gruppe $\Hgr$ in Complexe von Elementen, die nach einem Doppelmodul aequivalent sind, welcher aus zwei in $\Hgr$ enthaltenen Gruppen $\Ggr$ und $\Ggr'$ gebildet wird. Die Untersuchung des speciellen Falles, wo $\Ggr'=\Ggr$ ist, führt direct zur Ermittlung eines Charakters jeder zweifach transitiven Gruppe von Permutationen (§ 5).

Eine zweite Methode, um die Charaktere einer Gruppe zu berechnen, ergiebt sich aus der Theorie der Composition der Charaktere, die ich bei einer anderen Gelegenheit entwickeln werde.

\section*{§ 1.}

Sei $\Hgr$ eine Gruppe der Ordnung $h$, und sei
\[
  \Theta=\prod_\lambda \Phi_\lambda^{f_\lambda}
  \qquad(\lambda=0,1,\ldots,l-1)
\]
ihre in Primfactoren zerlegte Gruppendeterminante. Sei $\Ggr$ eine in $\Hgr$ enthaltene Gruppe der Ordnung $g=\frac{h}{n}$, und sei
\[
  \mathrm H=\prod_\kappa \Psi_\kappa^{e_\kappa}
  \qquad(\kappa=0,1,\ldots,k-1)
\]
ihre Gruppendeterminante. Setzt man in $\Theta$ alle Variabelen $x_R=0$ ausser denen, deren Indices die Elemente von $\Ggr$ sind, so wird (\emph{Gruppencharaktere} § 7)
\[
\tag{1.}
  \Theta=\mathrm H^n,
\]
und folglich wird auch jeder Primfactor $\Phi_\lambda$ von $\Theta$ ein Product von Primfactoren von $\mathrm H$,
\[
\tag{2.}
  \Phi_\lambda=\prod_\kappa \Psi_\kappa^{r_{\kappa\lambda}}.
\]
Hier ist $r_{\kappa\lambda}=0$ zu setzen, wenn $\Psi_\kappa$ nicht in $\Phi_\lambda$ aufgeht. Ersetzt man in dieser Gleichung $x_E$ durch $x_E+u$, so ergiebt sich durch Vergleichung der Coefficienten von $u^{f_\lambda-1}$
\[
\tag{3.}
  \sum_\kappa r_{\kappa\lambda}\psi^{(\kappa)}(P)=\chi^{(\lambda)}(P),
\]
wo $\psi^{(\kappa)}$ der Charakter von $\Psi_\kappa$ und $\chi^{(\lambda)}$ der von $\Phi_\lambda$ ist, und wo $P$ ein Element von $\Ggr$ bedeutet. Mit Hülfe der Gleichungen
\[
  \sum_P \psi^{(\kappa)}(P^{-1})\psi^{(\kappa)}(P)=g,
  \qquad
  \sum_P \psi^{(\kappa)}(P^{-1})\psi^{(\mu)}(P)=0
\]
folgt daraus
\[
\tag{4.}
  \sum_P \psi^{(\kappa)}(P^{-1})\chi^{(\lambda)}(P)=g r_{\kappa\lambda}.
\]
Ist daher $R$ ein Element von $\Hgr$, so ist
\[
  g\sum_\lambda r_{\kappa\lambda}\chi^{(\lambda)}(R)
  =
  \sum_P\psi^{(\kappa)}(P)
  \left(\sum_\lambda \chi^{(\lambda)}(P^{-1})\chi^{(\lambda)}(R)\right).
\]
Nach \emph{Gruppencharaktere} § 3, (4.) ist aber
\[
  \sum_\lambda \chi^{(\lambda)}(P^{-1})\chi^{(\lambda)}(R)=0,
\]
ausser wenn $P$ mit $R$ conjugirt ist (in Bezug auf $\Hgr$). Dann aber ist die Summe gleich $\frac{h}{h_R}=\frac{h}{h_\rho}$, falls $R$ und $P$ Elemente der $\rho^{\mathrm{ten}}$ Classe von $\Hgr$ sind. Folglich ist
\[
\tag{5.}
  \sum_\lambda r_{\kappa\lambda}\chi^{(\lambda)}(R)
  =
  \frac{h}{g h_\rho}\sum_{(\rho)}\psi^{(\kappa)}(P),
\]
wo $P$ die Elemente der $\rho^{\mathrm{ten}}$ Classe durchläuft, die in $\Ggr$ enthalten sind. Setzt man $\psi(R)=0$, wenn $R$ nicht in $\Ggr$ enthalten ist, so kann man für $P$ auch alle mit $R$ conjugirten Elemente von $\Hgr$ setzen. Bei dieser Festsetzung ist aber besonders darauf zu achten, dass die Gleichung (3.) nur für die Elemente $P$ gilt, die der Gruppe $\Ggr$ angehören.

Wenn die $k$ Charaktere $\psi^{(\kappa)}$ einer Untergruppe $\Ggr$ bekannt sind, so besteht die neue Eigenschaft der $l$ Charaktere $\chi^{(\lambda)}$, welche die Gleichungen (3.), (4.) und (5.) in drei verschiedenen, aber aequivalenten Formen ausdrücken, darin, dass
\[
\tag{4a.}
  \frac{1}{g}\sum_P \psi(P^{-1})\chi(P)=r
\]
eine positive ganze Zahl ist.

Für einen bestimmten Werth von $\lambda$ sind die $k$ Zahlen $r_{\kappa\lambda}$, und für einen bestimmten Werth von $\kappa$ die $l$ Zahlen $r_{\kappa\lambda}$ nicht sämmtlich Null. Dies ergiebt sich aus der Gleichung (2.) in Verbindung mit der Relation (1.).

Ist $(0)$ die Hauptclasse, so erhält man, wenn man in (3.) und (5.) $P=R=E$ setzt, die Gleichungen
\[
\tag{6.}
  \sum_\kappa r_{\kappa\lambda}e_\kappa=f_\lambda,
  \qquad
  \sum_\lambda r_{\kappa\lambda}f_\lambda=\frac{h}{g}e_\kappa,
\]
die sich auch unmittelbar aus (2.) und (1.) ergeben. Mithin ist
\[
\tag{7.}
  e_\kappa r_{\kappa\lambda}\le f_\lambda,
  \qquad
  r_{\kappa\lambda}\le f_\lambda.
\]
Wählt man für den Hauptcharakter den Index $0$, so ist
\[
\tag{8.}
  r_{00}=1,\qquad r_{\kappa 0}=0,\qquad(\kappa>0)
\]
und, wenn man $r_{0\lambda}=r_\lambda$ setzt,
\[
\tag{9.}
  \sum_\rho g_\rho\chi_\rho^{(\lambda)}=g r_\lambda,
  \qquad
  \sum_\lambda r_\lambda\chi_\rho^{(\lambda)}
  =
  \frac{h g_\rho}{g h_\rho},
\]
wo $g_\rho$ die Anzahl der Elemente in der $\rho^{\mathrm{ten}}$ Classe von $\Hgr$ bezeichnet, die der Gruppe $\Ggr$ angehören.

Da die Werthe der Charaktere ganze algebraische Zahlen sind, so ist der letzten Formel zufolge $h g_\rho$ durch $g h_\rho$ theilbar. Diesen Satz kann man leicht direct beweisen: Durchläuft $H$ die $h$ Elemente von $\Hgr$, und ist $R$ ein bestimmtes Element der $\rho^{\mathrm{ten}}$ Classe, so sind die $h$ Elemente $H^{-1}RH$ die $h_\rho$ verschiedenen Elemente der $\rho^{\mathrm{ten}}$ Classe, jedes $\frac{h}{h_\rho}$ Mal gezählt. In $\Hgr$ giebt es daher $g_\rho\frac{h}{h_\rho}$ verschiedene Elemente $H$ der Art, dass $H^{-1}RH$ in $\Ggr$ enthalten ist. Sei $\mathfrak R$ der Complex dieser Elemente $H$. Ist $A^{-1}RA$ in $\Ggr$ enthalten, und ist $G$ ein Element von $\Ggr$, so ist auch
\[
  G^{-1}(A^{-1}RA)G=(AG)^{-1}R(AG)
\]
in $\Ggr$ enthalten. Ist also $A$ ein Element von $\mathfrak R$, so gehören auch alle Elemente des Complexes $A\Ggr$ dem Complexe $\mathfrak R$ an. Folglich zerfällt $\mathfrak R$ in eine Anzahl Complexe $A\Ggr+B\Ggr+C\Ggr+\cdots$, von denen je zwei theilerfremd sind. Mithin ist die Ordnung $\frac{h g_\rho}{h_\rho}$ von $\mathfrak R$ durch $g$ theilbar. Die Zahl $\frac{h g_\rho}{g h_\rho}$ giebt an, wie viele unter den $n$ Complexen
\[
\tag{10.}
  A_0\Ggr+A_1\Ggr+\cdots+A_{n-1}\Ggr=\Hgr
\]
der Bedingung $A_\nu\Ggr=R A_\nu\Ggr$ genügen.

Nach Formel (3.) ist
\[
  \sum_P \chi^{(\lambda)}(P^{-1})\chi^{(\mu)}(P)
  =
  \sum_{\alpha,\beta} r_{\alpha\lambda}r_{\beta\mu}
  \left(\sum_P\psi^{(\alpha)}(P^{-1})\psi^{(\beta)}(P)\right)
\]
und mithin
\[
\tag{11.}
  \sum_P \chi^{(\lambda)}(P^{-1})\chi^{(\mu)}(P)
  =
  g\sum_\kappa r_{\kappa\lambda}r_{\kappa\mu},
\]
wo $P$ nur die Elemente von $\Ggr$ durchläuft, oder
\[
\tag{12.}
  \sum_\rho g_\rho\chi_{\rho'}^{(\lambda)}\chi_{\rho}^{(\mu)}
  =
  g\sum_\kappa r_{\kappa\lambda}r_{\kappa\mu}.
\]
Nun ist aber, wenn $R$ alle Elemente von $\Hgr$ durchläuft,
\[
  \sum_R \chi^{(\lambda)}(R^{-1})\chi^{(\lambda)}(R)=h.
\]
Da $\chi(R)$ und $\chi(R^{-1})$ conjugirte complexe Grössen sind, so ist ihr Product eine reelle positive Grösse. Mithin ist
\[
\tag{13.}
  \sum_\kappa r_{\kappa\lambda}^{2}\le \frac{h}{g}.
\]
Setzt man aber in (11.) $\mu=\lambda$ und summirt dann nach $\lambda$, so erhält man
\[
\tag{14.}
  \sum_{\kappa,\lambda}r_{\kappa\lambda}^{2}
  =
  \sum_\lambda \frac{h g_\lambda}{g h_\lambda},
\]
wo rechts über die $l$ Classen von $\Hgr$ zu summiren ist. Mit $g$ multiplicirt ist diese Zahl gleich der Anzahl der Lösungen der Gleichung $QR=RQ$, falls $Q$ die $g$ Elemente von $\Ggr$, und $R$ die $h$ Elemente von $\Hgr$ durchläuft.

Sind $\Ggr$ und $\Ggr'$ zwei Untergruppen von $\Hgr$, so mögen die Zahlen, die für die Gruppe $\Ggr$ mit $g,g_\rho,r_\lambda$ bezeichnet worden sind, für die Gruppe $\Ggr'$ mit $g',g'_\rho,r'_\lambda$ bezeichnet werden. Dann ist nach (2.)
\[
  gg'\sum_\lambda r_\lambda r'_\lambda
  =
  \sum_{\rho,\rho'}g_\rho g'_{\rho'}
  \left(\sum_\lambda \chi_\rho^{(\lambda)}\chi_{\rho'}^{(\lambda)}\right).
\]
Die letzte nach $\lambda$ genommene Summe ist Null, ausser wenn $(\rho)=(\rho')$ ist, dann aber gleich $\frac{h}{h_\rho}$. Mithin ist
\[
  gg'\sum_\lambda r_\lambda r'_\lambda
  =
  h\sum_\rho\frac{g_\rho g'_\rho}{h_\rho}.
\]
In meiner Arbeit \emph{Über die Congruenz nach einem aus zwei endlichen Gruppen gebildeten Doppelmodul}, \textsc{Crelle}'s Journal Bd. 101, habe ich § 2, (8.) für die Anzahl $(\Hgr:\Ggr,\Ggr')$ der Classen, worin die Elemente von $\Hgr$ nach dem Doppelmodul $(\Ggr,\Ggr')$ zerfallen, die Relation aufgestellt
\[
\tag{15.}
  \frac{gg'}{h}(\Hgr:\Ggr,\Ggr')
  =
  \sum_\rho\frac{g_\rho g'_\rho}{h_\rho}.
\]
Daraus ergiebt sich die Formel
\[
\tag{16.}
  \sum_\lambda r_\lambda r'_\lambda=(\Hgr:\Ggr,\Ggr'),
  \qquad
  \sum_\lambda r_\lambda^{2}=(\Hgr:\Ggr,\Ggr).
\]

Setzt man alle Variabelen $x_R=0$, deren Index $R$ nicht in $\Ggr$ enthalten ist, so enthält $\Phi_\lambda$ den Factor $\Psi_0^{r_\lambda}$. Ist $\Ggr'$ eine Untergruppe von $\Ggr$, so kann man jene Werthe in dem Ausdruck (2.) von $\Phi_\lambda$ einsetzen, dann geht $\Psi_0^{r_\lambda}$ in $\Psi_0^{\prime\,r'_\lambda}$ über. Folglich ist
\[
\tag{17.}
  r'_\lambda\ge r_\lambda,\quad\hbox{wenn }\Ggr'<\Ggr
\]
ist, d. h. wenn $\Ggr'$ in $\Ggr$ enthalten ist.

Ist $g<h$, so giebt es in $\Hgr$ stets Classen, von denen kein Element in $\Ggr$ enthalten ist. Denn unter der Voraussetzung (10.) sind $A_\nu\Ggr A_\nu^{-1}=\Ggr_\nu$, die $n$ mit $\Ggr$ conjugirten Gruppen, die nicht verschieden zu sein brauchen. Sei $t_\mu$ die Anzahl der Elemente von $\Hgr$, die in genau $\mu$ dieser $n$ Gruppen enthalten sind. Dann ist
\[
  t_0+t_1+t_2+\cdots+t_n=h.
\]
Andererseits enthält der Complex $\Ggr_0+\Ggr_1+\cdots+\Ggr_{n-1}$ $gn=h$ Elemente, falls man die mehrfach vorkommenden auch mehrfach zählt, also jedes der $t_\mu$ Elemente, die in genau $\mu$ jener $n$ Gruppen vorkommen, $\mu$-fach. Demnach ist
\[
  t_1+2t_2+\cdots+nt_n=h
\]
und folglich
\[
  t_0=t_2+2t_3+\cdots+(n-1)t_n.
\]
Da das Hauptelement $E$ in allen $n$ Gruppen vorkommt, so ist $t_n>0$. Ist also $n>1$, so ist $t_0>0$. Ist $R$ eins dieser $t_0$ Elemente, und durchläuft $H$ die $h$ Elemente von $\Hgr$, so kommt $R$ in keiner der Gruppen $H\Ggr H^{-1}$ vor, also ist keins der mit $R$ conjugirten Elemente $H^{-1}RH$ in $\Ggr$ enthalten. Nach Gleichung (5.) ist daher
\[
  \sum_\lambda r_{\kappa\lambda}\chi^{(\lambda)}(R)=0.
\]
Ist also $k\ge l$, so verschwinden in der Matrix
\[
\tag{18.}
  r_{\kappa\lambda}\qquad(\kappa=0,1,\ldots,k-1;\ \lambda=0,1,\ldots,l-1)
\]
alle Determinanten $l^{\mathrm{ten}}$ Grades.

\section*{§ 2.}

Die $kl$ Zahlen $r_{\kappa\lambda}$ lassen sich näher bestimmen, wenn $\Ggr$ eine invariante Untergruppe von $\Hgr$ ist. Für diesen Fall habe ich, \emph{Gruppencharaktere} § 7, folgende Sätze entwickelt: Sei $S$ ein festes Element von $\Hgr$, und $P$ ein veränderliches Element von $\Ggr$. Da $S^{-1}\Ggr S=\Ggr$ ist, so durchläuft $S^{-1}PS$ gleichzeitig mit $P$ die $g$ Elemente von $\Ggr$. Setzt man $\psi^{(\kappa)}(S^{-1}PS)=\psi^{(\kappa')}(P)$, so ist $\psi^{(\kappa')}(P)$ ein Charakter von $\Ggr$. Ist er von $\psi^{(\kappa)}(P)$ verschieden, so nenne ich ihn zu $\psi^{(\kappa)}(P)$ conjugirt (in Bezug auf $\Hgr$). Conjugirte Charaktere haben denselben Grad $e_\kappa=e_{\kappa'}$ und unterscheiden sich nur durch die Anordnung ihrer Werthe. Sei $s_\kappa=s_{\kappa'}=\cdots$ die Anzahl der verschiedenen Charaktere von $\Ggr$, die mit $\psi^{(\kappa)}$ und folglich auch unter einander conjugirt sind. Jedem Charakter $\psi^{(\kappa)}$ von $\Ggr$ entsprechen ein oder mehrere Charaktere $\chi^{(\lambda)},\chi^{(\lambda')},\ldots$ von $\Hgr$, so dass
\[
\tag{1.}
  \psi^{(\kappa)}(P)+\psi^{(\kappa')}(P)+\cdots
  =
  \frac{s_\kappa e_\kappa}{f_\lambda}\chi^{(\lambda)}(P)
  =
  \frac{s_\kappa e_\kappa}{f_{\lambda'}}\chi^{(\lambda')}(P)
  =\cdots
\]
ist für alle $g$ Elemente $P$ von $\Ggr$. Umgekehrt entspricht jedem Charakter $\chi^{(\lambda)}$ von $\Hgr$ mindestens ein Charakter $\psi^{(\kappa)}$ von $\Ggr$, so dass die Gleichung (1.) erfüllt wird, und wenn ihm mehrere entsprechen, so sind je zwei derselben conjugirt. Ist $\chi^{(\lambda)}$ gegeben, so sind $\chi^{(\lambda')},\chi^{(\lambda'')},\ldots$ dadurch bestimmt, dass die $g$ Werthe $\chi^{(\lambda)}(P)$ den $g$ Werthen $\chi^{(\lambda')}(P)$ proportional sind. Durchläuft $R$ die $h$ Elemente von $\Hgr$, und sind $\chi^{(\lambda)}$ und $\chi^{(\mu)}$ verschiedene Charaktere, so können nach \emph{Gruppencharaktere} § 3 die $h$ Werthe $\chi^{(\lambda)}(R)$ den $h$ Werthen $\chi^{(\mu)}(R)$ nicht proportional sein.

Setzt man nun den Ausdruck (1.) für $\chi^{(\lambda)}(P)$ in die Formel (3.), § 1 ein, so erkennt man, dass stets $r_{\kappa\lambda}=0$ ist, ausser wenn sich $\psi^{(\kappa)}$ und $\chi^{(\lambda)}$ entsprechen. In diesem Falle aber ist
\[
\tag{2.}
  r_{\kappa\lambda}=r_{\kappa'\lambda}=\cdots=\frac{f_\lambda}{s_\kappa e_\kappa}
\]
von Null verschieden, wo $\psi^{(\kappa)},\psi^{(\kappa')},\ldots$ die $s_\kappa$ unter einander conjugirten, dem $\chi^{(\lambda)}$ entsprechenden Charaktere von $\Ggr$ sind. Demnach reducirt sich die Gleichung (2.), § 1 auf
\[
\tag{3.}
  \Phi_\lambda=(\Psi_\kappa\Psi_{\kappa'}\cdots)^{\frac{f_\lambda}{s_\kappa e_\kappa}}
\]
und aus (6.), § 1 folgt
\[
\tag{4.}
  f_\lambda^2+f_{\lambda'}^2+\cdots=ns_\kappa e_\kappa^2,
  \qquad
  r_{\kappa\lambda}^2+r_{\kappa\lambda'}^2+\cdots=\frac{n}{s_\kappa}.
\]
Folglich ist $f_\lambda$ durch $s_\kappa e_\kappa$ theilbar, und der Quotient $\frac{f_\lambda}{s_\kappa e_\kappa}$ ist ein gemeinsamer Divisor der $g$ ganzen algebraischen Zahlen $\chi^{(\lambda)}(P)$. Dass $n$ durch $s_\kappa$ theilbar ist, habe ich schon \emph{Gruppencharaktere} § 7 gezeigt.

Ist $R$ nicht in $\Ggr$ enthalten, so gehört auch kein mit $R$ conjugirtes Element $P$ der Gruppe $\Ggr$ an. Mithin folgt aus (5.), § 1
\[
\tag{5.}
  f_\lambda\chi^{(\lambda)}(R)+f_{\lambda'}\chi^{(\lambda')}(R)+\cdots=0,
\]
wenn $R$ nicht in $\Ggr$ enthalten ist.

\section*{§ 3.}

Wir sind von den $l$ Primfactoren $\Phi$ der Determinante $\Theta$ der Gruppe $\Hgr$ zu den $k$ Primfactoren $\Psi$ der Determinante $\mathrm H$ der Untergruppe $\Ggr$ gelangt, indem wir alle Variabelen $x_R=0$ setzten, deren Index $R$ ein in $\Ggr$ nicht enthaltenes Element von $\Hgr$ ist. Ich will jetzt zeigen, wie man durch eine andere Construction von den $k$ Functionen $\Psi$ zu den $l$ Functionen $\Phi$ aufsteigen kann.

Sei $X$ eine zur Gruppe $\Ggr$ gehörige Matrix des Grades $e$. Ihre Elemente sind lineare Functionen der $g$ Variabelen $x_P$, und ihre Determinante verschwindet nicht identisch. Sie ist durch folgende Eigenschaft charakterisirt: Ersetzt man $x_P$ durch $y_P$ oder $z_P$, so möge $X$ in $Y$ oder $Z$ übergehen. Ist dann
\[
  z_Q=\sum_P x_{P^{-1}}y_{PQ},
\]
so ist $Z=XY$. Ersetzt man in $X$ jede der $g$ Variabelen $x_P$ durch $x_{APB^{-1}}$, wo $A$ und $B$ zwei Elemente von $\Hgr$ sind, so erhält man eine Matrix, die ich mit $X_{A,B}$ bezeichne. Die $n$ Elemente $A_0,A_1,\ldots,A_{n-1}$ mögen ein vollständiges Restsystem von $\Hgr(\mathrm{mod.}\,\Ggr)$ bilden, so dass
\[
\tag{1.}
  \Hgr=A_0\Ggr+A_1\Ggr+\cdots+A_{n-1}\Ggr
  =
  \Ggr A_0^{-1}+\Ggr A_1^{-1}+\cdots+\Ggr A_{n-1}^{-1}.
\]
Dann betrachte ich die $n^2$ Matrizen $e^{\mathrm{ten}}$ Grades, die man aus $X_{A,B}$ erhält, indem man für $A$ und $B$ jedes der $n$ Elemente $A_0,A_1,\ldots,A_{n-1}$ setzt, und bilde aus ihnen eine Matrix $ne^{\mathrm{ten}}$ Grades $(X_{A,B})$. Ersetzt man $x_R$ durch $y_R$ oder $z_R$, so gehe $X_{A,B}$ in $Y_{A,B}$ oder $Z_{A,B}$ über. Sind $A,B$ und $N$ Elemente von $\Hgr$, so geht die Matrix $X_{A,N}Y_{N,B}$ aus $X$ hervor, indem man $x_Q$ durch
\[
  \sum_P x_{AP^{-1}N^{-1}}y_{NPQB^{-1}}
\]
ersetzt. Hier sind $P$ und $Q$ wie oben Elemente von $\Ggr$. Nun sind die Elemente von $X$ lineare Functionen der $g$ Variabelen $x_Q$. Daher gehen die Elemente der Matrix
\[
  \sum_N X_{A,N}Y_{N,B}
  \qquad(N=A_0,A_1,\ldots,A_{n-1})
\]
aus $X$ hervor, indem man $x_Q$ durch
\[
  \sum_{N,P}x_{AP^{-1}N^{-1}}y_{NPQB^{-1}}
  =
  \sum_R x_{AR^{-1}}y_{RQB^{-1}}
\]
ersetzt. Durchläuft $P$ die $g$ Elemente von $\Ggr$ und $N$ die $n$ Elemente $A_0,A_1,\ldots,A_{n-1}$, so durchläuft $R=NP$ nach (1.) die $gn=h$ verschiedenen Elemente von $\Hgr$. Setzt man also jetzt, wenn $R$ und $S$ Elemente von $\Hgr$ sind,
\[
  z_S=\sum_R x_{R^{-1}}y_{RS},
\]
so ist die letzte Summe gleich $z_{AQB^{-1}}$, und mithin ist
\[
  \sum_N X_{A,N}Y_{N,B}=Z_{A,B}
\]
oder
\[
  (X_{A,B})(Y_{A,B})=(Z_{A,B}).
\]
Folglich ist $(X_{A,B})$ eine zur Gruppe $\Hgr$ gehörige Matrix des Grades $ne$, und daher ist ihre Determinante ein Product von Primfactoren der Gruppendeterminante $\Theta$
\[
\tag{2.}
  |(X_{A,B})|=\prod_\lambda \Phi_\lambda^{r_\lambda}.
\]

In jeder zu einer Gruppe gehörigen Matrix ist $x_E$ mit der Hauptmatrix multiplicirt. Ersetzt man also darin $x_E$ durch $x_E+u$, so tritt nur zu jedem Elemente der Diagonale das Glied $u$ hinzu. Daher ist der Coefficient von $u^{ne-1}$ in der Determinante $ne^{\mathrm{ten}}$ Grades (2.) gleich der Summe der Diagonalelemente.

Ist $\Psi(x)$ ein Primfactor $e^{\mathrm{ten}}$ Grades von $\mathrm H$, so kann man eine zu $\Ggr$ gehörige Matrix $X$ finden, deren Determinante gleich $\Psi(x)$ ist (\emph{Über die Darstellung der endlichen Gruppen durch lineare Substitutionen}, Sitzungsberichte 1897). Der Coefficient von $u^{e-1}$ in $\Psi(x+u\varepsilon)$ ist $\sum_P\psi(P)x_P$. Für die Matrix $X_{N,N}$ ist daher die Summe der Diagonalelemente gleich
\[
  \sum_P\psi(P)x_{NPN^{-1}}
  =
  \sum_R\psi(N^{-1}RN)x_R.
\]
Hier können $P$ und $R=NPN^{-1}$ alle Elemente von $\Hgr$ durchlaufen, wenn man wie oben festsetzt, dass $\psi(R)=0$ ist, wenn $R$ der Gruppe $\Ggr$ nicht angehört. Durch Vergleichung der Coefficienten von $u^{ne-1}$ in der Gleichung (2.) erhält man daher, wenn $r_\lambda=r_{\kappa\lambda}$ für $\Psi=\Psi_\kappa$ gesetzt wird,
\[
\tag{3.}
  \sum_N\psi^{(\kappa)}(N^{-1}RN)
  =
  \sum_\lambda r_{\kappa\lambda}\chi^{(\lambda)}(R).
\]
Ist $P$ ein Element von $\Ggr$, so ist $\psi(P^{-1}SP)=\psi(S)$. Ist nämlich $S$ ein Element von $\Ggr$, so ist dies die Gleichung \emph{Gruppencharaktere} § 5, (2.). Ist aber $S$ nicht in $\Ggr$ enthalten, so gehört auch $P^{-1}SP$ nicht der Gruppe $\Ggr$ an und beide Seiten der Gleichung sind Null. Daher ist
\[
  g\sum_N\psi(N^{-1}RN)
  =
  \sum_{N,P}\psi(P^{-1}N^{-1}RNP)
  =
  \sum_S\psi(S^{-1}RS),
\]
wo $S=NP$ die $h$ Elemente von $\Hgr$ durchläuft, also
\[
\tag{4.}
  g\sum_\lambda r_{\kappa\lambda}\chi^{(\lambda)}(R)
  =
  \sum_S\psi^{(\kappa)}(S^{-1}RS).
\]
Ist $R$ ein Element der $\rho^{\mathrm{ten}}$ Classe in $\Hgr$, so stellt $S^{-1}RS$ jedes der $h_\rho$ verschiedenen Elemente dieser Classe $\frac{h}{h_\rho}$ Mal dar. Mithin ist
\[
\tag{5.}
  \sum_\lambda r_{\kappa\lambda}\chi^{(\lambda)}(R)
  =
  \frac{h}{g h_\rho}\sum_{(\rho)}\psi^{(\kappa)}(P),
\]
wo $P$ die $h_\rho$ Elemente der $\rho^{\mathrm{ten}}$ Classe durchläuft, oder auch nur die unter ihnen, die in $\Ggr$ enthalten sind. Da die $kl$ Zahlen $r_{\kappa\lambda}$ durch diese Gleichungen vollständig bestimmt sind, so sind sie mit den in den Gleichungen
\[
\tag{6.}
  \Phi_\lambda=\prod_\kappa\Psi_\kappa^{r_{\kappa\lambda}}
\]
auftretenden Exponenten identisch. In Folge der Ungleichheiten (7.) § 1 ist die $e_\kappa^{\mathrm{te}}$ Potenz des Ausdrucks
\[
\tag{7.}
  |(X^{(\kappa)}_{A,B})|=\prod_\lambda\Phi_\lambda^{r_{\kappa\lambda}}
\]
ein Divisor der Gruppendeterminante $\Theta$.

Ist $\Ggr$ eine invariante Untergruppe von $\Hgr$, so wird unter Anwendung der obigen Bezeichnungen
\[
\tag{8.}
  |(X^{(\kappa)}_{A,B})|^{s_\kappa e_\kappa}
  =
  \Phi_\lambda^{f_\lambda}\Phi_{\lambda'}^{f_{\lambda'}}\cdots,
\]
also ein Divisor von $\Theta$, der zu dem complementären Divisor theilerfremd ist.

\section*{§ 4.}

In dem besonders bemerkenswerthen Falle $\kappa=0$ setze ich zur Vereinfachung der Darstellung, wenn
\[
  \mathfrak A=P+Q+R+\cdots
\]
ein Complex von Elementen ist,
\[
\tag{1.}
  x_{\mathfrak A}=x_P+x_Q+x_R+\cdots.
\]
Dann ist die Matrix des $n^{\mathrm{ten}}$ Grades
\[
\tag{2.}
  (x_{A\Ggr B^{-1}})\qquad(A,B=A_0,A_1,\ldots,A_{n-1})
\]
eine zu $\Hgr$ gehörige Matrix, und ihre Determinante ist
\[
\tag{3.}
  |x_{A\Ggr B^{-1}}|=\prod_\lambda \Phi_\lambda^{r_\lambda}.
\]
Jeder zu $\Hgr$ gehörigen Matrix entspricht eine Darstellung der Gruppe $\Hgr$ oder einer mit $\Hgr$ meroedrisch isomorphen Gruppe durch lineare Substitutionen. Der Matrix (2.) entspricht die Darstellung einer mit $\Hgr$ isomorphen Gruppe durch Permutationen von $n$ Symbolen, die ich in meiner Arbeit \emph{Über endliche Gruppen} § 4 (Sitzungsberichte 1895) entwickelt habe.

Ist $\Ggr$ eine invariante Untergruppe von $\Hgr$, so ist, da $e_0=s_0=1$ ist,
\[
\tag{4.}
  |x_{AB^{-1}\Ggr}|=\prod_\lambda\Phi_\lambda^{r_\lambda}
  =
  \Phi_\nu^{f_\nu}\Phi_{\nu'}^{f_{\nu'}}\cdots,
\]
wo $\chi^{(\nu)},\chi^{(\nu')},\ldots$ die Charaktere von $\Hgr$ sind, die dem Charakter $\psi^{(0)}$ von $\Ggr$ entsprechen. Die linke Seite ist die Determinante der Gruppe $\frac{\Hgr}{\Ggr}$, die ich \emph{Darstellung} § 1 betrachtet habe. Auch dort habe ich gezeigt, dass die von Null verschiedenen Exponenten $r_\lambda$ gleich $f_\lambda$ sind. Aus jener Gleichung ergiebt sich noch eine bemerkenswerthe Folgerung: Die Indices $\lambda=\nu,\nu',\ldots$ sind dadurch charakterisirt, dass $\chi^{(\lambda)}(P)$ für alle $g$ Elemente $P$ der Gruppe $\Ggr$ denselben Werth hat. Da aber die linke Seite die $h$ Grössen $x_R$ nur in den $n$ linearen Verbindungen $x_{\mathfrak A}$ enthält, so gilt dasselbe von jeder der Primfunctionen $\Phi_\nu,\Phi_{\nu'},\ldots$, und folglich hat auch, wenn $N$ ein festes Element von $\Hgr$ und $P$ ein veränderliches Element von $\Ggr$ ist, $\chi^{(\lambda)}(NP)$ für alle Elemente $P$ von $\Ggr$ denselben Werth. Demnach ergiebt sich der Satz:

\emph{I. Damit ein Charakter von $\Hgr$ zu der Gruppe $\frac{\Hgr}{\Ggr}$ gehöre, ist nothwendig und hinreichend, dass er für alle Elemente von $\Ggr$ denselben Werth hat. Dann hat er auch gleiche Werthe für je zwei Elemente von $\Hgr$, die mod. $\Ggr$ aequivalent sind.}

Man kann diesen Satz auch aus der Formel (\emph{Gruppencharaktere} § 5) ableiten
\[
\tag{5.}
  h_\beta\chi(A)\chi(B)=f\sum_{(S)}\chi(AS),
\]
worin $S$ die $h_\beta$ mit $B$ conjugirten Elemente durchläuft. Da mithin $\chi(S)=\chi(B)$ ist, so kann man diese Gleichung auch in der Form
\[
  \chi(R)\sum_{(\Ggr)}\chi(S)=f\sum_{(\Ggr)}\chi(RS)
\]
schreiben. Ist $\Ggr$ eine invariante Untergruppe von $\Hgr$, so enthält sie entweder kein Element der $\rho^{\mathrm{ten}}$ Classe oder alle. Setzt man für $(\Ggr)$ der Reihe nach sämmtliche Classen, deren Elemente in $\Ggr$ enthalten sind, so findet man durch Summation der entsprechenden Gleichungen
\[
  \chi(R)\sum\chi(P)=f\sum\chi(P'),
\]
wo $P$ die Elemente der Gruppe $\Ggr$ durchläuft, und $P'$ die des Complexes $R\Ggr$. Ist daher $\sum\chi(P)=0$, so ist auch für jedes Element $R$ von $\Hgr$ $\sum\chi(P')=0$. Ist ferner $S\sim R(\mathrm{mod.}\,\Ggr)$, so ist $R\Ggr=S\Ggr$ und folglich auch
\[
  \chi(S)\sum\chi(P)=f\sum\chi(P'),
\]
wo $P'$ dieselben Elemente durchläuft, wie oben. Wenn also $\sum\chi(P)$ von Null verschieden ist, so muss sein
\[
\tag{6.}
  \chi(R)=\chi(S),\qquad \hbox{falls } R\sim S(\mathrm{mod.}\,\Ggr).
\]
Der letzte Fall tritt sicher ein, wenn die $g$ Werthe $\chi(P)$ alle einander gleich, also alle gleich $\chi(E)=f$ sind. Dann gehört demnach der Charakter $\chi=\chi^{(\nu)}$ zur Gruppe $\frac{\Hgr}{\Ggr}$, die Function $\Phi_\nu$ enthält die $h$ Variabelen $x_R$ nur in den $n$ linearen Verbindungen $x_{\mathfrak A}$ und wird bis auf einen Zahlenfactor dem entsprechenden Primfactor der Gruppe $\frac{\Hgr}{\Ggr}$ gleich, wenn man darin setzt
\[
\tag{7.}
  x_R=x_S,\qquad \hbox{falls } R\sim S(\mathrm{mod.}\,\Ggr).
\]
Gehört aber der Charakter $\chi=\chi^{(\lambda)}$ nicht zur Gruppe $\frac{\Hgr}{\Ggr}$, so ist die Gleichung (6.) nicht für je zwei aequivalente Elemente erfüllt, daher ist $\sum_P\chi(P)=0$ und folglich auch für jedes Element $N$ von $\Hgr$
\[
\tag{8.}
  \sum_P\chi(NP)=0.
\]

Der Satz I lässt sich theilweise umkehren. Aus der Formel (5.) ergiebt sich, wenn $\chi(A)=\chi(B)=f$ ist, die Gleichung
\[
  h_\beta f=\sum\chi(AS).
\]
Da $\chi(R)$ eine Summe von $f$ Einheitswurzeln ist, so ist die rechte Seite eine Summe von $h_\beta f$ Einheitswurzeln. Eine solche Summe kann aber nur dann gleich $h_\beta f$ sein, wenn jedes Glied gleich 1 ist. Mithin ist $\chi(AS)=f$. Ist also $\chi(A)=f$ und $\chi(B)=f$, so ist auch $\chi(AB)=f$. Folglich bilden alle Elemente $R$ von $\Hgr$, für die $\chi(R)=f$ ist, eine Gruppe $\Ggr$. Enthält diese das Element $B$, so enthält sie auch alle mit $B$ conjugirten Elemente $S$, weil $\chi(S)=\chi(B)$ ist. Daher ist $\Ggr$ eine invariante Untergruppe von $\Hgr$.

\emph{II. Ist $\chi(R)$ ein Charakter $f^{\mathrm{ten}}$ Grades der Gruppe $\Hgr$, so bilden alle Elemente $R$ von $\Hgr$, für die $\chi(R)=f$ ist, eine invariante Untergruppe $\Ggr$ von $\Hgr$, und der Charakter $\chi$ gehört zu der Gruppe $\frac{\Hgr}{\Ggr}$.}

Mit Hülfe dieser beiden Sätze kann man, wenn die Charaktere einer Gruppe bekannt sind, ihre invarianten Untergruppen sämmtlich angeben. Dies Verfahren ist in dem speciellen Falle, wo $\Hgr$ eine commutative Gruppe ist, schon von \textsc{Weber} angegeben.

Unter den $h$ linearen Functionen
\[
  \frac{h}{f}\xi_R=\sum_S\chi(S)x_{RS}
\]
sind $f^2$ unter einander unabhängige, und durch diese $f^2$ Verbindungen der $h$ Variabelen $x_R$ lässt sich $\Phi$ ausdrücken. Durchläuft $P$ die $g$ Elemente von $\Ggr$ und $N$ ein vollständiges Restsystem von $\Hgr(\mathrm{mod.}\,\Ggr)$, so durchläuft $S=NP$ die $h$ Elemente von $\Hgr$. Daher ist
\[
  \frac{h}{f}\xi_R=\sum_{N,P}\chi(NP)x_{RNP},
\]
und folglich ist $\xi_R=0$, wenn die Voraussetzung (7.) gemacht wird. Wird $\Phi$ auf irgend eine Art als Function von $f^2$ unabhängigen Variabelen dargestellt, so sind diese lineare Verbindungen der Variabelen $\xi_R$ und haben daher dieselbe Eigenschaft. Es gilt also der Satz:

\emph{III. Ist $\Phi$ ein Primfactor $f^{\mathrm{ten}}$ Grades von der Determinante der Gruppe $\Hgr$, ist $\Ggr$ eine invariante Untergruppe von $\Hgr$, und setzt man in $\Phi$ stets $x_R=x_S$, falls $R\sim S(\mathrm{mod.}\,\Ggr)$ ist, so wird die Function $\Phi$, falls ihr Charakter $\chi$ zu $\frac{\Hgr}{\Ggr}$ gehört, einem Primfactor dieser Gruppe gleich; wenn aber $\chi$ nicht zu $\frac{\Hgr}{\Ggr}$ gehört, so verschwindet jede der $f^2$ unabhängigen Variabelen, durch die sich $\Phi$ darstellen lässt.}

Bringt man nun die Gruppenmatrix $X$ auf die reducirte Form $L^{-1}XL$ (\emph{Darstellung}, § 5), so geht diese durch die Annahme (7.) in die reducirte Form der Matrix der Gruppe $\frac{\Hgr}{\Ggr}$ über, der Rang von $L^{-1}XL$ wird $\frac{h}{g}=n$, und von ihren Unterdeterminanten des Grades $n$ ist nur eine von Null verschieden. Folglich wird auch der Rang der Gruppenmatrix $X$ gleich $n$ und jede Unterdeterminante $n^{\mathrm{ten}}$ Grades von $X$ wird bis auf einen constanten Factor gleich der Determinante der Gruppe $\frac{\Hgr}{\Ggr}$.

\section*{§ 5.}

Die Formel (5.), § 1 und die darin enthaltene Formel (9.), § 1
\[
\tag{1.}
  \sum_\lambda r_{\kappa\lambda}\chi^{(\lambda)}(R)
  =
  \frac{h}{g h_\rho}\sum_{(\rho)}\psi^{(\kappa)}(P),
  \qquad
  \sum_\lambda r_\lambda\chi_\rho^{(\lambda)}
  =
  \frac{h g_\rho}{g h_\rho},
\]
sind besonders dazu geeignet, aus den Charakteren der Untergruppe $\Ggr$ Charaktere der Gruppe $\Hgr$ abzuleiten. Mit Hülfe der letzteren Formel ist es mir, wie ich bei einer anderen Gelegenheit darlegen will, gelungen, die Charaktere der symmetrischen Gruppe des Grades $n$ allgemein zu bestimmen. Eine besonders einfache Anwendung dieser Gleichung bildet der folgende Satz:

\emph{Enthält die Gruppe $\Hgr$ der Ordnung $h$ die Gruppe $\Ggr$ der Ordnung $g$, besteht die $\rho^{\mathrm{te}}$ Classe conjugirter Elemente in $\Hgr$ aus $h_\rho$ Elementen, und gehören davon $g_\rho$ der Gruppe $\Ggr$ an, so besteht die nothwendige und hinreichende Bedingung dafür, dass die Grössen}
\[
\tag{1.}
  \chi_\rho=\frac{h g_\rho}{g h_\rho}-1
\]
\emph{einen Charakter von $\Hgr$ bilden, darin, dass die Anzahl der Classen, worin die Elemente von $\Hgr$ nach dem Doppelmodul $(\Ggr,\Ggr)$ zerfallen, gleich zwei ist.}

Ist $\chi_\rho$ ein Charakter, so ist
\[
\tag{2.}
  \sum_\rho h_\rho\chi_\rho\chi_{\rho'}=h,
\]
also
\[
  \sum_\rho h_\rho\left(\frac{h g_\rho}{g h_\rho}-1\right)^2=h,
  \qquad
  \sum_\rho\frac{g_\rho^2}{h_\rho}=\frac{2g^2}{h}.
\]
Nach der Formel (15.), § 1 ist folglich
\[
\tag{3.}
  (\Hgr:\Ggr,\Ggr)=2.
\]
Das Hauptelement $E$ repraesentirt (modd. $\Ggr,\Ggr$) den Complex $\Ggr E\Ggr=\Ggr$. Ist $L$ ein Element von $\Hgr$, das nicht in $\Ggr$ enthalten ist, so bilden die mit $L$ aequivalenten Elemente von $\Hgr$ den Complex $\Ggr L\Ggr$. Da $(\Hgr:\Ggr,\Ggr)=2$ ist, so ist folglich
\[
\tag{4.}
  \Hgr=\Ggr+\Ggr L\Ggr.
\]
Zerlegt man also $\Hgr(\mathrm{mod.}\,\Ggr)$ in $n$ verschiedene Complexe
\[
  \Hgr=\Ggr+P\Ggr+Q\Ggr+R\Ggr+\cdots,
\]
so giebt es in $\Ggr$ ein solches Element $G$, dass $G(P\Ggr)=Q\Ggr$ ist.

Endlich kann man die gefundene Bedingung auch so ausdrücken: Ist $\Dgr$ der grösste gemeinsame Divisor aller mit $\Ggr$ conjugirten Untergruppen von $\Hgr$, so lässt sich immer $\frac{\Hgr}{\Dgr}$ als transitive Gruppe von Permutationen von $n$ Symbolen in der Art darstellen, dass die Untergruppe $\frac{\Ggr}{\Dgr}$ von allen Permutationen gebildet wird, die ein bestimmtes Symbol ungeändert lassen. Die obige Bedingung besteht nun darin, dass diese Gruppe von Permutationen zweifach transitiv ist. Wird die dem Elemente $R$ entsprechende Permutation in ihre cyclischen Factoren zerlegt, so ist $1+\chi(R)$ nach (10.), § 1 gleich der Anzahl der Cyclen ersten Grades oder gleich der Anzahl der Symbole, die jene Permutation ungeändert lässt.

Dass die Bedingung (3.) auch hinreichend ist, ergiebt sich aus der Formel (16.), § 1, wonach $\sum r_\lambda^2=2$ ist. Da $r_0=1$ ist, so ist folglich eine und nur eine Zahl $r_\nu=1$, jede der anderen $l-2$ Zahlen $r_\lambda=0$. Ist dann $\chi_\rho^{(\nu)}=\chi_\rho$, so ist nach (9.), § 1
\[
  \chi_\rho^{(0)}+\chi_\rho^{(\nu)}
  =
  \frac{h g_\rho}{g h_\rho},
  \qquad
  \chi_\rho=\frac{h g_\rho}{g h_\rho}-1.
\]

Man kann aber auch direct beweisen, dass unter der Bedingung (4.) die Grössen $\chi_\rho$ den Gleichungen genügen, die zur Berechnung der Charaktere dienen. Von den Elementen der $\rho^{\mathrm{ten}}$ Classe sind $h_\rho-g_\rho$ nicht in $\Ggr$ enthalten. Von diesen $h_\rho-g_\rho$ Elementen finden sich in jedem der $n-1$ Complexe $P\Ggr,Q\Ggr,R\Ggr,\ldots$ gleich viele, also $\frac{h_\rho-g_\rho}{n-1}$. Denn seien
\[
  PA,\qquad PB,\qquad PC,\ldots
\]
die Elemente der $\rho^{\mathrm{ten}}$ Classe, die dem Complexe $P\Ggr$ angehören, so dass $A,B,C,\ldots$ Elemente von $\Ggr$ sind. Nun giebt es in $\Ggr$ ein solches Element $G$, dass $GP\Ggr=Q\Ggr$ ist. Dann sind die Elemente
\[
  GPAG^{-1},\ GPBG^{-1},\ GPCG^{-1},\ldots
\]
unter einander verschieden und den obigen Elementen conjugirt, also auch in der $\rho^{\mathrm{ten}}$ Classe enthalten. Endlich gehören sie dem Complexe $Q\Ggr$ an. Dieser enthält demnach nicht weniger Elemente der $\rho^{\mathrm{ten}}$ Classe wie $P\Ggr$, und da dasselbe umgekehrt gilt, so enthält jeder der $n-1$ Complexe $P\Ggr,Q\Ggr,R\Ggr,\ldots$ gleich viele Elemente der $\rho^{\mathrm{ten}}$ Classe.

Nun seien $(\alpha')$ und $(\beta)$ irgend zwei gleiche oder verschiedene Classen, und sei $A$ ein Element von $(\alpha')$, also $A^{-1}$ ein Element von $(\alpha)$, und $B$ ein Element von $(\beta)$. Sind $A$ und $B$ beide in $\Ggr$ enthalten, so ist auch $A^{-1}B$ in $\Ggr$ enthalten. Ist von diesen beiden Elementen das eine in $\Ggr$ enthalten, das andere nicht, so ist $A^{-1}B$ nicht in $\Ggr$ enthalten. Sind beide nicht in $\Ggr$ enthalten, so ist $A^{-1}B$ in $\Ggr$ enthalten oder nicht, je nachdem $A$ und $B$ beide demselben Complexe $P\Ggr$ (oder $Q\Ggr$, oder $R\Ggr,\ldots$) angehören oder nicht. Durchläuft daher $A$ die $h_\alpha$ Elemente von $(\alpha')$ und $B$ die $h_\beta$ Elemente von $(\beta)$, so sind von den $h_\alpha h_\beta$ Elementen $A^{-1}B$
\[
  g_\alpha g_\beta+(n-1)\frac{h_\alpha-g_\alpha}{n-1}\frac{h_\beta-g_\beta}{n-1}
\]
in $\Ggr$ enthalten.

Ist $C$ ein Element der Classe $(\gamma)$, so sind $\frac{h_{\alpha\beta\gamma'}}{h_\gamma}$ von den $h_\alpha h_\beta$ Elementen $A^{-1}B$ gleich $C$. Daher sind von ihnen $\frac{h_{\alpha\beta\gamma'}}{h_\gamma}g_\gamma$ in der Classe $(\gamma)$ und zugleich in der Gruppe $\Ggr$ enthalten, und folglich gehören von jenen $h_\alpha h_\beta$ Elementen
\[
  \sum_\gamma \frac{h_{\alpha\beta\gamma'}}{h_\gamma}g_\gamma
  =
  g_\alpha g_\beta+\frac{1}{n-1}(h_\alpha-g_\alpha)(h_\beta-g_\beta)
\]
der Gruppe $\Ggr$ an. Diese Gleichung lässt sich mit Hülfe der Beziehung
\[
  \sum_\gamma h_{\alpha\beta\gamma'}=h_\alpha h_\beta
\]
leicht in
\[
\tag{5.}
  h_\alpha h_\beta\chi_\alpha\chi_\beta
  =
  f\sum_\gamma h_{\alpha\beta\gamma'}\chi_\gamma
\]
umformen, wo $\chi_\rho$ durch die Gleichung (1.) definirt ist. Aus diesen Relationen folgt aber in Verbindung mit (2.), dass die Grössen $\chi_\rho$ einen Charakter von $\Hgr$ bilden.

Für die Gruppe des Grades $\frac12 p(p^2-1)$, die ich \emph{Gruppencharaktere} §§ 9, 10 untersucht habe, ergiebt sich aus der Formel (1.) der Charakter $p^{\mathrm{ten}}$ Grades. Denn sie lässt sich als eine zweifach transitive Gruppe des Grades $p+1$ darstellen.

\begin{center}
\small Ausgegeben am 21. Juli.\\[1.5em]
Berlin, gedruckt in der Reichsdruckerei.
\end{center}
% ===== END WORK 057 =====

\clearpage

% ===== BEGIN WORK 058: 058_ber_die_Composition_der_Charaktere_einer =====
\begin{center}
{\Large Über die Composition der Charaktere einer Gruppe}\\[0.75em]
Von G. \textsc{Frobenius}.
\end{center}

Um die Berechnung der Charaktere einer Gruppe zu erleichtern, habe ich in meiner letzten Arbeit (Sitzungsberichte 1898) Relationen abgeleitet, die zwischen den Charakteren einer Gruppe und denen ihrer Untergruppen bestehen. Eine andere Methode, die demselben Zwecke dient, ergiebt sich aus dem Satze, den ich in dieser Arbeit entwickeln will. Danach lässt sich das Product zweier Charaktere einer Gruppe als eine lineare Verbindung ihrer Charaktere darstellen, deren Coefficienten positive ganze Zahlen sind. Diese Coefficienten, die ich mit $f_{\varkappa\lambda\mu}$ bezeichne, haben ähnliche Eigenschaften wie die Zahlen $h_{\alpha\beta\gamma}$, die ich in meiner Arbeit \emph{Über Gruppencharaktere} (Sitzungsberichte 1896) eingeführt habe. Es ist mir zwar nicht gelungen, die Bedeutung der Zahlen $f_{\varkappa\lambda\mu}$ für eine gegebene Gruppe zu erforschen. Aber schon die Gewissheit, dass zwischen den Charakteren einer Gruppe Relationen der angegebenen Art existiren, gestattet in vielen Fällen, aus einem oder mehreren bekannten Charakteren neue abzuleiten.

\section*{§ 1.}

Aus zwei linearen Substitutionen
\begin{equation*}\tag{a}
u_\alpha=\sum_\beta a_{\alpha\beta}v_\beta\qquad(\alpha,\beta=1,2,\cdots f)
\end{equation*}
und
\begin{equation*}\tag{a'}
u'_\gamma=\sum_\delta a'_{\gamma\delta}v'_\delta\qquad(\gamma,\delta=1,2,\cdots f')
\end{equation*}
kann man eine dritte ableiten
\begin{equation*}\tag{A}
u_\alpha u'_\gamma=\sum_{\beta,\delta}a_{\alpha\beta}a'_{\gamma\delta}v_\beta v'_\delta,
\end{equation*}
indem man die $ff'$ Producte $u_\alpha u'_\gamma$ in irgend einer Reihenfolge mit $U_\lambda$ $(\lambda=1,2,\cdots ff')$, die Producte $v_\alpha v'_\gamma$ in derselben Reihenfolge mit $V_\lambda$ bezeichnet. Nennt man nach dem Vorgange von \textsc{Dedekind} die Summe der Diagonalelemente einer Substitution oder Matrix ihre \emph{Spur}, so ist die Spur von (A) gleich dem Producte der Spuren von (a) und (a'). In derselben Weise bilde man aus den Matrizen (b) und (b') der Grade $f$ und $f'$ die Matrix (B) des Grades $ff'$, und aus (c) und (c') die Matrix (C). Ist dann $(c)=(a)(b)$ und $(c')=(a')(b')$, so ist auch $(C)=(A)(B)$, wie aus der oben angegebenen Entstehung von (A) aus (a) und (a') unmittelbar ersichtlich ist.

Seien $\Phi_\nu$ $(\nu=0,1,\cdots k-1)$ die $k$ Primfactoren der Determinante der Gruppe $\Hgr$, und $\chi^{(\nu)}$ ihre Charaktere. Den Elementen $A,B,C,\cdots$ von $\Hgr$ mögen in der primitiven Darstellung von $\Hgr$ durch lineare Substitutionen, die zu $\Phi_\varkappa$ gehört, die Matrizen $(a),(b),(c),\cdots$ entsprechen, in der zu $\Phi_\lambda$ gehörigen Darstellung die Matrizen $(a'),(b'),(c'),\cdots$. Ist dann $AB=C$, so ist $(a)(b)=(c)$ und $(a')(b')=(c')$ und mithin auch $(A)(B)=(C)$. Folglich ist $(A)x_A+(B)x_B+(C)x_C+\cdots$ eine zur Gruppe $\Hgr$ gehörige Matrix, und ihre Determinante ist
\begin{equation*}\tag{1.}
\prod_\mu \Phi_{\mu'}^{\,f_{\varkappa\lambda\mu}},
\end{equation*}
wo $f_{\varkappa\lambda\mu}$ eine positive ganze Zahl oder Null ist, und mit $\Phi_{\mu'}$ die zu $\Phi_\mu$ conjugirte complexe Primfunction bezeichnet werden soll.

Da die Spur von (a) gleich $\chi^{(\varkappa)}(A)$, die von (a') gleich $\chi^{(\lambda)}(A)$ ist, so ist die von (A) gleich $\chi^{(\varkappa)}(A)\chi^{(\lambda)}(A)$. Aus der Formel (1.) erhält man aber für diese Spur den Ausdruck
\[
\chi^{(\varkappa)}(A)\chi^{(\lambda)}(A)
 =\sum_\mu f_{\varkappa\lambda\mu}\chi^{(\mu')}(A)
 =\sum_\mu f_{\varkappa\lambda\mu}\chi^{(\mu)}(A^{-1}),
\]
oder auch, weil $\mu'$ zugleich mit $\mu$ die Werthe $0,1,\cdots k-1$ durchläuft,
\begin{equation*}\tag{2.}
\chi^{(\varkappa)}(R)\chi^{(\lambda)}(R)
=\sum_\mu f_{\varkappa\lambda\mu'}\chi^{(\mu)}(R),
\qquad
\chi_\rho^{(\varkappa)}\chi_\rho^{(\lambda)}
=\sum_\mu f_{\varkappa\lambda\mu'}\chi_\rho^{(\mu)}.
\end{equation*}
Diese Formel enthält die Regeln, nach denen die \emph{Composition} der Charaktere erfolgt. Setzt man
\begin{equation*}\tag{3.}
h f_{\varkappa\varkappa'}=\sum_R \chi^{(\varkappa)}(R)\chi^{(\varkappa')}(R)=h,
\qquad
h f_{\varkappa\lambda}=\sum_R \chi^{(\varkappa)}(R)\chi^{(\lambda)}(R)=0
\end{equation*}
(wo $\lambda$ von $\varkappa'$ verschieden ist), so ergiebt sich mit Hülfe dieser Relationen
\begin{equation*}\tag{4.}
h f_{\varkappa\lambda\mu}=\sum_R \chi^{(\varkappa)}(R)\chi^{(\lambda)}(R)\chi^{(\mu)}(R),
\end{equation*}
oder, wenn $\rho$ die $k$ Klassen conjugirter Elemente durchläuft,
\begin{equation*}\tag{5.}
h f_{\varkappa\lambda\mu}
=\sum_\rho h_\rho\chi_\rho^{(\varkappa)}\chi_\rho^{(\lambda)}\chi_\rho^{(\mu)}.
\end{equation*}
Mithin bleibt $f_{\varkappa\lambda\mu}$ bei allen Vertauschungen der Indices ungeändert, und da $R^{-1}$ zugleich mit $R$ die $h$ Elemente von $\Hgr$ durchläuft, so ist
\begin{equation*}\tag{6.}
f_{\varkappa'\lambda'\mu'}=f_{\varkappa\lambda\mu}.
\end{equation*}

Dass die rechte Seite der Gleichung (4.) eine ganze Zahl ist, kann man leicht direct erkennen. Denn diese Summe ist eine ganze ganzzahlige Function einer primitiven $h^{\ten}$ Einheitswurzel $\vartheta$, und sie bleibt ungeändert, wenn man $\vartheta$ durch irgend eine conjugirte Grösse $\vartheta^n$ ersetzt, wo $n$ zu $h$ theilerfremd ist. Denn um die Substitution von $\vartheta$ durch $\vartheta^n$ auszuführen, braucht man nur in jedem Gliede der Summe $R$ durch $R^n$ zu ersetzen (\emph{Primfactoren}, § 12). Dabei bleibt die Summe ungeändert, weil $R^n$ zugleich mit $R$ die $h$ Elemente von $\Hgr$ durchläuft. Auf diesem Wege erkennt man aber nicht, dass jene ganze Zahl positiv, und auch nicht, dass sie durch $h$ theilbar ist.

Da $\chi^{(\varkappa)}(E)=f_\varkappa=f_{\varkappa'}$ der Grad von $\Phi_\varkappa$ ist, so folgt aus (2.)
\begin{equation*}\tag{7.}
f_\varkappa f_\lambda=\sum_\mu f_{\varkappa\lambda\mu}f_\mu.
\end{equation*}
Mithin ist
\begin{equation*}\tag{8.}
f_{\varkappa\lambda\mu}\leq \frac{f_\varkappa f_\lambda}{f_\mu},
\end{equation*}
wo man für $f_\mu$ die grösste der drei Zahlen $f_\varkappa,f_\lambda,f_\mu$ nehmen kann, und $f_{\varkappa\lambda\mu}\leq f_\varkappa$.

Setzt man
\begin{equation*}\tag{9.}
\sum_\rho f_{\varkappa\lambda\rho'}f_{\rho\mu\nu}=f_{\varkappa\lambda\mu\nu},
\end{equation*}
so ist
\[
h^2 f_{\varkappa\lambda\mu\nu}
=\sum_{R,S}\chi^{(\varkappa)}(R)\chi^{(\lambda)}(R)\chi^{(\mu)}(S)\chi^{(\nu)}(S)
\left(\sum_\rho\chi^{(\rho')}(R)\chi^{(\rho)}(S)\right).
\]
Nun ist aber
\[
\sum_\rho\chi^{(\rho')}(R)\chi^{(\rho)}(S)=0,
\]
ausser wenn $S$ mit $R$ conjugirt ist, dann aber gleich $\frac{h}{h_R}$. Ist $R$ gegeben, so tritt der letzte Fall für $h_R$ verschiedene Werthe von $S$ ein, für die $\chi(S)=\chi(R)$ ist. Folglich ist
\begin{equation*}\tag{10.}
h f_{\varkappa\lambda\mu\nu}
=\sum_R \chi^{(\varkappa)}(R)\chi^{(\lambda)}(R)\chi^{(\mu)}(R)\chi^{(\nu)}(R).
\end{equation*}
Auch diese Zahl bleibt demnach bei jeder Vertauschung der Indices ungeändert, was aus der Formel (9.) nicht ersichtlich ist. In ähnlicher Art ist das Zeichen $f_{\varkappa\lambda\mu\ldots}$ für beliebig viele Indices zu erklären, aber nicht für nur einen Index, wo $f_\varkappa$ den Grad von $\Phi_\varkappa$ bedeutet. Wie schon oben bemerkt, ist $f_{\varkappa\varkappa'}=1$ und $f_{\varkappa\lambda}=0$, falls $\chi_\varkappa$ und $\chi_\lambda$ nicht inverse Charaktere sind. Ferner ist $f_{\varkappa\lambda 0}=f_{\varkappa\lambda}$, $f_{\varkappa\lambda\mu0}=f_{\varkappa\lambda\mu}$.

Ist $\chi_\varkappa$ ein Charakter ersten Grades, also $f_\varkappa=1$, so ist $f_{\varkappa\lambda\mu}$ nach (8.) nicht grösser als $\frac{f_\lambda}{f_\mu}$ oder $\frac{f_\mu}{f_\lambda}$, also gleich Null. Nur wenn $f_\lambda=f_\mu$ ist, kann $f_{\varkappa\lambda\mu}=1$ sein. Nach (7.) ist $f_\lambda=\sum_\mu f_{\varkappa\lambda\mu}f_\mu$, und folglich muss der letzte Fall bei gegebenem $\lambda$ für einen, aber auch nur einen Werth von $\mu$ eintreten. Ist also $f_\varkappa=1$, so ist für jeden Werth von $\lambda$ von den $k$ Zahlen $f_{\varkappa\lambda\mu}$ die eine gleich 1, die anderen $k-1$ sind Null; und ist $f_{\varkappa\lambda\mu}=1$, so ist $f_\lambda=f_\mu$. Das Product aus einem Charakter $f^{\ten}$ Grades $\chi^{(\lambda)}(R)$ und einem Charakter ersten Grades $\chi^{(\varkappa)}(R)$ ist also ein Charakter $f^{\ten}$ Grades $\chi^{(\mu)}(R)=\chi^{(\varkappa)}(R)\chi^{(\lambda)}(R)$, wie ich auf einem anderen Wege schon, \emph{Primfactoren}, § 3, gezeigt habe. Die Formel (2.) ist als eine Verallgemeinerung dieses Satzes anzusehen.

Mit Hülfe der Formel
\begin{equation*}\tag{11.}
\sum_\varkappa \chi^{(\varkappa')}(R)\chi^{(\varkappa)}(R)=\frac{h}{h_R}
\end{equation*}
ergiebt sich leicht
\begin{equation*}\tag{12.}
\sum_\lambda f_{\varkappa\lambda\lambda'}=\sum_\rho \chi_\rho^{(\varkappa)},\qquad
\sum_\mu f_{\varkappa\lambda\mu\mu'}=\sum_\rho \chi_\rho^{(\varkappa)}\chi_\rho^{(\lambda)},\qquad
\sum_\nu f_{\varkappa\lambda\mu\nu\nu'}=\sum_\rho \chi_\rho^{(\varkappa)}\chi_\rho^{(\lambda)}\chi_\rho^{(\mu)}.
\end{equation*}
Daher ist
\[
\sum_{\lambda,\mu}f_{\lambda\lambda'\mu\mu'}
=\sum_\rho\left(\sum_\lambda \chi_\rho^{(\lambda')}\chi_\rho^{(\lambda)}\right)
=\sum_\rho\frac{h}{h_\rho}
\]
und mithin nach (9.) und (6.)
\begin{equation*}\tag{13.}
\sum_{\varkappa,\lambda,\mu} f_{\varkappa\lambda\mu}^{\,2}
=\sum_\rho\frac{h}{h_\rho}.
\end{equation*}
Demnach haben die Zahlen $f_{\varkappa\lambda\mu}$ verhältnissmässig kleine Werthe.

Dass die Summen (12.) rationale ganze Zahlen sind, kann man in derselben Weise, wie es oben für die Summe (5.) gezeigt ist, direct einsehen. Wie die Formeln (12.) ergeben, sind diese Zahlen positiv, und sie lassen sich durch die Zahlen $f_{\varkappa\lambda\mu}$ ausdrücken.

Die hier eingeführten Zahlen $f_{\varkappa\lambda\mu}$ haben mit den früher benutzten Zahlen $h_{\alpha\beta\gamma}$ manche Eigenschaften gemeinsam. Z. B. ist die Determinante $k^{\ten}$ Grades
\begin{equation*}\tag{14.}
\left|\left(\sum_\mu f_{\varkappa\lambda'\mu}x_\mu\right)-f_{\varkappa\lambda'}u\right|
=\prod_\lambda\left(\left(\sum_\varkappa \chi_\lambda^{(\varkappa)}x_\varkappa\right)-u\right)
\qquad(\varkappa,\lambda=0,1,\cdots k-1).
\end{equation*}
Ihre Indices $\varkappa,\lambda,\mu$ beziehen sich auf die $k$ Primfactoren der Gruppendeterminante, während sich die Indices $\alpha,\beta,\gamma$ auf die $k$ Klassen conjugirter Elemente beziehen, worin die $h$ Elemente von $\Hgr$ zerfallen. Auch die hier entwickelten Formeln lassen sich dadurch verallgemeinern, dass man, wie in meiner letzten Arbeit, zu den Charakteren von $\Hgr$ die Charaktere einer Untergruppe hinzunimmt.

\section*{§ 2.}

Aus zwei linearen Substitutionen (a) und (a') von $f$ und $f'$ Variabeln haben wir eine lineare Substitution (A) von $ff'$ Variabeln gebildet. Die Wurzeln ihrer charakteristischen Gleichung sind die $ff'$ Producte, die man erhält, indem man jede der $f$ Wurzeln der charakteristischen Gleichung von (a) mit jeder der $f'$ Wurzeln der charakteristischen Gleichung von (a') multiplicirt. Dies folgt leicht aus der Bemerkung, dass, wenn $(c)^{-1}(a)(c)=(b)$ und $(c')^{-1}(a')(c')=(b')$ ist, auch $(C)^{-1}(A)(C)=(B)$ sein muss.

Bilden die Matrizen $(a),(b),(c)\cdots$ die primitive Darstellung von $\Hgr$, die der Primfunction $\Phi(x)$ entspricht, so erhält man die charakteristische Function von (a), indem man in $\Phi(x+\varepsilon u)$ $x_A=-1$, die anderen Variabelen $x_R=0$ setzt. Ist sie gleich
\[
F_A(u)=(u-\alpha_1)(u-\alpha_2)\cdots(u-\alpha_f),
\]
so sind $\alpha_1,\alpha_2,\cdots,\alpha_f$ $m^{\te}$ Wurzeln der Einheit, wenn $m$ die Ordnung des Elementes $A$ ist (\emph{Primfactoren}, § 12). Ihre Summe ist
\begin{equation*}\tag{1.}
\chi(A)=\alpha_1+\alpha_2+\cdots+\alpha_f.
\end{equation*}
Ist $\Phi=\Phi_\varkappa$, so möge $F_A(u)$ mit $F_A^{(\varkappa)}(u)$ bezeichnet werden oder mit $F_\alpha^{(\varkappa)}(u)$, wenn $A$ zur $\alpha^{\ten}$ Classe gehört. Aus der Formel (2.) § 1 ergiebt sich dann in Verbindung mit der Gleichung (1.) der Satz:
\begin{quote}\itshape
Die Wurzeln der Gleichung
\begin{equation*}\tag{2.}
\prod_\mu\left(F_\alpha^{(\mu)}(u)\right)^{f_{\varkappa\lambda\mu'}}
\end{equation*}
vom Grade $f_\varkappa f_\lambda$ werden erhalten, indem man jede der $f_\varkappa$ Wurzeln der Gleichung $F_\alpha^{(\varkappa)}(u)=0$ mit jeder der $f_\lambda$ Wurzeln der Gleichung $F_\alpha^{(\lambda)}(u)=0$ multiplicirt.
\end{quote}
Entwickelt man die logarithmische Ableitung des Ausdrucks (2.) nach absteigenden Potenzen von $u$, so ist der Coefficient von $u^{-n-1}$
\[
\sum_\mu f_{\varkappa\lambda\mu'}\chi^{(\mu)}(A^n)
=\chi^{(\varkappa)}(A^n)\chi^{(\lambda)}(A^n).
\]
Daraus ist ersichtlich, dass die in jenem Satze ausgesprochene Beziehung um nichts allgemeiner ist als die Formel (2.), § 1.

In der Function
\begin{equation*}\tag{3.}
(-1)^f F_\alpha^{(\varkappa)}(-u)
=(u+\alpha_1)(u+\alpha_2)\cdots(u+\alpha_f)
\end{equation*}
ist der Coefficient von $u^{f-1}$ gleich $\chi_\alpha^{(\varkappa)}$. Das constante Glied ist ein Charakter ersten Grades (\emph{Primfactoren}, § 12). Allgemeiner ist der Coefficient von $u^{f-n}$, für den ich die \emph{Primfactoren}, § 4 (8.) eingeführte Bezeichnung wähle, gleich einem Ausdruck von der Form
\begin{equation*}\tag{4.}
\FrobP_n(A)=\frac{1}{n!}\chi(A,A,\cdots A)=\sum_\lambda s_{n\lambda}\chi^{(\lambda)}(A),
\end{equation*}
wo die Grössen $s_{n\lambda}$ positive ganze Zahlen sind, die von $\alpha$ unabhängig sind.

Um dies zu beweisen, bilde man aus der Substitution (a) $n$ verschiedene Substitutionen
\begin{equation*}\tag{5.}
u_\alpha^{(\nu)}=\sum_\beta a_{\alpha\beta}v_\beta^{(\nu)}
\qquad(\alpha,\beta=1,2,\cdots f),
\end{equation*}
wo
\[
u_1^{(\nu)},u_2^{(\nu)},\cdots,u_f^{(\nu)}
\qquad(\nu=1,2,\cdots n)
\]
$nf$ unabhängige Variabele sind. Die $\binom{f}{n}=g$ Determinanten $n^{\ten}$ Grades, die sich aus ihnen bilden lassen, bezeichne man in irgend einer Reihenfolge mit $U_1,U_2,\cdots,U_g$, die aus den Variabeln
\[
v_1^{(\nu)},v_2^{(\nu)},\cdots,v_f^{(\nu)}
\qquad(\nu=1,2,\cdots n)
\]
analog gebildeten Determinanten mit $V_1,V_2,\cdots,V_g$. Nach den Gleichungen (5.) sind dann $U_1,U_2,\cdots,U_g$ lineare Functionen von $V_1,V_2,\cdots,V_g$, deren Coefficienten die $g^2$ Unterdeterminanten $n^{\ten}$ Grades der Matrix $a_{\alpha\beta}$ sind. Die so erhaltene lineare Substitution möge jetzt mit (A) bezeichnet werden. Die Wurzeln ihrer charakteristischen Gleichung sind die Producte von je $n$ der Wurzeln $\alpha_1,\alpha_2,\cdots,\alpha_f$, ihre Spur ist die Summe der Hauptunterdeterminanten $n^{\ten}$ Grades der Matrix (a), also gleich dem Ausdruck $\FrobP_n(A)$. Entsprechen in derselben Weise den Matrizen $f^{\ten}$ Grades $(b),(c),\cdots$ die Matrizen $g^{\ten}$ Grades $(B),(C),\cdots$, so ist, falls $(c)=(a)(b)$ ist, auch $(C)=(A)(B)$. Folglich ist $(A)x_A+(B)x_B+(C)x_C+\cdots$ eine zu $\Hgr$ gehörige Matrix. Diese ist auch von \textsc{Moliên}, \emph{Ueber die Invarianten der linearen Substitutionsgruppen}, § 5 (Sitzungsberichte 1897) betrachtet worden. Ihre Determinante ist ein Ausdruck von der Form $\prod_\lambda\Phi_\lambda^{s_{\varkappa\lambda}^{n}}$, woraus sich die Relation (4.) ergiebt. Ferner sind die Wurzeln der Gleichung
\[
\prod_\lambda\left(F_\alpha^{(\lambda)}(u)\right)^{s_{\varkappa\lambda}^{n}}=0
\]
vom Grade $\binom{f_\varkappa}{n}$ die Producte von je $n$ der $f_\varkappa$ Wurzeln $\alpha_1,\alpha_2,\cdots$ der Gleichung $F_\alpha^{(\varkappa)}(u)=0$.

Die Ausdrücke (4.) haben die Form
\begin{equation*}\tag{6.}
2\FrobP_2(R)=\chi(R)^2-\chi(R^2),\qquad
6\FrobP_3(R)=\chi(R)^3-3\chi(R)\chi(R^2)+2\chi(R^3),
\end{equation*}
allgemein
\begin{equation*}\tag{7.}
\FrobP_n(R)=\sum (-1)^{\beta+\delta+\cdots}
\frac{\chi(R)^\alpha}{1^\alpha\alpha!}
\frac{\chi(R^2)^\beta}{2^\beta\beta!}
\frac{\chi(R^3)^\gamma}{3^\gamma\gamma!}
\frac{\chi(R^4)^\delta}{4^\delta\delta!}\cdots,
\end{equation*}
wo die Summe über alle positiven Lösungen der Gleichung
\begin{equation*}\tag{8.}
\alpha+2\beta+3\gamma+4\delta+\cdots=n
\end{equation*}
zu erstrecken ist.

Wie \textsc{Moliên} a. a. O. § 2 gezeigt hat, gilt eine der Formel (4.) analoge Formel für den Coefficienten $\zeta_n(R)$ von $u^{-f-n}$ in der Entwickelung der Function $\frac{1}{F_\alpha^{(\varkappa)}(u)}$. Dieser ist gleich der Summe
\begin{equation*}\tag{9.}
\zeta_n(R)=\sum
\frac{\chi(R)^\alpha}{1^\alpha\alpha!}
\frac{\chi(R^2)^\beta}{2^\beta\beta!}
\frac{\chi(R^3)^\gamma}{3^\gamma\gamma!}
\frac{\chi(R^4)^\delta}{4^\delta\delta!}\cdots,
\end{equation*}
erstreckt über die positiven Lösungen der Gleichung (8).

\section*{§ 3.}

Zum Schluss theile ich noch einige Beispiele von Gruppen mit, deren Charaktere ich mit Hülfe der Methoden berechnet habe, die ich hier und in meiner letzten Arbeit entwickelt habe.

Die Elemente der symmetrischen Gruppe des Grades $n=6$ sind die 720 Permutationen von 6 Symbolen $a,b,c,d,e,f$, die der alternirenden Gruppe die 360 geraden Permutationen. In der ersten Spalte der beifolgenden Tabelle ist jede Classe conjugirter Elemente durch eine in ihre Cyklen zerlegte Permutation repräsentirt. Durch Composition mit dem Charakter ersten Grades $\chi^{(1)}$ entstehen aus den Charakteren $\chi^{(2)},\chi^{(4)},\chi^{(6)},\chi^{(8)}$ die Charaktere $\chi^{(3)},\chi^{(5)},\chi^{(7)},\chi^{(9)}$, während $\chi^{(10)}$ ungeändert bleibt. Besonders bemerkenswerth ist der Charakter $\chi=\chi^{(2)}$. Wenn $\lambda$ Symbole von der Permutation $R$ nicht versetzt werden (wenn $R$ $\lambda$ Cyklen ersten Grades enthält), so ist $\chi(R)=\lambda-1$.

Nach den Entwickelungen meiner letzten Arbeit (§ 5) kann man die Regel, nach der $\chi_2$ zu berechnen ist, auch so ausdrücken: Die Permutationen von $\Hgr$, die ein bestimmtes Symbol ungeändert lassen, bilden eine Gruppe $\Ggr$ des Grades $g=\frac{h}{n}$. Enthält die $\rho^{\te}$ Classe $h_\rho$ Elemente, und sind davon $g_\rho$ in $\Ggr$ enthalten, so ist
\begin{equation*}\tag{1.}
\chi_\rho=\frac{h g_\rho}{g h_\rho}-1.
\end{equation*}
Aus derselben Formel findet man den Charakter $\chi^{(6)}$, indem man für $\Ggr$ die in $\Hgr$ enthaltene Gruppe der Ordnung $g=72$ nimmt. Diese imprimitive Gruppe erhält man, indem man die $n=6$ Symbole in $s=2$ Systeme von je $r=3$ Symbolen theilt, $abc,def$, und dann die $s$ Systeme und die $r$ Symbole jedes Systems auf alle möglichen Arten permutirt. Die Ordnung dieser Gruppe ist $g=(r!)^s s!=(3!)^2 2!=72$.

Endlich kann man mittelst der Formel (1.) auch den Charakter $\chi^{(4)}$ berechnen, indem man für $\Ggr$ die dreifach transitive Gruppe der Ordnung 120 nimmt, die der symmetrischen Gruppe des Grades 5 isomorph ist. Statt dessen kann man $\chi^{(4)}$ aus $\chi^{(2)}$ durch den bekannten Isomorphismus von $\Hgr$ in sich ableiten, wodurch sich die Classen 1 und 7, 4 und 10, 5 und 9 unter einander vertauschen, während die übrigen ungeändert bleiben. Durch diesen Isomorphismus geht auch $\chi^{(8)}$ in $\chi^{(9)}=\chi^{(8)}\chi^{(1)}$ über.

\begin{center}
\emph{Symmetrische Gruppe des Grades 6.}\\[0.4em]
$h=720$.\\[0.4em]
{\scriptsize
\resizeMath{
\begin{array}{ll|rrrrrrrrrrr|r}
 & & \chi^{(0)}&\chi^{(1)}&\chi^{(2)}&\chi^{(3)}&\chi^{(4)}&\chi^{(5)}&\chi^{(6)}&\chi^{(7)}&\chi^{(8)}&\chi^{(9)}&\chi^{(10)}&h_\alpha\\ \hline
 &\chi_0&1&1&5&5&5&5&9&9&10&10&16&1\\
(ab)&\chi_1&1&-1&3&-3&-1&1&3&-3&2&-2&0&15\\
(ab)(cd)&\chi_2&1&1&1&1&1&1&1&1&-2&-2&0&45\\
(abcd)&\chi_3&1&-1&1&-1&1&-1&-1&1&0&0&0&90\\
(abc)&\chi_4&1&1&2&2&-1&-1&0&0&1&1&-2&40\\
(abc)(de)&\chi_5&1&-1&0&0&-1&1&0&0&-1&1&0&120\\
(abcde)&\chi_6&1&1&0&0&0&0&-1&-1&0&0&1&144\\
(ab)(cd)(ef)&\chi_7&1&-1&-1&1&3&-3&3&-3&-2&2&0&15\\
(abcd)(ef)&\chi_8&1&1&-1&-1&-1&-1&1&1&0&0&0&90\\
(abcdef)&\chi_9&1&-1&-1&1&0&0&0&0&1&-1&0&120\\
(abc)(def)&\chi_{10}&1&1&-1&-1&2&2&0&0&1&1&-2&40
\end{array}}}
\end{center}

\begin{center}
\emph{Alternirende Gruppe des Grades 6.}\\[0.4em]
$h=360$.\\[0.4em]
{\scriptsize
\resizeMath{
\begin{array}{ll|rrrrrrr|r}
 & & \chi^{(0)}&\chi^{(2)}&\chi^{(4)}&\chi^{(6)}&\chi^{(8)}&\chi^{(10)}&\chi^{(\overline{10})}&h_\alpha\\ \hline
 &\chi_0&1&5&5&9&10&8&8&1\\
(ab)(cd)&\chi_2&1&1&1&1&-2&0&0&45\\
(abc)&\chi_4&1&2&-1&0&1&-1&-1&40\\
(abcde)&\chi_6&1&0&0&-1&0&\half(1+\sqrt5)&\half(1-\sqrt5)&72\\
(acebd)&\chi_{\overline6}&1&0&0&-1&0&\half(1-\sqrt5)&\half(1+\sqrt5)&72\\
(abcd)(ef)&\chi_8&1&-1&-1&1&0&0&0&90\\
(abc)(def)&\chi_{10}&1&-1&2&0&1&-1&-1&40
\end{array}}}
\end{center}

\section*{§ 4.}

Die binären Gruppen des Tetraeders, Oktaeders und Ikosaeders haben die Ordnungen $h=24,48$ und 120. Die Classen sind meist schon durch die Ordnungen ihrer Elemente bestimmt, die sich in der ersten Spalte jeder Tabelle finden. Nur wo die Ordnung zur Definition der Classe nicht ausreicht, findet sich in jener Spalte ein Element der Classe, und zwar bedeutet in der Gruppe des Tetraeders (Oktaeders, Ikosaeders) das Zeichen $R$ ($S,T$) ein beliebiges Element der Ordnung 6 (8, 10). Die Classe (4) der Oktaedergruppe enthält jedes Element der Ordnung 4, das sich nicht als Quadrat eines Elementes der Ordnung 8 darstellen lässt; die übrigen Elemente der Ordnung 4 bilden die durch das Zeichen $S^2$ charakterisirte Classe (5).

Jede der drei Gruppen enthält eine invariante Untergruppe $\Fgr=E+F$ der Ordnung 2. Das Element $F$ der Ordnung 2 ist mit jedem Elemente der Gruppe vertauschbar und bildet für sich allein die Classe (1). Die Gruppe $\frac{\Hgr}{\Fgr}$ ist die Gruppe der Ordnung 12 (24, 60), deren Charaktere ich \emph{Gruppencharaktere}, § 8 berechnet habe. Aus diesen Tabellen kann man unmittelbar die Werthe der Charaktere entnehmen, wofür $\chi_0=\chi_1$ ist, und die zur Gruppe $\frac{\Hgr}{\Fgr}$ gehören. Für diese Charaktere ist allgemein $\chi(RF)=\chi(R)$, für alle anderen aber $\chi(RF)=-\chi(R)$.

Die interessanteste Eigenschaft dieser Gruppen besteht darin, dass sie eigentliche Charaktere zweiten Grades besitzen. (In der Tabelle für die Oktaedergruppe, \emph{Gruppencharaktere}, § 8, gehört der Charakter $\chi^{(3)}$ des Grades $f=2$ zur Gruppe $\frac{\Hgr}{\Ggr}$, wo $\Ggr$ die invariante Untergruppe der Ordnung 4 von $\Hgr$ ist). Ihre Werthe ergeben sich sofort aus den bekannten Darstellungen dieser Gruppen durch binäre lineare Substitutionen mit Hülfe der Formel \emph{Darstellung}, § 4 (5). Für die Oktaedergruppe giebt es genau zwei nicht aequivalente Darstellungen dieser Art, die sich aber nur durch das Vorzeichen von $\sqrt2$ unterscheiden (denn der Charakter $\chi^{(4)}$ gehört zur Gruppe $\frac{\Hgr}{\Fgr}$).

Dasselbe gilt von der Ikosaedergruppe. Für die Tetraedergruppe aber giebt es drei nicht aequivalente Darstellungen. Ihre Charaktere $\chi^{(4)},\chi^{(5)}$ und $\chi^{(6)}$ ergeben sich aus einem von ihnen durch Composition mit den drei Charakteren ersten Grades. Auf diese Weise findet man die Charaktere dieser drei Gruppen alle ohne Benutzung der Zahlen $h_{\alpha\beta\gamma}$. Denn die noch fehlenden Charaktere $\chi^{(7)}$ der Oktaedergruppe und $\chi^{(7)}$ und $\chi^{(8)}$ der Ikosaedergruppe ergeben sich aus den zwischen den Charakteren bestehenden quadratischen Relationen.

\clearpage

\begin{center}
\emph{Tetraeder.}\\[0.25em]
$h=24$.\\[0.25em]
{\scriptsize
\resizeMath{
\begin{array}{ll|rrrrrrr|r}
 & & \chi^{(0)}&\chi^{(1)}&\chi^{(2)}&\chi^{(3)}&\chi^{(4)}&\chi^{(5)}&\chi^{(6)}&h_\alpha\\ \hline
1&\chi_0&1&1&1&3&2&2&2&1\\
2&\chi_1&1&1&1&3&-2&-2&-2&1\\
4&\chi_2&1&1&1&-1&0&0&0&6\\
R^2&\chi_3&1&\rho&\rho^2&0&-1&-\rho&-\rho^2&4\\
R^4&\chi_4&1&\rho^2&\rho&0&-1&-\rho^2&-\rho&4\\
R^5&\chi_5&1&\rho&\rho^2&0&1&\rho&\rho^2&4\\
R&\chi_6&1&\rho^2&\rho&0&1&\rho^2&\rho&4
\end{array}}}
\end{center}

\begin{center}
\emph{Oktaeder.}\\[0.25em]
$h=48$.\\[0.25em]
{\scriptsize
\resizeMath{
\begin{array}{ll|rrrrrrrr|r}
 & & \chi^{(0)}&\chi^{(1)}&\chi^{(2)}&\chi^{(3)}&\chi^{(4)}&\chi^{(5)}&\chi^{(6)}&\chi^{(7)}&h_\alpha\\ \hline
1&\chi_0&1&1&3&3&2&2&2&4&1\\
2&\chi_1&1&1&3&3&2&-2&-2&-4&1\\
3&\chi_2&1&1&0&0&-1&-1&-1&1&8\\
6&\chi_3&1&1&0&0&-1&1&1&-1&8\\
4&\chi_4&1&-1&1&-1&0&0&0&0&12\\
S^2&\chi_5&1&1&-1&-1&2&0&0&0&6\\
S^3&\chi_6&1&-1&-1&1&0&\sqrt2&-\sqrt2&0&6\\
S&\chi_7&1&-1&-1&1&0&-\sqrt2&\sqrt2&0&6
\end{array}}}
\end{center}

\begin{center}
\emph{Ikosaeder.}\\[0.25em]
$h=120$.\\[0.25em]
{\scriptsize
\resizeMath{
\begin{array}{ll|rrrrrrrrr|r}
 & & \chi^{(0)}&\chi^{(1)}&\chi^{(2)}&\chi^{(3)}&\chi^{(4)}&\chi^{(5)}&\chi^{(6)}&\chi^{(7)}&\chi^{(8)}&h_\alpha\\ \hline
1&\chi_0&1&5&4&3&3&2&2&4&6&1\\
2&\chi_1&1&5&4&3&3&-2&-2&-4&-6&1\\
4&\chi_2&1&1&0&-1&-1&0&0&0&0&30\\
3&\chi_3&1&-1&1&0&0&-1&-1&1&0&20\\
6&\chi_4&1&-1&1&0&0&1&1&-1&0&20\\
T^4&\chi_5&1&0&-1&\half(1+\sqrt5)&\half(1-\sqrt5)&\half(-1+\sqrt5)&\half(-1-\sqrt5)&-1&1&12\\
T^2&\chi_6&1&0&-1&\half(1-\sqrt5)&\half(1+\sqrt5)&\half(-1-\sqrt5)&\half(-1+\sqrt5)&-1&1&12\\
T^3&\chi_7&1&0&-1&\half(1-\sqrt5)&\half(1+\sqrt5)&\half(1+\sqrt5)&\half(1-\sqrt5)&1&-1&12\\
T&\chi_8&1&0&-1&\half(1+\sqrt5)&\half(1-\sqrt5)&\half(1-\sqrt5)&\half(1+\sqrt5)&1&-1&12
\end{array}}}
\end{center}
% ===== END WORK 058 =====

\clearpage

% ===== BEGIN WORK 059: 059_ber_die_Darstellung_der_endlichen_Gruppen =====
\frobtitle{Über die Darstellung der endlichen Gruppen\\durch lineare Substitutionen. II.}

Den $h$ verschiedenen Elementen $A,B,C,\cdots$ einer endlichen Gruppe $\Hgr$ seien $h$ homogene lineare Substitutionen $a,b,c,\cdots$ von $n$ Variabelen so zugeordnet, dass immer, wenn $AB=C$ ist, auch $ab=c$ ist. Dann bilden diese Substitutionen eine Gruppe, die mit $\Hgr$ holoedrisch oder meroedrisch isomorph ist, oder, anders ausgedrückt, die Substitutionen bilden eine \emph{Darstellung} der Gruppe $\Hgr$. Unter den Zeichen $a,b,c,\cdots$ kann man auch die Matrizen der Substitutionen verstehen, unter $ab$ die aus $a$ und $b$ zusammengesetzte Matrix.

Aus der gegebenen Darstellung kann man eine neue ableiten, indem man in den Substitutionen andere Variabele einführt. Dies Verfahren kommt darauf hinaus, dass man die Matrizen $a,b,c,\cdots$ durch $p^{-1}ap,p^{-1}bp,p^{-1}cp,\cdots$ ersetzt, wo $p$ eine beliebige Matrix von nicht verschwindender Determinante ist. Zwei solche Darstellungen habe ich in dem ersten Theile dieser Arbeit (Sitzungsberichte 1897, im Folgenden mit $D.$ citirt, § 2) als \emph{aequivalent} bezeichnet.

Kennt man für dieselbe Gruppe $\Hgr$ eine zweite Darstellung durch die Matrizen $a',b',c',\cdots$ des Grades $n'$, so bilden die Matrizen
\[
\begin{pmatrix}a&0\\0&a'\end{pmatrix},\quad
\begin{pmatrix}b&0\\0&b'\end{pmatrix},\quad
\begin{pmatrix}c&0\\0&c'\end{pmatrix},\cdots
\]
eine neue Darstellung des Grades $n+n'$. Dabei ist auch der Fall nicht ausgeschlossen, dass die zweite Darstellung mit der ersten identisch oder aequivalent ist. In der nämlichen Weise kann man aus mehreren bekannten, gleichen oder verschiedenen Darstellungen eine neue ableiten. Diese Gruppe von Substitutionen kann unter Umständen der Gruppe $\Hgr$ holoedrisch isomorph sein, trotzdem die gegebenen Darstellungen ihr nur meroedrisch isomorph waren; und dies ist der Grund, weshalb bei der Untersuchung aller Darstellungen einer gegebenen Gruppe $\Hgr$ auch die nicht ausgeschlossen werden dürfen, die in Wirklichkeit eine mit $\Hgr$ meroedrisch isomorphe Gruppe darstellen.

Jede Darstellung, die in der oben erörterten Art aus mehreren erhalten wird, nenne ich eine \emph{zerfallende} oder \emph{zerlegbare}, und jede Darstellung, die einer zerfallenden aequivalent ist, eine \emph{imprimitive} oder \emph{reducibele}. Ist eine Darstellung aber keiner zerlegbaren aequivalent, so nenne ich sie eine \emph{primitive} oder \emph{irreducibele} (vergl. die frühere vorläufige Definition $D.$ § 5).

Betrachtet man aequivalente Darstellungen nicht als verschieden, so giebt es nur eine endliche Anzahl verschiedener primitiver Darstellungen einer Gruppe $\Hgr$ durch lineare Substitutionen oder ihre Matrizen. Diese Zahl $k$ ist gleich der Anzahl der Classen conjugirter Elemente, worin die Elemente von $\Hgr$ zerfallen. Jede imprimitive Darstellung ist einer anderen aequivalent, die in lauter primitive Darstellungen zerfällt, wobei aber jede einzelne der $k$ primitiven Darstellungen mehrfach auftreten kann. Und zwar ist eine solche Zerlegung nur in einer Art möglich.

Seien $x_A,x_B,x_C,\cdots$ $h$ unabhängige Variabele, und sei
\[
\Theta=|x_{PQ^{-1}}|=\Phi^f{\Phi'}^{f'}{\Phi''}^{f''}\cdots
\]
die Determinante der Gruppe $\Hgr$, und $\Phi,\Phi',\Phi'',\cdots$ ihre verschiedenen Primfactoren. Bilden die Matrizen $n^{\ten}$ Grades $a,b,c,\cdots$ eine Darstellung von $\Hgr$, so nenne ich ($D.$ § 2) die Matrix $n^{\ten}$ Grades $ax_A+bx_B+cx_C+\cdots$ die der Darstellung von $\Hgr$ entsprechende Matrix oder eine zur Gruppe $\Hgr$ gehörige Matrix. Ihre $n^2$ Elemente sind lineare Functionen der $h$ Variabelen $x_A,x_B,x_C,\cdots$, ihre Determinante
\[
|ax_A+bx_B+cx_C+\cdots|=\Phi^s{\Phi'}^{s'}{\Phi''}^{s''}\cdots
\]
ist ein Product von Primfactoren der Gruppendeterminante. Umgekehrt entspricht jedem solchen Producte eine und nur eine Darstellung von $\Hgr$, d. h. zwei Darstellungen, deren entsprechende Matrizen gleiche Determinanten haben, sind aequivalent. Einem jeden der $k$ Primfactoren $\Phi$ der Gruppendeterminante $\Theta$ entspricht eine der $k$ primitiven Darstellungen, die ich mit $[\Phi]$ bezeichnen will. Einem Producte $\Phi^s{\Phi'}^{s'}\cdots$ entspricht eine Darstellung, die in $s$ Darstellungen $[\Phi]$, $s'$ Darstellungen $[\Phi']$, $\cdots$ zerfällt. Durch die Untersuchung der Determinante der Matrix, die einer gegebenen Darstellung entspricht, kann man daher erkennen, ob die Darstellung eine primitive ist oder nicht, und im letzteren Falle, in welche primitive Darstellungen sie zerlegt werden kann.

Die hier entwickelten durch ihre Einfachheit ausgezeichneten Resultate bilden den Abschluss meiner allgemeinen Untersuchungen über die Gruppendeterminante. Auf einem anderen Wege hat sie \textsc{Molien} in der $D.$ § 4 citirten Arbeit erhalten.

Zum Schluss gehe ich auf die Herstellung der primitiven Darstellungen näher ein. Um eine solche zu erhalten, braucht man nur eine Lösung eines gewissen Systems von linearen und quadratischen Gleichungen zu berechnen. Jede solche Lösung nenne ich daher ein die Darstellung \emph{determinirendes Werthsystem}.

\section*{§ 1.}

In der Matrix $f^{\ten}$ Grades $u$ seien die Elemente
\[
u_{\alpha\beta}\qquad(\alpha,\beta=1,2,\cdots f)
\]
$f^2$ von einander unabhängige Variabele. Seien $v_{\alpha\beta}$ $f^2$ andere Variabele, und sei
\[
w_{\alpha\beta}=u_{\alpha1}v_{1\beta}+u_{\alpha2}v_{2\beta}+\cdots+u_{\alpha f}v_{f\beta}.
\]
Geht dann die Matrix $u$ in $v$ oder $w$ über, falls man $u_{\alpha\beta}$ durch $v_{\alpha\beta}$ oder $w_{\alpha\beta}$ ersetzt, so ist $w=uv$.

Die $n^2$ Elemente $x_{\kappa\lambda}$ der Matrix $n^{\ten}$ Grades $X$ seien lineare Functionen der $f^2$ Variabelen $u_{\alpha\beta}$. Sie gehe in $Y$ oder $Z$ über, wenn man $u_{\alpha\beta}$ durch $v_{\alpha\beta}$ oder $w_{\alpha\beta}$ ersetzt. Wir wollen untersuchen, wie jene linearen Functionen beschaffen sein müssen, damit $Z=XY$ sei, wenn $w=uv$ ist. Ich beschränke mich dabei auf den Fall, wo die Determinante $|X|$ von Null verschieden ist.

Hat $X$ jene Eigenschaft, so hat sie auch $PXP^{-1}$, wo $P$ eine Matrix von $n^2$ constanten Elementen und $|P|$ von Null verschieden ist. Ist ferner $0$ die Matrix $f^{\ten}$ Grades, deren Elemente sämmtlich verschwinden, und ist $n=fg$ ein Vielfaches von $f$, so hat
\[
U=\begin{pmatrix}
u&0&0&\cdots\\
0&u&0&\cdots\\
0&0&u&\cdots\\
\cdot&\cdot&\cdot&\cdots
\end{pmatrix}
\]
die verlangte Eigenschaft, und ebenso $P^{-1}UP=X$. Ich will nun umgekehrt zeigen: Ist $X$ eine beliebige Matrix der betrachteten Art, so muss $n=fg$ ein Vielfaches von $f$ sein, und man kann eine constante Matrix $P$ so bestimmen, dass $PXP^{-1}=U$ wird.

Ist $B$ die Hauptmatrix des Grades $g$, so kann man die Matrix $U$ des Grades $fg$ durch eine gewisse Umstellung der Zeilen und die gleiche Umstellung der Spalten auf die Form
\[
V=\begin{pmatrix}
u_{11}B&u_{12}B&\cdots&u_{1f}B\\
u_{21}B&u_{22}B&\cdots&u_{2f}B\\
\cdot&\cdot&\cdots&\cdot\\
u_{f1}B&u_{f2}B&\cdots&u_{ff}B
\end{pmatrix}
\]
bringen. Die Umstellung der Spalten erfolgt durch Composition von $U$ mit einer Matrix $Q$, bei der in jeder Zeile und in jeder Spalte ein Element gleich $1$, die anderen gleich $0$ sind. Führt man dann in $UQ$ die nämliche Umstellung der Zeilen aus, so erhält man $Q'UQ$, wo $Q'$ die zu $Q$ conjugirte Matrix ist. Da aber $Q$ eine orthogonale Matrix ist, so ist $Q'=Q^{-1}$, also $V=Q^{-1}UQ$.

Da die Elemente der Matrix $X$ lineare Functionen von $u_{11},u_{12},\cdots,u_{ff}$ sind, so kann man $X=\sum A_{\alpha\beta}u_{\alpha\beta}$ setzen, wo $A_{\alpha\beta}$ eine constante Matrix $n^{\ten}$ Grades ist. Nun soll $XY=Z$ sein, also
\[
\Bigl(\sum A_{\alpha\beta}u_{\alpha\beta}\Bigr)
\Bigl(\sum A_{\gamma\delta}v_{\gamma\delta}\Bigr)
 =\sum A_{\alpha\gamma}w_{\alpha\gamma}
 =\sum A_{\alpha\gamma}u_{\alpha\beta}v_{\beta\gamma}.
\]
Daher ist $A_{\alpha\beta}A_{\gamma\delta}=0$, wenn $\beta$ von $\gamma$ verschieden ist, dagegen $A_{\alpha\beta}A_{\beta\gamma}=A_{\alpha\gamma}$. Da $A_{11}^2=A_{11}$ ist, so ist
\[
|A_{11}+rE|=r^{n-g_1}(1+r)^{g_1}.
\]
Wäre $g_1=0$, so wäre $A_{11}=0$, also auch $A_{\alpha\beta}=A_{\alpha1}A_{11}A_{1\beta}=0$. Daher ist $g_1>0$, und man kann $P$ so bestimmen, dass $P^{-1}A_{11}P$ eine Matrix wird, worin $c_{11}=\cdots=c_{g_1,g_1}=1$, alle anderen Elemente gleich Null sind. Wir denken uns $X$ durch $P^{-1}XP$ ersetzt, nehmen also an, dass $A_{11}$ selbst jene Matrix ist. Nun ist $A_{11}A_{22}=A_{22}A_{11}=0$, und mithin sind in $A_{22}$ die Elemente der ersten $g_1$ Zeilen und Spalten sämmtlich Null, so dass man $A_{22}$ auch als eine Matrix des Grades $n-g_1$ betrachten und als solche transformiren kann. Folglich lässt sich eine Matrix $n^{\ten}$ Grades $P$, worin die Elemente der ersten $g_1$ Zeilen und Spalten mit den entsprechenden von $A_{11}$ übereinstimmen, so bestimmen, dass in $P^{-1}A_{22}P$ $c_{g_1+1,g_1+1}=\cdots=c_{g_1+g_2,g_1+g_2}=1$ wird, alle anderen Elemente aber verschwinden. Dann sind in $A_{12}=A_{11}A_{12}A_{22}$ nur die Elemente nicht nothwendig Null, welche die Zeilen $1,2,\cdots,g_1$ mit den Spalten $g_1+1,g_1+2,\cdots,g_1+g_2$ gemeinsam haben. Nach Ausführung dieser Transformationen ist
\[
A_{11}=\begin{pmatrix}
B_{11}&N_{12}&N_{13}&\cdots\\
N_{21}&N_{22}&N_{23}&\cdots\\
N_{31}&N_{32}&N_{33}&\cdots\\
\cdot&\cdot&\cdot&\cdots
\end{pmatrix},\quad
A_{12}=\begin{pmatrix}
N_{11}&B_{12}&N_{13}&\cdots\\
N_{21}&N_{22}&N_{23}&\cdots\\
N_{31}&N_{32}&N_{33}&\cdots\\
\cdot&\cdot&\cdot&\cdots
\end{pmatrix},\quad\cdots,
\]
wo $B_{\alpha\beta}$ eine Matrix von $g_\alpha$ Zeilen und $g_\beta$ Spalten ist, und $N_{\alpha\beta}$ eine Matrix derselben Art mit lauter verschwindenden Elementen. $B_{\alpha\alpha}$ ist die Hauptmatrix des Grades $g_\alpha$. Aus $A_{\alpha\beta}A_{\beta\gamma}=A_{\alpha\gamma}$ folgt $B_{\alpha\beta}B_{\beta\gamma}=B_{\alpha\gamma}$. So ist $B_{12}B_{21}=B_{11}$. Wäre daher $g_2<g_1$, so müsste die Determinante von $B_{11}$ verschwinden. Ebenso folgt aus $B_{21}B_{12}=B_{22}$, dass $g_2\geq g_1$ ist. Mithin ist $g_1=g_2=g_3=\cdots=g$, und $n=fg$. Denn wäre $n>fg$, so beständen die letzten Zeilen und Spalten von $X$ aus lauter Nullen. Ist $B$ die Hauptmatrix des Grades $g$, so ist $B_{\alpha\alpha}=B$ und $B_{\alpha\beta}B_{\beta\alpha}=B$. Folglich sind $B_{\alpha\beta}$ und $B_{\beta\alpha}$ reciproke Matrizen des Grades $g$, und ihre Determinanten sind von Null verschieden. Sei $N$ eine Matrix von $g^2$ verschwindenden Elementen. Dann ist die Determinante der Matrix $n^{\ten}$ Grades
\[
L=\begin{pmatrix}
B_{11}&N&N&\cdots\\
N&B_{12}&N&\cdots\\
N&N&B_{13}&\cdots\\
\cdot&\cdot&\cdot&\cdots
\end{pmatrix}
\]
von Null verschieden, und es ist
\[
L^{-1}=\begin{pmatrix}
B_{11}&N&N&\cdots\\
N&B_{21}&N&\cdots\\
N&N&B_{31}&\cdots\\
\cdot&\cdot&\cdot&\cdots
\end{pmatrix}.
\]
Die Matrix $LA_{\alpha\beta}L^{-1}$ unterscheidet sich von $A_{\alpha\beta}$ nur dadurch, dass an Stelle von $B_{\alpha\beta}$ tritt $B_{1\alpha}B_{\alpha\beta}B_{\beta1}=B_{11}=B$. Daher ist $LXL^{-1}=V$, und damit ist die aufgestellte Behauptung bewiesen.

\section*{§ 2.}

Ist $(x_{PQ^{-1}})$ die Matrix der Gruppe $h^{\ten}$ Ordnung $\Hgr$, und ist
\begin{equation}\tag{1.}
z_R=\sum x_Py_Q\qquad(PQ=R),
\end{equation}
so ist $(z_{PQ^{-1}})=(x_{PQ^{-1}})(y_{PQ^{-1}})$. Sei
\[
X=(x_{\kappa\lambda})\qquad(\kappa,\lambda=1,2,\cdots,n)
\]
eine Matrix $n^{\ten}$ Grades, deren Determinante nicht verschwindet, und deren $n^2$ Elemente $x_{\kappa\lambda}$ lineare Functionen der $h$ Variabelen $x_R$ sind. Sie gehe in $Y$ oder $Z$ über, falls man $x_R$ durch $y_R$ oder $z_R$ ersetzt. Ist dann $Z=XY$ unter der Bedingung (1.), so heisst $X$ eine zur Gruppe $\Hgr$ gehörige Matrix.

Seien $\Phi,\Phi',\cdots$ verschiedene Primfactoren der Determinante $\Theta$ der Gruppe $\Hgr$, seien $f,f',\cdots$ ihre Grade, $\chi(R),\psi(R),\cdots$ ihre Charaktere, $\chi'(R)=\chi(R^{-1}),\psi'(R)=\psi(R^{-1}),\cdots$ die conjugirten Charaktere. Dann ist nach \emph{Primfactoren}, § 8
\begin{equation}\tag{2.}
\Phi\!\left(\eps+u\frac{f}{h}\chi'\right)=(1+u)^f,
\qquad
\Phi\!\left(\eps+u\frac{f'}{h}\psi'\right)=1,\cdots .
\end{equation}
Setzt man $x_R=\frac{f}{h}\chi'(R)$, so sei $X=A$, setzt man $x_R=\frac{f'}{h}\psi'(R)$, so sei $X=B$, u. s. w. Dann ist
\[
\sum \frac{f}{h}\chi'(P)\frac{f}{h}\chi'(Q)=\frac{f}{h}\chi'(R),
\qquad
\sum \frac{f}{h}\chi'(P)\frac{f'}{h}\psi'(Q)=0\qquad(PQ=R),
\]
und mithin $A^2=A$, $B^2=B$, $AB=BA=0$. Da
\[
\sum \frac{f}{h}\chi'(P)x_Q=\sum x_P\frac{f}{h}\chi'(Q)\qquad(PQ=R)
\]
ist, so ist $A$ (und ebenso $B,C,\cdots$) mit jeder Matrix $X$ vertauschbar. Weil endlich
\[
\frac{f}{h}\chi'(R)+\frac{f'}{h}\psi'(R)+\cdots=\eps_R
\]
ist, wenn die Summe über alle $k$ Charaktere von $\Hgr$ erstreckt wird, so ist
\begin{equation}\tag{3.}
A+B+C+\cdots=E.
\end{equation}
Mithin ist
\[
(E+uAX)(E+uBX)(E+uCX)\cdots=E+u(A+B+C+\cdots)X=E+uX.
\]
Alle anderen Glieder in der Entwicklung des symbolischen Productes nach Potenzen von $u$ verschwinden, weil $AXBX=ABXX=0$ ist. Folglich ist auch
\begin{equation}\tag{4.}
|E+uAX|\,|E+uBX|\cdots=|E+uX|.
\end{equation}
Diese Determinante aber ist, weil $X$ eine zur Gruppe $\Hgr$ gehörige Matrix ist, ein Product von Potenzen der Primfunctionen $\Phi(\eps+ux),\Phi'(\eps+ux),\cdots$. Daher ist auch
\[
|E+uAX|=\Phi(\eps+ux)^s\Phi'(\eps+ux)^t\cdots.
\]
Es muss aber $t=0$ sein. Denn setzt man $x_R=\frac{f'}{h}\psi'(R)$, also $X=B$, so wird $AX=AB=0$, also die Determinante links gleich $1$. Rechts aber ist $\Phi'(\eps+u\frac{f'}{h}\psi')=(1+u)^{f'}$. Folglich ist
\begin{equation}\tag{5.}
|E+uAX|=\Phi(\eps+ux)^s,
\qquad
|E+uBX|=\Phi'(\eps+ux)^{s'},\cdots
\end{equation}
und mithin nach (4.)
\begin{equation}\tag{6.}
|X|=\Phi(x)^s\Phi'(x)^{s'}\cdots,
\end{equation}
und
\[
n=fs+f's'+\cdots.
\]
Setzt man in der Gleichung (5.) $x_R=\frac{f}{h}\chi'(R)$, also $X=A$, so erhält man
\begin{equation}\tag{7.}
|E+uA|=(1+u)^{fs},
\qquad |uE-A|=(u-1)^{fs}u^{n-fs}.
\end{equation}
Da $A^2=A$ ist, so sind die Elementartheiler dieser Determinante alle linear. Mithin ist der Rang von $A$ gleich $r=fs$, und die Hauptunterdeterminanten $r^{\ten}$ Grades von $A$ sind nicht alle Null, da ihre Summe gleich $1$ ist. Ist $s=0$, so ist $r=0$ und folglich $A=0$. In der Summe (3.) kommen daher nur die Matrizen $A,B,C,\cdots$ wirklich vor, welche aus den Charakteren $\chi$ solcher Primfunctionen $\Phi$ gebildet sind, die in $|X|$ aufgehen. Ist z. B. $|X|$ eine Potenz von $\Phi$, so ist $A=E$. Sei
\begin{equation}\tag{8.}
|a_{\kappa\lambda}|\qquad(\kappa,\lambda=\alpha_1,\alpha_2,\cdots,\alpha_r)
\end{equation}
eine von Null verschiedene Hauptunterdeterminante $r^{\ten}$ Grades von $A$. Sei $r'=f's'$ der Rang von $B$ und sei
\[
|b_{\kappa\lambda}|\qquad(\kappa,\lambda=\beta_1,\beta_2,\cdots,\beta_{r'})
\]
eine von Null verschiedene Hauptunterdeterminante $r'^{\ten}$ Grades von $B$ u. s. w.

Sei $M$ die Matrix des Grades $n=r+r'+\cdots$, deren $n$ Zeilen aus der Zeile
\[
a_{\lambda,\alpha_1},\cdots,a_{\lambda,\alpha_r},
 b_{\lambda,\beta_1},\cdots,b_{\lambda,\beta_{r'}},\cdots
\]
erhalten werden, indem man $\lambda=1,2,\cdots,n$ setzt. Ebenso sei $L'$ die Matrix
\[
a_{\alpha_1,\kappa},\cdots,a_{\alpha_r,\kappa},
 b_{\beta_1,\kappa},\cdots,b_{\beta_{r'},\kappa},\cdots,
\]
und $L$ die conjugirte Matrix. Dann bilde ich die Matrix $LXM$. Sind $\rho$ und $\sigma$ zwei der Indices $1,2,\cdots,r$, so ist das $\sigma^{\te}$ Element der $\rho^{\ten}$ Zeile
\[
\sum_{\kappa,\lambda}a_{\alpha\kappa}x_{\kappa\lambda}a_{\lambda\beta},
\]
wo $\alpha=\alpha_\rho$ und $\beta=\alpha_\sigma$ ist. Dies ist ein Element der Matrix $AXA=XAA=XA=AX$, also gleich
\begin{equation}\tag{9.}
\xi_{\alpha\beta}=\sum_\lambda a_{\alpha\lambda}x_{\lambda\beta}
=\sum_\lambda x_{\alpha\lambda}a_{\lambda\beta}
\end{equation}
und geht aus $x_{\alpha\beta}$ hervor, indem man darin $x_R$ durch
\begin{equation}\tag{10.}
\xi_R=\sum_S \frac{f}{h}\chi'(RS^{-1})x_S
\end{equation}
ersetzt. Ist $\rho$ eine der Zahlen $1,2,\cdots,r$, und $\sigma$ eine der Zahlen $1,2,\cdots,r'$, so ist das $(r+\sigma)^{\te}$ Element der $\rho^{\ten}$ Zeile
\[
\sum_{\kappa,\lambda}a_{\alpha\kappa}x_{\kappa\lambda}b_{\lambda\beta},
\]
wo $\alpha=\alpha_\rho$, $\beta=\beta_\sigma$ ist. Dies ist ein Element der Matrix $AXB=XAB=0$. Folglich zerfällt $X$ in Theilmatrizen der Grade $r,r',\cdots$, deren erste von den $r^2$ Elementen
\[
(\xi_{\kappa\lambda})\qquad(\kappa,\lambda=\alpha_1,\alpha_2,\cdots,\alpha_r)
\]
gebildet wird. Darin ist $x_E$ mit der Matrix
\[
N_1=(a_{\kappa\lambda})\qquad(\kappa,\lambda=\alpha_1,\alpha_2,\cdots,
\alpha_r)
\]
multiplicirt, deren Determinante nicht verschwindet. Ebenso sei
\[
N_2=(b_{\kappa\lambda})\qquad(\kappa,\lambda=\beta_1,\beta_2,\cdots,
\beta_{r'})
\]
und
\[
N=\begin{pmatrix}
N_1&0&0&\cdots\\
0&N_2&0&\cdots\\
0&0&N_3&\cdots
\end{pmatrix}.
\]
Dann zerfällt auch $LXMN^{-1}=Z$ in Theilmatrizen der Grade $r,r',\cdots$, deren erste ist
\begin{equation}\tag{11.}
(\xi_{\kappa\lambda})(a_{\kappa\lambda})^{-1}
\qquad(\kappa,\lambda=\alpha_1,\alpha_2,\cdots,\alpha_r).
\end{equation}
In der Matrix $Z$ ist daher $x_E$ mit der Hauptmatrix $E$ multiplicirt. Setzt man nun $x_R=\eps_R$, also $X=E$, so wird auch $Z=E$, und mithin ist $LMN^{-1}=E$. Daher sind $|L|$ und $|M|$ von Null verschieden, und es ist $MN^{-1}=L^{-1}$. Folglich ist $Z=LXL^{-1}$ eine zur Gruppe $\Hgr$ gehörige Matrix, und ebenso jede der Theilmatrizen, wie (11.), worin sie zerfällt.

Die Matrix $A$ hat den Rang $r$. Nach \emph{Primfactoren}, § 11 verhalten sich daher die Determinanten $r^{\ten}$ Grades der Matrix $AX=XA$, wie die entsprechenden der Matrix $A$, unterscheiden sich also nur durch constante Factoren von einander. Nun ist nach (5.) und (7.) die Summe der Hauptunterdeterminanten $r^{\ten}$ Grades von $AX$ gleich $\Phi(x)^s$ und von $A$ gleich $1$. Folglich ist die Determinante der Matrix (11.) gleich $\Phi(x)^s$.

\section*{§ 3.}

Die Matrix $(x_{PQ^{-1}})$ der Gruppe $\Hgr$ habe ich \emph{Darstellung}, § 5 in eine aequivalente transformirt, die zerfällt in $f$ einander gleiche Theilmatrizen
\begin{equation}\tag{1.}
(u_{\alpha\beta})\qquad(\alpha,\beta=1,2,\cdots,f)
\end{equation}
des Grades $f$, in $f'$ einander gleiche Theilmatrizen
\begin{equation}\tag{2.}
(u'_{\alpha\beta})\qquad(\alpha,\beta=1,2,\cdots,f')
\end{equation}
des Grades $f'$, u. s. w. Die $f^2+f'^2+\cdots=h$ Elemente $u_{\alpha\beta},u'_{\alpha\beta},\cdots$ dieser Matrizen sind $h$ von einander unabhängige lineare Functionen der $h$ Variabelen $x_R$. Ihre Determinanten sind
\[
|u_{\alpha\beta}|=\Phi(x),\qquad |u'_{\alpha\beta}|=\Phi'(x),\cdots.
\]
In derselben Weise kann die im vorigen Paragraphen betrachtete zur Gruppe $\Hgr$ gehörige Matrix $X$ in eine aequivalente transformirt werden, die in $s$ Theilmatrizen (1.), in $s'$ Theilmatrizen (2.) u. s. w. zerfällt.

Ist $|X|$ durch mehrere verschiedene Primfunctionen $\Phi,\Phi',\cdots$ theilbar, so kann $X$ nach § 2 in eine zerfallende Matrix transformirt werden, worin jede Theilmatrix eine Potenz einer Primfunction zur Determinante hat. Es genügt daher, die Behauptung für den Fall zu beweisen, dass $|X|=\Phi^s$ eine Potenz einer Primfunction $\Phi$ ist.

Die $h$ Variabelen $x_R$ sind lineare Functionen der $h$ Variabelen $u_{\alpha\beta},u'_{\alpha\beta},\cdots$. Drückt man die Elemente $x_{\kappa\lambda}$ von $X$ durch diese aus, so hängt $\Phi$, also auch $|X|$ nur von den $f^2$ Variabelen $u_{\alpha\beta}$ ab. Dass aber auch jedes der $n^2$ Elemente $x_{\kappa\lambda}$ nur von diesen Grössen abhängt, ergiebt sich aus den Entwicklungen des § 2. Nach diesen ist im vorliegenden Falle $A=E$, also $X=XA$, demnach $x_{\kappa\lambda}=\xi_{\kappa\lambda}$, und mithin kann $x_{\kappa\lambda}$ als Function der $h$ Variabelen $\xi_R$ dargestellt werden. Unter diesen befinden sich genau $f^2$ linear unabhängige. Andererseits lässt sich $\Phi$ durch die $f^2$ von einander unabhängigen Variabelen $u_{\alpha\beta}$ darstellen, und nicht durch eine lineare Substitution in eine Function von weniger als $f^2$ Variabelen transformiren. Folglich sind die $h$ Grössen $\xi_R$, und mithin auch die $n^2$ Grössen $x_{\kappa\lambda}$ lineare Functionen der $f^2$ Variabelen $u_{\alpha\beta}$.

Ersetzt man $x_R$ durch $y_R$ oder $z_R$, so gehe $u=(u_{\alpha\beta})$ in $(v)$ oder $(w)$, und $X$ in $Y$ oder $Z$ über. Dann ist $w=uv$ und $Z=XY$, falls $z_R$ durch die Gleichung (1.), § 2 definirt ist. Die Elemente $x_{\kappa\lambda}$ von $X$ sind also solche lineare Functionen der Elemente $u_{\alpha\beta}$ von $u$, dass $Z=XY$ ist, wenn $w=uv$ ist. Demnach ergiebt sich die Möglichkeit der behaupteten Transformation von $X$ aus den Entwicklungen des § 1.

Man hätte bei diesem Beweise auch die Sätze des § 2 entbehren können und nur den Sätzen des § 1 eine etwas allgemeinere Fassung zu geben brauchen. Statt von einer Matrix $u$ kann man von mehreren Matrizen $u,u',\cdots$ der Grade $f,f',\cdots$ ausgehen, deren $f^2+f'^2+\cdots$ Elemente lauter von einander unabhängige Variabele sind. Die Elemente der Matrix $X$ sind dann lineare Functionen aller dieser Variabelen, und es ist $Z=XY$, wenn gleichzeitig $w=uv,w'=u'v',\cdots$ ist. Indessen ist die in § 2 ausgeführte Transformation auch an sich von Interesse, weil dazu nur die Kenntniss der Charaktere erforderlich ist.

\section*{§ 4.}

Jeder Primfactor $f^{\ten}$ Grades $\Phi$ der Gruppendeterminante $\Theta$ lässt sich als eine Determinante $f^{\ten}$ Grades darstellen, deren Elemente $f^2$ von einander unabhängige lineare Functionen der $h$ Variabelen $x_R$ sind. Diese Darstellung kann so gewählt werden, dass die Elemente der Determinante eine zur Gruppe $\Hgr$ gehörige Matrix bilden. Die $f^{\te}$ Potenz von $\Phi$ habe ich (\emph{Darstellung}, § 3) durch eine Determinante des Grades $f^2$ ausgedrückt. Diese Darstellung ist dadurch ausgezeichnet, dass ihre Elemente alle aus der linearen Function $\sum_R\chi(R)x_R$ erhalten werden, indem man die $h$ Variabelen $x_R$ in besonderer Weise permutirt. Jede Determinante des Grades $f^2$ gebildet aus den Elementen der Matrix
\begin{equation}\tag{1.}
\sum_R\chi(R)x_{RPQ^{-1}}
\end{equation}
ist bis auf einen constanten Factor gleich $\Phi^f$.

Ich will nun zeigen, dass auch die Determinante $f^{\ten}$ Grades, die gleich $\Phi$ ist, auf eine ähnliche Form gebracht werden kann. Auch zur Bestimmung dieser Determinante genügt es, eine einzige lineare Function $\sum a_{R^{-1}}x_R$ zu berechnen. Dann ist jede Determinante $f^{\ten}$ Grades der Matrix
\begin{equation}\tag{2.}
x_{P,Q}=\sum_R a_{R^{-1}}x_{PRQ^{-1}}
\end{equation}
bis auf einen constanten Factor gleich $\Phi$. Während aber der Charakter $\chi(R)$ durch die Primfunction $\Phi$ vollständig bestimmt ist, können die $h$ Grössen $a_R$ auf unendlich viele Arten gewählt werden, abgesehen von dem Falle $f=1$, wo beide Darstellungen zusammenfallen. Dies liegt daran, dass $\chi(PQ)=\chi(QP)$ ist, während die Grössen $a_R$ dieser Bedingung nicht genügen.

Um zu dieser Darstellung von $\Phi$ zu gelangen, brauche ich einige Hülfssätze über Matrizen, die zwar bekannt sind, aber hier kurz entwickelt werden sollen. Zur Bequemlichkeit der Darstellung bezeichne ich mit den Zeichen $A,B,C,\cdots$ nicht nur Matrizen $h^{\ten}$ Grades, sondern zugleich bilineare Formen, deren Coefficienten die Elemente jener Matrizen sind, wie in meiner Arbeit \emph{Über lineare Substitutionen und bilineare Formen}, \textsc{Crelle}'s Journal Bd. 84. Die folgenden Sätze gelten für eine beliebige Matrix $K$ (a. a. O. § 13), sollen aber hier nur für eine solche abgeleitet werden, deren charakteristische Determinante $|uE-K|$ lauter lineare Elementartheiler hat. Dann kann man eine Substitution $P$ (d. h. eine Matrix $P$ von nicht verschwindender Determinante) so bestimmen, dass die bilineare Form $P^{-1}KP$ die Normalform
\[
\rho(u_1v_1+\cdots+u_fv_f)+\rho'(u_{f+1}v_{f+1}+\cdots+u_{f+f'}v_{f+f'})+\cdots
\]
annimmt. Hier sind $\rho,\rho',\cdots$ die verschiedenen Wurzeln jener charakteristischen Gleichung, und zwar $\rho$ eine $f$ fache, $\rho'$ eine $f'$ fache, u. s. w. Da die für $u=\rho$ verschwindenden Elementartheiler von $|uE-K|$ alle linear sind, so fängt die Entwicklung
\begin{equation}\tag{3.}
(uE-K)^{-1}=\frac{A}{u-\rho}+\cdots
\end{equation}
mit der $(-1)^{\ten}$ Potenz von $u-\rho$ an. Aus der Normalform ist ersichtlich, dass $P^{-1}AP=u_1v_1+\cdots+u_fv_f$ ist. Mithin hat $A$ den Rang $f$, genügt der Gleichung
\begin{equation}\tag{4.}
A^2=A,
\end{equation}
und es ist
\begin{equation}\tag{5.}
|E+\mu A|=(\mu+1)^f.
\end{equation}
Folglich ist die Summe der Hauptunterdeterminanten $f^{\ten}$ Grades von $A$ gleich $1$. In der Matrix $A$, worin die Unterdeterminanten $(f+1)^{\ten}$ Grades sämmtlich verschwinden, können daher die Hauptunterdeterminanten $f^{\ten}$ Grades nicht alle Null sein. Ist ferner
\[
(uE-K)^{-1}=\frac{B}{u-\rho'}+\cdots,
\]
so ist $P^{-1}BP=u_{f+1}v_{f+1}+\cdots+u_{f+f'}v_{f+f'}$, und folglich ist
\begin{equation}\tag{6.}
AB=BA=0.
\end{equation}
Ohne die Normalform zu benutzen, kann man zu diesen Sätzen durch die folgenden Betrachtungen gelangen. Indem man die Gleichung (3.) einmal quadrirt, das andere Mal nach $u$ differentiirt, erhält man die beiden Relationen
\[
(uE-K)^{-2}=\frac{A^2}{(u-\rho)^2}+\cdots,
\qquad
-(uE-K)^{-2}=-\frac{A}{(u-\rho)^2}+\cdots,
\]
aus denen sich sofort die Eigenschaft (4.) ergiebt. (Vergl. \textsc{Weierstrass}, Monatsber. 1858, S. 215).

Aus der identischen Gleichung
\[
(uE-K)\bigl((uE-K)^{-1}-(vE-K)^{-1}\bigr)(vE-K)=(vE-K)-(uE-K)=(v-u)E
\]
folgt die Relation
\[
-\frac{(uE-K)^{-1}-(vE-K)^{-1}}{u-v}=(uE-K)^{-1}(vE-K)^{-1}.
\]
Entwickelt man diese Ausdrücke nach aufsteigenden Potenzen von $u-\rho$ und $v-\rho'$, so erhält man
\[
-\frac{1}{(u-\rho)-(v-\rho')+(\rho-\rho')}
\left(\frac{A}{u-\rho}+\cdots-\frac{B}{v-\rho'}-\cdots\right)
=\frac{AB}{(u-\rho)(v-\rho')}+\cdots.
\]
Sind $u-\rho$ und $v-\rho'$ hinlänglich klein, so kommen in der Entwicklung des ersten Factors der linken Seite nur positive Potenzen von $u-\rho$ und $v-\rho'$ vor, und folglich findet sich links keine negative Potenz von $u-\rho$ mit einer negativen von $v-\rho'$ multiplicirt. Daher ist $AB=0$.

Da $(uE-K)^{-1}$ eine echt gebrochene rationale Function von $u$ ist, so ist
\begin{equation}\tag{7.}
(uE-K)^{-1}=\frac{A}{u-\rho}+\frac{B}{u-\rho'}+\cdots.
\end{equation}
Entwickelt man beide Seiten dieser Gleichung nach absteigenden Potenzen von $u$, so ergiebt sich durch Vergleichung der Coefficienten
\begin{equation}\tag{8.}
E=A+B+C+\cdots,
\qquad
K=\rho A+\rho'B+\rho''C+\cdots.
\end{equation}
Dass endlich der Rang von $A$ gleich $f$ ist, kann man so einsehen. Ist $\psi(u)=(u-\rho)(u-\rho')(u-\rho'')\cdots$, so ist $\psi(K)=0$ die Gleichung niedrigsten Grades, der $K$ genügt. Ist dann
\[
\frac{\psi(u)-\psi(v)}{u-v}=\psi(u,v),
\]
so ist auch
\[
\frac{\psi(u)E-\psi(K)}{uE-K}=\psi(u,K),
\]
also
\[
(uE-K)^{-1}=\frac{\psi(u,K)}{\psi(u)}.
\]
Daher ist
\[
A=\frac{\psi(\rho,K)}{\psi'(\rho)}=g(K),
\]
wo
\[
\psi'(\rho)g(v)=\psi(\rho,v)=\frac{\psi(v)}{v-\rho}=(v-\rho')(v-\rho'')\cdots
\]
ist. Mithin ist $g(\rho)=1$, $g(\rho')=0$, $g(\rho'')=0,\cdots$.

Nun folgt aus (3.)
\[
\bigl((\rho E-K)+(u-\rho)E\bigr)\left(\frac{A}{u-\rho}+\cdots\right)=E,
\]
und mithin ist
\begin{equation}\tag{9.}
(\rho E-K)A=0.
\end{equation}
Da die Determinante $|uE-K|$ $f$ lineare Elementartheiler $u-\rho$ hat, so ist der Rang der Matrix $\rho E-K$ gleich $h-f$. Nach Gleichung (9.) kann daher der Rang der Matrix $A$ höchstens gleich $f$ sein. Er muss gleich $f$ sein, wenn sich unter den verschiedenen Spalten von $A$ $f$ unabhängige Lösungen der linearen Gleichungen
\[
(\rho e_{\alpha1}-k_{\alpha1})v_1+(\rho e_{\alpha2}-k_{\alpha2})v_2+\cdots+(\rho e_{\alpha h}-k_{\alpha h})v_h=0
\]
finden, wenn sich also zeigen lässt: Ist $X$ irgend eine Lösung der Gleichung $(\rho E-K)X=0$, so kann man $Y$ so bestimmen, dass $X=AY$ ist. Nun folgt aus $KX=\rho X$, dass $K^2X=\rho KX=\rho^2X,\cdots,K^nX=\rho^nX$, also auch $g(K)X=g(\rho)X$ ist, wenn $g(v)$ eine ganze Function von $v$ ist. Ist, wie oben, $(v-\rho)\psi'(\rho)g(v)=\psi(v)$, so ist $g(K)=A$, $g(\rho)=1$, also $AX=X$, womit die Behauptung bewiesen ist.

\section*{§ 5.}

In der Gruppenmatrix $X=(x_{PQ^{-1}})$ gebe ich den $h$ unabhängigen Variabelen $x_R$ solche constanten Werthe $x_R=k_R$, dass die $f+f'+\cdots$ Wurzeln $\rho,\rho',\cdots$ der Gleichung $\Phi(k-u\eps)\Phi'(k-u\eps)\cdots=0$ alle unter einander verschieden sind. Ist dann $\rho$ eine Wurzel der Gleichung $f^{\ten}$ Grades $\Phi(k-u\eps)=0$, so hat die charakteristische Determinante der Matrix $K=(k_{PQ^{-1}})$ $f$ lineare Elementartheiler $u-\rho$. Ist nun, nach Potenzen von $u-\rho$ (oder $u-\rho',\cdots$) entwickelt
\begin{equation}\tag{1.}
(uE-K)^{-1}=\frac{A}{u-\rho}+\cdots=\frac{B}{u-\rho'}+\cdots,
\end{equation}
so bieten die Matrizen $A,B,\cdots$ dieselben Symmetrieverhältnisse dar wie $X$, d. h. man kann $A=(a_{PQ^{-1}})$, $B=(b_{PQ^{-1}})\cdots$ setzen. Dann ist $A^2=A$, $AB=0$, also
\begin{equation}\tag{2.}
\sum a_Pa_Q=a_R\qquad(PQ=R)
\end{equation}
und
\begin{equation}\tag{3.}
\sum a_Pb_Q=0\qquad(PQ=R).
\end{equation}
Die Matrix $A$ hat den Rang $f$ und ist eine ganze Function von $K$, $A=g(K)$, und zwar ist $g(\rho)=1$, für jede andere Wurzel aber $g(\rho')=0$. Sind also $\rho,\rho_1,\cdots,\rho_{f-1}$ die Wurzeln der Gleichung $\Phi(k-u\eps)=0$, so sind $g(\rho)=1,g(\rho_1)=0,\cdots,g(\rho_{f-1})=0$ die der Gleichung $\Phi(a-u\eps)=0$, wie ich \emph{Primfactoren}, § 3 gezeigt habe.

Demnach ist
\begin{equation}\tag{4.}
\Phi(a+u\eps)=u^{f-1}(u+1),
\qquad
\Phi'(a+u\eps)=u^{f'},
\qquad
\Phi''(a+u\eps)=u^{f''},\cdots .
\end{equation}
Ist $x_{PQ}=x_{QP}$, so ist $\Phi(x)$ gleich der $f^{\ten}$ Potenz einer linearen Function der Variabelen $x_R$. Der ersten Gleichung nach können daher die Grössen $a_R$ jene Eigenschaft nicht besitzen, falls $f>1$ ist.

Nun ist die Summe der Wurzeln der Gleichung $\Phi(x-u\eps)=0$ gleich $\sum_R\chi(R)x_R$. Daher ist, wenn $\psi$ irgend einen der $k-1$ von $\chi$ verschiedenen Charaktere bezeichnet,
\begin{equation}\tag{5.}
\sum_R\chi(R)a_R=1,
\qquad
\sum_R\psi(R)a_R=0.
\end{equation}
Ein System von $h$ Grössen $a_R$, das den Gleichungen (2.) und (5.) genügt, nenne ich ein die Primfunction $\Phi$ oder die entsprechende primitive Darstellung der Gruppe $\Hgr$ determinirendes Werthsystem.

Durchläuft $R$ die $h_\rho$ Elemente der $\rho^{\ten}$ Classe, so sei
\[
\sum_{(\rho)} a_R=h_\rho a_\rho.
\]
Da $\chi(R)$ für diese $h_\rho$ Elemente denselben Werth $\chi_\rho$ hat, so folgt aus (5.)
\[
\sum_\rho h_\rho\chi_\rho a_\rho=1,
\qquad
\sum_\rho h_\rho\psi_\rho a_\rho=0.
\]
Durch diese $k$ Gleichungen sind aber die $k$ Grössen $a_\rho$ vollständig bestimmt. Nach \emph{Gruppencharaktere}, § 3 genügen ihnen die Werthe
\[
a_\rho=\frac1h\chi_{\rho'}.
\]
Mithin ist
\begin{equation}\tag{6.}
\sum_{(\rho)}a_R=\frac{h_\rho}{h}\chi_{\rho'},
\qquad
\sum_S a_{S^{-1}RS}=\chi(R^{-1}).
\end{equation}
Z. B. ist
\begin{equation}\tag{7.}
a_E=\frac{f}{h}.
\end{equation}
Die $k$ linearen Gleichungen (5.) sind den $k$ linearen Gleichungen (6.) völlig aequivalent. Den Gleichungen (2.) und (5.), durch die wir ein die Primfunction $\Phi$ determinirendes Werthsystem $a_R$ definirt haben, sind aber auch die Gleichungen (2.) und (4.) aequivalent. Denn oben haben wir (5.) aus (4.) erhalten. Benutzt man aber die Bezeichnungen \emph{Primfactoren}, § 1, (6.), so ist nach (2.) $a_R^{(2)}=a_R$, also auch $a_R^{(n)}=a_R$. Nun ist die Summe der $n^{\ten}$ Potenzen der Wurzeln der Gleichung $\Phi(a-u\eps)=0$ gleich
\[
\sum_R\chi(R)a_R^{(n)}=\sum_R\chi(R)a_R=1,
\]
und da durch die Potenzsummen die Wurzeln völlig bestimmt sind, so ist eine derselben $1$, die andere $f-1$ Null. Ebenso erkennt man, dass die $f'$ Wurzeln der Gleichung $\Phi'(a-u\eps)=0$ alle verschwinden. Daher ist die Determinante $h^{\ten}$ Grades
\[
|a_{PQ^{-1}}+u\eps_{PQ^{-1}}|=(u+1)^f u^{h-f}.
\]
In Folge der Gleichung $A^2=A$ sind ihre Elementartheiler alle linear. Mithin hat die Matrix $A$ den Rang $f$ und die Matrix $A-E$ den Rang $h-f$.

Ist $H$ ein festes, $R$ ein veränderliches Element von $\Hgr$, und setzt man (unabhängig von der in Gleichung (3.) benutzten Bezeichnung) $b_R=a_{H^{-1}RH}$, so genügen die $h$ Grössen $b_R$ den Gleichungen (2.) und (5.), bilden also ein die Function $\Phi$ determinirendes Werthsystem. Sind allgemeiner $v_R$ $h$ unabhängige Variabele, ist $V=(v_{PQ^{-1}})$ und $V^{-1}AV=B=(b_{PQ^{-1}})$, so bilden die $h$ Grössen $b_R$ ein solches Werthsystem (\emph{Primfactoren}, § 1, (9.)). Wie ich jetzt zeigen will, erhält man jedes solche Werthsystem auf diese Weise, d. h. indem man für $V$ jede Gruppenmatrix von nicht verschwindender Determinante setzt.

Sind in der Gruppenmatrix $X=(x_{PQ^{-1}})$ die $h$ Grössen $x_R$ unabhängige Variabele, so kann man ($D.$ § 5) eine constante Matrix $L$ (deren $h^2$ Elemente von den Variabelen $x_R$ unabhängig sind, und deren Determinante nicht verschwindet) so bestimmen, dass $L^{-1}XL=Z$ zerfällt in $f$ Matrizen $Z_1,Z_2,\cdots,Z_f$ des Grades $f$, die in den Elementen übereinstimmen, in $f'$ Matrizen $Z_{f+1},\cdots,Z_{f+f'}$ des Grades $f'$, u. s. w. Die $f^2+f'^2+\cdots=h$ Elemente der Matrizen $Z_1,Z_{f+1},\cdots$ sind $h$ von einander unabhängige lineare Functionen der $h$ Variabelen $x_R$. Wenn daher $C$ irgend eine Matrix ist, die in derselben Weise zerfällt, wie $Z$, und wofür $C_1=C_2=\cdots=C_f$, $C_{f+1}=\cdots=C_{f+f'}$, $\cdots$ ist, so kann $LCL^{-1}$ aus $X$ erhalten werden, indem man den Variabelen $x_R$ bestimmte Werthe ertheilt. Die charakteristische Determinante von $Z_1$ ist $\Phi(x-u\eps)$, die von $Z_{f+1}$ ist $\Phi'(x-u\eps)$ u. s. w. Die Zeichen $L,Z,C$ bedeuten hier keine Gruppenmatrizen.

Setzt man $x_R=a_R$, so mögen $Z$ und $Z_\lambda$ in $C$ und $C_\lambda$ übergehen. Da $A^2=A$ ist, so sind die Elementartheiler von $|uE-A|$ alle linear, also auch die von $|uE-C|$. Zerfällt aber eine Determinante in Theile, so sind ihre Elementartheiler die der einzelnen Theile zusammengenommen. Folglich sind auch die Elementartheiler von
\[
|uE_1-C_1|=(u-1)u^{f-1}
\]
linear. Nun sei $B=(b_{PQ^{-1}})$ eine zweite Matrix, deren Elemente den Bedingungen (2.) und (6.) genügen. Dann zerfällt $L^{-1}BL=D$ entsprechend in die Theile $D_1,D_2,\cdots$, und $uE_1-C_1$ und $uE_1-D_1$ stimmen in den Elementartheilern überein. Daher kann man eine Matrix $M_1$ von nicht verschwindender Determinante so bestimmen, dass $M_1^{-1}C_1M_1=D_1$ ist. Stimmen $M_2,\cdots,M_f$ in den Elementen mit $M_1$ überein, so ist auch $M_2^{-1}C_2M_2=D_2$, u. s. w. Ebenso bestimme man $M_{f+1},M_{f+2},\cdots$, und setze aus diesen Theilen die Matrix $M$ zusammen. Dann ist $LML^{-1}=K$ eine Gruppenmatrix, und aus $M^{-1}CM=D$ folgt
\[
(LML^{-1})^{-1}(LCL^{-1})(LML^{-1})=LDL^{-1}
\]
oder $K^{-1}AK=B$.

Nunmehr ist es auch leicht, die allgemeinste Gruppenmatrix $V$ anzugeben, die $A$ in $B$ transformirt. Sind nämlich $u_R$ $h$ unabhängige Variabele, und ist $U=(u_{PQ^{-1}})$, so ist
\begin{equation}\tag{8.}
V=AUB+(E-A)U(E-B).
\end{equation}
Denn aus den Gleichungen $A^2=A$ und $B^2=B$ folgt $AV=AUB$ und $VB=AUB$, also $AV=VB$. Ist ferner $K$ irgend eine gegebene Gruppenmatrix, die der Bedingung $K^{-1}AK=B$ genügt, so kann man den Variabelen $u_R$ solche Werthe geben, dass $V=K$ wird. Denn dann ist $AK=KB$, $(E-A)K=K(E-B)$. Setzt man also $U=K$, so wird
\[
\begin{aligned}
V&=A(KB)+(E-A)(K(E-B))\\
&=A(AK)+(E-A)((E-A)K)\\
&=AK+(E-A)K=K.
\end{aligned}
\]
Nun habe ich oben die Existenz einer Gruppenmatrix $K$ bewiesen, deren Determinante nicht verschwindet, und die der Bedingung $AK=KB$ genügte. Für $U=K$ ist daher $|V|$ von Null verschieden, und folglich kann $|V|$ für unbestimmte $u_R$ nicht verschwinden.

Mit Hülfe des eben entwickelten Satzes kann man aus den quadratischen und linearen Gleichungen (2.) und (5.) allgemeinere lineare Relationen zwischen den Grössen $a_R$ herleiten. Die Gruppenmatrix
\[
C=\left(\frac{f}{h}\chi'(PQ^{-1})\right)
\]
ist mit jeder Gruppenmatrix vertauschbar und genügt der Gleichung $C^2=C$. Setzt man also $AC=L=(l_{PQ^{-1}})$, so ist $L^2=L$. Auch den Gleichungen (5.) genügen die Grössen $l_R$, wie sich aus \emph{Primfactoren}, § 5, (6.) oder § 8 leicht ergiebt. Daher kann man eine Gruppenmatrix $K$ so bestimmen, dass $A=K^{-1}ACK=K^{-1}AKC$ wird, oder wenn man $K^{-1}AK=M$ setzt, $MC=A$.

Mithin ist $AC=MC^2=MC=A$. Ist $\Psi$ ein von $\Phi$ verschiedener Primfactor von $\Theta$, $\psi$ sein Charakter und $D=\left(\frac{f'}{h}\psi'(PQ^{-1})\right)$, so ist $CD=0$, und daher auch $AD=(AC)D=A(CD)=0$. Ist $b_R$ irgend ein die Primfunction $\Psi$ determinirendes Werthsystem und $B=(b_{PQ^{-1}})$, so ist $B=BD$, und mithin $AB=(AC)(BD)=(AB)(CD)=0$. Ist $X$ eine beliebige Gruppenmatrix, so kann man $B$ durch $XBX^{-1}$ ersetzen. Daher ist auch $A(XBX^{-1})=0$ und folglich
\begin{equation}\tag{9.}
AXB=0.
\end{equation}
Die $h$ Grössen $a_R$ genügen demnach den Gleichungen $AC=A$ oder
\begin{equation}\tag{10.}
\sum_R\chi(RS^{-1})a_R=\frac{h}{f}a_S
\end{equation}
und $AD=0$, oder
\begin{equation}\tag{11.}
\sum_R\psi(RS^{-1})a_R=0.
\end{equation}
Die letzteren sind, wie die obige Herleitung der Relation $AD=0$ aus $AC=0$ zeigt, eine Folge der ersteren. Diese Gleichungen (10.) aber, oder $(E-C)A=0$, enthalten, da der Rang der Matrix $E-C$ gleich $h-f^2$ ist, genau $h-f^2$ unabhängige homogene lineare Relationen zwischen den $h$ Grössen $a_R$. Die $k$ Gleichungen (5.) sind aber unter den Gleichungen (10.) und (11.) enthalten, falls man zu diesen noch die nicht homogene Gleichung (7.) hinzunimmt. In der That ist die Anzahl der aufgestellten unabhängigen linearen Gleichungen zwischen den $h$ Grössen $a_R$ $h-f^2+1\geq k$, weil $h-k=(f^2-1)+(f'^2-1)+(f''^2-1)+\cdots$ ist.

Die Gleichungen (5.) lassen sich auf die Form (6.) bringen. In derselben Weise lassen sich die Gleichungen (10.) und (11.) in
\begin{equation}\tag{12.}
f\sum_{(\rho)}a_{PRQ}=h_\rho\chi_{\rho'}a_{PQ}
\end{equation}
transformiren, wo $R$ die $h_\rho$ Elemente der $\rho^{\ten}$ Classe durchläuft.

\section*{§ 6.}

Ist $X=(x_{PQ^{-1}})$ irgend eine Gruppenmatrix, so setze ich $\overline{X}=(x_{Q^{-1}P})$. Ist dann $Z=XY$, so ist $\overline{Z}=\overline{Y}\,\overline{X}$. Die beiden Matrizen $X$ und $\overline{Y}$ sind mit einander vertauschbar. Seien $u_1,\cdots,u_f$ die Wurzeln der Gleichung $\Phi(x-u\eps)=0$, und $v_1,\cdots,v_f$ die der Gleichung $\Phi(y-u\eps)=0$. Ist dann $g(u,v)$ eine ganze Function von $u$ und $v$, so sind die Wurzeln der charakteristischen Gleichung der Matrix $g(X,\overline{Y})$ die $f^2$ Grössen
\[
g(u_1,v_1),g(u_1,v_2),\cdots,g(u_1,v_f),g(u_2,v_1),\cdots,g(u_f,v_f)
\]
und die $f'^2+f''^2+\cdots$ Grössen, die aus den anderen Primfactoren $\Phi',\Phi'',\cdots$ von $\Theta$ in analoger Weise gebildet sind (\emph{Primfactoren}, § 10).

Daher sind die Wurzeln der charakteristischen Gleichung der Matrix $X\overline{A}$ gleich $u_1,u_2,\cdots,u_f,0,0,\cdots,0$, und folglich ist
\begin{equation}\tag{1.}
|E+uX\overline{A}|=(1+u_1u)\cdots(1+u_fu)=\Phi(\eps+ux).
\end{equation}
Die Elemente der Matrix $X\overline{A}=\overline{A}X$ sind
\begin{equation}\tag{2.}
\begin{aligned}
x_{P,Q}&=\sum_R x_{PR^{-1}}a_{Q^{-1}R}
       =\sum_R a_{R^{-1}P}x_{RQ^{-1}}\\
      &=\sum_R a_{Q^{-1}R^{-1}P}x_R
       =\sum_R a_{R^{-1}}x_{PRQ^{-1}}.
\end{aligned}
\end{equation}
Ersetzt man in der Matrix $X=(x_{PQ^{-1}})$ $P$ durch $P^{-1}$ und $Q$ durch $Q^{-1}$, so erhält man die Matrix $(x_{P^{-1}Q})$. Diese geht also aus jener hervor, indem man die Zeilen und die Spalten in gleicher Weise unter einander vertauscht. Daher sind beide Matrizen einander ähnlich. Dasselbe Resultat ergiebt sich aus der leicht zu beweisenden Identität
\[
(y_{PQ})(x_{PQ^{-1}})=(x_{P^{-1}Q})(y_{PQ}).
\]
Vertauscht man dann in $(x_{P^{-1}Q})$ die Zeilen mit den Spalten, so erhält man die Matrix $\overline{X}=(x_{Q^{-1}P})$. Für die Matrix $\overline{X}$ ist daher die Gesammtheit aller Unterdeterminanten $f^{\ten}$ Grades dieselbe wie für $X$. Mithin ist auch der Rang von $\overline{A}$ gleich $f$. Nach \emph{Primfactoren}, § 11 verhalten sich daher die Unterdeterminanten $f^{\ten}$ Grades der Matrix $X\overline{A}=\overline{A}X$ wie die entsprechenden Unterdeterminanten $f^{\ten}$ Grades der Matrix $\overline{A}$. Nun ist nach (1.) die Summe aller Unterdeterminanten $f^{\ten}$ Grades von $X\overline{A}$ gleich $\Phi(x)$. Mithin ist jede Unterdeterminante $f^{\ten}$ Grades von $X\overline{A}$ gleich der entsprechenden von $\overline{A}$, multiplicirt mit $\Phi(x)$,
\begin{equation}\tag{3.}
|x_{P,Q}|=\Phi(x)|a_{Q^{-1}P}|
\qquad(P=A_1,\cdots,A_f;\\ Q=B_1,\cdots,B_f).
\end{equation}
In der Matrix $f^{\ten}$ Grades
\[
(x_{P,Q})_f\qquad(P=A_1,\cdots,A_f;\ Q=B_1,\cdots,B_f)
\]
ist $x_E$ mit der Matrix $(a_{Q^{-1}P})_f$ multiplicirt. Ist deren Determinante von Null verschieden, so ist, wie ich jetzt zeigen will,
\begin{equation}\tag{4.}
(x_{P,Q})_f(a_{Q^{-1}P})_f^{-1}
\qquad(P=A_1,\cdots,A_f;\ Q=B_1,\cdots,B_f)
\end{equation}
eine zur Gruppe $\Hgr$ gehörige Matrix.

Setzt man $E-A=B$, so ist der Rang von $B$ gleich $h-f=g$, und es ist
\[
\sum a_Pa_Q=a_R,
\qquad
\sum b_Pb_Q=b_R,
\qquad
\sum a_Pb_Q=\sum b_Pa_Q=0\qquad(PQ=R).
\]
Die beiden Determinanten $f^{\ten}$ und $g^{\ten}$ Grades
\begin{equation}\tag{5.}
|a_{Q^{-1}P}|\\[-0.5em]
\qquad(P=A_1,\cdots,A_f;\ Q=B_1,\cdots,B_f)
\end{equation}
und
\[
|b_{Q^{-1}P}|
\qquad(P=C_1,\cdots,C_g;\ Q=D_1,\cdots,D_g)
\]
seien von Null verschieden. Sei $M$ die Matrix $h^{\ten}$ Grades, deren Zeilen man aus der Zeile
\[
a_{B_1^{-1}R},\cdots,a_{B_f^{-1}R},b_{D_1^{-1}R},\cdots,b_{D_g^{-1}R}
\]
erhält, indem man für $R$ die $h$ Elemente von $\Hgr$ setzt. Sei $L'$ die Matrix, deren Zeilen man in derselben Weise aus
\[
a_{R^{-1}A_1},\cdots,a_{R^{-1}A_f},b_{R^{-1}C_1},\cdots,b_{R^{-1}C_g}
\]
erhält, und sei $L$ die zu $L'$ conjugirte Matrix.

Ich bilde nun die Matrix $LXM$. Seien $\alpha$ und $\beta$ zwei der Indices $1,2,\cdots,f$, und $\gamma$ und $\delta$ zwei der Indices $1,2,\cdots,g$. Dann ist in jener Matrix das $\beta^{\te}$ Element der $\alpha^{\ten}$ Zeile
\[
\sum_{R,S}a_{R^{-1}A_\alpha}x_{RS^{-1}}a_{B_\beta^{-1}S}
=\sum_{R,S}a_{R^{-1}P}x_{RS^{-1}}a_{Q^{-1}S},
\]
falls man $A_\alpha=P$, $B_\beta=Q$ setzt. Dies ist ein Element der Matrix $\overline{A}X\overline{A}=X\overline{A}\,\overline{A}=X\overline{A}$, also gleich $x_{P,Q}$. Das $(f+\delta)^{\te}$ Element der $\alpha^{\ten}$ Zeile ist, falls man $A_\alpha=P$, $D_\delta=Q$ setzt,
\[
\sum_{R,S}a_{R^{-1}P}x_{RS^{-1}}b_{Q^{-1}S}.
\]
Dies ist ein Element der Matrix $\overline{A}X\overline{B}=X\overline{A}\,\overline{B}=0$, verschwindet also. Ebenso verschwindet das $\beta^{\te}$ Element der $(f+\gamma)^{\ten}$ Zeile. Endlich ist das $(f+\delta)^{\te}$ Element der $(f+\gamma)^{\ten}$ Zeile, falls man $C_\gamma=P$, $D_\delta=Q$ setzt,
\[
\sum_{R,S}b_{R^{-1}P}x_{RS^{-1}}b_{Q^{-1}S}.
\]
Dies ist ein Element der Matrix $\overline{B}X\overline{B}=X\overline{B}\,\overline{B}=X\overline{B}=X-X\overline{A}$, also gleich $x_{PQ^{-1}}-x_{P,Q}$.

Die Elemente der Matrix $LXM$ sind lineare Functionen der $h$ Variabelen $x_R$. Sie kann daher als eine lineare Verbindung von $h$ constanten Matrizen aufgefasst werden, deren jede mit einer der $h$ Variabelen $x_R$ multiplicirt ist. Speciell ist $x_E$ mit der Matrix
\[
\begin{pmatrix}
a_{B_\beta^{-1}A_\alpha}&0\\
0&b_{D_\delta^{-1}C_\gamma}
\end{pmatrix}=N
\]
multiplicirt, deren Determinante nicht verschwindet. Die Matrix $LXMN^{-1}$ zerfällt in eine Matrix $f^{\ten}$ Grades
\[
(x_{P,Q})(a_{Q^{-1}P})^{-1}
\qquad(P=A_1,\cdots,A_f;\ Q=B_1,\cdots,B_f)
\]
und eine Matrix $g^{\ten}$ Grades
\[
(x_{PQ^{-1}}-x_{P,Q})(\eps_{PQ^{-1}}-a_{Q^{-1}P})^{-1}
\qquad(P=C_1,\cdots,C_g;\ Q=D_1,\cdots,D_g).
\]
Giebt man in dieser Darstellung den Variabelen $x_R$ den Werth $\eps_R$, setzt man also $X=E$, so erhält man $LMN^{-1}=E$. Daher sind die Determinanten von $L$ und $M$ von Null verschieden, und es ist $MN^{-1}=L^{-1}$.

Nun ist aber $LXL^{-1}$ eine zur Gruppe $\Hgr$ gehörige Matrix. Da sie in zwei Matrizen $f^{\ten}$ und $g^{\ten}$ Grades zerfällt, so hat jeder der beiden Theile dieselbe Eigenschaft. Demnach ist (4.) eine zur Gruppe $\Hgr$ gehörige Matrix, deren Determinante gleich $\Phi(x)$ ist.

Ist also $\Phi$ einer der $k$ Primfactoren der Determinante der Gruppe $\Hgr$, so braucht man nur irgend ein jene Function determinirendes Werthsystem $a_R$ zu kennen, um die der Primfunction $\Phi$ entsprechende primitive Darstellung der Gruppe $\Hgr$ durch lineare Substitutionen angeben zu können.
% ===== END WORK 059 =====

\clearpage

% ===== BEGIN WORK 060: 060_ber_die_Charaktere_der_symmetrischen_Grup =====
\frobtitle{Über die Charaktere der symmetrischen Gruppe}

Die sämmtlichen Permutationen von $n$ Symbolen $1,2,\cdots,n$ sind die Elemente einer Gruppe $\Hgr$ der Ordnung $h=n!$, die man die symmetrische nennt. Eine Permutation kann stets und nur in einer Weise als Product von cyklischen Vertauschungen dargestellt werden, von denen nicht zwei ein Symbol gemeinsam haben. Damit zwei Permutationen in Bezug auf $\Hgr$ conjugirt seien, ist nothwendig und hinreichend, dass sie aus gleich vielen Cyklen derselben Ordnung zusammengesetzt sind. Bestehen die Permutationen der $\rho^{\ten}$ Classe $(\rho=0,1,\cdots,k-1)$ aus $\alpha$ Cyklen der Ordnung $1$, $\beta$ Cyklen der Ordnung $2$, $\gamma$ Cyklen der Ordnung $3$, $\cdots$, so ist ihre Anzahl nach \textsc{Cauchy} gleich
\begin{equation}\tag{1.}
 h_\rho=\frac{n!}{1^\alpha\alpha!\,2^\beta\beta!\,3^\gamma\gamma!\cdots}.
\end{equation}
Die Anzahl $k$ der Classen conjugirter Elemente, worin die $h$ Permutationen zerfallen, ist gleich der Anzahl der Lösungen der Gleichung
\begin{equation}\tag{2.}
 n=\alpha+2\beta+3\gamma+\cdots
\end{equation}
durch positive ganze Zahlen $\alpha,\beta,\gamma,\cdots$, die auch Null sein können, oder auch gleich der Anzahl der verschiedenen Zerlegungen der Zahl
\begin{equation}\tag{3.}
 n=n_1+n_2+\cdots
\end{equation}
in eine unbestimmte Anzahl positiver Theile.

Um die $k$ Charaktere der Gruppe $\Hgr$
\[
 \chi_\rho^{(\kappa)}\qquad(\rho=0,1,\cdots,k-1;\ \kappa=0,1,\cdots,k-1)
\]
zu berechnen, benutze ich $k$ bestimmte Untergruppen von $\Hgr$. Nachdem die Zahl $n$ in die Summanden (3.) zerlegt ist, bilde man alle Permutationen, die nur die ersten $n_1$ Symbole unter sich vertauschen, ebenso nur die folgenden $n_2$ Symbole, u. s. w. Diese Permutationen bilden eine Gruppe $\Ggr$ der Ordnung
\begin{equation}\tag{4.}
 g=n_1!\,n_2!\cdots.
\end{equation}
Jede Permutation $R$ von $\Ggr$ zerfällt in eine Permutation $R_1$ unter den ersten $n_1$ Symbolen, eine Permutation $R_2$ unter den folgenden $n_2$ Symbolen, u. s. w. Besteht $R_\nu$ aus $\alpha_\nu$ Cyklen der Ordnung $1$, $\beta_\nu$ Cyklen der Ordnung $2$, $\cdots$, so ist
\begin{equation}\tag{5.}
 n_1=\alpha_1+2\beta_1+\cdots,\\qquad n_2=\alpha_2+2\beta_2+\cdots,\quad\cdots,
\end{equation}
und wenn $R$ in $\Hgr$ der $\rho^{\ten}$ Classe angehört,
\begin{equation}\tag{6.}
 \alpha=\alpha_1+\alpha_2+\cdots,\qquad \beta=\beta_1+\beta_2+\cdots.
\end{equation}
Die Anzahl der in $\Ggr$ enthaltenen Permutationen dieser Classe ist daher
\[
 g_\rho=\sum \frac{n_1!}{1^{\alpha_1}\alpha_1!\,2^{\beta_1}\beta_1!\cdots}
           \frac{n_2!}{1^{\alpha_2}\alpha_2!\,2^{\beta_2}\beta_2!\cdots}\cdots,
\]
die Summe erstreckt über alle Lösungen der Gleichungen (5.) und (6.). Haben diese Gleichungen keine Lösung, so ist $g_\rho=0$. Sonst ist, da
\begin{equation}\tag{7.}
 \frac{h}{h_\rho}=1^\alpha\alpha!\,2^\beta\beta!\,3^\gamma\gamma!\cdots
\end{equation}
ist,
\[
 \frac{h g_\rho}{g h_\rho}=\sum
 \frac{\alpha!}{\alpha_1!\alpha_2!\cdots}
 \frac{\beta!}{\beta_1!\beta_2!\cdots}\cdots.
\]
Seien $x_1,x_2,\cdots,x_m$ unabhängige Variabele, deren Anzahl $m$ nicht kleiner ist als die Anzahl der Theile $n_1,n_2,\cdots$ von $n$. Dann ist diese Summe, ausgedehnt über alle Lösungen der Gleichungen (5.) und (6.), gleich dem Coefficienten von $x_1^{n_1}x_2^{n_2}\cdots$ in der Entwicklung des Productes
\begin{equation}\tag{8.}
 (x_1+x_2+\cdots+x_m)^\alpha
 (x_1^2+x_2^2+\cdots+x_m^2)^\beta
 (x_1^3+x_2^3+\cdots+x_m^3)^\gamma\cdots
\end{equation}
nach Potenzen der Veränderlichen.

Dieses für die folgende Untersuchung grundlegende Ergebniss lässt sich auch ohne jede Rechnung so einsehen: Sei $R=C_1C_2\cdots C_s$ eine Permutation der $\rho^{\ten}$ Classe, die aus $s$ Cyklen von je $c_1,c_2,\cdots,c_s$ Symbolen besteht, so dass
\begin{equation}\tag{9.}
 n=c_1+c_2+\cdots+c_s
\end{equation}
ist. Nach den Entwicklungen in § 1 meiner Arbeit \emph{Über Relationen zwischen den Charakteren einer Gruppe und denen ihrer Untergruppen}, Sitzungsberichte 1898, die ich im Folgenden mit $U.$ citiren werde, giebt die Zahl $\frac{h g_\rho}{g h_\rho}$ an, wie viele der $\frac{h}{g}$ Untergruppen von $\Hgr$, die mit $\Ggr$ conjugirt sind, die Permutation $R$ enthalten. Man findet diese Gruppen, indem man die $n$ Symbole auf alle möglichen Arten in Systeme von je $n_1,n_2,\cdots$ Symbolen theilt, und die Symbole jedes dieser Systeme in jeder Weise unter sich vertauscht. Die Anzahl der so erhaltenen Gruppen ist
\[
 \frac{h}{g}=\frac{n!}{n_1!n_2!\cdots},
\]
sie brauchen aber nicht alle unter einander verschieden zu sein. Damit $R$ in einer dieser Gruppen vorkomme, muss es möglich sein, aus einigen der Cyklen $C_1,C_2,\cdots,C_s$ eine Permutation von $n_1$ Symbolen zu bilden, aus einigen anderen eine solche von $n_2$ Symbolen, u. s. w. Dies ist auf so viele Weisen ausführbar, als man die Zahlen $n_1,n_2,\cdots$ als Summen der Zahlen $c_1,c_2,\cdots,c_s$ darstellen kann. Diese Anzahl ist aber gleich dem Coefficienten von $x_1^{n_1}x_2^{n_2}\cdots$ in der Entwicklung des Productes
\begin{equation}\tag{10.}
 (x_1^{c_1}+x_2^{c_1}+\cdots+x_m^{c_1})(x_1^{c_2}+x_2^{c_2}+\cdots+x_m^{c_2})\cdots
 (x_1^{c_s}+x_2^{c_s}+\cdots+x_m^{c_s}),
\end{equation}
das, abgesehen von der gewählten Bezeichnung, mit dem Ausdrucke (8.) übereinstimmt. Sind unter den Zahlen $n_1,n_2,\cdots$ oder unter den Zahlen $c_1,c_2,\cdots$ mehrere einander gleiche, so sind zwei Zerlegungen auch dann als verschieden anzusehen, wenn sie sich nur durch die Bezeichnung der darin vorkommenden Zahlen unterscheiden. Denn wenn auch $c_1=c_2$ ist, so ist doch $C_1$ von $C_2$ verschieden.

Nach der Formel $U.$ § 1, (9.) lässt sich $\frac{h g_\rho}{g h_\rho}$ als eine lineare Verbindung der Charaktere von $\Hgr$ mit ganzzahligen Coefficienten darstellen (oder anders ausgedrückt, ist das System der $k$ Zahlen $\frac{h g_\rho}{g h_\rho}$ durch den Modul theilbar, der von den $k^2$ Werthen der $k$ Charaktere gebildet wird). Jede solche Verbindung will ich einen zusammengesetzten Charakter von $\Hgr$ nennen. In der Entwicklung der ganzen symmetrischen Function $n^{\ten}$ Grades (8.) nach Potenzen von $x_1,x_2,\cdots,x_m$ ist daher jeder Coefficient ein zusammengesetzter Charakter von $\Hgr$, d. h. der Werth eines solchen Charakters für die durch die Exponenten $\alpha,\beta,\gamma\cdots$ bestimmte Classe $(\rho)$.

\section*{§ 2.}

Das Product aller Differenzen der $m$ Variabelen $x_1,x_2,\cdots,x_m$ sei
\[
 \Delta(x_1,x_2,\cdots,x_m)=(x_2-x_1)(x_3-x_1)(x_3-x_2)\cdots(x_m-x_{m-1}).
\]
In der Entwicklung dieser Function vom Grade $\frac12m(m-1)$ hat jedes Glied die Gestalt
\[
 \pm x_1^{k_1}x_2^{k_2}\cdots x_m^{k_m},
\]
wo $k_1,k_2,\cdots,k_m$ die Zahlen $0,1,\cdots,m-1$ in irgend einer Reihenfolge sind. Je nachdem diese Permutation eine positive oder eine negative ist, ist der Coefficient des Gliedes gleich $+1$ oder $-1$. Auch wenn $k_1,k_2,\cdots,k_m$ irgend $m$ positive oder negative Zahlen sind, benutze ich im Folgenden das Zeichen
\begin{equation}\tag{1.}
 [k_1,k_2,\cdots,k_m]=\sign.\,\Delta(k_1,k_2,\cdots,k_m).
\end{equation}
Es hat den Werth $0$, wenn zwei der $m$ Zahlen einander gleich sind, sonst den Werth $+1$ oder $-1$, je nachdem unter den Differenzen $k_2-k_1,k_3-k_1,k_3-k_2,\cdots$ eine gerade oder ungerade Anzahl negativ sind, je nachdem in der Folge $k_1,k_2,\cdots,k_m$ eine gerade oder ungerade Anzahl von Inversionen vorkommt.

Dann ist in dem Producte
\begin{equation}\tag{2.}
 (x_1+x_2+\cdots+x_m)^\alpha(x_1^2+x_2^2+\cdots+x_m^2)^\beta(x_1^3+x_2^3+\cdots+x_m^3)^\gamma\cdots\Delta(x_1,x_2,\cdots,x_m)
\end{equation}
jeder Coefficient eine ganzzahlige lineare Verbindung gewisser Coefficienten des Ausdrucks (8.), § 1 und mithin ebenfalls ein zusammengesetzter Charakter von $\Hgr$. Wie ich darauf gekommen bin, die symmetrische Function (8.), § 1 durch Multiplication mit dem Differenzenproducte in eine alternirende zu verwandeln, geht aus dem folgenden Beweise deutlich hervor. In dieser alternirenden Function bezeichne ich den Coefficienten von
\[
 x_1^{\lambda_1}x_2^{\lambda_2}\cdots x_m^{\lambda_m}\quad\text{mit}\quad [\lambda_1,\lambda_2,\cdots,\lambda_m]\chi_\rho^{(\lambda)}.
\]
Das Symbol $(\lambda)$ bezeichnet das System der $m$ Exponenten $\lambda_1,\lambda_2,\cdots,\lambda_m$, deren Summe
\begin{equation}\tag{3.}
 \lambda_1+\lambda_2+\cdots+\lambda_m=n+\frac12m(m-1)
\end{equation}
ist, und das Symbol $(\rho)$ die durch die Zahlen $\alpha,\beta,\gamma\cdots$ der Zerlegung (2.), § 1 bestimmte Classe conjugirter Elemente von $\Hgr$. Ebenso wie die Classen denke ich mir die verschiedenen Lösungen der Gleichung (3.), wobei es auf die Reihenfolge der Summanden nicht ankommt, ganz willkürlich numerirt und mit $(0),(1),\cdots$ allgemein mit $(\lambda)$ bezeichnet.

Für ein bestimmtes $\lambda$ bilden, wie ich jetzt zeigen werde, die $k$ Zahlen
\[
 \chi_\rho^{(\lambda)}\qquad(\rho=0,1,\cdots,k-1)
\]
einen (einfachen) Charakter von $\Hgr$, für verschiedene $\lambda$ verschiedene Charaktere. Ist $m=n$, so giebt es gerade $k$ Systeme von je $n$ verschiedenen Zahlen, die der Bedingung
\begin{equation}\tag{4.}
 \kappa_1+\kappa_2+\cdots+\kappa_n=\frac12 n(n+1)
\end{equation}
genügen. Denn ordnet man in einer Lösung die Summanden so, dass
\begin{equation}\tag{5.}
 \kappa_1<\kappa_2<\cdots<\kappa_n
\end{equation}
ist, so entspricht ihr eine Zerlegung der Zahl $n$ in die Theile
\[
 \kappa_1\leqq \kappa_2-1\leqq \kappa_3-2\leqq\cdots\leqq\kappa_n-n+1
\]
und umgekehrt. Mithin liefert die alternirende Function
\begin{equation}\tag{6.}
\begin{aligned}
 &(x_1+x_2+\cdots+x_n)^\alpha(x_1^2+x_2^2+\cdots+x_n^2)^\beta(x_1^3+x_2^3+\cdots+x_n^3)^\gamma\cdots\Delta(x_1,x_2,\cdots,x_n)\\
 &=\sum_{(\kappa)}[\kappa_1,\kappa_2,\cdots,\kappa_n]\chi_\rho^{(\kappa)}x_1^{\kappa_1}x_2^{\kappa_2}\cdots x_n^{\kappa_n}
\end{aligned}
\end{equation}
die Werthe der $k$ Charaktere der symmetrischen Gruppe $\Hgr$. Man erhält sie schon alle, indem man in dieser Summe von $n!k$ Gliedern nur die $k$ betrachtet, worin die Exponenten der Bedingung (5.) genügen.

Zwischen den Charakteren $\psi_\rho^{(\kappa)}$ einer Gruppe $\Hgr$ bestehen die Relationen
\begin{equation}\tag{7.}
 \sum_\rho h_\rho\psi_\rho^{(\kappa)}\psi_\rho^{(\kappa')}=h,
 \qquad
 \sum_\rho h_\rho\psi_\rho^{(\kappa)}\psi_\rho^{(\lambda)}=0,
\end{equation}
falls $(\lambda)$ von dem zu $(\kappa)$ inversen (conjugirten complexen) Charakter $(\kappa')$ verschieden ist. Nun sei
\[
 \chi_\rho=\sum_\kappa r_\kappa\psi_\rho^{(\kappa)}
\]
ein zusammengesetzter Charakter, seien also die Coefficienten $r_\kappa$ positive oder negative ganze Zahlen. Dann ist, weil $\psi_\rho^{(\kappa')}=\psi_\rho^{(\kappa)}$ ist,
\[
 \sum_\rho h_\rho\chi_\rho\chi_\rho'=h\sum_\kappa r_\kappa^2.
\]
Kann man also zeigen, dass diese Summe gleich $h$ ist, so muss eine der Zahlen $r_\kappa$ gleich $\pm1$, die $k-1$ anderen gleich $0$ sein. Ist überdies $\chi_0$ eine positive ganze Zahl, falls $(0)$ die Hauptclasse ist, so folgt daraus, dass $\chi_\rho$ ein (einfacher) Charakter ist. Kennt man mehrere solche Charaktere $\chi_\rho^{(\kappa)},\chi_\rho^{(\lambda)},\cdots$ und ist
\[
 \sum_\rho h_\rho\chi_\rho^{(\kappa)}\chi_\rho^{(\lambda)}=0,
\]
so sind sie von einander verschieden.

Für die symmetrische Gruppe vereinfachen sich diese Formeln noch dadurch, dass die in der Formel (6.) auftretenden Charaktere $\chi_\rho^{(\kappa)}$ nur reelle Werthe haben, also den inversen Charakteren gleich sind.

\section*{§ 3.}

Um die aufgestellte Behauptung zu beweisen, führe ich ein zweites System von $m$ unabhängigen Variabelen $y_1,y_2,\cdots,y_m$ ein und bilde die Summe
\begin{equation}\tag{1.}
\begin{aligned}
 &\sum_\rho \frac{h_\rho}{h}(x_1+\cdots+x_m)^\alpha(x_1^2+\cdots+x_m^2)^\beta\cdots
 (y_1+\cdots+y_m)^\alpha(y_1^2+\cdots+y_m^2)^\beta\cdots\\
 &=\sum \frac{1}{1^\alpha\alpha!}(x_1+\cdots+x_m)^\alpha(y_1+\cdots+y_m)^\alpha
 \frac{1}{2^\beta\beta!}(x_1^2+\cdots+x_m^2)^\beta(y_1^2+\cdots+y_m^2)^\beta\cdots .
\end{aligned}
\end{equation}
Beschränkt man die Exponenten $\alpha,\beta,\gamma,\cdots$ zunächst nicht durch die Bedingung
\begin{equation}\tag{2.}
 \alpha+2\beta+3\gamma+\cdots=n,
\end{equation}
sondern ertheilt jedem die Werthe von $0$ bis $\infty$, so ist die Summe dieser Reihe
\[
 e^{(x_1+\cdots+x_m)(y_1+\cdots+y_m)+\frac12(x_1^2+\cdots+x_m^2)(y_1^2+\cdots+y_m^2)+\cdots}.
\]
Führt man in jedem Gliede die Multiplication aus, so wird der Exponent gleich
\[
 -l(1-x_1y_1)-l(1-x_1y_2)-l(1-x_2y_1)-\cdots-l(1-x_my_m).
\]
Folglich ist die Summe gleich dem reciproken Werthe des Productes
\[
 (1-x_1y_1)(1-x_1y_2)(1-x_2y_1)\cdots(1-x_my_m).
\]
Multiplicirt man daher den betrachteten Ausdruck noch mit
\[
 \Delta(x_1,x_2,\cdots,x_m)\Delta(y_1,y_2,\cdots,y_m),
\]
so wird er nach einer bekannten Formel von \textsc{Cauchy} gleich der Determinante $m^{\ten}$ Grades
\[
 \left|\frac{1}{1-x_\mu y_\nu}\right|\qquad(\mu,\nu=1,2,\cdots,m),
\]
oder ausgerechnet
\[
 \sum[\mu_1,\mu_2,\cdots,
\mu_m]\frac{1}{(1-x_{\mu_1}y_1)(1-x_{\mu_2}y_2)\cdots(1-x_{\mu_m}y_m)}.
\]
Nun ist
\[
 \frac{1}{1-xy}=\sum_{\lambda=0}^\infty x^\lambda y^\lambda,
\]
und folglich ist die Determinante gleich
\[
 \sum_\lambda\sum_\mu[\mu_1,\mu_2,\cdots,
\mu_m]x_{\mu_1}^{\lambda_1}y_1^{\lambda_1}x_{\mu_2}^{\lambda_2}y_2^{\lambda_2}\cdots x_{\mu_m}^{\lambda_m}y_m^{\lambda_m},
\]
oder wenn man in jedem Gliede die Factoren in geeigneter Weise umstellt,
\begin{equation}\tag{3.}
 \sum_{\kappa,\lambda}[\kappa_1,\kappa_2,\cdots,
\kappa_m][\lambda_1,\lambda_2,\cdots,
\lambda_m]x_1^{\kappa_1}x_2^{\kappa_2}\cdots x_m^{\kappa_m}y_1^{\lambda_1}y_2^{\lambda_2}\cdots y_m^{\lambda_m}.
\end{equation}
Jeder der $m$ Exponenten $\lambda_1,\lambda_2,\cdots,\lambda_m$ durchläuft die Werthe von $0$ bis $\infty$. Sind zwei dieser Exponenten einander gleich, so hat das Glied den Coefficienten $0$. Für jedes System von $m$ verschiedenen Exponenten durchlaufen $\kappa_1,\kappa_2,\cdots,\kappa_m$ die $m!$ Permutationen von $\lambda_1,\lambda_2,\cdots,
\lambda_m$.

Nimmt man in der Summe (1.) nur die Glieder, die der Bedingung (2.) genügen, so hat man sich in der Reihe (3.) auf die Glieder zu beschränken, worin
\begin{equation}\tag{4.}
 \lambda_1+\lambda_2+\cdots+\lambda_m=n+\frac12m(m-1)
\end{equation}
ist. Daher ist diese endliche Reihe nach (6.), § 2 gleich
\[
 \sum_{\kappa,\lambda}\left(\sum_\rho\frac{h_\rho}{h}\chi_\rho^{(\kappa)}\chi_\rho^{(\lambda)}\right)
 [\kappa_1,\cdots,
\kappa_m][\lambda_1,\cdots,
\lambda_m]x_1^{\kappa_1}\cdots x_m^{\kappa_m}y_1^{\lambda_1}\cdots y_m^{\lambda_m},
\]
wo über alle Werthsysteme $\kappa,\lambda$ zu summiren ist, die der Bedingung (4.) und der analogen Bedingung
\[
 \kappa_1+\kappa_2+\cdots+\kappa_m=n+\frac12m(m-1)
\]
genügen. Wenn sich also $\kappa_1,\kappa_2,\cdots,
\kappa_m$ von $\lambda_1,\lambda_2,\cdots,
\lambda_m$ nicht durch die Reihenfolge allein unterscheiden, so ist
\[
 \sum_\rho h_\rho\chi_\rho^{(\kappa)}\chi_\rho^{(\lambda)}=0,
\]
im anderen Falle aber ist die Summe gleich $h$. Damit ist der wesentliche Theil des Beweises erledigt, und es ist nur noch zu zeigen, dass die Zahlen $\chi_0^{(\kappa)}=f^{(\kappa)}$ positiv sind.

Für die Hauptclasse $(0)$ ist $\alpha=n,\beta=0,\gamma=0,\cdots$, und daher ist
\begin{equation}\tag{5.}
 (x_1+x_2+\cdots+x_m)^n\Delta(x_1,x_2,\cdots,x_m)=\sum[\lambda_1,\cdots,
\lambda_m]f^{(\lambda)}x_1^{\lambda_1}\cdots x_m^{\lambda_m}.
\end{equation}
Die linke Seite ist gleich
\[
 \left(\sum \frac{n!}{\mu_1!\cdots\mu_m!}x_1^{\mu_1}\cdots x_m^{\mu_m}\right)
 \left(\sum[\kappa_1,\cdots,
\kappa_m]x_1^{\kappa_1}\cdots x_m^{\kappa_m}\right),
\]
wo $n=\mu_1+\cdots+
\mu_m$ ist, und $\kappa_1,
\kappa_2,
\cdots,
\kappa_m$ eine der $m!$ Permutationen der Zahlen $0,1,
\cdots,m-1$ bedeutet. Die Function $\frac{1}{\mu!}=\frac{1}{\Pi(\mu)}$ ist Null, wenn $\mu$ negativ ist. Das Product der beiden Summen ist gleich $n!$ mal
\[
 \sum[\kappa_1,
\cdots,
\kappa_m]\frac{1}{(\lambda_1-
\kappa_1)!}\cdots\frac{1}{(\lambda_m-
\kappa_m)!}x_1^{\lambda_1}\cdots x_m^{\lambda_m}.
\]
In dieser Summe ist der Coefficient von $x_1^{\lambda_1}\cdots x_m^{\lambda_m}$ gleich der Determinante $m^{\ten}$ Grades
\[
 \frac{f^{(\lambda)}}{h}=\left|\frac{1}{(\lambda_\mu-
u+1)!}\right|\qquad(\mu,
\nu=1,2,
\cdots,m).
\]
Multiplicirt man die Elemente der $\mu^{\ten}$ Zeile mit $\lambda_\mu!$, und dividirt man die der $\nu^{\ten}$ Spalte durch $(\nu-1)!$, so erhält man
\[
 \frac{\lambda_1!\cdots\lambda_m!}{0!\cdots(m-1)!}\frac{f^{(\lambda)}}{h}=\left|\binom{\lambda_\mu}{\nu-1}\right|.
\]
Das Zeichen
\[
 \binom{x}{n}=\frac{x(x-1)\cdots(x-n+1)}{1\cdot2\cdots n}
\]
brauche ich auch in dem Falle, wo $x$ eine Variabele ist. Die Zeilen der letzten Determinante sind die Werthe der Functionen
\[
 \binom{x}{0},\binom{x}{1},\cdots,\binom{x}{m-1}\quad\text{für}\quad x=\lambda_1,\lambda_2,\cdots,
\lambda_m.
\]
Die Grade dieser Functionen sind $0,1,
\cdots,m-1$. Sind also $\lambda_1,
\lambda_2,
\cdots,
\lambda_m$ unabhängige Variabele, so ist diese Determinante gleich $\Delta(\lambda_1,
\lambda_2,
\cdots,
\lambda_m)$ bis auf einen constanten Divisor $\Delta(0,1,
\cdots,m-1)$, den man erhält, indem man $\lambda_1=0,
\lambda_2=1,
\cdots,
\lambda_m=m-1$ setzt. Da
\[
 \Delta(0,1,
\cdots,m-1)=0!1!\cdots(m-1)!
\]
ist, so ist
\begin{equation}\tag{6.}
 f^{(\lambda)}=\frac{n!\Delta(\lambda_1,
\lambda_2,
\cdots,
\lambda_m)}{\lambda_1!\lambda_2!\cdots\lambda_m!},
\end{equation}
also positiv, wenn $\lambda_1<\lambda_2<\cdots<\lambda_m$ ist.

\section*{§ 4.}

Um die $k$ verschiedenen Charaktere zu erhalten, müssen wir $m=n$ setzen. Dann ist $\chi^{(\kappa)}$ durch ein System von $n$ Zahlen
\begin{equation}\tag{1.}
 (\kappa)=\kappa_1,
\kappa_2,
\cdots,
\kappa_n
\end{equation}
charakterisirt, die den Bedingungen
\begin{equation}\tag{2.}
 \kappa_1+
\kappa_2+\cdots+
\kappa_n=\frac12n(n+1),\qquad 0\leqq\kappa_1<\kappa_2<\cdots<\kappa_n
\end{equation}
genügen. Was die Beziehung zwischen den beiden Entwicklungen (2.) und (6.), § 2 betrifft, so ist
\[
 [\lambda_1,
\lambda_2,
\cdots,
\lambda_m]\chi_\rho^{(\lambda)}=[\kappa_1,
\kappa_2,
\cdots,
\kappa_n]\chi_\rho^{(\kappa)},
\]
wenn $(\kappa)$ in dem Falle $m<n$ das System der $n$ Zahlen
\begin{equation}\tag{3.}
 \kappa_1,
\kappa_2,
\cdots,
\kappa_n=0,1,
\cdots,
\mu-1,
\mu+
\lambda_1,
\cdots,
\mu+
\lambda_m\qquad(m+
\mu=n)
\end{equation}
ist. Umgekehrt ist in dem Falle $m>n$
\begin{equation}\tag{4.}
 \lambda_1,
\lambda_2,
\cdots,
\lambda_m=0,1,
\cdots,
\mu-1,
\mu+
\kappa_1,
\cdots,
\mu+
\kappa_n\qquad(n+
\mu=m).
\end{equation}
Um also von $m$ unabhängig zu sein, könnte man die von Null verschiedenen unter den Zahlen
\[
 \lambda_1\leqq\lambda_2-1\leqq\lambda_3-2\leqq\cdots\leqq\lambda_m-m+1
\]
zur Definition des Charakters $\chi$ verwenden. Es erweist sich aber als praktischer, dazu die folgende Charakteristik zu benutzen.

Die $n$ Zahlen (1.) sind alle $\leqq2n-1$, denn den Maximalwerth $2n-1$ erreicht $\kappa_n$, wenn die $n-1$ anderen Zahlen die Minimalwerthe $\kappa_1=0,
\kappa_2=1,
\cdots,
\kappa_{n-1}=n-2$ haben. Seien $a_1,
 a_2,
\cdots,
 a_n$ die Zahlen $0,1,
\cdots,n-1$ in irgend einer Reihenfolge. Dann sind
\[
 n-1-a_1,
\qquad n-1-a_2,
\cdots,
 n-1-a_n
\]
dieselben Zahlen. Von diesen mögen $n-r$ unter den Zahlen (1.) vorkommen, etwa $n-1-a_{r+1},
\cdots,n-1-a_n$. Die übrigen $r$ der Zahlen (1.) sind $\geqq n$, und mögen mit $n+b_1,
\cdots,n+b_r$ bezeichnet werden. Dann nenne ich
\begin{equation}\tag{5.}
 (\kappa)=\left(\begin{matrix}a_1&a_2&\cdots&a_r\\ b_1&b_2&\cdots&b_r\end{matrix}\right)
\end{equation}
die Charakteristik von $\chi^{(\kappa)}$. Damit sie völlig bestimmt sei, möge
\begin{equation}\tag{6.}
 a_1<a_2<\cdots<a_r,\qquad b_1<b_2<\cdots<b_r,
\end{equation}
sein. Die Gleichung (2.) lautet jetzt
\[
 (n-1-a_{r+1})+\cdots+(n-1-a_n)+(n+b_1)+\cdots+(n+b_r)=\frac12n(n+1).
\]
Ferner ist
\[
 a_1+\cdots+a_r+a_{r+1}+\cdots+a_n=0+1+\cdots+(n-1)=\frac12n(n-1),
\]
und mithin
\begin{equation}\tag{7.}
 a_1+a_2+\cdots+a_r+b_1+b_2+\cdots+b_r=n-r.
\end{equation}
Die linke Seite ist $\geqq2(0+1+\cdots+(r-1))=r(r-1)$, und folglich ist die Zahl $r$, die ich als Rang der Charakteristik (5.) bezeichne,
\begin{equation}\tag{8.}
 r\leqq\sqrt n.
\end{equation}
Der Grad des Charakters $\chi^{(\kappa)}$ wird nun
\begin{equation}\tag{9.}
 f^{(\kappa)}=\frac{n!\Delta(\kappa_1,
\kappa_2,
\cdots,
\kappa_n)}{\kappa_1!\kappa_2!\cdots\kappa_n!}
 =\frac{n!\Delta(a_1,
\cdots,a_r)\Delta(b_1,
\cdots,b_r)}{a_1!\cdots a_r!\,b_1!\cdots b_r!\prod(a_\alpha+b_\beta+1)},
\end{equation}
wo in dem letzten Producte $\alpha$ und $\beta$ die Werthe $1,2,\cdots,r$ durchlaufen.

Bei der Durchführung dieser Umrechnung wollen wir, da $f^{(\kappa)}$ positiv ist, die Vorzeichen ausser Acht lassen. Dann ist zu zeigen, dass
\[
\begin{aligned}
 &a_1!\cdots a_r!\,b_1!\cdots b_r!\prod(a_\alpha+b_\beta+1)
 \Delta(-a_{r+1}-1,\cdots,-a_n-1,b_1,\cdots,b_r)\\
 &=\pm (n-a_{r+1}-1)!\cdots(n-a_n-1)!(n+b_1)!\cdots(n+b_r)!
 \Delta(a_1,\cdots,a_r)\Delta(b_1,\cdots,b_r)
\end{aligned}
\]
ist. Nun ist
\[
 \frac{(b_\beta+n)!}{b_\beta!}=\prod_\alpha(a_\alpha+b_\beta+1)\prod_\gamma(a_\gamma+b_\beta+1),
\]
wo $\gamma$ die Werthe $r+1,
\cdots,n$ durchläuft, und
\[
 \pm\Delta(-a_{r+1}-1,
\cdots,-a_n-1,b_1,
\cdots,b_r)=\Delta(a_{r+1},
\cdots,a_n)\Delta(b_1,
\cdots,b_r)
 \prod_{\gamma,
\beta}(a_\gamma+b_\beta+1).
\]
Daher ist noch zu beweisen, dass
\[
 a_1!\cdots a_r!\Delta(a_{r+1},\cdots,a_n)=(n-1-a_{r+1})!\cdots(n-1-a_n)!\Delta(a_1,
\cdots,a_r)
\]
ist. Nun ist (falls in den Producten der Factor $a_\alpha-a_\alpha$ weggelassen wird)
\[
 (a_\alpha-a_1)(a_\alpha-a_2)\cdots(a_\alpha-a_n)=(a_\alpha-0)(a_\alpha-1)\cdots(a_\alpha-n+1)=\pm a_\alpha!(n-1-a_\alpha)!.
\]
Wenn man also die obige Gleichung mit $(n-1-a_1)!\cdots(n-1-a_r)!$ multiplicirt, so erhält man
\[
 \Delta(a_1,\cdots,a_r)^2\prod(a_\alpha-a_\gamma)\Delta(a_{r+1},\cdots,a_n)
 =\pm0!1!\cdots(n-1)!\Delta(a_1,\cdots,a_r).
\]
Diese Formel ergiebt sich aber aus der Gleichung
\[
\begin{aligned}
0!1!\cdots(n-1)!&=\Delta(0,1,\cdots,n-1)\\
&=\pm\Delta(a_1,\cdots,a_r,a_{r+1},\cdots,a_n)\\
&=\pm\Delta(a_1,\cdots,a_r)\Delta(a_{r+1},\cdots,a_n)\prod(a_\gamma-a_\alpha).
\end{aligned}
\]
Der Ausdruck (9.) kann auf die Form einer Determinante $r^{\ten}$ Grades
\begin{equation}\tag{9.}
 f^{(\kappa)}=\frac{n!}{a_1!\cdots a_r!\cdots b_1!\cdots b_r!}
 \left|\frac{1}{a_\alpha+b_\beta+1}\right|\qquad(\alpha,
\beta=1,2,
\cdots,r)
\end{equation}
gebracht werden.

Die Charakteristik (5.) lässt sich auch aus den $m$ Zahlen ableiten, die oben mit $\lambda_1,
\lambda_2,
\cdots,
\lambda_m$ bezeichnet sind. $r$ derselben sind $\geqq m$, und gleich $m+b_1,m+b_2,
\cdots,m+b_r$, und von den Zahlen $0,1,
\cdots,m-1$ fehlen unter ihnen die $r$ Zahlen $m-1-a_1,m-1-a_2,m-1-a_r$. Um $\chi_\rho^{(\lambda)}$ mit Hülfe der Entwicklung des Ausdrucks (2.), § 2 zu berechnen, braucht man darin nur $m\geqq a_r+1$ zu wählen.

Den der Charakteristik (5.) entsprechenden Charakter werde ich auch mit
\[
 \chi_\rho^{(\kappa)}=\chi\left(\begin{matrix}a_1&a_2&\cdots&a_r\\ b_1&b_2&\cdots&b_r\end{matrix}\right)_\rho
\]
bezeichnen. Setzt man z. B. in der Formel (2.), § 2 $m=2$, so erhält man
\begin{equation}\tag{10.}
\begin{aligned}
 &(1-x)(1+x)^\alpha(1+x^2)^\beta(1+x^3)^\gamma\cdots\\
 &=1+x\chi\left(\begin{matrix}1\\ n-2\end{matrix}\right)
 +x^2\chi\left(\begin{matrix}0&1\\0&n-3\end{matrix}\right)
 +x^3\chi\left(\begin{matrix}0&1\\1&n-4\end{matrix}\right)+\cdots\\
 &\quad +x^\lambda\chi\left(\begin{matrix}0&1\\ \lambda-2&n-\lambda-1\end{matrix}\right)+\cdots,
\end{aligned}
\end{equation}
falls $\lambda<\frac12(n+1)$ ist. Ist $n+1$ gerade und $\lambda=\frac12(n+1)$, so ist der Coefficient von $x^\lambda$ gleich $0$. Dann kehren dieselben Coefficienten in umgekehrter Reihenfolge, aber mit entgegengesetztem Vorzeichen wieder. Daher ist z. B.
\begin{equation}\tag{11.}
\begin{aligned}
\chi\left(\begin{matrix}1\\ n-2\end{matrix}\right)&=\alpha-1,\qquad
\chi\left(\begin{matrix}0&1\\0&n-3\end{matrix}\right)=\frac12\alpha(\alpha-3)+\beta,\\
\chi\left(\begin{matrix}0&1\\1&n-4\end{matrix}\right)&=\frac16\alpha(\alpha-1)(\alpha-5)+(\alpha-1)\beta+
\gamma,\\
\chi\left(\begin{matrix}0&1\\2&n-5\end{matrix}\right)&=\frac1{24}\alpha(\alpha-1)(\alpha-2)(\alpha-7)
 +\frac12(\alpha-1)(\alpha-2)\beta+\frac12\beta(\beta-3)+(\alpha-1)\gamma+
\delta.
\end{aligned}
\end{equation}
Der letzte Ausdruck hat demnach für alle Lösungen der Gleichung $\alpha+2\beta+3\gamma+4\delta=7$ den Werth $0$.

\section*{§ 5.}

In besonders einfacher Weise lassen sich die Charaktere ersten Ranges berechnen. Dazu gehört der Charakter $\chi(R)=\alpha-1$, der ($U.$, § 5) jeder zweifach transitiven Gruppe von Substitutionen zukommt. Seine charakteristische Function ist
\begin{equation}\tag{1.}
 F(x)=(1-x)^{\alpha-1}(1-x^2)^\beta(1-x^3)^\gamma\cdots.
\end{equation}
Denn die Summe der Wurzeln der Gleichung $F(x)=0$ ist $\chi(R)=\alpha-1$, und die Summe ihrer $x^{\ten}$ Potenzen ist $\chi(R^x)$, nämlich gleich der Anzahl der Symbole, die $R^x$ ungeändert lässt, vermindert um $1$. Nach den Resultaten des § 2 meiner Arbeit \emph{Über die Composition der Charaktere einer Gruppe}, Sitzungsberichte 1899, ist der Coefficient von $(-x)^{n-1-
\lambda}$ in der Entwicklung von $F(x)$ nach Potenzen von $x$ eine lineare Verbindung der Charaktere von $\Hgr$ mit ganzen positiven Zahlencoefficienten. Um diese zu ermitteln, mache ich $F(x)$ homogen und bilde die Summe
\[
 \sum_\rho\frac{h_\rho}{h}(u-v)^\alpha(u^2-v^2)^\beta(u^3-v^3)^\gamma\cdots
 (x_1+
\cdots+x_n)^\alpha(x_1^2+
\cdots+x_n^2)^\beta(x_1^3+
\cdots+x_n^3)^\gamma\cdots,
\]
worin
\[
 \frac{h_\rho}{h}=\frac{1}{1^\alpha\alpha!\,2^\beta\beta!\,3^\gamma\gamma!\cdots}
\]
ist. Sieht man von der Summationsbeschränkung
\[
 n=\alpha+2\beta+3\gamma+
\cdots
\]
zunächst ab, so erhält man, wie in § 3, für die unendliche Reihe den Ausdruck
\[
 \left(\frac{(1-x_1v)\cdots(1-x_nv)}{(1-x_1u)\cdots(1-x_nu)}-1\right).
\]
Dabei ist, um Weiterungen zu vermeiden, das Werthsystem $\alpha=\beta=\gamma=\cdots=0$ bei der Summation ausgeschlossen worden. Dieser Ausdruck ist mit
\[
 \frac{\Delta(x_1,x_2,
\cdots,x_n)}{u-v}
\]
zu multipliciren, und dann ist in seiner Entwicklung nach steigenden Potenzen von $u$ und $v$ der Coefficient von $u^\lambda(-v)^{n-1-
\lambda}$ zu bestimmen. Ersetzt man $u$ und $v$ durch ihre reciproken Werthe, und setzt man
\[
 f(x)=(x-x_1)(x-x_2)\cdots(x-x_n),
\]
so ist in der Entwicklung des Ausdrucks
\[
 \frac{\Delta(x_1,x_2,
\cdots,x_n)}{v-u}\left(\frac{f(v)}{f(u)}\frac{v^n}{u^n}\right)
\]
nach fallenden Potenzen von $u$ und $v$ der Coefficient von
\[
 (-1)^{n-1-
\lambda}u^{-n-1-
\lambda}v^\lambda\qquad(\lambda=0,1,
\cdots,n-1)
\]
zu ermitteln. Nun ist
\[
 \frac1{v-u}\left(\frac{f(v)}{f(u)}-\frac{v^n}{u^n}\right)=\frac1{u^n}\frac{u^n-v^n}{u-v}+
 \sum_\nu\frac{f(v)}{f'(x_\nu)(v-x_\nu)(u-x_\nu)},
\]
und darin ist der Coefficient von $u^{-n-1-
\lambda}$, mit $\Delta(x_1,x_2,
\cdots,x_n)$ multiplicirt, gleich
\[
 \Delta(x_1,x_2,
\cdots,x_n)\sum_\nu\frac{f(v)x_\nu^{n+\lambda}}{f'(x_\nu)(v-x_\nu)}
 =\sum_\nu\Delta(x_1,
\cdots,x_{\nu-1},v,x_{\nu+1},
\cdots,x_n)x_\nu^{n+
\lambda}.
\]
In dieser Summe ist der Coefficient $(-1)^{n-1-
\lambda}v^\lambda x_1^{\kappa_1}x_2^{\kappa_2}\cdots x_n^{\kappa_n}$ zu berechnen, wo $\kappa_1,
\kappa_2,
\cdots,
\kappa_n$ den Bedingungen (2.), § 4 genügen. Dieser ist Null, ausser wenn $\kappa_1,
\kappa_2,
\cdots,
\kappa_{n-1}\leqq n-1$ sind, und $\kappa_n=n+
\lambda$. Daher hat diese Zerlegung den Rang $r=1$, und es ist
\[
 (\kappa)=\left(\begin{matrix}n-1-
\lambda\\ \lambda\end{matrix}\right).
\]
Endlich ist der Coefficient von
\[
 (-1)^{n-1-
\lambda}v^\lambda x_1^0x_2^1\cdots x_\lambda^{\lambda-1}x_{\lambda+1}^{\lambda+1}\cdots x_{n-1}^{n-1}x_n^{n+
\lambda}
\]
in $\Delta(x_1,
\cdots,x_{n-1},v)x_n^{n+
\lambda}$ gleich $+1$.

Daher ist
\begin{equation}\tag{2.}
\begin{aligned}
 F(x)&=\chi\left(\begin{matrix}0\\ n-1\end{matrix}\right)
 -x\chi\left(\begin{matrix}1\\ n-2\end{matrix}\right)
 +x^2\chi\left(\begin{matrix}2\\ n-3\end{matrix}\right)\cdots\\
 &\quad+(-x)^\lambda\chi\left(\begin{matrix}\lambda\\ n-1-
\lambda\end{matrix}\right)+\cdots+(-x)^{n-2}\chi\left(\begin{matrix}n-2\\1\end{matrix}\right)+(-x)^{n-1}\chi\left(\begin{matrix}n-1\\0\end{matrix}\right).
\end{aligned}
\end{equation}
Die Charaktere ersten Ranges sind also
\begin{equation}\tag{3.}
\begin{aligned}
\chi\left(\begin{matrix}0\\ n-1\end{matrix}\right)&=1,\qquad
\chi\left(\begin{matrix}1\\ n-2\end{matrix}\right)=\alpha-1,
\qquad
\chi\left(\begin{matrix}2\\ n-3\end{matrix}\right)=\binom{\alpha-1}{2}-\beta,\\
\chi\left(\begin{matrix}3\\ n-4\end{matrix}\right)&=\binom{\alpha-1}{3}-(\alpha-1)\beta+
\gamma,\\
\chi\left(\begin{matrix}4\\ n-5\end{matrix}\right)&=\binom{\alpha-1}{4}-\binom{\alpha-1}{2}\beta+(
\alpha-1)\gamma+
\binom{\beta}{2}-\delta,
\cdots.
\end{aligned}
\end{equation}
Ferner ist
\begin{equation}\tag{4.}
 \chi\left(\begin{matrix}n-1\\0\end{matrix}\right)=(-1)^{\beta+
\delta+\cdots}=(-1)^{n-s},
\end{equation}
wenn $s=\alpha+
\beta+
\gamma+
\delta+
\cdots$ die Anzahl der Cyklen von $R$ ist, und allgemein
\[
 \chi\left(\begin{matrix}b\\a\end{matrix}\right)_\rho=(-1)^{\beta+
\delta+
\cdots}\chi\left(\begin{matrix}a\\b\end{matrix}\right)_\rho.
\]

\section*{§ 6.}

Ist $\Hgr$ die symmetrische und $\Ggr$ die alternirende Gruppe des Grades $n$, und ist $T$ irgend eine ungerade Permutation, so ist $\Hgr=\Ggr+
\Ggr T$, wo der Complex $\Ggr$ die $\frac12n!$ geraden, und der Complex $\Ggr T=T\Ggr$ die $\frac12n!$ ungeraden Permutationen enthält. Diese beiden Complexe sind die Elemente der Gruppe $\frac{\Hgr}{\Ggr}$ der Ordnung $2$. Sie hat zwei Charaktere, $1,1$ und $1,-1$. Folglich hat die Gruppe $\Hgr$ einen Charakter
\begin{equation}\tag{1.}
 \chi_\rho^{(1)}=\chi\left(\begin{matrix}n-1\\0\end{matrix}\right)_\rho=(-1)^{\beta+
\delta+
\cdots}=(-1)^{n-s},
\end{equation}
der für die Elemente von $\Ggr$ den Werth $+1$, für die von $\Ggr T$ den Werth $-1$ hat. Durch Multiplication mit diesem Charakter ersten Grades ergiebt sich aus jedem Charakter $\chi^{(\kappa)}$ ein anderer $\chi^{(\lambda)}$. Zwei solche Charaktere, für die
\begin{equation}\tag{2.}
 \chi^{(\lambda)}=\chi^{(\kappa)}\chi^{(1)},\qquad
 \chi^{(\kappa)}=\chi^{(\lambda)}\chi^{(1)}
\end{equation}
ist, nenne ich associirte. Hat $\chi^{(\kappa)}(R)$ für jede ungerade Permutation $R$ den Werth $0$, so ist dieser Charakter sich selbst associirt. Sonst sind zwei associirte Charaktere von einander verschieden. Wir charakterisiren $\chi^{(\lambda)}$ durch ein System von $n$ Zahlen, die den Bedingungen
\begin{equation}\tag{3.}
 \frac12n(n+1)=\lambda_1+
\lambda_2+
\cdots+
\lambda_n,
\qquad 0\leqq\lambda_1<\lambda_2<\cdots<\lambda_n
\end{equation}
genügen und $\chi^{(\kappa)}$ durch eine analoge Zerlegung (2.), § 4, und wir nennen auch diese beiden Zerlegungen associirte. Ich will nun zeigen, dass dann die Zahlen
\begin{equation}\tag{4.}
 \kappa_1,
\kappa_2,
\cdots,
\kappa_n,
 2n-1-
\lambda_1,
 2n-1-
\lambda_2,
\cdots,
 2n-1-
\lambda_n
\end{equation}
abgesehen von der Reihenfolge mit den Zahlen $0,1,
\cdots,2n-1$ übereinstimmen. Zunächst sind auf diese Art die Zerlegungen (3.) paarweise einander zugeordnet. Denn $\kappa_1,
\kappa_2,
\cdots,
\kappa_n$ sind $n$ der $2n$ Zahlen $0,1,
\cdots,2n-1$. Bezeichnet man die $n$ übrigen mit $2n-1-
\lambda_1,2n-1-
\lambda_2,
\cdots,2n-1-
\lambda_n$, so ist
\[
 \sum(\kappa)+
\sum(2n-1-
\lambda)=0+1+
\cdots+2n-1=n(2n-1),
\]
und daher $\sum\kappa=
\sum\lambda$. Dass dann die Beziehungen (2.) bestehen, geht aus der Gleichung
\begin{equation}\tag{5.}
 \sum_\rho\frac{h_\rho}{h}\chi_\rho^{(1)}\chi_\rho^{(\kappa)}\chi_\rho^{(\lambda)}=1
\end{equation}
hervor. Um sie zu beweisen, bilde ich die unendliche Reihe
\[
\begin{aligned}
 &\sum (-1)^{\beta+
\delta+
\cdots}\frac1{\alpha!}(x_1+
\cdots+x_n)^\alpha(y_1+
\cdots+y_n)^\alpha
 \frac1{2^\beta\beta!}(x_1^2+
\cdots+x_n^2)^\beta(y_1^2+
\cdots+y_n^2)^\beta\cdots\\
 &=e^{(x_1+
\cdots+x_n)(y_1+
\cdots+y_n)-\frac12(x_1^2+
\cdots+x_n^2)(y_1^2+
\cdots+y_n^2)+
\cdots}\\
 &=\prod(1+x_\mu y_\nu).
\end{aligned}
\]
Dann ist die Summe (5.) gleich dem Coefficienten von $x_1^{\kappa_1}\cdots x_n^{\kappa_n}y_1^{\lambda_1}\cdots y_n^{\lambda_n}$ in der Entwicklung des Productes
\[
 \Delta(x_1,
\cdots,x_n)\Delta(y_1,
\cdots,y_n)\prod(1+x_\mu y_\nu).
\]
Ersetzt man $y_1,
\cdots,y_n$ durch $-\frac1{y_1},
\cdots,-\frac1{y_n}$, so ist dies der Coefficient von
\[
\pm x_1^{\kappa_1}\cdots x_n^{\kappa_n}
y_1^{2n-1-\lambda_1}\cdots y_n^{2n-1-\lambda_n}
\]
in der Entwicklung des Ausdrucks
\[
 \Delta(x_1,
\cdots,x_n)\Delta(y_1,
\cdots,y_n)\prod(y_\nu-x_\mu)=\Delta(x_1,
\cdots,x_n,y_1,
\cdots,y_n).
\]
Dieser ist aber gleich $0$, wenn nicht die $2n$ Zahlen (4.) alle unter einander verschieden sind, also abgesehen von der Reihenfolge mit den Zahlen $0,1,
\cdots,2n-1$ übereinstimmen. In diesem Falle ist der Coefficient $\pm1$. Auch ohne auf die genaue Bestimmung des Zeichens einzugehen, ergiebt sich daraus die Gleichung (5.), weil die Summe links nur einen der beiden Werthe $0$ oder $+1$ haben kann.

Die Beziehung zwischen zwei associirten Charakteren lässt sich bequemer darstellen, wenn man $\chi^{(\kappa)}$ durch die Charakteristik (5.), § 4 definirt. Seien $b_1,
\cdots,b_r,b_{r+1},
\cdots,b_n$ die Zahlen $0,1,
\cdots,n-1$. Sind dann $\kappa_1,
\cdots,
\kappa_n$ die Zahlen $n-1-a_{r+1},
\cdots,n-1-a_n,n+b_1,
\cdots,n+b_r$, so sind $2n-1-
\lambda_1,
\cdots,2n-1-
\lambda_n$ die Zahlen $n-1-a_1,
\cdots,n-1-a_r,n+b_{r+1},
\cdots,n+b_n$, und $\lambda_1,
\cdots,
\lambda_n$ selbst die Zahlen $n-1-b_{r+1},
\cdots,n-1-b_n,n+a_1,
\cdots,n+a_r$. Daher sind
\begin{equation}\tag{6.}
 (\kappa)=\left(\begin{matrix}a_1&\cdots&a_r\\ b_1&\cdots&b_r\end{matrix}\right),
 \qquad
 (\lambda)=\left(\begin{matrix}b_1&\cdots&b_r\\ a_1&\cdots&a_r\end{matrix}\right)
\end{equation}
associirte Charakteristiken, oder es ist
\begin{equation}\tag{7.}
 \chi\left(\begin{matrix}b_1&\cdots&b_r\\ a_1&\cdots&a_r\end{matrix}\right)_\rho
 =(-1)^{n-s}\chi\left(\begin{matrix}a_1&\cdots&a_r\\ b_1&\cdots&b_r\end{matrix}\right)_\rho.
\end{equation}
Eine Charakteristik ist folglich stets und nur dann sich selbst associirt, wenn $a_1=b_1,
\cdots,a_r=b_r$ ist. Dann ergiebt sich aber aus der Gleichung (7.), § 4
\begin{equation}\tag{8.}
 n=(2a_1+1)+
\cdots+(2a_r+1)
\end{equation}
eine Darstellung von $n$ als Summe von lauter verschiedenen ungeraden Zahlen. Umgekehrt entspricht jeder solchen Darstellung eine sich selbst associirte Charakteristik
\begin{equation}\tag{9.}
 (\kappa)=\left(\begin{matrix}a_1&\cdots&a_r\\ a_1&\cdots&a_r\end{matrix}\right),
 \qquad
 a_1+a_2+
\cdots+a_r=\frac12(n-r).
\end{equation}
Die Anzahl $v$ der sich selbst associirten Charaktere von $\Hgr$ ist also gleich der Anzahl der Darstellungen von $n$ als Summe von lauter verschiedenen ungeraden Zahlen. Zufolge der bekannten Formel
\[
 \prod(1+x^{2\alpha-1})=1:\prod(1-(-1)^{\alpha-1}x^\alpha)
\]
ist diese Anzahl gleich der Differenz der Anzahl der Zerlegungen
\begin{equation}\tag{10.}
 n=c_1+c_2+c_3+
\cdots,
\end{equation}
wofür $\sum(c-1)$ gerade ist, und der Anzahl derer, wofür $\sum(c-1)$ ungerade ist, also gleich der Differenz zwischen der Anzahl der geraden und der ungeraden Classen. Die $k$ Classen zerfallen demnach in
\begin{equation}\tag{11.}
 \frac12(k+v)=u+v\quad\text{gerade und}\quad \frac12(k-v)=u\quad\text{ungerade Classen.}
\end{equation}
Jeder Darstellung (8.) entspricht eine Zerlegung
\[
 n=1+(2a_1+1)+
\cdots+(2a_{r-1}+1)+2a_r,
\]
wofür $\sum(c-1)$ ungerade ist. Daher ist $u>v$, nur für $n=3$ ist $u=v$.

Die hier entwickelten Sätze sind von Wichtigkeit bei der Berechnung der Charaktere der alternirenden Gruppe, die sich, wie ich nächstens zeigen werde, aus ihren Beziehungen zu den Charakteren von Untergruppen ebenfalls vollständig bestimmen lassen.

\section*{§ 7.}

Ich hebe noch einige specielle Fälle hervor, worin sich gewisse Werthe von Charakteren in einfacher Weise bestimmen lassen. Die Substitutionen der $\rho^{\ten}$ Classe mögen aus $s$ Cyklen von je $c_1,c_2,
\cdots,c_s$ Symbolen bestehen. Dann ist
\begin{equation}\tag{1.}
 n=c_1+c_2+
\cdots+c_s,
\end{equation}
und es sei
\begin{equation}\tag{2.}
 c_1\leqq c_2\leqq\cdots\leqq c_s.
\end{equation}
Nun ist $[\kappa_1,
\kappa_2,
\cdots,
\kappa_n]\chi_\rho^{(\kappa)}$ der Coefficient von $x_1^{\kappa_1}x_2^{\kappa_2}\cdots x_n^{\kappa_n}$ in der Entwicklung des Productes
\begin{equation}\tag{3.}
 (x_1^{c_1}+
\cdots+x_n^{c_1})\cdots(x_1^{c_s}+
\cdots+x_n^{c_s})\Delta(x_1,
\cdots,x_n).
\end{equation}
Ist die Charakteristik von $\chi^{(\kappa)}$
\begin{equation}\tag{4.}
 (\kappa)=\left(\begin{matrix}a_1&a_2&\cdots&a_r\\ b_1&b_2&\cdots&b_r\end{matrix}\right),
\end{equation}
so ist
\[
 \kappa_1,
\kappa_2,
\cdots,
\kappa_n=n-1-a_{r+1},
\cdots,n-1-a_n,n+b_1,
\cdots,n+b_r.
\]
Irgend ein Glied in der Entwicklung von $\Delta(x_1,
\cdots,x_n)$ ist $\pm x_1^{k_1}x_2^{k_2}\cdots x_n^{k_n}$, wo $k_1,k_2,
\cdots,k_n$ die Zahlen $0,1,
\cdots,n-1$ sind. Irgend ein Glied in der Entwicklung des Productes
\begin{equation}\tag{5.}
 (x_1^{c_1}+
\cdots+x_n^{c_1})\cdots(x_1^{c_s}+
\cdots+x_n^{c_s})
\end{equation}
ist $x_{\sigma_1}^{l_1}x_{\tau_1}^{l_2}\cdots x_{\omega}^{l_{s'}}$, wo $l_1=c_{\sigma_1}+c_{\sigma_2}+\cdots$, $l_2=c_{\tau_1}+c_{\tau_2}+
\cdots$, ist und
\[
 \sigma_1,
\sigma_2,
\cdots\qquad \tau_1,
\tau_2,
\cdots
\]
zusammen die Zahlen $1,2,
\cdots,s$ sind. Folglich ist $s'\leqq s$ und nur dann $s'=s$, wenn $l_1,l_2,
\cdots,l_s$ abgesehen von der Reihenfolge mit $c_1,c_2,
\cdots,c_s$ übereinstimmen. Multiplicirt man diese beiden Glieder, so werden $s'$ der Exponenten $k_1,k_2,
\cdots,k_n$ vermehrt. Erhält man so das Glied $x_1^{\kappa_1}x_2^{\kappa_2}\cdots x_n^{\kappa_n}$, so werden $r$ von den $n$ Exponenten $k_1,k_2,
\cdots,k_n$, die $\leqq n-1$ sind, durch diese Vermehrung grösser als $n-1$, nämlich gleich
\[
 n+b_1,
\cdots,n+b_r.
\]
Daher muss $s'\geqq r$ sein. Dies ist nicht möglich, wenn $s<r$ ist, und folglich ist in diesem Falle $\chi_\rho^{(\kappa)}=0$.

\emph{Ist die Anzahl $s$ der Cyklen, aus denen die Substitutionen der $\rho^{\ten}$ Classe bestehen, kleiner als der Rang $r$ der Charakteristik $(\kappa)$, so ist $\chi_\rho^{(\kappa)}=0$.}

Sei jetzt $s=r$. Dann müssen genau $s$ der Exponenten $k_1,k_2,
\cdots,k_n$ vermehrt werden, damit sie gleich $n+b_1,
\cdots,n+b_s$ werden, und die übrigen $n-s$ Exponenten müssen mit $n-1-a_{s+1},
\cdots,n-1-a_n$ übereinstimmen. Ferner muss $s'=s$ sein und $l_1,l_2,
\cdots,l_s=c_1,c_2,
\cdots,c_s$. Die $s$ Exponenten $n-1-a_1,
\cdots,n-1-a_s$ werden jeder um eine der Zahlen $c_1,
\cdots,c_s$ vermehrt, und so ergeben sich die Zahlen $n+b_1,
\cdots,n+b_s$. Daher ist
\begin{equation}\tag{6.}
 c_1-1=a_\alpha+b_\lambda,
\qquad c_2-1=a_\beta+b_\mu,
\quad\cdots\quad c_s-1=a_\vartheta+b_\tau,
\end{equation}
wo $\alpha,
\beta,
\cdots,
\vartheta$ und $\lambda,
\mu,
\cdots,
\tau$ zwei Permutationen der Zahlen $1,2,
\cdots,s$ sind. In einer anderen Reihenfolge mögen diese $s$ Gleichungen lauten
\[
 c_\xi-1=a_\gamma+b_1,
\qquad c_\eta-1=a_\delta+b_2,
\quad\cdots\quad c_\zeta-1=a_\nu+b_s.
\]
Zu dem Gliede
\[
 [n+b_1,
\cdots,n+b_s,n-1-a_{s+1},
\cdots,n-1-a_n]\chi_\rho^{(\kappa)}x_1^{n+b_1}\cdots x_s^{n+b_s}x_{s+1}^{n-1-a_{s+1}}\cdots x_n^{n-1-a_n}
\]
von (3.) erhält man also einen Beitrag, indem man das Glied
\[
 x_1^{n-1-a_\gamma}\cdots x_s^{n-1-a_\nu}x_{s+1}^{n-1-a_{s+1}}\cdots x_n^{n-1-a_n}
\]
von $\Delta(x_1,
\cdots,x_n)$ mit dem Gliede $x_1^{c_\xi}\cdots x_s^{c_\zeta}$ von (5.) multiplicirt. Daher ist
\[
 \chi_\rho^{(\kappa)}=[-a_{s+1}-1,
\cdots,-a_n-1,b_1,
\cdots,b_s]
 \sum[-a_{s+1},
\cdots,-a_n,-a_\gamma,
\cdots,-a_\nu]
\]
oder, wenn man wieder die Reihenfolge (6.) herstellt,
\[
 \chi_\rho^{(\kappa)}=\sum[-a_{s+1}-1,
\cdots,-a_n-1,b_\lambda,
\cdots,b_\tau][-a_{s+1},
\cdots,-a_n,-a_\alpha,
\cdots,-a_\vartheta],
\]
die Summe erstreckt über alle Lösungen der Gleichungen (6.). Das Product der beiden Vorzeichen ist gleich
\[
 \sign.(a_{s+1}-a_\alpha)\cdots(a_n-a_\alpha)\cdots(a_{s+1}-a_\vartheta)\cdots(a_n-a_\vartheta)
 [b_\lambda,
\cdots,b_\tau][a_\alpha,
\cdots,a_\vartheta](-1)^{\frac12s(s-1)}.
\]
Das letztere Vorzeichen ist, da sich $\alpha,
\beta,
\cdots,
\vartheta$ von $1,2,
\cdots,s$ nur durch die Reihenfolge unterscheiden, gleich
\[
 \sign.(a_{s+1}-a_1)
\cdots(a_n-a_1)
\cdots(a_{s+1}-a_s)
\cdots(a_n-a_s).
\]
Da $a_1<a_2<\cdots<a_s$ ist, und $a_1,
\cdots,a_s,a_{s+1},
\cdots,a_n=0,1,
\cdots,n-1$ sind, so kommen unter den Zahlen $a_{s+1},
\cdots,a_n$ die Zahlen $0,1,
\cdots,a_1-1$ vor, und demnach sind unter den $n-s$ Differenzen $a_{s+1}-a_1,
\cdots,a_n-a_1$ genau $a_1$ negative. Ebenso kommen unter den Zahlen $a_{s+1},
\cdots,a_n$ die Zahlen $0,1,
\cdots,a_2-1$ vor mit Ausnahme von $a_1$, und demnach sind unter den $n-s$ Differenzen $a_{s+1}-a_2,
\cdots,a_n-a_2$ genau $a_2-1$ negative.

Daher ist das letzte Vorzeichen gleich
\[
 (-1)^{a_1+(a_2-1)+(a_3-2)+\cdots+(a_s-s+1)}.
\]
Da endlich $b_1<b_2<\cdots<b_s$ ist, so ist $[b_\lambda,b_\mu,
\cdots,b_\tau]=[\lambda,
\mu,
\cdots,
\tau]$, und folglich ist
\begin{equation}\tag{7.}
 \chi_\rho^{(\kappa)}=(-1)^{a_1+a_2+
\cdots+a_s}\sum[\alpha,
\beta,
\cdots,
\vartheta][\lambda,
\mu,
\cdots,
\tau],
\end{equation}
die Summe erstreckt über alle Lösungen der Gleichungen (6.). Ein Glied dieser Summe hat den Werth $+1$ oder $-1$, je nachdem $\lambda,
\mu,
\cdots,
\tau$ eine positive oder negative Permutation von $\alpha,
\beta,
\cdots,
\vartheta$ ist.

Ist z. B.
\begin{equation}\tag{8.}
 c_1-1=a_1+b_1,
\qquad c_2-1=a_2+b_2,
\quad\cdots\quad c_s-1=a_s+b_s,
\end{equation}
so ist damit eine Lösung der Gleichungen (6.) gegeben. Eine andere können sie nicht haben. Denn ist $a_1+b_1=c_1-1=a_\alpha+b_\lambda$, so muss, weil $a_\alpha\geqq a_1$ und $b_\lambda\geqq b_1$ ist, $a_\alpha=a_1$ und $b_\lambda=b_1$ sein und mithin $\alpha=1$ und $\lambda=1$, u. s. w. Daher ist unter der Bedingung (8.)
\begin{equation}\tag{9.}
 \chi_\rho^{(\kappa)}=(-1)^{a_1+a_2+
\cdots+a_s}.
\end{equation}
Jeder Charakter $\chi^{(\kappa)}$ hat also für mindestens eine Classe den Werth $\pm1$. Ist noch specieller $a_1=b_1,
\cdots,a_s=b_s$, also noch (8.)
\begin{equation}\tag{10.}
 c_1=2a_1+1,
\qquad c_2=2a_2+1,
\quad\cdots\quad c_s=2a_s+1,
\end{equation}
so ist
\begin{equation}\tag{11.}
 \chi_\rho^{(\kappa)}=(-1)^{\frac12(n-s)}=(-1)^{\frac12(p-1)},
\end{equation}
wo
\begin{equation}\tag{12.}
 p=c_1c_2\cdots c_s
\end{equation}
ist. Denn dieses Product ist
\[
 (1+2a_1)\cdots(1+2a_s)\equiv 1+2(a_1+
\cdots+a_s)\pmod{4}.
\]
Ist aber $r=s$ und $a_1=b_1,
\cdots,a_s=b_s$, ohne dass die Bedingung (10.) erfüllt ist, so ist
\[
 \chi_\rho^{(\kappa)}=(-1)^{\frac12(n-s)}\sum[\alpha,
\beta,
\cdots,
\vartheta][\lambda,
\mu,
\cdots,
\tau],
\]
die Summe ausgedehnt über alle Lösungen der Gleichungen
\[
 c_1-1=a_\alpha+a_\lambda,
\qquad c_2-1=a_\beta+a_\mu,
\quad\cdots\quad c_s-1=a_\vartheta+a_\tau.
\]
Diesen Gleichungen genügt man nicht, indem man
\[
 \alpha=\lambda,
\qquad \beta=\mu,
\quad\cdots\quad \vartheta=\tau
\]
setzt, weil sonst nach (2.) $\alpha=
\lambda=1$, $\beta=
\mu=2$, $\cdots$, $\vartheta=
\tau=s$ wäre. Daher entspricht jeder Lösung eine andere
\[
 c_1-1=a_\lambda+a_\alpha,
\qquad c_2-1=a_\mu+a_\beta,
\quad\cdots\quad c_s-1=a_\tau+a_\vartheta,
\]
die von ihr verschieden ist, und folglich ist $\chi_\rho^{(\kappa)}$ eine gerade Zahl. Auch für alle anderen Classen ist, wie ich zeigen werde, $\chi_\rho^{(\kappa)}$ gerade, mit Ausnahme der einen Classe (10.).

Um noch ein anderes Beispiel zu behandeln, will ich $\chi_\rho^{(\kappa)}$ für den Fall berechnen, wo die Permutationen der $\rho^{\ten}$ Classe aus einem Cyklus von $c$ Symbolen und aus $n-c$ Cyklen von je einem Symbole bestehen, und zwar nach der in § 3 zur Bestimmung von $f^{(\lambda)}$ benutzten Methode. Bedient man sich auch der dort gebrauchten Bezeichnungen, so ist
\[
\begin{aligned}
 &\sum[\lambda_1,
\cdots,
\lambda_n]\chi_\rho^{(\lambda)}x_1^{\lambda_1}\cdots x_n^{\lambda_n} \\
 &\quad=(x_1^c+
\cdots+x_n^c)(x_1+
\cdots+x_n)^{n-c}\Delta(x_1,
\cdots,x_n)\\
 &\quad=(x_1^c+
\cdots+x_n^c)
 \left(\sum\frac{(n-c)!}{\mu_1!\cdots\mu_n!}x_1^{\mu_1}\cdots x_n^{\mu_n}\right)\\
 &\qquad\cdot
 \left(\sum[\kappa_1,
\cdots,
\kappa_n]x_1^{\kappa_1}\cdots x_n^{\kappa_n}\right),
\end{aligned}
\]
wo $\mu_1+
\cdots+
\mu_n=n-c$ ist, und $\kappa_1,
\cdots,
\kappa_n$ die Zahlen $0,1,
\cdots,n-1$ sind. Ist $\lambda_1<\lambda_2<\cdots<\lambda_n$, so besteht $\chi_\rho^{(\lambda)}:(n-c)!$ aus $n$ Theilen. Den ersten erhält man, indem man von der Summe $x_1^c+
\cdots+x_n^c$ nur das erste Glied nimmt. Er ist gleich
\[
 \sum[\kappa_1,
\kappa_2,
\cdots,
\kappa_n]\frac1{(\lambda_1-c-
\kappa_1)!}\frac1{(\lambda_2-
\kappa_2)!}\cdots\frac1{(\lambda_n-
\kappa_n)!},
\]
also nach § 3 gleich
\[
 \frac{\Delta(\lambda_1-c,
\lambda_2,
\cdots,
\lambda_n)}{(\lambda_1-c)!\lambda_2!\cdots\lambda_n!}.
\]
Dividirt man durch
\[
 f^{(\lambda)}=\frac{n!\Delta(\lambda_1,
\lambda_2,
\cdots,
\lambda_n)}{\lambda_1!\lambda_2!\cdots\lambda_n!}
\]
und multiplicirt man mit
\[
 h_\rho=\frac{n!}{c(n-c)!},
\]
so wird
\[
 -\frac{c^2h_\rho\chi_\rho^{(\lambda)}}{f^{(\lambda)}}
\]
eine Summe von $n$ Gliedern, deren erstes
\[
 \frac{(\lambda_1-c-
\lambda_1)(\lambda_1-c-
\lambda_2)\cdots(\lambda_1-c-
\lambda_n)\lambda_1(\lambda_1-1)\cdots(\lambda_1-c+1)}{(\lambda_1-
\lambda_2)\cdots(\lambda_1-
\lambda_n)}
\]
ist. Setzt man
\[
 \varphi(x)=(x-
\lambda_1)\cdots(x-
\lambda_n),
\quad
 \psi(x)=(x-c-
\lambda_1)\cdots(x-c-
\lambda_n)x(x-1)\cdots(x-c+1),
\]
so ist also
\[
 -\frac{c^2h_\rho\chi_\rho^{(\lambda)}}{f^{(\lambda)}}=
 \frac{\psi(\lambda_1)}{\varphi'(\lambda_1)}+
\cdots+
\frac{\psi(\lambda_n)}{\varphi'(\lambda_n)},
\]
oder gleich dem Coefficienten von $x^{-1}$ in der Entwicklung von $\frac{\psi(x)}{\varphi(x)}$ nach absteigenden Potenzen von $x$. Ist
\[
 (\lambda)=\left(\begin{matrix}a_1&a_2&\cdots&a_r\\ b_1&b_2&\cdots&b_r\end{matrix}\right),
\]
so erweitere man diesen Bruch mit
\[
 (x-(n-1-a_1))\cdots(x-(n-1-a_r))\,(x-c-(n-1-a_1))\cdots(x-c-(n-1-a_r)).
\]
Nun sind die Zahlen
\[
 \lambda_1,
\cdots,
\lambda_n,
 n-1-a_1,
\cdots,n-1-a_r
\]
gleich
\[
 0,1,
\cdots,n-1,
 n+b_1,
\cdots,n+b_r.
\]
Ersetzt man nun $x$ durch $x+n$, so wird der Bruch gleich
\[
 \frac{(x+a_1+1)\cdots(x+a_r+1)(x-c-b_1)\cdots(x-c-b_r)x(x-1)\cdots(x-c+1)}{(x-b_1)\cdots(x-b_r)(x-c+a_1+1)\cdots(x-c+a_r+1)}.
\]
Ist also
\begin{equation}\tag{13.}
 f(x)=\frac{(x-b_1)\cdots(x-b_r)}{(x+a_1+1)\cdots(x+a_r+1)},
\end{equation}
so ist
\begin{equation}\tag{14.}
\begin{aligned}
 -\frac{c^2h_\rho\chi_\rho^{(\lambda)}}{f^{(\lambda)}}
 &=\left[\frac{f(x-c)x(x-1)\cdots(x-c+1)}{f(x)}\right]_{x^{-1}}\\
 &=\left[\frac{\varphi(x-c)x(x-1)\cdots(x-c+1)}{\varphi(x)}\right]_{x^{-1}},
\end{aligned}
\end{equation}
d. h. gleich dem Coefficienten von $x^{-1}$ in der Entwicklung dieser rationalen Function nach absteigenden Potenzen von $x$, und diese Formel bleibt auch gültig, wenn man unter Anwendung der Bezeichnungen (3.) oder (4.), § 4
\begin{equation}\tag{15.}
 \varphi(x)=(x-
\lambda_1)(x-
\lambda_2)\cdots(x-
\lambda_m)
\end{equation}
setzt. Für $c=2$ ist daher
\begin{equation}\tag{16.}
 \frac{h_\rho\chi_\rho^{(\lambda)}}{f^{(\lambda)}}=\frac12\left(\sum b(b+1)-\sum a(a+1)\right).
\end{equation}
Dass dieser Ausdruck eine ganze Zahl ist, habe ich schon in § 2 meiner Arbeit \emph{Über Gruppencharaktere}, Sitzungsberichte 1896, bewiesen.
% ===== END WORK 060 =====

\clearpage

% ===== BEGIN WORK 061: 061_ber_die_Charaktere_der_alternirenden_Grup =====
\frobtitle{Über die Charaktere der alternirenden Gruppe}

Die Charaktere der alternirenden Gruppe der Permutationen von $n$ Symbolen lassen sich durch ein Verfahren bestimmen, ähnlich dem, das ich (Sitzungsberichte 1900) zur Berechnung der Charaktere der symmetrischen Gruppe benutzt habe.

Die alternirende Gruppe $\Hgr$ des Grades $n$ besteht aus den geraden oder positiven Permutationen der symmetrischen Gruppe $\Hgr'$. Eine Permutation, die aus $s$ Cyklen $C_1,C_2,\ldots,C_s$ von je $c_1,c_2,\ldots,c_s$ Symbolen besteht, ist gerade oder ungerade, je nachdem $\sum(c-1)=n-s$ gerade oder ungerade ist. Die Gruppen $\Hgr$ und $\Hgr'$ haben die Ordnungen $h=\frac12 n!$ und $h'=2h$, ihre Elemente mögen in $k$ und $k'$ Classen zerfallen. Auch in $\Hgr$ können zwei Permutationen nur dann conjugirt sein, wenn sie aus gleich vielen Cyklen derselben Ordnung bestehen. Es fragt sich aber, ob alle solche Permutationen auch in $\Hgr$ nur eine Classe bilden. Denn zwei gerade Permutationen $R$ und $S$ sind in $\Hgr$ stets und nur dann conjugirt, wenn es eine gerade Permutation $P$ giebt, die der Bedingung $P^{-1}RP=S$ genügt.

Sind $R$ und $S$ in $\Hgr'$ conjugirt, so ist diese Bedingung stets erfüllt, wenn $R$ mit einer negativen Permutation $T$ vertauschbar ist. Denn in $\Hgr'$ giebt es eine solche Permutation $Q$, dass $Q^{-1}RQ=S$ ist. Ferner ist $(TQ)^{-1}R(TQ)=Q^{-1}(T^{-1}RT)Q=Q^{-1}RQ=S$, und von den beiden Permutationen $Q$ und $TQ$ ist immer die eine gerade. In diesem Falle ist also die Classe von $\Hgr$ durch die Zahlen $c_1,c_2,\ldots,c_s$ vollständig bestimmt.

Ist aber $R$ mit keiner negativen Permutation vertauschbar, so können $R$ und $T^{-1}RT=S$ in $\Hgr$ nicht conjugirt sein, falls $T$ ungerade ist. Denn sonst gäbe es eine solche positive Permutation $P$, dass $P^{-1}SP=R$ wäre, und daher wäre die negative Permutation $TP$ mit $R$ vertauschbar.

Sind $P_1,P_2,\ldots,P_h$ die $h$ positiven Permutationen, so sind $TP_1,TP_2,\ldots,TP_h$ die $h$ negativen. Daher sind
\[
 P_\lambda^{-1}RP_\lambda,\qquad P_\lambda^{-1}T^{-1}RTP_\lambda\qquad(\lambda=1,2,\ldots,h)
\]
alle Permutationen, die mit $R$ in $\Hgr'$ conjugirt sind. Sie bilden in $\Hgr'$ eine, in $\Hgr$ aber zwei Classen, die durch $R$ und $T^{-1}RT=S$ repraesentirt werden. Zwei solche \emph{conjugirte Classen} von $\Hgr$ haben dieselbe Ordnung.

Damit aber $R$ nur mit positiven Permutationen vertauschbar sei, ist nothwendig und hinreichend, dass die Ordnungen der Cyklen von $R$ lauter verschiedene ungerade Zahlen sind. Denn jeder einzelne Cyklus $C$ von $R$ ist eine mit $R$ vertauschbare Permutation, und zwar eine negative, wenn seine Ordnung $c$ gerade ist. Ist aber $c$ ungerade, und sind zwei dieser Cyklen, etwa $(1,2,\ldots,c)$ und $(c+1,c+2,\ldots,2c)$ von gleicher Ordnung, so ist $R$ mit der Permutation $(1,c+1)(2,c+2)\cdots(c,2c)$ vertauschbar, die aus einer ungeraden Anzahl $c$ von Transpositionen besteht, also ungerade ist.

Wenn aber die Ordnungen $c_1,c_2,\ldots,c_s$ der Cyklen $C_1,C_2,\ldots,C_s$ von $R$ lauter verschiedene ungerade Zahlen sind, so ist $R$ nur mit positiven Permutationen vertauschbar. Diese bilden eine Gruppe vertauschbarer Elemente
\begin{equation}\tag{1.}
 C_1^{\gamma_1}C_2^{\gamma_2}\cdots C_s^{\gamma_s}\qquad(\gamma_1=0,1,\ldots,c_1-1,\ldots,\gamma_s=0,1,\ldots,c_s-1)
\end{equation}
der Ordnung
\begin{equation}\tag{2.}
 p=c_1c_2\cdots c_s
\end{equation}
und sind alle gerade, weil die Basiselemente $C_1,C_2,\ldots,C_s$ alle gerade sind.

In der symmetrischen Gruppe $\Hgr'$ sei $u$ die Anzahl der ungeraden, $u+v$ die der geraden Classen. Dann ist $v$ die Anzahl der Classen von $\Hgr'$, deren Permutationen aus lauter Cyklen verschiedener ungerader Ordnungen bestehen. In der Gruppe $\Hgr$ zerfällt jede dieser $v$ Classen von $\Hgr'$ in zwei conjugirte Classen, jede der $u$ anderen geraden Classen von $\Hgr'$ bleibt aber auch in $\Hgr$ eine Classe. Daher ist
\begin{equation}\tag{3.}
 k=u+2v,	\qquad k'=2u+v,
\end{equation}
und da $u>v$ ist, so ist auch $k'>k$.

\section*{§ 2.}

Mit Hülfe der Regeln, die ich in § 1 meiner Arbeit \emph{Über Relationen zwischen den Charakteren einer Gruppe und denen ihrer Untergruppen}, Sitzungsberichte 1898 (im Folgenden mit U. citirt), entwickelt habe, lassen sich die $k$ Charaktere von $\Hgr$ zum grössten Theile aus denen von $\Hgr'$ ableiten. Setzt man in einem Charakter $\chi(S)$ von $\Hgr'$ für $S$ nur die $h$ Elemente $R$ von $\Hgr$, so ist $\chi(R)$ eine lineare Verbindung von den $k$ Charakteren $\phi^{(\varkappa)}(R)$ der Gruppe $\Hgr$,
\[
 \chi(R)=\sum_{\varkappa} r_\varkappa\phi^{(\varkappa)}(R)
\]
mit positiven ganzzahligen Coefficienten $r_\varkappa$, und mithin ist
\[
 \sum_R \chi(R)\chi(R^{-1})=h\sum_\varkappa r_\varkappa^2.
\]
Für die symmetrische Gruppe sind die beiden (im Allgemeinen conjugirten complexen) Werthe $\chi(R)$ und $\chi(R^{-1})$ gleich. Durchläuft $S$ die $2h$ Elemente von $\Hgr'$, so ist
\[
 \sum_S \chi(S)^2=2h.
\]
Ist $T$ eine bestimmte negative Permutation, so durchlaufen $R$ und $RT$ zusammen die $2h$ Elemente $S$ von $\Hgr'$, und mithin ist
\[
 \sum_R \chi(R)^2+\sum_R \chi(RT)^2=2h.
\]
Der zu $\chi(S)$ associirte Charakter von $\Hgr'$ hat die Werthe $\omega(R)=\chi(R)$, $\omega(RT)=-\chi(RT)$. Ist nun $\chi(S)$ nicht sich selbst associirt, so sind $\chi(S)$ und $\omega(S)$ verschiedene Charaktere von $\Hgr'$, und folglich ist
\[
 \sum_S \chi(S)\omega(S)=0,\qquad \sum_R \chi(R)^2=\sum_R\chi(RT)^2,
\]
also
\[
 \sum_R\chi(R)\chi(R^{-1})=h,	\qquad \sum_\varkappa r_\varkappa^2=1.
\]
Daher ist von den $k$ positiven Zahlen $r_\varkappa$ eine gleich $1$, die anderen gleich $0$, und mithin ist $\chi(R)$ ein Charakter von $\Hgr$. So entspringt aus jedem der $u$ Paare associirter Charaktere von $\Hgr'$ ein Charakter von $\Hgr$, der die Bedingung
\begin{equation}\tag{1.}
 \chi(T^{-1}RT)=\chi(R)
\end{equation}
erfüllt, d. h. der sich selbst conjugirt ist.

Ist aber $\chi(S)$ sich selbst associirt, so ist $\chi(RT)=0$,
\[
 \sum_R \chi(R)^2=2h,	\qquad \sum_\varkappa r_\varkappa^2=2.
\]
Folglich sind zwei der Zahlen $r_\varkappa$ gleich $1$, die anderen gleich $0$, und $\chi(R)$ ist gleich der Summe von zwei verschiedenen Charakteren $\phi(R)$ und $\psi(R)$ von $\Hgr$. Nach U. § 2 sind dies, da $\Hgr$ eine invariante Untergruppe von $\Hgr'$ ist, zwei conjugirte Charaktere, also
\begin{equation}\tag{2.}
 \psi(R)=\phi(T^{-1}RT),\qquad \chi(R)=\phi(R)+\psi(R)=\phi(R)+\phi(T^{-1}RT).
\end{equation}
Gehört $R$ zu einer der $u$ nicht zerfallenden Classen von $\Hgr'$, so sind $R$ und $T^{-1}RT$ auch in $\Hgr$ conjugirt und mithin ist für ein solches Element $R$
\begin{equation}\tag{u.}
 \phi(R)=\psi(R)=\frac12\chi(R).
\end{equation}
Insbesondere gilt dies für $R=E$. Die Grade der beiden conjugirten Charaktere $\phi$ und $\psi$ sind also einander gleich und gleich dem halben Grade von $\chi$. Sind aber die Ordnungen der Cyklen von $R$ lauter verschiedene ungerade Zahlen, gehört also $R$ zu einer der $v$ zerfallenden Classen von $\Hgr'$, so können die Werthe von $\phi(R)$ und $\psi(R)$ verschieden sein, und zwar ist, wenn man $T^{-1}RT=S$ setzt,
\begin{equation}\tag{v.}
 \phi(R)=\psi(S),	\qquad \phi(S)=\psi(R).
\end{equation}
Aus jedem der $v$ sich selbst associirten Charaktere von $\Hgr'$ entspringen also zwei verschiedene, und zwar conjugirte Charaktere von $\Hgr$.

Die so erhaltenen $u+2v=k$ Charaktere von $\Hgr$ sind alle unter einander verschieden. Denn seien $\phi(R)=\chi(R)$ und $\overline{\phi}(R)=\overline{\chi}(R)$ zwei der ersten $u$ Charaktere von $\Hgr$, und sei, wenn $S$ die $2h$ Elemente von $\Hgr'$ durchläuft, $\overline{\chi}(S)$ sowohl von $\chi(S)$ wie von dem zu $\chi(S)$ associirten Charakter verschieden. Dann ist
\[
 \sum_R \chi(R)\overline{\chi}(R)+\sum_R \chi(RT)\overline{\chi}(RT)=0,
\]
und wenn man $\chi(R)$ durch den associirten Charakter ersetzt, auch
\[
 \sum_R \chi(R)\overline{\chi}(R)-\sum_R\chi(RT)\overline{\chi}(RT)=0
\]
und folglich
\[
 \sum_R \phi(R)\overline{\phi}(R^{-1})=0,
\]
also ist $\overline{\phi}(R)$ von $\phi(R)$ verschieden.

Ist aber $\overline{\phi}(R)=\overline{\chi}(R)$ einer der $u$ ersten, und $\phi(R)$ einer der $2v$ letzten Charaktere von $\Hgr$, der aus $\chi(R)=\phi(R)+\psi(R)$ entspringt, so ist $\chi(RT)=0$, also
\[
 \sum_S \overline{\chi}(S^{-1})\chi(S)=0,	\qquad \sum_R\overline{\phi}(R^{-1})(\phi(R)+\psi(R))=0.
\]
Nun sind $\phi$ und $\psi$ verschieden. Wäre daher $\overline{\phi}$ einem dieser beiden Charaktere gleich, so wäre diese Summe gleich $h$. Endlich erkennt man in derselben Weise, dass je zwei der letzten $2v$ Charaktere von $\Hgr$ verschieden sind.

Um also die $u+2v=k$ verschiedenen Charaktere von $\Hgr$ zu bestimmen, ist nur noch jeder der $v$ sich selbst associirten Charaktere $\chi$ von $\Hgr'$ in zwei conjugirte Charaktere $\phi$ und $\psi$ zu spalten, und zwar nur noch für die $v$ Permutationen $R$, die aus lauter Cyklen von verschiedenen ungeraden Ordnungen bestehen.

Die $k'=2u+v$ Charaktere von $\Hgr'$ bestehen aus $u$ Paaren associirter Charaktere und $v$ Charakteren, die sich selbst associirt sind. Die $k=u+2v$ Charaktere von $\Hgr$ bestehen aus $v$ Paaren conjugirter Charaktere und $u$ Charakteren, die sich selbst conjugirt sind.

\section*{§ 3.}

Ich verweise zunächst auf die Beispiele $n=4$ und $n=5$ in § 3 meiner Arbeit \emph{Über Gruppencharaktere} und $n=6$ in § 4 meiner Arbeit \emph{Über die Composition der Charaktere einer Gruppe}, Sitzungsberichte 1896 und 1899. In diesen Fällen ist $v=1$, ebenso für $n=7$. Ich gebe daher hier noch das Beispiel $n=8$, wo $v=2$ ist.

In der ersten Spalte der beifolgenden Tabelle ist die Classe einer Permutation durch die Formel $n=1\cdot\alpha+2\cdot\beta+\cdots$ gegeben, d. h. dadurch, dass sie aus $\alpha$ Cyklen der Ordnung $1$, $\beta$ Cyklen der Ordnung $2$, u. s. w. besteht. Kommt nur ein Cyklus einer bestimmten Ordnung vor, so ist der Multiplicator $\alpha,\beta,\cdots$ weggelassen. Erst kommen die $u=10$ ungeraden Classen, dann die $u+v=12$ geraden.

Von jedem der $u=10$ Paare associirter Charaktere ist nur der eine angegeben. Der andere entsteht daraus durch Multiplication mit dem Charakter ersten Grades $\chi\ch{7}{0}$, der dem Hauptcharakter $\chi\ch{0}{7}=1$ associirt ist. In den beiden letzten Spalten stehen die $v=2$ sich selbst associirten Charaktere.

Die Charaktere sind gegeben durch die Formeln (11.), § 4 und (3.), § 5 meiner Arbeit \emph{Über die Charaktere der symmetrischen Gruppe} und durch die folgenden:
\begin{align}\tag{1.}
 \chi\ch{0\quad 2}{0\quad n-4}&=\frac13\alpha(\alpha-2)(\alpha-4)-\gamma,\\
 \chi\ch{0\quad 2}{1\quad n-5}&=\frac18\alpha(\alpha-1)(\alpha-3)(\alpha-6)+\frac12(\alpha-1)(\alpha-2)\beta-\frac12\beta(\beta-1)-\delta,\\
 \chi\ch{1\quad 2}{0\quad n-5}&=\frac1{12}\alpha(\alpha-1)(\alpha-4)(\alpha-5)+\beta(\beta-2)-(\alpha-1)\gamma,\\
 \chi\ch{0\quad 3}{0\quad n-5}&=\frac18\alpha(\alpha-2)(\alpha-3)(\alpha-5)-\frac12\alpha(\alpha-3)\beta-\frac12\beta(\beta-1)+\delta,\\
 \chi\ch{1\quad 2}{1\quad n-6}&=\frac1{24}\alpha(\alpha-1)(\alpha-2)(\alpha-5)(\alpha-7)+\frac16\alpha(\alpha-1)(\alpha-5)\beta\notag\\
 &\quad+\frac12(\alpha-1)\beta(\beta-1)-\frac12(\alpha-1)(\alpha-2)\gamma+\beta\gamma-(\alpha-1)\delta.
\end{align}
In der alternirenden Gruppe zerfällt die Classe $(1+7)$ in zwei Classen, die ich durch ein angehängtes Zeichen $+$ oder $-$ unterscheide. Ebenso zerfällt der Charakter $\chi\ch{03}{03}$ in zwei. Dasselbe gilt von der Classe $(3+5)$ und dem ihr entsprechenden Charakter $\chi\ch{12}{12}$.

\begin{landscape}
\begin{center}
\small
\textbf{Symmetrische Gruppe.}\\
$h=8!=40320.$\\[0.5em]
\setlength{\tabcolsep}{3.5pt}
\renewcommand{\arraystretch}{1.15}
\begin{tabular}{c r|rrrrrrrrrrrr}
\toprule
$1\cdot\alpha+2\cdot\beta+\cdots$ & $h_\rho$ & $\ch{7}{0}$ & $\ch{1}{6}$ & $\ch{2}{5}$ & $\ch{3}{4}$ & $\ch{01}{05}$ & $\ch{01}{14}$ & $\ch{01}{23}$ & $\ch{02}{04}$ & $\ch{02}{13}$ & $\ch{12}{03}$ & $\ch{03}{03}$ & $\ch{12}{12}$\\
\midrule
$1\cdot6+2$ & 28 & $-1$ & 5 & 9 & 5 & 10 & 10 & 4 & 16 & 10 & 4 & 0 & 0\\
$1\cdot4+4$ & 420 & $-1$ & 3 & 3 & 1 & 2 & $-2$ & $-2$ & 0 & $-4$ & 0 & 0 & 0\\
$1\cdot3+2+3$ & 1120 & $-1$ & 2 & 0 & $-1$ & 1 & 1 & 1 & $-2$ & 1 & $-2$ & 0 & 0\\
$1\cdot2+2\cdot3$ & 420 & $-1$ & 1 & $-3$ & $-3$ & 2 & 2 & 0 & 0 & $-2$ & 4 & 0 & 0\\
$1\cdot2+6$ & 3360 & $-1$ & 1 & 0 & 0 & $-1$ & $-1$ & 0 & 0 & 1 & 1 & 0 & 0\\
$1+2+5$ & 4032 & $-1$ & 0 & $-1$ & 0 & 0 & 0 & $-1$ & 1 & 0 & $-1$ & 0 & 0\\
$1+3+4$ & 3360 & $-1$ & 0 & 0 & 1 & $-1$ & 1 & 1 & 0 & $-1$ & 0 & 0 & 0\\
$2\cdot2+4$ & 1260 & $-1$ & $-1$ & $-1$ & 1 & 2 & $-2$ & 2 & 0 & 0 & 0 & 0 & 0\\
$2+3\cdot2$ & 1120 & $-1$ & $-1$ & 0 & 2 & 1 & 1 & $-2$ & $-2$ & 1 & 1 & 0 & 0\\
$8$ & 5040 & $-1$ & $-1$ & 1 & $-1$ & 0 & 0 & 0 & 0 & 0 & 0 & 0 & 0\\
\midrule
$1\cdot8$ & 1 & 1 & 7 & 21 & 35 & 20 & 28 & 14 & 64 & 70 & 56 & 90 & 42\\
$1\cdot5+3$ & 112 & 1 & 4 & 6 & 5 & 5 & 1 & $-1$ & 4 & $-5$ & $-4$ & 0 & $-6$\\
$1\cdot4+2\cdot2$ & 210 & 1 & 3 & 1 & $-5$ & 4 & 4 & 2 & 0 & 2 & 0 & $-6$ & 2\\
$1\cdot3+5$ & 1344 & 1 & 2 & 1 & 0 & 0 & $-2$ & $-1$ & $-1$ & 0 & 1 & 0 & 2\\
$1\cdot2+2+4$ & 2520 & 1 & 1 & $-1$ & $-1$ & 0 & 0 & 0 & 0 & 0 & 0 & 2 & $-2$\\
$1\cdot2+3\cdot2$ & 1120 & 1 & 1 & 0 & 2 & $-1$ & 1 & 2 & $-2$ & 1 & $-1$ & 0 & 0\\
$1+2\cdot2+3$ & 1680 & 1 & 0 & $-2$ & 1 & 1 & 1 & $-1$ & 0 & $-1$ & 0 & 0 & 2\\
$2\cdot4$ & 105 & 1 & $-1$ & $-3$ & 3 & 4 & $-4$ & 6 & 0 & $-2$ & 8 & $-6$ & $-6$\\
$2+6$ & 3360 & 1 & $-1$ & 0 & 0 & 1 & $-1$ & 0 & 0 & 1 & $-1$ & 0 & 0\\
$4\cdot2$ & 1260 & 1 & $-1$ & 1 & $-1$ & 0 & 0 & 2 & 0 & $-2$ & 0 & 2 & 2\\
\midrule
$1+7$ & 5760 & 1 & 0 & 0 & 0 & $-1$ & 0 & 0 & 1 & 0 & 0 & $-1$ & 0\\
$3+5$ & 2688 & 1 & $-1$ & 1 & 0 & 0 & 1 & $-1$ & $-1$ & 0 & 1 & 0 & $-1$\\
\bottomrule
\end{tabular}
\end{center}
\end{landscape}

\begin{landscape}
\begin{center}
\small
\textbf{Alternirende Gruppe.}\\
$h=\frac12 8! =20160.$\\[0.5em]
\setlength{\tabcolsep}{2.7pt}
\renewcommand{\arraystretch}{1.15}
\begin{tabular}{c r|rrrrrrrrrrrrrr}
\toprule
$1\cdot\alpha+2\cdot\beta$ & $h_\rho$ & $\ch{0}{7}$ & $\ch{1}{6}$ & $\ch{2}{5}$ & $\ch{3}{4}$ & $\ch{01}{05}$ & $\ch{01}{14}$ & $\ch{01}{23}$ & $\ch{02}{04}$ & $\ch{02}{13}$ & $\ch{12}{03}$ & $\ch{03}{03}_{+}$ & $\ch{03}{03}_{-}$ & $\ch{12}{12}_{+}$ & $\ch{12}{12}_{-}$\\
\midrule
$1\cdot8$ & 1 & 1 & 7 & 21 & 35 & 20 & 28 & 14 & 64 & 70 & 56 & 45 & 45 & 21 & 21\\
$1\cdot5+3$ & 112 & 1 & 4 & 6 & 5 & 5 & 1 & $-1$ & 4 & $-5$ & $-4$ & 0 & 0 & $-3$ & $-3$\\
$1\cdot4+2\cdot2$ & 210 & 1 & 3 & 1 & $-5$ & 4 & 4 & 2 & 0 & 2 & 0 & $-3$ & $-3$ & 1 & 1\\
$1\cdot3+5$ & 1344 & 1 & 2 & 1 & 0 & 0 & $-2$ & $-1$ & $-1$ & 0 & 1 & 0 & 0 & 1 & 1\\
$1\cdot2+2+4$ & 2520 & 1 & 1 & $-1$ & $-1$ & 0 & 0 & 0 & 0 & 0 & 0 & 1 & 1 & $-1$ & $-1$\\
$1\cdot2+3\cdot2$ & 1120 & 1 & 1 & 0 & 2 & $-1$ & 1 & 2 & $-2$ & 1 & $-1$ & 0 & 0 & 0 & 0\\
$1+2\cdot2+3$ & 1680 & 1 & 0 & $-2$ & 1 & 1 & 1 & $-1$ & 0 & $-1$ & 0 & 0 & 0 & 1 & 1\\
$2\cdot4$ & 105 & 1 & $-1$ & $-3$ & 3 & 4 & $-4$ & 6 & 0 & $-2$ & 8 & $-3$ & $-3$ & $-3$ & $-3$\\
$2+6$ & 3360 & 1 & $-1$ & 0 & 0 & 1 & $-1$ & 0 & 0 & 1 & $-1$ & 0 & 0 & 0 & 0\\
$4\cdot2$ & 1260 & 1 & $-1$ & 1 & $-1$ & 0 & 0 & 2 & 0 & $-2$ & 0 & 1 & 1 & 1 & 1\\
\midrule
$(1+7)_{+}$ & 2880 & 1 & 0 & 0 & 0 & $-1$ & 0 & 0 & 1 & 0 & 0 & $\frac12(-1+\sqrt{-7})$ & $\frac12(-1-\sqrt{-7})$ & 0 & 0\\
$(1+7)_{-}$ & 2880 & 1 & 0 & 0 & 0 & $-1$ & 0 & 0 & 1 & 0 & 0 & $\frac12(-1-\sqrt{-7})$ & $\frac12(-1+\sqrt{-7})$ & 0 & 0\\
$(3+5)_{+}$ & 1344 & 1 & $-1$ & 1 & 0 & 0 & 1 & $-1$ & $-1$ & 0 & 1 & 0 & 0 & $\frac12(-1+\sqrt{-15})$ & $\frac12(-1-\sqrt{-15})$\\
$(3+5)_{-}$ & 1344 & 1 & $-1$ & 1 & 0 & 0 & 1 & $-1$ & $-1$ & 0 & 1 & 0 & 0 & $\frac12(-1-\sqrt{-15})$ & $\frac12(-1+\sqrt{-15})$\\
\bottomrule
\end{tabular}
\end{center}
\end{landscape}

In derselben Art wie in diesen Beispielen findet man für jeden Werth von $n$ die Charaktere der alternirenden Gruppe $\Hgr$. Seien
\begin{equation}\tag{2.}
 c_1=2a_1+1,	\quad c_2=2a_2+1,	\quad \ldots,	\quad c_s=2a_s+1\qquad(c_1<c_2<\cdots<c_s)
\end{equation}
$s$ verschiedene ungerade Zahlen, deren Summe
\begin{equation}\tag{3.}
 c_1+c_2+\cdots+c_s=n
\end{equation}
ist. Der Classe (2.) der symmetrischen Gruppe $\Hgr'$ entspricht der sich selbst associirte Charakter
\begin{equation}\tag{4.}
 \chi\ch{a_1\ a_2\ \cdots\ a_s}{a_1\ a_2\ \cdots\ a_s},
\end{equation}
und umgekehrt diesem Charakter die Classe (2.). Auf diese Art sind die $v$ sich selbst associirten Charaktere (4.) der Gruppe $\Hgr'$ und die $v$ Classen (2.) einander gegenseitig eindeutig zugeordnet. Gehört $R$ der Classe (2.) an, und ist $T$ irgend eine negative Permutation, so repraesentiren $R$ und $T^{-1}RT=S$ die beiden conjugirten Classen, worin die Classe (2.) in der alternirenden Gruppe $\Hgr$ zerfällt. Ebenso zerfällt der Charakter (4.) von $\Hgr'$ in zwei conjugirte Charaktere von $\Hgr$,
\begin{equation}\tag{5.}
 \chi(P)=\phi(P)+\psi(P),
\end{equation}
wo
\[
 \psi(P)=\phi(T^{-1}PT)
\]
ist. Dann ist für jede durch $Q$ repraesentirte Classe von $\Hgr$
\begin{equation}\tag{6.}
 \phi(Q)=\psi(Q)=\frac12\chi(Q)
\end{equation}
mit Ausnahme der beiden Classen $(R)$ und $(S)$. Ich habe schon in § 2 gezeigt, dass diese Gleichung für die $u$ Classen von $\Hgr$ gilt, deren Cyklen nicht den Bedingungen (2.) genügen. Sie gilt aber auch für $2v-2$ der übrigen Classen. Dagegen ist
\begin{equation}\tag{7.}
\begin{aligned}
 \phi(R)&=\frac12(\eps+\sqrt{\eps p}),&\qquad \psi(R)&=\frac12(\eps-\sqrt{\eps p}),\\
 \phi(S)&=\frac12(\eps-\sqrt{\eps p}),&\qquad \psi(S)&=\frac12(\eps+\sqrt{\eps p}),
\end{aligned}
\end{equation}
wo
\begin{equation}\tag{8.}
 p=c_1c_2\cdots c_s=\frac{h}{h_R},	\qquad \eps=(-1)^{\frac12(p-1)}=(-1)^{\frac12(n-s)}
\end{equation}
ist.

Die cyklische Permutation $C=(0,1,2,\ldots,a,a+1,\ldots,2a-1,2a)$ der Ordnung $c=2a+1$ wird durch die Permutation
\[
 G=(1,2a)(2,2a-1)\cdots(a,a+1)
\]
in
\[
 G^{-1}CG=(0,2a,2a-1,\ldots,a+1,a,\ldots,2,1)=C^{-1}
\]
transformirt. $G$ besteht aus $a=\frac12(c-1)$ Transpositionen. Daher wird $R=C_1C_2\cdots C_s$ durch eine Permutation $H$, die aus $\frac12\sum(c-1)=\frac12(n-s)$ Transpositionen besteht in $H^{-1}RH=R^{-1}$ transformirt. Nun sind unter der Voraussetzung (2.) die Permutationen, die $R$ in $R^{-1}$ transformiren, entweder alle gerade oder alle ungerade. Daher gehören $R$ und $R^{-1}$ in $\Hgr$ zu derselben Classe oder zu zwei verschiedenen conjugirten Classen, je nachdem $\eps=+1$ oder $-1$ ist. Im ersten Falle ist $\phi(R)=\phi(R^{-1})$ reell, im anderen sind $\phi(R)$ und $\phi(R^{-1})=\psi(R)$ conjugirte complexe Grössen.

\section*{§ 4.}

Dem Beweise des oben aufgestellten Satzes schicke ich zwei Bemerkungen voraus, die für beliebige Gruppen gelten.

1. Seien $\Hgr_1$ und $\Hgr_2$ zwei Gruppen der Ordnungen $h_1$ und $h_2$, die nur das Hauptelement gemeinsam haben. Ferner sei jedes Element von $\Hgr_1$ mit jedem von $\Hgr_2$ vertauschbar. Dann ist $\Hgr=\Hgr_1\Hgr_2=\Hgr_2\Hgr_1$ eine Gruppe der Ordnung $h=h_1h_2$. Zerfallen die Elemente von $\Hgr_1$ und $\Hgr_2$ in $k_1$ und $k_2$ Classen, so zerfallen die von $\Hgr$ in $k=k_1k_2$ Classen. Sind $A_1,B_1,\ldots,R_1,\ldots$ ($A_2,B_2,\ldots,R_2,\ldots$) Repraesentanten der $k_1$ ($k_2$) Classen von $\Hgr_1$ ($\Hgr_2$), so repraesentiren $A_1A_2,A_1B_2,A_2B_1,A_2B_2,\ldots,R_1R_2,\ldots$ die $k_1k_2=k$ Classen von $\Hgr$. Ist $\chi_1(R_1)$ ($\chi_2(R_2)$) ein Charakter von $\Hgr_1$ ($\Hgr_2$), und ist $R=R_1R_2$, so ist $\chi(R)=\chi_1(R_1)\chi_2(R_2)$ ein Charakter von $\Hgr$. Setzt man für $\chi_1$ ($\chi_2$) der Reihe nach jeden der $k_1$ ($k_2$) verschiedenen Charaktere von $\Hgr_1$ ($\Hgr_2$), so erhält man die $k_1k_2=k$ verschiedenen Charaktere von $\Hgr$. Diese Sätze ergeben sich ohne Weiteres aus den Eigenschaften, durch die ich in § 1 meiner Arbeit \emph{Über die Darstellung der endlichen Gruppen durch lineare Substitutionen}, Sitzungsberichte 1897, die Charaktere definirt habe.

2. Zwischen den Charakteren $\chi$ einer Gruppe $\Hgr$ und den Charakteren $\psi$ einer ihrer Untergruppen $\Ggr$ besteht nach U. § 1 (5.) die Relation
\begin{equation}\tag{1.}
 \sum_\lambda r_{\varkappa\lambda}\chi^{(\lambda)}(R)=\frac{h}{g h_\rho}\sum_{P\in(\rho)}\psi^{(\varkappa)}(P).
\end{equation}
Die positiven ganzen Zahlen $r_{\varkappa\lambda}$ sind von der Classe $(\rho)$ des Elementes $R$ unabhängig. Die $g_\rho$ in $\Ggr$ enthaltenen Elemente $P$ der $\rho^{\ten}$ Classe von $\Hgr$ vertheilen sich auf mehrere, etwa auf $m$, verschiedene Classen von $\Ggr$. Sie seien durch $P_1,P_2,\ldots,P_m$ repraesentirt. Besteht die Classe von $P_\mu$ aus $l_\mu$ Elementen, so ist
\[
 \sum_{(\rho)}\psi(P)=l_1\psi(P_1)+l_2\psi(P_2)+\cdots+l_m\psi(P_m).
\]
Die Gruppe der Elemente von $\Ggr$, die mit $P$ vertauschbar sind, ist enthalten in der Gruppe der Elemente von $\Hgr$, die mit $P$ vertauschbar sind. In dem Falle, den ich im Folgenden zu betrachten habe, sind diese beiden Gruppen identisch für jedes der $g_\rho$ Elemente $P$ von $\Ggr$, die einer bestimmten Classe $(\rho)$ von $\Hgr$ angehören, oder jedes mit $P$ vertauschbare Element von $\Hgr$ ist in $\Ggr$ enthalten. Die Anzahl der mit $P$ vertauschbaren Elemente von $\Hgr$ ist $h/h_\rho$, die Anzahl der mit $P_\mu$ vertauschbaren Elemente von $\Ggr$ ist $g/l_\mu$. Daher ist
\[
 l_1=l_2=\cdots=l_m=\frac{g_\rho}{m},	\qquad \frac{h}{h_\rho}=\frac{g}{l_\mu}=\frac{gm}{g_\rho},
\]
und nach (1.)
\begin{equation}\tag{2.}
 \psi^{(\varkappa)}(P_1)+\psi^{(\varkappa)}(P_2)+\cdots+\psi^{(\varkappa)}(P_m)=\sum_\lambda r_{\varkappa\lambda}\chi^{(\lambda)}(R).
\end{equation}

\section*{§ 5.}

Ich wende mich jetzt zu dem Beweise der Regel, die ich in § 3 zur Berechnung der Charaktere der alternirenden Gruppe $\Hgr$ des Grades $n$ aufgestellt habe. Ich nehme dabei an, diese Regel sei für jede alternirende Gruppe $\Hgr_1$ des Grades $n_1<n$ bereits bewiesen.

Ich benutze die Bezeichnungen (2.) und (3.), § 3, setze aber, falls $n$ ungerade ist, $s>1$ voraus. Ich theile die $n$ Symbole in die $c_1$ ersten und die $n-c_1$ letzten und bilde die symmetrische und die alternirende Gruppe $\Hgr'_1$, $\Hgr_1$ und $\Hgr'_2$, $\Hgr_2$ für jene $c_1$ und für diese $n-c_1$ Symbole. Ist $c_1=1$, so bedürfen die folgenden Entwickelungen einer Modification, die ich ihrer Einfachheit halber übergehe.

Dann ist $\Ggr=\Hgr_1\Hgr_2$ eine Untergruppe der alternirenden Gruppe $\Hgr$ des Grades $n$. $R_1$ bestehe aus einem Cyklus der $c_1$ ersten Symbole, $R_2$ aus $s-1$ Cyklen von je $c_2,\ldots,c_s$ der letzten $n-c_1$ Symbole. $T_1$ ($T_2$) sei eine bestimmte negative Permutation von $\Hgr'_1$ ($\Hgr'_2$). Ferner sei $T_1^{-1}R_1T_1=S_1$ und $T_2^{-1}R_2T_2=S_2$. Ist dann
\[
 p_1=c_1,	\quad p_2=c_2\cdots c_s,	\quad \eps_1=(-1)^{\frac12(p_1-1)},	\quad \eps_2=(-1)^{\frac12(p_2-1)},
\]
so giebt es in $\Hgr_1$ zwei conjugirte Charaktere $\phi_1(P_1)$ und $\psi_1(P_1)=\phi_1(T_1^{-1}P_1T_1)$, wofür
\[
 \phi_1(R_1)=\psi_1(S_1)=\frac12(\eps_1+\sqrt{\eps_1p_1}),	\quad
 \phi_1(S_1)=\psi_1(R_1)=\frac12(\eps_1-\sqrt{\eps_1p_1}),
\]
und in $\Hgr_2$ zwei conjugirte Charaktere $\phi_2(P_2)$ und $\psi_2(P_2)=\phi_2(T_2^{-1}P_2T_2)$, wofür
\[
 \phi_2(R_2)=\psi_2(S_2)=\frac12(\eps_2+\sqrt{\eps_2p_2}),	\quad
 \phi_2(S_2)=\psi_2(R_2)=\frac12(\eps_2-\sqrt{\eps_2p_2})
\]
ist, während für jede andere Classe $\phi_1(Q_1)=\psi_1(Q_1)$ und $\phi_2(Q_2)=\psi_2(Q_2)$ eine rationale ganze Zahl ist. Daraus ergeben sich nach § 4 vier verschiedene Charaktere von $\Ggr$, $\phi_1\phi_2$, $\phi_1\psi_2$, $\phi_2\psi_1$, $\phi_2\psi_2$. Ist
\[
 R_1R_2=R,	\qquad R_1S_2=S,	\qquad T_1T_2=T,
\]
so ist $T$ eine positive Permutation, und es ist
\[
 T^{-1}RT=S_1S_2,	\qquad T^{-1}ST=S_1R_2,	\qquad T_2^{-1}RT_2=S.
\]
Daher vereinigen sich die beiden Classen $(R_1R_2)$ und $(S_1S_2)$ von $\Ggr$ zu einer Classe $(R)$ von $\Hgr$, und die beiden Classen $(R_1S_2)$ und $(S_1R_2)$ zu einer Classe $(S)$, die beiden Classen $(R)$ und $(S)$ aber von $\Hgr$ sind verschieden, weil $T_2$ ungerade ist.

Die Classe $(R)$ kann nur solche Elemente von $\Ggr$ enthalten, die einer der beiden Classen $(R_1R_2)$ oder $(S_1S_2)$ angehören. Denn jede Permutation $P$ der Classe $(R)$ besteht aus $s$ Cyklen $C_1,C_2,\ldots,C_s$ von je $c_1,c_2,\ldots,c_s$ Symbolen. Soll sie der Gruppe $\Ggr$ angehören, so muss $P=P_1P_2$ sein, wo $P_1$ nur die ersten $c_1$, $P_2$ nur die letzten $n-c_1$ Symbole versetzt. Dies ist aber nur möglich, wenn $P_1=C_1$ und $P_2=C_2\cdots C_s$ ist, weil $c_1$ die kleinste der Zahlen $c_1,c_2,\ldots,c_s$ ist, also nicht als Summe von einigen der Zahlen $c_2,\ldots,c_s$ dargestellt werden kann.

Das ist der Grund, aus dem ich die $n$ Symbole gerade in dieser Weise in zwei Abtheilungen von $n_1=c_1$ und $n_2=n-c_1$ Symbolen getheilt habe. Denselben Zweck erreicht man aber auch, wenn man $n_1=c_2$, $n_2=c_1+c_3+\cdots+c_s$ setzt, und man hat dann noch den Vortheil, dass nicht $n_1=1$ sein kann.

Demnach gehört $P_1$ der Classe $(R_1)$ oder $(S_1)$ an, und $P_2$ der Classe $(R_2)$ oder $(S_2)$. Wie oben gezeigt, sind aber nur die Combinationen
\[
 (R)=(R_1R_2)+(S_1S_2),	\qquad (S)=(R_1S_2)+(S_1R_2)
\]
möglich.

Die mit $R$ vertauschbaren Elemente von $\Hgr$ sind nach (1.), § 1 sämmtlich in $\Ggr$ enthalten. Daher giebt es einen zusammengesetzten Charakter $\phi(P)$ von $\Hgr$, für den
\[
 \phi(R)=\frac14(\eps_1+\sqrt{\eps_1p_1})(\eps_2+\sqrt{\eps_2p_2})+\frac14(\eps_1-\sqrt{\eps_1p_1})(\eps_2-\sqrt{\eps_2p_2})
\]
ist. Setzt man
\[
 p=p_1p_2=c_1c_2\cdots c_s,	\qquad \eps=\eps_1\eps_2=(-1)^{\frac12(p-1)}=(-1)^{\frac12(n-s)},
\]
und ist $\psi(P)=\phi(T_2^{-1}PT_2)$ der zu $\phi(P)$ conjugirte Charakter, so ist
\[
 \phi(R)=\psi(S)=\frac12(\eps+\sqrt{\eps p}),	\qquad \phi(S)=\psi(R)=\frac12(\eps-\sqrt{\eps p}).
\]
Für jedes Element $Q_2$ von $\Hgr_2$, das der Classe $(R_2)$ nicht angehört, ist $\phi_2(Q_2)=\phi_2(T_2^{-1}Q_2T_2)$. Wendet man nun die Formel (1.), § 4 auf das Element $Q=R_1Q_2$ von $\Hgr$ an, so steht rechts neben jedem Gliede $\phi_1(R_1)\phi_2(Q_2)$ auch das Glied $\phi_1(T_1^{-1}R_1T_1)\phi_2(T_2^{-1}Q_2T_2)$ und die Summe dieser beiden Glieder ist eine rationale Zahl, die sich nicht ändert, wenn man $Q$ durch $T_2^{-1}QT_2$ ersetzt. Daher ist $\phi(Q)=\psi(Q)$ eine ganze rationale Zahl, und dasselbe gilt für ein Element $Q=Q_1R_2$ oder $Q=Q_1Q_2$.

Folglich hat $\Hgr$ einen zusammengesetzten Charakter $\phi-\psi=\vartheta$, der die Werthe
\begin{equation}\tag{1.}
 \vartheta(R)=\sqrt{\eps p},\qquad \vartheta(S)=-\sqrt{\eps p},\qquad \vartheta(Q)=0
\end{equation}
hat, falls $(Q)$ irgend eine Classe ist, die von den beiden bestimmten conjugirten Classen $(R)$ und $(S)$ verschieden ist. Für jede dieser beiden Classen ist $h_\rho=h/p$ und mithin
\[
 \sum_\rho h_\rho\vartheta_\rho\overline{\vartheta}_\rho=2h,
\]
falls $\overline{\vartheta}_\rho$ die zu $\vartheta_\rho$ conjugirte complexe Grösse bezeichnet. Ist also, wie in § 2, $\vartheta(P)=\sum r_\varkappa\phi^{(\varkappa)}(P)$, so ist $\sum r_\varkappa^2=2$, daher sind zwei der Zahlen $r_\varkappa$ gleich $\pm1$, die übrigen gleich $0$. Nun ist $\phi^{(\varkappa)}(E)$ eine positive ganze Zahl und $\vartheta(E)=0$. Folglich ist $\vartheta(P)$ die Differenz von zwei einfachen Charakteren von $\Hgr$. Die beiden Charaktere $\phi$ und $\psi$ von $\Hgr$, als deren Differenz $\vartheta$ oben erhalten ist, können zusammengesetzt sein. Von jetzt an bezeichne ich mit $\phi$ und $\psi$ die beiden einfachen Charaktere von $\Hgr$, deren Differenz $\phi-\psi=\vartheta$ ist.

Ist $\chi(P)$ einer der $u$ Charaktere von $\Hgr$, die aus einem Paare associirter Charaktere von $\Hgr'$ entspringen, so ist $\chi(R)=\chi(S)$, und nach (1.) ist mithin $\sum h_\rho\vartheta_\rho\chi_\rho=0$. Folglich ist keiner der beiden Charaktere $\phi$ oder $\psi$ gleich $\chi$, sondern $\phi$ und $\psi$ gehören zu den $2v$ Charakteren von $\Hgr$, die aus den $v$ sich selbst associirten Charakteren von $\Hgr'$ entspringen, und die in $v$ Paare conjugirter Charaktere zerfallen. Der zu $\phi$ conjugirte Charakter $\omega$ ist also von $\phi$ verschieden. Ist $\chi=\phi+\omega$, so ist $\chi(R)=\chi(S)$ und mithin $\sum h_\rho\vartheta_\rho\chi_\rho=0$ oder $\sum h_\rho(\phi_\rho-\psi_\rho)(\phi_\rho+\omega_\rho)=0$, und weil $\phi$ von $\psi$ und von $\omega$ verschieden ist, $\sum h_\rho\psi_\rho\omega_\rho=h$. Folglich ist $\omega=\psi$, oder $\phi$ und $\psi$ sind conjugirte Charaktere.

Das Paar conjugirter Classen $(R)$ und $(S)$ und das Paar conjugirter Charaktere $\phi$ und $\psi$ ordne ich einander zu. Dann entsprechen nach (1.) verschiedenen Paaren $(R),(S)$ verschiedene Paare $\phi,\psi$ und umgekehrt. Die Werthe eines Charakters sind ganze algebraische Zahlen. Daher ist $\chi(R)=\chi(S)=2\phi(R)-\sqrt{\eps p}$ eine ungerade (rationale) Zahl, dagegen $\chi(Q)=2\phi(Q)=2\psi(Q)$ eine gerade Zahl.

Um den Inductionsschluss anwenden zu können, ist bisher $s>1$ vorausgesetzt. Ist also $n$ ungerade, so bleibt ein Paar conjugirter Classen $(R),(S)$ übrig, deren Permutationen aus einem einzigen Cyklus von $n$ Symbolen bestehen, und mithin auch ein Paar conjugirter Charaktere $\phi,\psi$, deren Summe und Differenz ich mit $\chi$ und $\vartheta$ bezeichne. Ist dann $(\overline{R}),(\overline{S})$ irgend ein anderes Paar conjugirter Classen, und $\overline{\phi},\overline{\psi}$ das entsprechende Paar conjugirter Charaktere, und ist $\overline{\vartheta}=\overline{\phi}-\overline{\psi}$, so ist $\sum h_\rho\overline{\vartheta}_\rho\phi_\rho=0$ und mithin $\phi(\overline{R})=\phi(\overline{S})=\psi(\overline{R})=\psi(\overline{S})$. Ist also $(Q)$ irgend eine von $(R)$ und $(S)$ verschiedene Classe, die von der conjugirten verschieden ist, so ist $\vartheta(Q)=0$, $\phi(Q)=\psi(Q)=\frac12\chi(Q)$. Nach § 2 ist $\phi(R)=\psi(S)$ und $\phi(S)=\psi(R)$. Die Elemente $R$ und $R^{-1}$ gehören derselben Classe an, wenn $n\equiv1\pmod 4$ ist, zwei conjugirten Classen, wenn $n\equiv -1\pmod 4$ ist. Im ersten Falle sind daher $\phi(R)$ und $\phi(S)$ reelle Grössen, im zweiten conjugirte complexe. Nun ist $\sum h_\rho\phi_\rho\vartheta_\rho=h$, also wenn $\eps=(-1)^{\frac12(n-1)}=+1$ ist, $h_R(\phi(R)\vartheta(R)+\phi(S)\vartheta(S))=h$, also
\[
 \vartheta(R)^2=\frac{h}{h_R}=p=n.
\]
Ist aber $\eps=-1$, so ist $h_R(\phi(R)\vartheta(S)+\phi(S)\vartheta(R))=h$, also $\vartheta(R)^2=-p$. Mithin ist $\vartheta(R)=\sqrt{\eps p}$, $\vartheta(S)=-\sqrt{\eps p}$ und demnach gelten für diese Zuordnung dieselben Gesetze wie oben.

Durch die Bedingung, dass $\chi(R)=\chi(S)$ ungerade, sonst aber $\chi(Q)$ gerade ist, ist die Zuordnung zwischen einem Paar conjugirter Charaktere, deren Summe $\chi$ ist, und dem entsprechenden Paar conjugirter Classen $(R),(S)$ vollständig bestimmt, und zwar muss, wie oben behauptet, $\chi$ die Charakteristik
\[
 \ch{a_1\ a_2\ \cdots\ a_s}{a_1\ a_2\ \cdots\ a_s}
\]
haben, wenn
\[
 c_1=2a_1+1,	\quad c_2=2a_2+1,	\quad \cdots,	\quad c_s=2a_s+1
\]
die Ordnungen der Cyklen von $R$ sind. Denn nach der Formel (11.) § 7 meiner Arbeit \emph{Über die Charaktere der symmetrischen Gruppe} ist
\[
 \chi\ch{a_1\ \cdots\ a_s}{a_1\ \cdots\ a_s}(R)=\eps=(-1)^{\frac12(p-1)}=(-1)^{\frac12(n-s)}
\]
ungerade. Damit sind die $k$ Charaktere von $\Hgr$ vollständig bestimmt.

Entspricht die Classe $(\varkappa)$ und die conjugirte dem Charakter $\chi^{(\varkappa)}$ und dem conjugirten, so ist nach der Formel (9.), § 4 der eben citirten Arbeit und der Formel (6.), § 3 dieser Arbeit
\[
 f^{(\varkappa)}=\frac12\frac{n!\,\Delta(a_1,\ldots,a_s)^2}{(a_1!\cdots a_s!)^2\prod(a_\alpha+a_\beta+1)}.
\]
Nun ist
\[
 \frac{h}{h_\varkappa}=(2a_1+1)\cdots(2a_s+1)
\]
und folglich ist
\begin{equation}\tag{2.}
 \frac{h_\varkappa}{f^{(\varkappa)}}=\left(\frac{a_1!\cdots a_s!\prod'(a_\alpha+a_\beta+1)}{\Delta(a_1,\ldots,a_s)}\right)^2
\end{equation}
das Quadrat einer ganzen Zahl. Das Product $\prod'$ ist nur über die $\frac12s(s-1)$ Paare verschiedener Indices zu erstrecken. Die in $\chi^{(\varkappa)}_\varkappa$ auftretende Quadratwurzel kann daher aus
\[
 \frac{\eps h}{h_\varkappa}\quad\hbox{oder}\quad \frac{\eps h}{f^{(\varkappa)}}
\]
gezogen werden. Sie kann in Ausnahmefällen rational sein, z. B. wenn $n=9$ und $R$ aus einem Cyklus von $9$ Symbolen besteht.
% ===== END WORK 061 =====

\clearpage

% ===== BEGIN WORK 070: 070_Theorie_der_hyperkomplexen_Gr_en =====
\frobtitle{Theorie der hyperkomplexen Größen}

Die Methoden, die ich in meiner Arbeit \emph{Über die Primfactoren der Gruppendeterminante}, Sitzungsberichte 1896 (im folgenden \emph{Gr.} zitiert) zur Erforschung der Eigenschaften der Determinante einer endlichen Gruppe entwickelt habe, reichen auch zur Untersuchung eines beliebigen aus $n$ Grundzahlen gebildeten Systems hyperkomplexer Größen aus. Mit solchen Größen beschäftigt sich Hr. \textsc{Molien} in seiner grundlegenden Abhandlung \emph{Über Systeme höherer complexer Zahlen}, Math. Ann. Bd. 41 (im folgenden \emph{Mol.} zitiert).

Das Verständnis seiner Arbeit hat er dadurch etwas erschwert, daß er nach Möglichkeit die Rechnung unterdrückt und den gedanklichen Inhalt der Beweise rein darzustellen versucht hat. So bin ich nicht darüber ins klare gekommen, ob die von ihm entwickelten Hülfsmittel zu einem strengen Beweise des Satzes 25 ausreichen. Umgekehrt hat er einen in § 5 gemachten Fehlschluß in einer Notiz im 42. Bande der Math. Ann. durch eine recht umständliche Rechnung berichtigt. Indem ich an Stelle der Gleichung, die er die \textsc{Killing}sche nennt, eine lineare Verbindung $S(x)+T(y)$ der in verschiedenen Variabeln ausgedrückten Matrizen der beiden antistrophen Gruppen benutze, gelingt es mir, gerade diesen Beweis erheblich zu vereinfachen.

Außerdem aber findet sich, wie ich vor kurzem bemerkt habe, in seinem Beweise des Satzes 19 eine nicht unwesentliche Lücke. Aber ungeachtet dieser kleinen Mängel bedeutet seine bahnbrechende, gedankenreiche Arbeit, die er trotz eines ziemlich unvollkommenen Rüstzeuges mit unablässiger Beharrlichkeit und durchdringendem Scharfsinn durchgeführt hat, einen der wichtigsten Fortschritte auf dem Teilgebiete der Algebra, das man als Gruppentheorie bezeichnet. Aus den angeführten Gründen halte ich es für angemessen, die Untersuchung mit den Hülfsmitteln, die ich namentlich in meiner Arbeit \emph{Über vertauschbare Matrizen}, Sitzungsberichte 1896, entwickelt habe, von neuem aufzunehmen.

Die Methoden von \textsc{Lie} werden in dieser rein algebraischen Arbeit gar nicht benutzt. Dagegen hat meine Darstellung manche Berührungspunkte mit der von \textsc{Dedekind} in seiner Abhandlung \emph{Zur Theorie der aus $n$ Haupteinheiten gebildeten complexen Grössen}, Göttinger Nachrichten 1885 (im folgenden \emph{Ded.} zitiert).

Die Grundlage meiner Untersuchung bildet die Formel (2.) § 5. Mit ihrer Hülfe zeige ich in § 6 direkt, daß die in den verschiedenen Primfaktoren der Gruppendeterminante auftretenden Variabeln alle von einander unabhängig sind, und umgehe dadurch die Sätze 3, 4 und 5 von \textsc{Molien} sowie den Beweis des Satzes 25. Und mit derselben Formel beweise ich in § 7, daß die Elementarteiler der Determinante einer einfachen Gruppe alle linear sind. Auf diesen beiden Ergebnissen aber beruht die ganze Entwicklung.

Zur Klassifikation der Gruppen, wozu bisher nur unzureichende Ansätze gemacht sind, erweisen sich als die geeignetsten Invarianten weniger die Exponenten und Grade der Elementarteiler, worin die Determinanten der beiden antistrophen Gruppen zerfallen, als vielmehr (§ 9) die nämlichen Zahlen für die bisher kaum beachtete parastrophe Matrix $R(\xi)$, insbesondere der Rang von $R(\xi)$; bei nicht kommutativen Gruppen aber vor allem die elementaren Invarianten der Matrix $uR(\xi)+vR'(\xi)$ und der von $2n$ Variabeln abhängenden Matrizen $S(x)+T(y)$ und $R(\xi)+R'(\eta)$, und zwar sind diese Invarianten auch dann von Bedeutung, wenn etwa die Determinante von $R(\xi)$ oder von $R(\xi)+R'(\eta)$ identisch verschwindet.

Erst nach Vollendung dieser Untersuchung bin ich auf die ausgezeichnete Abhandlung des Hrn. \textsc{Cartan}, \emph{Sur les groupes bilinéaires et les systèmes de nombres complexes}, Ann. de Toulouse, tome XII, 1898, aufmerksam geworden, worin er die Resultate von \textsc{Molien} ableitet, ohne, wie es scheint, seine Arbeit zu kennen. Mit dessen Methoden und mit den hier benutzten hat der von Hrn. \textsc{Cartan} eingeschlagene Weg nicht das geringste gemeinsam. Die Transformation der Basis, der Ausgangspunkt und das Endziel seiner Untersuchung, ist von mir so lange wie irgend möglich vermieden worden (§ 9). Die von jeder Darstellung der Gruppe unabhängigen, invarianten Eigenschaften, mit denen ich beginne, ergeben sich bei ihm erst am Schluß durch Deutung einer Normalform der Gruppe, erhalten durch eine lange Reihe von Umformungen, deren Ziel erst am Ende der Entwicklung klar wird. Der Unterschied zwischen den beiden Methoden ist also derselbe, wie der zwischen dem Verfahren von \textsc{Weierstrass} und dem von \textsc{Kronecker} in der Theorie der Scharen von bilinearen Formen. Eine besonders beachtenswerte Formel des Hrn. \textsc{Cartan} (§ 65, (37.)), die sich bei \textsc{Molien} nicht findet, hatte ich (§ 12, (5.)) in der einfachsten Art durch die Zerlegung der Determinante $|S(x)+T(y)|$ in Primfaktoren erhalten.

\section*{§ 1.}

Die Koordinaten $x_1,x_2,\ldots,x_n$ der aus $n$ Grundzahlen $\eps_1,\eps_2,\ldots,\eps_n$ gebildeten hyperkomplexen Größe
\[
\fnum{1} x=\eps_1x_1+\eps_2x_2+\cdots+\eps_nx_n
\]
können alle reellen und komplexen Werte annehmen. Die Gesamtheit dieser Größen, die sich durch Addition und Multiplikation mit gewöhnlichen Größen reproduzieren, nenne ich eine Gruppe $(\eps)$, wenn außerdem noch das Produkt von je zweien wieder dem Systeme angehört, und für die Multiplikation von je dreien das assoziative Gesetz $(xy)z=x(yz)$ gilt.

Die Grundzahlen können durch lineare Relationen miteinander verknüpft sein. Die Anzahl der unabhängigen unter ihnen nenne ich die Ordnung der Gruppe. Diesen Fall kann man leicht auf den zurückführen, wo die Grundzahlen unabhängig sind. Ich nehme daher an, daß $n$ die Ordnung der betrachteten Gruppe $(\eps)$ ist.

Erfolgt die Multiplikation der Grundzahlen nach dem Gesetze
\[
\fnum{2} \eps_\beta\eps_\gamma=\sum_\alpha a_{\alpha\beta\gamma}\eps_\alpha,
\]
so ergeben sich zwischen den hier auftretenden Koeffizienten aus dem assoziativen Prinzip und der Unabhängigkeit der Grundzahlen die Relationen
\[
\fnum{3} \sum_\kappa a_{\kappa\alpha\beta}a_{\gamma\kappa\delta}
      =\sum_\kappa a_{\kappa\beta\delta}a_{\gamma\alpha\kappa}.
\]
Ist dann $x=yz$ das Produkt der beiden Größen
\[
 y=\sum_\lambda\eps_\lambda y_\lambda,\qquad
 z=\sum_\lambda\eps_\lambda z_\lambda,
\]
so ist
\[
\fnum{4} x_\alpha=\sum_{\beta,\gamma}a_{\alpha\beta\gamma}y_\beta z_\gamma,
\]
oder, wenn $\xi_1,\xi_2,\ldots,\xi_n$ Variable sind, deren System ich mit $\xi$ bezeichne und einen Parameter nenne,
\[
\fnum{5} \xi(x)=x(\xi)=\sum_\alpha \xi_\alpha x_\alpha
  =F(\xi,y,z)=\sum_{\alpha,\beta,\gamma}a_{\alpha\beta\gamma}\xi_\alpha y_\beta z_\gamma.
\]
Setzt man
\[
\fnum{6} r_{\alpha\beta}(\xi)=\sum_\kappa a_{\kappa\alpha\beta}\xi_\kappa,\quad
s_{\alpha\beta}(y)=\sum_\lambda a_{\alpha\lambda\beta}y_\lambda,\quad
t_{\alpha\beta}(z)=\sum_\mu a_{\alpha\beta\mu}z_\mu,
\]
also
\[
s_{\alpha\beta}(y)=\frac{\partial x_\alpha}{\partial z_\beta},\qquad
t_{\alpha\beta}(z)=\frac{\partial x_\alpha}{\partial y_\beta},
\]
so ist demnach die trilineare Form
\[
\fnum{7} F(\xi,y,z)=\sum_{\beta,\gamma}r_{\beta\gamma}(\xi)y_\beta z_\gamma
=\sum_{\alpha,\gamma}s_{\alpha\gamma}(y)\xi_\alpha z_\gamma
=\sum_{\alpha,\beta}t_{\alpha\beta}(z)\xi_\alpha y_\beta .
\]
Die aus jenen drei Systemen linearer Funktionen gebildeten Matrizen $n$ten Grades bezeichne ich mit
\[
 R(\xi)=R_\xi=(r_{\alpha\beta}(\xi)),\qquad
 S(y)=S_y=(s_{\alpha\beta}(y)),\qquad
 T(z)=T_z=(t_{\alpha\beta}(z)).
\]

Unter $f(x)$ verstehe ich stets eine Funktion der Koordinaten $x_1,x_2,\ldots,x_n$ der Größe $x$, unter $S(x)$ oder $S_x$ eine Matrix, deren Elemente (lineare) Funktionen dieser Koordinaten sind. Wenn es sich um eine Funktion der hyperkomplexen Größe $x$ selbst handelt, werde ich mich des Zeichens $f((x))$ bedienen.

Die Bedingungen (3.) des assoziativen Prinzips sind bequemer zu erörtern, wenn man sie dadurch zusammenfaßt, daß man sie mit Variabeln multipliziert und addiert. So erhält man, wenn $x$ und $y$ zwei beliebige Größen sind (\emph{Ded.} (36.)),
\[
\fnum{8} \sum_\lambda s_{\alpha\lambda}(x)t_{\lambda\beta}(y)
      =\sum_\lambda t_{\alpha\lambda}(y)s_{\lambda\beta}(x),
\]
oder einfacher
\[
\fnum{9} S_xT_y=T_yS_x,
\]
und
\[
\fnum{10} \sum_\lambda r_{\alpha\lambda}(\xi)s_{\lambda\beta}(x)
       =\sum_\lambda t_{\lambda\alpha}(x)r_{\lambda\beta}(\xi),
\]
oder
\[
\fnum{11} R_\xi S_x=T'_xR_\xi,
\]
wo $T'$ die zu $T$ konjugierte Matrix ist.

Ich nenne $S(x)$ und $|S(x)|$ die Gruppenmatrix und die Gruppendeterminante, $T(x)$ und $|T(x)|$ die antistrophe Gruppenmatrix und Gruppendeterminante, $R(\xi)$ und $|R(\xi)|$ die parastrophe Matrix und Determinante. Der in der Formel (9.) enthaltene Satz läßt sich in folgender Art aussprechen und umkehren:

\emph{Die Gruppenmatrix $S(x)$ ist mit der antistrophen Matrix $T(y)$ vertauschbar. Verschwindet die parastrophe Determinante nicht identisch, so läßt sich jede von $x$ unabhängige Matrix, die für jeden Wert von $x$ mit $S(x)$ ($T(x)$) vertauschbar ist, auf die Form $T(y)$ ($S(y)$) bringen.}

Denn sei der Parameter $\xi$ so gewählt, daß die Determinante der Matrix $R(\xi)=R=(r_{\alpha\beta})$ nicht verschwindet, und sei $U=(u_{\alpha\beta})$ irgend eine von $x$ unabhängige Matrix $n$ten Grades, die für jedes $x$ mit $S(x)$ vertauschbar ist. Dann kann man $y$ so bestimmen, daß für $\lambda=1,2,\ldots,n$
\[
 \sum_\kappa \xi_\kappa u_{\kappa\lambda}
 =\sum_\alpha r_{\lambda\alpha}y_\alpha
 =\sum_\kappa \xi_\kappa t_{\kappa\lambda}(y)
\]
(nach (7.)) ist. Nun ist $SU=US$, also
\[
 \sum_\lambda s_{\alpha\lambda}(x)u_{\lambda\beta}
 =\sum_\lambda u_{\alpha\lambda}s_{\lambda\beta}(x),
\]
und mithin
\[
 \sum_\lambda\xi_\lambda s_{\lambda\alpha}(x)u_{\alpha\beta}
 =\sum_{\lambda,\mu}\xi_\lambda u_{\lambda\mu}s_{\mu\beta}(x)
 =\sum_{\lambda,\mu}\xi_\lambda t_{\lambda\mu}(y)s_{\mu\beta}(x).
\]
also nach (8.)
\[
 \sum_{\mu,\lambda}\xi_\lambda s_{\lambda\alpha}(x)u_{\alpha\beta}
 =\sum_{\mu,\lambda}\xi_\lambda s_{\lambda\alpha}(x)t_{\alpha\beta}(y).
\]
Nun ist aber nach (7.)
\[
 \sum_\kappa \xi_\kappa s_{\kappa\lambda}(x)=\sum_\alpha r_{\alpha\lambda}x_\alpha.
\]
Setzt man dies ein, so ergibt sich durch Vergleichung der Koeffizienten von $x_\alpha$
\[
 \sum_\lambda r_{\alpha\lambda}u_{\lambda\beta}
 =\sum_\lambda r_{\alpha\lambda}t_{\lambda\beta}(y),
\]
und mithin, weil $|R|$ von Null verschieden ist, $u_{\alpha\beta}=t_{\alpha\beta}(y)$, $U=T(y)$.

Ist speziell $U=E=(e_{\alpha\beta})$ die Hauptmatrix, so kann man eine Größe $y=e$ so bestimmen, daß $T(e)=E$ wird (\emph{Ded.} (44.)). Nach (11.) ist dann auch $S(e)=E$. Diese Größe
\[
\fnum{12} e=e_1\eps_1+e_2\eps_2+\cdots+e_n\eps_n
\]
heißt die Haupteinheit. Ihre Koordinaten genügen den Bedingungen
\[
\fnum{13} s_{\alpha\beta}(e)=t_{\alpha\beta}(e)=e_{\alpha\beta}
\]
oder
\[
\fnum{14} \sum_\lambda a_{\alpha\lambda\beta}e_\lambda=e_{\alpha\beta},\qquad
\sum_\mu a_{\alpha\beta\mu}e_\mu=e_{\alpha\beta}.
\]
Ist also $S(x)=S(y)$ oder $T(x)=T(y)$ oder $S(x)=T(y)$, so ist $x=y$; ist $R(\xi)=R(\eta)$ oder $R(\xi)=R'(\eta)$, so ist $\xi=\eta$.

\section*{§ 2.}

Ist $x=yz$, also
\[
\fnum{1} x_\alpha=\sum_{\lambda,\mu}a_{\alpha\lambda\mu}y_\lambda z_\mu
=\sum_\mu s_{\alpha\mu}(y)z_\mu=\sum_\lambda t_{\alpha\lambda}(z)y_\lambda,
\]
so ist
\[
s_{\alpha\beta}(x)=\sum_\lambda a_{\alpha\lambda\beta}x_\lambda
 =\sum_{\kappa,\lambda,\mu}a_{\lambda\kappa\mu}a_{\alpha\lambda\beta}y_\kappa z_\mu
 =\sum_\lambda s_{\alpha\lambda}(y)s_{\lambda\beta}(z)
\]
und
\[
t_{\alpha\beta}(x)=\sum_\lambda a_{\alpha\beta\lambda}x_\lambda
 =\sum_{\kappa,\lambda,\mu}a_{\lambda\kappa\mu}a_{\alpha\beta\lambda}y_\kappa z_\mu
 =\sum_\lambda t_{\lambda\beta}(z)t_{\alpha\lambda}(y).
\]
Mit $f(yz)$ oder $S(yz)$ bezeichne ich das, was aus $f(x)$ oder $S(x)$ wird, falls man darin die $n$ Variabeln $x_\alpha$ durch die Koordinaten (1.) des Produktes $yz$ ersetzt. Dann erhält man
\[
\fnum{2} S(yz)=S(y)S(z),\qquad T(yz)=T(z)T(y).
\]

Ist $S(x)=x_1E_1+x_2E_2+\cdots+x_nE_n$, so ist folglich
\[
\fnum{3} E_\beta E_\gamma=\sum_\alpha a_{\alpha\beta\gamma}E_\alpha,
\]
und da nur dann $S(x)=0$ ist, wenn $x=0$ ist, so sind die $n$ konstanten Matrizen $E_1,E_2,\ldots,E_n$ linear unabhängig. Sie bilden mithin eine Darstellung (vergl. § 16) der Gruppe $(\eps)$, und aus dieser folgt, daß die Voraussetzung der Unabhängigkeit der Grundzahlen mit den quadratischen Gleichungen (2.) § 1 verträglich ist, deren Koeffizienten den Gleichungen (3.) § 1 und den Ungleichheiten (1.) § 3 genügen, oder daß aus diesen Bedingungen keine linearen Beziehungen zwischen den Grundzahlen fließen. Erst auf Grund dieser Gewißheit sind Beweise zulässig, wie der an Formel (3.) § 3 anknüpfende, worin die hyperkomplexen Größen selbst benutzt werden.

Ist ferner $T(x)=x_1F_1+x_2F_2+\cdots+x_nF_n$, so ist
\[
\fnum{4} F_\beta F_\gamma=\sum_\alpha a_{\alpha\gamma\beta}F_\alpha.
\]
Da für Matrizen das assoziative Gesetz gilt, so bilden $F_1,F_2,\ldots,F_n$ die Basis einer Gruppe, der antistrophen Gruppe $(\eps')$. Eine andere Darstellung von $(\eps)$ liefert die zu $T$ konjugierte Matrix $T'$. Man könnte daher auch $T'$ als die Matrix der Gruppe $(\eps)$ und $S'$ als die der antistrophen Gruppe $(\eps')$ bezeichnen.

Ist $\xi$ ein Parameter, und setzt man
\[
\fnum{5^*} \zeta_\mu=\sum_{\kappa,\lambda}a_{\kappa\lambda\mu}\xi_\kappa y_\lambda
=\sum_\lambda r_{\lambda\mu}(\xi)y_\lambda
=\sum_\kappa s_{\kappa\mu}(y)\xi_\kappa,
\]
so wird
\[
 r_{\alpha\beta}(\zeta)=\sum_\mu a_{\mu\alpha\beta}\zeta_\mu
 =\sum_{\kappa,\mu}\xi_\kappa s_{\kappa\mu}(y)a_{\mu\alpha\beta}.
\]
Dies ist der Koeffizient von $z_\beta$ in
\[
 \sum_\beta r_{\alpha\beta}(\zeta)z_\beta
 =\sum_{\kappa,\mu}\xi_\kappa s_{\kappa\mu}(y)t_{\mu\alpha}(z)
 =\sum_{\kappa,\mu}\xi_\kappa t_{\kappa\mu}(z)s_{\mu\alpha}(y)
 =\sum_{\mu,\beta}s_{\mu\alpha}(y)r_{\mu\beta}(\xi)z_\beta .
\]
Daher ist
\[
\fnum{5} r_{\alpha\beta}(\zeta)=\sum_\mu s_{\mu\alpha}(y)r_{\mu\beta}(\xi)
\]
oder
\[
\fnum{5} R(\zeta)=S'(y)R(\xi).
\]
Setzt man endlich
\[
\fnum{6^*} \eta_\lambda=\sum_{\kappa,\mu}a_{\kappa\lambda\mu}\xi_\kappa z_\mu
=\sum_\mu r_{\lambda\mu}(\xi)z_\mu
=\sum_\kappa t_{\kappa\lambda}(z)\xi_\kappa,
\]
so wird
\[
 r_{\alpha\beta}(\eta)=\sum_\lambda a_{\lambda\alpha\beta}\eta_\lambda
 =\sum_{\kappa,\lambda}\xi_\kappa t_{\kappa\lambda}(z)a_{\lambda\alpha\beta}
\]
und daher
\[
 \sum_\alpha r_{\alpha\beta}(\eta)y_\alpha
 =\sum_{\kappa,\lambda}\xi_\kappa t_{\kappa\lambda}(z)s_{\lambda\beta}(y)
 =\sum_{\kappa,\lambda}\xi_\kappa s_{\kappa\lambda}(y)t_{\lambda\beta}(z)
 =\sum_{\alpha,\lambda}r_{\alpha\lambda}(\xi)t_{\lambda\beta}(z)y_\alpha,
\]
mithin
\[
\fnum{6} r_{\alpha\beta}(\eta)=\sum_\lambda r_{\alpha\lambda}(\xi)t_{\lambda\beta}(z)
\]
oder
\[
\fnum{6} R(\eta)=R(\xi)T(z).
\]

Bei passend geänderter Bezeichnung ergibt sich daraus:

\emph{Wenn die parastrophe Determinante nicht identisch verschwindet, und für den Wert $\xi$ des Parameters die Determinante der Matrix $R_\xi=R=(r_{\alpha\beta})$ von Null verschieden ist, so geht $R_\xi$ durch die Substitution $\xi_\alpha=\sum_\beta r_{\alpha\beta}x_\beta$ in $RT_x$ über, und durch die konjugierte Substitution $\xi_\beta=\sum_\alpha r_{\alpha\beta}x_\alpha$ in $S'_xR$.}

\emph{Unter jener Bedingung ist die Gruppendeterminante der antistrophen Gruppendeterminante identisch gleich, und bis auf den Faktor $|R|$ auch gleich der Funktion, in welche die parastrophe Determinante durch irgend eine jener beiden konjugierten Substitutionen übergeht; und alle drei Determinanten stimmen in den Elementarteilern überein.}

Sind $\xi,\eta,\zeta,\tau$ beliebige Parameter, so folgt aus (9.) § 1, daß $R_\eta^{-1}R_\xi$ mit $R_\tau^{-1}R'_\zeta$ vertauschbar ist, oder daß die Matrix
\[
 R_\zeta R_\eta^{-1}R_\xi R_\tau^{\prime -1}
\]
von $\zeta$ unabhängig ist.

\section*{§ 3.}

Wenn $|R(\xi)|$ nicht identisch verschwindet, so verschwindet keine der beiden Determinanten $|S(x)|$ oder $|T(x)|$ identisch, und es gibt eine Zahl $e$, wofür $S(e)=T(e)=E$ ist. Wenn aber identisch $|R(\xi)|=0$ ist, so füge ich zu den bisherigen Voraussetzungen die weitere hinzu, daß keine der beiden Determinanten
\[
\fnum{1} |S(x)|,\qquad |T(x)|
\]
identisch Null ist. Wählt man $z$ so, daß $|S(z)|$ und $|T(z)|$ beide von Null verschieden sind, so kann man (vergl. z. B. \emph{Mol.} Satz 1), weil $|S(z)|$ von Null verschieden ist, eine Zahl $e$ bestimmen, die der Gleichung $ze=z$ genügt. Ist ferner $x$ eine beliebige Größe, so kann man, weil $|T(z)|$ von Null verschieden ist, $y$ so bestimmen, daß $yz=x$ wird. Dann ist auch $xe=x$ und $zex=zx$, also weil $|S(z)|$ von Null verschieden ist, $ex=x$. Da ferner nach (2.) § 2 $S(z)=S(z)S(e)$ und $T(z)=T(e)T(z)$ ist, so ist $S(e)=T(e)=E$.

Der Satz des § 1 bleibt richtig, wenn weder $|S(x)|$ noch $|T(x)|$ identisch verschwindet, auch wenn identisch $|R(\xi)|=0$ ist. Denn setzt man $\sum_\kappa u_{\alpha\kappa}e_\kappa=y_\alpha$, so ist
\[
 \sum_{\beta,\kappa}u_{\alpha\beta}s_{\beta\kappa}(x)e_\kappa
 =\sum_{\gamma,\kappa}s_{\alpha\gamma}(x)u_{\gamma\kappa}e_\kappa,
\]
also
\[
 \sum_\beta u_{\alpha\beta}x_\beta=\sum_\gamma s_{\alpha\gamma}(x)y_\gamma=\sum_\beta t_{\alpha\beta}(y)x_\beta,
\]
und mithin $u_{\alpha\beta}=t_{\alpha\beta}(y)$.

Ist $|S(x)|$ für eine bestimmte Größe $x$ von Null verschieden, so kann man $y$ so bestimmen, daß $xy=e$ wird. Dann ist nach (2.) § 2 $T(y)T(x)=T(e)=E$. Ist also $|S(x)|$ für den Wert $x$ von Null verschieden, so ist auch $|T(x)|$ von Null verschieden und umgekehrt. Ist also $|S(x)|=0$, so ist auch $|T(x)|=0$. Jeder Primfaktor $\Phi(x)$ der Gruppendeterminante $|S(x)|$ ist daher auch in der antistrophen Determinante $|T(x)|$ enthalten und umgekehrt. In den beiden Zerlegungen
\[
\fnum{2} |S(x)|=\prod \Phi(x)^s,\qquad |T(x)|=\prod \Phi(x)^t
\]
treten genau dieselben Primfaktoren $\Phi(x)$ auf, doch können die Exponenten $s$ und $t$ verschieden sein, aber nur, wenn identisch $|R(\xi)|=0$ ist.

Das einfachste Beispiel für diese Möglichkeit liefert, wie mir \textsc{Molien} mitgeteilt hat, die einzige nicht kommutative Gruppe der Ordnung $n=3$, für die
\[
 x_1=y_1z_1,\qquad x_2=y_2z_2,\qquad x_3=y_1z_3+y_3z_2
\]
ist. Demnach ist
\[
 R(\xi)=
 \begin{pmatrix}\xi_1&0&\xi_3\\0&\xi_2&0\\0&\xi_3&0\end{pmatrix},\quad
 S(y)=
 \begin{pmatrix}y_1&0&0\\0&y_2&0\\0&y_3&y_1\end{pmatrix},\quad
 T(z)=
 \begin{pmatrix}z_1&0&0\\0&z_2&0\\z_3&0&z_2\end{pmatrix},
\]
also, in Elementarteiler zerlegt,
\[
 |S(x)|=x_1^2x_2,\qquad |T(x)|=x_1x_2^2,\qquad
 |uR(\xi)+vR'(\xi)|=-(\xi_1+\xi_2)\xi_3^2uv(u+v).
\]

Aus der Gleichung (2.) § 2 folgt durch wiederholte Anwendung
\[
 S_x^\nu=S(x^\nu),\qquad T_x^\nu=T(x^\nu).
\]
Unter $u,v,w$ verstehe ich hier stets gewöhnliche (nicht hyperkomplexe) Größen. Ist $g(u)$ eine ganze Funktion der Variabeln $u$, also $g((x))$ eine ganze Funktion der hyperkomplexen Größe $x$ selbst, so ist mithin
\[
\fnum{3} g(S_x)=S(g((x))),\qquad g(T_x)=T(g((x))).
\]
Wenn demnach einer der Ausdrücke $g((x))$, $g(S_x)$ oder $g(T_x)$ für alle Werte von $x$ verschwindet, so verschwinden auch die beiden anderen (\emph{Mol.} § 4). Die Gleichung niedrigsten Grades, der eine Matrix $S_x=S$ genügt, $g(S)=0$, wird ihre reduzierte Gleichung genannt. Die Funktion $g(u)$ erhält man, indem man die charakteristische Determinante von $S$, die Determinante der Matrix
\[
\fnum{4} uE-S(x)=S(ue-x),
\]
durch den größten gemeinsamen Teiler ihrer Unterdeterminanten $(n-1)$ten Grades dividiert. Nach (3.) genügen $S$ und $T$ derselben reduzierten Gleichung $g(S)=0$ und $g(T)=0$, und diese ist zugleich die Gleichung niedrigsten Grades, der eine variable Größe $x$ genügt, $g((x))=0$. Wenn also auch die Determinanten der beiden antistrophen Gruppen verschieden sein können, so stimmen doch ihre ersten Elementarteiler immer überein.

Die reduzierte Funktion $g(u)$ verschwindet für jede der charakteristischen Wurzeln von $S$, d. h. der Wurzeln der charakteristischen Gleichung $|uE-S|=0$. Damit ist aufs neue bewiesen, daß die beiden antistrophen Gruppendeterminanten genau dieselben Primfaktoren enthalten.

\section*{§ 4.}

Ist $\Theta(x)=|S(x)|$ die Gruppendeterminante, so ist nach (2.) § 2 $\Theta(yz)=\Theta(y)\Theta(z)$. In jedem Primfaktor $\Phi(x)$ von $\Theta(x)$ denke ich mir den konstanten Faktor so gewählt, daß $\Phi(e)=1$ ist. Ist dann $\Phi(x)$ ein solcher Faktor oder auch ein Produkt von mehreren, so ist auch
\[
\fnum{1} \Phi(yz)=\Phi(y)\Phi(z).
\]
Umgekehrt muß jede homogene Funktion $\Phi(x)$, welche diese Eigenschaft besitzt, ein Produkt von Primfaktoren von $\Theta(x)$ sein (\emph{Gr.} § 1). Denn bestimmt man $y$ so, daß $xy=\Theta(x)e$ wird, so werden $y_1,y_2,\ldots,y_n$ aus $n$ linearen Gleichungen mit der Determinante $\Theta(x)$ gefunden, und sind daher ganze Funktionen von $x_1,x_2,\ldots,x_n$. Ist $r$ der Grad von $\Phi(x)$, so ist dann nach (1.)
\[
 \Phi(x)\Phi(y)=\Theta(x)^r.
\]
Die Wurzeln $u_1,u_2,\ldots,u_r$ der Gleichung $\Phi(ue-x)=0$ nenne ich die charakteristischen Wurzeln von $\Phi(x)$. Dann ist
\[
\fnum{2} \Phi(ue-x)=u^r-\Phi_1(x)u^{r-1}+\Phi_2(x)u^{r-2}-\cdots\pm\Phi_r(x)
=(u-u_1)(u-u_2)\cdots(u-u_r).
\]
Die Summe der $r$ Wurzeln, $\Phi_1(x)$ oder
\[
\fnum{3} u_1+u_2+\cdots+u_r=\chi(x)=\sum_\alpha\chi_\alpha x_\alpha,
\]
nenne ich die Spur von $\Phi(x)$. Durch das System ihrer $n$ Koeffizienten, das ich ebenfalls mit $\chi$ bezeichne und den für $\Phi(x)$ charakteristischen Parameter nenne, ist die Funktion $\Phi(x)$ vollständig bestimmt (\emph{Mol.} § 3; \emph{Gr.} § 3). Denn ist
\[
 g(u)=a(u+v_1)(u+v_2)\cdots(u+v_m)
\]
eine ganze Funktion von $u$, so ist auch
\[
 y=g((x))=a(x+v_1e)(x+v_2e)\cdots(x+v_me)
\]
und mithin nach (1.)
\[
 \Phi(y)=a^r\Phi(x+v_1e)\Phi(x+v_2e)\cdots\Phi(x+v_me).
\]
Setzt man hierin
\[
 x=u_ie,
\]
so erhält man
\[
 \Phi(y)=g(u_1)g(u_2)\cdots g(u_r).
\]
Ersetzt man nun $g(u)$ durch $g(u)-v$, so erkennt man, daß $g(u_1),g(u_2),\ldots,g(u_r)$ die charakteristischen Wurzeln von $\Phi(y)$ sind. Daher ist
\[
\fnum{4} \chi(g((x)))=g(u_1)+g(u_2)+\cdots+g(u_r)
\]
und speziell
\[
\fnum{5} \chi(x^\nu)=u_1^\nu+u_2^\nu+\cdots+u_r^\nu,
\]
also für $\nu=0$
\[
\fnum{6} r=\chi(e)=\sum_\alpha\chi_\alpha e_\alpha.
\]
Nach den bekannten Relationen zwischen den Koeffizienten einer Gleichung und den Potenzsummen ihrer Wurzeln ist daher identisch für Primfunktionen ersten Grades
\[
\fnum{7} \chi(x)\chi(y)-\chi(xy)=0,
\]
für solche zweiten Grades
\[
\fnum{8} \chi(x)\chi(y)\chi(z)-\chi(x)\chi(yz)-\chi(y)\chi(zx)-\chi(z)\chi(xy)+\chi(xyz)+\chi(xzy)=0.
\]
Das allgemeine Gesetz für die Bildung dieser Gleichungen habe ich \emph{Gr.} § 3 auseinandergesetzt.

Insbesondere ist die Anzahl der verschiedenen linearen Primfunktionen gleich der Anzahl der Lösungen der Gleichungen (2.) § 1 durch gewöhnliche Zahlen (\emph{Ded.} (19.)). Wie ich in § 6 zeigen werde, sind diese Lösungen alle linear unabhängig.

Ist $\Theta(x)$ für einen bestimmten Wert $x$ von Null verschieden, so ist nach (1.) § 4
\[
 \Phi(x^{-1}yx+ue)=\Phi(x^{-1})\Phi(y+ue)\Phi(x)=\Phi(y+ue)
\]
und mithin $\chi(x^{-1}yx)=\chi(y)$, oder wenn man $y$ durch $xy$ ersetzt (\emph{Mol.} Satz 14; \emph{Gr.} § 3, (13.)),
\[
\fnum{9} \chi(xy)=\chi(yx).
\]
Die Matrix dieser symmetrischen bilinearen Form $F(\chi,x,y)$ ist
\[
\fnum{10} R_\chi=R'_\chi.
\]
Nennt man also $\xi$ einen symmetrischen Parameter, wenn es den Gleichungen $r_{\alpha\beta}(\xi)=r_{\beta\alpha}(\xi)$ oder $R_\xi=R'_\xi$ genügt, so ist $\chi$ ein solcher Parameter.

Eine Größe $x$ heißt eine invariante Größe der Gruppe $(\eps)$, wenn sie mit jeder Größe $y$ dieser Gruppe vertauschbar ist, $xy=yx$. Demnach ist
\[
 \sum_\gamma a_{\alpha\gamma\beta}x_\gamma=\sum_\gamma a_{\alpha\beta\gamma}x_\gamma,\qquad
 t_{\alpha\beta}(x)=s_{\alpha\beta}(x).
\]
Wird umgekehrt die Veränderlichkeit von $x$ durch die Gleichung $S(x)=T(x)$ beschränkt, so ist $x$ eine invariante Größe von $(\eps)$.

Auch die allgemeinere Gleichung $S(x)=T(y)$ oder
\[
 x_1E_1+\cdots+x_nE_n=y_1F_1+\cdots+y_nF_n
\]
kann nach (14.) § 1 nur bestehen, wenn $x=y$ eine invariante Größe ist. Die von den invarianten Größen gebildete Gruppe ist also der größte gemeinsame Divisor der beiden antistrophen Gruppen (3.) und (4.) § 2.

Ist $y$ unbeschränkt veränderlich, $x$ aber eine invariante Variable, so ist $S(x)=T(x)$ mit $S(y)$ vertauschbar, und daher lassen sich die charakteristischen Wurzeln $u_1,u_2,\ldots$ der Matrix $S(x)$ den charakteristischen Wurzeln $v_1,v_2,\ldots$ der Matrix $S(y)$ so zuordnen, daß $u_1v_1,u_2v_2,\ldots$ die charakteristischen Wurzeln von $S(x)S(y)=S(xy)$ sind. Mithin ist auch
\[
 \Phi(ue-xy)=(u-u_1v_1)(u-u_2v_2)\cdots(u-u_rv_r),
\]
wo $u_1,u_2,\ldots,u_r$ nur von $x$ abhängen, $v_1,v_2,\ldots,v_r$ nur von $y$. Setzt man $y=e$, so erkennt man, daß $u_1,u_2,\ldots,u_r$ die charakteristischen Wurzeln von $\Phi(x)$ sind, setzt man $x=e$, daß $v_1,v_2,\ldots,v_r$ die von $\Phi(y)$ sind.

Bei der Bestimmung der Art, wie diese Wurzeln in der obigen Zerlegung einander zugeordnet sind, beschränke ich mich auf den Fall, wo $\Phi(y)$ eine Primfunktion ist. Dann ist $\Phi(ue-y)$ als Funktion von $u$ irreduzibel, d. h. kann nicht als Produkt ganzer Funktionen von $u$ dargestellt werden, deren Koeffizienten rationale Funktionen der $n$ unbeschränkt veränderlichen Größen $y_1,y_2,\ldots,y_n$ sind. Da $u_1$ von diesen unabhängig ist, so ist die Funktion $\Phi(ue-u_1y)$ der Variabeln $u$ in demselben Sinne irreduzibel. Die Funktion $\Phi(ue-xy)$ hat mit ihr den Faktor $u-u_1v_1$ gemeinsam, und ist folglich mit ihr identisch $\Phi(ue-xy)=\Phi(ue-u_1y)$. Mithin ist $\chi(xy)=u_1\chi(y)$, also für $y=e$, $\chi(x)=ru_1$. Setzt man $u=0$, $y=e$, so wird $\Phi(x)=u_1^r$. So ergibt sich der Satz (\emph{Gr.} § 6):

\emph{Ist $\chi(x)$ die Spur des Primfaktors $r$ten Grades $\Phi(x)$ der Gruppendeterminante, ist $y$ eine unbeschränkt veränderliche, $x$ aber eine invariante Größe, so ist}
\[
\fnum{11} \chi(x)\chi(y)=r\chi(xy)=\chi(e)\chi(xy),
\]
\emph{und}
\[
\fnum{12} \Phi(x)=\left(\frac1r\chi(x)\right)^r,\qquad
 \Phi(ue-x)=\left(u-\frac1r\chi(x)\right)^r
\]
\emph{ist die $r$te Potenz einer linearen Funktion, und die $r$ charakteristischen Wurzeln von $\Phi(x)$ sind alle gleich $\frac1r\chi(x)$.}

\section*{§ 5.}

Ist $t=xy$, so ist $t+ux=x(y+ue)$ und folglich, wenn $\Phi(x)$ ein Faktor einer Potenz von $\Theta(x)$ ist, $\Phi(t+ux)=\Phi(x)\Phi(y+ue)$. Vergleicht man darin die Koeffizienten von $u^{r-1}$, so erhält man
\[
 \sum_\alpha \frac{\partial\Phi(x)}{\partial x_\alpha}t_\alpha=\Phi(x)\chi(y),
\]
also weil $t_\alpha=\sum_\beta s_{\alpha\beta}(x)y_\beta$ ist,
\[
 \sum_\alpha\frac{\partial\Phi(x)}{\partial x_\alpha}s_{\alpha\beta}(x)=\Phi(x)\chi_\beta .
\]
Auf diese Weise, und durch die Zerlegung $yx+ux=(y+ue)x$, erhält man (\emph{Gr.} § 5, (3.))
\[
\fnum{1} \sum_\kappa\frac{\partial\Phi(x)}{\partial x_\kappa}s_{\kappa\alpha}(x)
=\sum_\kappa\frac{\partial\Phi(x)}{\partial x_\kappa}t_{\kappa\alpha}(x)
=\Phi(x)\chi_\alpha .
\]
Setzt man also in der Formel $(5^*.)$ § 2 $\xi_\kappa=\frac{\partial\Phi}{\partial x_\kappa}$, so wird $\zeta_\alpha=\Phi(x)\chi_\alpha$, und nach (5.) und (6.) § 2 ist daher (\emph{Gr.} § 4, (1.))
\[
\fnum{2} S'(x)R\!\left(\frac{\partial\Phi}{\partial x}\right)
=R\!\left(\frac{\partial\Phi}{\partial x}\right)T(x)
=\Phi(x)R(\chi),
\]
wo die Matrix
\[
 R\!\left(\frac{\partial\Phi}{\partial x}\right)
 =\left(\sum_\kappa a_{\kappa\alpha\beta}\frac{\partial\Phi}{\partial x_\kappa}\right)
\]
ist. Geht diese in
\[
 R_1u^{r-1}+R_2u^{r-2}+\cdots+R_r
\]
über, falls man $x$ durch $ue-x$ ersetzt, so ist nach (2.) § 4
\[
 (R_1u^{r-1}+\cdots+R_r)(Eu-T)
=R_\chi(u^r-\Phi_1u^{r-1}+\Phi_2u^{r-2}-\cdots\pm\Phi_r).
\]
Vergleicht man die Koeffizienten von $u^r,u^{r-1},\ldots,u^0$, so erhält man $r+1$ Gleichungen. Multipliziert man diese rechts mit $T^r,T^{r-1},\ldots,T^0$ und addiert sie, so ergibt sich (\emph{Gr.} § 4, (5.))
\[
\fnum{3} R_\chi(T^r-\Phi_1T^{r-1}+\Phi_2T^{r-2}-\cdots\pm\Phi_r)=0
\]
oder einfacher
\[
 R_\chi\Phi(xE-eT)=0.
\]
In derselben Weise ergibt sich $\Phi(xE-eS')R_\chi=0$, also wenn man die konjugierten Matrizen nimmt, nach (10.) § 4
\[
\fnum{3} R_\chi\Phi(xE-eS)=0,\qquad R_\chi\Phi(xE-eT)=0.
\]
Nach (3.) § 3 kann man für $\Phi(xE-eS)$ auch $S(\Phi(x-((x))e))$ schreiben, d. h. $S(y)$, wo $y=g((x))$ und $g(u)=\Phi(x-ue)$ ist.

Zu der Formel (2.), worauf die folgende Entwicklung wesentlich beruht, gelangt man auch, indem man $t=xyz$ (oder $yzx$) setzt, und in der Gleichung
\[
 \sum_\kappa\frac{\partial\Phi(x)}{\partial x_\kappa}t_\kappa=\Phi(x)\chi(yz),
\]
worin
\[
 t_\kappa=\sum_{\lambda,\mu,\alpha,\beta}a_{\kappa\lambda\mu}x_\lambda a_{\mu\alpha\beta}y_\alpha z_\beta,\qquad
 \chi(yz)=\sum_{\alpha,\beta}r_{\alpha\beta}(\chi)y_\alpha z_\beta
\]
ist, die Koeffizienten von $y_\alpha z_\beta$ vergleicht.

\section*{§ 6.}

Die kleinste Anzahl unabhängiger linearer Verbindungen der Variabeln, wodurch sich eine Funktion oder ein System von Funktionen ausdrücken läßt, nenne ich seinen \emph{linearen Rang}, oder auch nur seinen \emph{Rang}, sobald eine Verwechslung mit dem Begriffe des Ranges einer Determinante oder Matrix ausgeschlossen ist. Der lineare Rang $m$ einer quadratischen Funktion
\[
\fnum{1} \sum_{\alpha,\beta}r_{\alpha\beta}(\chi)x_\alpha x_\beta
=F(\chi,x,x)=\chi(x^2)=\Phi_1^2-2\Phi_2
\]
ist gleich dem Range der aus ihren Koeffizienten gebildeten symmetrischen Matrix $R_\chi$, und gleich dem linearen Range des Systems der $n$ linearen Funktionen
\[
\fnum{2} \sum_\beta r_{\alpha\beta}(\chi)x_\beta,
\]
den halben Ableitungen der quadratischen Funktion. Durch irgend $m$ gerade dieser linearen Funktionen, die von einander unabhängig sind, läßt sich $\chi(x^2)$ darstellen. Ebenso kann man ihre Kovariante, die symmetrische bilineare Form
\[
\fnum{3} \sum_{\alpha,\beta}r_{\alpha\beta}(\chi)x_\alpha y_\beta
=F(\chi,x,y)=\chi(xy)=\chi(yx)
\]
durch jene $m$ Variabeln ausdrücken und durch die $m$, die unter den $n$ Variabeln $\sum_\beta r_{\alpha\beta}(\chi)y_\beta$ von einander unabhängig sind.

Die trilineare Funktion $\chi(xyz)=\chi((xy)z)$ ist eine bilineare Funktion der Koordinaten von $xy$ und von $z$, läßt sich also durch die $m$ Variabeln ausdrücken, die unter den $n$ Variabeln $\sum_\beta r_{\alpha\beta}z_\beta$ von einander unabhängig sind. Sie ist gleich $\chi(x(yz))=\chi(z(xy))=\chi((zx)y)$, also auch nur abhängig von den Variabeln $\sum_\beta r_{\alpha\beta}x_\beta$ und den Variabeln $\sum_\beta r_{\alpha\beta}y_\beta$. Dasselbe gilt von $\chi(xyzt)=\chi(yztx)=\chi(ztxy)=\chi(txyz)$. Daher ist auch
\[
 \chi(x^\nu)=u_1^\nu+u_2^\nu+\cdots+u_r^\nu
\]
nur von den Variabeln (2.) abhängig, also auch das Produkt $\Phi(x)=u_1u_2\cdots u_r$, das eine ganze Funktion jener Potenzsummen ist. Umgekehrt lassen sich durch die linearen Verbindungen der Variabeln $x_1,x_2,\ldots,x_n$, von denen $\Phi(x)$ abhängt, auch die in der Entwicklung von $\Phi(ue-x)$ auftretenden Funktionen $\Phi_1(x),\Phi_2(x),\ldots$ ausdrücken, also auch $\chi(x^2)$.

\emph{I. Ist $\Phi(x)$ ein Produkt von Primfaktoren der Gruppendeterminante, und $\chi(x)$ die Spur von $\Phi(x)$, so ist der lineare Rang von $\Phi(x)$ gleich dem der quadratischen Funktion $\chi(x^2)$, also gleich dem Range der Matrix $R_\chi$, und die Funktion $\Phi(x)$ läßt sich durch die Ableitungen von $\chi(x^2)$ ausdrücken.}

Jetzt seien $\Phi,\Phi',\Phi'',\ldots$ die $k$ verschiedenen Primfaktoren von $\Theta$, und $\chi,\chi',\chi'',\ldots$ ihre Spuren. Sind dann $C,C',C'',\ldots$ irgendwelche Matrizen $n$ten Grades, so will ich zeigen: Eine Relation
\[
\fnum{4} CR_\chi+C'R_{\chi'}+C''R_{\chi''}+\cdots=0
\]
kann nur bestehen, wenn einzeln
\[
 CR_\chi=C'R_{\chi'}=C''R_{\chi''}=\cdots=0
\]
ist. Denn nach (2.) § 5 ist
\[
 R_\chi=R\!\left(\frac{\partial l(\Phi)}{\partial x}\right)T(x)
\]
und mithin, da $|T(x)|$ von Null verschieden ist,
\[
 CR\!\left(\frac{\partial l(\Phi)}{\partial x}\right)
+C'R\!\left(\frac{\partial l(\Phi')}{\partial x}\right)
+C''R\!\left(\frac{\partial l(\Phi'')}{\partial x}\right)+\cdots=0.
\]
Nun sei $\Psi=\Phi'\Phi''\cdots$ das Produkt der $k-1$ von $\Phi$ verschiedenen Primfunktionen. Multipliziert man mit $\Phi\Psi$, so erkennt man, daß die Elemente der Matrix $\Psi(x)CR\!\left(\frac{\partial\Phi}{\partial x}\right)$ alle durch die Funktion $r$ten Grades $\Phi(x)$ teilbar sind, also weil $\Phi$ und $\Psi$ teilerfremd sind, auch die der Matrix $CR\!\left(\frac{\partial\Phi}{\partial x}\right)$. Da diese aber nur vom $(r-1)$ten Grade sind, so müssen sie Null sein. Mithin ist
\[
 0=CR\!\left(\frac{\partial l(\Phi)}{\partial x}\right)T(x)=CR_\chi.
\]
Zu diesem Ergebnis kann man auch mittels der Formel (3.) § 5 gelangen. Nach dieser ist $R_\chi\Phi(xE-eS)=0$, und weil je zwei ganze Funktionen der Matrix $S$ mit einander vertauschbar sind, auch
\[
 R_\chi\Psi(xE-eS)=R_\chi\Phi''(xE-eS)\Phi'(xE-eS)\cdots=0.
\]
Folglich ist auch
\[
 CR_\chi\bigl(\Psi(xE-eS)+\Phi(xE-eS)\bigr)=0.
\]
Die Determinante der in der Klammer stehenden Matrix, einer Funktion von $S$, ist von Null verschieden, weil die Funktion $\Psi(x-ue)+\Phi(x-ue)$ der Variabeln $u$ für keine Wurzel der charakteristischen Gleichung $\Theta(x-ue)=0$ verschwindet. Folglich muß $CR_\chi=0$ sein.

Sind $m,m',m'',\ldots$ die Rangzahlen der Matrizen $R_\chi,R_{\chi'},R_{\chi''},\ldots$, so sind unter den Variabeln $\sum_\alpha r_{\alpha\beta}(\chi)x_\beta$ genau $m$ unabhängig, unter den Variabeln $\sum_\alpha r_{\alpha\beta}(\chi')x_\beta$ genau $m'$ unabhängig, unter den Variabeln $\sum_\alpha r_{\alpha\beta}(\chi'')x_\beta$ genau $m''$, usw. Die so erhaltenen $m+m'+m''+\cdots$ linearen Verbindungen der $n$ Variabeln $x_1,x_2,\ldots,x_n$ sind aber alle unter einander unabhängig. Denn eine in $x_1,x_2,\ldots,x_n$ identische Beziehung
\[
 \sum_\beta c_\beta\left(\sum_\alpha r_{\alpha\beta}(\chi)x_\alpha\right)
+\sum_\beta c'_\beta\left(\sum_\alpha r_{\alpha\beta}(\chi')x_\alpha\right)
+\sum_\beta c''_\beta\left(\sum_\alpha r_{\alpha\beta}(\chi'')x_\alpha\right)+\cdots=0
\]
kann nur bestehen, wenn die $k$ Teilsummen einzeln verschwinden. Die $n$ Gleichungen
\[
 \sum_\alpha c_\alpha r_{\alpha\beta}(\chi)
+\sum_\alpha c'_\alpha r_{\alpha\beta}(\chi')
+\sum_\alpha c''_\alpha r_{\alpha\beta}(\chi'')+\cdots=0
\]
haben nämlich die Gestalt (4.), falls $C$ eine Matrix ist, worin eine Zeile aus den Elementen $c_1,c_2,\ldots,c_n$ besteht, während die Elemente der anderen Zeilen verschwinden.

Da $\sum_\alpha x_\alpha\chi_\alpha=\sum_\alpha e_\alpha\bigl(\sum_\beta r_{\alpha\beta}(\chi)x_\beta\bigr)$ eine lineare Verbindung der $n$ Variabeln (2.) ist, so sind auch die $k$ Funktionen
\[
\fnum{5} \chi(x),\quad \chi'(x),\quad \chi''(x),\quad\ldots
\]
linear unabhängig.

Ferner läßt sich $\chi(x^2)$ als eine quadratische Funktion von $m$ unabhängigen Variabeln mit nicht verschwindender Determinante ausdrücken, $\chi'(x^2)$ als eine Funktion von $m'$ Variabeln vom Range $m'$, usw., und diese $m+m'+m''+\cdots$ Variabeln sind alle unter einander unabhängig. Sind also $c,c',c'',\ldots$ Konstanten, so ist der Rang der quadratischen Funktion
\[
 c\chi(x^2)+c'\chi'(x^2)+c''\chi''(x^2)+\cdots
\]
gleich der Summe der Rangzahlen der einzelnen Summanden. Der Rang von $c\chi(x^2)$ ist $0$ oder $m$, je nachdem $c=0$ ist, oder nicht.

\emph{II. Der Rang der Matrix $cR_\chi+c'R_{\chi'}+c''R_{\chi''}+\cdots$ ist gleich der Summe der Rangzahlen der Matrizen $cR_\chi,c'R_{\chi'},c''R_{\chi''},\ldots$.}

\emph{Der lineare Rang eines Produktes von Primfaktoren der Gruppendeterminante ist gleich der Summe der Rangzahlen ihrer verschiedenen Primfaktoren.}

\section*{§ 7.}

Die Spuren der Determinanten $|S(x)|$ und $|T(x)|$ bezeichne ich mit
\[
\fnum{1} \sigma(x)=\sum_\alpha \sigma_\alpha x_\alpha=\sum_\kappa s_{\kappa\kappa}(x),\qquad
\tau(x)=\sum_\alpha \tau_\alpha x_\alpha=\sum_\kappa t_{\kappa\kappa}(x).
\]
Ihre Koeffizienten sind
\[
\fnum{2} \sigma_\alpha=\sum_\kappa a_{\kappa\alpha\kappa},\qquad
\tau_\alpha=\sum_\kappa a_{\kappa\kappa\alpha},
\]
und nach (2.) § 3 ist
\[
\fnum{3} \sigma=s\chi+s'\chi'+s''\chi''+\cdots=\sum s\chi,
\qquad
\tau=t\chi+t'\chi'+t''\chi''+\cdots=\sum t\chi.
\]
Daher haben nach Satz II, § 6 die Matrizen $R_\sigma$ und $R_\tau$ beide den Rang $m+m'+m''+\cdots$.

Eine Gruppe $(\eps)$, wofür die Determinante der symmetrischen Matrix
\[
\fnum{4} R_\sigma=P=(p_{\alpha\beta})
\]
von Null verschieden ist, nenne ich eine \textsc{Dedekind}sche Gruppe (\textsc{Ded}. (27.)), da \textsc{Dedekind} zuerst wenigstens für die kommutativen Gruppen die Bedeutung dieser Bedingung erkannt hat. Für eine solche ist nach dem Satze des § 2 $\Theta(x)=|S(x)|=|T(x)|$, also $\sigma=\tau$, und der Rang von $R_\sigma$
\[
 n=m+m'+m''+\cdots=\sum m.
\]
Eine \textsc{Dedekind}sche Gruppe kann also auch als eine solche definiert werden, deren Ordnung dem linearen Range ihrer Determinante gleich ist. Ist
\[
 \psi=\chi+\chi'+\chi''+\cdots=\sum\chi
\]
die Spur des Produkts der $k$ verschiedenen Primfaktoren von $\Theta$
\[
 \Psi(x)=\Phi\Phi'\Phi''\cdots=\prod\Phi(x),
\]
so hat nach Satz II, § 6 auch $R_\psi$ den Rang $\sum m=n$. Demnach ist $|R_\psi|$ von Null verschieden. Nach (3.) § 5 ist aber
\[
 R_\psi\Psi(Se-Ex)=0,
\]
und folglich ist
\[
\fnum{5} \Psi(Se-Ex)=0.
\]
Da die Funktion
\[
\fnum{6} g(u)=\Psi(ue-x)=\prod\Phi(ue-x)
\]
keine mehrfachen Faktoren hat, so ist mithin $g(S)=0$ die reduzierte Gleichung für die Matrix $S$ (vergl. \textsc{Mol}. Satz 24).

Dies kann man auch so einsehen: Nach (2.) § 5 ist
\[
 S'R\!\left(\frac{\partial\Theta}{\partial x}\right)=\Theta R_\sigma.
\]
Nun ist $\Theta S'^{-1}=\overline{S'}$ die zu $S'$ adjungierte Matrix, und demnach ist (\textsc{Ded}. (63.))
\[
\fnum{7} \overline{S'}=R\!\left(\frac{\partial\Theta}{\partial x}\right)R_\sigma^{-1}.
\]
Die Elemente der Matrix $\overline{S'}$ sind die Unterdeterminanten $(n-1)$ten Grades von $\Theta$. Diese sind also lineare Verbindungen der Ableitungen von $\Theta$ nach den $n$ Variabeln $x_1,x_2,\ldots,x_n$. Umgekehrt sind die Ableitungen einer Determinante lineare Verbindungen ihrer Unterdeterminanten. Daher ist der größte gemeinsame Divisor der Unterdeterminanten gleich dem der Ableitungen, also wenn man $x$ durch $ue-x$ ersetzt, gleich
\[
 \prod\bigl(\Phi(ue-x)\bigr)^{s-1}.
\]
Nach der in § 3 erwähnten Regel erhält man aber die reduzierte Gleichung $g(S)=0$, indem man
\[
 \Theta(ue-x)=\prod\bigl(\Phi(ue-x)\bigr)^s
\]
durch jenen Ausdruck dividiert, und folglich ist $g(u)$ die Funktion (6.). Die Wurzeln $u_1,u_2,\ldots,u_p$ der Gleichung $g(u)=0$ sind die $p=r+r'+r''+\cdots$ verschiedenen unter den $n$ Wurzeln $u_1,u_2,\ldots,u_n$ der Gleichung $\Theta(ue-x)=0$.

\section*{§ 8.}

Sind $A$ und $B$ zwei vertauschbare Matrizen $n$ten Grades, so lassen sich ihre charakteristischen Wurzeln $a_1,a_2,\ldots,a_n$ und $b_1,b_2,\ldots,b_n$ einander so zuordnen, daß $f(a_1,b_1),f(a_2,b_2),\ldots,f(a_n,b_n)$ die charakteristischen Wurzeln der Matrix $f(A,B)$ sind, und diese Zuordnung ist von der Wahl der ganzen Funktion $f(u,v)$ unabhängig. In § 3 der vorstehenden Arbeit habe ich ferner bewiesen:

\emph{I. Zerfallen die charakteristischen Determinanten von zwei (oder mehr) mit einander vertauschbaren Matrizen $A$ und $B$ in lauter lineare Elementarteiler, so hat jede Matrix $f(A,B)$ dieselbe Eigenschaft.}

\emph{II. Ist außerdem immer $b_\kappa=b_\lambda$, falls $a_\kappa=a_\lambda$ ist, so ist $B$ eine ganze Funktion von $A$.}

Nun ist
\[
\fnum{1} g(u,x)=\prod\Phi(ue-x)
\]
eine ganze Funktion $p$ten Grades von $u$, deren Koeffizienten ganze Funktionen der $n$ Variabeln $x_1,x_2,\ldots,x_n$ sind. Die $p$ Wurzeln $u_1,u_2,\ldots,u_p$ der Gleichung $g(u,x)=0$ sind alle unter einander verschieden. Ist $X=S(x)$, so sind die Elementarteiler von $|uE-X|$ alle linear und $X$ genügt der Gleichung $g(X,x)=0$.

Diese in $x_1,x_2,\ldots,x_n$ identische Gleichung wird daher auch durch jedes spezielle Wertsystem der Variabeln befriedigt. Für ein solches brauchen $u_1,u_2,\ldots,u_p$ nicht alle verschieden, die Elementarteiler von $|uE-X|$ nicht alle linear, und es braucht $g(X,x)=0$ nicht die reduzierte Gleichung für $X$ zu sein. Wenn aber $x$ so gewählt wird, daß $u_1,u_2,\ldots,u_p$ verschieden sind, so ist $g(X,x)=0$ die reduzierte Gleichung für die Matrix $X$. Denn mehr als $p$ verschiedene Wurzeln kann die Gleichung $|uE-X|=0$ nicht haben, und ist $\psi(X)=0$, so muß $\psi(u)$ für jede Wurzel dieser charakteristischen Gleichung verschwinden. Da dann die Gleichung $g(u,x)=0$ keine mehrfache Wurzel hat, so sind die Elementarteiler von $|uE-X|$ alle linear.

Sei $x$ so gewählt wie eben, und sei $y$ eine mit $x$ vertauschbare Größe. Dann ist auch die Matrix $Y=S(y)$ mit $X=S(x)$ vertauschbar. Daher ist die Determinante
\[
\fnum{2} |uX+vY+wE|=|S(ux+vy+we)|=\prod\bigl(\Phi(ux+vy+we)\bigr)^s
\]
ein Produkt von linearen Faktoren $uu_\alpha+vv_\alpha+w$, worin $u_\alpha$ und $v_\alpha$ zugeordnete charakteristische Wurzeln der beiden vertauschbaren Matrizen $X$ und $Y$ sind, und mithin ist auch
\[
\fnum{3} \Phi(ux+vy+we)=\prod_\alpha^{r}(uu_\alpha+vv_\alpha+w)
\]
ein Produkt linearer Funktionen von $u,v,w$. Setzt man $v=0$ oder $u=0$, so erkennt man: den $r$ charakteristischen Wurzeln $u_1,\ldots,u_r$ von $\Phi(x)$ sind die $r$ Wurzeln $v_1,\ldots,v_r$ von $\Phi(y)$ in einer bestimmten Reihenfolge zugeordnet, und folglich auch den $p$ Wurzeln $u_1,\ldots,u_p$ von $g(u,x)=0$ die $p$ Wurzeln $v_1,\ldots,v_p$ von $g(v,y)=0$. Da ferner $u_1,\ldots,u_p$ verschieden sind, so zeigen die Formeln (2.) und (3.), daß je zwei gleichen Wurzeln $u_\kappa=u_\lambda$ unter den $n$ charakteristischen Wurzeln $u_1,\ldots,u_n$ von $X$ gleiche Wurzeln $v_\kappa=v_\lambda$ von $Y$ zugeordnet sind.

Ist $l$ eine Konstante, so sind $v_1-lu_1,\ldots,v_p-lu_p$ die Wurzeln der Gleichung $g(w,y-lx)=0$. Ist $l$ nicht gleich einem der Brüche $\frac{v_\alpha-v_\beta}{u_\alpha-u_\beta}$, so sind jene $p$ Wurzeln alle verschieden, und folglich sind die Elementarteiler der charakteristischen Determinante der Matrix $Z=Y-lX$ alle linear. Da $X$ und $Z$ vertauschbar sind, und $Y=lX+Z$ eine Funktion von $X$ und $Z$ ist, so sind nach Satz I auch die Elementarteiler von $|vE-Y|$ alle linear. Nach Satz II ist daher $Y$ eine ganze Funktion von $X$ (vergl. \textsc{Mol}. Satz 27).

Die Gleichung niedrigsten Grades, der $x$ genügt, hat den Grad $p$. Daher sind $x^0,x^1,\ldots,x^{p-1}$ linear unabhängig, jede ganze Funktion von $x$ aber von diesen $p$ abhängig. Die linearen Gleichungen $xy-yx=0$ zwischen den $n$ Unbekannten $y_1,y_2,\ldots,y_n$ haben daher genau $p$ unabhängige Lösungen. Die Matrix ihrer Koeffizienten $S(x)-T(x)$ hat daher den Rang $n-p$.

\emph{III. Sind für eine bestimmte Größe $x$ einer \textsc{Dedekind}schen Gruppe die $r+r'+r''+\cdots$ Wurzeln der Gleichung $\prod\Phi(ue-x)=0$ alle verschieden, so hat die Matrix $S(x)-T(x)$ den Rang $n-(r+r'+r''+\cdots)$. Dann ist jede mit $x$ vertauschbare Größe $y$ eine ganze Funktion von $x$, und es sind für jede solche Größe $y$ die Elementarteiler von $|vE-S(y)|$ alle linear.}

\section*{§ 9.}

Die Theorie der \textsc{Dedekind}schen Gruppen werde ich in den §§ 13 bis 16 zu Ende führen. Jetzt kehre ich zur allgemeinen Theorie zurück. An Stelle der Basis $\eps_1,\eps_2,\ldots,\eps_n$ der Gruppe $(\eps)$ kann man eine neue Basis durch eine lineare Substitution
\[
\fnum{1} \eps_\beta=\sum_\alpha c_{\alpha\beta}\bar\eps_\alpha
\]
einführen, worin die Determinante der Matrix $C=(c_{\alpha\beta})$ von Null verschieden ist. Ist dann
\[
\fnum{2} \sum \eps_\alpha x_\alpha=\sum\bar\eps_\alpha\bar x_\alpha,
\]
so ist
\[
\fnum{3} \bar x_\alpha=\sum_\beta c_{\alpha\beta}x_\beta.
\]
Die Größen $\frac{\partial\Phi}{\partial x_\alpha},\xi_\alpha,\chi_\alpha,\sigma_\alpha$ sind den Grundzahlen $\eps_\alpha$ kogredient, die Größen $y_\alpha,z_\alpha$ sind den Koordinaten $x_\alpha$ kogredient, den Grundzahlen $\eps_\alpha$ kontragredient. Ferner ist
\[
\fnum{4}
\begin{aligned}
S(x)&=C^{-1}\bar S(\bar x)C,
& T(x)&=C^{-1}\bar T(\bar x)C,\\
R(\xi)&=C'\bar R(\bar\xi)C,
& R'(\xi)&=C'\bar R'(\bar\xi)C.
\end{aligned}
\]
Demnach sind die Exponenten und die Grade der Elementarteiler, worin die Determinanten der Matrizen $R(\xi),S(x),T(x)$ zerfallen, Invarianten der Gruppe. Dasselbe gilt aber für die Determinanten der von je $2n$ Variabeln abhängenden Matrizen
\[
\fnum{5} S(x)+T(y),\qquad R(\xi)+R'(\eta),
\]
insbesondere auch für die Elementarteiler der Determinante der Matrix der bilinearen Form $uF(\xi,y,z)+vF(\xi,z,y)$ von $y$ und $z$
\[
\fnum{6} uR(\xi)+vR'(\xi).
\]
Zur Erläuterung dieser Bemerkung betrachte ich einige Beispiele aus der Arbeit des Hrn. \textsc{Study}, \emph{Über Systeme von complexen Zahlen}, Göttinger Nachrichten 1889. Ist (\textsc{St}. IX)
\[
\fnum{7} R(\xi)=
\begin{pmatrix}
\xi_1&\xi_2&\xi_3&\xi_4\\
\xi_2&\xi_4&\xi_4&0\\
\xi_3&-\xi_4&c\xi_4&0\\
\xi_4&0&0&0
\end{pmatrix},
\]
so sind jene Elementarteiler
\[
\fnum{8} \xi_4^3\xi_4(u+v)(u+v)\bigl((u-v)^2+c(u+v)^2\bigr),
\]
worin der letzte Faktor für $c=0$ ein quadratischer Elementarteiler ist, sonst in zwei verschiedene lineare Elementarteiler zerfällt. Diese Formel zeigt, daß $c$ eine (nicht numerische) Invariante der Gruppe ist.

Ist ferner (\textsc{St}. XIV)
\[
\fnum{9} R(\xi)=
\begin{pmatrix}
\xi_1&\xi_2&\xi_3&\xi_4\\
\xi_2&0&\xi_4&0\\
\xi_3&-\xi_4&0&0\\
\xi_4&0&0&0
\end{pmatrix},
\]
so sind die Elementarteiler von $|uR(\xi)+vR'(\xi)|$
\[
\fnum{10} \xi_4^2\xi_4^2(u+v)(u+v)(u-v)(u-v),
\]
also wesentlich von den obigen (für $c=0$) verschieden. Doch will ich hier auf die Klassifikation der Gruppen nicht näher eingehen.

Es ist möglich, daß die Größen (2.) auch dann noch eine Gruppe $(\mathfrak S)$ bilden, wenn die Koordinaten gewissen linearen Relationen unterworfen werden. Dann wird $(\mathfrak S)$ eine Untergruppe von $(\eps)$ genannt. Ihre Grundzahlen $\mathfrak s_1,\mathfrak s_2,\ldots$ sind lineare Verbindungen der $n$ Grundzahlen $\eps_\alpha$, und damit $(\mathfrak S)$ eine Gruppe sei, muß jedes der Produkte $\mathfrak s_\alpha\mathfrak s_\beta$ eine lineare Verbindung der Größen $\mathfrak s_1,\mathfrak s_2,\ldots$ sein.

Ist auch jedes der Produkte $\eps_\alpha\mathfrak s_\beta$ eine lineare Verbindung der Größen $\mathfrak s_1,\mathfrak s_2,\ldots$, so heißt $(\mathfrak S)$ eine invariante Untergruppe\footnote{Ist $\mathfrak H$ eine endliche Gruppe der Ordnung $n$ und $\mathfrak G$ eine Untergruppe der Ordnung $n/m$, ist $R$ ein Element von $\mathfrak H$ und $P$ ein Element von $\mathfrak G$, so bilden die hyperkomplexen Größen $\sum R x_R$ eine Gruppe $(\eps)$ der Ordnung $n$, und die Größen $\sum P x_P$ eine Untergruppe $(\mathfrak S)$. Sie ist nicht etwa eine invariante Untergruppe von $(\eps)$, wenn $\mathfrak G$ eine solche von $\mathfrak H$ ist. Ist aber $\mathfrak H=\mathfrak G+\mathfrak G A+\mathfrak G B+\cdots$, so erhält man dann eine invariante Untergruppe der Ordnung $n-m$ von $(\eps)$, wenn man die Veränderlichkeit der $n$ Koordinaten $x_R$ durch die $n$ Gleichungen $\sum_P x_P=0$, $\sum_P x_{AP}=0$, $\sum_P x_{BP}=0$, $\cdots$ einschränkt. Die ihr komplementäre mit $(\eps)$ homomorphe Gruppe $(\eta)$ steht zu der Gruppe $\mathfrak H/\mathfrak G$ in derselben Beziehung wie $(\eps)$ zu $\mathfrak H$.} von $(\eps)$. Dann definieren die Formeln (2.) § 1 auch dann noch eine Gruppe $(\eta)$, wenn man die $n$ Grundzahlen $\eps_\alpha$ den linearen Relationen $\mathfrak s_1=0,\mathfrak s_2=0,\ldots$ unterwirft, oder wenn man je zwei Größen von $(\eps)$ als gleich mod. $(\mathfrak S)$ betrachtet, deren Differenz der Gruppe $(\mathfrak S)$ angehört. Sie wird die der invarianten Untergruppe $(\mathfrak S)$ komplementäre, mit $(\eps)$ homomorphe Gruppe genannt (Begleitendes Zahlensystem bei \textsc{Molien}).

Um dies einzusehen, transformiere man die Basis so, daß $\mathfrak s_1,\mathfrak s_2,\ldots$ die Grundzahlen $\eps_{m+1},\ldots,\eps_n$ werden. Dann wird die Untergruppe $(\mathfrak S)$ durch die $m$ linearen Gleichungen $x_1=0,\ldots,x_m=0$ zwischen den Koordinaten bestimmt. Damit $(\mathfrak S)$ eine Untergruppe sei, ist notwendig und hinreichend, daß $a_{\alpha\beta\gamma}=0$ ist, falls $\alpha\leq m$, $\beta>m$ und $\gamma>m$ ist; damit $(\mathfrak S)$ eine invariante Untergruppe sei, daß $a_{\alpha\beta\gamma}=0$ ist, falls $\alpha\leq m$ und auch nur eine der beiden Zahlen $\beta$ oder $\gamma>m$ ist (\textsc{Mol}. Satz 2). Dann definieren die Gleichungen (2.) § 1 auch dann noch eine Gruppe, wenn $\eps_{m+1}=\cdots=\eps_n=0$ gesetzt wird. Denn in den Formeln (4.) § 1 hängen $x_1,\ldots,x_m$ nur von $y_1,\ldots,y_m$ und $z_1,\ldots,z_m$ ab. In den beiden Matrizen $S(x)$ und $T(x)$ verschwinden daher alle Elemente, welche die ersten $m$ Zeilen mit den letzten $n-m$ Spalten gemeinsam haben. Demnach ist
\[
 S=\begin{pmatrix}S_1&0\\ S_{21}&S_2\end{pmatrix},
 \qquad
 T=\begin{pmatrix}T_1&0\\ T_{21}&T_2\end{pmatrix},
\]
wo z. B. die Teilmatrix $S_1$ nur die ersten $m$ Zeilen und Spalten von $S$ enthält. Aus $S(y)S(z)=S(yz)$ folgt daher $S_1(y)S_1(z)=S_1(yz)$, und da $S_1(x)=x_1E_1+\cdots+x_mE_m$ nur von $x_1,\ldots,x_m$ abhängt, so definieren die Gleichungen
\[
 \eps_\beta\eps_\gamma=\sum_{\alpha=1}^{m}a_{\alpha\beta\gamma}\eps_\alpha\qquad(\beta,\gamma=1,2,\ldots,m)
\]
eine Gruppe $(\eta)$, die aus $(\eps)$ hervorgeht, indem man $\eps_{m+1}=\cdots=\eps_n=0$ setzt. Weil $|S|=|S_1||S_2|$ und $|T|=|T_1||T_2|$ ist, so genügt $(\eta)$ auch den Bedingungen (1.) § 3. Ferner ist die Determinante der Gruppe $(\eps)$ durch die Determinante jeder mit $(\eps)$ homomorphen Gruppe $(\eta)$ teilbar, und dasselbe gilt von den mit $(\eps)$ und $(\eta)$ antistrophen Gruppen.

Ist außerdem auch $a_{\alpha\beta\gamma}=0$, wenn $\alpha>m$ und eine der beiden Zahlen $\beta$ oder $\gamma\leq m$ ist, so zerfällt $(\eps)$ in zwei Gruppen, deren jede sowohl eine invariante Untergruppe von $(\eps)$ wie eine mit $(\eps)$ homomorphe Gruppe ist.

\section*{§ 10.}

Sind $m,m',m'',\ldots$ die Rangzahlen der $k$ Matrizen $R_\chi,R_{\chi'},R_{\chi''},\ldots$, so hängt $\Phi$ nur von den $m$ unabhängigen unter den $n$ linearen Funktionen $\sum_\beta r_{\alpha\beta}(\chi)x_\beta$ ab, $\Phi'$ nur von den $m'$ unabhängigen unter den Funktionen $\sum_\beta r_{\alpha\beta}(\chi')x_\beta$ usw., und diese $m+m'+m''+\cdots$ Variabeln sind alle unter einander unabhängig. Man kann daher die Grundzahlen so wählen, daß $\Phi$ nur von $x_1,\ldots,x_m$ abhängt, $\Phi'$ nur von $x_{m+1},\ldots,x_{m+m'}$, usw. Dann hängt auch $\chi(x^2)=\sum r_{\alpha\beta}(\chi)x_\alpha x_\beta$ nur von $x_1,\ldots,x_m$ ab. Daher ist $r_{\alpha\beta}(\chi)=0$, wenn $\alpha$ oder $\beta>m$ ist, und da $R_\chi$ den Rang $m$ hat, so ist die Determinante $m$ten Grades
\[
 |r_{\kappa\lambda}(\chi)|\qquad(\kappa,\lambda=1,2,
 \ldots,m)
\]
von Null verschieden. Nach § 6 hängt ferner auch die trilineare Funktion
\[
 \chi(xyz)=\sum r_{\alpha\delta}(\chi)a_{\delta\beta\gamma}x_\alpha y_\beta z_\gamma
 =\sum r_{\kappa\lambda}(\chi)a_{\lambda\beta\gamma}x_\kappa y_\beta z_\gamma
\]
nur von $x_1,\ldots,x_m$, $y_1,\ldots,y_m$, $z_1,\ldots,z_m$ ab. Ist daher $\beta$ oder $\gamma>m$, so ist $\sum_{\lambda=1}^m r_{\kappa\lambda}(\chi)a_{\lambda\beta\gamma}=0$ und mithin $a_{\lambda\beta\gamma}=0$. Demnach entsprechen den $k$ verschiedenen Primfaktoren von $\Theta$ $k$ mit $(\eps)$ homomorphe Gruppen.

Oder noch einfacher: Alle Größen $x$, deren Koordinaten den linearen Gleichungen $\sum_\beta r_{\alpha\beta}(\chi)x_\beta=0$ genügen, bilden eine invariante Untergruppe $(\mathfrak S)$ von $(\eps)$. Denn diese Gleichungen drücken aus, daß $\chi(tx)=0$ ist für jede Größe $t$. Ist $y$ eine beliebige Größe von $(\eps)$, und ersetzt man $t$ durch $yt$ oder $ty$, so erhält man $\chi(t(xy))=0$ oder $\chi(t(yx))=0$. Gehört also $x$ dem Komplexe $(\mathfrak S)$ an, so gehören ihm auch $xy$ und $yx$ an. Folglich ist $(\mathfrak S)$ eine invariante Untergruppe von $(\eps)$.

Eine Gruppe $(\eps)$ der Ordnung $r$ heißt einfach (\textsc{Molien}: Ursprüngliches Zahlensystem), wenn sie keine invariante Untergruppe hat (außer $(\eps)$), also wenn sie keine homomorphe Gruppe hat, deren Ordnung $<n$ ist.

Nach der obigen Entwicklung darf die Determinante $\Theta$ einer einfachen Gruppe keinen Primfaktor $\Phi$ enthalten, dessen linearer Rang $m<n$ ist. Nach Satz II, § 6 darf daher $\Theta$ nicht zwei verschiedene Primfaktoren enthalten, muß also $\Theta=\Phi^s$ eine Potenz einer Primfunktion und mithin $\sigma=s\chi$ sein. Ferner muß der Rang $m$ von $\Phi$, also der Rang der Matrix $R_\sigma=sR_\chi$, gleich $m=n=rs$, also muß $(\eps)$ eine \textsc{Dedekind}sche Gruppe sein.

Ist umgekehrt $|R_\sigma|$ von Null verschieden, und $\Theta=\Phi^s$ eine Potenz einer Primfunktion, so ist $(\eps)$ eine einfache Gruppe, so ist sie nicht einer Gruppe $(\eta)$ homomorph, deren Ordnung $n'<n$ ist (\textsc{Mol}. Satz 23). Denn die Determinante von $(\eta)$ ist ein Divisor von $\Theta$, also gleich $\Theta'=\Phi^{s'}$, und hat daher denselben Rang wie $\Phi$, nämlich $n$. Der Rang $n$ von $\Theta'$ kann aber nicht größer als die Ordnung $n'$ von $(\eta)$ sein.

Da $S(x)$ und $T(y)$ vertauschbare Matrizen sind, so zerfällt die Determinante (\emph{Gr.} § 10, (1.))
\[
 |uS(x)+vT(y)+wE|
\]
in ein Produkt linearer Funktionen von $u,v,w$. Ist $uu_1+vv_1+w$ einer davon, so erkennt man, indem man $v=0$ ($u=0$) setzt, daß $u_1$ ($v_1$) eine charakteristische Wurzel von $\Phi(x)$ ($\Phi(y)$) ist. Betrachtet man $y$ als konstant, $x$ aber als variabel, so ist $\Phi(ux+we)$ als Funktion von $w$ irreduzibel, also auch $\Phi(ux+(vv_1+w)e)$. Mit dieser Funktion hat die Determinante den Faktor $uu_1+vv_1+w$ gemeinsam, also alle $r$ Faktoren $uu_\alpha+vv_1+w$. Betrachtet man dann $x$ als konstant und $y$ als variabel, so erkennt man, daß die Determinante jeden der $r^2$ Faktoren
\[
 uu_\alpha+vv_\beta+w\qquad(\alpha,\beta=1,2,\ldots,r)
\]
und jeden gleich oft enthält. Mithin ist
\[
 |uS(x)+vT(y)+wE|=\left(\prod(uu_\alpha+vv_\beta+w)\right)^c
\]
und demnach $s=rc$. Speziell ist (\emph{Gr.} § 10, (8.))
\[
\fnum{1} |S(x)-T(x)+wE|=w^s\prod_{\beta>\alpha}\bigl(w^2-(u_\beta-u_\alpha)^2\bigr)^c.
\]
Die Elementarteiler der charakteristischen Determinanten der beiden vertauschbaren Matrizen $S(x)$ und $T(x)$ sind alle linear. Nach Satz I, § 8 sind es also auch die Elementarteiler der Determinante (1.), und insbesondere die $s$ für $w=0$ verschwindenden. Folglich hat die Matrix $S(x)-T(x)$ den Rang $n-s$.

Nach Satz III, § 8 hat sie aber den Rang $n-r$. Daher ist (\textsc{Mol}. Satz 29)
\[
\fnum{2} r=s,
\qquad m=n=r^2,
\]
mithin $c=1$ und
\[
\fnum{3} |uS(x)+vT(y)+wE|=\prod_{\alpha,\beta}^{r}(uu_\alpha+vv_\beta+w).
\]
\emph{Die Determinante einer einfachen Gruppe ist eine Potenz einer Primfunktion, deren Exponent gleich dem Grade der Primfunktion ist. Der lineare Rang dieser Funktion, welcher mit der Ordnung der Gruppe übereinstimmt, ist gleich dem Quadrate ihres Grades.}

\section*{§ 11.}

Sind für eine bestimmte Größe $h$ der einfachen Gruppe $(\eps)$ der Ordnung $n=r^2$ die $r$ Wurzeln der Gleichung $\Phi(ue-h)=0$ alle verschieden, und ist $a$ eine bestimmte dieser $r$ Wurzeln, so hat die Determinante $|S(ue-h)|$ lauter lineare Elementarteiler, und folglich die Matrix $S(ae-h)$ den Rang $n-r$. Daher haben (\textsc{Mol}. § 8) die linearen Gleichungen $(ae-h)t=0$ zwischen den Koordinaten $t_1,\ldots,t_n$ der unbekannten Größe $t$ genau $r$ unabhängige Lösungen, $t=t^{(1)},t^{(2)},\ldots,t^{(r)}$.

Ist $x$ eine variable Größe, so ist auch $t^{(\kappa)}x$ eine Lösung, und mithin ist
\[
 t^{(\kappa)}x=\sum_\lambda x_{\kappa\lambda}t^{(\lambda)}.
\]
Da die $r$ Lösungen $t^{(\lambda)}$ unabhängig sind, so sind die Größen $x_{\kappa\lambda}=f_{\kappa\lambda}(x)$ lineare Funktionen der Koordinaten $x_1,\ldots,x_n$. Sind $y$ und $z$ zwei andere Größen, und setzt man $y_{\kappa\lambda}=f_{\kappa\lambda}(y)$, $z_{\kappa\lambda}=f_{\kappa\lambda}(z)$, so ist auch
\[
 t^{(\kappa)}y=\sum_\lambda y_{\kappa\lambda}t^{(\lambda)},\qquad
 t^{(\kappa)}z=\sum_\lambda z_{\kappa\lambda}t^{(\lambda)}
\]
und mithin
\[
 t^{(\kappa)}yz=\sum_\mu y_{\kappa\mu}t^{(\mu)}z
 =\sum_{\lambda,\mu}y_{\kappa\mu}z_{\mu\lambda}t^{(\lambda)}.
\]
Ist also $x=yz$, so ist
\[
\fnum{1} x_{\kappa\lambda}=\sum_\mu y_{\kappa\mu}z_{\mu\lambda}.
\]
Unter den $n=r^2$ linearen Funktionen $x_{\kappa\lambda}=f_{\kappa\lambda}(x)$ seien $m$ unabhängig. Ist dann $x=yz$, so hängen diese $m$ Verbindungen von $x_1,\ldots,x_n$ nach (1.) von den nämlichen $m$ Verbindungen von $y_1,\ldots,y_n$ und denen von $z_1,\ldots,z_n$ ab. Folglich hat $(\eps)$ eine homomorphe Gruppe der Ordnung $m$. Da aber $(\eps)$ einfach ist, so ist $m=n=r^2$, und mithin kann man $r^2$ neue Grundzahlen $\eps_{\kappa\lambda}$ einführen, so daß
\[
 \sum \eps_\alpha x_\alpha=\sum \eps_{\kappa\lambda}x_{\kappa\lambda}
\]
wird. Für diese gehen die Formeln (4.) § 1 in (1.) und die Relationen (2.) § 1 in
\[
\fnum{2} \eps_{\alpha\beta}\eps_{\delta\gamma}=\eps_{\alpha\gamma},
\qquad
\eps_{\alpha\delta}\eps_{\beta\gamma}=0\quad(\beta\ne\delta)
\]
über.

Den $k$ verschiedenen Primfaktoren von $\Theta$ entsprechen, wie am Anfang des § 10 gezeigt ist, $k$ mit $(\eps)$ homomorphe Gruppen. Die erste hat die Ordnung $m$, ihre Determinante ist ein Teiler von $\Theta$ und hängt nur von $x_1,\ldots,x_m$ ab. Da $\Phi',\Phi'',\ldots$ von $x_1,\ldots,x_m$ unabhängig sind, so ist ihre Determinante eine Potenz von $\Phi$, hat also, ebenso wie $\Phi$, den linearen Rang $m$. Weil ihre Ordnung ebenfalls gleich $m$ ist, so ist sie nach § 10 eine einfache Gruppe. Folglich ist ihre Determinante gleich $\Phi^r$, und der Rang von $\Phi$ ist $m=r^2$.

\emph{I. Der lineare Rang jedes Primfaktors der Gruppendeterminante ist gleich dem Quadrate seines Grades.}

Ich bezeichne jetzt die $k$ verschiedenen Funktionen $\Phi$ mit $\Phi_1,\Phi_2,\ldots,\Phi_k$. In der Gruppendeterminante $S$ verschwinden alle Elemente $s_{\alpha\beta}$, welche die ersten $m_1$ Zeilen mit den letzten $n-m_1$ Spalten gemeinsam haben, ebenso alle Elemente, welche die folgenden $m_2$ Zeilen mit den ersten $m_1$ und den letzten $n-m_1-m_2$ Spalten gemeinsam haben, ferner die, welche die folgenden $m_3$ Zeilen mit den ersten $m_1+m_2$ und den letzten $n-m_1-m_2-m_3$ Spalten gemeinsam haben usw., während über die letzten $n-m_1-\cdots-m_k$ Zeilen nichts bekannt ist. Dasselbe gilt von der antistrophen Matrix $T$, also auch von $uS(x)+vT(y)+wE$. Bilden die ersten $m_1$ Zeilen und Spalten von $S$ ($T$) die Matrix $S_1$ ($T_1$), die folgenden $m_2$ Zeilen und Spalten die Matrix $S_2$ ($T_2$), $\ldots$, so sei
\[
 U=\begin{pmatrix}
 S_1&0&\cdots&0\\
 0&S_2&\cdots&0\\
 \cdot&\cdot&\cdots&\cdot\\
 0&0&\cdots&S_k
 \end{pmatrix},
 \qquad
 V=\begin{pmatrix}
 T_1&0&\cdots&0\\
 0&T_2&\cdots&0\\
 \cdot&\cdot&\cdots&\cdot\\
 0&0&\cdots&T_k
 \end{pmatrix}.
\]
Dann ist
\[
 S=\begin{pmatrix}U&0\\ U_0&S_0\end{pmatrix},
 \qquad
 T=\begin{pmatrix}V&0\\ V_0&T_0\end{pmatrix}.
\]
Daraus ergibt sich der Satz von \textsc{Molien}:

\emph{II. Jede Gruppe ist einer \textsc{Dedekind}schen Gruppe homomorph, deren Determinante jeden Primfaktor ihrer Determinante $\Theta$ enthält, und deren Ordnung gleich dem linearen Range von $\Theta$ ist.}

Dieser Rang $m$ ist gleich dem Range der quadratischen Form $\sigma(x^2)=\sum p_{\alpha\beta}x_\alpha x_\beta$, und um jene Gruppe zu erhalten, braucht man nur die $m$ unabhängigen unter den $n$ Funktionen $\sum_\beta p_{\alpha\beta}x_\beta$ als neue Variable einzuführen; also in der Formel (4.) § 9 die Substitution $C$ so zu wählen, daß sie die Form $\sigma(x^2)$ in eine Funktion von $m$ Variabeln, etwa in eine Summe von $m$ Quadraten transformiert. Es ist mir aber nicht gelungen, diesen Satz direkt, d. h. ohne Benutzung der Primfunktionen $\Phi_\kappa$, zu beweisen.

Ferner ist
\[
 |S|=|S_1|\cdots|S_k||S_0|,
 \qquad
 |T|=|T_1|\cdots|T_k||T_0|
\]
und
\[
 |uS(x)+vT(y)+wE|=|uS_1+vT_1+wE_1|\cdots|uS_k+vT_k+wE_k|\,|uS_0+vT_0+wE_0|.
\]
Da $S_\kappa$ die Matrix einer einfachen Gruppe ist, so ist nach (3.) § 10
\[
 \Psi_\kappa=|uS_\kappa(x)+vT_\kappa(y)+wE_\kappa|
 =\prod(uu_\alpha+vv_\beta+w),
\]
wo $u_1,\ldots,u_r$ die charakteristischen Wurzeln von $\Phi_\kappa(x)$ sind, $v_1,\ldots,v_r$ die von $\Phi_\kappa(y)$. Diese Resultate bleiben unverändert, wenn die Grundzahlen von $(\eps)$ nicht in der oben beschriebenen Weise, sondern beliebig gewählt werden.

Auch $\Psi_0=|uS_0(x)+vT_0(y)+wE_0|$ zerfällt in lineare Faktoren $uu_0+vv_0+w$. Hier ist $u_0$ eine charakteristische Wurzel eines Primfaktors $\Phi_\kappa(x)$ und $v_0$ eine solche eines Primfaktors $\Phi_\lambda(y)$. Aus der Irreduzibilität von $\Phi_\kappa(ue-x)$ und $\Phi_\lambda(ve-y)$ folgt dann, daß $\Psi_0$ auch alle $r_\kappa r_\lambda$ Faktoren des Produktes $\prod(uu_\alpha+vv_\beta+w)$ enthält, worin $u_\alpha$ die $r_\kappa$ charakteristischen Wurzeln von $\Phi_\kappa(x)$ und $v_\beta$ die $r_\lambda$ von $\Phi_\lambda(y)$ durchläuft. Ist $\kappa=\lambda$, so tritt also der Faktor $\Psi_\kappa$ wiederholt auf. Wie das Beispiel des § 3 zeigt, kann aber auch $\kappa$ von $\lambda$ verschieden sein.

Wählt man für $x$ eine invariante Größe, so sind die $r_\kappa$ Größen $u_\alpha$ nach (12.) § 4 alle gleich $\frac1{r_\kappa}\chi^{(\kappa)}(x)$. Die Formeln (2.) und (3.) § 8 zeigen aber, daß dann der Wurzel $v_\beta$ die Wurzel $\frac1{r_\lambda}\chi^{(\lambda)}(x)$ und nur diese zugeordnet ist. Demnach kann $\kappa$ nur dann von $\lambda$ verschieden sein, wenn für eine invariante Größe $x$
\[
\fnum{1} \frac1{r_\kappa}\chi^{(\kappa)}(x)=\frac1{r_\lambda}\chi^{(\lambda)}(x)\qquad(S_x=T_x)
\]
ist.

Nun sei $w=0$, $u=v=1$ und
\[
\fnum{2} \Psi_{\kappa\lambda}(x,y)=\prod(u_\alpha+v_\beta),
\]
wo $u_\alpha$ die $r_\kappa$ charakteristischen Wurzeln von $\Phi_\kappa(x)$ durchläuft, und $v_\beta$ die $r_\lambda$ von $\Phi_\lambda(y)$. Dann ist $\Psi_{\kappa\lambda}$ eine irreduzibele Funktion von $x_1,\ldots,x_n,y_1,\ldots,y_n$, in bezug auf jene vom Grade $r_\kappa$, in bezug auf diese vom Grade $r_\lambda$, und es ist
\[
\fnum{3} \Psi(x,y)=|S(x)+T(y)|=\prod\bigl(\Psi_{\kappa\lambda}(x,y)\bigr)^{c_{\kappa\lambda}},
\]
wo die ganze Zahl $c_{\kappa\kappa}>0$, aber $c_{\kappa\lambda}\geq0$ ist, und nur unter der Bedingung (1.) $c_{\kappa\lambda}>0$ sein kann. Da
\[
 \Psi_{\kappa\lambda}(x,0)=\Phi_\kappa(x)^{r_\lambda},
 \qquad
 \Psi_{\kappa\lambda}(0,y)=\Phi_\lambda(y)^{r_\kappa}
\]
ist, so ergeben sich für die in den Zerlegungen
\[
\fnum{4} |S(x)|=\prod\Phi_\lambda(x)^{s_\lambda},
\qquad
|T(x)|=\prod\Phi_\lambda(x)^{t_\lambda}
\]
auftretenden Exponenten die Ausdrücke
\[
\fnum{5} s_\lambda=c_{\lambda1}r_1+c_{\lambda2}r_2+\cdots+c_{\lambda k}r_k,
\qquad
 t_\lambda=c_{1\lambda}r_1+c_{2\lambda}r_2+\cdots+c_{k\lambda}r_k,
\]
so daß stets
\[
\fnum{6} s_\lambda\geq r_\lambda,
\qquad t_\lambda\geq r_\lambda
\]
ist. Die beiden Funktionen $|S(x)|$ und $|T(x)|$ haben denselben linearen Rang
\[
\fnum{7} m=r_1^2+r_2^2+\cdots+r_k^2,
\]
dieser ist gleich dem Range der Matrix $R_\sigma$ und dem von $R_\tau$ und gleich der Summe der Rangzahlen der $k$ Matrizen $R(\chi^{(\kappa)})$ oder der $k$ Primfunktionen $\Phi_\kappa(x)$. Dagegen ist die Ordnung der Gruppe $(\eps)$
\[
\fnum{8} n=r_1s_1+\cdots+r_ks_k=r_1t_1+\cdots+r_kt_k=\sum c_{\kappa\lambda}r_\kappa r_\lambda,
\]
wo $r_\kappa r_\lambda$, falls $\kappa$ von $\lambda$ verschieden ist, den Koeffizienten $c_{\kappa\lambda}+c_{\lambda\kappa}$ hat.

\section*{§ 13.}

Ich nehme jetzt die in den §§ 7 und 8 begonnene Theorie der \textsc{Dedekind}schen Gruppen wieder auf. Für eine solche ist $m=n$, also nach (6.), (7.) und (8.) § 12 $s_\kappa=t_\kappa=r_\kappa$.

\emph{I. Der Exponent der in der Determinante einer \textsc{Dedekind}schen Gruppe aufgehenden Potenz einer Primfunktion ist dem Grade dieser Funktion gleich.}

Die Gruppe $(\eps)$ zerfällt in $k$ einfache Gruppen der Ordnungen $r^2,r'^2,\ldots$, in jeder derselben zerfällt $S(x)$ in $r$ identische Matrizen $r$ten Grades, und die $r^2$ Elemente $x_{\kappa\lambda}$ einer solchen Matrix sind unter einander und von den Variabeln der anderen einfachen Gruppen unabhängig. In meiner Arbeit \emph{Über die Darstellung der endlichen Gruppen durch lineare Substitutionen} II., Sitzungsberichte 1899, § 5, habe ich daraus den Satz abgeleitet:

\emph{II. Sind $a$ und $b$ zwei bestimmte Größen einer \textsc{Dedekind}schen Gruppe, und stimmen die elementaren Invarianten der beiden Matrizen $S(ue-a)$ und $S(ue-b)$ überein, so gibt es in der Gruppe eine Größe $c$, wofür $|S(c)|$ von Null verschieden ist, und die der Bedingung $c^{-1}ac=b$ genügt.}

Die Determinante einer \textsc{Dedekind}schen Gruppe ist
\[
\fnum{1} \Theta(x)=|S(x)|=|T(x)|=\prod\Phi^r,
\]
ihre Spur ist
\[
\fnum{2} \sigma=\tau=r\chi+r'\chi'+r''\chi''+\cdots=\sum r\chi,
\]
demnach ist auch
\[
\fnum{3} P=R_\sigma=rR_\chi+r'R_{\chi'}+r''R_{\chi''}+\cdots=\sum rR_\chi.
\]
Die Matrix $R_\chi$ hat den Rang $r^2$, die Matrix $r'R_{\chi'}+r''R_{\chi''}+\cdots=P-rR_\chi$ nach Satz II, § 6 den Rang $r'^2+r''^2+\cdots=n-r^2$. Die Determinante $|uP-rR_\chi|$ verschwindet daher für $u=0$ mindestens von der Ordnung $n-r^2$, für $u=1$ mindestens von der Ordnung $r^2$, und da sie nicht für mehr als $r$ Werte verschwindet, so ist nicht nur
\[
\fnum{4} |uP-rR_\chi|=|P|u^{n-r^2}(u-1)^{r^2},
\]
sondern es sind auch die Elementarteiler dieser Determinante alle linear (weil sie sonst für $u=0$ oder $1$ von höherer Ordnung verschwinden würde). Folglich genügt die Matrix $H=rP^{-1}R_\chi$ der Gleichung $H^2=H$, also ist
\[
\fnum{5} R_\chi P^{-1}R_\chi=\frac1r R_\chi.
\]
Fügt man daher in (3.) links den Faktor $R_\chi P^{-1}$ hinzu, so erhält man
\[
 0=r'(R_\chi P^{-1})R_{\chi'}+r''(R_\chi P^{-1})R_{\chi''}+\cdots,
\]
und mithin ist nach (4.) § 6
\[
\fnum{6} R_\chi P^{-1}R_{\chi'}=0.
\]
Bestehen zwischen $\xi_1,\ldots,\xi_n$ und $x_1,\ldots,x_n$ die Gleichungen
\[
\fnum{7} \xi_\alpha=\sum_\beta p_{\alpha\beta}x_\beta,
\]
so nenne ich den Parameter $\xi$ und die Größe $x$ konjugiert. Dann ist
\[
\fnum{8} R_\xi=S_x'P=PT_x,
\qquad
R_\xi'=PS_x=T_x'P.
\]
Ist also $R_\xi=R_\xi'$, so ist $S_x=T_x$, einem symmetrischen Parameter ist eine invariante Größe konjugiert und umgekehrt.

Die dem charakteristischen Parameter $\xi=\chi$ konjugierte invariante Größe $x=h$ nenne ich eine für die Primfunktion $\Phi$ charakteristische Größe. Dann ist
\[
\fnum{9} \chi_\alpha=\sum_\beta p_{\alpha\beta}h_\beta
\]
und
\[
\fnum{10} R_\chi=PS_h=PT_h.
\]
Aus (5.) und (10.) folgt $\frac1r S(h)=S(h)^2=S(h^2)$ und $0=S(h)S(h')=S(hh')$, also (\emph{Gr.} § 5 (9.))
\[
\fnum{11} h^2=\frac1r h,
\qquad hh'=0
\]
und aus (3.)
\[
\fnum{12} e=rh+r'h'+r''h''+\cdots=\sum rh,
\]
so daß der Parameter $\sigma$ und die Haupteinheit $e$ konjugiert sind. Ferner ist
\[
 rR_\chi S_h=R_\chi,
 \qquad
 r\sum_\lambda r_{\alpha\lambda}(\chi)s_{\lambda\beta}(h)=r_{\alpha\beta}(\chi).
\]
Multipliziert man mit $e_\alpha e_\beta$ und summiert nach $\alpha$ und $\beta$, so erhält man (vergl. \textsc{Ded}. (13.))
\[
\fnum{13} \sum_\alpha \chi_\alpha h_\alpha=1,
\qquad
\sum_\alpha \chi'_{\alpha}h_\alpha=0,
\]
oder wenn man
\[
 h(\xi)=h_1\xi_1+\cdots+h_n\xi_n
\]
setzt,
\[
\fnum{14} \chi(h)=h(\chi)=1,
\qquad
\chi(h')=h(\chi')=0,
\]
oder endlich, wenn $P^{-1}=Q=(q_{\alpha\beta})$ ist,
\[
\fnum{15} \sum p_{\alpha\beta}h_\alpha h_\beta=\sum q_{\alpha\beta}\chi_\alpha\chi_\beta=1,
\qquad
\sum p_{\alpha\beta}h_\alpha h'_{\beta}=\sum q_{\alpha\beta}\chi_\alpha\chi'_{\beta}=0.
\]

\section*{§ 14.}

Seien $u_1,\ldots,u_r$ die charakteristischen Wurzeln von $\Phi(x)$, $v_1,\ldots,v_r$ die von $\Phi(y)$. Ist dann $g(u,v)$ eine ganze Funktion von $u$ und $v$, so sind nach § 12 die $n$ charakteristischen Wurzeln der Matrix $g(S_x,T_y)$ die $r^2$ Größen $g(u_\alpha,v_\beta)$ und die $r'^2,r''^2,\ldots$ Größen, die aus $\Phi',\Phi''$ in der nämlichen Weise erhalten werden.

Ist $x$ eine invariante Größe, so ist nach dem Satze des § 4
\[
 u_1=\cdots=u_r=\frac1r\chi(x)=c.
\]
Nach (13.) § 13 ist daher die charakteristische Determinante von $S_h=T_h$ gleich $(u-\frac1r)^{r^2}u^{n-r^2}$, und die charakteristischen Wurzeln von $(S(x)-cE)T(h)=S((x-ce)h)$ sind alle Null. Mithin verschwindet eine Potenz dieser Matrix. Da sie aber mit jeder Gruppenmatrix $S(y)$ vertauschbar ist, so sind nach Satz III, § 8 die Elementarteiler ihrer charakteristischen Determinante alle linear. Daher verschwindet die erste Potenz, und mithin ist
\[
\fnum{1} xh=ch=h\frac1r\chi(x).
\]
Multipliziert man also (12.) § 13 mit $x$, so erhält man
\[
\fnum{2} x=rch+r'c'h'+r''c''h''+\cdots
=h\chi(x)+h'\chi'(x)+h''\chi''(x)+\cdots,
\]
wo die Koeffizienten $rc=\chi(x)$, $r'c'=\chi'(x)$, $\ldots$ gewöhnliche Größen sind.

Die $k$ Größen $h,h',h'',\ldots$, die nach (5.) § 6 oder nach (11.) § 13 voneinander linear unabhängig sind, bilden also ein vollständiges System von Lösungen der linearen Gleichungen $S(x)=T(x)$. Unter den $n^2$ Gleichungen $s_{\alpha\beta}(x)=t_{\alpha\beta}(x)$ sind demnach $n-k$ unabhängige. Wenn nun die Größe $x$ außerdem den Gleichungen $\chi(x)=\chi'(x)=\cdots=0$ genügt, so ist nach (2.) $x=0$.

Dies kann man auch so einsehen: Ist $xy=yx$, so ist nach (11.) § 4 $r\chi(xy)=\chi(x)\chi(y)$. Ist also $\chi(x)=0$, so ist auch $\chi(xy)=0$ oder $\sum_\beta r_{\alpha\beta}(\chi)x_\beta=0$. Gilt dies auch für $\chi',\chi'',\ldots$, so ergeben sich $r^2+r'^2+\cdots=n$ unabhängige lineare Gleichungen, und mithin ist $x_\beta=0$.

\emph{I. Unter den $n^2$ linearen Funktionen $s_{\alpha\beta}(x)-t_{\alpha\beta}(x)$ der $n$ Variabeln $x_1,\ldots,x_n$ sind $n-k$ unabhängige. Sie bilden zusammen mit den $k$ Funktionen $\chi(x),\chi'(x),\ldots$ $n$ unabhängige Funktionen. Die $k$ charakteristischen Größen $h,h',\ldots$ bilden ein vollständiges System unabhängiger Lösungen der Gleichungen $s_{\alpha\beta}(x)=t_{\alpha\beta}(x)$.}

\emph{II. Unter den linearen Funktionen $r_{\alpha\beta}(\xi)-r_{\beta\alpha}(\xi)$ der $n$ Variabeln $\xi_1,\ldots,\xi_n$ sind $n-k$ unabhängige. Sie bilden zusammen mit den $k$ Funktionen $h(\xi),h'(\xi),\ldots$ $n$ unabhängige Funktionen. Die $k$ charakteristischen Parameter $\chi,\chi',\ldots$ bilden ein vollständiges System unabhängiger Lösungen der Gleichungen $r_{\alpha\beta}(\xi)=r_{\beta\alpha}(\xi)$.}

Ein spezieller Fall dieser Sätze ist das elegante Kriterium, das Hr. \textsc{Molien} in Satz 9 und 10, § 3 für die Einfachheit einer Gruppe angibt. Jede einfache Gruppe ist nach § 10 eine \textsc{Dedekind}sche, und damit eine solche einfach sei, ist notwendig und hinreichend, daß $k=1$ ist.

\emph{III. Damit eine Gruppe einfach sei, ist notwendig und hinreichend, daß die Determinante $n$ten Grades $\left|\sum_\kappa \sigma_\kappa a_{\kappa\alpha\beta}\right|$ von Null verschieden ist, und daß $\xi_\kappa=\sigma_\kappa$ die einzige Lösung der linearen Gleichungen}
\[
 \sum_\kappa(a_{\kappa\alpha\beta}-a_{\kappa\beta\alpha})\xi_\kappa=0
\]
\emph{ist.}

Oder auch:

\emph{IV. Damit eine Gruppe einfach sei, ist notwendig und hinreichend, daß die Determinante $n$ten Grades $\left|\sum_\kappa \sigma_\kappa a_{\kappa\alpha\beta}\right|$ von Null verschieden ist, und daß $x_\lambda=e_\lambda$ die einzige Lösung der linearen Gleichungen}
\[
 \sum_\lambda(a_{\alpha\lambda\beta}-a_{\alpha\beta\lambda})x_\lambda=0
\]
\emph{ist.}

Daß hier jeder Satz in zwei verschiedenen Gestalten erscheint, hat seinen Grund darin, daß für die Gruppe selbst zwei konjugierte Formen neben einander gestellt werden können. Sind den Parametern $\xi,\eta,\zeta$ die Größen $x,y,z$ konjugiert, so sei
\[
 F(\xi,y,z)=G(x,\eta,\zeta)=\sum b_{\alpha\beta\gamma}x_\alpha\eta_\beta\zeta_\gamma.
\]
Dann ist
\[
\fnum{3} a_{\alpha\beta\gamma}=\sum_\lambda b_{\beta\lambda\alpha}p_{\lambda\gamma}
=\sum_\lambda b_{\gamma\alpha\lambda}p_{\lambda\beta}.
\]
Durch diese Transformation geht also die Gruppe $(\eps)$ in eine äquivalente Gruppe $(\bar\eps)$ über, die konjugierte Gruppe. Bedient man sich für diese der Bezeichnungen (4.) § 9, so ergibt sich, falls in jenen Formeln $C=P$ gesetzt wird,
\[
\fnum{4} R_x=\bar T_xP,
\qquad R'_x=\bar S_xP,
\qquad S_x=\bar R'_xP,
\qquad T_x=\bar R_xP.
\]
Ferner ist $\bar\sigma_\alpha=e_\alpha$. Ebenso, wie $\bar a_{\alpha\beta\gamma}=b_{\alpha\beta\gamma}$ gesetzt ist, werde $\bar p_{\alpha\beta}=q_{\alpha\beta}$ gesetzt. Dann ist $q_{\alpha\beta}=\sum_\kappa e_\kappa b_{\kappa\alpha\beta}$ und mithin
\[
 \sum_\lambda p_{\alpha\lambda}q_{\lambda\beta}
 =\sum_{\kappa,\lambda}p_{\alpha\lambda}e_\kappa b_{\kappa\lambda\beta}
 =\sum_\kappa e_\kappa a_{\beta\kappa\alpha}=e_{\alpha\beta},
\]
also $Q=P^{-1}$. Folglich ist die konjugierte Gruppe von $(\bar\eps)$ wieder die ursprüngliche Gruppe,
\[
\fnum{5} b_{\alpha\beta\gamma}=\sum_\lambda a_{\beta\lambda\alpha}q_{\lambda\gamma}
=\sum_\lambda a_{\gamma\alpha\lambda}q_{\lambda\beta}.
\]

\section*{§ 15.}

In § 6, (5.) ist gezeigt, daß die $k$ Funktionen $\chi(x),\chi'(x),\ldots$ linear unabhängig sind. Nach Satz I, § 14 ist dies auch dann noch der Fall, wenn $x$ nicht unbeschränkt veränderlich ist, sondern nur die invarianten Größen von $(\eps)$ durchläuft. Diese bilden eine Untergruppe $(\eta)$ von $(\eps)$, deren Ordnung $k$ ist (im allgemeinen keine invariante Untergruppe). Als Grundzahlen $\eta_1,\eta_2,\ldots,\eta_k$ von $(\eta)$ kann man die $k$ Größen $rh,r'h',\ldots$ wählen. Dann ist
\[
\fnum{1} \eta_\kappa^2=\eta_\kappa,
\qquad \eta_\kappa\eta_\lambda=0.
\]
Folglich ist $(\eta)$ eine \textsc{Dedekind}sche kommutative Gruppe, deren Determinante, falls man die Größen $\eta_1,\ldots,\eta_k$ durch $\eps_1,\ldots,\eps_n$ ausdrückt, gleich $\prod \frac1r\chi(x)$ ist. In der Tat sind diese $k$ linearen Faktoren, worin die Veränderlichkeit der Koordinaten durch die Bedingungen $s_{\alpha\beta}(x)=t_{\alpha\beta}(x)$ eingeschränkt ist, linear unabhängig.

\emph{I. Ist $k$ die Anzahl der verschiedenen Primfaktoren der Determinante einer \textsc{Dedekind}schen Gruppe, so bilden deren invariante Größen eine kommutative Gruppe der Ordnung $k$, die ebenfalls eine \textsc{Dedekind}sche Gruppe ist. Ihre Determinante ist $\prod\frac1r\chi(x)$, während die Determinante der gegebenen Gruppe für eine invariante Variable $x$ gleich $\prod\left(\frac1r\chi(x)\right)^{r^2}$ ist.}

Sind die Grade $r=r'=\cdots=1$, so ist $k=n$, also ist nach Satz I, § 14 identisch $s_{\alpha\beta}(x)=t_{\alpha\beta}(x)$, $a_{\alpha\beta\gamma}=a_{\alpha\gamma\beta}$.

\emph{II. Die Determinante einer \textsc{Dedekind}schen Gruppe zerfällt stets, aber auch nur dann in lauter lineare Faktoren, wenn die Gruppe eine kommutative ist.}

Damit die Determinante einer beliebigen Gruppe in lauter lineare Faktoren zerfalle, ist nach Satz II, § 12 notwendig und hinreichend, daß die ihr homomorphe \textsc{Dedekind}sche Gruppe der Ordnung $m=\sum r^2$ diese Eigenschaft besitzt, demnach eine kommutative Gruppe ist. Dieselbe wird erhalten, indem man die $m$ unabhängigen unter den Funktionen $\sum_\beta p_{\alpha\beta}t_\beta$ als Koordinaten einer Größe $t$ einführt. Setzt man also $t=yz$, so müssen diese bilinearen Funktionen von $y_1,\ldots,y_n,z_1,\ldots,z_n$ bei Vertauschung von $y$ und $z$ ungeändert bleiben. Ist $x$ eine dritte Variable, so hat demnach die Funktion
\[
 \sum_{\alpha,\beta}p_{\alpha\beta}x_\alpha t_\beta=\sigma(xt)=\sigma(xyz)
\]
dieselbe Eigenschaft. Diese bleibt nach § 6 stets bei einer zyklischen Vertauschung von $x,y,z$ ungeändert, in dem betrachteten Falle also bei jeder Vertauschung. Sie ist die Spur der Matrix $S(xyz)=S(x)S(y)S(z)$, also gleich
\[
\fnum{2} \sigma(xyz)=\sum_{\kappa,\lambda,\mu}s_{\lambda\mu}(x)s_{\mu\kappa}(y)s_{\kappa\lambda}(z)
=\sum a_{\lambda\alpha\mu}a_{\mu\beta\kappa}a_{\kappa\gamma\lambda}x_\alpha y_\beta z_\gamma.
\]
Benutzt man die antistrophe Gruppe, so tritt $\tau(xyz)$ an Stelle von $\sigma(xyz)$. Demnach ergibt sich der Satz (vergl. \textsc{Cartan}, \emph{I. Thèse, Sur la structure des groupes de transformations finis et continus}, Paris 1894, p. 48, (5.)):

\emph{III. Damit die Determinante einer Gruppe in lauter lineare Faktoren zerfalle, ist notwendig und hinreichend, daß die trilineare Funktion $\sigma(xyz)$ (oder $\tau(xyz)$) bei Vertauschung von $y$ und $z$ ungeändert bleibt, mithin eine symmetrische Funktion der drei Reihen von Variabeln ist, daß also die Ausdrücke}
\[
\fnum{3} \sum_{\kappa,\lambda,\mu}a_{\lambda\alpha\mu}a_{\mu\beta\kappa}a_{\kappa\gamma\lambda}
\quad\text{oder}\quad
\sum_{\kappa,\lambda,\mu}a_{\lambda\mu\alpha}a_{\mu\kappa\beta}a_{\kappa\lambda\gamma}
\]
\emph{auch bei einer Transposition von $\alpha,\beta,\gamma$ ungeändert bleiben.}

Da $\sigma(xyz)$ auch die Spur von $S(x)S(yz)$ ist, so können statt der Summen (3.) auch die Ausdrücke
\[
\fnum{4} \sum_{\kappa,\lambda,\mu}a_{\kappa\lambda\mu}a_{\mu\alpha\kappa}a_{\lambda\beta\gamma}
\quad\text{oder}\quad
\sum_{\kappa,\lambda,\mu}a_{\kappa\lambda\mu}a_{\lambda\kappa\alpha}a_{\mu\beta\gamma}
\]
genommen werden. Mittels der Formeln (1.) und (2.) § 7 erhält man endlich die Ausdrücke
\[
\fnum{5}\begin{gathered}
 \sum_{\kappa,\lambda,\mu}a_{\kappa\lambda\kappa}a_{\lambda\alpha\mu}a_{\mu\beta\gamma}
 \quad\text{oder}\quad
 \sum_{\kappa,\lambda,\mu}a_{\kappa\kappa\lambda}a_{\lambda\alpha\mu}a_{\mu\beta\gamma},\\
 \sum_{\kappa,\lambda,\mu}a_{\kappa\lambda\kappa}a_{\lambda\mu\alpha}a_{\mu\beta\gamma}
 \quad\text{oder}\quad
 \sum_{\kappa,\lambda,\mu}a_{\kappa\kappa\lambda}a_{\lambda\mu\alpha}a_{\mu\beta\gamma}.
\end{gathered}
\]

Mit Hülfe der obigen Sätze läßt sich die Zerlegung der Determinante einer \textsc{Dedekind}schen Gruppe in ihre Primfaktoren ausführen. Die linearen Gleichungen $s_{\alpha\beta}(x)-t_{\alpha\beta}(x)=0$ haben $k$ unabhängige Lösungen. Nun stelle man die quadratischen Gleichungen zwischen den Unbekannten $\chi_1,\chi_2,\ldots,\chi_n$ auf, die aus
\[
\fnum{6} \chi(x)\chi(y)=\chi(e)\chi(xy)
\]
erhalten werden, indem man für $x$ der Reihe nach jene $k$ Lösungen, für $y$ die $n$ Grundzahlen setzt. Sie liefern in Verbindung mit den linearen Gleichungen $r_{\alpha\beta}(\chi)=r_{\beta\alpha}(\chi)$ für die Verhältnisse $\chi_1,\chi_2,\ldots,\chi_n$ $k$ verschiedene Wertsysteme. Den konstanten Faktor wähle man jedesmal so, daß $\sum q_{\alpha\beta}\chi_\alpha\chi_\beta=1$ und $\sum\chi_\alpha e_\alpha$ positiv wird. Dann wird es gleich einer positiven ganzen Zahl $r$, dem Grade der Primfunktion $\Phi$, die durch den charakteristischen Parameter $\chi$ bestimmt ist.

Oder man bezeichne $k$ unabhängige Lösungen der Gleichungen $s_{\alpha\beta}(x)=t_{\alpha\beta}(x)$ mit $g,g',\ldots$ oder auch mit $\eta_1,\eta_2,\ldots,\eta_k$. Diese bilden die Basis einer kommutativen \textsc{Dedekind}schen Gruppe der Ordnung $k$, es ist also $\eta_\beta\eta_\gamma=\sum_\alpha c_{\alpha\beta\gamma}\eta_\alpha$, wo $c_{\alpha\beta\gamma}=c_{\alpha\gamma\beta}$ ist. Ihre Determinante ist ein Produkt von $k$ unabhängigen linearen Funktionen, und gehe, wenn die Basis $\eps_1,\eps_2,\ldots,\eps_n$ wieder eingeführt wird, in $\prod\psi(x)$ über, wo $\psi(x)=\sum\psi_\alpha x_\alpha$ ist. Dann bestimme man, da $\psi_1,\psi_2,\ldots,\psi_n$ einen willkürlichen konstanten Faktor enthalten, die Verhältnisse der Unbekannten $\chi_1,\chi_2,\ldots,\chi_n$ aus den linearen Gleichungen
\[
 r_{\alpha\beta}(\chi)-r_{\beta\alpha}(\chi)=0,
 \qquad g(\chi)=g(\psi),
 \qquad g'(\chi)=g'(\psi),\ldots,
\]
und daraus die wirklichen Werte wie oben.

\section*{§ 16.}

Sind $E_1,E_2,\ldots,E_n$ $n$ Matrizen $m$ten Grades, die den Bedingungen
\[
\fnum{1} E_\beta E_\gamma=\sum_\alpha a_{\alpha\beta\gamma}E_\alpha
\]
genügen, so bilden sie eine Darstellung der Gruppe $(\eps)$. Sind diese $n$ Matrizen nicht linear unabhängig (vergl. dagegen \textsc{Mol}. S. 126, Bedingung 1), so stellen sie eine mit $(\eps)$ homomorphe Gruppe dar.

Ist dann
\[
\fnum{2} (x_{\kappa\lambda})=\sum x_\alpha E_\alpha,
\]
so ist, falls $x=yz$ ist,
\[
 (x_{\kappa\lambda})=(y_{\kappa\lambda})(z_{\kappa\lambda}).
\]
Daher nenne ich die Matrix (2.) eine zur Gruppe $(\eps)$ gehörige Matrix. Ist ihre Determinante von Null verschieden, so ist sie nach (1.) § 4
\[
\fnum{3} |x_{\kappa\lambda}|=\prod\Phi(x)^s
\]
ein Produkt von Primfaktoren der Gruppendeterminante $\Theta$.

Ist $C$ eine konstante Matrix $m$ten Grades, deren Determinante von Null verschieden ist, so bilden auch die $n$ Matrizen $C^{-1}E_1C,C^{-1}E_2C,\ldots,C^{-1}E_nC$ eine Darstellung von $(\eps)$, die der ersten äquivalent genannt wird. Kann man $C$ nicht so wählen, daß diese $n$ Matrizen die Gestalt
\[
 \begin{pmatrix}E_1'&0\\ E_1^{(0)}&E_1''\end{pmatrix},
 \quad
 \begin{pmatrix}E_2'&0\\ E_2^{(0)}&E_2''\end{pmatrix},
 \quad \ldots \quad
 \begin{pmatrix}E_n'&0\\ E_n^{(0)}&E_n''\end{pmatrix}
\]
annehmen, so nenne ich die Darstellung primitiv oder irreduzibel, ist dies aber möglich, so nenne ich sie imprimitiv oder reduzibel, sind auch $E_1^{(0)},\ldots,E_n^{(0)}=0$, zerfallend oder zerlegbar.

Die Methoden, die ich in meiner Arbeit \emph{Über die Darstellung der endlichen Gruppen durch lineare Substitutionen}, II., Sitzungsberichte 1899, entwickelt habe, lassen sich (mit der am Ende des § 3 angedeuteten Modifikation) unmittelbar auf beliebige \textsc{Dedekind}sche Gruppen übertragen.

Die $r^2$ linearen Funktionen von $x_1,\ldots,x_n$, die ich in § 11 mit $x_{\kappa\lambda}$ bezeichnet habe, bilden eine zur Gruppe $(\eps)$ gehörige Matrix $r$ten Grades, deren Determinante $|x_{\kappa\lambda}|=\Phi(x)$ ist. Die entsprechende Darstellung von $(\eps)$ bezeichne ich mit $[\Phi]$. Die so erhaltenen, den $k$ verschiedenen Primfaktoren der Gruppendeterminante $\Theta$ entsprechenden Darstellungen sind die sämtlichen primitiven Darstellungen der Gruppe $(\eps)$.

Eine mit der Darstellung (2.) äquivalente zerfällt dann in $s$ primitive Darstellungen $[\Phi]$, $s'$ Darstellungen $[\Phi']$, usw. Der Koeffizient von $u^{m-1}$ in der Determinante $|x_{\kappa\lambda}+ue_{\kappa\lambda}|$
\[
 \varphi(x)=\sum s\chi(x)
\]
heißt die Spur der Darstellung (2.). Damit zwei Darstellungen einer \textsc{Dedekind}schen Gruppe äquivalent sind, ist notwendig und hinreichend, daß ihre Spuren übereinstimmen.
% ===== END WORK 070 =====

\clearpage

% ===== BEGIN WORK 071: 071_Theorie_der_hyperkomplexen_Gr_en_II =====
\frobtitle{Theorie der hyperkomplexen Größen II}

In meiner Arbeit \emph{Theorie der hyperkomplexen Größen}, im folgenden mit \emph{H.} zitiert, behandle ich im Anschluß an meine Untersuchungen über die Determinanten der endlichen Gruppen vorzugsweise die Eigenschaften der \textsc{Dedekind}schen Gruppen. Wie Hr. \textsc{Molien} gezeigt hat (\emph{H.} § 11, II), ist jede Gruppe $(\eps)$ mit Haupteinheit einer \textsc{Dedekind}schen Gruppe $(\grS)$ homomorph, deren Determinante durch jeden Primfaktor der Determinante der ganzen Gruppe teilbar ist. Ist $(\eta)$ die größte invariante Untergruppe von $(\eps)$, die aus lauter \emph{Wurzeln der Null} besteht -- ich nenne sie das \emph{Radikal} von $(\eps)$ --, so ist $(\eps)$ mod. $(\eta)=(\grS)$, d. h. $(\grS)$ ist die Gruppe $(\eps)$, falls man darin je zwei Größen als gleich betrachtet, deren Differenz in $(\eta)$ enthalten ist.

Wenn von den $n$ unabhängigen Grundzahlen $\eps_1,\ldots,\eps_n$ einer Gruppe $(\eps)$ die ersten $m$ die Basis einer Gruppe $(\grS)$, die letzten $n-m$ die einer Gruppe $(\eta)$ bilden, so nenne ich $(\eps)=(\grS)+(\eta)$ die \emph{Summe} dieser beiden Untergruppen. Ist eine von ihnen $(\eta)$ eine invariante Untergruppe von $(\eps)$, so ist die andere $(\grS)$ eine mit $(\eps)$ homomorphe Gruppe, weil
\[
 \eps_1x_1+\cdots+\eps_mx_m+\eps_{m+1}x_{m+1}+\cdots+\eps_nx_n
 \equiv \eps_1x_1+\cdots+\eps_mx_m \pmod{\eta}
\]
ist. Wenn beide Untergruppen invariante sind, so verschwindet das Produkt aus jeder Größe von $(\grS)$ und jeder von $(\eta)$, und $(\eps)$ \emph{zerfällt} in die beiden Gruppen $(\grS)$ und $(\eta)$. (\emph{H.} § 9.)

Eine in $(\eps)$ enthaltene Gruppe $(\eta)$ heißt eine \emph{invariante Untergruppe} von $(\eps)$, wenn $xy$ und $yx$ Größen von $(\eta)$ sind, falls $y$ irgendeine Größe von $(\eta)$ und $x$ irgendeine Größe von $(\eps)$ ist. Ist $n-m$ die Ordnung von $(\eta)$, so ist dann $(\eps)$ (mod. $\eta$) einer Gruppe $(\grS)$ der Ordnung $m$ homomorph. Es braucht aber nicht immer eine Untergruppe $(\grS)$ von $(\eps)$ zu geben der Art, daß $(\eps)=(\grS)+(\eta)$ ist.

Hr. \textsc{Cartan} hat in seiner Arbeit \emph{Sur les groupes bilinéaires et les systèmes de nombres complexes}, Ann. de Toulouse, tome XII, 1898, (im folgenden mit \emph{C.} zitiert), dem Satze von \textsc{Molien} eine präzisere Fassung gegeben, die ich zur Vervollständigung meiner Darstellung aus den Ergebnissen meiner ersten Arbeit herleiten will:

\begin{quote}\itshape
Jede Gruppe mit Haupteinheit ist die Summe ihres Radikals und einer {\normalfont\textsc{Dedekind}}schen Gruppe, deren Determinante durch jeden Primfaktor der Determinante der ganzen Gruppe teilbar ist.
\end{quote}

Während aber die invariante Untergruppe $(\eta)$, die ich das Radikal von $(\eps)$ nenne, völlig bestimmt ist, kann die Untergruppe $(\grS)$ meist auf unendlich viele Arten gewählt werden.

Wenn eine Potenz einer Zahl verschwindet, so nenne ich sie eine \emph{Wurzel der Null} (Pseudo-Nul bei \textsc{Cartan}, Nr. 21, nilpotent bei \textsc{Peirce}). Eine Gruppe, die aus lauter Wurzeln der Null besteht, nenne ich eine \emph{Wurzelgruppe}.

Wenn eine Gruppe $(\eps)$ die Summe einer \textsc{Dedekind}schen Gruppe $(\grS)$ der Ordnung $m$ und einer Wurzelgruppe $(\eta)$ der Ordnung $n-m$ ist, die eine invariante Untergruppe von $(\eps)$ ist, so ist $(\eta)$ das Radikal von $(\eps)$. Denn da $(\eps)$ und $(\grS)$ homomorph sind, so ist (\emph{H.} § 9) die Determinante der Gruppe $(\eps)$ durch die der Gruppe $(\grS)$ teilbar. Der lineare Rang der Determinante einer \textsc{Dedekind}schen Gruppe $(\grS)$ ist ihrer Ordnung $m$ gleich. Folglich ist der lineare Rang der Determinante von $(\eps)$, $\bar m\geq m$. Das Radikal $(\bar\eta)$ von $(\eps)$ hat die Ordnung $n-\bar m$. Da $(\eta)$ in $(\bar\eta)$ enthalten ist, so ist $n-m\leq n-\bar m$. Mithin ist $m=\bar m$ und $(\eta)=(\bar\eta)$.

Ehe ich zum Beweise des \textsc{Cartan}schen Satzes (§ 5) übergehe, will ich die Haupteigenschaften der Wurzelgruppen kurz herleiten.

\section*{§ 2.}

Wenn keine der beiden Determinanten $|S(x)|$ und $|T(x)|$ identisch verschwindet, so besitzt die Gruppe $(\eps)$ eine \emph{Haupteinheit} $e$, und für jede Größe $x$ in $(\eps)$ ist
\begin{equation}\tag{1.}
 ex=xe=x.
\end{equation}

Umgekehrt hat Hr. \textsc{Cartan} (\emph{C.} 15) gezeigt: Wenn für jede Größe $x$ $ex=x$ ist, so ist $S(e)=E$, also $|S(e)|=1$, und wenn $xe=x$ ist, so ist $T(e)=E$, also $|T(e)|=1$. Denn sucht man $x$ so zu bestimmen, daß $ex=0$ ist, so erhält man für die Koordinaten $x_1,\ldots,x_n$ $n$ homogene lineare Gleichungen mit der Matrix $S(e)$. Dann ist aber $x=ex=0$, also genügt diesen Gleichungen nur das Wertsystem $x_1=\cdots=x_n=0$, und folglich ist ihre Determinante $|S(e)|$ von Null verschieden. Weil nun $S(e)^2=S(e^2)=S(e)$ ist, so ist $S(e)=E$.

Sind $|S(x)|=\Theta$ und $|T(x)|=\Theta'$ die beiden antistrophen Determinanten der Gruppe $\eps$, so kann man zwei Größen $y$ und $z$, deren Koordinaten ganze Funktionen von $x_1,\ldots,x_n$ sind, so bestimmen, daß $xy=\Theta e$ und $zx=\Theta'e$ wird. Dann ist $T(y)T(x)=T(xy)=\Theta T(e)$, und mithin $|T(y)|\Theta'=\Theta^n$, und ebenso $|S(z)|\Theta=\Theta'{}^n$. Folglich ist jeder Primfaktor der einen der beiden Determinanten $\Theta$ und $\Theta'$ auch in der andern enthalten (\emph{H.} § 3).

Da ich vielfach auch Gruppen ohne Haupteinheit zu betrachten habe, so schicke ich darüber die folgenden Bemerkungen voraus: Bestehen zwischen $(n-1)^3$ Größen $a_{\alpha\beta\gamma}$ $(\alpha,\beta,\gamma=1,2,\ldots,n-1)$ die Relationen
\begin{equation}\tag{2.}
 \sum_{\kappa} a_{\kappa\alpha\beta}a_{\gamma\kappa\delta}
 =\sum_{\kappa} a_{\kappa\beta\delta}a_{\gamma\alpha\kappa},
\end{equation}
so gibt es immer eine Gruppe $(\bar\eps)$ mit $n-1$ linear unabhängigen Grundzahlen $\eps_1,\ldots,\eps_{n-1}$, zwischen denen die Beziehungen
\begin{equation}\tag{3.}
 \eps_{\beta}\eps_{\gamma}=\sum_{\alpha}a_{\alpha\beta\gamma}\eps_{\alpha}
\end{equation}
bestehen. Besitzt $(\bar\eps)$ eine Haupteinheit, so habe ich dies \emph{H.} § 2 bewiesen (vergl. auch \textsc{Study} in der Enzyklopädie). Im andern Falle setze man
\begin{equation}\tag{4.}
 a_{\alpha0\alpha}=a_{\alpha\alpha0}=1\qquad(\alpha=0,1,\ldots,n-1),
\end{equation}
sonst aber $a_{\alpha\beta\gamma}=0$, falls einer der Indizes $0$ ist. Dann gelten die Gleichungen (2.) auch für die $n^3$ Größen $a_{\alpha\beta\gamma}$ $(\alpha,\beta,\gamma=0,1,\ldots,n-1)$. Bezeichnet man ferner für ein bestimmtes $\lambda$ die Matrix $(a_{\alpha\lambda\beta})$ mit $E_{\lambda}$, so ist $E_0=E$ und
\begin{equation}\tag{5.}
 E_{\beta}E_{\gamma}=\sum_{\alpha}a_{\alpha\beta\gamma}E_{\alpha}.
\end{equation}
Wäre ferner $\sum c_{\lambda}E_{\lambda}=0$, so wäre $\sum c_{\lambda}a_{\alpha\lambda0}=0$, also $c_{\alpha}=0$. Demnach bilden $E_1,\ldots,E_{n-1}$ eine Darstellung der Gruppe $(\bar\eps)$ der Ordnung $n-1$ durch Matrizen des Grades $n$.

Diese Gruppe $(\bar\eps)$ ist eine invariante Untergruppe einer Gruppe $(\eps)$ mit der Basis $\eps_0,\eps_1,\ldots,\eps_{n-1}$, für die $\eps_0$ die Haupteinheit ist. In ihren beiden antistrophen Matrizen $S(x)$ und $T(x)$ ist $s_{00}(x)=t_{00}(x)=x_0$, aber $s_{0\beta}(x)=t_{0\beta}(x)=0$ (und $s_{\alpha0}(x)=t_{\alpha0}(x)=x_{\alpha}$). Läßt man die erste Zeile und Spalte weg, und setzt in den so erhaltenen Matrizen $(n-1)^{\ten}$ Grades $x_0=0$, so erhält man die beiden antistrophen Matrizen $\bar S(x)$ und $\bar T(x)$ von $(\bar\eps)$. Da von ihren Determinanten mindestens eine identisch verschwindet, so ist eine der beiden Determinanten $|S(x)|$ oder $|T(x)|$ durch $x_0^2$ teilbar.

Sei umgekehrt $(\eps)$ eine Gruppe der Ordnung $n$ mit der Haupteinheit $e$, die eine invariante Untergruppe $(\bar\eps)$ der Ordnung $n-1$ enthält. Sei $\eps_1,\ldots,\eps_{n-1}$ eine Basis von $(\bar\eps)$, $\eps_0$ eine in $(\bar\eps)$ nicht enthaltene Größe von $(\eps)$. Dann sind $|S(x)|$ und $|T(x)|$ beide durch $x_0$ teilbar. Je nachdem keine dieser beiden Determinanten oder mindestens eine von beiden durch $x_0^2$ teilbar ist, besitzt $(\bar\eps)$ eine Haupteinheit $(\bar e)$ oder nicht. Wäre $e$ in $(\bar\eps)$ enthalten, so wäre es nach der Definition einer invarianten Untergruppe auch $e\eps_0=\eps_0$. Daher kann man für $\eps_0$ immer die Haupteinheit $e$ wählen.

\section*{§ 3.}

Sind $A$ und $B$ zwei Matrizen $n^{\ten}$ Grades, so sind die charakteristischen Funktionen von $AB$ und $BA$ einander gleich. Denn ist $|B|$ von Null verschieden, so ist $B^{-1}(BA)B=AB$ der Matrix $BA$ ähnlich. Sind also die Elemente von $A$ und $B$ variable Größen, so gilt die Gleichung $|uE-AB|=|uE-BA|$ für alle Werte dieser Veränderlichen, wofür $|B|$ nicht Null ist, und folglich gilt sie identisch.

Sind $y$ und $z$ zwei Größen der Gruppe $(\eps)$, so sind die Koordinaten ihres Produktes $x=yz$
\begin{equation}\tag{1.}
 x_{\alpha}=\sum_{\beta,\gamma}a_{\alpha\beta\gamma}y_{\beta}z_{\gamma}
 =\sum_{\gamma}s_{\alpha\gamma}(y)z_{\gamma}.
\end{equation}
Ist nun $S(y)=0$, genügen also $y_1,\ldots,y_n$ den linearen Gleichungen $s_{\alpha\gamma}(y)=0$, so ist $yz=0$, falls $z$ eine beliebige Größe von $(\eps)$ ist. Besitzt also $(\eps)$ eine Haupteinheit, so ist $y=0$, in jedem Falle aber ist $y^2=0$.

Ist allgemeiner $x$ eine Größe von $(\eps)$, für die $S(x)^m=S(x^m)=0$ ist, so ist $x^mz=0$ und $x^{m+1}=0$. Damit also $x$ eine Wurzel der Null sei, ist notwendig und hinreichend, daß die Matrix $S(x)$ (oder $T(x)$) eine Wurzel der Null ist, daß mithin die charakteristischen Wurzeln von $S(x)$ alle verschwinden. Ist $zy$ eine Wurzel der Null, so ist auch $yz$ eine solche. Denn die charakteristischen Wurzeln der Matrix $S(zy)=S(z)S(y)$ sind alle Null, und folglich auch die der Matrix $S(y)S(z)=S(yz)$.

In einer Gruppe $(\eps)$ heißt eine Größe $y$ eine \emph{Wurzelgröße}, wenn $yz$ stets eine Wurzel der Null ist, falls $z$ eine beliebige Größe von $(\eps)$ ist. Sind $x$ und $z$ irgend zwei Größen von $(\eps)$, so ist dann auch $y(zx)=(yz)x$ eine Wurzel der Null, demnach auch $x(yz)$. Ist also $y$ eine Wurzelgröße von $(\eps)$, so sind es auch $xy$, $yz$ und $xyz$.

Dann verschwinden die charakteristischen Wurzeln der Matrix $S(yz)$, und folglich auch ihre Summe $\sigma(yz)$. Sei umgekehrt $\sigma(yz)=0$ für jede Größe $z$ von $(\eps)$. Ersetzt man $z$ durch $z(yz)^{\kappa-1}$, so erhält man $\sigma((yz)^\kappa)=0$. Demnach verschwindet (\emph{H.} § 4, (5.)) die Summe der $\kappa^{\ten}$ Potenzen der charakteristischen Wurzeln der Matrix $S(yz)$, also auch diese Wurzeln selbst, und folglich ist $yz$ eine Wurzel der Null. Die Koordinaten $y_1,\ldots,y_n$ der Wurzelgrößen $y$ werden mithin gefunden durch Auflösung der homogenen linearen Gleichungen, die man erhält, indem man die Ableitungen der bilinearen Form $\sigma(yz)$ nach $z_1,\ldots,z_n$ (oder die der quadratischen Form $\sigma(y^2)$ nach $y_1,\ldots,y_n$) gleich Null setzt. Daher reproduzieren sich die Wurzelgrößen von $(\eps)$ durch Addition und durch Multiplikation mit gewöhnlichen Größen. Ferner sind $yz$ und $zy$ Wurzelgrößen, wenn $y$ eine solche ist. Folglich bilden die Wurzelgrößen von $(\eps)$ eine invariante Untergruppe, die ich das Radikal von $(\eps)$ nenne.

\begin{quote}\itshape
I. Das Radikal einer Gruppe $(\eps)$ wird von allen Größen $y$ gebildet, deren Koordinaten für jede Größe $z$ den linearen Gleichungen $\sigma(yz)=0$ (oder $\tau(yz)=0$) genügen.
\end{quote}

Daß auch bei Gruppen ohne Haupteinheit (\emph{H.} § 4, (9.))
\begin{equation}\tag{2.}
 \sigma(xy)=\sigma(yx)
\end{equation}
ist, ergibt sich aus (2.), § 2, indem man $\gamma=\delta$ setzt und nach $\gamma$ summiert (\textsc{Mol}. § 3, (4.)).

Eine Wurzelgruppe kann auch als eine Gruppe definiert werden, die ihrem Radikal gleich ist. Nennt man eine von Null verschiedene Größe $x$, die der Gleichung $x^2=x$ genügt, eine \emph{Einheit}, so hat \textsc{Peirce}, \emph{Linear Associative Algebra}, Nr. 51 (American Journ. tome IV) gezeigt:

\begin{quote}\itshape
II. Damit eine Gruppe eine Wurzelgruppe sei, ist notwendig und hinreichend, daß sie keine Einheit enthält.
\end{quote}

Denn ist $m$ der kleinste Exponent, für den $x^m=0$ ist, so kann nicht $x=x^2$ sein, weil sonst $x^{m-1}=x^m=0$ wäre.

Sei ferner $\psi((x))=0$ die \emph{reduzierte} Gleichung, der die Größe $x$ der Gruppe $(\eps)$ genügt. Dies ist die Gleichung niedrigsten Grades der Form $px^0+qx+rx^2+\cdots=0$ oder $px+qx^2+rx^3+\cdots=0$, je nachdem $(\eps)$ eine Haupteinheit $x^0=e$ hat oder nicht. Sei $\psi(u)=(u-a)^\alpha\chi(u)$ und $\chi(a)$ von Null verschieden. Man bestimme (\emph{Über vertauschbare Matrizen}, Sitzungsberichte 1896, § 3) eine ganze Funktion $f(u)$ so, daß $f(u)-1$ durch $(u-a)^\alpha$ teilbar ist. Besitzt $(\eps)$ eine Haupteinheit, so sei ferner $f(u)$ durch $\chi(u)$ teilbar, ist dies aber nicht der Fall und ist $a$ von Null verschieden, so sei $f(u)$ durch $u\chi(u)$ teilbar. Dann ist $f(u)^2-f(u)$ durch $\psi(u)$ teilbar, $f(u)$ selbst aber nicht. Ist also $y=f((x))$, so ist $y^2=y$ und $y$ von Null verschieden. Ist demnach $(\eps)$ keine Wurzelgruppe, so enthält sie eine Einheit.

Hat $(\eps)$ eine Haupteinheit, und ist
\[
 \psi(u)=(u-a)^\alpha(u-b)^\beta(u-c)^\gamma\cdots,
\]
so mögen den verschiedenen Wurzeln $a,b,c,\ldots$ von $\psi(u)$ in der angegebenen Weise die ganzen Funktionen $f(u),g(u),h(u),\ldots$ entsprechen. Dann ist $f(u)g(u)$ und $f(u)+g(u)+h(u)+\cdots-1$ durch $\psi(u)$ teilbar, und mithin ist
\begin{equation}\tag{3.}
 f((x))^2=f((x)),\qquad f((x))g((x))=0
\end{equation}
und
\begin{equation}\tag{4.}
 f((x))+g((x))+h((x))+\cdots=e.
\end{equation}

Auf diese Art läßt sich (\emph{C.} 17) die Haupteinheit $e$ in so viele unabhängige Einheiten zerlegen, als die Gleichung $\psi(u)=0$ verschiedene Wurzeln hat.

\begin{quote}\itshape
III. Wenn die Primfaktoren der Determinante einer Gruppe mit Haupteinheit alle linear sind, so besteht ihr Radikal aus allen in der Gruppe enthaltenen Wurzeln der Null.
\end{quote}

Denn sind $x$ und $y$ zwei Größen der Gruppe $(\eps)$, so ist
\[
 |S(ux+vy+we)|=\prod (uu_\alpha+vv_\alpha+w)
\]
ein Produkt von linearen Faktoren. Setzt man $v=0$ oder $u=0$, so erkennt man, daß hierin $u_1,\ldots,u_n$ die charakteristischen Wurzeln von $S(x)$ sind und $v_1,\ldots,v_n$ die von $S(y)$. Ist $y$ eine Wurzel der Null, so ist $v_1=\cdots=v_n=0$. Ist $z=g((y))$, so sind daher die charakteristischen Wurzeln von $S(z)$ alle gleich $g(0)$. Ist also $v$ von Null verschieden und $z=u(y+ve)^{-1}$, so sind sie alle gleich $\frac{u}{v}$. Folglich ist
\[
 |S(x+u(y+ve)^{-1})|=\prod\left(u_\alpha+\frac{u}{v}\right),\qquad |S(y+ve)|=v^n,
\]
also weil
\[
 S(x+u(y+ve)^{-1})S(y+ve)=S(xy+vx+ue)
\]
ist
\[
 |S(xy+vx+ue)|=\prod(u_\alpha v+u).
\]
Da beide Seiten dieser Gleichung ganze Funktionen von $v$ sind, so gilt sie auch für den bisher ausgeschlossenen Wert $v=0$. Demnach ist $|S(xy+ue)|=u^n$, folglich ist (\emph{C.} 28) $xy$ eine Wurzel der Null und $y$ eine Wurzelgröße von $(\eps)$.

Die Umkehrung dieses Satzes ergibt sich aus den in § 5 erhaltenen Resultaten, wonach, wenn dort $r>1$ ist, $\eps_{12}^2=0$ ist, aber $\eps_{12}$ dem Radikale von $(\eps)$ nicht angehört.

\section*{§ 4.}

\begin{quote}\itshape
I. In jeder Wurzelgruppe gibt es eine von Null verschiedene Größe $x$, die den Gleichungen $xy=yx=0$ genügt, falls $y$ eine beliebige Größe der Gruppe ist.
\end{quote}

Für diesen Satz gibt Hr. \textsc{Cartan} (Nr. 31--34) einen sehr scharfsinnigen, aber etwas umständlichen Beweis, den ich hier durch einen wesentlich einfacheren ersetzen will. Man kann den Satz auch so aussprechen:

\begin{quote}\itshape
II. Jede Wurzelgruppe enthält eine invariante Untergruppe der Ordnung 1.
\end{quote}

Daß beide Sätze identisch sind, folgt aus einem von Hrn. \textsc{Cartan} (\emph{C.} 29) viel benutzten Lemma, worin $x$ und $y$ zwei Größen einer Gruppe mit oder ohne Haupteinheit bedeuten:

\begin{quote}\itshape
III. Ist $xy=ax$ (oder $yx=ax$), wo $a$ eine gewöhnliche Größe ist, und ist $y$ eine Wurzel der Null, so ist entweder $x=0$ oder $a=0$.
\end{quote}

Nach Voraussetzung gibt es eine solche Zahl $k$, daß $y^k=0$ ist, also auch eine solche Zahl $l$, daß $xy^l=0$ ist. Sei nicht $x=0$, und sei $m$ die kleinste Zahl dieser Art. Dann ist $m>0$, und $xy^m=0$, aber $xy^{m-1}$ (worunter für $m=1$ $x$ zu verstehen ist) von Null verschieden. Aus $xy^m=axy^{m-1}$ folgt daher $a=0$. Zu demselben Resultate gelangt man, indem man die Gleichung $x(y-a)=0$ rechts mit $y^{m-1}+y^{m-2}a+\cdots+a^{m-1}$ multipliziert. Der Satz läßt sich so verallgemeinern:

\begin{quote}\itshape
IV. Sind $x_1,x_2,\ldots,x_m$ Größen einer Wurzelgruppe, und ist das Produkt $x_1x_2\cdots x_m$ von Null verschieden, so sind die $m$ Größen
\[
 x_1,
 \quad x_1x_2,
 \quad x_1x_2x_3,
 \ldots
 \quad x_1x_2\cdots x_m
\]
linear unabhängig.
\end{quote}

Denn sei
\[
 ax_1\cdots x_l+bx_1\cdots x_lx_{l+1}+cx_1\cdots x_lx_{l+1}x_{l+2}+\cdots=0,
\]
wo $a$ der erste von Null verschiedene Koeffizient ist. Dann ist
\[
 bx_{l+1}+cx_{l+1}x_{l+2}+\cdots=-y
\]
eine Größe der Wurzelgruppe $(\eta)$, also eine Wurzel der Null, und es ist
\[
 x_1\cdots x_l(a-y)=0,
\]
also $x_1\cdots x_l=0$ und mithin auch $x_1\cdots x_m=0$.

\begin{quote}\itshape
V. In einer Wurzelgruppe der Ordnung $n-1$ verschwindet das Produkt von je $n$ Größen.
\end{quote}

Denn wäre $x_1\cdots x_n$ von Null verschieden, so wären die $n$ Größen
\[
 x_1,
 \quad x_1x_2,
 \quad x_1x_2x_3,
 \ldots
 \quad x_1x_2\cdots x_n
\]
linear unabhängig, also die Ordnung der Gruppe $\geq n$.

Nach diesem Satze gibt es für eine Wurzelgruppe $(\eta)$ eine solche invariante Zahl $m$, daß das Produkt von je $m$ Größen von $(\eta)$ verschwindet, aber nicht das von je $m-1$. Ist $x$ ein nicht verschwindendes Produkt von $m-1$ Größen von $(\eta)$, so ist $xy$ und $yx$ ein Produkt von $m$ Größen, und folglich Null. Durch wiederholte Anwendung dieses Satzes ergibt sich (\emph{C.} 31):

\begin{quote}\itshape
VI. Die Grundzahlen $\eta_1,\eta_2,\ldots,\eta_{n-1}$ einer Wurzelgruppe der Ordnung $n-1$ kann man so wählen, daß das Produkt $\eta_\alpha\eta_\beta$ eine lineare Verbindung der Grundzahlen ist, deren Index $>\alpha$ und $>\beta$ ist.
\end{quote}

Nach dem Satze V gibt es ferner für eine Wurzelgruppe $(\eta)$ eine invariante Zahl $l$ $(\leq m)$, die kleinste der Art, daß die $l^{\ten}$ Potenz jeder Größe der Gruppe verschwindet. Für eine kommutative Gruppe ist $l=m$. Ferner gilt der Satz:

\begin{quote}\itshape
VII. Gibt es in einer Wurzelgruppe, deren Ordnung $n-1$ ist, $n-1$ Größen, deren Produkt $x_1x_2\cdots x_{n-1}$ von Null verschieden ist, so ist auch $x_1^{n-1}$ von Null verschieden, und die Gruppe ist die aus allen ganzen Funktionen von $x_1$ gebildete kommutative Gruppe.
\end{quote}

Denn nach Satz IV sind die $n-1$ Größen
\[
 x_1,
 \quad x_1x_2,
 \quad x_1x_2x_3,
 \ldots
 \quad x_1x_2\cdots x_{n-1}
\]
linear unabhängig, bilden also eine Basis der Gruppe $(\eta)$. Daher ist
\[
 x_\kappa=a_{\kappa1}x_1+a_{\kappa2}x_1x_2+\cdots+a_{\kappa,n-1}x_1x_2\cdots x_{n-1}\quad(\kappa=1,2,\ldots,n-1).
\]
Da das Produkt von je $n$ Größen der Gruppe $(\eta)$ verschwindet, so erhält man durch Multiplikation dieser $n-1$ Gleichungen
\[
 x_1x_2\cdots x_{n-1}=a_{11}a_{21}\cdots a_{n-1,1}x_1^{n-1},
\]
und mithin ist $x_1^{n-1}$ von Null verschieden. Daher bilden auch, wenn man $x_1=x$ setzt,
\[
 x,
\quad x^2,
\ldots,
\quad x^{n-1}
\]
eine Basis von $(\eta)$. Ergänzt man sie durch $x^0=e$ zu einer Gruppe $(\eps)$ der Ordnung $n$, die aus den Größen $z=zx^0+z_1x+\cdots+z_{n-1}x^{n-1}$ besteht, so enthält deren Determinante $|S(z)|=z^n$ nur einen Elementarteiler.

\section*{§ 5.}

Ich wende mich jetzt zum Beweise des Satzes von \textsc{Cartan}. Sei $(\eps)$ eine Gruppe der Ordnung $n$ mit Haupteinheit, sei $(\eta)$ ihr Radikal, $n-m$ seine Ordnung. Dann ist $m$ der lineare Rang der Determinante von $(\eps)$
\begin{equation}\tag{1.}
 \Theta(x)=|S(x)|=\prod \Phi^{s}.
\end{equation}
Jeder ihrer $k$ Primfaktoren kann durch passende Wahl der Grundzahlen auf die Form einer Determinante
\begin{equation}\tag{2.}
 \Phi(x)=|x_{\alpha\beta}|\qquad(\alpha,\beta=1,2,\ldots,r)
\end{equation}
gebracht werden. Die
\begin{equation}\tag{3.}
 m=r^2+r'{}^2+r''{}^2+\cdots
\end{equation}
Elemente
\begin{equation}\tag{4.}
 x_{\alpha\beta},\qquad x'_{\alpha\beta},\qquad x''_{\alpha\beta},\ldots
\end{equation}
der $k$ verschiedenen Determinanten $\Phi,\Phi',\Phi'',\ldots$ der Grade $r,r',r'',\ldots$ sind lauter unabhängige Variable. Die Determinante der mit $(\eps)$ homomorphen \textsc{Dedekind}schen Gruppe $(\grS)$ ist $\prod \Phi^r$, der Exponent $r$ von $\Phi$ ist dem Grade des Primfaktors gleich.

Man kann die $n$ Grundzahlen von $(\eps)$ so wählen, daß $\eps_{m+1},\ldots,\eps_n$ die Basis des Radikals $(\eta)$ bilden, und daß die Größen (4.) die Koordinaten von $\eps_1,\ldots,\eps_m$ werden. Diese $m$ Grundzahlen bezeichne ich daher mit
\begin{equation}\tag{5.}
 \eps_{\alpha\beta},\qquad \eps'_{\alpha\beta},\qquad \eps''_{\alpha\beta},\ldots,
\end{equation}
jene $n-m$ mit $\eta_1,\ldots,\eta_{n-m}$, so daß
\begin{equation}\tag{6.}
 x=\sum x_{\alpha\beta}\eps_{\alpha\beta}
 +\sum x'_{\alpha\beta}\eps'_{\alpha\beta}
 +\sum x''_{\alpha\beta}\eps''_{\alpha\beta}
 +\cdots+\sum y_\rho\eta_\rho
\end{equation}
eine beliebige Größe von $(\eps)$ wird. Die charakteristische Determinante von $S(x)$ ist dann
\begin{equation}\tag{7.}
 |S(ue-x)|=\prod |u e_{\alpha\beta}-x_{\alpha\beta}|^{s}.
\end{equation}
Zwischen den Grundzahlen (5.) bestehen die Relationen
\begin{equation}\tag{8.}
 \eps_{\alpha\beta}\eps_{\beta\gamma}\equiv\eps_{\alpha\gamma},
 \qquad
 \eps_{\alpha\delta}\eps_{\beta\gamma}\equiv0,
 \qquad
 \eps_{\alpha\delta}\eps'_{\beta\gamma}\equiv0
 \pmod{\eta},
\end{equation}
in der zweiten sind $\beta$ und $\delta$ als verschieden vorausgesetzt. Es ist zu zeigen, daß sich die Grundzahlen (5.) mod. $(\eta)$ so abändern lassen, daß diese Gleichungen absolut gelten. Sei
\[
 z=a_1\eps_{11}+\cdots+a_r\eps_{rr}+a_{r+1}\eps'_{11}+\cdots+a_{r+r'}\eps'_{r'r'}+\cdots,
\]
oder wenn man
\begin{equation}\tag{9.}
 \eps_{11}=\eps_1,
 \ldots,
 \eps_{rr}=\eps_r,
 \qquad
 \eps'_{11}=\eps_{r+1},
 \ldots,
 \eps'_{r'r'}=\eps_{r+r'},\ldots
\end{equation}
setzt, $z=\sum \eps_\lambda a_\lambda$. Die
\begin{equation}\tag{10.}
 p=r+r'+r''+\cdots
\end{equation}
Koordinaten $a_\lambda$ können in irgend einer bestimmten Weise gewählt werden, nur sollen sie alle untereinander verschieden sein. Dann ist
\[
 \varphi(u)=|S(ue-z)|=(u-a_1)^s\cdots(u-a_r)^s(u-a_{r+1})^{s'}\cdots(u-a_{r+r'})^{s'}\cdots,
\]
und wenn $f(u)$ eine ganze Funktion von $u$ ist, nach (8.)
\[
 f((z))\equiv \sum \eps_\lambda f(a_\lambda)\pmod{\eta}.
\]
Wie in § 3, (3.) bestimme ich nun $p$ ganze Funktionen $f_\lambda(u)$ so, daß $(u-a_\lambda)^s f_\lambda(u)$ durch $\varphi(u)$ und $f_\lambda(u)-1$ durch $(u-a_\lambda)^s$ teilbar ist. Dann ist
\begin{equation}\tag{11$^{*}$.}
 f_\lambda((z))^2=f_\lambda((z)),
 \qquad
 f_\kappa((z))f_\lambda((z))=0
\end{equation}
und
\begin{equation}\tag{12$^{*}$.}
 \sum f_\lambda((z))=e.
\end{equation}
Ferner ist $f_\lambda((z))\equiv \eps_\lambda$ (mod. $\eta$), und daher kann man $\eps_\lambda$ (mod. $\eta$) so abändern, daß $f_\lambda(z)=\eps_\lambda$ wird. Dann ist
\begin{equation}\tag{11.}
 \eps_{\alpha\alpha}^2=\eps_{\alpha\alpha},
 \qquad
 \eps_{\alpha\alpha}\eps_{\beta\beta}=0,
 \qquad
 \eps_{\alpha\alpha}\eps'_{\beta\beta}=0
\end{equation}
und
\begin{equation}\tag{12.}
 \sum \eps_\lambda=e.
\end{equation}

Nun ist nach (8.)
\[
 \eps_{\alpha\beta}\equiv\eps_{\alpha\alpha}\eps_{\alpha\beta}\eps_{\beta\beta},
 \qquad
 \eps'_{\alpha\beta}\equiv\eps'_{\alpha\alpha}\eps'_{\alpha\beta}\eps'_{\beta\beta},\ldots
 \pmod{\eta}.
\]
Ist $\alpha=\beta$, so gelten die Gleichungen absolut. Ist $\alpha$ von $\beta$ verschieden, so ersetze man $\eps_{\alpha\beta}$ durch $\eps_{\alpha\alpha}\eps_{\alpha\beta}\eps_{\beta\beta}$. Dadurch wird diese Grundzahl (mod. $\eta$) so abgeändert, daß nach (11.) jetzt
\begin{equation}\tag{13.}
 \eps_{\alpha\alpha}\eps_{\alpha\beta}=\eps_{\alpha\beta},
 \qquad
 \eps_{\alpha\beta}\eps_{\beta\beta}=\eps_{\alpha\beta}
\end{equation}
und
\begin{equation}\tag{14.}
 \eps_{\alpha\beta}\eps_{\delta\gamma}=0,
 \qquad
 \eps_{\alpha\beta}\eps'_{\delta\gamma}=0
\end{equation}
ist, falls in der ersten Gleichung $\delta$ von $\beta$ verschieden ist.

Die Grundzahlen $\eps_{\alpha\alpha},\eps'_{\alpha\alpha},\ldots$ ändere ich nicht weiter ab, wohl aber die Grundzahlen $\eps_{\alpha\beta},\eps'_{\alpha\beta},\ldots$, doch so, daß die bereits gewonnenen Relationen bestehen bleiben. Ist $\eta$ eine beliebige Größe von $(\eps)$, so nenne ich
\begin{equation}\tag{15.}
 \eps_\kappa\eta\eps_\lambda=\eta_{\kappa\lambda}
\end{equation}
eine Größe vom Typus $(\kappa,\lambda)$. Sie genügt den Gleichungen
\begin{equation}\tag{16.}
 \eps_\kappa\eta_{\kappa\lambda}=\eta_{\kappa\lambda},
 \qquad
 \eta_{\kappa\lambda}\eps_\lambda=\eta_{\kappa\lambda},
\end{equation}
und wenn $\mu$ von $\kappa$, und $\nu$ von $\lambda$ verschieden ist,
\begin{equation}\tag{17.}
 \eps_\mu\eta_{\kappa\lambda}=0,
 \qquad
 \eta_{\kappa\lambda}\eps_\nu=0.
\end{equation}
Ist $\eta$ eine Größe der invarianten Untergruppe $(\eta)$, so gehört auch $\eta_{\kappa\lambda}$ dem Radikal an. Dann kann man $\eps_{\alpha\beta}$ um eine beliebige Wurzelgröße $\eta_{\alpha\beta}$ vom Typus $(\alpha,\beta)$ abändern, ohne daß die Gleichungen (13.) und (14.) sich ändern.

Nun ist $\eps_{\alpha\beta}\eps_{\beta\alpha}=\eps_{\alpha\alpha}$, also
\begin{equation}\tag{18.}
 \eps_{\alpha\beta}\eps_{\beta\alpha}=\eps_{\alpha\alpha}-\eta_{\alpha\alpha},
\end{equation}
wo $\eta_{\alpha\alpha}$ eine Größe von $(\eta)$ ist, und zwar nach (13.) eine solche vom Typus $(\alpha,\alpha)$. Daher ist die rechte Seite gleich $\eps_{\alpha\alpha}(e-\eta_{\alpha\alpha})$. Aus der Gleichung
\[
 e-x^l=(e-x)(x^0+x+\cdots+x^{l-1})
\]
folgt
\[
 e=(e-x)(x^0+x+\cdots+x^{l-1}),
\]
falls $x^l=0$ ist. Ist nun $\eta_{\alpha\alpha}^l=0$, so multipliziere man die Gleichung (18.) rechts mit
\[
 \eta_{\alpha\alpha}^{0}+\eta_{\alpha\alpha}+\cdots+\eta_{\alpha\alpha}^{l-1}.
\]
Dann erhält man
\[
 \eps_{\alpha\beta}(\eps_{\beta\alpha}+\eta_{\beta\alpha})=\eps_{\alpha\alpha},
\]
wo
\[
 \eta_{\beta\alpha}=\eps_{\beta\alpha}(\eta_{\alpha\alpha}+\eta_{\alpha\alpha}^2+\cdots+\eta_{\alpha\alpha}^{l-1})
\]
eine Wurzelgröße vom Typus $(\beta,\alpha)$ ist. Ändert man also $\eps_{\beta\alpha}$ um $\eta_{\beta\alpha}$ ab, so wird
\begin{equation}\tag{19.}
 \eps_{\alpha\beta}\eps_{\beta\alpha}=\eps_{\alpha\alpha}.
\end{equation}
Dann ist aber auch immer (\emph{C.} 52.)
\begin{equation}\tag{20.}
 \eps_{\beta\alpha}\eps_{\alpha\beta}=\eps_{\beta\beta}.
\end{equation}
Denn zunächst ist $\eps_{\beta\alpha}\eps_{\alpha\beta}=\eps_{\beta\beta}-\eta_{\beta\beta}$, und daraus ergibt sich wie eben $\eps_{\beta\alpha}(\eps_{\alpha\beta}+\eta_{\alpha\beta})=\eps_{\beta\beta}$, also wenn man links mit $\eps_{\alpha\beta}$ multipliziert, $\eps_{\alpha\beta}+\eta_{\alpha\beta}=\eps_{\alpha\beta}$ und mithin $\eta_{\alpha\beta}=0$.

Nun wähle man $\eps_{12},\ldots,\eps_{1r}$ irgendwie (\emph{C.} 67) den bisher aufgestellten Relationen gemäß, alsdann $\eps_{21},\ldots,\eps_{r1}$ so, daß $\eps_{1\alpha}\eps_{\alpha1}=\eps_{11}$, und folglich auch $\eps_{\alpha1}\eps_{1\alpha}=\eps_{\alpha\alpha}$ ist. Dann gilt die Gleichung
\begin{equation}\tag{21.}
 \eps_{\alpha\beta}=\eps_{\alpha1}\eps_{1\beta},
\end{equation}
falls $\alpha=1$ oder $\beta=1$ oder $\alpha=\beta$ ist, also für die Fälle, wo bereits über $\eps_{\alpha\beta}$ verfügt ist. Sind aber $\alpha$ und $\beta$ verschieden und beide $>1$, so ist $\eps_{\alpha\beta}=\eps_{\alpha1}\eps_{1\beta}+\eta_{\alpha\beta}$, wo $\eta_{\alpha\beta}$ eine Wurzelgröße vom Typus $(\alpha,\beta)$ ist. Daher kann man $\eps_{\alpha\beta}$ so abändern, daß die Gleichung (21.) erfüllt wird. Dann ist
\[
 \eps_{\alpha\beta}\eps_{\beta\gamma}
 =\eps_{\alpha1}(\eps_{1\beta}\eps_{\beta1})\eps_{1\gamma}
 =\eps_{\alpha1}\eps_{11}\eps_{1\gamma}
 =\eps_{\alpha1}\eps_{1\gamma}
 =\eps_{\alpha\gamma}.
\]
Die Gleichungen (8.) gelten jetzt alle absolut, und mithin sind die $m$ Größen (5.) die Grundzahlen einer Gruppe $(\grS)$.

Nach (12.) ist
\[
 \eta=(\eps_1+\cdots+\eps_p)\eta(\eps_1+\cdots+\eps_p)=\sum(\eps_\kappa\eta\eps_\lambda).
\]
Daher kann man die Wurzelgrößen alle aus den $(n-m)p^2$ Größen $\eps_\kappa\eta_i\eps_\lambda$ zusammensetzen. Wählt man aus diesen ein System unabhängiger aus, so gehört auch jede der Grundzahlen des Radikals zu einem bestimmten Typus $(\kappa,\lambda)$. In der Untergruppe $(\grS)$ hat $\eps_{\alpha\beta}$, falls $\alpha$ und $\beta$ zwei der Zahlen von 1 bis $r$ sind, den Typus $(\alpha,\beta)$. Dagegen enthält sie z. B. keine Größe vom Typus $(1,r+1)$.

Werden die $n$ Grundzahlen von $(\eps)$ in der angegebenen Weise gewählt, so lassen sich nach (17.) alle Größen $\xi_{\kappa\lambda}$ vom Typus $(\kappa,\lambda)$ nur durch die Grundzahlen ausdrücken, die denselben Typus haben. Die Grundzahlen vom Typus $(1,1)$ bilden die Basis einer Gruppe, wofür $\eps_1=\eps_{11}$ die Haupteinheit und zugleich die einzige Einheit ist. Die übrigen Grundzahlen vom Typus $(1,1)$ bilden die Basis ihres Radikals. Sind $\alpha$ und $\beta$ zwei der Zahlen $1,2,\ldots,r$, so ist von den beiden Gleichungen
\[
 \eps_{\alpha\beta}\xi_{\beta\lambda}=\xi_{\alpha\lambda},
 \qquad
 \eps_{\beta\alpha}\xi_{\alpha\lambda}=\xi_{\beta\lambda}
\]
jede die Folge der anderen. Daher erhält man (\emph{C.} 57) alle Größen vom Typus $(\beta,\lambda)$, indem man $\eps_{\beta\alpha}$ mit allen Größen vom Typus $(\alpha,\lambda)$ multipliziert.

Ferner bilden alle Grundzahlen vom Typus
\[
\begin{gathered}
 (1,1),(1,r+1),(1,r+r'+1),\ldots\ (r+1,1)(r+1,r+r),(r+1,r+r'+1),\ldots \\
 (r+r'+1,1),(r+r'+1,r+1),(r+r'+1,r+r'+1),\ldots
\end{gathered}
\]
zusammengenommen die Basis einer Gruppe $(\eps')$, die $k$ unter ihnen enthaltenen Größen von $(\grS)$, $\eps_1,\eps_{r+1},\eps_{r+r'+1},\ldots$, bilden die Basis einer kommutativen \textsc{Dedekind}schen Gruppe $(\grS')$, die übrigen die Basis einer Wurzelgruppe $(\eta')$, die eine invariante Untergruppe von $(\eps')$ ist. Nach § 1 ist daher $(\eta')$ das Radikal von $(\eps')$. Da $(\grS')$ eine kommutative Gruppe ist, so zerfällt die Determinante von $(\eps')$ in lauter lineare Faktoren. Wie sich aus den Relationen zwischen den Grundzahlen dieser Gruppe $(\eps')$ die zwischen den Grundzahlen von $(\eps)$ ableiten lassen, hat Hr. \textsc{Cartan} (\emph{C.} 57--60) so ausführlich dargelegt, daß ich darauf hier nicht zurückkommen will.
% ===== END WORK 071 =====

\end{document}
